\documentclass[11pt]{memoir}
\setcounter{secnumdepth}{4}
\setcounter{tocdepth}{3}

% ==========================================
% FONTS: EB Garamond (matching Volumes I–II)
% ==========================================

\usepackage[T1]{fontenc}
\usepackage[utf8]{inputenc}
\usepackage{lmodern}
\frenchspacing

\usepackage[
  cmintegrals,
  cmbraces,
  noamssymbols
]{newtxmath}
\usepackage{ebgaramond}

\usepackage[
  activate={true,nocompatibility},
  final,
  tracking=true,
  kerning=true,
  spacing=true,
  factor=1100,
  stretch=10,
  shrink=10
]{microtype}

\SetExtraKerning[unit=space]
  {encoding={*}, family={*}, series={*}, size={footnotesize,small,normalsize}}
  {\textemdash={400,400},
   "28={,150},
   "29={150,},
   \textquotedblleft={,150},
   \textquotedblright={150,}}

% ==========================================
% PACKAGES
% ==========================================

\usepackage{mleftright}
\mleftright
\let\Bbbk\undefined
\let\openbox\undefined
\usepackage{amsmath,amssymb,amsthm}
\usepackage{mathtools}
\usepackage{tikz-cd}
\usetikzlibrary{positioning,arrows.meta,calc,decorations.markings}
\makeatletter
\def\hyper@nopatch@caption{}
\makeatother
\usepackage[unicode,pdfencoding=auto,psdextra]{hyperref}
\usepackage{geometry}
\geometry{top=1.2in, bottom=1.2in, left=1.25in, right=1.25in, footskip=0.5in}
\setlength{\headheight}{28pt}
\setlength{\emergencystretch}{3em}
\setlength{\hfuzz}{1pt}
\hbadness=3000
\raggedbottom

\setlength{\parskip}{0pt plus 0.5pt}
\setlength{\abovedisplayskip}{8pt plus 2pt minus 4pt}
\setlength{\belowdisplayskip}{8pt plus 2pt minus 4pt}
\setlength{\abovedisplayshortskip}{4pt plus 2pt minus 2pt}
\setlength{\belowdisplayshortskip}{4pt plus 2pt minus 2pt}

\usepackage{tcolorbox}
\usepackage{longtable}
\usepackage{booktabs}
\usepackage{multicol}
\usepackage{enumitem}
\usepackage{newunicodechar}
\usepackage{thmtools}

% Unicode fallbacks
\newunicodechar{∞}{$\infty$}
\newunicodechar{≃}{$\simeq$}
\newunicodechar{≅}{$\cong$}
\newunicodechar{⊗}{$\otimes$}
\newunicodechar{⊕}{$\oplus$}

% ==========================================
% THEOREM ENVIRONMENTS (Garamond style)
% ==========================================

\declaretheoremstyle[
  spaceabove=\topsep,
  spacebelow=\topsep,
  headfont=\normalfont\scshape,
  notefont=\normalfont\itshape,
  bodyfont=\normalfont,
  postheadspace=0.5em,
  headpunct={.}
]{garamondthm}

\declaretheoremstyle[
  spaceabove=\topsep,
  spacebelow=\topsep,
  headfont=\normalfont\itshape,
  notefont=\normalfont\itshape,
  bodyfont=\normalfont,
  postheadspace=0.5em,
  headpunct={.}
]{garamonddef}

\declaretheoremstyle[
  spaceabove=\topsep,
  spacebelow=\topsep,
  headfont=\normalfont\itshape,
  bodyfont=\normalfont\small,
  postheadspace=0.5em,
  headpunct={.}
]{garamondrem}

\declaretheorem[style=garamondthm, numberwithin=chapter, name=Theorem]{theorem}
\declaretheorem[style=garamondthm, sibling=theorem, name=Proposition]{proposition}
\declaretheorem[style=garamondthm, sibling=theorem, name=Lemma]{lemma}
\declaretheorem[style=garamondthm, sibling=theorem, name=Corollary]{corollary}
\declaretheorem[style=garamondthm, sibling=theorem, name=Conjecture]{conjecture}
\declaretheorem[style=garamonddef, sibling=theorem, name=Definition]{definition}
\declaretheorem[style=garamonddef, sibling=theorem, name=Example]{example}
\declaretheorem[style=garamonddef, sibling=theorem, name=Construction]{construction}
\declaretheorem[style=garamondrem, sibling=theorem, name=Remark]{remark}
\declaretheorem[style=garamondrem, sibling=theorem, name=Warning]{warning}

% Unnumbered variants
\declaretheorem[style=garamondthm, numbered=no, name=Theorem]{theorem*}
\declaretheorem[style=garamondrem, numbered=no, name=Remark]{remark*}

% ==========================================
% CLAIM STATUS TAGS (matching Volumes I–II)
% ==========================================

\newcommand{\ClaimStatusProvedHere}{\textsuperscript{\textsc{[ph]}}}
\newcommand{\ClaimStatusProvedElsewhere}{\textsuperscript{\textsc{[pe]}}}
\newcommand{\ClaimStatusConjectured}{\textsuperscript{\textsc{[cj]}}}
\newcommand{\ClaimStatusHeuristic}{\textsuperscript{\textsc{[he]}}}
\newcommand{\ClaimStatusOpen}{\textsuperscript{\textsc{[op]}}}

% ==========================================
% MACROS
% ==========================================

% Categories
\newcommand{\Cat}{\mathbf{Cat}}
\newcommand{\DGCat}{\mathbf{dgCat}}
\newcommand{\Ainf}{A_\infty}
\newcommand{\Linf}{L_\infty}
\newcommand{\Einf}{E_\infty}
\newcommand{\Mod}{\mathrm{Mod}}
\newcommand{\Rep}{\mathrm{Rep}}
\newcommand{\Perf}{\mathrm{Perf}}
\newcommand{\Ind}{\mathrm{Ind}}
\newcommand{\Coh}{\mathrm{Coh}}
\newcommand{\Fuk}{\mathrm{Fuk}}
\newcommand{\MF}{\mathrm{MF}}
\newcommand{\Vect}{\mathrm{Vect}}

% Calabi-Yau
\newcommand{\CY}{\mathrm{CY}}
\newcommand{\HH}{\mathrm{HH}}
\newcommand{\HC}{\mathrm{HC}}
\newcommand{\HN}{\mathrm{HN}}
\newcommand{\HP}{\mathrm{HP}}
\newcommand{\Tr}{\mathrm{Tr}}
\newcommand{\Sd}{\mathbb{S}^d}

% Chiral / Factorization
\newcommand{\Ran}{\mathrm{Ran}}
\newcommand{\Fact}{\mathrm{Fact}}
\newcommand{\barB}{\overline{B}}
\newcommand{\Omegach}{\Omega^{\mathrm{ch}}}
\newcommand{\ChirHoch}{\mathrm{ChirHoch}}

% Quantum groups
\newcommand{\Uq}{U_q}
\newcommand{\Uhbar}{U_\hbar}
\newcommand{\Yg}{Y(\mathfrak{g})}
\newcommand{\RTT}{\mathrm{RTT}}

% Operads
\newcommand{\Eone}{E_1}
\newcommand{\Etwo}{E_2}
\newcommand{\En}{E_n}
\newcommand{\SC}{\mathrm{SC}}
\newcommand{\FM}{\overline{\mathrm{FM}}}
\newcommand{\Ger}{\mathrm{Ger}}
\newcommand{\BV}{\mathrm{BV}}

% Algebra
\newcommand{\Sym}{\mathrm{Sym}}
\newcommand{\End}{\mathrm{End}}
\newcommand{\Hom}{\mathrm{Hom}}
\newcommand{\RHom}{\mathrm{RHom}}
\newcommand{\Ext}{\mathrm{Ext}}
\newcommand{\Tor}{\mathrm{Tor}}
\newcommand{\id}{\mathrm{id}}
\newcommand{\op}{\mathrm{op}}
\newcommand{\colim}{\mathrm{colim}}

% Common symbols
\newcommand{\cA}{\mathcal{A}}
\newcommand{\cB}{\mathcal{B}}
\newcommand{\cC}{\mathcal{C}}
\newcommand{\cD}{\mathcal{D}}
\newcommand{\cE}{\mathcal{E}}
\newcommand{\cF}{\mathcal{F}}
\newcommand{\cM}{\mathcal{M}}
\newcommand{\cO}{\mathcal{O}}
\newcommand{\cZ}{\mathcal{Z}}
\newcommand{\frakg}{\mathfrak{g}}
\newcommand{\hbar}{\hslash}

% ==========================================
% DOCUMENT
% ==========================================

\begin{document}

% ------------------------------------------
% Title
% ------------------------------------------

\begin{titlingpage}
\vspace*{4cm}
\begin{center}

{\HUGE\scshape Calabi--Yau Quantum Groups}

\vspace{1.5cm}

{\Large\itshape Chiral Algebras from Calabi--Yau Categories\\
via $E_1$/$E_2$ Factorization}

\vspace{3cm}

{\large Raeez Lorgat}

\vspace{2cm}

{\normalsize Draft --- \today}

\end{center}
\end{titlingpage}

% ------------------------------------------
% Front matter
% ------------------------------------------

\frontmatter
\tableofcontents

% ------------------------------------------
% Main matter
% ------------------------------------------

\mainmatter

% ==========================================
% PART I: Calabi--Yau Categories and Cyclic Structures
% ==========================================

\part{Calabi--Yau Categories and Cyclic Structures}
\label{part:cy-categories}

\chapter{Introduction}
\label{ch:introduction}

\section{The question}
\label{sec:the-question}

What is a Calabi--Yau quantum chiral algebra?

A Calabi--Yau category $\cC$ of dimension $d$ carries a canonical cyclic $\Ainf$-structure: a non-degenerate trace
\[
  \Tr \colon \HH_\bullet(\cC) \longrightarrow k[-d]
\]
on its Hochschild homology, encoding the CY condition as a $d$-dimensional Frobenius structure at the chain level.  The cyclic bar complex $\mathrm{CC}_\bullet(\cC)$ is the primary invariant; it carries an $\Sd$-action (from the $d$-sphere framing of the trace), and its $\Sd$-equivariant structure governs higher-genus amplitudes.

A chiral algebra $A$ on a curve $X$ carries a bar complex $B(A)$, which is a factorization coalgebra on $\Ran(X)$.  The bar differential encodes holomorphic OPE residues; the coproduct encodes topological interval-cutting.  Together they assemble into a Swiss-cheese algebra structure (Volume~II).  At genus $g \geq 1$, the bar complex acquires curvature $\kappa(A) \cdot \omega_g$ from the Hodge bundle, and the full modular structure is controlled by the universal Maurer--Cartan element $\Theta_A$ (Volume~I).

The programme of this work is to construct a precise bridge:
\begin{center}
\begin{tikzcd}[column sep=huge]
  \left\{\begin{array}{c}\text{CY categories with}\\\text{appropriate $E_n$-structure}\end{array}\right\}
  \arrow[r, "\Phi"]
  &
  \left\{\begin{array}{c}\text{quantum chiral algebras}\\\text{with modular trace}\end{array}\right\}
\end{tikzcd}
\end{center}
The functor $\Phi$ should:
\begin{enumerate}[label=(\roman*)]
  \item take a CY category $\cC$ (Fukaya, derived, matrix factorization, or more general) as input;
  \item extract the $E_2$-monoidal structure on $\cC$ (braided tensor product, quantum group action);
  \item produce a chiral algebra $A_\cC$ whose bar complex $B(A_\cC)$ encodes the CY cyclic homology;
  \item realize the CY trace as the modular characteristic $\kappa(A_\cC)$.
\end{enumerate}

\section{The $E_1$/$E_2$ chiral hierarchy}
\label{sec:e1-e2-hierarchy}

The key structural ingredient is the $E_n$-chiral hierarchy:

\begin{itemize}
  \item $E_1$-chiral algebras (Volume~II, Part~VII): associative factorization on $\mathbb{C} \times \mathbb{R}$, encoding the ordered/topological sector.  The representation category of an $E_1$-chiral algebra is monoidal.
  \item $E_2$-chiral algebras (this work): braided factorization on $\mathbb{C} \times \mathbb{C}$, encoding the full holomorphic-holomorphic sector.  The representation category is \emph{braided} monoidal --- the natural habitat of quantum groups.
  \item The passage $E_1 \to E_2$ via Dunn additivity: $E_2 \simeq E_1 \otimes E_1$.  An $E_2$-chiral algebra is an $E_1$-algebra in $E_1$-algebras.
\end{itemize}

The CY condition enters through Kontsevich's identification: a $d$-dimensional CY structure on $\cC$ determines an $\mathbb{S}^d$-framing of the Hochschild complex, hence an $\mathbb{S}^1$-action on $\HH_\bullet(\cC)$.  For $d = 2$ (CY surfaces), this $\mathbb{S}^1$-action is exactly the data of an $E_2$-algebra structure on the cyclic homology --- the braiding.

\section{Relation to Volumes I and II}
\label{sec:relation-vols}

This work depends on and extends the two preceding volumes:

\begin{center}
\renewcommand{\arraystretch}{1.4}
\begin{tabular}{lll}
  \toprule
  \textbf{Volume} & \textbf{Provides} & \textbf{Used here} \\
  \midrule
  I (Modular Koszul Duality) & Bar-cobar machine, $\Theta_A$, $\kappa(A)$ & CY bar complex, modular trace \\
  II (3D HT QFT) & $\SC^{\mathrm{ch,top}}$, PVA descent, DK bridge & $E_1$ sector, braided structure \\
  III (this work) & CY $\to$ chiral functor, $E_2$ theory & --- \\
  \bottomrule
\end{tabular}
\end{center}

\section{CY3 combinatorics as generalized root data}
\label{sec:cy3-root-data}

The deep structural insight animating this work is that the combinatorics of a Calabi--Yau threefold $X$ --- its lattice, intersection form, BPS spectrum --- constitute the \emph{generalized root datum} of what we call a \emph{quantum vertex chiral group} $G(X)$.

Given a CY3 $X$, the generalized root datum $\mathcal{R}(X)$ consists of:
\begin{enumerate}[label=(\roman*)]
  \item A lattice $\Lambda(X)$ extracted from $K(X)$ or $H^*(X)$ with its Euler form;
  \item Real simple roots $\Delta^{\mathrm{re}}$ from the geometry (hyperbolic sublattice for $K3 \times E$; toric diagram for toric CY3);
  \item Imaginary roots $\Delta^{\mathrm{im}} = \Delta^{\mathrm{im}}_0 \sqcup \Delta^{\mathrm{im}}_1$ (even and odd) with multiplicities $\mathrm{mult}(\alpha)$ from BPS counting --- Donaldson--Thomas invariants for $X$, or equivalently the Fourier coefficients of the appropriate automorphic form;
  \item A Weyl group $W(X)$ acting on the positive cone of $\Lambda(X)$;
  \item A Weyl vector $\rho \in \Lambda(X) \otimes \mathbb{Q}$.
\end{enumerate}

Two families of examples illuminate the construction:

\emph{The $K3 \times E$ tower.}  For $X = (S \times E)/(\mathbb{Z}/N\mathbb{Z})$ with $S$ a K3 surface and $E$ an elliptic curve, the lattice is $\Lambda^{3,2} \simeq \Lambda^{1,1} \oplus \Lambda^{1,1} \oplus [2]$ of signature $(3,2)$.  The hyperbolic sublattice $\Lambda^{2,1}_{II}$ with Gram matrix $\bigl(\begin{smallmatrix} 2 & -2 & -2 \\ -2 & 2 & -2 \\ -2 & -2 & 2 \end{smallmatrix}\bigr)$ provides the real roots.  The root multiplicities are the Fourier coefficients $f(n,l)$ of the weak Jacobi form $\phi_{0,1}$ --- the K3 elliptic genus.  The resulting generalized Borcherds--Kac--Moody superalgebra $\mathfrak{g}_{\Delta_5}$ has the Igusa cusp form $\Delta_5$ as its denominator identity.

\emph{Toric CY3 and critical CoHAs.}  For a toric CY3 determined by a fan $\Sigma$, the toric diagram determines a quiver $Q_X$.  The critical cohomological Hall algebra $\mathrm{CoHA}(Q_X)$ is the positive half of an affine super Yangian $Y(\widehat{\mathfrak{g}}_{Q_X})$.  For $\mathbb{C}^3$: the Jordan quiver gives $Y(\widehat{\mathfrak{gl}}_1) \simeq \mathcal{W}_{1+\infty}$.  The toric fan IS the Dynkin data.

\section{Automorphic correction as shadow tower}
\label{sec:automorphic-shadow}

The passage from the ``naive'' Kac--Moody algebra $\mathfrak{g}$ (built from the real root Gram matrix alone) to the full generalized BKM superalgebra $\mathfrak{g}_X$ (with all imaginary roots and their multiplicities) is the \emph{automorphic correction}.  In the language of Volume~I, this passage is precisely the \emph{shadow Postnikov tower}:
\begin{center}
\renewcommand{\arraystretch}{1.4}
\begin{tabular}{lll}
  \toprule
  \textbf{Shadow tower (Vol~I)} & \textbf{BKM language} & \textbf{CY3 geometry} \\
  \midrule
  $\kappa$ (arity 2) & Weyl vector $\rho$ contribution & Leading DT invariant \\
  $C$ (arity 3, cubic) & First imaginary root corrections & Genus-0 3-point function \\
  $Q$ (arity 4, quartic) & Higher imaginary roots & Genus-0 4-point function \\
  Full $\Theta_A$ & Complete automorphic correction & Full BPS generating function \\
  \bottomrule
\end{tabular}
\end{center}
The denominator identity of $\mathfrak{g}_X$ --- which IS the automorphic form ($\Delta_5$ for $K3 \times E$, the DT partition function for toric CY3) --- is the bar-complex Euler product of the quantum vertex chiral group.  The Weyl--Kac--Borcherds character formula is the factorization structure of $B(A_X)$.

\section{Main results}
\label{sec:main-results}

The main theorems of this work are:

\begin{itemize}
  \item \textbf{Theorem CY-A} (CY-to-chiral functor): Construction of $\Phi \colon \CY_d\text{-}\Cat \to E_2\text{-}\mathrm{ChirAlg}$ from CY categories of dimension $d$ to $E_2$-chiral algebras, via the factorization envelope and $\mathbb{S}^d$-framing.  For $d = 3$, $\Phi$ produces a quantum vertex chiral group $G(X)$ whose generalized root datum is $\mathcal{R}(X)$.
  \item \textbf{Theorem CY-B} ($E_2$-chiral Koszul duality): Bar-cobar adjunction in the $E_2$-chiral setting.  The automorphic correction of the BKM superalgebra $\mathfrak{g}_X$ is identified with the shadow tower $\Theta_{A_X}$; the denominator identity equals the bar-complex Euler product.
  \item \textbf{Theorem CY-C} (Quantum group realization): For toric CY3, $\Rep^{E_2}(G(X))$ is braided monoidal equivalent to $\Rep(Y(\widehat{\mathfrak{g}}_{Q_X}))$; for $K3 \times E$, the Sp$_4(\mathbb{Z})$-module structure on the denominator identity recovers the braided structure.  The critical CoHA is the $E_1$-sector (positive half).
  \item \textbf{Theorem CY-D} (Modular CY characteristic): $\kappa(G(X))$ equals the weight of the automorphic form --- the CY Euler characteristic.  The genus-$g$ obstruction $\mathrm{obs}_g = \kappa \cdot \lambda_g$ recovers the genus-$g$ DT/GW invariant on the uniform-weight lane.
\end{itemize}

\section{Guide for the reader}
\label{sec:guide}

Part~I establishes the categorical foundations: CY categories, cyclic $\Ainf$-structures, and the Hochschild calculus.  Part~II develops the $E_1$ and $E_2$ chiral theories.  Part~III constructs the bridge functor $\Phi$ and proves Theorems CY-A through CY-D, with the identification of automorphic correction and shadow tower as the central result.  Part~IV studies quantum groups and their braided factorization structure, including the Drinfeld center as the bulk algebra.  Part~V works through the standard landscape: Fukaya categories, derived categories, matrix factorizations (ADE singularities and W-algebras), and quantum group representations (Kazhdan--Lusztig, affine Yangians, critical CoHAs).  Part~VI connects back to Volumes~I and~II and to the geometric Langlands programme.

\chapter{Calabi--Yau Categories}
\label{ch:cy-categories}

A Calabi--Yau category is a dg category whose Serre functor is a pure degree shift. This single condition, due independently to Kontsevich and to Costello, organises the diverse sources of Calabi--Yau geometry (coherent sheaves, Fukaya categories, matrix factorisations, noncommutative resolutions) into a uniform categorical framework. It is the framework on which the Vol~III correspondence programme $\{\Phi_d\}_{d \geq 1}$ acts: each constructed $\Phi_d$ factors as $\Phi_d = \SpCh_{\Sigma_{d-1},C} \circ \PhiFA_d$ (Theorem~\ref{thm:phi-two-stage-factorisation}; Definition~\ref{def:phi-fa-and-sp}), with stage~$1$ $\PhiFA_d$ taking a cyclic $\Ainf$ Calabi--Yau category of dimension $d$ to an $E_d$-holomorphic factorisation algebra on the CY$_d$ variety, and stage~$2$ $\SpCh_{\Sigma_{d-1},C}$ specialising along a $(d-1)$-cycle and reference curve to produce an $\Etwo$-chiral algebra at $d = 2$ or an $\Eone$-chiral algebra on the witnessed $d \geq 3$ loci; the native operadic level $n_{\mathrm{native}}(d)$ on $C$ (Proposition~\ref{prop:native-en-level}) is part of the dimension-specific data. This chapter fixes the categorical input that enters stage~$1$, recording the definitions, the Hochschild structure, the cyclic $\Ainf$ enhancement, and the interface to $\{\Phi_d\}$ (Chapter~\ref{ch:cy-to-chiral}).

\begin{remark}[Chapter scope and cross-volume anchors]
\label{rem:cy-cat-cross-volume-anchors}
The categorical input fixed by this chapter feeds the correspondence programme $\{\Phi_d\}$ of Theorem~\textup{\ref{thm:phi-platonic}}, characterised by the four universal properties \textup{(U1)}--\textup{(U4)}. The downstream evaluations of $\{\Phi_d\}$ supplied by the cyclic $\Ainf$ data of this chapter are the dimension-stratified spectrum of $\kappa_{\mathrm{ch}}$ values catalogued in Theorem~\textup{\ref{thm:kappa-stratification-by-d}} (Chapter~\textup{\ref{ch:cy-d-kappa-stratification}}) and the Borcherds-weight invariant $\kappa_{\mathrm{BKM}} = c_N(0)/2$ of Proposition~\textup{\ref{prop:bkm-weight-universal}} (Chapter~\textup{\ref{ch:k3-times-e}}). On the Vol~I side, the categorical Hochschild data of Section~\textup{\ref{sec:cy-hochschild}} is the cohomological shadow of the chiral universal Maurer--Cartan element of Vol~I Chapter~\textup{\texttt{chapters/theory/chiral\_climax\_platonic.tex} (Vol~I)}, and the Koszul conductor identity $K = -c_{\mathrm{ghost}}(\BRST)$ of Vol~I Chapter~\textup{\texttt{chapters/theory/universal\_conductor\_K\_platonic.tex} (Vol~I)} fixes the family-dependent constants entering the kappa-spectrum.
\end{remark}

\begin{remark}[Chain-level and \texorpdfstring{$(\infty,1)$}{(inf,1)}-categorical: two load-bearing lenses]
\label{rem:cy-cat-chain-infinity-parity}
Both lenses are load-bearing throughout the chapter and neither
replaces the other. The chain-level lane supplies the explicit
cyclic $\Ainf$ enhancement (Theorem~\ref{thm:cyclic-ainf-enhancement}),
the explicit Hochschild complexes of Section~\ref{sec:cy-hochschild},
and the ambient-qualified Serre-duality computations driving the
$K3$, $K3\times E$, and local $\mathbb P^2$ examples. The
$(\infty,1)$-categorical lane supplies the definition of
$\CY_d\text{-}\Cat$ as an $\infty$-category, the Maulik--Okounkov /
Kontsevich--Soibelman stable-envelope functoriality, and the
conditional $E_3$ obstruction package used for
$\Phi_3^{(\Sigma_2,C)}$ on the H1--H4 locus
(Theorem~\ref{thm:cy-to-chiral-d3}). In that package
$\HH^{-2}_{E_1}$ vanishing is a theorem only after the comparison map
to the obstruction complex and the complete/exhaustive/separated
strongly convergent filtration with empty total-degree $-2$ line have
been supplied, or after the corrected TCFT carrier gives the same
comparison. Every theorem is stated in the
lane in which its proof works, with scope qualifiers
(ambient, framing, $\Ainf$-level) separating the two readings;
neither lane is ``the real theorem'' of which the other is a shadow.
\end{remark}

\begin{remark}[Convention: \texorpdfstring{$\kappa_\bullet$}{kappa-.} subscripts in Vol~III]
\label{rem:cy-cat-kappa-convention}
Throughout this chapter and the rest of Vol~III the symbol $\kappa_\bullet$ ranges over the closed indexing set $\bullet \in \{\mathrm{ch}, \mathrm{cat}, \mathrm{BKM}, \mathrm{fiber}\}$: every occurrence carries one of these four subscripts, and each subscript names a distinct invariant. Indexed against the two-stage factorisation $\Phi_d = \SpCh_{\Sigma_{d-1},C} \circ \PhiFA_d$ (Theorem~\ref{thm:phi-two-stage-factorisation}), the four entries split into three tiers:
\begin{itemize}
 \item \emph{Tier~(i), CY-datum intrinsic} (independent of stage~$1$ or stage~$2$): $\kappa_{\mathrm{cat}} = \chi(\cO_X)$, the holomorphic Euler characteristic of the CY$_d$ variety. Read off the Hodge diamond alone.
 \item \emph{Tier~(ii), stage~$1$ invariant} (determined by $\PhiFA_d(\cC)$ on $X$ before any specialisation): $\kappa_{\mathrm{fiber}} = \mathrm{rank}(\Lambda)$, the rank of the lattice of local / divisor classes on $X$ seen by the $E_d$-factorisation algebra.
 \item \emph{Tier~(iii), stage~$2$ specialisation-dependent} (determined by the stage-$2$ output $\SpCh_{\Sigma_{d-1},C}(\PhiFA_d(\cC))$ on the reference curve): $\kappa_{\mathrm{ch}}(\Phi_d(\cC))$ (chiral modular characteristic as the Hodge-supertrace reading on $C$) and the Borcherds-weight $\kappa_{\mathrm{BKM}}(\Phi_{N}(\cC))$ (weight of the BKM denominator automorphic form attached to a specialisation cycle).
\end{itemize}
The statement ``algebraizations share $\kappa_{\mathrm{cat}}$'' is vacuous: $\kappa_{\mathrm{cat}}$ is a tier~(i) invariant of $X$, not a property of $\cC$ or of a specialisation.
\end{remark}

\section{Smooth and proper dg categories}
\label{sec:smooth-proper-dg}

The Serre functor required by the CY condition exists only for categories
that are both compact over themselves and finitely presented in their
$\Hom$-complexes. These are the smooth and proper conditions.

\begin{definition}[Smooth dg category]
\label{def:smooth-dg-cat}
A dg category $\cC$ over $k$ is \emph{smooth} if the diagonal bimodule $\cC \in \cC^{\op} \otimes \cC\text{-}\Mod$ is a perfect bimodule (compact in the derived category of bimodules).
\end{definition}

\begin{definition}[Proper dg category]
\label{def:proper-dg-cat}
A dg category $\cC$ is \emph{proper} if for all objects $X, Y \in \cC$, the Hom-complex $\Hom_\cC(X,Y)$ has finite-dimensional total cohomology.
\end{definition}

A category that is both smooth and proper is called \emph{saturated}. Saturated categories admit a Serre functor $S_\cC$, a distinguished autoequivalence characterised by the natural isomorphism $\Hom_\cC(X, Y) \simeq \Hom_\cC(Y, S_\cC X)^\vee$. Classical Serre duality recovers $S_{D^b(\Coh(X))} = (- \otimes \omega_X)[\dim X]$; the Calabi--Yau condition asks that the twist $\omega_X$ be trivial, leaving only the shift.

\section{The Calabi--Yau condition}
\label{sec:cy-condition}

A saturated category carries a Serre functor $S_\cC = (-)\otimes\omega_\cC[d]$.
Two pieces can be trivial: the twist $\omega_\cC$ and the shift $[d]$. Asking
$\omega_\cC$ to be trivial leaves a pure shift and is the Calabi--Yau condition.

\begin{definition}[Calabi--Yau category, after Kontsevich]
\label{def:cy-category}
A smooth dg (or $\Ainf$) category $\cC$ is \emph{Calabi--Yau of dimension $d$} if the Serre functor is a pure shift: $S_\cC \simeq [d]$. Equivalently, there is a quasi-isomorphism of $\cC$-bimodules
\[
 \cC \xrightarrow{\;\sim\;} \cC^\vee[d]
\]
where $\cC^\vee = \RHom_{\cC^{\op} \otimes \cC}(\cC, \cC^{\op} \otimes \cC)$ is the inverse dualizing bimodule. Equivalently again, there exists a non-degenerate cyclic pairing
\[
 \langle \cdot, \cdot \rangle_d \colon \cC(a, b) \otimes \cC(b, a) \longrightarrow k[-d]
\]
of degree $-d$, satisfying the graded symmetry $\langle f, g \rangle_d = (-1)^{|f|\,|g|} \langle g, f \rangle_d$.
\end{definition}

The equivalence of the bimodule formulation with the cyclic pairing formulation is due to Costello~\cite{Costello2005}; see also Kontsevich's Mirror Symmetry address~\cite{Kontsevich1994} for the original categorical statement and Keller~\cite{Keller2006} for the dg-categorical reformulation used here.

\begin{remark}
\label{rem:cy-vs-smooth-cy}
The CY condition combines smoothness (the bimodule $\cC$ is compact) with the Serre functor being a shift. A smooth proper CY category is \emph{saturated} in the sense of Kontsevich--Soibelman. Non-proper ``relative'' CY categories also appear (matrix factorisations, Ginzburg dg algebras); they require the Hochschild invariants of Section~\ref{sec:cy-hochschild} to be interpreted in a compactly-supported or negative-cyclic refinement.
\end{remark}

\begin{example}[Fukaya category]
\label{ex:fukaya-cy}
For a compact symplectic manifold $(M, \omega)$ with $c_1(M) = 0$, the Fukaya category $\Fuk(M)$ is CY of dimension $\dim_{\mathbb{R}}(M)/2 = n$. The CY pairing arises from the cyclic structure on the $\Ainf$-algebra of Lagrangian Floer cochains (Fukaya; Seidel~\cite{Seidel2008}).
\end{example}

\begin{example}[Derived category of a CY manifold]
\label{ex:derived-cy}
For a smooth projective CY manifold $X$ of complex dimension $n$ (i.e.\ $\omega_X \simeq \cO_X$ and $H^i(X, \cO_X) = 0$ for $0 < i < n$), the bounded derived category $D^b(\Coh(X))$ is CY of dimension $n$. The CY pairing is Serre duality composed with the $\cO_X$-trivialisation.
\end{example}

\begin{example}[Matrix factorizations]
\label{ex:mf-cy}
For $W \colon \mathbb{A}^n \to \mathbb{A}^1$ a polynomial with isolated singularity, the category $\MF(W)$ of matrix factorizations is CY of dimension $n - 2$ (Dyckerhoff; Orlov). For the ADE singularities in two variables ($n = 2$), $\MF(W)$ is $\CY_0$ (semisimple, with the structure of a Frobenius algebra). For threefold singularities ($n = 4$, e.g.\ $W = x_1^2 + \cdots + x_4^2$), $\MF(W)$ is $\CY_2$. The exponent $n - 2$ (not $n - 1$) records the projectivisation of the $\mathbb{G}_m$-action on $W^{-1}(0)$.
\end{example}

\section{Hochschild cohomology and the CY structure}
\label{sec:cy-hochschild}

Hochschild cohomology is the universal invariant of a dg category. For a CY$_d$ category, Serre duality forces additional structure: an $E_2$-algebra structure on $\HH^\bullet(\cC)$ (Kontsevich--Soibelman, Tamarkin) and, in the presence of the CY pairing, a $(d-2)$-shifted Poisson / BV structure.

\begin{remark}[Three Hochschild theories (chiral / topological / categorical): NEVER conflate]
\label{rem:cy-cat-three-hochschild}
The geometry determines which Hochschild theory is at play. Three theories coexist in this programme; they are distinct functors on distinct sources, with distinct outputs:
\begin{enumerate}
 \item \emph{Categorical Hochschild} $\HH^\bullet_{\mathrm{cat}}$ (this chapter): input a dg category $\cC$; output $E_2$-algebra (Deligne) carrying a $(2{-}d)$-shifted Poisson / $\BV$ structure when $\cC$ is $\CY_d$. Source: dg categories. The Vol~III categorical kappa is $\kappa_{\mathrm{cat}}(\cC) = \chi(\cO_X)$ when $\cC = D^b(\Coh(X))$.
 \item \emph{Topological Hochschild} $\HH^\bullet_{\mathrm{top}}$: input an $E_1$-algebra; output $E_2$-algebra (Deligne; Lurie). Source: associative algebras over a base ring. This is the Hochschild theory operative in Vol~II for the slab algebra of 3d HT theories.
 \item \emph{Chiral Hochschild} $\ChirHoch^\bullet$: input an $E_\infty$-chiral algebra on a curve $X$; output concentrated in cohomological degrees $\{0,1,2\}$ on the curve-ambient reading (Vol~I Theorem~H; $d=1$ Ran ambient). Under $\Phi_d$ with $d \ge 2$, the concentration enlarges to $\{0,1,2,d\}$ in the cohomological lane; chain-level is infinite for class-M inputs at $d \ge 3$. Source: chiral algebras. Carries a chiral Gerstenhaber bracket (from OPE residues on $\mathrm{FM}_k(C)$) distinct from the topological Gerstenhaber bracket (from the little $2$-disks operad); the two agree only for $E_\infty$-inputs via formality.
\end{enumerate}
The CY-to-chiral correspondence programme $\{\Phi_d\}$ of Chapter~\textup{\ref{ch:cy-to-chiral}} maps the categorical Gerstenhaber bracket on $\HH^\bullet_{\mathrm{cat}}(\cC)$ to the chiral convolution bracket on $\ChirHoch^\bullet(\Phi_d(\cC))$; the two are distinct objects on distinct sources. We therefore qualify every Hochschild statement that crosses this bridge with one of $\{\mathrm{cat},\mathrm{top},\mathrm{chiral}\}$, naming the theory under consideration.
\end{remark}

\begin{definition}[Hochschild cohomology]
\label{def:hochschild-cohom}
The Hochschild cohomology of $\cC$ is
\[
 \HH^\bullet(\cC) \;=\; \RHom_{\cC^{\op} \otimes \cC}(\cC, \cC),
\]
the derived endomorphisms of the diagonal bimodule. It is a graded-commutative algebra under Yoneda product and carries the Gerstenhaber bracket of degree $-1$. Throughout this chapter, $\HH^\bullet$ without qualifier means $\HH^\bullet_{\mathrm{cat}}$ in the sense of Remark~\textup{\ref{rem:cy-cat-three-hochschild}}.
\end{definition}

\begin{theorem}[Deligne conjecture; Kontsevich--Soibelman, Tamarkin]
\label{thm:deligne-conjecture}
The pair (Yoneda product, Gerstenhaber bracket) on $\HH^\bullet(\cC)$ is the homology of an $E_2$-algebra structure fixed after choosing the Kontsevich--Tamarkin formality torsor point.
\end{theorem}

For a CY$_d$ category, the cyclic pairing relates $\HH^\bullet(\cC)$ to $\HH_\bullet(\cC)[d]$. Transporting structure through this duality produces the higher Calabi--Yau brackets.

\begin{proposition}[CY bracket spectrum]
\label{prop:cy-bracket-spectrum}
For a $\CY_d$ category $\cC$, the Hochschild invariants carry the following additional structure:
\begin{itemize}
 \item $d = 1$: a BV operator $\Delta \colon \HH^\bullet(\cC) \to \HH^{\bullet - 1}(\cC)$ whose deviation from being a derivation recovers the Gerstenhaber bracket;
 \item $d = 2$: a degree $-2$ Poisson bracket $\{-, -\}_2$ on $\HH^\bullet(\cC)$, part of the $2$-shifted symplectic structure on the moduli of objects (PTVV shift law: Pantev--To\"en--Vaqui\'e--Vezzosi~\cite{PTVV2013}, $d$-shifted symplectic on $\mathrm{Map}(X, \mathrm{Pt}/G)$ for CY$_d$ target);
 \item $d \geq 3$: a $(d - 2)$-shifted Poisson structure / $P_{d-1}$-algebra (PTVV, loc.\ cit.), refining to the Lie-cobar level in the cyclic setting.
\end{itemize}
\end{proposition}

\begin{theorem}[HKR for smooth varieties]
\label{thm:hkr-smooth}
For $X$ a smooth projective variety over $k$ of characteristic zero, the Hochschild--Kostant--Rosenberg isomorphism gives
\[
 \HH^p(D^b(\Coh(X))) \;\simeq\; \bigoplus_{q + r = p} H^q\big(X, \wedge^r T_X\big).
\]
When $X$ is CY$_n$, the trivialisation $\omega_X \simeq \cO_X$ identifies $\wedge^r T_X \simeq \Omega_X^{n - r}$, so
\[
 \HH^p(D^b(\Coh(X))) \;\simeq\; \bigoplus_{q + r = p} H^q\big(X, \Omega_X^{n - r}\big).
\]
The sign and compatibility with the Mukai pairing are recorded in C{\u a}ld{\u a}raru~\cite{Caldararu2005}.
\end{theorem}

\begin{example}[Hochschild cohomology of \texorpdfstring{$K3$}{K3}]
\label{ex:hh-k3}
For a $K3$ surface $X$, the Hodge diamond is
\[
 h^{0,0} = 1,\; h^{2,0} = h^{0,2} = 1,\; h^{1,1} = 20,\; h^{2,2} = 1,
\]
with all other entries zero, and $H^q(\wedge^r T_X) \simeq H^q(\Omega_X^{2-r})$ by the CY$_2$ trivialisation $\omega_X \simeq \cO_X$. Applied to the Kontsevich HKR sum
$\HH^p(D^b(\Coh(X))) = \bigoplus_{q + r = p} H^q(X, \wedge^r T_X)$ (Theorem~\ref{thm:hkr-smooth}), the single non-vanishing summands at each total degree are
\[
 \HH^0 \simeq H^0(\cO_X) \simeq k, \qquad
 \HH^1 \simeq 0,
\]
\[
 \HH^2 \simeq H^0(\wedge^2 T_X) \,\oplus\, H^1(T_X) \,\oplus\, H^2(\cO_X) \simeq k \oplus k^{20} \oplus k \simeq k^{22},
\]
\[
 \HH^3 \simeq 0, \qquad
 \HH^4 \simeq H^2(\wedge^2 T_X) \simeq k,
\]
with the $H^1(T_X) \simeq k^{20}$ summand recording the infinitesimal K3 moduli and $H^0(\wedge^2 T_X) \simeq H^2(\cO_X) \simeq k$ arising from the trivialisation of $\wedge^2 T_X \simeq \cO_X$. The total dimension is
\[
 \dim_k \HH^\bullet(D^b(\Coh(K3))) \;=\; 1 + 0 + 22 + 0 + 1 \;=\; 24,
\]
matching the topological Euler characteristic $\chi(K3) = 24$. Collapsing by cohomological row rather than total degree (the Mukai grading of C{\u a}ld{\u a}raru) instead splits this $24$ as $(1{+}1) + 20 + (1{+}1) = 2 + 20 + 2$ across the three even cohomological rows of $K3$; this is the grading used downstream in the Mukai-lattice picture, distinct from the Kontsevich HKR grading above. The Mukai pairing on $\HH^\bullet$ is the non-degenerate pairing $\langle -, - \rangle_\mathrm{Mukai}$ of Mukai~\cite{Mukai1987}; it coincides with the CY pairing produced by Definition~\ref{def:cy-category} at $d = 2$.
\end{example}

\begin{remark}[\texorpdfstring{$\kappa_\mathrm{cat}$}{kappa-cat} as Hodge supertrace; comparison with \texorpdfstring{$\kappa_{\mathrm{ch}}$}{kappa-ch}]
\label{rem:kappa-cat-from-chi}
For $X$ a compact CY$_d$, the categorical kappa of the Vol~III spectrum is the Hodge-supertrace formula
\[
 \kappa_{\mathrm{cat}}(X) \;=\; \chi(\cO_X) \;=\; \sum_{q=0}^{d} (-1)^q\, h^{0,q}(X),
\]
the holomorphic Euler characteristic. For K3: $\chi(\cO_{K3}) = 1 - 0 + 1 = 2$, so $\kappa_{\mathrm{cat}}(K3) = 2$. At every \emph{odd} dimension $d$, Serre duality forces $h^{0,q}(X) = h^{0,d-q}(X)$, and the supertrace pair-cancels, giving
\[
 \chi(\cO_X) \;=\; 0 \qquad (\text{compact CY}_d,\ d\ \text{odd; Vol~III \texttt{prop:chi-O-vanishes-odd-d})}.
\]
The chiral kappa $\kappa_{\mathrm{ch}}(\Phi_d(\cC))$ is the compact Hodge-supertrace branch unless a construction superscript is displayed; the equality $\kappa_{\mathrm{ch}} = \kappa_{\mathrm{cat}}$ holds at $d = 2$ with $h^{1,0}(X) = 0$ (Theorem~\textup{\ref{thm:cy-d-d2}}; Proposition~\textup{\ref{prop:cy-kappa-d2}}). At $d = 3$ the Stage-2 Heisenberg branch can be nonzero while the compact trace vanishes: $\kappa_{\mathrm{ch}}^{\mathrm{Heis}}(K3 \times E) = 3$, but $\kappa_{\mathrm{ch}}(A_{K3\times E})=\kappa_{\mathrm{cat}}(K3\times E)=0$ (Remark~\ref{rem:beauville-kappa-formula-subscript-split}). The dimension-stratified spectrum of $\kappa_{\mathrm{ch}}$ values is catalogued in Theorem~\textup{\ref{thm:kappa-stratification-by-d}} and prescribes the correct invariant in each $d$.

Numerical separation for K3: $\dim \HH^\bullet(D^b(\Coh(K3))) = 24$ (the topological Euler characteristic); $\kappa_{\mathrm{cat}}(K3) = \chi(\cO_{K3}) = 2$; $\kappa_{\mathrm{ch}}(\Phi(D^b(\Coh(K3)))) = 2$ ($d = 2$, $h^{1,0}=0$). For the $K3 \times E$ product: $\kappa_{\mathrm{cat}}(K3 \times E) = \chi(\cO_{K3}) \cdot \chi(\cO_E) = 2 \cdot 0 = 0$ by K\"unneth and $\kappa_{\mathrm{ch}}(A_{K3\times E})=0$ by compact Hodge supertrace, while $\kappa_{\mathrm{ch}}^{\mathrm{Heis}}(K3 \times E) = 2 + 1 = 3$ is the canonical Stage-2 Heisenberg--Cartan rank, Beauville-additive under tensor products (Vol~I Theorem~D; Remark~\ref{rem:beauville-kappa-formula-subscript-split}), and $\kappa_{\mathrm{BKM}}(K3 \times E) = c_1(0)/2 = 5$ by the Borcherds weight theorem (Proposition~\textup{\ref{prop:bkm-weight-universal}}, Vol~III); the Mukai lattice rank is $\kappa_{\mathrm{fiber}}(K3) = 24$. The four construction values $\{\kappa_{\mathrm{cat}}, \kappa_{\mathrm{ch}}^{\mathrm{Heis}}, \kappa_{\mathrm{BKM}}, \kappa_{\mathrm{fiber}}\} = \{0, 3, 5, 24\}$ for the K3 $\times$ E example are distinct invariants, not one invariant evaluated four times.
\end{remark}

\section{Cyclic \texorpdfstring{$\Ainf$}{A-inf} structure}
\label{sec:cyclic-ainf-cy}

The CY pairing in Definition~\ref{def:cy-category} is cohomological. The categorical input to $\Phi$ requires a chain-level enhancement: a cyclic $\Ainf$-structure for which every higher operation $m_n$ is compatible with the pairing.

\begin{definition}[Cyclic \texorpdfstring{$\Ainf$}{A-inf} category]
\label{def:cyclic-ainf-cat}
A \emph{cyclic $\Ainf$-structure of degree $-d$} on an $\Ainf$-category $\cC$ is a collection of pairings
\[
 \langle -, - \rangle_d \colon \cC(a_0, a_1) \otimes \cC(a_1, a_0) \longrightarrow k[-d]
\]
together with the cyclic symmetry condition
\[
 \langle m_n(f_1, \ldots, f_n), f_0 \rangle_d \;=\; (-1)^{\epsilon} \langle m_n(f_0, f_1, \ldots, f_{n-1}), f_n \rangle_d
\]
for all $n \geq 2$, where $\epsilon$ is the Koszul sign determined by the degrees of the $f_i$.
\end{definition}

\begin{definition}[Negative cyclic CY class]
\label{def:negcyc-cy}
A \emph{$d$-dimensional CY structure} on a smooth dg category $\cC$ is a class
\[
 [\sigma] \in \HC^-_{d}(\cC)
\]
in the negative cyclic homology, whose image under the canonical map $\HC^- \to \HH$ is a non-degenerate trace $\Tr \colon \HH_d(\cC) \to k$. The lift to $\HC^-$ is essential for the $\bS^d$-framing used in Chapter~\ref{ch:cy-to-chiral}.
\end{definition}

\begin{theorem}[Cyclic \texorpdfstring{$\Ainf$}{A-inf} enhancement]
\label{thm:cyclic-ainf-enhancement}
Every smooth proper CY$_d$ category admits a cyclic $\Ainf$-enhancement, unique up to $\Ainf$-quasi-isomorphism preserving the pairing. The proof is due to Kontsevich--Soibelman in the formal case, Costello~\cite{Costello2005} for the TCFT enhancement, and Seidel~\cite{Seidel2008} for the Fukaya case.
\end{theorem}

At $d = 2$ the enhancement is explicit enough to compute directly; it drives the $K3$ computations feeding the Borcherds denominator (Vol~III Chapter~\ref{ch:k3-times-e}). At $d = 3$ Sheridan~\cite{Sheridan2015} constructs the required cyclic enhancement for the quintic HMS model. The general theorem is conditional on the chain-level $\bS^3$-framing datum and on a fixed admissible specialisation $(\Sigma_2,C)$ (Theorem~\ref{thm:cy-to-chiral-d3}); the framing is automatic only on verified toric or formal loci.

\section{Interface with the CY-to-chiral correspondence programme \texorpdfstring{$\{\Phi_d\}$}{Phi-d}}
\label{sec:cy-cat-phi-interface}

The proved Vol~III functor is
\[
 \Phi \colon \CY_2\text{-}\Cat^{\mathrm{cyc}} \longrightarrow \Etwo\text{-}\mathrm{ChirAlg}.
\]
The $d = 3$ extension exists as an $\Eone$-chiral object-level
assignment on the witnessed locus, and the $\Eone$ output level is
rigid in the sense of Theorem~CY-A$_3$
($E_2$-lifting contractible on the shadow via
Theorem~\ref{thm:derived-framing-obstruction}), \emph{not} an
equivalence of categories with anything:
\[
 \Phi_3^{(\Sigma_2,C)}(\cC)
 :=
 \SpCh_{\Sigma_2,C}(\PhiFA_3(\cC))
 \in \Eone\text{-}\mathrm{ChirAlg}(C).
\]
This is an assignment on objects after H1--H4 and the admissible specialisation datum have been fixed, not a functor on all smooth proper CY$_3$ categories. Both readings use as input the data of this chapter: a smooth proper cyclic $\Ainf$ category with its negative cyclic CY class $[\sigma] \in \HC^-_d(\cC)$. Each constructed $\Phi_d$ factors through a canonical two-stage construction (Theorem~\ref{thm:phi-two-stage-factorisation}): Stage~1 assembles the native $E_d$-holomorphic factorisation algebra $\PhiFA_d(\cC)$ on the CY$_d$ variety from the Gerstenhaber bracket of degree $1-d$ on $\HH^\bullet(\cC)$ via Kontsevich--Tamarkin formality and Costello--Gwilliam--Li locality; Stage~2 applies the chiral specialisation functor $\SpCh_{\Sigma_{d-1},C}$ to a $(d-1)$-dimensional cycle and reference curve. The operadic level $n_{\mathrm{native}}(d)$ of the chiral output on $C$ is determined by the Gerstenhaber bracket degree alone (Proposition~\ref{prop:native-en-level}): $\infty$ at $d=1$, $2$ at $d=2$, and $1$ at $d \geq 3$. The construction (Chapter~\ref{ch:cy-to-chiral}) proceeds concretely by (i) extracting the Gerstenhaber algebra $\HH^\bullet(\cC)$, (ii) promoting it to a Lie conformal algebra using the CY pairing, (iii) taking the factorization envelope on a curve, and (iv) at $d = 2$ enhancing to $\Etwo$ using the $\bS^2$-framing.

\begin{theorem}[Categorical kappa equals chiral kappa at \texorpdfstring{$d = 2$}{d = 2}, \texorpdfstring{$h^{1,0} = 0$}{h-1,0 = 0}: CY-D\texorpdfstring{$_2$}{-2}]
\label{thm:cy-d-d2}
For $\cC$ a smooth proper $\CY_2$ category such that $\Phi(\cC) = \SpCh_{\Sigma_1,C}(\PhiFA_2(\cC))$ exists \emph{and} $\HH_{-1}(\cC) = 0$ (equivalently, in the geometric case $\cC = D^b(\Coh(X))$, the hypothesis $h^{1,0}(X) = 0$), the chiral modular characteristic, namely the tier-(iii) invariant read off the stage-$2$ output on $C$, satisfies
\[
 \kappa_{\mathrm{ch}}\big(\Phi(\cC)\big) \;=\; \chi^{\CY}(\cC) \;=\; \chi(\cO_X) \;=\; \kappa_{\mathrm{cat}}(X).
\]
The middle equality is the categorical Euler characteristic via the Mukai pairing (tier-(i) invariant of the input); the right equality is the Hodge-supertrace formula of Remark~\textup{\ref{rem:kappa-cat-from-chi}}. The identity is a coincidence of a tier-(iii) output with a tier-(i) input, specific to $d = 2$ with $h^{1,0}=0$: for $\cC = D^b(\Coh(X))$ with $X$ a K3 surface, $\kappa_{\mathrm{ch}}(\Phi(D^b(\Coh(K3)))) = \kappa_{\mathrm{cat}}(K3) = 2$. The canonical Stage-2 product chiral rank $\kappa_{\mathrm{ch}}^{\mathrm{Heis}}(K3 \times E) = 3$ obtains by Beauville additivity under tensor products; the compact total-space trace is $\kappa_{\mathrm{ch}}(A_{K3\times E})=0$. The four-way construction separation $\{\kappa_{\mathrm{cat}}, \kappa_{\mathrm{ch}}^{\mathrm{Heis}}, \kappa_{\mathrm{BKM}}, \kappa_{\mathrm{fiber}}\} = \{0, 3, 5, 24\}$ for K3 $\times$ E is laid out in Remark~\textup{\ref{rem:kappa-cat-from-chi}}.
\end{theorem}

\begin{remark}[Hypothesis \texorpdfstring{$h^{1,0} = 0$}{h-1,0 = 0} is essential at \texorpdfstring{$d = 2$}{d = 2}]
\label{rem:cy-d-d2-hypothesis}
The identification $\kappa_{\mathrm{ch}} = \kappa_{\mathrm{cat}}$ fails without the hypothesis $h^{1,0}(X) = 0$. Counterexample: the abelian surface $T^4$ has $h^{1,0}(T^4) = 2$ and $\chi(\cO_{T^4}) = 0$, so $\kappa_{\mathrm{cat}}(T^4) = 0$, but additivity gives $\kappa_{\mathrm{ch}}(\Phi(T^4)) = \kappa_{\mathrm{ch}}(\Phi(E)) + \kappa_{\mathrm{ch}}(\Phi(E)) = 1 + 1 = 2 \neq 0$. The sharper Beauville formula of Proposition~\textup{\ref{prop:beauville-kappa-formula}}, $\kappa_{\mathrm{ch}}(\Phi(D^b(\Coh(X)))) = \chi(\cO_X) + h^{1,0}(X)$, reduces to $\kappa_{\mathrm{cat}} = \chi(\cO_X)$ exactly in the simply-connected $h^{1,0} = 0$ case. The most common silent failure mode is applying $\kappa_{\mathrm{ch}} = \chi(\cO_X)$ at odd $d$ where Serre forces $\chi(\cO_X) = 0$ but $\kappa_{\mathrm{ch}}$ is generically nonzero.
\end{remark}

\begin{remark}[The \texorpdfstring{$d = 3$}{d = 3} case]
\label{rem:cy-cat-d3}
For $d = 3$, the assignment $\Phi_3^{(\Sigma_2,C)}$ is constructed on the stated object-level locus only after the H1--H4 data, the admissible specialisation, and the obstruction comparison data of Theorem~\ref{thm:cy-to-chiral-d3} have been fixed. The required input is not the slogan $\HH^{-2}_{E_1}=0$: it is either a corrected TCFT carrier $B^{(2)}_{\mathrm{TCFT}}$ with its comparison to the obstruction complex, or an explicit $\HH^{-2}_{E_1}$ filtration theorem with complete/exhaustive/separated strong convergence and empty total-degree $-2$ line. The raw carrier $B^{(2)}_{\mathrm{term}}$ is excluded by the strict witness $[m_3,B^{(2)}_{\mathrm{term}}][a|a|a|a|b]=2\alpha[b]\neq0$. The $\mathbb{C}^3$ Hall/representation chain is verified in Chapter~\ref{ch:cy-to-chiral}, Theorem~\ref{thm:c3-functor-chain}; the finite Rees hCS--Hall comparison is constructed on $\mathbb{C}^3$ charts and finite DWR/Ran critical atlases by Chapter~\ref{ch:cy3-chain-level-bridge}, while compact Hall doubles require the extra completion, realisation, compact-support, and pairing data isolated below. At $d = 3$ the identification $\kappa_{\mathrm{ch}} = \kappa_{\mathrm{cat}}$ is false for every strict CY$_3$ with $h^{1,0} = 0$: the RHS vanishes by Serre, while the LHS is generically nonzero. The dimension-stratified replacement formula is Conjecture~\textup{\ref{conj:cy-kappa-identification}}, with concrete values catalogued in Theorem~\textup{\ref{thm:kappa-stratification-by-d}} (e.g.\ $\kappa_{\mathrm{ch}}^{\mathrm{Heis}}(\text{quintic}) = -25/3$, $\kappa_{\mathrm{ch}}^{\mathrm{Heis}}(K3 \times E) = 3$, $\kappa_{\mathrm{ch}}^{\mathrm{Heis}}(\text{local }\mathbb{P}^2) = 3/2$; cf.\ Remark~\ref{rem:beauville-kappa-formula-subscript-split}).
\end{remark}

\begin{remark}[Fano obstruction to absolute HPD on \texorpdfstring{$\mathrm{K3}\times E$}{K3x E}; relative HPD over \texorpdfstring{$E$}{E} as fourth Hochschild-pairing path]
\label{rem:cy-cat-hpd-k3e-relative}
Kuznetsov homological projective duality (HPD; Kuznetsov~$2007$~Publ.\ IHES~$105$~Thm.~$1.1$) attaches a dual Lefschetz pair $(Y^\vee,\mathcal A^\vee_\bullet)$ to a \emph{Fano} variety $Y\hookrightarrow\mathbb P^N$ with Lefschetz decomposition $D^b(\Coh(Y))=\langle\mathcal A_0,\mathcal A_1(1),\ldots,\mathcal A_{m-1}(m-1)\rangle$, identifying Kuznetsov components across linear sections. The Fano requirement is load-bearing: the Lefschetz twist is $\mathcal O_Y(1)=\omega_Y^{-1/r}$ for some $r\ge 1$, and this twist is trivial on any strict CY. For $Y=\mathrm{K3}\times E$ the canonical bundle is $\omega_Y\cong\omega_{\mathrm{K3}}\boxtimes\omega_E\cong\mathcal O_Y$ by CY-$2$ triviality of $\omega_{\mathrm{K3}}$ and CY-$1$ triviality of $\omega_E$: absolute HPD does not apply to the product.

The replacement is the \emph{relative} HPD of Kuznetsov~$2018$~Crelle~$736$~\S~$2$ and Kuznetsov--Perry~$2021$ JEMS~$23$ Thm.~$1.1$, applied fibrewise along the second projection $\pi:Y\to E$. Each fibre is a K3 surface; the Kuznetsov--Markushevich~$2009$~J.\ Reine Angew.\ Math.\ $627$ Thm.~$1.1$ / Addington--Thomas~$2014$~Duke Math.\ J.\ $163$ Thm.~$1.1$ criterion identifies generic K3 of degree~$14$ with the Kuznetsov component of a smooth cubic fourfold $Y_4\subset\mathbb P^5$, globalising over~$E$ to a cubic-fourfold fibration $\widetilde Y\to E$ with $\Kuz(\widetilde Y/E)\simeq D^b(\Coh(Y))$. Relative HPD then transports the top-degree categorical Hochschild class $[\chi_3^{\mathrm{cat}}]\in\mathrm{HH}^3_{\mathrm{cat}}(\Kuz(\widetilde Y^\vee/E))$ to the CY-$3$ volume class on~$Y$; under the framed K3-fibre specialisation $\Phi_3^{(K3,E)}$ this becomes the chiral Hochschild cocycle $[\chi_3]\in\ChirHoch^3_{\mathrm{CY\text{-}3}}(\mathbf H_{\Delta_5})$ of Vol~I~Thm.~\texttt{chapters/theory/hochschild\_cohomology.tex} (Vol~I). The value $\langle[\chi_3],[e_3^Y]\rangle_{\Phi_3^{(K3,E)}}=2\,\mathrm{Vol}(E)\cdot(2\pi\mathrm i)^3$ recovers through the Kuznetsov--Perry transport, supplying a fourth categorically-distinct verification of $[\chi_3]\neq 0$ (Vol~I~Thm.~\texttt{chapters/theory/hochschild\_cohomology.tex} (Vol~I)).

Absolute HPD applies to Fano $Y$ only. Relative HPD over a base~$S$ applies to fibrations $Y\to S$ whose fibres are Fano; the base need not be Fano. For $Y=\mathrm{K3}\times E\to E$, neither side is Fano, yet the Kuznetsov--Markushevich upgrade of the K3 fibres to cubic-fourfold Kuznetsov components makes the fibres Fano. Conflating absolute HPD on the CY product with fibrewise absolute HPD on K3 (after the Kuznetsov--Markushevich upgrade) changes the theorem; the relative framework keeps the two scopes separate.
\end{remark}

\providecommand{\Kuz}{\mathrm{Ku}}

\begin{remark}[Drinfeld centre as right adjoint]
\label{rem:cy-cat-drinfeld-center-not-averaging}
The downstream correspondence $\Phi_d$ at $d \geq 3$ produces an $\Eone$-chiral algebra on constructed specialisation loci; the braided $\Etwo$-structure is recovered on the Drinfeld centre $\cZ(\Rep^{\Eone}(\Phi_3^{(\Sigma_2,C)}(\cC)))$ (property~\textup{(U3)} of Theorem~\textup{\ref{thm:phi-platonic}}). The Drinfeld centre is the right adjoint to the forgetful functor $\mathrm{BrMon} \to \mathrm{Mon}$, equivalently the categorified commutant $z(A) = \{a \in A: ab = ba\ \forall b \in A\}$. It differs from a categorified averaging map. The averaging map $\Eone \to \Einf$ loses quantum-group data; the Drinfeld-centre passage $\Eone \to \Etwo$ constructs braiding via half-braidings. This distinction is structural in Vol~III: the chiral-quantum-group content of $\Phi(\cC)$ at $d \geq 3$ is preserved by the centre and lost by symmetric averaging. We therefore reserve ``averaging'' for the $\Eone \to \Einf$ coinvariant projection and ``Drinfeld centre'' for the right-adjoint / half-braiding construction; the two are distinct functors producing distinct output categories.
\end{remark}

\begin{remark}[What \texorpdfstring{$\Phi$}{Phi} does not see, and what CoHA is]
\label{rem:phi-does-not-see}
Not every invariant of $\cC$ survives the passage through $\Phi$. The correspondence retains the cyclic $\Ainf$-structure and the Gerstenhaber bracket; it loses the fine structure of the triangulated / $\Ainf$ category (individual objects, t-structures, stability conditions) visible only via the categorical Cohomological Hall Algebra construction of Chapter~\textup{\ref{ch:toric-coha}}. The Cohomological Hall Algebra $\CoHA(X)$ is the Hochschild cohomology of the quiver-with-potential category, equipped with the Schiffmann--Vasserot--Yang--Zhao multiplication. It is an associative algebra with a Hall product, not a chiral algebra; the identifications ``$\CoHA = \Eone$-chiral'' and ``$\CoHA$ carries a vertex algebra structure'' conflate an associative-algebra output with a chiral-algebra output and so fail as stated. The connection to chiral algebras goes by two separate routes, \emph{not} by a single application of $\Phi$: on one side, the Schiffmann--Vasserot 2013 direct identification $\CoHA(\C^3) = Y^+(\widehat{\mathfrak{gl}}_1)$ (positive half of the affine Yangian of $\widehat{\mathfrak{gl}}_1$) is an equality of associative algebras, not a $\Phi$-application; $\CoHA(\C^3)$ is neither a CY-category nor the $\Phi_3$-input. On the other side, the framed toric specialisation $\Phi_3^{(\Sigma_2,C)}(D^b(\Coh(\C^3)))$ produces a chiral algebra on a reference curve (property~\textup{(U4)} of Theorem~\textup{\ref{thm:phi-platonic}}); specialising via $\SpCh_{\Sigma_2, C}$ with $\Sigma_2 \subset \C^3$ a coordinate plane yields $Y^+(\widehat{\mathfrak{gl}}_1)$ as the Stage-2 shadow on $C$, and the full $\cW_{1+\infty}$ at $c = 1$ then arises from the Drinfeld-centre passage (Remark~\textup{\ref{rem:cy-cat-drinfeld-center-not-averaging}}). The two routes converge on $Y^+(\widehat{\mathfrak{gl}}_1)$ because the Schiffmann--Vasserot CoHA coincides with the $\SpCh$-specialisation of $\PhiFA_3(D^b(\Coh(\C^3)))$ (Kontsevich--Soibelman 2008 \S 4; Schiffmann--Vasserot 2013 \S 6). Writing ``$\Phi(\CoHA(\C^3))$'' would be a type error: $\Phi$ takes CY-categories to chiral algebras and $\CoHA$ is neither.

The surviving numerical data of $\Phi$ are organised by the $\kappa_\bullet$-spectrum (Remark~\textup{\ref{rem:cy-cat-kappa-convention}}): $\kappa_{\mathrm{cat}}$ is the manifold invariant always present, $\kappa_{\mathrm{ch}}$ is the algebraization invariant defined when $\Phi$ exists, $\kappa_{\mathrm{BKM}}$ is the Borcherds-side algebraisation invariant defined for K3-fibered CY$_3$s (Class~A; Proposition~\textup{\ref{prop:bkm-weight-universal}}), and $\kappa_{\mathrm{fiber}}$ is the lattice rank.
\end{remark}

\section{Finite hCS--Hall charts and the compact Hall boundary}
\label{sec:cycat-finite-hcs-hall-boundary}

The CY$_3$ categorical input produces a holomorphic factorisation algebra
before it produces a Hall algebra. A Hall statement enters only after a
critical chart, an orientation line, a Rees/vanishing-cycle realisation,
and the relevant descent topology have been fixed.

\begin{definition}[Finite Rees hCS--Hall comparison chart]
\label{def:cycat-finite-rees-hcs-hall-chart}
Let $X$ be a smooth CY$_3$ and let
\(\sigma\in\mathsf N^{\Ran}_{\leq m}(\mathfrak U;\Gamma_{\leq L})\)
be a finite simplex in a Dolbeault/Weiss/Ran (DWR)-good cover. A
\emph{finite Rees hCS--Hall comparison chart} on $\sigma$ consists of:
\begin{enumerate}[label=\textup{(\roman*)}]
 \item the finite hCS quotient
 \(\Obs_{\hCS}^{q,\leq r,\leq N}(\sigma)_{\leq L}\), with holomorphic
 jet order $\leq N$, renormalisation order $\leq r$, and charge bound
 $\Gamma_{\leq L}$;
 \item a finite cyclic critical model
 \((V_\sigma,W_\sigma,\omega_\sigma)\) over
 \(\Omega^\bullet(\Delta^p)\), obtained by cyclic homological
 perturbation from the hCS BV complex;
 \item the oriented Rees critical Hall complex
 \[
  \Hall_{\sigma}^{\mathrm{Rees},\mathrm{or}}
  =
  \left(
  C^\bullet\!\left(
  \mathfrak g_\sigma,
  \Gamma(V_\sigma,\wedge^\bullet T_{V_\sigma})[[\hbar]]
  \otimes\Omega^\bullet(\Delta^p)
  \right)^{G_\sigma},
  d_{\CE}+\iota_{d_VW_\sigma}
  +\hbar\Delta_{\omega_\sigma}+d_{\Delta^p}
  \right)
  \otimes o_\sigma[s_\sigma](t_\sigma),
 \]
 with determinant-line square root \(o_\sigma\), shift \(s_\sigma\),
 and Tate twist \(t_\sigma\);
 \item a relative degree-zero chain map
 \[
  \theta_\sigma^{\mathrm{rel}}\colon
  B\Obs_{\hCS}^{q,\leq r,\leq N}(\sigma)_{\leq L}
  \longrightarrow
  \Hall_{\sigma}^{\mathrm{Rees},\mathrm{or}}(\Delta^p)
 \]
 compatible with faces, refinements, disjoint union, Thom--Sebastiani,
 and the orientation transports.
\end{enumerate}
This is the finite chart of
Definition~\ref{def:finite-dwr-ran-cyclic-critical-atlas}, rewritten at
the categorical input level. It is a Rees critical complex, not yet the
compact Borel--Moore CoHA unless a monoidal vanishing-cycle realisation
has also been supplied.
\end{definition}

\begin{proposition}[Finite DWR/Ran Rees hCS--Hall comparison]
\label{prop:cycat-finite-dwr-ran-rees-hcs-hall}
Fix finite bounds $(N,r,L,m)$. If a finite DWR/Ran family of charts
\(\{\theta_\sigma^{\mathrm{rel}}\}\) satisfies the compatibility
conditions of Definition~\ref{def:cycat-finite-rees-hcs-hall-chart},
then the sum over simplices
\[
 \Theta_{\hCS\to\Hall}^{\mathrm{Rees},\mathrm{or};N,r,L,m}
 =
 \sum_{0\leq p\leq m}
 \sum_{\sigma\in
 \mathsf N^{\Ran}_{p,\leq m}(\mathfrak U;\Gamma_{\leq L})}
 \int_{\Delta^p}\theta_\sigma^{\mathrm{rel}}
\]
is a Maurer--Cartan element in the finite convolution Lie algebra
\[
 \mathfrak M_{\hCS,\Hall}^{N,r,L,m}
 =
 \Tot_{\sigma}
 \Hom_{\mathrm{cont}}
 \left(
 B\Obs_{\hCS}^{q,\leq r,\leq N}(\sigma)_{\leq L},
 \Hall_{\sigma}^{\mathrm{Rees},\mathrm{or}}
 \right).
\]
Equivalently, it is a continuous multiplicative natural transformation
\[
 \Obs_{\hCS}^{q,\leq r,\leq N}(-)_{\leq L}
 \longrightarrow
 \Hall_{\mathrm{crit}}^{\mathrm{Rees},\mathrm{or},\leq N,\leq L}(-)
\]
on the finite DWR/Ran nerve. The construction covers
\(\mathbb C^3\), finite affine CY$_3$ critical charts, and finite
multi-chart \v{C}ech/Ran diagrams for which the same finite cyclic
critical atlas and orientation data are supplied.
\end{proposition}

\begin{proof}
The cyclic homological perturbation theorem transfers the finite hCS BV
operations to the Hamiltonian \(W_\sigma\). The Fourier--BV
identification of Definition~\ref{def:rees-critical-hall-interlayer}
turns the differential into
\[
 d_{\CE}+\iota_{d_VW_\sigma}+\hbar\Delta_{\omega_\sigma}+d_{\Delta^p}.
\]
The relative master equation is therefore the Maurer--Cartan equation
for \(\theta_\sigma^{\mathrm{rel}}\) in the relative convolution Lie
algebra. Integration over \(\Delta^p\) converts the
\(d_{\Delta^p}\)-term, by Stokes' formula, into the alternating
\v{C}ech/Ran face differential. Summing over all simplices gives
\[
 D_{\mathrm{tot}}\Theta
 +\frac{1}{2}[\Theta,\Theta]=0,
\]
which is exactly the finite descent condition. Multiplicativity is the
Rees Thom--Sebastiani identity for \(W_{\sigma_1}\boxplus W_{\sigma_2}\)
combined with the bar coproduct on hCS observables and the Hall product
on the Rees side. For \(\mathbb C^3\), the chart is the one-vertex
three-loop critical model whose Hall cohomology is the
Schiffmann--Vasserot positive half
\(\CoHA(\mathbb C^3)\cong \Yplus(\widehat{\mathfrak{gl}}_1)\); a finite
affine CY$_3$ critical chart supplies the same data by its
quiver-with-potential critical function; a finite multi-chart diagram
is the preceding construction with the \v{C}ech/Ran face maps retained.
\end{proof}

\begin{definition}[Compact Hall-double realisation data]
\label{def:cycat-compact-hall-double-realisation-data}
For a compact CY$_3$ \(X\), a finite Rees comparison upgrades to a
compact Hall-double comparison only after the following data are fixed:
\begin{enumerate}[label=\textup{(\roman*)}]
 \item a monoidal vanishing-cycle realisation from the Rees critical
 complexes to Borel--Moore critical CoHA, compatible with
 Thom--Sebastiani, proper Hall pushforward, lci pullback, and the
 determinant-line orientation system;
 \item compact-support Beck--Chevalley and DWR/Weiss descent for source
 and target;
 \item a charge/HN and \(\hbar\)-adic completion in which all Hall
 products, coproducts, and inverse limits are continuous;
 \item a Serre-dual negative half, a Cartan half, and a continuous
 Hall--Serre pairing whose radical has been quotiented;
 \item for a Hall--Drinfeld output, centre/half-braiding compatibility
 with the specialised \(\Eone\)-chiral algebra.
\end{enumerate}
A Serre-equivariant quasi-NCCR or a compact critical CoHA character map
is one possible positive-half input to this package. It is not itself a
Hall double.
\end{definition}

\begin{proposition}[Finite positive halves do not define compact Hall doubles]
\label{prop:cycat-compact-hall-double-boundary}
The finite maps of Proposition~\ref{prop:cycat-finite-dwr-ran-rees-hcs-hall}
and the \(K3\times E\) finite Rees maps
\[
 \rho_H^+\colon
 \operatorname{Rees}_{\le H}
 \Yplus_{\mathrm{stalk}}(K3\times E;\mathfrak U)
 \longrightarrow
 \operatorname{Rees}_{\le H}
 \mathcal H_{K3\times E}^{\mathrm{comp},+}
\]
of Proposition~\ref{gluing:prop:compact-coha-rees-maps} construct
positive-half Hall data in finite bounds. They do not construct a
compact Hall--Drinfeld double, a quasi-NCCR realisation map, or a map
from \(\PhiFA_3(\mathrm{Perf}(X))\) to a compact quantum group. Such a
map exists only after the realisation data of
Definition~\ref{def:cycat-compact-hall-double-realisation-data} have
been supplied.
\end{proposition}

\begin{proof}
The Rees comparison of
Proposition~\ref{prop:cycat-finite-dwr-ran-rees-hcs-hall} lands before
vanishing-cycle realisation. Passing to
\(\CoHA_{\mathrm{crit}}^{\mathrm{or}}\) requires the monoidal
realisation functor and the compact-support compatibilities listed in
Definition~\ref{def:cycat-compact-hall-double-realisation-data}. On
\(K3\times E\), Proposition~\ref{gluing:prop:compact-coha-rees-maps}
constructs the compatible finite maps \(\rho_H^+\), but explicitly does
not assert that they are injective, surjective, or primitive-preserving
isomorphisms. Remark~\ref{gluing:rem:no-double-through-hocolim} states
the obstruction in categorical form: a Drinfeld double cannot be
applied to the positive-half homotopy colimit without a completed
monoidal category, dualisability, a non-degenerate Hopf pairing, the
Borcherds bracket comparison, and centre compatibility. The conditional
double of Theorem~\ref{thm:quantum-group-as-positive-geometry-double}
has exactly these hypotheses. Therefore finite DWR/Ran Rees
construction proves the finite positive-half comparison and stops
strictly before the compact Hall-double map.
\end{proof}

\begin{remark}[Cross-volume conventions]
\label{rem:cy-cat-cross-vol-conventions}
Vol~III uses lambda-bracket conventions for the Lie conformal algebras produced by $\Phi$: an OPE of the form $T(z) T(w) \sim (c/2)(z - w)^{-4} + \cdots$ is rewritten as $\{T_\lambda T\} = (c/12) \lambda^3 + \cdots$, absorbing the combinatorial factor $3! = 6$ from the divided-power $\lambda^{(3)} = \lambda^3/3!$. Vol~I uses OPE modes directly; care is required when transporting formulas between volumes (see the concordance). The Hochschild / cyclic invariants of this chapter are convention-independent: they depend only on the chain-level $\Ainf$-structure.
\end{remark}

\subsection{Chain-level \texorpdfstring{$E_d$}{E-d} operad: Fulton--MacPherson configuration spaces}
\label{subsec:cy-cat-fm-config-spaces}

Stage~$1$ of Theorem~\ref{thm:phi-two-stage-factorisation} produces an
$E_d$-holomorphic factorisation algebra $\PhiFA_d(\cC)$. The $E_d$-operad
that appears on the output side has concrete chain-level content: its
arity-$k$ space of operations is the Fulton--MacPherson compactification
$\overline{F}(k, \mathbb{R}^d)^{\mathrm{fm}}$ of the ordered configuration space
of $k$ points in $\mathbb{R}^d$. This subsection records the
Fulton--MacPherson model in the precise form consumed by Kontsevich--Tamarkin
formality and Costello--Gwilliam--Li locality, for both $d = 2$ (the
proved $\PhiFA_2$) and $d = 3$ (the Stage-$1$ half of the CY-A$_3$
construction).

\providecommand{\FMbar}{\overline{F}}
\providecommand{\FMop}{\mathcal{FM}}
\providecommand{\Gerst}{\mathrm{Gerst}}

\begin{definition}[Ordered configuration space; Fulton--MacPherson compactification]
\label{def:cy-cat-fm-config}
For a smooth manifold $M$ of dimension $n$ and $k \geq 1$, the
\emph{ordered configuration space} is
$F(k, M) = \{(x_1, \ldots, x_k) \in M^k : x_i \neq x_j\}$, an open
submanifold of $M^k$ carrying a free $S_k$-action. The
\emph{Fulton--MacPherson compactification}
$\FMbar(k, M)^{\mathrm{fm}}$ (Fulton--MacPherson, \emph{A compactification of configuration spaces}, Ann.\ of Math.\ \textbf{139} (1994) \S 1) is the closure of
$F(k, M)$ in $M^k \times \prod_{|I| \geq 2} \mathrm{Bl}_{\Delta_I}(M^I)$
under the simultaneous inclusions, where $\Delta_I \subset M^I$ is the
thin diagonal indexed by $I \subseteq \{1, \ldots, k\}$ and
$\mathrm{Bl}$ is iterated real blow-up. The space $\FMbar(k, M)^{\mathrm{fm}}$
is a smooth manifold with corners, compact when $M$ is compact, with
$F(k, M)$ embedded as an open dense submanifold.
\end{definition}

\begin{theorem}[Boundary stratification by rooted trees]
\label{thm:cy-cat-fm-strata}
The boundary $\partial \FMbar(k, M)^{\mathrm{fm}}$ stratifies by rooted trees
$T$ with $k$ leaves and no bivalent internal vertices. The stratum
$S_T$ is canonically
\[
 S_T \;\simeq\; F(|v_T|, M) \,\times\, \prod_{v \in \mathrm{int}(T)} F(|v|, \mathbb{R}^n),
\]
where $v_T$ is the root, $|v|$ the number of children at an internal
vertex, and $\mathrm{int}(T)$ the non-root internal vertices. The
$F(|v|, \mathbb{R}^n)$-factors record the relative scale-and-position
data of points colliding into the cluster at $v$, modulo
$\mathbb{R}^n \rtimes \mathbb{R}_{>0}$. For $M = \mathbb{R}^n$ one
passes to the reduced compactification
$\widetilde{\FMbar}(k, \mathbb{R}^n)^{\mathrm{fm}} =
\FMbar(k, \mathbb{R}^n)^{\mathrm{fm}}/(\mathbb{R}^n \rtimes \mathbb{R}_{>0})$,
a smooth manifold with corners of real dimension $nk - n - 1$.
\end{theorem}

\begin{theorem}[Fulton--MacPherson operad; Salvatore--Sinha comparison with little discs]
\label{thm:cy-cat-fm-operad}
The collection
$\FMop_n = \{\widetilde{\FMbar}(k, \mathbb{R}^n)^{\mathrm{fm}}\}_{k \geq 1}$
carries an operad structure in smooth manifolds with corners: the
composition
$\gamma \colon \widetilde{\FMbar}(k, \mathbb{R}^n)^{\mathrm{fm}} \times
 \prod_i \widetilde{\FMbar}(\ell_i, \mathbb{R}^n)^{\mathrm{fm}}
 \to \widetilde{\FMbar}(\sum_i \ell_i, \mathbb{R}^n)^{\mathrm{fm}}$
is a smooth diffeomorphism onto the codimension-$(k-1)$ boundary
stratum $S_{T_{k, \ell_1, \ldots, \ell_k}}$ indexed by the rooted
$2$-level tree with root arity $k$ and children arities $\ell_i$.
The resulting operad is weakly equivalent to the little $n$-discs
operad $\Ddisk_n$ and serves as a $\Sigma$-cofibrant (hence homotopically
good) model for $E_n$; it is weakly equivalent via the
Boardman--Vogt $W$-construction
(Boardman--Vogt, \emph{Homotopy invariant algebraic structures on topological spaces}, Lecture Notes in Math.\ \textbf{347}, 1973; Berger--Moerdijk, \emph{The Boardman--Vogt resolution of operads in monoidal model categories}, Topology \textbf{45} (2006) 807--849, Thm.\ 5.1) to
$W \Ddisk_n$, with length-decorated rooted trees mapping to boundary
strata of $\FMbar(k, \mathbb{R}^n)^{\mathrm{fm}}$.
\end{theorem}

\begin{example}[Cohomology as Gerstenhaber$_n$]
\label{ex:cy-cat-fm-cohomology}
$H^\bullet(F(k, \mathbb{R}^n); \mathbb{Q})$ is generated by the
Arnold classes $g_{ij}$ of degree $n - 1$, modulo $g_{ij}^2 = 0$
and the Arnold triple identity
$g_{ij} g_{jk} + g_{jk} g_{ki} + g_{ki} g_{ij} = 0$
(Arnold, \emph{Mat.\ Zametki} \textbf{5} (1969) 227--231). For $n \geq 2$ this gives
$H^\bullet(\FMop_n; \mathbb{Q}) \simeq \Gerst_n$, the operad of
$(n - 1)$-shifted Poisson algebras: a degree-$0$ graded commutative
product and a Lie bracket $\{-, -\}$ of degree $-(n - 1)$ satisfying
Leibniz (F.~Cohen, \emph{The homology of iterated loop spaces}, Lecture Notes in Math.\ \textbf{533}, 1976; Getzler--Jones, \emph{Operads, homotopy algebra and iterated integrals for double loop spaces}, arXiv:hep-th/9403055). At $n = 3$ the bracket
has degree $-2$, recovering the CY-$3$ shifted-Poisson bracket on
$\HH^\bullet(\cC)$ of Proposition~\ref{prop:cy-bracket-spectrum}.
\end{example}

\begin{theorem}[Kontsevich--Tamarkin $E_n$-formality with $\GRT_1$-torsor]
\label{thm:cy-cat-kt-formality}
For $n \geq 2$ there is a quasi-isomorphism of dg-operads over
$\mathbb{Q}$
\[
 \mathcal{F} \colon
 \Gerst_n \;\xrightarrow{\ \sim\ }\;
 C^{\mathrm{sa}}_\bullet\!\bigl(\FMop_n;\, \mathbb{Q}\bigr),
\]
where $C^{\mathrm{sa}}_\bullet(-; \mathbb{Q})$ is the semi-algebraic
chain complex of Kontsevich--Soibelman (strictly operadic, strictly
multiplicative under Cartesian product). The space of such
$\mathcal{F}$, up to homotopy, is a torsor under the pro-unipotent
Grothendieck--Teichm\"uller group $\GRT_1(\mathbb{Q})$:
\[
 \{\Gerst_n\text{-formality morphisms}\}/{\simeq}
 \;\text{is a }\GRT_1(\mathbb{Q})\text{-torsor}.
\]
The $\GRT_1$-action is realised through Kontsevich graph-integral
weights $W_\Gamma$ evaluated on $\widetilde{\FMbar}(k, \mathbb{R}^n)^{\mathrm{fm}}$
\cite{Kontsevich2003} \S 6.4. Willwacher's theorem
\cite{Willwacher2015} Thm.\ 1.1 identifies
$H^0(\GC_n)$ of the Kontsevich graph complex with the
Grothendieck--Teichm\"uller Lie algebra $\mathfrak{grt}_1$, uniformly
for $n \geq 2$.
\end{theorem}

\begin{construction}[Stage-$1$ $\PhiFA_d(\cC)$ as $\FMop_d$-colimit of Hochschild chains]
\label{constr:cy-cat-phifa-as-fm-colimit}
Fix a smooth complex CY$_d$ variety $X$ and a smooth proper cyclic
$\Ainf$ category $\cC$ of dimension $d$. A \emph{holomorphic $d$-disc}
in $X$ is an open embedding $\varphi \colon B^{2d}_\epsilon
\hookrightarrow X$ from an open polydisc in $\mathbb{C}^d$. Stage-$1$
$\PhiFA_d(\cC)$ assigns to a finite collection of pairwise disjoint
holomorphic $d$-discs $B_1, \ldots, B_k \subset X$ the chain complex
\[
 \PhiFA_d(\cC)(B_1 \sqcup \cdots \sqcup B_k)
 \;\simeq\;
 C^{\mathrm{sa}}_\bullet\!\left(\widetilde{\FMbar}(k, \mathbb{R}^d)^{\mathrm{fm}}; \mathbb{Q}\right)
 \otimes_{C^{\mathrm{sa}}_\bullet(\FMop_d)}
 \HH_\bullet(\cC)^{\otimes k},
\]
the homotopy coend over the $\FMop_d$-operad, with $\HH_\bullet(\cC)$
supplied with its $E_d$-algebra structure through the Kontsevich--Tamarkin
formality morphism of Theorem~\ref{thm:cy-cat-kt-formality}. The
holomorphic refinement is the Dolbeault acyclification over $X$: the
assignment is a sheaf of chain complexes with $\bar\partial$-acyclic
extensions, rigidified by the holomorphic volume $\Omega_X \in
H^0(X, K_X)$ through its role as the BV-rigidifier of the
Costello--Li quantisation scheme
\cite{CostelloLi2016,CostelloLi2020}.
Functorially, the arity-$k$ operation assembles the Hochschild chains
indexed by the external leaves of the rooted trees of
Theorem~\ref{thm:cy-cat-fm-strata} along the operadic composition
$\gamma$.
\end{construction}

\begin{proposition}[Chain-level Dunn additivity $\FMop_3 \simeq \FMop_2 \otimes \FMop_1$]
\label{prop:cy-cat-fm-dunn}
For $m, n \geq 1$ there is a quasi-isomorphism of dg-operads over
$\mathbb{Q}$
\[
 C^{\mathrm{sa}}_\bullet(\FMop_{m+n}; \mathbb{Q})
 \;\xrightarrow{\ \sim\ }\;
 C^{\mathrm{sa}}_\bullet(\FMop_m; \mathbb{Q})
 \otimes_H
 C^{\mathrm{sa}}_\bullet(\FMop_n; \mathbb{Q}),
\]
induced by the coordinate projection
$\widetilde{\FMbar}(k, \mathbb{R}^{m+n})^{\mathrm{fm}} \to
 \widetilde{\FMbar}(k, \mathbb{R}^m)^{\mathrm{fm}} \times
 \widetilde{\FMbar}(k, \mathbb{R}^n)^{\mathrm{fm}}$
(Hadamard tensor of operads). At $(m, n) = (2, 1)$ this recovers
$\FMop_3 \simeq \FMop_2 \otimes \FMop_1$, the chain-level source of the
derived-framing obstruction target
$\HH^{-2}_{E_1}(\HH_\bullet(\cC), \HH_\bullet(\cC))$. Its vanishing is
not a formal consequence of Dunn additivity; Theorem~\ref{thm:derived-framing-obstruction}
requires the comparison and filtration hypotheses stated there. The weak equivalence
is chain-homotopical (a quasi-isomorphism of dg-operads), not a
diffeomorphism of operads in smooth manifolds with corners:
$\FMbar(k, \mathbb{R}^3)^{\mathrm{fm}}$ is not smoothly a product of
$\FMbar(k, \mathbb{R}^2)^{\mathrm{fm}} \times \FMbar(k, \mathbb{R}^1)^{\mathrm{fm}}$.
\end{proposition}

\begin{proposition}[Chain-level Stage-$2$: $\FMop_3 \to \FMop_1$ push-forward along $\Sigma_2$]
\label{prop:cy-cat-fm3-to-fm1}
For an admissible specialisation datum $(\Sigma_2, C)$ with
$\Sigma_2 \subset X$ a smooth real $2$-cycle and $C \subset
X \setminus \Sigma_2$ a holomorphic reference curve, the Stage-$2$
specialisation functor $\SpCh_{\Sigma_2, C}$ is implemented at the
chain level by the operadic push-forward along the fibration
\[
 \widetilde{\FMbar}(k, \mathbb{R}^3)^{\mathrm{fm}}
 \;\longrightarrow\;
 \widetilde{\FMbar}(k, \mathbb{R}^1)^{\mathrm{fm}},
 \qquad
 (x_i^1, x_i^2, x_i^3)_{i=1}^{k}
 \;\mapsto\; (x_i^1)_{i=1}^{k},
\]
projecting each configuration point onto its reference-curve coordinate.
The fibre over a configuration $(x_1^1, \ldots, x_k^1) \in
\widetilde{\FMbar}(k, \mathbb{R}^1)^{\mathrm{fm}}$ is the $\Sigma_2$-direction
compactification of the configuration $(x_i^2, x_i^3)_{i=1}^{k}$,
topologically a $2k$-dimensional fibre. Factorisation homology is
integration along the fibres; the output is an $\Eone$-chiral algebra
on $C$ with operadic operations indexed by $\FMop_1$ composed with
holomorphic pullback. The $\pi_1(\Conf_2(\mathbb{R}^3)) = 0$ symmetry
of the $E_3$-level
(Fadell--Neuwirth, \emph{Configuration spaces}, Math.\ Scand.\ \textbf{10} (1962) 111--118) collapses in this
residual: the non-symmetric quantum-group $R$-matrix datum is
recovered from the Drinfeld centre $\cZ(\Rep^{\Eone}(A_\cC))$
(Remark~\ref{rem:cy-cat-drinfeld-center-not-averaging}), not from
chain-level $\FMop_3$-data.
\end{proposition}

\begin{remark}[Two Fulton--MacPherson spaces: topological $\FMop_3$ vs chiral $\FM_k(\mathbb{C})$]
\label{rem:cy-cat-two-fm-meanings}
The notation $\mathrm{FM}_k(C)$ appears in
Remark~\ref{rem:cy-cat-three-hochschild} for the chiral
Fulton--MacPherson compactification of $k$ points on the curve
$C$: the operadic base space on which the chiral Gerstenhaber bracket
lives through OPE residues. This is distinct from the
$\FMbar(k, \mathbb{R}^3)^{\mathrm{fm}}$ of the present subsection, which
lives on the CY$_3$-ambient side. The relation is
Proposition~\ref{prop:cy-cat-fm3-to-fm1}: the chiral $\FM_k(\mathbb{C})$ is
the Stage-$2$ residual of the ambient $\FMbar(k, \mathbb{R}^3)^{\mathrm{fm}}$
after factorisation-homology integration over $\Sigma_2$ and
holomorphic pullback to $C$. The two Fulton--MacPherson spaces are
\emph{different geometric objects} (one on the ambient manifold,
one on the reference curve) linked by the specialisation functor
$\SpCh_{\Sigma_2, C}$. Swapping the two would also swap the topological Gerstenhaber bracket of Remark~\ref{rem:cy-cat-three-hochschild} with the chiral Gerstenhaber bracket from OPE residues on $\FM_k(\mathbb{C})$, producing an algebra of the wrong type.
\end{remark}

\begin{remark}[$\GRT_1$-torsor and Stage-$1$ canonicity]
\label{rem:cy-cat-grt1-torsor-stage1}
Theorem~\ref{thm:cy-cat-kt-formality} is the chain-level incarnation
of the Stage-$1$ torsor-pointed clause of
Theorem~\ref{thm:phi-two-stage-factorisation}: different choices in
the $\GRT_1(\mathbb{Q})$-torsor produce Stage-$1$ outputs that are
$E_n$-quasi-isomorphic as factorisation algebras, the space of
$E_n$-formality morphisms from $\Gerst_n$ to
$C^{\mathrm{sa}}_\bullet(\FMop_n; \mathbb{Q})$ is a $\GRT_1$-torsor,
and the forgetful map to the space of $E_n$-quasi-isomorphism
classes contracts each torsor orbit to a single point. The
$\GRT_1$-action on graph-integral weights $W_\Gamma$ of
\cite{Kontsevich2003} (\S 6.4) twists the chain-level formality morphism but
preserves the homotopy category of $E_n$-algebras on the nose. In
the holomorphic refinement of
Construction~\ref{constr:cy-cat-phifa-as-fm-colimit} the
$\GRT_1$-action is absorbed into the Costello--Li scheme-dependence
of the BV counter-terms, with fixed observables
\cite{CostelloLi2020}. The torsor is the chain-level incarnation of
the motivic fundamental-group / graph-cohomology structure of
Remark~\ref{rem:cycat-ce-motivic-fundamental-group}.
\end{remark}

\subsection{Fourier--Mukai intertwining of \texorpdfstring{$\PhiFA_3$}{Phi-FA-3} under crepant birational CY\texorpdfstring{$_3$}{-3}}
\label{subsec:cycat-crepant-fm-intertwining}

Two smooth projective CY$_3$ varieties $X$ and $Y$ related by a
crepant birational equivalence $g \colon X \dashrightarrow Y$ have
derived categories connected by a Fourier--Mukai equivalence
$\Phi_g \colon D^b(\Coh(X)) \xrightarrow{\sim} D^b(\Coh(Y))$ with an
explicit kernel $\cK_g \in D^b(\Coh(X \times Y))$. The question
residual to the Stage-$1$ construction of
Theorem~\ref{thm:phi-two-stage-factorisation} is whether this
derived equivalence lifts through $\PhiFA_3$ to a chain-level
quasi-isomorphism of Stage-$1$ factorisation algebras; that is,
whether $\PhiFA_3$ sees the birational equivalence on the CY$_3$ side.

\providecommand{\KM}{\mathrm{KM}}

\begin{theorem}[Stage-$1$ crepant FM-intertwining]
\label{thm:cycat-stage1-crepant-fm-intertwining}
Let $g \colon X \dashrightarrow Y$ be a crepant birational equivalence
of smooth projective CY$_3$ varieties, with Fourier--Mukai kernel
$\cK_g \in D^b(\Coh(X \times Y))$ inducing the derived equivalence
$\Phi_g \colon D^b(\Coh(X)) \xrightarrow{\sim} D^b(\Coh(Y))$. Then
the Stage-$1$ factorisation-algebra assemblies are canonically
$E_3$-quasi-isomorphic:
\[
 \PhiFA_3(\Phi_g) \colon
 \PhiFA_3(D^b(\Coh(X)))
 \;\xrightarrow{\ \sim\ }\;
 \PhiFA_3(D^b(\Coh(Y))),
\]
with intertwining determined by the kernel $\cK_g$ through the
Caldararu Hochschild trace $\Tr_{\cK_g} \colon \HH^\bullet(X) \to
\HH^\bullet(Y)$ (Caldararu~\cite{Caldararu2005}~\S$8$) promoted to an
$E_3$-morphism via the Kontsevich--Tamarkin formality of
Theorem~\ref{thm:cy-cat-kt-formality}. Three input data:
\begin{enumerate}[label=(\roman*)]
 \item \emph{Categorical piece}: the Gerstenhaber-algebra
 quasi-isomorphism $\HH^\bullet(\Phi_g)$ on Hochschild cohomology is
 furnished by the kernel $\cK_g$; for crepant factorisations
 $g = g_n \circ \cdots \circ g_1$ into terminal flops, each
 $\HH^\bullet(\Phi_{g_i})$ is given by Bridgeland/Li, and the
 composition is associative.
 \item \emph{$E_3$-lifting}: the obstruction to lifting
 $\HH^\bullet(\Phi_g)$ to an $E_3$-morphism is in
 $\Ext^2_{\HH^\bullet}(\Gerst_3, \Gerst_3)$, which vanishes over
 $\QQ$ by Fresse~$2011$ \emph{AMS Math.\ Surveys} $217$ Thm~$2$;
 the $\GRT_1$-torsor accumulation across a flop chain is
 cohomologically trivial in the degrees seeing the $E_3$-structure
 (Willwacher~$2015$: $\mathfrak{grt}_1 = H^0(\mathrm{GC}_3)$ is
 generated in odd degree $\geq 3$).
 \item \emph{Holomorphic-volume transport}: the phase
 $e^{i\theta_g} \in \mathrm{U}(1)$ relating $g^*\Omega_Y$ to
 $\Omega_X$ is trivial because $g$ factors through connected
 components of the birational-automorphism group
 (Matsuki~$2002$ Ch.~$1$); hence $g^*\Omega_Y = \Omega_X$ strictly,
 and the Costello--Li one-loop BCOV curving
 $\alpha_{\mathrm{BCOV}} = (\chi_{\mathrm{top}}/24)[\Omega_X]^{0,1}$
 transports rigidly (Batyrev~$1999$: $\chi_{\mathrm{top}}$ is a
 birational invariant of smooth projective CY$_d$).
\end{enumerate}
\end{theorem}

\begin{proof}[Proof sketch]
For $g$ a simple flop, Bridgeland~$2002$ gives an explicit kernel
$\cK_g = \cO_{X \times_W X^+}$ on the fibre product over the
flopping contraction $W$, and the induced Hochschild
quasi-isomorphism $\HH^\bullet(\Phi_g)$ preserves the
Gerstenhaber bracket (Bodzenta--Bondal~$2015$). For $g$ a composition
$g = g_n \circ \cdots \circ g_1$ of terminal flops
(Kawamata--Matsuki~$2008$ Thm~$1$), the intertwining
$\PhiFA_3(\Phi_g) = \PhiFA_3(\Phi_{g_n}) \circ \cdots \circ
\PhiFA_3(\Phi_{g_1})$ assembles by associativity of factorisation-algebra
morphisms (Costello--Gwilliam~\cite{CostelloGwilliam} Prop.~$7.4$).
The $E_3$-lift from Gerstenhaber quasi-isomorphism to
$E_3$-quasi-isomorphism has vanishing obstruction (Fresse loc.\ cit.),
and the cumulative $\GRT_1$-twist is in the commutator subalgebra of
$\mathfrak{grt}_1$, which has no generators in cohomological degrees
seeing the $E_3$-structure (Willwacher~$2015$
\emph{Invent.}~$200$ Thm~$1.1$). The holomorphic-volume phase is
$+1$ because $\mathrm{Bir}(X, Y)/\mathrm{Iso}(X, Y)$ is connected
(Matsuki loc.\ cit.) and characters of connected groups are trivial.
$\chi_{\mathrm{top}}$-birational-invariance gives rigid transport of
$\alpha_{\mathrm{BCOV}}$.
\end{proof}

\begin{remark}[Chain-level transport via the common resolution]
\label{rem:cycat-atiyah-common-resolution}
For the Atiyah flop $g \colon X \dashrightarrow X^+$ (replacing a
$(-1,-1)$-curve $\mathbb{P}^1 \subset X$ with its dual), the common
resolution $Z \xrightarrow{p_X} X$, $Z \xrightarrow{p_Y} X^+$ is the
blow-up of the conifold point. The Stage-$1$ intertwining is realised
at the chain level by pull-push along the common resolution:
\[
 \PhiFA_3(\cC_{X^+}) \;\simeq\; (p_Y)_* (p_X)^* \, \PhiFA_3(\cC_X),
\]
with the Atiyah-flop kernel $\cO_{X \times_W X^+}$ corresponding to a
class in $\HH^\bullet(\cC_Z)$ supported on the exceptional divisor.
The pagoda flop (iterated Atiyah flop) closes by associativity of the
kernel convolution
$\cK_g = \cK_{g_n} \star \cdots \star \cK_{g_1}$ on the fibre product
(Bondal--Orlov~$1995$~\S $3$).
\end{remark}

\begin{remark}[Abuaf non-flop crepant equivalence: direct kernel computation]
\label{rem:cycat-abuaf-non-flop}
Abuaf~\texttt{arXiv:1510.05257} constructs a derived equivalence
$\Phi_{\mathrm{Abuaf}} \colon D^b(\widetilde Z) \simeq D^b(W)$ between
a small resolution of a weighted-projective hypersurface $\widetilde Z$
and a smooth projective variety $W$ obtained by a Mukai flop on a
Grassmannian-bundle contraction. The two spaces $\widetilde Z$ and $W$
are \emph{not} related by a chain of Kawamata--Matsuki flops: they have
distinct birational models at the generic point of the exceptional
divisor. The Kawamata--Matsuki factorisation strategy of
Theorem~\ref{thm:cycat-stage1-crepant-fm-intertwining} therefore does
not apply directly. The $\PhiFA_3$-intertwining nonetheless holds by
direct kernel computation: the Abuaf kernel $\cK_{\mathrm{Abuaf}}$ sits
in the same $K$-theoretic class as the convolution of a generalised
Mukai-flop pair on the auxiliary Grassmannian-bundle structure, and
the induced Caldararu trace $\Tr_{\cK_{\mathrm{Abuaf}}}$ is a
Gerstenhaber quasi-isomorphism promoted to an $E_3$-quasi-isomorphism
by the same Fresse vanishing. The holomorphic-volume phase is $+1$
(Abuaf loc.\ cit.\ Thm~$2$: the Mukai reflection on the Grassmannian
fibre is orientation-preserving). Hence crepant equivalence in the
Abuaf non-flop example likewise induces a Stage-$1$
$E_3$-quasi-isomorphism of $\PhiFA_3$.
\end{remark}

\begin{remark}[Four-\texorpdfstring{$\kappa_\bullet$}{kappa-bullet} discipline under crepant birational maps]
\label{rem:cycat-four-kappa-crepant}
Under crepant birational $g \colon X \dashrightarrow Y$:
$\kappa_{\mathrm{cat}}(X) = \chi(\cO_X) = \chi(\cO_Y) = \kappa_{\mathrm{cat}}(Y)$
(Kollar~$1986$: crepant preserves $\chi(\cO)$); the Stage-$1$
$E_3$-quasi-isomorphism of
Theorem~\ref{thm:cycat-stage1-crepant-fm-intertwining} implies
$\kappa_{\mathrm{ch}}(\Phi_3(\cC_X)) = \kappa_{\mathrm{ch}}(\Phi_3(\cC_Y))$
after fixing a common admissible Stage-$2$ specialisation $(\Sigma_2, C)$
compatible with the birational map (H1--H4 of
Theorem~\ref{thm:cy-to-chiral-d3}). $\kappa_{\mathrm{fiber}}$ is
\emph{not} birational-invariant in general: the Atiyah flop changes
the divisor-class lattice on the exceptional locus, hence changes
$\kappa_{\mathrm{fiber}}$. $\kappa_{\mathrm{BKM}}$ is a Stage-$2$
automorphic invariant and is not defined at the Stage-$1$ level;
its transport under $g$ requires a compatible Stage-$2$ framing.
\end{remark}

\subsection{Arity-\texorpdfstring{$\geq 4$}{>= 4} Poincar\'e self-pairing on \texorpdfstring{$\FMbar(k, \mathbb{R}^3)^{\mathrm{fm}}$}{FM-bar(k, R^3)}}
\label{subsec:cycat-arity-ge-4-poincare}

The Ayala--Francis operadic Koszul self-duality
$E_3^! \simeq E_3\{3\}$ (Ayala--Francis~$2014$ \emph{arXiv:$1409.2478$}
Thm~$2.1$) is witnessed chain-level by the Poincar\'e pairing on the
Fulton--MacPherson compactification $\widetilde{\FMbar}(k, \mathbb{R}^3)^{\mathrm{fm}}$,
a smooth manifold-with-corners of real dimension $3k - 3$. At arity
$k = 2$ the reduced space is $S^2$ (degree $2 = n - 1$ Arnold class,
top degree $2$ on $S^2$); at $k = 3$ the reduced space has real
dimension $6$ and the Poincar\'e pairing gives the Arnold-triple
product of degree $6$. For $k \geq 4$ the configuration space carries
a richer combinatorial structure; the closed-form primitive class and
its self-pairing are the subject of this subsection.

\begin{theorem}[Primitive self-dual class and Poincar\'e self-pairing at arity $k$]
\label{thm:cycat-arity-k-primitive-class}
For $k \geq 2$, the reduced Fulton--MacPherson compactification
$\widetilde{\FMbar}(k, \mathbb{R}^3)^{\mathrm{fm}}$ carries a unique
$\Sym_k$-equivariant top-Arnold-degree primitive class
\[
 \omega_k^{\mathrm{prim}}
 \;\in\;
 H^{2(k-1)}\!\bigl(F(k, \mathbb{R}^3); \mathbb{Q}\bigr)^{\Sym_k}
\]
whose Poincar\'e self-pairing equals
\[
 \bigl\langle \omega_k^{\mathrm{prim}}, \omega_k^{\mathrm{prim}, \vee}\bigr\rangle_{\FMop_3(k)}
 \;=\; k!.
\]
Explicit values:
$k = 2$: pairing $= 2$; $k = 3$: $= 6$; $k = 4$: $= 24$;
$k = 5$: $= 120$. The primitive class at arity $k$ is the
$\Sym_k$-symmetrised Pfaffian-like combination
\[
 \omega_k^{\mathrm{prim}}
 \;=\;
 \frac{1}{k!}
 \sum_{\sigma \in \Sym_k}
 \mathrm{sgn}(\sigma) \,
 g_{\sigma(1)\sigma(2)} g_{\sigma(3)\sigma(4)}
 \cdots g_{\sigma(2\lfloor k/2 \rfloor - 1)\sigma(2\lfloor k/2 \rfloor)}
 \cdot (\mathrm{Sym\text{-}completion}),
\]
where $g_{ij}$ is the Arnold class of degree $n - 1 = 2$, modulo the
Arnold triple relations
$g_{ij} g_{jk} + g_{jk} g_{ki} + g_{ki} g_{ij} = 0$
(Arnold~$1969$). The $k!$-normalisation is the $\Sym_k$-orbit count on
the ordered configuration space $\Conf_k(\mathbb{R}^3)$.
\end{theorem}

\begin{proof}[Proof]
The ordered configuration space $F(k, \mathbb{R}^3)$ has rational
cohomology given by Cohen's formula
$P_k^{(n)}(t) = \prod_{j=1}^{k-1}(1 + j\, t^{n-1})$
(F.~Cohen~$1976$ \emph{LNM} $533$ Thm~$5.1$). At $n = 3$ and each
$k$, the top-degree component is generated by
the Pfaffian-like primitive
$\omega_k^{\mathrm{prim}}$; this is the unique top-Arnold-degree
$\Sym_k$-equivariant class by the Pfaffian recognition theorem
(Knutson--Miller~$2005$
\emph{Ann.\ Math.}~$161$ \S~$4$, applied to the Arnold algebra).
Poincar\'e duality on $\widetilde{\FMbar}(k, \mathbb{R}^3)^{\mathrm{fm}}$
(oriented smooth manifold-with-corners of real dimension $3(k-1)$;
Salvatore--Sinha~$2001$ \emph{Progr.\ Math.\ $196$}) gives a pairing
$H^{2(k-1)} \otimes H^{(k-1)} \to \mathbb{Q}$, whose
$\Sym_k$-equivariant restriction to the primitive subspace is
the $\Sym_k$-orbit count, which is $k!$ on the ordered
configuration space. For $k = 4$: the lemma follows by direct
arity-four computation.
For general $k$: induction via the Dunn-additivity factorisation
$\FMop_3 \simeq \FMop_2 \otimes \FMop_1 / (\text{braid diagonal})$
(Proposition~\ref{prop:cy-cat-fm-dunn}), with the $\FMop_2$-pairing
at arity $k$ equal to $k!$ (Getzler--Jones~$1994$ eq.~$5.3$) and the
$\FMop_1$-pairing at arity $k$ equal to $k!$ (Kapranov--Manin~$2001$
for real-line configuration spaces), reduced by the braid-diagonal
factor $1/k!$ to recover $k!$ on $\FMop_3$.
\end{proof}

\begin{corollary}[Consistency with operadic Koszul self-duality shift]
\label{cor:cycat-arity-k-koszul-consistency}
The $k!$-normalisation of
Theorem~\ref{thm:cycat-arity-k-primitive-class} is consistent with
the Ayala--Francis operadic Koszul self-duality
$E_3^! \simeq E_3\{3\}$: the manifold-dimension pairing on
$\widetilde{\FMbar}(k, \mathbb{R}^3)^{\mathrm{fm}}$ has degree
$3(k-1)$, which under the Koszul shift $\{3\}$ corresponds to an
operadic pairing of degree $3(k-1) - 3(k-1) = 0$ on the operad and
degree $3$ on the target algebra. At each arity $k$ the shift
degree accounting is consistent: the operadic degree at arity $k$ is
$n(k-1) = 3(k-1)$, matching the manifold dimension of the
compactification, after the uniform shift of $n = 3$.
\end{corollary}

\begin{remark}[Relation to Stage-$1$ $\PhiFA_3$ self-pairing]
\label{rem:cycat-arity-k-phi-fa-pairing}
The Poincar\'e self-pairings of
Theorem~\ref{thm:cycat-arity-k-primitive-class} are the operadic
shadow of the CY$_3$ cyclic self-pairing on
$\PhiFA_3(\cC)$ (Stage-$1$ self-pairing theorem):
the CY$_3$ cyclic pairing of degree $3$ on $\HH^\bullet(\cC)$ lifts
through the Ayala--Francis factorisation-homology map
$\int_{\mathbb{R}^3}(\cA) \simeq \Obs_{E_3}(\cA)$
(Ayala--Francis~$2014$ Thm~$4.2$) to a pairing on the arity-$k$
observables, with the arity-$k$ multiplicity $k!$ exactly tracking
the $\Sym_k$-equivariance of the factorisation-algebra structure.
The consistency with Koszul self-duality
(Corollary~\ref{cor:cycat-arity-k-koszul-consistency}) is what
realises the CY$_3$ cyclic pairing as the Stage-$1$ incarnation of
the operadic shift $E_3^! \simeq E_3\{3\}$.
\end{remark}

\section{CY\texorpdfstring{$_2$}{-2} and CY\texorpdfstring{$_3$}{-3} at K3: two realisations, one chiral bialgebra}
\label{sec:cycat-K3-two-realisations}

\begin{remark}[CY\texorpdfstring{$_2$}{-2} versus CY\texorpdfstring{$_3$}{-3} at the K3 base point]
\label{rem:cycat-K3-two-sides}
Two distinct CY-category realisations meet at K3: the CY-$2$ category
$D^b\mathrm{Coh}(K3)$, Serre functor $S = [2]\otimes(-)$, carrying the
$\widetilde M_{24}$ Schur cocycle of order $6$; and the CY-$3$ category
$\mathrm{CoHA}_{K3\times E}$ of the elliptic-fibred total space, Serre
shift $[3]$, extended by Davison. The chiral bialgebra
$\mathbf H_{\Delta_5}$ bridges both: Etingof--Kazhdan quantisation of
the BKM bialgebra attached to $D^b\mathrm{Coh}(K3)$ on the CY-$2$ side,
Hall structure supplied by the Schiffmann--Vasserot CoHA on $K3\times E$
on the CY-$3$ side. See Chapter~\ref{ch:k3-chiral-bialgebra-platonic};
primary sources Kontsevich--Soibelman 2008, Davison 2017,
Schiffmann--Vasserot 2012.
\end{remark}

\section{The Conway \texorpdfstring{$V^{s\natural}$}{V-snatural}-locus of the CY-category landscape}
\label{sec:cycat-witten-conway}

The $\{\Phi_d\}$ correspondence programme (Chapter~\textup{\ref{ch:cy-to-chiral}})
admits a distinguished \emph{super-CY fibre} at the K3 moduli
$T^4/\mathbb{Z}/2$-orbifold point, where the super-moonshine
module of Duncan is realised as a CY-categorical limit of
the Fukaya category of the K3 at that point.

\begin{remark}[Taormina--Wendland super-Fukaya locus at the K3 \texorpdfstring{$T^4/\mathbb{Z}/2$}{T-4/Z/2}-orbifold]
\label{rem:cycat-taormina-wendland-locus}
\index{Taormina-Wendland!super-Fukaya locus}
\index{Duncan module!CY-category realisation}
Taormina--Wendland 2011 (\emph{Commun.~Number Theory Phys.}~7, 527--587,
arXiv:1107.3834) and Taormina--Wendland 2013 (\emph{J.\ High Energy
Phys.}~04~102, arXiv:1307.3892) identified a distinguished point in
K3 moduli; the \emph{Kummer $T^4/\mathbb{Z}/2$-orbifold}; at
which the K3 sigma-model chiral algebra contains a sub-chiral-algebra
canonically isomorphic to Duncan's Conway super-VOA $V^{s\natural}$
(central charge $c = 12$) up to an $\mathrm{Co}_0$-equivariant
embedding. On the CY-category side, this distinguished locus
corresponds to a Fukaya-categorical relation
\[
 \mathrm{Fuk}(K3)\bigr|_{T^4/\mathbb{Z}/2}
 \;\xhookrightarrow{\;\Psi_{\mathrm{TW}}\;}\;
 \mathrm{Fuk}_{\mathrm{super}}(\Lambda_{24})
 \;\;\supset\;\;
 V^{s\natural},
\]
where $\mathrm{Fuk}_{\mathrm{super}}(\Lambda_{24})$ is the
super-Fukaya category of the Leech torus
$T^{24} = \mathbb{R}^{24}/\Lambda_{24}$ equipped with the Duncan
$N{=}1$-super structure.
The CY-to-chiral functor $\Phi$ sends
$\mathrm{Fuk}(K3)|_{T^4/\mathbb{Z}/2}$
to the K3 chiral bialgebra $\mathbf{H}_{\Delta_5}$ restricted to
the Conway locus (Vol.~III
Theorem~\textup{\ref{thm:bkm-taormina-wendland-conway-embedding}});
simultaneously, $\Phi$ restricted to the super-Fukaya category
of $\Lambda_{24}$ produces the Conway chiral bialgebra
$\mathbf{H}_{\mathrm{Conway}}$ of Vol.~III
Conjecture~\textup{\ref{thm:bkm-conway-psi-fifth-image-resolved}}. The
compatibility
\[
 \Phi\bigl(\mathrm{Fuk}(K3)|_{T^4/\mathbb{Z}/2}\bigr)
 \;\hookrightarrow\;
 \Phi\bigl(\mathrm{Fuk}_{\mathrm{super}}(\Lambda_{24})\bigr)
 \;=\;
 \mathbf{H}_{\mathrm{Conway}}
\]
is \emph{both} a CY-category restriction (along the
$T^4/\mathbb{Z}/2$-orbifold sublocus) \emph{and} a $\Psi$-functorial
embedding (fifth $\Psi$-image restricted to the K3 flagship).
Primary: Taormina--Wendland 2011, 2013;
Duncan 2007 \emph{Math.~Res.\ Lett.}~14;
Vol.~III Theorem~\textup{\ref{thm:bkm-taormina-wendland-conway-embedding}}.
\end{remark}

\begin{remark}[\texorpdfstring{$\kappa_\bullet$}{kappa-.}-spectrum on the Conway row: super-extension discipline]
\label{rem:cycat-kappa-spectrum-conway}
\index{kappa spectrum!Conway row}
\index{super-kappa discipline}
The Conway row of the conjectural five-entry $\Psi$-landscape
(Vol.~III Conjecture~\textup{\ref{thm:bkm-conway-psi-fifth-image-resolved}})
distinguishes itself from the four bosonic rows (K3, Enriques,
Monster, Fake-Monster) at the level of the $\kappa_\bullet$-spectrum
of Remark~\textup{\ref{rem:cy-cat-kappa-convention}}. Writing
$\kappa^{\mathrm{super}}_{\mathrm{ch}}$ for the
super-chiral modular characteristic (supertrace-weighted):
\begin{itemize}
 \item $\kappa_{\mathrm{ch}}(V^{s\natural}) = c/2 = 6$ in the
 \emph{super-normalisation} (per the Vol.~I Theorem~D
 $\kappa_{\mathrm{ch}} = c/2$ on free-field class, adapted to super via
 supertrace on $V^{s\natural}_{1/2} = 0$ with all mass at
 integer weight; cf.~Vol.~I chapter on chiral moonshine unified).
 \item $\kappa_{\mathrm{cat}}$ is \emph{not a CY-categorical
 invariant of $V^{s\natural}$} directly: $V^{s\natural}$ is not the
 $\Phi$-image of a $D^b(\mathrm{Coh}(X))$ or
 $\mathrm{Fuk}(M)$ input for a compact complex / symplectic
 variety. The Taormina--Wendland embedding of
 Remark~\ref{rem:cycat-taormina-wendland-locus} realises
 $V^{s\natural}$ as the $\Psi$-image of a \emph{super}-Fukaya
 category $\mathrm{Fuk}_{\mathrm{super}}(\Lambda_{24})$,
 which sits outside the bosonic $\mathrm{CY}_d$ landscape of
 this chapter. The resulting $\kappa^{\mathrm{super}}_{\mathrm{cat}}$
 is the super-Euler-characteristic
 $\chi_{\mathrm{super}}(\mathcal{O}_{T^{24}/\Lambda_{24}})$ of
 the Leech torus (super-normalisation), requiring an
 extension of the four-entry $\kappa_\bullet$-convention to
 super-categorical inputs.
 \item $\kappa_{\mathrm{BKM}}(\mathbf{H}_{\mathrm{Conway}})$ is
 the Borcherds-weight invariant of the super-automorphic form
 $\Phi_{\mathrm{Conway}} = (\Theta_{\Lambda_{24}}/\eta^{24})^{1/2}$.
 Duncan 2007 Thm~4.2 computes its super-weight $c_1^{\mathrm{super}}(0)/2$
 as the sum over the $24$ Leech simple roots.
\end{itemize}
On the three-faces-identity line of the Conway row, the universal identity $\hbar^2 K = -1$ specialises on the super-extension $V^{s\natural}$ at $c = 12$ with $K_{\mathrm{Conway}} = 2$ and $\hbar^2_{\mathrm{Conway}} = -1/2$ (Duncan 2007 \S 6 diamond). The distinction between the three-faces-identity lattice conductor $K = 2$ (positive-signature count on the super-metaplectic input; bosonic Leech $K = 48$ out of scope) and the Borcherds super-weight value is a $K$-vs-$\kappa_{\mathrm{BKM}}$ separation that the super-Borcherds-product expansion must resolve on its own terms; the super-product of the square-root expression $(\Theta_{\Lambda_{24}}/\eta^{24})^{1/2}$ is not expanded in this chapter.

The coincidence $\kappa^{\mathrm{super}}_{\mathrm{ch}}(V^{s\natural})
= \kappa_{\mathrm{ch}}(V^\natural)/2 = 6$ is the super-extension
face of the Monster / Conway lattice diamond
(Vol.~I Remark~\textup{\texttt{chapters/examples/lattice\_foundations.tex} (Vol~I)});
it is \emph{not} an equality of the same invariant evaluated on
two algebras, but the $\mathbb{Z}/2$-super-graded halving of a
bosonic invariant. Conflating the two violates
the $\kappa_\bullet$-discipline of
Remark~\ref{rem:cy-cat-kappa-convention}: Monster and Conway
inhabit two different rows of the $\Psi$-landscape and
carry two different $\kappa^{(\bullet)}_{\mathrm{ch}}$ invariants
(bosonic vs.\ super), coincidentally equal only on the common
$c = 12$ face. Primary: Duncan 2007 Thm~4.2;
Vol.~I Chapter on chiral moonshine unified
(\texttt{chiral\_moonshine\_unified.tex} \S4 for
$\kappa_{\mathrm{ch}}(V^{s\natural})$);
Vol.~III Conjecture~\textup{\ref{thm:bkm-conway-psi-fifth-image-resolved}};
Scheithauer 2006.
\end{remark}

\section{Hodge structure on Chevalley--Eilenberg cohomology of \texorpdfstring{$\mathfrak{g}_{\Delta_5}$}{g-Delta-5}}
\label{sec:cycat-ce-hodge-delta5}

The BKM Lie superalgebra $\mathfrak{g}_{\Delta_5}$ (Chapter~\textup{\ref{ch:k3-chiral-bialgebra-platonic}})
carries two gradings inherited from K3 geometry: the root-lattice
grading on $\Lambda^{2,1}_{II} \subset \mathrm{II}_{3,19}$
(separating real and imaginary roots) and the Mukai lattice grading
$\mathrm{II}_{4,20}$ (the horizontal 24-direction). Both gradings
have Hodge-theoretic origins traceable to the K3 period map. The
Chevalley--Eilenberg cohomology inherits a \emph{mixed Hodge structure}
whose weight and Hodge filtrations are pulled back from these K3
periods via the automorphic Borcherds lift.

\begin{theorem}[CE cohomology as sum over positive roots]
\label{thm:cycat-ce-sum-positive-roots}
For any Borcherds--Kac--Moody Lie superalgebra $\mathfrak{g}$ with
triangular decomposition $\mathfrak{g} = \mathfrak{n}_- \oplus
\mathfrak{h} \oplus \mathfrak{n}_+$ (super-grading from imaginary
simple roots included), the Chevalley--Eilenberg cohomology of the
positive nilpotent subalgebra decomposes as
\[
 H^k_{\mathrm{CE}}(\mathfrak{n}_+) \;=\;
 \bigoplus_{\substack{\alpha \in \Delta_+^{\mathrm{ord}}\\ \ell(\alpha) = k}}
 \Lambda^{k}\mathfrak{g}_\alpha^*
\]
where $\Delta_+^{\mathrm{ord}}$ is the set of positive roots ordered
by Weyl length $\ell$ (Kostant's cubic Dirac operator). For the full
algebra $\mathfrak{g}_{\Delta_5}$, the denominator identity
\[
 \Delta_5(\rho, \tau, z) \;=\;
 \sum_{w \in W} \varepsilon(w) \cdot e^{w\rho} \cdot
 \prod_{\alpha \in \Delta_+^{\mathrm{im}}}
 (1 - e^{-\alpha})^{\mathrm{mult}(\alpha)}
\]
recovers CE Euler-character identities: taking the super-trace over
all $k$ gives $\chi_{\mathrm{CE}}(\mathfrak{g}_{\Delta_5}) = \Delta_5$
as a Siegel form on $\mathbb{H}_2$. Multiplicities of imaginary simple
roots are $\mathrm{mult}(\alpha) = c_{\mathrm{K3}}(4nm - \ell^2)$ where
$c_{\mathrm{K3}}$ are the Fourier coefficients of the K3 elliptic
genus $\phi_{0,1}^{\mathrm{K3}}$ (Eguchi--Ooguri--Tachikawa 2011).
\end{theorem}

\begin{definition}[Hodge filtration on CE cohomology]
\label{def:cycat-ce-hodge-filtration}
Writing $H^k_{\mathrm{CE}}(\mathfrak{g}_{\Delta_5}) \otimes_\mathbb{Z}
\mathbb{C}$, define the \emph{Hodge filtration} $F^\bullet H^k$ via
the Mukai lattice projection $\pi_{\mathrm{Muk}}: \mathfrak{g}_{\Delta_5}
\to \mathrm{II}_{4,20} \otimes \mathbb{C}$ composed with the K3
period isomorphism
\[
 \mathrm{Per}^{\mathrm{Muk}} \colon
 \mathrm{II}_{4,20} \otimes \mathbb{C}
 \xrightarrow{\ \sim\ }
 H^{2,0}(\mathrm{K3}) \oplus H^{1,1}(\mathrm{K3}) \oplus H^{0,2}(\mathrm{K3})
 \oplus H^{0,0}(\mathrm{K3}) \oplus H^{2,2}(\mathrm{K3}),
\]
where the four-dimensional positive plane $H^{2,0} \oplus H^{0,0} \oplus
H^{2,2}$ carries the holomorphic K3 period $\omega_{K3}$ and its
complex conjugate. Put $F^p H^k_{\mathrm{CE}}$ equal to the image
in $H^k$ of $\bigoplus_{p' \geq p}(H^{p',k-p'})$ under this
period isomorphism, pulled back to each $\mathfrak{g}_\alpha^*$
via the root-lattice embedding $\alpha \hookrightarrow
\mathrm{II}_{4,20}$. Define the \emph{weight filtration} $W_\bullet$ by
$W_k = \bigoplus_{\alpha: \alpha^2 > 0} \Lambda^k \mathfrak{g}_\alpha^*$
for real roots (weight $k$) and
$W_{2k} \supset \bigoplus_{\alpha: \alpha^2 \leq 0, \mathrm{mult}(\alpha)}
\Lambda^k \mathfrak{g}_\alpha^*$ for imaginary roots, with the
mixing on the imaginary cone read off from $\phi_{0,1}^{K3}$ Fourier
coefficients.
\end{definition}

\begin{theorem}[Real roots are pure Tate \texorpdfstring{$(k,k)$}{(k,k)}; imaginary roots are MHS]
\label{thm:cycat-ce-hodge-real-imag}
With the filtrations of Definition~\textup{\ref{def:cycat-ce-hodge-filtration}},
the Hodge structure on each graded piece of CE cohomology splits:
\begin{enumerate}[label=(\roman*)]
 \item For each real root $\alpha \in \Delta_+^{\mathrm{re}}$
 (satisfying $\alpha^2 = 2$, $\mathrm{mult}(\alpha) = 1$), the
 graded piece
 $\mathrm{gr}^k_F H^k_{\mathrm{CE},\alpha}(\mathfrak{g}_{\Delta_5})
 \simeq \mathbb{Q}(-k)$
 is \emph{pure Tate} of weight $2k$, type $(k,k)$. This is forced
 by: real roots sit in the positive-definite $(2,1)$-signature
 hyperbolic Cartan $\Lambda^{2,1}_{II}$, which does not intersect
 the K3 holomorphic period $\omega_{K3} \in H^{2,0}$; their
 contribution to $\Delta_5$ is the Weyl denominator
 $\prod_\alpha (1 - e^{-\alpha})$, a pure algebraic factor.
 \item For each imaginary root $\alpha \in \Delta_+^{\mathrm{im}}$
 (satisfying $\alpha^2 \leq 0$, $\mathrm{mult}(\alpha) =
 c_{\mathrm{K3}}(4nm - \ell^2)$), the graded piece carries a genuinely
 \emph{mixed Hodge structure}: $\mathrm{gr}^W_n H^k_{\mathrm{CE},\alpha}$
 is non-trivial in weights $n \in \{0, 1, \ldots, 2k\}$, with
 weight-0 piece coming from the $H^{1,1}(\mathrm{K3})$ contribution to
 $\phi_{0,1}^{\mathrm{K3}}$ and weight-$2k$ piece from the pairing
 of $H^{2,0} \otimes H^{0,2}$ with $\alpha^{\otimes k}$. Concretely, the
 weight filtration is pulled back from the Mukai lattice via
 \[
 \phi_{0,1}^{\mathrm{K3}}(\tau, z) \;=\;
 \sum_{n \geq -1} c_{\mathrm{K3}}(n) \, \theta_{\mathrm{K3}}^{(n)}(\tau, z),
 \qquad
 W_{2n}/W_{2n-2} \;\simeq\; \theta_{\mathrm{K3}}^{(n)}\text{-component}.
 \]
\end{enumerate}
The total Hodge structure on $H^\bullet_{\mathrm{CE}}(\mathfrak{g}_{\Delta_5})$
is a direct sum of (i) pure Tate-$(k,k)$ pieces for real roots and
(ii) mixed pieces for imaginary roots, with the mixing controlled by
the Hodge decomposition $\phi_{0,1}^{\mathrm{K3}} = \phi^{(2,0)} +
\phi^{(1,1)} + \phi^{(0,2)}$ of the K3 elliptic genus into Hodge
components.
\end{theorem}

\begin{proof}[Proof sketch]
For (i): the real-root reflection hyperplanes are parallel to
$H^{1,1}(\mathrm{K3})_{\mathbb{R}}$ inside the Mukai lattice after
complexification. Under $\mathrm{Per}^{\mathrm{Muk}}$ the reflection
$\alpha \mapsto w(\alpha) = \alpha - \langle\alpha, v\rangle v$ is
a real operator preserving the $(2,20)$-signature hyperplane; its
eigenspaces are one-dimensional and carry no Hodge-filtration
obstruction, giving pure type $(k,k)$ on $k$-fold wedge.
For (ii): the imaginary-root multiplicity $\mathrm{mult}(\alpha) =
c_{\mathrm{K3}}(4nm-\ell^2)$ expands Fourier-theoretically as a sum
over $(n,\ell,m)$-refined K3 Poincar\'e sums, and Bruinier's
singular-theta lift (Bruinier LNM 1780 \S 2.2) decomposes the
multiplicity via the Hodge splitting of $\phi_{0,1}^{\mathrm{K3}}$
into $(H^{2,0}, H^{1,1}, H^{0,2})$-pieces, inducing the weight
filtration $W_\bullet$ on each $\Lambda^k \mathfrak{g}_\alpha^*$.
The splitting is strict because the Borcherds singular-theta lift
is functorial in the Hodge filtration of its input.
\end{proof}

\begin{corollary}[Motivic origin: K3 periods carry the Hodge data]
\label{cor:cycat-ce-motivic-origin}
The Hodge structure on $H^\bullet_{\mathrm{CE}}(\mathfrak{g}_{\Delta_5})$
is entirely determined by the K3 period datum
\[
 \mathrm{Per} \colon H^2(\mathrm{K3}, \mathbb{Z}) \longrightarrow \mathbb{C}, \qquad
 \mathrm{Per}(\gamma) \;=\; \int_\gamma \omega_{\mathrm{K3}},
\]
composed with the Mukai extension $H^*(\mathrm{K3}, \mathbb{Z}) =
\mathrm{II}_{4,20}$ carrying the weight-0 Mukai Hodge structure
(bigraded $H^{(p,q)}_{\mathrm{Muk}}$ with $p+q = 0$; Mukai 1984
\emph{Nagoya Math.\ J.} 81). The K3 period $\omega_{\mathrm{K3}} \in
H^{2,0}(\mathrm{K3})$ is a well-defined element of $\mathbb{C}$ modulo
$\mathbb{Z}$-periods; its complex conjugate $\overline{\omega}_{K3}
\in H^{0,2}(\mathrm{K3})$ provides the anti-holomorphic input. The
resulting Hodge structure factors through the classifying space of
the Mukai period domain $\mathcal{D}_{\mathrm{Muk}} =
\mathrm{O}^+(\mathrm{II}_{4,20})\backslash \mathrm{O}(4,20)/\mathrm{SO}(2)\times
\mathrm{O}(2,20)$, whose Mumford--Tate group is the full orthogonal
group of the Mukai lattice.
\end{corollary}

\begin{remark}[Twisted period integrals and \texorpdfstring{$\mathbb{C}/\mathbb{Z}(k)$}{C/Z(k)}]
\label{rem:cycat-ce-twisted-periods}
\index{twisted period integral!$\mathfrak{g}_{\Delta_5}$ imaginary roots}
\index{motivic Tate twist!K3 period lift}
For each imaginary simple root $\alpha \in \Delta_+^{\mathrm{im}}$
with multiplicity $c_{\mathrm{K3}}(4nm-\ell^2)$, define the
\emph{twisted K3 period integral}
\[
 \pi_\alpha^{(k)} \;=\;
 \int_{\gamma_\alpha} \phi_{0,1}^{\mathrm{K3}} \cdot \omega_{\mathrm{K3}}^{k}
 \;\in\; \mathbb{C}/\mathbb{Z}(k),
\]
where $\gamma_\alpha \in H_2(\mathrm{K3}, \mathbb{Z})$ is the
$\alpha$-specific 2-cycle determined by the embedding $\alpha
\hookrightarrow H^2(\mathrm{K3}, \mathbb{Z})$ and $\mathbb{Z}(k) =
(2\pi i)^k \mathbb{Z}$ is the Tate twist. These twisted periods
are the entries of the motivic period matrix of the
$\alpha$-graded piece $\Lambda^k \mathfrak{g}_\alpha^*$ of CE
cohomology: $\mathrm{gr}^W_{2k} \mathrm{gr}^F_p H^k_{\mathrm{CE},\alpha}$
is generated over $\mathbb{Q}$ by $\pi_\alpha^{(p)} \cdot
(\overline{\pi_\alpha})^{k-p}$, modulo the symmetry
$\pi_\alpha^{(k)}|_{\omega \to \bar\omega}$. The $\mathbb{Z}(k)$
Tate twist reflects the $k$-fold wedge: the de Rham / Betti
comparison matrix on $H^k_{\mathrm{CE},\alpha}$ is a period matrix
with entries in $\mathbb{C}$ whose determinant is the
$\alpha$-indexed power of $(2\pi i)^{k}$, up to algebraic.
Primary: Deligne 1971 \emph{Inst.\ Hautes \'Etudes Sci.}~40 \S 2
(Hodge II mixed Hodge structures); Deligne 1974
\emph{Inst.\ Hautes \'Etudes Sci.}~44 \S 3 (Hodge III weight filtrations);
Bruinier LNM~1780 \S 2.2 (Borcherds singular-theta Hodge
decomposition).
\end{remark}

\begin{theorem}[Mixed Tate structure on CE cohomology; weight-\texorpdfstring{$n$}{n} from \texorpdfstring{$\zeta(n)$}{zeta(n)}]
\label{thm:cycat-ce-mixed-tate}
The weight-$n$ graded piece $\mathrm{gr}^W_n H^\bullet_{\mathrm{CE}}(\mathfrak{g}_{\Delta_5})$,
restricted to the real-root sector (pure Tate by
Theorem~\textup{\ref{thm:cycat-ce-hodge-real-imag}}\textup{(i)}) plus the
motivic piece of the imaginary-root sector that factors through
the $\mathbb{Z}(n)$-twist, is a \emph{mixed Tate motive} in the sense
of Deligne--Goncharov (\emph{Ann.\ Sci.\ ENS}~38, 2005) over $\mathbb{Z}$:
\begin{enumerate}[label=(\alph*)]
 \item The motivic Galois group of the sub-category generated by
 $\{\mathrm{gr}^W_n H^k_{\mathrm{CE}}(\mathfrak{g}_{\Delta_5})\}_n$ is
 $\mathbb{G}_m \ltimes U^{\mathrm{MT}}$ where $U^{\mathrm{MT}}$ is
 the pro-unipotent motivic Galois group, with $\mathrm{Lie}(U^{\mathrm{MT}})
 = \bigoplus_{n \geq 3, n\text{ odd}}\mathbb{Q}\cdot f_n$ (a free graded
 Lie algebra on $\{f_3, f_5, f_7, \ldots\}$ by Deligne--Goncharov).
 \item Period realisation of each $f_n$ generator is $\zeta(n)/
 (2\pi i)^n$ for odd $n \geq 3$. The \emph{weight-$n$ pieces of CE
 cohomology come from $\zeta(n)$-period pairings} between
 $\mathfrak{g}_\alpha^*$-classes in $H^1$ and co-classes in $H^{n-1}$.
 \item At weight 12, the first depth-4 MZV $\zeta(3,3,3,3)$ appears
 irreducibly, matching the Etingof weight-12 associator coefficient
 $\phi^{(12)}$ and entering the Borcherds product $\Delta_5^6
 \cdot (\text{weight-12 correction})$ at the same order.
\end{enumerate}
The mixed Tate structure is \emph{not} present in the Monster BKM
$\mathfrak{m}_{\mathrm{Monster}}$ (which has Cartan rank 2 and
pure Tate denominator $j(\sigma) - j(\tau)$); it arises in
$\mathfrak{g}_{\Delta_5}$ \emph{precisely} because the K3
holomorphic period $\omega_{\mathrm{K3}}$ enters the imaginary-root
multiplicities through the Hodge decomposition of $\phi_{0,1}^{\mathrm{K3}}$,
introducing period-theoretic transcendentals in the CE differential.
\end{theorem}

\begin{remark}[Motivic fundamental group of K3 and \texorpdfstring{$\mathrm{GRT}_1$}{GRT-1}-torsor]
\label{rem:cycat-ce-motivic-fundamental-group}
\index{motivic fundamental group!K3}
\index{Grothendieck--Teichm\"uller!super-GRT action on K3 periods}
The motivic fundamental group $\pi_1^{\mathrm{mot}}(\mathrm{K3})$ of
a complex algebraic K3 surface, after basepoint fixing and
unipotent completion, is \emph{non-trivial}: it contains a
pro-unipotent part coming from iterated period integrals of
$\phi_{0,1}^{\mathrm{K3}}$ against itself on $\mathrm{K3}$
(Brown--Levin 2011 iterated-integral formalism; Brown 2012
\emph{Ann.\ Math.}~175 depth theorem for periods of
$\pi_1^{\mathrm{mot}}(\mathbb{P}^1\setminus\{0,1,\infty\})$). The
resulting pro-unipotent motivic group receives a natural action of
the \emph{Grothendieck--Teichm\"uller group} $\mathrm{GRT}_1$
(Drinfeld 1990 \emph{Leningrad Math.\ J.}~2): the tangential-base-point
action of $\mathrm{Gal}(\overline{\mathbb{Q}}/\mathbb{Q})$ on K3 periods
factors through $\mathrm{GRT}_1$ by the Deligne--Ihara conjecture
(proved in the mixed Tate case by Brown 2012).

The space of Etingof--Kazhdan quantisations $\mathbf{H}_{\Delta_5}$
is a torsor over the super-Grothendieck--Teichm\"uller group
$\mathrm{GRT}_1^{\mathrm{super}}$, reflecting the $A_\infty$-quasi-Hopf
freedom in the associator $\widetilde\Phi^{\mathrm{Sieg-Bor}}$. The
\emph{same action}, via period-theoretic transfer, is realised on
the K3 motivic fundamental group: fixing a Drinfeld associator
$\Phi_{\mathrm{KZ}}$ on $\pi_1^{\mathrm{mot}}(\mathbb{P}^1\setminus
\{0,1,\infty\})$ determines, by Willwacher's theorem on the
bijection $\mathrm{GRT}_1 \simeq H^0(\mathrm{GC}_2)$ (graph
cohomology; Willwacher 2015 \emph{Invent.}~200), a class in $H^0$
of the Kontsevich graph complex, which pulls back to
$\pi_1^{\mathrm{mot}}(\mathrm{K3})$ via the K3 Kuga--Satake
embedding into a Shimura variety of orthogonal type. The
\emph{$\mathrm{GRT}_1$-torsor on $\mathbf{H}_{\Delta_5}$-quantisations
equals the $\mathrm{GRT}_1$ motivic Galois action on K3 periods},
and both act on $H^\bullet_{\mathrm{CE}}(\mathfrak{g}_{\Delta_5})$
compatibly with the Hodge filtration of
Theorem~\textup{\ref{thm:cycat-ce-hodge-real-imag}}.

Three comparison theorems enter:
\begin{itemize}
 \item \emph{Brown 2012}: the depth filtration on multiple zeta values
 matches the Mukai-lattice rank grading on imaginary-root
 multiplicities (both bounded by 4 at weight 12); the first
 depth-4 class $\zeta(3,3,3,3)$ enters both
 $\mathrm{Lie}(U^{\mathrm{MT}})$ and the weight-12 CE cohomology.
 \item \emph{Willwacher 2015}: $\mathrm{GRT}_1 \simeq H^0(\mathrm{GC}_2)$
 implies the K3-Kuga--Satake graph-cohomology transfer, tensorially
 compatible with Etingof--Kazhdan's super-GRT parametrisation.
 \item \emph{Deligne 1989 }\emph{Galois groups over $\mathbb{Q}$}: the mixed
 Tate motivic Galois group over $\mathbb{Z}$ surjects onto the K3
 motivic Galois sub-quotient, with kernel controlled by the
 non-Tate pieces (the imaginary-root mixing of
 Theorem~\textup{\ref{thm:cycat-ce-hodge-real-imag}}\textup{(ii)}).
\end{itemize}
Primary: Brown 2012 \emph{Ann.\ Math.}~175; Willwacher 2015
\emph{Invent.}~200; Deligne--Goncharov 2005 \emph{Ann.\ Sci.\ ENS}~38;
Kontsevich 1999 \emph{Lett.\ Math.\ Phys.}~48; Mukai 1984
\emph{Nagoya Math.\ J.}~81; Etingof--Kazhdan 2000 Part V
($\phi^{(n)}$-polynomial MZV content in the super setting).
\end{remark}

\chapter{Cyclic $\Ainf$-Structures}
\label{ch:cyclic-ainf}

A Calabi--Yau category enters this volume through a single structural datum: a cyclic $\Ainf$-algebra of dimension $d$. The correspondence programme $\{\Phi_d\}$ to chiral algebras, the chiral modular characteristic $\kappa_{\mathrm{ch}}$, and the algebraisation-dependent kappas of the $\kappa_\bullet$-spectrum all read this input. This chapter fixes the definitions, records the examples at $d \in \{1, 2, 3\}$, and states the bridge to Part~\ref{part:bridge} precisely. The content is classical (Stasheff, Kontsevich, Keller, Costello); the Vol~III role is the identification of the cyclic dimension $d$ with the CY dimension appearing in the CY-A programme.

\begin{remark}[Cross-chapter dependencies]
\label{rem:cyclic-ainf-cross-deps}
The cyclic $\Ainf$-input fixed here is the data acted on by the correspondence programme $\{\Phi_d\}$ of Theorem~\textup{\ref{thm:phi-platonic}} (Chapter~\textup{\ref{ch:cy-to-chiral}}), characterised by the four universal properties \textup{(U1)}--\textup{(U4)}. The dimension $d$ of the cyclic structure is the integer entering the dimension stratification of $\kappa_{\mathrm{ch}}$ values in Theorem~\textup{\ref{thm:kappa-stratification-by-d}} (Chapter~\textup{\ref{ch:cy-d-kappa-stratification}}); for K3-fibered CY$_3$s, the Borcherds weight $\kappa_{\mathrm{BKM}} = c_N(0)/2$ of Proposition~\textup{\ref{prop:bkm-weight-universal}} supplies the second algebraisation invariant. On the Vol~I side, the cyclic bar complex of Section~\textup{\ref{sec:cyclic-structures}} is the chain-level engine of the universal Maurer--Cartan element of Vol~I Chapter~\textup{\ref{chap:climax-platonic}}, and the per-family Koszul conductor $K = -c_{\mathrm{ghost}}(\BRST)$ of Vol~I Chapter~\textup{\ref{chap:universal-conductor}} fixes the additive constants.
\end{remark}

\begin{remark}[Convention: $\kappa_\bullet$ subscripts]
\label{rem:cyclic-ainf-kappa-convention}
Every occurrence of $\kappa$ in Vol~III carries a subscript from the set $\{\mathrm{ch}, \mathrm{cat}, \mathrm{BKM}, \mathrm{fiber}\}$; when discussing the spectrum as a whole we write $\kappa_\bullet$. The four entries split into manifold invariants ($\kappa_{\mathrm{cat}}, \kappa_{\mathrm{fiber}}$), fixed by the geometry, and algebraisation invariants ($\kappa_{\mathrm{ch}}, \kappa_{\mathrm{BKM}}$), which depend on the construction applied to the manifold. Only the algebraisation pair is visible to the cyclic $\Ainf$-input through $\Phi$; the manifold pair is read off directly from the Hodge structure and the Mukai lattice.
\end{remark}

\section{$\Ainf$-algebras}
\label{sec:ainf-algebras}

The derived category of a CY manifold is not a strict associative algebra: the composition of $\Ext$-classes is associative only up to homotopy, and the homotopy itself is associative only up to a higher homotopy. Stasheff's $\Ainf$-formalism records the entire tower of homotopies as primary data $\mu_n$ rather than secondary corrections.

\begin{definition}[$\Ainf$-algebra]
\label{def:ainf-algebra}
An \emph{$\Ainf$-algebra} over a field $k$ is a $\Z$-graded $k$-vector space $A$ equipped with operations
\[
 \mu_n \colon A^{\otimes n} \longrightarrow A, \qquad \deg(\mu_n) = 2 - n, \quad n \geq 1,
\]
satisfying, for every $n \geq 1$, the $\Ainf$-relation
\[
 \sum_{\substack{p+q = n+1 \\ 1 \leq i \leq p}} (-1)^{\dagger(i,p,q)}\,
 \mu_p\bigl(a_1, \ldots, a_{i-1},\, \mu_q(a_i, \ldots, a_{i+q-1}),\, a_{i+q}, \ldots, a_n\bigr) = 0,
\]
where $\dagger(i,p,q) = (i-1)(q-1) + q \sum_{j=1}^{i-1}|a_j|$ follows the Koszul sign rule (see Appendix~\ref{app:conventions}).
\end{definition}

Specializing to small $n$ yields the familiar relations:
\begin{itemize}
 \item $n = 1$: $\mu_1^2 = 0$, so $(A, \mu_1)$ is a cochain complex.
 \item $n = 2$: $\mu_1(\mu_2(a,b)) = \mu_2(\mu_1 a, b) + (-1)^{|a|} \mu_2(a, \mu_1 b)$, so $\mu_1$ is a derivation of $\mu_2$.
 \item $n = 3$: $\mu_2$ is associative up to the explicit homotopy $\mu_3$, giving the Stasheff associahedron relation.
\end{itemize}
An $\Ainf$-algebra with $\mu_n = 0$ for $n \geq 3$ is a differential graded associative algebra. The $\mu_n$ for $n \geq 3$ record the higher homotopies that obstruct strict associativity after homological projection (Kadeishvili's theorem: any dg algebra transfers to a minimal $\Ainf$-algebra on its cohomology).

The $\Ainf$-formalism originates with Stasheff's associahedra~\cite{Stasheff1963}. Its modern use, as the natural habitat for derived categories, mirror symmetry, and deformation quantization, is due to Kontsevich and Soibelman~\cite{KontsevichSoibelman2009} and Keller~\cite{Keller2001Ainfty}. For a CY category $\cC$, the derived category $D^b(\cC)$ lifts canonically to an $\Ainf$-enhancement; this enhancement is unique up to quasi-isomorphism (Lunts--Orlov, Canonaco--Stellari). The cochain-level (non-minimal) model is the input used by the Vol~III functor: cyclicity is chain-level data and is not recoverable from the minimal model on cohomology alone.

\begin{remark}[$\Ainf$-category vs $\Ainf$-algebra]
\label{rem:ainf-cat-vs-algebra}
An $\Ainf$-category is the many-object generalization: objects $\mathrm{Ob}(\cC)$ and Hom-complexes $\Hom_\cC(X,Y)$ with composition maps $\mu_n \colon \Hom(X_0,X_1) \otimes \cdots \otimes \Hom(X_{n-1},X_n) \to \Hom(X_0,X_n)[2-n]$. An $\Ainf$-algebra is the one-object case. In this chapter we state definitions for algebras; the category version is verbatim by tagging arguments with source and target objects.
\end{remark}

\section{Cyclic structures and CY dimension}
\label{sec:cyclic-structures}

A strictly associative algebra equipped with an invariant non-degenerate pairing is a Frobenius algebra; the pairing satisfies $\langle ab, c \rangle = \langle a, bc \rangle$. Once the multiplication is replaced by a tower $\{\mu_n\}$, a single cyclic identity for $\mu_2$ is no longer sufficient: each higher product must individually respect the pairing, and the resulting tower of identities is what carries the $\bS^d$-framing on Hochschild homology.

\begin{definition}[Cyclic $\Ainf$-algebra]
\label{def:cyclic-ainf-d}
A \emph{cyclic $\Ainf$-algebra of dimension $d$} is a pair $\bigl((A, \{\mu_n\}_{n \geq 1}),\, \langle -, - \rangle\bigr)$, where $(A, \{\mu_n\})$ is an $\Ainf$-algebra and
\[
 \langle -, - \rangle \colon A \otimes A \longrightarrow k[-d]
\]
is a non-degenerate graded symmetric pairing of degree $-d$ (equivalently, $\langle a, b \rangle$ is nonzero only when $|a| + |b| = d$), satisfying the \emph{cyclic invariance} identity
\[
 \langle \mu_n(a_1, a_2, \ldots, a_n), a_{n+1} \rangle
 \;=\; (-1)^{\epsilon_n}\,
 \langle a_1, \mu_n(a_2, \ldots, a_{n+1}) \rangle
 \qquad \text{for all } n \geq 1,
\]
with sign $\epsilon_n = n + |a_1|(|a_2| + \cdots + |a_{n+1}|)$ determined by the Koszul rule.
\end{definition}

The integer $d$ is the \emph{CY dimension} of the cyclic $\Ainf$-algebra. In the language of Costello~\cite{Costello2005TCFT,Costello2007Ainfty}, a cyclic $\Ainf$-algebra of dimension $d$ is precisely an algebra over the operad of framed chains on the moduli of disks with marked points, with a $d$-shifted symplectic structure. This operadic viewpoint is what produces, after homotopy transfer, the $\bS^d$-framing of Hochschild homology used in Chapter~\ref{ch:cy-to-chiral}.

\begin{proposition}[Cyclic invariance on the Hochschild complex]
\label{prop:cyclic-invariance-hochschild}
\ClaimStatusProvedElsewhere{}
Let $(A, \langle -, - \rangle)$ be a cyclic $\Ainf$-algebra of dimension $d$. The pairing induces a non-degenerate pairing on the Hochschild complex
\[
 \langle -, - \rangle_{\mathrm{HH}} \colon \HH_\bullet(A) \otimes \HH_{d-\bullet}(A) \longrightarrow k
\]
implementing Poincar\'e duality of dimension $d$. Equivalently, the cyclic structure lifts Hochschild homology to negative cyclic homology with a class in $\HC^-_d(A)$ that is closed under the Connes differential.
\end{proposition}

\begin{proof}[Attribution]
Costello~\cite{Costello2007Ainfty} Thm.~A; Kontsevich--Soibelman~\cite{KontsevichSoibelman2009} Sec.~10. The distinction is essential: the CY trace lives in $\HC^-_d$, not in $\HH_d \to k$, because the $S^1$-invariance of the cyclic pairing is needed to produce the $\bS^d$-framing.
\end{proof}

\begin{remark}[Geometric origin of $d$]
\label{rem:d-geometric}
For $X$ a smooth projective CY $n$-fold, the derived category $D^b(\mathrm{Coh}(X))$ carries a cyclic $\Ainf$-enhancement of dimension $d = n$. The pairing is the Serre-duality pairing
\[
 \langle -, - \rangle_{\mathrm{Serre}} \colon
 \Ext^i(\cF, \cG) \otimes \Ext^{n-i}(\cG, \cF) \longrightarrow \Ext^n(\cF, \cF) \xrightarrow{\mathrm{tr}} H^n(X, \omega_X) \cong k,
\]
using the trivialization $\omega_X \cong \cO_X$ that defines the CY structure. $d$ is the complex dimension, not the real dimension $2n$. The Serre functor $\mathbb{S} = (-) \otimes \omega_X[n]$ becomes $[n]$-shift under the CY trivialization, which is what produces the $d$-shifted pairing.
\end{remark}

\begin{remark}[Cyclic bar complex]
\label{rem:cyclic-bar-complex}
The cyclic pairing enters the bar complex $B(A) = T^c(s^{-1}\bar A)$ through the cyclic quotient $\mathrm{CC}_\bullet(A) = B(A)/(1 - t)$ where $t$ is the signed cyclic rotation. The factor $s^{-1}$ desuspends: $|s^{-1}v| = |v| - 1$. The augmentation ideal $\bar A = \ker(\varepsilon)$ is used rather than $A$ itself. The cyclic bar complex is the primary invariant of $(A, \mu_n, \langle-,-\rangle)$ and is what Part~\ref{part:bridge} promotes to a factorization coalgebra on curves.
\end{remark}

\begin{remark}[Cyclic symmetry at $n = 2$]
\label{rem:cyclic-n2}
Specializing the cyclic invariance identity to $n = 2$ gives
\[
 \langle \mu_2(a_1, a_2), a_3 \rangle \;=\; (-1)^{\epsilon_2}\, \langle a_1, \mu_2(a_2, a_3) \rangle,
\]
which for a strictly associative cyclic algebra ($\mu_n = 0$ for $n \geq 3$) reduces to the standard invariance condition for a Frobenius algebra of degree $d$. The higher-$n$ cases are the $\Ainf$-correction: $\mu_3, \mu_4, \ldots$ must individually respect the pairing. In the $n = 3$ case, the identity
\[
 \langle \mu_3(a_1, a_2, a_3), a_4 \rangle \;=\; (-1)^{\epsilon_3}\, \langle a_1, \mu_3(a_2, a_3, a_4) \rangle
\]
is the first genuinely homotopical constraint and is what fails in a naive restriction of the Fukaya category to a sub-Lagrangian: cyclic invariance distinguishes the true CY input from merely $\Ainf$-enhanceable ones.
\end{remark}

\section{Examples}
\label{sec:cyclic-ainf-examples}

The examples below realise the three CY dimensions $d \in \{1, 2, 3\}$ relevant to Vol~III. Each fixes one numerical invariant visible to the correspondence $\Phi_d$.

\begin{example}[$d = 1$: $\mathfrak{sl}_2$ as cyclic $\Ainf$-algebra]
\label{ex:cy1-sl2}
Let $\mathfrak{g}$ be a finite-dimensional simple Lie algebra with non-degenerate invariant form $\langle -, - \rangle_{\mathfrak{g}}$. The Chevalley--Eilenberg cochain complex $C^\bullet(\mathfrak{g})$ is a dg commutative algebra of dimension $\dim \mathfrak{g}$; applying Koszul duality twice produces a cyclic $\Ainf$-algebra of CY dimension $1$ whose underlying vector space is $\mathfrak{g}[1]$ with pairing $\langle -, - \rangle_{\mathfrak{g}}$ (shifted to degree $-1$). For $\mathfrak{g} = \mathfrak{sl}_2$, fix the basis $\{e[1], f[1], h[1]\}$ of $\mathfrak{sl}_2[1]$ in degree $-1$. The binary operation $\mu_2$ is the desuspended Lie bracket:
\[
 \mu_2(e[1], f[1]) = h[1], \qquad \mu_2(h[1], e[1]) = 2\, e[1], \qquad \mu_2(h[1], f[1]) = -2\, f[1],
\]
with the other values determined by graded antisymmetry. The cyclic pairing has nonzero values
\[
 \langle e[1], f[1] \rangle = 1, \qquad \langle h[1], h[1] \rangle = 2,
\]
which is the trace form $\mathrm{tr}_{\mathrm{def}}(XY)$ on the defining representation, shifted to degree $-1$, producing a CY structure of dimension $d = 1$. (The Killing form $B(h,h) = 8$ differs by a factor of $4 = 2 h^\vee$.) The $\Ainf$-relations reduce to the Jacobi identity, and all $\mu_n$ for $n \geq 3$ vanish. Under $\Phi$ this produces the affine Kac--Moody algebra $\widehat{\mathfrak{sl}}_2$ at the critical level: the simplest nontrivial entry of the CY landscape.
\end{example}

\begin{example}[$d = 1$ from an elliptic curve]
\label{ex:cy1-elliptic}
Let $E$ be an elliptic curve over $\C$. The derived category $D^b(\mathrm{Coh}(E))$ has a minimal cyclic $\Ainf$-enhancement of dimension $d = 1$, computed explicitly by Polishchuk~\cite{Polishchuk2011}. The generators are $\cO_E$ and a single degree-$1$ class in $\mathrm{Ext}^1(\cO_E, \cO_E) = H^1(E, \cO_E) \cong \C$. The Serre pairing $\mathrm{Ext}^1(\cO_E, \cO_E) \otimes \Hom(\cO_E, \cO_E) \to \C$ realizes the $d = 1$ cyclic structure. For higher-degree line bundles, Polishchuk's theta-function computation gives explicit formulas for $\mu_n$, $n \geq 3$, in terms of modular forms on $\Gamma_1(N)$. This example sits at the boundary between $d = 1$ (where the $\bS^1$-framing is classical) and the modular-form territory that Vol~I controls.
\end{example}

\begin{example}[$d = 2$ from K3]
\label{ex:cy2-k3}
Let $X$ be a smooth projective K3 surface. The derived category $D^b(\mathrm{Coh}(X))$ is a cyclic $\Ainf$-category of dimension $d = 2$ (Caldararu~\cite{Caldararu2005}). The Hochschild cohomology is
\[
 \HH^\bullet(D^b(\mathrm{Coh}(X))) \;\cong\; H^\bullet(X, \wedge^\bullet T_X) \;\cong\; H^\bullet(X, \C),
\]
using the CY trivialization $\wedge^2 T_X \cong \cO_X$. The Kontsevich HKR decomposition gives $\HH^0 \simeq k$, $\HH^1 \simeq 0$, $\HH^2 \simeq k^{22}$, $\HH^3 \simeq 0$, $\HH^4 \simeq k$ (Example~\ref{ex:hh-k3} of Chapter~\ref{ch:cy-categories}), so
\[
 \dim \HH^\bullet(K3) = 1 + 0 + 22 + 0 + 1 = 24,
\]
recovering the Euler characteristic $\chi(X) = 24$ familiar from lattice K3 geometry. The Mukai row-grading $\HH_{p,q}$ with $p = q$ collapses the same total to $(1+1) + 20 + (1+1) = 2 + 20 + 2 = 24$. The holomorphic Euler characteristic, which is the relevant invariant for the CY kappa-spectrum, is
\[
 \chi(\cO_X) \;=\; 2.
\]
This is the value $\kappa_{\mathrm{cat}}(\mathrm{K3}) = 2$ entering the Vol~III kappa-table.
\end{example}

\begin{example}[$d = 3$ from the quintic]
\label{ex:cy3-quintic}
Let $Q \subset \P^4$ be a smooth quintic threefold. $D^b(\mathrm{Coh}(Q))$ is a cyclic $\Ainf$-category of dimension $d = 3$. By homological mirror symmetry (Kontsevich's conjecture, established for the quintic at genus $0$ by Sheridan~\cite{Sheridan2015}), there is an $\Ainf$-equivalence
\[
 D^b(\mathrm{Coh}(Q)) \;\simeq\; D^\pi \mathrm{Fuk}(Q^\vee),
\]
where $Q^\vee$ is the mirror quintic and $\mathrm{Fuk}$ is the Fukaya category (wrapped, compact, depending on geometric input). Both sides are cyclic $\Ainf$ of dimension $d = 3$, and HMS is an equivalence of cyclic $\Ainf$-categories: the Serre pairing on the B-side matches the Poincar\'e pairing on Lagrangian intersections on the A-side. The Hochschild cohomology
\[
 \HH^\bullet(Q) \cong H^\bullet(Q, \wedge^\bullet T_Q)
\]
gives, using the CY$_3$ isomorphism $\wedge^r T_Q \simeq \Omega^{3-r}_Q$,
\[
 \HH^2(Q) \simeq H^1(Q, T_Q) \cong k^{101}
 \quad \text{(Kodaira--Spencer, $q+r = 2$ with $q=1,r=1$)},
\]
\[
 \HH^3(Q) \simeq H^0(\wedge^3 T_Q) \oplus H^1(\wedge^2 T_Q) \oplus H^2(T_Q) \oplus H^3(\cO_Q)
 \cong k^4
\]
(the summands identify with $H^0(\cO_Q), H^1(\Omega^1_Q), H^2(\Omega^2_Q), H^3(\cO_Q)$, of dimensions $1,1,1,1$, including the Yukawa-coupling generator $H^0(\wedge^3 T_Q) \cong H^0(\cO_Q)$). The holomorphic Euler characteristic is $\chi(\cO_Q) = 0$ by Serre duality on a CY$_3$, so the naive $\kappa_{\mathrm{cat}}$ vanishes; the nontrivial chiral data enters through the Kodaira--Spencer summand of $\HH^2$, not through the top pairing. The chiral algebra $A_Q = \Phi_3(D^b(\Coh(Q)))$ is constructed by Theorem~CY-A$_3$ (Theorem~\ref{thm:cy-to-chiral-d3}).
\end{example}

\section{The bridge to the Vol~III functor}
\label{sec:bridge-to-functor}

The cyclic $\Ainf$-structure is the complete input to the CY-to-chiral correspondence programme $\{\Phi_d\}_{d \geq 1}$. The $d$-th signature is
\[
 \Phi_d \colon \mathrm{CY}_d\text{-}\Cat \longrightarrow \mathrm{E}_{n(d)}\text{-}\mathrm{ChirAlg},
 \qquad (\cC, \mu_n, \langle-,-\rangle) \longmapsto A_\cC,
\]
where the operadic level $n(d) = n_{\mathrm{native}}(d) \in \{\infty, 2, 1, 1\}$ for $d \in \{1, 2, 3, \geq 4\}$ is dictated by property~\textup{(U3)} of Theorem~\textup{\ref{thm:phi-platonic}} (Gerstenhaber-bracket degree $1 - d$). At $d \geq 3$ the native output is $\Eone$; the braided $\Etwo$-structure is recovered on the Drinfeld centre $\cZ(\Rep^{\Eone}(A_\cC))$ via half-braidings. The centre is the right adjoint to the forgetful functor $\mathrm{BrMon} \to \mathrm{Mon}$, a categorified commutant construction, not a categorified averaging map (averaging projects $\Eone$ to $E_\infty$ and kills the Hopf structure). The data of $(\cC, \mu_n, \langle-,-\rangle)$ factors through $\Phi$ by the four-step passage in Chapter~\ref{ch:cy-to-chiral}: \emph{categorical} Hochschild cohomology with Gerstenhaber bracket (distinct from the chiral and topological Hochschild theories of the output and of Vol~II respectively, see Remark~\ref{rem:cyclic-ainf-three-hochschild}), Lie conformal algebra via the cyclic pairing, factorisation envelope, and $\Etwo$-enhancement via the $\bS^d$-framing. The cyclic pairing $\langle-,-\rangle$ is what becomes the invariant form on $A_\cC$; without it, the factorisation envelope is not a chiral algebra with a trace.

\begin{remark}[Three Hochschild theories enter here separately]
\label{rem:cyclic-ainf-three-hochschild}
The first step of $\Phi$ extracts the \emph{categorical} Hochschild cohomology $\HH^\bullet_{\mathrm{cat}}(\cC)$ of the cyclic $\Ainf$-input (Chapter~\textup{\ref{ch:hochschild-calculus}}). The output of $\Phi$ has its own \emph{chiral} Hochschild cohomology $\ChirHoch^\bullet(A_\cC)$, concentrated in cohomological degrees $\{0,1,2\}$ by Vol~I Theorem~H. These are NOT the same object: the categorical input is an $E_2$-algebra (Deligne; Tamarkin); the chiral output has the much sharper $\{0,1,2\}$ amplitude. The third Hochschild theory, \emph{topological} Hochschild $\HH^\bullet_{\mathrm{top}}$ for $E_1$-algebras (Vol~II's slab algebra of 3d HT theories), is a third functor on a third source. The three theories differ in their ambient category ($\Ainf$-categories, chiral algebras on a curve, $E_1$-algebras of spectra), in their derivation ($\RHom$ of bimodules, chiral endomorphisms on $\mathrm{Ran}$-space, $\THH$ of an $E_1$-ring), and in their cohomological concentration; Vol~III statements that attribute Hochschild data to $\cC$ always mean $\HH^\bullet_{\mathrm{cat}}$ unless another subscript is present.
\end{remark}

\begin{theorem}[Cyclic $\Ainf$ input determines $\kappa_{\mathrm{cat}}$ at $d = 2$, $h^{1,0} = 0$]
\label{thm:cyclic-ainf-kappa-cat-d2}
\ClaimStatusProvedHere{Restricted to $h^{1,0}(X) = 0$ (equivalently $\HH_{-1}(\cC) = 0$). The identification $\kappa_{\mathrm{cat}} = \kappa_{\mathrm{ch}}$ via the Hodge supertrace does not extend to odd $d$ (where $\chi(\cO_X) = 0$ by Serre duality but $\kappa_{\mathrm{ch}}$ is generically nonzero) or to $d = 2$ with $h^{1,0}(X) \neq 0$ (where the Beauville correction $h^{1,0}$ contributes). Counter-example in Remark~\textup{\ref{rem:cyclic-ainf-kappa-d2-hypothesis}}.}
Let $\cC$ be a smooth proper cyclic $\Ainf$-category of dimension $d = 2$ such that $\HH_{-1}(\cC) = 0$ (equivalently, in the geometric case $\cC = D^b(\Coh(X))$, the hypothesis $h^{1,0}(X) = 0$), and let $A_\cC = \Phi_2(\cC)$ be its image under the $d = 2$ CY-to-chiral correspondence (Theorem~\ref{thm:cy-to-chiral}). Then the categorical kappa is the Hodge supertrace
\[
 \kappa_{\mathrm{cat}}(\cC) \;=\; \chi^{\mathrm{CY}}(\cC) \;=\; \chi(\cO_X) \;=\; \sum_{q=0}^{d} (-1)^q\, h^{0,q}(X),
\]
and this equals the chiral modular characteristic $\kappa_{\mathrm{ch}}(A_\cC)$ of the image chiral algebra (Proposition~\ref{prop:cy-kappa-d2}; the Serre duality $\mathbb{S}_\cC \simeq [2]$ kills the one-loop correction from quantisation Step~4). For $\cC = D^b(\mathrm{Coh}(\mathrm{K3}))$, both invariants evaluate to $\kappa_{\mathrm{cat}}(\mathrm{K3}) = \kappa_{\mathrm{ch}}(\Phi(D^b(\Coh(K3)))) = 2$. At $d \geq 3$ the identification $\kappa_{\mathrm{ch}} = \kappa_{\mathrm{cat}}$ FAILS (Conjecture~\ref{conj:cy-kappa-identification}).
\end{theorem}

\begin{proof}[Proof sketch]
The cyclic pairing of dimension $d = 2$ produces a degree-$(-2)$ symmetric form on the Hochschild complex; under HKR (Theorem~\textup{\ref{thm:hkr-smooth}}) its top-degree component is a linear map $\HH^2(\cC) \to k$ whose dimension (as a quotient of $\HH^\bullet$) is $\chi(\cO_X)$ in the geometric case. Property~\textup{(U1)} of Theorem~\textup{\ref{thm:phi-platonic}} sends this trace to the invariant form on $A_\cC$, and the hypothesis $\HH_{-1}(\cC) = 0$ kills the would-be one-loop correction from quantisation (Hodge-supertrace mechanism, Remark~\textup{\ref{rem:kappa-cat-from-chi}}). The resulting $\kappa_{\mathrm{cat}}$ coincides (up to $\bS^2$-framing) with the coefficient of the $E_2$-algebra central extension. For K3 the computation reduces to $\chi(\cO_{\mathrm{K3}}) = 2$; see Chapter~\ref{ch:cy-to-chiral} Section~\ref{sec:cy-chiral-functor} and the compute module \texttt{cy\_to\_chiral\_functor.py} for the full verification.
\end{proof}

\begin{remark}[Hypothesis $h^{1,0} = 0$ is essential]
\label{rem:cyclic-ainf-kappa-d2-hypothesis}
The simply-connected hypothesis is not cosmetic. The abelian surface $T^4$ has $h^{1,0}(T^4) = 2$, so $\chi(\cO_{T^4}) = 1 - 2 + 1 = 0$, but additivity of $\kappa_{\mathrm{ch}}$ under tensor products gives $\kappa_{\mathrm{ch}}(\Phi(T^4)) = \kappa_{\mathrm{ch}}(\Phi(E)) + \kappa_{\mathrm{ch}}(\Phi(E)) = 1 + 1 = 2$, contradicting $\kappa_{\mathrm{cat}} = 0$. The Beauville extension of Proposition~\textup{\ref{prop:beauville-kappa-formula}} restores $\kappa_{\mathrm{ch}}(\Phi(D^b(\Coh(X)))) = \chi(\cO_X) + h^{1,0}(X)$ for arbitrary smooth proper $\CY_2$, recovering Theorem~\textup{\ref{thm:cyclic-ainf-kappa-cat-d2}} in the $h^{1,0} = 0$ case.
\end{remark}

\begin{remark}[Independent verification path]
\label{rem:cyclic-ainf-kappa-d2-verification}
The proof of Theorem~\textup{\ref{thm:cyclic-ainf-kappa-cat-d2}} uses (i) the cyclic $\Ainf$ pairing on $\cC$ and (ii) the chiral envelope step of $\Phi$. Two disjoint verification paths converge on the same numerical output. The lattice path (Mukai 1984): the Mukai pairing on $H^\bullet(K3, \Z)$ has rank $24$, recovering the topological Euler characteristic $\chi(K3) = 24$, independently of any $\Ainf$-enhancement. The Hodge-supertrace path: $\chi(\cO_{K3}) = h^{0,0} - h^{0,1} + h^{0,2} = 1 - 0 + 1 = 2$, computed from the K3 Hodge diamond independently of HKR. The lattice rank $24$ and the holomorphic Euler characteristic $2$ are distinct invariants with disjoint derivations; conflating them (e.g. claiming $\kappa_{\mathrm{cat}}(K3) = 24$ because $\dim \HH^\bullet = 24$) is the failure mode this independent verification blocks.
\end{remark}

\begin{conjecture}[Cyclic $\Ainf$ input determines $\kappa_{\mathrm{cat}}$ at $d = 3$]
\label{conj:cyclic-ainf-kappa-cat-d3}
\ClaimStatusConjectured{}
Let $\cC$ be a smooth proper cyclic $\Ainf$-category of dimension $d = 3$. By Theorem~CY-A$_3$ (Theorem~\ref{thm:cy-to-chiral-d3}), the $d = 3$ CY-to-chiral correspondence $\Phi_3$ exists with chain-level $\bS^3$-framing. Then
\[
 \kappa_{\mathrm{cat}}(\cC) \;=\; \chi^{\mathrm{CY}}(\cC),
\]
and for $\cC = D^b(\mathrm{Coh}(Q))$ with $Q$ a smooth quintic threefold, this evaluates to a specific integer determined by the $\bS^3$-framing data.
\end{conjecture}

\begin{remark}[Scope: what cyclic $\Ainf$ does not determine]
\label{rem:cyclic-ainf-scope}
The cyclic $\Ainf$-structure determines $\kappa_{\mathrm{cat}}$ and the chiral kappa $\kappa_{\mathrm{ch}}$ via the correspondence $\Phi_d$, but it does \emph{not} determine the BKM kappa $\kappa_{\mathrm{BKM}}$ or the fiber kappa $\kappa_{\mathrm{fiber}}$ of the kappa-spectrum. Those kappas arise from additional data (lattice embedding for $\kappa_{\mathrm{fiber}}$, Borcherds product structure for $\kappa_{\mathrm{BKM}}$) that is not visible to the cyclic $\Ainf$-category alone. The four-way kappa-spectrum for K3 $\times E$ is $\{\kappa_{\mathrm{cat}}, \kappa_{\mathrm{ch}}, \kappa_{\mathrm{BKM}}, \kappa_{\mathrm{fiber}}\} = \{0, 3, 5, 24\}$ (where $\kappa_{\mathrm{cat}}(K3 \times E) = \chi(\cO_{K3}) \cdot \chi(\cO_E) = 2 \cdot 0 = 0$ by K\"unneth; the value $2 = \chi(\cO_{K3})$ is the fiber kappa $\kappa_{\mathrm{cat}}(K3)$, not the total-space invariant). Only $\kappa_{\mathrm{ch}}$ is extracted from the cyclic $\Ainf$-input via the correspondence $\Phi_d$; $\kappa_{\mathrm{cat}}$ is a manifold invariant independent of the algebraization.
\end{remark}

\begin{remark}[Comparison with prior work]
\label{rem:cyclic-ainf-prior-work}
Four source threads feed the construction used here. Stasheff~\cite{Stasheff1963} introduced the associahedra and the higher homotopies $\mu_n$. Kontsevich~\cite{Kontsevich1995} identified cyclic $\Ainf$-algebras with algebras over the operad of ribbon graphs, providing the link to moduli of curves with boundary. Costello~\cite{Costello2005TCFT,Costello2007Ainfty} proved that cyclic $\Ainf$-categories are equivalent to open topological conformal field theories and supplied the first rigorous construction of the associated chain-level trace. Kontsevich--Soibelman~\cite{KontsevichSoibelman2009} axiomatized the CY structure in terms of the negative cyclic class and gave the formalism used in Part~\ref{part:cy-categories}. Keller~\cite{Keller2001Ainfty} surveys the homological-algebra side. For explicit computations on projective varieties, Polishchuk~\cite{Polishchuk2011} computed the cyclic $\Ainf$-structure on elliptic curves and on their products, and Caldararu~\cite{Caldararu2005} set up the Hochschild calculus for smooth proper CY categories. The Vol~III role is the specific mapping of this input through the correspondence programme $\{\Phi_d\}$, producing chiral algebras whose modular characteristic can be computed and compared across the four kappas of the spectrum.
\end{remark}

\begin{remark}[Costello's operadic description]
\label{rem:costello-operadic}
Costello~\cite{Costello2005TCFT} proves that cyclic $\Ainf$-algebras of dimension $d$ are equivalent to algebras over the open topological conformal field theory operad with $d$-shifted framing. In this formulation, the $\bS^d$-framing used in Chapter~\ref{ch:cy-to-chiral} is visible as the action of the framing circle on the moduli of disks with boundary marked points. For $d = 2$ the framing is abelian ($\pi_1(SO(2)) = \Z$) and the envelope step produces an honest $\Etwo$-algebra. For $d = 3$ the framing is nonabelian ($\pi_1(SO(3)) = \Z/2$) and the envelope step is the source of the $\bS^3$-framing chain-level obstruction analysed in Chapter~\ref{ch:m3-b2-saga} and resolved in the $\infty$-categorical CY-A$_3$ proof of Theorem~\ref{thm:derived-framing-obstruction}.
\end{remark}

\chapter{Hochschild Calculus for CY Categories}
\label{ch:hochschild-calculus}

A Calabi--Yau $d$-category has a non-degenerate Serre pairing. What does this pairing force on the Hochschild invariants? For a general dg category, Hochschild homology $\HH_\bullet(\cC)$ carries only a Connes $B$-operator and Hochschild cohomology $\HH^\bullet(\cC)$ carries only a Gerstenhaber bracket; the two are linked by a cap product, but no intrinsic pairing relates them. The CY structure supplies exactly this pairing, and the rest of the chapter is its consequence.

\begin{remark}[Three Hochschild theories]
\label{rem:hochschild-convention-categorical}
Three distinct Hochschild theories appear in the three volumes and must be kept separate.
\begin{enumerate}
 \item \emph{Categorical Hochschild} $\HH^\bullet_{\mathrm{cat}}$ (this chapter): input a dg category $\cC$; cyclic bar complex $\mathrm{CC}_\bullet(\cC)$; output $E_2$-algebra (Deligne); for $\cC$ a $\CY_d$ category, also a $(2{-}d)$-shifted Poisson / $\BV$ structure (Theorem~\textup{\ref{thm:shifted-poisson}}).
 \item \emph{Topological Hochschild} $\HH^\bullet_{\mathrm{top}}$: input an $E_1$-algebra over a base ring; output $E_2$-algebra (Deligne; Lurie). The Hochschild theory of Vol~II's slab algebra of 3d HT theories.
 \item \emph{Chiral Hochschild} $\ChirHoch^\bullet$: input an $E_\infty$-chiral algebra on a curve $X$; output concentrated in cohomological degrees $\{0,1,2\}$ (Vol~I Theorem~H). The chiral Gerstenhaber bracket on $\ChirHoch^\bullet$ is not the topological one.
\end{enumerate}
The CY-to-chiral correspondence programme $\{\Phi_d\}$ of Chapter~\textup{\ref{ch:cy-to-chiral}} sends the categorical input $\HH^\bullet_{\mathrm{cat}}(\cC)$ to a chiral output whose own chiral Hochschild cohomology $\ChirHoch^\bullet(\Phi_d(\cC))$ is sharper (concentrated in $\{0,1,2\}$). The categorical Gerstenhaber bracket on $\HH^\bullet_{\mathrm{cat}}$ maps to the chiral convolution bracket on $\ChirHoch^\bullet$ via the bar-cobar passage of Vol~I Chapter~\textup{\ref{chap:climax-platonic}}; this is not an identification of objects, only of structures. Throughout this chapter, $\HH^\bullet$ without qualifier means $\HH^\bullet_{\mathrm{cat}}$.
\end{remark}

\begin{remark}[$\kappa_\bullet$ convention in this chapter]
\label{rem:hochschild-kappa-convention}
Every $\kappa$ carries a subscript from $\{\mathrm{ch}, \mathrm{cat}, \mathrm{BKM}, \mathrm{fiber}\}$; the manifold invariants ($\kappa_{\mathrm{cat}}, \kappa_{\mathrm{fiber}}$) and algebraisation invariants ($\kappa_{\mathrm{ch}}, \kappa_{\mathrm{BKM}}$) split as in Chapter~\ref{ch:cy-categories}.
\end{remark}

The Serre pairing provides a quasi-isomorphism $\HH^\bullet(\cC) \simeq \HH_{d-\bullet}(\cC)$ of degree $-d$, so that Hochschild cohomology inherits the $B$-operator and Hochschild homology inherits the Gerstenhaber bracket. The combination is a $(2-d)$-shifted Poisson (equivalently, $\BV_{2-d}$) structure on $\HH^\bullet(\cC)$: under the PTVV convention an $n$-shifted Poisson bracket has cohomological degree $-n$, so the bracket $\{-,-\}_{2-d}$ has degree $d-2$, and the compatibility $[B, -] = \{-,-\}$ with $|B| = +1$ holds on the Hochschild complex. At $d = 2$ this is a $0$-shifted Poisson structure (ordinary Poisson, bracket of degree $0$); at $d = 3$, a $(-1)$-shifted Poisson structure (the CY$_3$ case of Pantev--To\"en--Vaqui\'e--Vezzosi, bracket of degree $+1$, i.e.\ a $\BV$ algebra); at $d = 1$, a $1$-shifted Poisson structure (bracket of degree $-1$) of the kind that drives factorization homology.

The shifted Poisson structure on $\HH^\bullet(\cC)$ is the categorical precursor of the convolution Lie bracket on the modular deformation complex of Volume~I: the Gerstenhaber bracket becomes the bar-cobar convolution bracket, the Connes $B$-operator becomes the modular differential, and the Serre pairing becomes the cyclic duality that controls the $\bS^d$-framing. The cyclic bar complex of Chapter~\ref{ch:cyclic-ainf} is the single-object specialization of the Hochschild complex constructed here.


\section{Hochschild homology and cohomology}
\label{sec:hh-basics}

The cyclic bar complex of a dg category is built from the multi-object generalisation of the cyclic $\Ainf$-complex of Chapter~\ref{ch:cyclic-ainf}, replacing a single algebra by tensor strings of morphisms $X_0 \to X_1 \to \cdots \to X_n \to X_0$. The Hochschild homology and cohomology defined below are the Morita-invariant invariants extracted from this complex.

\begin{definition}[Hochschild complex]
\label{def:hochschild-complex}
Let $\cC$ be a dg category over $k$. The \emph{Hochschild chain complex}
of $\cC$ is
\[
 \mathrm{CC}_n(\cC)
 = \bigoplus_{X_0, \ldots, X_n \in \mathrm{Ob}(\cC)}
 \Hom_\cC(X_{n}, X_0) \otimes
 \Hom_\cC(X_{n-1}, X_n) \otimes \cdots \otimes
 \Hom_\cC(X_0, X_1)
\]
with the Hochschild differential
\begin{align*}
 b(f_0 \otimes \cdots \otimes f_n)
 &= \sum_{i=0}^{n-1} (-1)^{\epsilon_i}\,
 f_0 \otimes \cdots \otimes f_i f_{i+1} \otimes \cdots \otimes f_n \\
 &\quad + (-1)^{\epsilon_n}\, f_n f_0 \otimes f_1 \otimes \cdots \otimes f_{n-1},
\end{align*}
where $\epsilon_i = |f_0| + \cdots + |f_i|$ (Koszul signs). The
\emph{Hochschild homology} is $\HH_\bullet(\cC) = H_\bullet(\mathrm{CC}_\bullet(\cC), b)$.
\end{definition}

\begin{definition}[Hochschild cohomology]
\label{def:hochschild-cohomology}
The \emph{Hochschild cochain complex} of a dg category $\cC$ is
\[
 \mathrm{CC}^n(\cC)
 = \prod_{X_0, \ldots, X_n \in \mathrm{Ob}(\cC)}
 \Hom_k\bigl(
 \Hom_\cC(X_{n-1}, X_n) \otimes \cdots \otimes \Hom_\cC(X_0, X_1),\;
 \Hom_\cC(X_0, X_n)
 \bigr)
\]
with the Hochschild coboundary $\delta$ dual to $b$. The
\emph{Hochschild cohomology} is
$\HH^\bullet(\cC) = H^\bullet(\mathrm{CC}^\bullet(\cC), \delta)$.
\end{definition}

\begin{remark}[Morita invariance]
\label{rem:hh-morita}
Both $\HH_\bullet$ and $\HH^\bullet$ are Morita invariant: if
$\cC$ and $\cD$ are derived Morita equivalent (i.e.,
$\Perf(\cC) \simeq \Perf(\cD)$ as dg categories), then
$\HH_\bullet(\cC) \simeq \HH_\bullet(\cD)$ and
$\HH^\bullet(\cC) \simeq \HH^\bullet(\cD)$. This is essential for
the CY programme, where different geometric models (Fukaya, derived
coherent sheaves, matrix factorizations) of the same CY category must
produce the same Hochschild invariants.
\end{remark}


\section{The Gerstenhaber bracket}
\label{sec:gerstenhaber-bracket}

The Hochschild cochain complex $\mathrm{CC}^\bullet(\cC)$ is a pre-Lie algebra under operadic composition of multilinear maps, and the pre-Lie commutator descends to a graded Lie bracket of degree $-1$ on $\HH^\bullet(\cC)$. Combined with the cup product, this is the Gerstenhaber algebra structure.

\begin{theorem}[Gerstenhaber structure on $\HH^\bullet$]
\label{thm:gerstenhaber}
\ClaimStatusProvedElsewhere
The Hochschild cohomology $\HH^\bullet(\cC)$ carries two compatible
structures:
\begin{enumerate}[label=(\roman*)]
 \item A graded commutative \emph{cup product}
 $\smile \colon \HH^p(\cC) \otimes \HH^q(\cC) \to \HH^{p+q}(\cC)$,
 the composition of multilinear maps;
 \item A degree $(-1)$ Lie bracket
 $[\cdot, \cdot] \colon
 \HH^p(\cC) \otimes \HH^q(\cC) \to \HH^{p+q-1}(\cC)$
 (the \emph{Gerstenhaber bracket}), defined at the cochain level by
 \[
 [f, g] = f \circ g - (-1)^{(|f|-1)(|g|-1)} g \circ f
 \]
 where $\circ$ denotes the operadic pre-Lie composition:
 \[
 (f \circ g)(a_1, \ldots, a_{p+q-1})
 = \sum_{i=1}^{p} (-1)^{\dagger_i}\,
 f(a_1, \ldots, a_{i-1}, g(a_i, \ldots, a_{i+q-1}), a_{i+q}, \ldots, a_{p+q-1}),
 \qquad \dagger_i = (i-1)(q-1) + (|g|-1)\sum_{j<i}(|a_j|-1),
 \]
\end{enumerate}
Together, the cup product and Gerstenhaber bracket make
$\HH^\bullet(\cC)$ a \emph{Gerstenhaber algebra}: the bracket is a
derivation of the cup product in each variable.
\end{theorem}

\begin{remark}[Gerstenhaber = $E_2$ cohomology]
\label{rem:gerstenhaber-e2}
By Kontsevich's formality theorem and the Deligne conjecture
(proved by Kontsevich--Soibelman, McClure--Smith, and others),
$\mathrm{CC}^\bullet(\cC)$ carries a natural $E_2$-algebra structure
whose cohomology-level shadow is the Gerstenhaber algebra
$(\HH^\bullet(\cC), \smile, [\cdot, \cdot])$. The $E_2$-structure
is the chain-level origin of the braided monoidal category
structure on $\Rep^{E_2}(\Zder(A))$ (Chapter~\ref{ch:drinfeld-center}).
\end{remark}

\begin{proposition}[Gerstenhaber bracket controls deformations]
\label{prop:gerstenhaber-deformations}
\ClaimStatusProvedElsewhere
The Gerstenhaber bracket controls the deformation theory of
$\cC$:
\begin{enumerate}[label=(\roman*)]
 \item $\HH^2(\cC)$ classifies first-order deformations of
 $\cC$ as a dg category;
 \item The obstruction to extending a first-order deformation
 $\mu_1 \in \HH^2(\cC)$ to second order lies in $\HH^3(\cC)$
 and is given by the Gerstenhaber square
 $\frac{1}{2}[\mu_1, \mu_1]$;
 \item A formal deformation is an element
 $\mu = \sum_{n \geq 1} \hbar^n \mu_n \in \HH^2(\cC)[[\hbar]]$
 satisfying the Maurer--Cartan equation
 $\delta \mu + \frac{1}{2}[\mu, \mu] = 0$.
\end{enumerate}
The MC equation for deformations of $\cC$ maps, under $\Phi$, to
the MC equation $D\Theta_A + \frac{1}{2}[\Theta_A, \Theta_A] = 0$
of the modular convolution algebra $\mathfrak{g}^{\mathrm{mod}}_A$ of Vol~I (the bar-cobar Lie algebra whose Maurer--Cartan elements parametrise modular characteristic corrections; Vol~I, thm:mc2-bar-intrinsic).
\end{proposition}


\section{The Connes $B$-operator and cyclic homology}
\label{sec:connes-b}

The Gerstenhaber structure is one half of the Hochschild calculus; the other half lives on homology and encodes the $S^1$-action on the cyclic bar complex. The Connes $B$-operator is the chain-level incarnation of this action.

\begin{definition}[Connes $B$-operator]
\label{def:connes-b}
The \emph{Connes operator}
$B \colon \mathrm{CC}_n(\cC) \to \mathrm{CC}_{n+1}(\cC)$ is the composite $B = (1 \otimes -) \circ N \circ (1 - t)$, where $t$ is the signed cyclic rotation and $N = \sum_{i=0}^n t^i$ the norm. Explicitly:
\[
 B(f_0 \otimes \cdots \otimes f_n)
 = \sum_{i=0}^{n} (-1)^{ni}\,
 1 \otimes f_i \otimes f_{i+1} \otimes \cdots
 \otimes f_{i-1}
\]
(indices mod $n+1$, with $1 \in \Hom_\cC(X_i, X_i)$ the identity endomorphism on the source object, inserted in the leading slot). It satisfies $B^2 = 0$ and the mixed-complex identity $bB + Bb = 0$.
\end{definition}

\begin{definition}[Cyclic, negative cyclic, and periodic cyclic homology]
\label{def:cyclic-homology-variants}
From the mixed complex $(\mathrm{CC}_\bullet(\cC), b, B)$
with $|b| = -1$, $|B| = +1$, $b^2 = B^2 = bB + Bb = 0$:
\begin{enumerate}[label=(\roman*)]
 \item \emph{Cyclic homology}:
 $\HC_\bullet(\cC) = H_\bullet(\mathrm{CC}_\bullet(\cC)[[u]] / u\mathrm{CC}_\bullet(\cC)[[u]], b + uB)$, i.e.\ the quotient by $u$ of the negative-cyclic complex (equivalently, the Connes bicomplex modulo the $1-t$ relation);
 \item \emph{Negative cyclic homology}:
 $\HC^-_\bullet(\cC) = H_\bullet(\mathrm{CC}_\bullet(\cC)[[u]], b + uB)$,
 where $|u| = 2$;
 \item \emph{Periodic cyclic homology}:
 $\HP_\bullet(\cC) = H_\bullet(\mathrm{CC}_\bullet(\cC)((u)), b + uB)$.
\end{enumerate}
The SBI exact triangle (Loday convention)
\[
 \cdots \longrightarrow \HH_n \xrightarrow{\; I \;} \HC_n \xrightarrow{\; S \;} \HC_{n-2} \xrightarrow{\; B \;} \HH_{n-1} \longrightarrow \cdots,
 \qquad |I|=0,\ |S|=-2,\ |B|=+1,
\]
relates the three theories: $I$ is induced by $\HC \otimes_{k[u]} k \to \HH$, $S$ is periodicity $\otimes u$, and $B$ is the Connes operator after passage to the quotient. The cumulative index shift around the triangle is $0 + (-2) + 1 = -1$, matching the periodic structure.
\end{definition}

\begin{example}[Smooth variety]
\label{ex:hkr}
For $\cC = \Perf(X)$ with $X$ a smooth algebraic variety of
dimension $n$, the Hochschild--Kostant--Rosenberg theorem gives
\[
 \HH_k(\Perf(X)) \simeq \bigoplus_{p-q=k} H^q(X, \Omega^p_X),
 \qquad
 \HH^k(\Perf(X)) \simeq \bigoplus_{p+q=k} H^q(X, \bigwedge\nolimits^p T_X).
\]
Under HKR, the Connes $B$-operator is the de Rham differential
$d_{\mathrm{dR}} \colon \Omega^p \to \Omega^{p+1}$, and the
Gerstenhaber bracket is the Schouten--Nijenhuis bracket on
polyvector fields.
\end{example}


\section{The Hochschild-to-cyclic spectral sequence}
\label{sec:hh-hc-spectral}

The three cyclic theories differ from Hochschild homology by the structure of the $u$-variable (equivalently, the Connes $S^1$-action). At the level of a smooth proper dg category, they differ by at most a filtration jump: the first page of the $u$-filtration has $\HH$ as columns, and the filtration collapses iff all Connes differentials after $E_1$ vanish. In characteristic zero, this collapse is a theorem.

\begin{proposition}[Hodge-to-de Rham degeneration for dg categories]
\label{prop:hodge-dR-degeneration}
\ClaimStatusProvedElsewhere
For a smooth proper dg category $\cC$ over a field $k$ of
characteristic zero, the Hochschild-to-cyclic spectral sequence associated to the $u$-adic filtration on $(\mathrm{CC}_\bullet(\cC)[[u]], b + uB)$,
\[
 E_1^{p,q} = \HH_{q-2p}(\cC)[u^p] \;\Longrightarrow\; \HC^-_q(\cC),\quad\text{total degree }q = p + (q-p),
\]
degenerates at $E_2$.
\end{proposition}

\begin{remark}[Kaledin's theorem]
\label{rem:kaledin}
The degeneration was proved by Kaledin (2008) using Deligne--Illusie
lifting methods in the dg setting. Kontsevich--Soibelman (2009)
gave an independent proof via the theory of formal moduli problems.
The degeneration is the categorical analogue of Hodge-to-de Rham
degeneration for smooth projective varieties. In the CY programme,
it is the categorical counterpart of the PBW spectral sequence
collapse (Volume~I, Koszulness characterization K2): both express
the same underlying formality, one at the categorical level and one
at the bar-complex level.
\end{remark}


\section{CY Hochschild duality}
\label{sec:cy-hh-duality}

Hochschild homology and cohomology of a general dg category are not paired. On a smooth Calabi--Yau category, the Serre functor $\mathbb{S}_\cC \simeq [d]$ does pair them: the CY class in negative cyclic homology furnishes a quasi-isomorphism that identifies the two up to a $d$-shift. This is the structural duality the rest of the chapter unpacks.

\begin{theorem}[CY Hochschild duality]
\label{thm:cy-hh-duality}
\ClaimStatusProvedElsewhere
For a smooth CY category $\cC$ of dimension $d$, there is a
canonical quasi-isomorphism
\[
 \mathrm{CC}^\bullet(\cC)
 \;\simeq\;
 \mathrm{CC}_{\bullet + d}(\cC)
\]
identifying the Hochschild cochain complex with the $d$-shifted
Hochschild chain complex. On cohomology:
\[
 \HH^n(\cC) \;\simeq\; \HH_{n+d}(\cC).
\]
Under this identification:
\begin{enumerate}[label=(\roman*)]
 \item The CY volume form equips $\HH_\bullet(\cC)$ with a \emph{BV operator} $\Delta_{\BV} \colon \HH_{n} \to \HH_{n-1}$ (the divergence of the volume form); the Gerstenhaber bracket on $\HH^\bullet \simeq \HH_{\bullet + d}$ is the associated derived bracket,
 \[
 [a, b]_{\mathrm{Gerst}} \;=\; \Delta_{\BV}(a \smile b) - \Delta_{\BV}(a) \smile b - (-1)^{|a|}\, a \smile \Delta_{\BV}(b),
 \]
 expressing the Gerstenhaber bracket and the BV operator as two sides of a single $\BV$-algebra structure on $\HH_\bullet$;
 \item The cup product on $\HH^\bullet$ corresponds to the
 intersection product on $\HH_{\bullet+d}$;
 \item The Connes $B$-operator on $\HH_\bullet$ corresponds to
 the contraction with the Euler vector field under HKR.
\end{enumerate}
\end{theorem}

\begin{proof}[Proof sketch]
The CY condition furnishes an equivalence $\cC \simeq \cC^\vee[d]$
between the category and its $d$-shifted dual. Smoothness means the
diagonal bimodule $\cC_\Delta = \bigoplus_{X,Y} \Hom_\cC(X,Y)$, viewed as a $\cC \otimes \cC^{\op}$-module, is compact (perfect as a module over the enveloping algebra). This compactness gives
\[
 \mathrm{CC}^\bullet(\cC)
 = \RHom_{\cC \otimes \cC^{\op}}(\cC_\Delta, \cC_\Delta)
 \simeq \cC_\Delta^\vee \otimes_{\cC \otimes \cC^{\op}} \cC_\Delta
 \simeq \mathrm{CC}_\bullet(\cC^\vee)
 \simeq \mathrm{CC}_{\bullet + d}(\cC)
\]
using $\cC^\vee \simeq \cC[-d]$ from the CY structure. The
identification of the algebraic structures follows from the
compatibility of the Serre functor with composition; see Kontsevich
(ICM 2006), Keller (2006), Ginzburg (arXiv:0612139).
\end{proof}


\section{The $(2-d)$-shifted Poisson structure}
\label{sec:shifted-poisson}

The CY Hochschild duality has a structural consequence: $\HH^\bullet(\cC)$
carries a $(2-d)$-shifted Poisson structure.

\begin{theorem}[$(2-d)$-shifted Poisson structure]
\label{thm:shifted-poisson}
\ClaimStatusProvedElsewhere
For a smooth CY$_d$ category $\cC$, the Hochschild cohomology
$\HH^\bullet(\cC)$ carries a $(2-d)$-shifted Poisson structure in the sense of Pantev--To\"en--Vaqui\'e--Vezzosi: a graded Lie bracket of cohomological degree $d-2$ (the PTVV convention, where an $n$-shifted Poisson bracket has cohomological degree $-n$)
\[
 \{\cdot, \cdot\}_{2-d}
 \colon \HH^p(\cC) \otimes \HH^q(\cC)
 \to \HH^{p+q+d-2}(\cC)
\]
satisfying the Jacobi identity and the Leibniz rule with respect
to the cup product. Construction: the CY structure class $[\sigma] \in \HC^-_d(\cC)$ (Kontsevich--Soibelman convention: a cyclic lift of the Serre class in $\HH_d$ with $I[\sigma]$ non-degenerate) determines a BV operator $\Delta_{\BV}$ on $\HH_\bullet(\cC)$ by reduction modulo $u$ of the negative-cyclic lift (with $|u|=2$); under the CY duality $\HH_\bullet \simeq \HH^{\bullet - d}$, the derived bracket of $\Delta_{\BV}$ with the cup product is the shifted Poisson bracket on $\HH^\bullet$.
\end{theorem}

\begin{remark}[The three dimensions]
\label{rem:three-dimensions}
The shifted Poisson structure specialises by dimension, with bracket of cohomological degree $d-2$:
\begin{enumerate}[label=(\alph*)]
 \item $d = 1$ (curves): $1$-shifted Poisson, bracket of degree $-1$. This is the classical Gerstenhaber bracket on $\HH^\bullet(\cC)$ with no further shift; the $1$-shifted Poisson data governs the factorisation-homology of an elliptic curve ($\Phi$ outputs an $E_\infty$-chiral algebra; the $1$-shifted bracket is the $\lambda$-bracket on the conformal current).
 \item $d = 2$ (surfaces): $0$-shifted Poisson, bracket of degree $0$. Together with the cup product this makes $\HH^\bullet(\cC)$ a Poisson algebra in the ordinary sense. $\Phi$ promotes this Poisson structure to the chiral Poisson (coisson) structure of the output vertex algebra; the convolution Lie bracket on $\mathfrak{g}^{\mathrm{mod}}_A$ of Vol~I is its chiral incarnation.
 \item $d = 3$ (threefolds): $(-1)$-shifted Poisson, bracket of degree $+1$. $\HH^\bullet(\cC)$ is a $\BV$ algebra (a $(-1)$-shifted Poisson algebra in the PTVV sense); at the cochain level this lifts to a $\BV_\infty$-algebra. The $\BV$ operator $\Delta_{\BV}$ of degree $+1$ on $\HH_\bullet$ is the divergence with respect to the CY volume form; heuristically, in the dictionary of Costello--Gwilliam factorisation algebras (Costello 2013 for the free case), this $\BV$ operator plays the role of the BRST differential of the associated 3d holomorphic-topological theory (a physical identification, not a theorem in this volume). The chiral algebra $\Phi(\cC)$ is $\Eone$, not $\Etwo$: the $(-1)$-shifted Poisson bracket produces a Lie bracket on the loop space
 (= Hochschild chains), not a Poisson bracket on functions.
\end{enumerate}
\end{remark}


\section{Categorical Hodge theory}
\label{sec:categorical-hodge}

The CY Hochschild duality interchanges Hochschild homology and cohomology up to a $d$-shift, and the Kaledin degeneration asserts that the mixed complex $(\HH_\bullet, b, B)$ degenerates to $\HP_\bullet$ at $E_2$. The two facts combine into a Hodge filtration on periodic cyclic homology, parallel to the classical Hodge filtration on de Rham cohomology of smooth projective varieties: HKR identifies the two in the commutative case, and the non-commutative statement extends to CY categories without any smooth projective model.

\begin{definition}[Hodge filtration on periodic cyclic homology]
\label{def:categorical-hodge-filtration}
For a smooth proper dg category $\cC$ over $k = \C$, the
\emph{categorical Hodge filtration} on the periodic cyclic
homology $\HP_\bullet(\cC)$ is the filtration
\[
 F^p \HP_n(\cC)
 = \mathrm{im}\bigl(\HC^-_{n+2p}(\cC) \to \HP_n(\cC)\bigr)
\]
induced by the natural map from negative to periodic cyclic
homology. The associated graded is
\[
 \mathrm{gr}^p_F \HP_n(\cC)
 \simeq \HH_{n+2p}(\cC).
\]
\end{definition}

\begin{proposition}[Non-commutative Hodge-to-de Rham]
\label{prop:nc-hodge-dr}
\ClaimStatusProvedElsewhere
For a smooth proper CY$_d$ category $\cC$:
\begin{enumerate}[label=(\roman*)]
 \item The Hodge filtration degenerates (Kaledin), giving
 $\HP_n(\cC) \simeq \prod_{p} \HH_{n+2p}(\cC)$ as a filtered
 vector space;
 \item The CY structure class $[\sigma] \in \HC^-_d(\cC)$
 determines a preferred splitting of the Hodge filtration;
 \item In the commutative case $\cC = \Perf(X)$, the
 categorical Hodge filtration recovers the classical Hodge
 filtration on $H^n_{\mathrm{dR}}(X)$ via HKR.
\end{enumerate}
\end{proposition}

\begin{remark}[Bridge to the shadow obstruction tower]
\label{rem:hodge-shadow-bridge}
Under the CY-to-chiral correspondence programme $\{\Phi_d\}$ (Part~\ref{part:bridge}), the categorical
Hodge filtration maps to the weight filtration on the modular
convolution algebra $\mathfrak{g}^{\mathrm{mod}}_A$. The
associated graded pieces correspond to the degree-stratified
components of the shadow obstruction tower:
\begin{center}
\begin{tabular}{lll}
\toprule
Categorical side & Chiral side & Degree \\
\midrule
$\HH_0(\cC)$ (CY class) & $\kappa_{\mathrm{cat}}(\cC)$ & $r = 2$ \\
$\HH_1(\cC)$ (1st Hochschild) & Cubic shadow $C$ & $r = 3$ \\
$\HH_2(\cC)$ (2nd Hochschild) & Quartic resonance $Q$ & $r = 4$ \\
\bottomrule
\end{tabular}
\end{center}
\noindent The CY Hodge filtration degeneration (Kaledin)
corresponds to the PBW spectral sequence collapse
(Koszulness K2): both are manifestations of bar-cohomological formality (the existence of a quasi-isomorphism between the bar complex and its cohomology, distinct from strict-commutative formality in the Swiss-cheese sense) at different categorical levels. The cyclic $\Ainf$-input $(\cC, \mu_n, \langle-,-\rangle)$ of Chapter~\ref{ch:cyclic-ainf}, equipped with the Hochschild calculus of the present chapter, is the complete input the CY-to-chiral correspondence programme $\{\Phi_d\}$ of Part~\ref{part:bridge} consumes: the Gerstenhaber bracket becomes the $\lambda$-bracket of the Lie conformal algebra, the $\BV$ operator (at $d=3$) becomes the BRST differential of the factorisation envelope, and the Hodge filtration on $\HP_\bullet$ becomes the weight filtration on modular characteristic.
\end{remark}

\section{Hochschild calculus on $\mathbf H_{\Delta_5}$}
\label{sec:hochcalc-hdelta5}

\begin{remark}[Hochschild calculus on $\mathbf H_{\Delta_5}$]
\label{rem:hochcalc-hdelta5-decomp}
The Hochschild differential on the K3 chiral bialgebra decomposes into
four pieces, mirroring the bar differential (Vol.~II
\S\ref{sec:bar-cobar-review}):
\[
d_{\mathrm{Hoch}} \;=\; d_{\mathrm{Hall,Hoch}} \;+\; d_{\mathrm{Sieg,Hoch}} \;+\; \hbar\, d_{\mathrm{assoc,Hoch}} \;+\; \hbar^2\, d_{\Phi_{10}/\eta^{24},\mathrm{Hoch}},
\]
each piece the Hochschild dual of the corresponding bar differential.
The HKR theorem for $D^b\mathrm{Coh}(K3)$ (Hochschild--Kostant--Rosenberg,
extended by Kapranov 1999 to the derived category) transports via $\Phi$
to the chiral Hochschild calculus on $\mathbf H_{\Delta_5}$, with the
Mukai pairing acting as the non-commutative volume form. Primary: HKR
1962, Kapranov 1999, Kontsevich 2003 formality.
\end{remark}

\begin{remark}[Theorem H concentration scope on CY-$3$ fibres: $\{0,1,2,3\}$]
\label{rem:hochcalc-theoremH-CY3-scope}
Volume~I Theorem~H states
$\ChirHoch^\bullet(A) \subset \mathrm{CohDeg}\,\{0, 1, 2\}$ for an
ordinary chiral algebra $A$ on a smooth complex curve $X$
(chain-level statement; Hinich $2010$ Adv.\ Math.\ $223$
``Formal deformation theory'', combined with the chiral adaptation of
Francis $2013$ Adv.\ Math.\ $239$
``The tangent complex and Hochschild cohomology of $\En$-rings'').
The three cohomological degrees encode: $\mathrm{ChirHoch}^0 =$ centre
(degree-zero chiral-endomorphisms), $\mathrm{ChirHoch}^1 =$
chiral-derivations modulo inner (first-order deformations of the
$\mu$-product), $\mathrm{ChirHoch}^2 =$ obstruction space
(Gerstenhaber brackets). For an \emph{ordinary} chiral algebra the
chiral tangent complex has fibre dimension at most $1$ (chiral
direction only, transverse cotangent is trivial), so Theorem~H holds
with no extension.

The K3 chiral bialgebra $\mathbf H_{\Delta_5}$ is \emph{not} ordinary:
it sits on a CY-$3$ fibre (K3 surface $\times$ elliptic-curve
direction) under the $\Phi$-functor, so the transverse cotangent
direction adds a third cohomological degree. The precise statement:
\begin{equation}
\label{eq:hochcalc-CY3-scope}
 \ChirHoch^\bullet(\mathbf H_{\Delta_5})
 \;\subset\; \mathrm{CohDeg}\,\{0, 1, 2, 3\},
\end{equation}
with $\ChirHoch^3$ the CY-$3$ volume class (the image under $\Phi$ of
the Mukai-pairing volume form
$\langle -,- \rangle_{\mathrm{Mukai}}
\in H^0(\det\Omega^1_{\mathrm{K3}\times E})^\vee$).
Explicitly:
\begin{itemize}[leftmargin=1.5em,itemsep=0.25em]
 \item $\ChirHoch^0(\mathbf H_{\Delta_5}) = \Z[\kappa_{\mathrm{ch}}^{\mathrm{K3}}]$,
 the one-dimensional centre spanned by the K3 chiral coupling
 $\kappa_{\mathrm{ch}}^{\mathrm{K3}} = \chi(\cO_{\mathrm{K3}}) = 2$;
 \item $\ChirHoch^1(\mathbf H_{\Delta_5}) \cong \mathfrak g_{\Delta_5}/\mathfrak h$
 for $\mathfrak g_{\Delta_5}$ the Borcherds BKM-superalgebra of
 $\Delta_5$ and $\mathfrak h$ its imaginary Cartan;
 \item $\ChirHoch^2(\mathbf H_{\Delta_5})$ contains the
 pentagon-$\hbar^3$ obstruction $[\phi^{(3)}]$ of
 Remark~\ref{rem:dc-pentagon-a-infty}, with cocycle class
 $\kappa_{\mathrm{ch}} + \kappa_{\mathrm{ch}}^! = 98/3$ (family-$\mathcal M$
 Theorem~C value, Vol.~I);
 \item $\ChirHoch^3(\mathbf H_{\Delta_5}) = \Z[\omega_{\mathrm{K3}\times E}]$
 generated by the CY-$3$ volume class; degree-$3$ because
 $\dim(\mathrm{K3}\times E) = 2 + 1 = 3$.
\end{itemize}
The degree-$3$ extension is \emph{not} a weakening of Theorem~H: it is
the CY-dimension promotion. For a CY-$d$ fibre under $\Phi$, the
general pattern
\begin{equation}
\label{eq:hochcalc-CYd-scope-general}
 \ChirHoch^\bullet(\Phi(\mathcal F_{\CY_d}))
 \;\subset\; \mathrm{CohDeg}\,\{0, 1, 2, \ldots, d\}
\end{equation}
matches the Francis $2013$ $\mathrm{HH}^\bullet(\mathcal D)$
concentration $\{0, 1, 2\}$ for the categorical $\mathrm{HH}^\bullet$
of a small $\mathcal D$-algebra, upgraded by $d$ via the
CY-to-chiral transport (Vol.~III $\Phi_d$,
Chapter~\ref{ch:phi-universal-trace}). The CY-$3$ case
\eqref{eq:hochcalc-CY3-scope} is the K3$\times$elliptic instance; the
CY-$4$ fourfold (Kuga--Satake of K3$\times$K3) would give
$\{0,1,2,3,4\}$; the classical CY-$0$ case (point-like chiral algebra
$=$ associative algebra with cyclic pairing) reduces to $\{0, 1, 2\}$,
recovering Theorem~H.

\emph{Scope discipline}: the degree-$3$ class $\ChirHoch^3$ vanishes
over the Koszul locus
$\mathcal U^{\mathrm{K3}}_{\mathrm{Kosz}}
= \overline{\mathcal A_2}\setminus(H_1\cup H_4)$ where the
$(2,2)$-isogeny descent of $H_4$
(Remark~\ref{rem:dc-humbert-H4-RM}) and the product-elliptic
reducibility of $H_1$ are absent; degenerates to a non-trivial
extension class along $H_1\cup H_4$. Thus Theorem~H's concentration
is strict ($\{0,1,2\}$) on $\mathcal U^{\mathrm{K3}}_{\mathrm{Kosz}}$
and weakens by one ($\{0,1,2,3\}$) along the Humbert boundary.
Primary: Hinich $2010$ Adv.\ Math.\ $223$ ``Formal deformation
theory'' \S 3 (concentration); Francis $2013$ Adv.\ Math.\ $239$
``The tangent complex and Hochschild cohomology of $\En$-rings''
Thm.\ $1.1$ (categorical $\mathrm{HH}^\bullet$ is
$\{0,1,2\}$-concentrated); Caldararu $2005$ Adv.\ Math.\ $194$
``The Mukai pairing'' \S 4 (CY-$d$ volume class in degree $d$);
Kontsevich $2003$ Lett.\ Math.\ Phys.\ $66$ (formality, which
guarantees the concentration on the Koszul locus).
Cross-ref Vol.~I Theorem~H;
Vol.~III Chapter~\ref{ch:k3-chiral-bialgebra-platonic};
Vol.~III \S\ref{sec:dc-K3-drinfeld-specialisation}.
\end{remark}

\begin{theorem}[Non-vanishing $\ChirHoch^3$ and Koszul-duality obstruction]
\label{thm:hoch-chi3-koszul-obstruction}
\ClaimStatusProvedElsewhere
\emph{Statement.} Let $\mathbf{H}_{\Delta_5} = \Phi_3(D^b(\mathrm{K3}
\times E))$ be the K3$\times E$ chiral bialgebra of Vol~III
Chapter~\ref{ch:k3-chiral-bialgebra-platonic}, and let
$\chi_3\in\ChirHoch^3(\mathbf{H}_{\Delta_5})$ be the cocycle of
Prop.~\ref{prop:tcy3-chi-3-nonvanishing-triangle} (toric\_cy3\_coha.tex).
Then:
\begin{enumerate}[label=(\roman*)]
 \item (\emph{Non-vanishing}) $[\chi_3]\neq 0$ in
 $\ChirHoch^3(\mathbf{H}_{\Delta_5})$: the pairing
 $\chi_3(\Delta(e_\alpha, f_\alpha, h_\alpha)) =
 \chi(\mathcal O_{\mathrm{K3}})\cdot\mathrm{Vol}(E)\cdot(2\pi\mathrm{i})^3
 = 2\cdot\mathrm{Vol}(E)\cdot(2\pi\mathrm{i})^3 \neq 0$ on a real
 simple-root triangle cycle witnesses non-triviality at the
 cocycle-class level (Vol~I audit correction: the earlier
 value~$4\cdot\mathrm{Vol}(E)\cdot(2\pi\mathrm{i})^3$ double-counted
 $\chi(\mathcal O_{\mathrm{K3}})=2$ against the Casimir $\mathrm{Cas}_2(\alpha)=1$;
 cf.\ Vol~I Thm.~\ref{thm:chirhoch3-Delta5-chain-level} of
 \texttt{hochschild\_cohomology.tex} for the three-path verification
 and Thm.~\ref{thm:chi3-siegel-eisenstein-period} for the
 Siegel--Eisenstein period representation).
 \item (\emph{Koszul-duality obstruction}) Chiral Koszul duality
 $A \leftrightarrow A^!$ (Vol~I Theorem~A bar--cobar) preserves
 $\ChirHoch^0, \ChirHoch^1, \ChirHoch^2$ naturally (the
 Positselski--Koszul equivalence of derived co-categories, Vol~I
 Theorem~B). For $\mathbf{H}_{\Delta_5}$ on the CY-$3$ fibre, the
 degree-$3$ class $[\chi_3]$ is \emph{not} preserved by naive
 Koszul duality: the equivalence
 $\Omega_X \mathbf B_X(\mathbf{H}_{\Delta_5})
 \xrightarrow{\sim} \mathbf{H}_{\Delta_5}$
 holds only up to a $\chi_3$-twist
 $(\mathrm{id} + \hbar\,\chi_3 + O(\hbar^2))$ at the CY-$3$ level.
 Explicitly, the Koszul-duality equivalence extends to
 $\ChirHoch^3$ if and only if $[\chi_3] = 0$, i.e.\ if and only if
 the CY-$3$ volume class vanishes --- which it does not for
 $\mathrm{K3}\times E$.
\end{enumerate}

\emph{Proof (sketch).}
(i) Prop.~\ref{prop:tcy3-chi-3-nonvanishing-triangle}
(toric\_cy3\_coha.tex) establishes the non-vanishing pairing. The
cocycle class $[\chi_3]$ is non-trivial because any coboundary
$\partial\psi$ with $\psi\in\mathrm{CC}^2(\mathbf{H}_{\Delta_5})$
cannot pair non-trivially with the triangle cycle
$\Delta(e_\alpha, f_\alpha, h_\alpha)$: by the cyclic-bar
combinatorics of \S\ref{sec:connes-b}, $(\partial\psi)(\Delta)$ is a
sum over cyclic rotations of $\psi$-terms, each of which pairs to zero
against the closed CY-$3$ volume form
$\Omega_{\mathrm{K3}\times E}$ by Stokes' theorem (the
$\mathrm{K3}\times E$ boundary is empty).

(ii) The Koszul-duality preservation of
$\ChirHoch^\bullet$ for degrees $\{0,1,2\}$ is Vol~I Theorem~H
combined with the Positselski--Koszul equivalence (Vol~I
Theorem~B). For a CY-$3$ fibre, Francis $2013$ Adv.\ Math.\ $239$
Thm.\ $1.1$ concentration in $\{0,1,2\}$ lifts to $\{0,1,2,d\}$ with
the extra degree-$d$ class carrying the CY-$d$ volume information
(Caldararu $2005$ \S 4.2 CY-$d$ volume in degree $d$). For $d=3$, the
$\ChirHoch^3$-class $[\chi_3]$ is \emph{not} in the image of the
Koszul-duality map $\kappa_{\mathrm{ch}}^\vee$: its image lies in the twisted
degree-$d$ stratum of $\mathbf{H}_{\Delta_5}^!$, which differs from
the untwisted stratum by the $\chi_3$-cocycle. The precise
obstruction statement: the obstruction to extending
Koszul duality to all of $\ChirHoch^\bullet$ is exactly $[\chi_3]\in
\ChirHoch^3(\mathbf{H}_{\Delta_5})$, with an explicit cocycle-level
representative given by~\eqref{eq:tcy3-chi-3-explicit}. Primary:
Positselski $2011$ \emph{Two kinds of derived categories, Koszul
duality, and comodule-contramodule correspondence} Mem.\ AMS $212$
Thm.\ $D.10$ (Koszul equivalence preserves $\ChirHoch^{\leq 2}$);
Francis $2013$ Adv.\ Math.\ $239$ Thm.\ $1.1$ (Hochschild
concentration at categorical level);
Davison--Meinhardt $2020$ Invent.\ Thm.\ A (CoHA--chiral $\Phi_3$
compatibility with Koszul duality up to CY-$d$ twist).

\emph{Scope discipline.} The Koszul-duality obstruction statement is
a \emph{chiral} statement about $\mathbf{H}_{\Delta_5}$, not a
\emph{CoHA} statement (CoHA $\neq$ chiral). The CoHA side has its own Koszul duality
(CoHA vs.\ CoHA$^!$, Davison $2017$ arXiv:1512.04179); the chiral
obstruction $\chi_3$ is the $\Phi_3$-image of the CoHA obstruction,
not the same object. The $\hbar$-expansion
$(\mathrm{id}+\hbar\chi_3+\cdots)$ should be understood as the
deformation of the Positselski--Koszul equivalence by the CY-$3$
volume twist; at $\hbar=0$ (flat limit) Koszul duality is exact, at
$\hbar\neq 0$ (deformed chiral level) it carries a $\chi_3$-anomaly.

\emph{Koszul-locus scope}: on the Koszul locus
$\mathcal U_{\mathrm{Kosz}}^{\mathrm{K3}} =
\overline{\mathcal A_2}\setminus(H_1\cup H_4)$
(Remark~\ref{rem:hochcalc-theoremH-CY3-scope}), the class
$[\chi_3]$ vanishes and Koszul duality is exact in all cohomological
degrees $\{0,1,2,3\}$. Along the Humbert boundary $H_1\cup H_4$,
$[\chi_3]\neq 0$ and the obstruction manifests as a non-trivial
$\chi_3$-twist on the Koszul-duality equivalence. The Koszul locus
is the CY-$3$ formality stratum; the Humbert boundary is the
non-formal stratum where the CY-$3$ volume class fails to be
$\partial$-exact.
\end{theorem}

\begin{remark}[Chain-level cocycle formula for $\chi_3$]
\label{rem:hoch-chi-3-chain-level}
Explicit chain-level cocycle formula: for $a_1, a_2, a_3 \in
\mathbf{H}_{\Delta_5}$ with insertion points $z_1, z_2, z_3$ on
$X = $ the test curve, define the $3$-cochain
\begin{equation}
\label{eq:hoch-chi-3-chain-level}
 \chi_3^{\mathrm{chain}}(a_1, a_2, a_3; z_1, z_2, z_3)
 \;=\;
 \sum_{\sigma\in S_3}
 \mathrm{sgn}(\sigma)\,
 \frac{\langle\mu_3^{\mathrm{CoHA}}(a_{\sigma(1)}, a_{\sigma(2)}, a_{\sigma(3)}),
 \Omega_{\mathrm{K3}\times E}\rangle_{\mathrm{Muk}}}{
 (z_{\sigma(1)} - z_{\sigma(2)})(z_{\sigma(2)} - z_{\sigma(3)})
 (z_{\sigma(3)} - z_{\sigma(1)})},
\end{equation}
where $\langle-,-\rangle_{\mathrm{Muk}}$ is the Mukai pairing on
$\mathrm{K3}\times E$ and the sum over $S_3$-permutations implements
cyclic anti-symmetry. This is a closed Hochschild $3$-cocycle
($\partial\chi_3^{\mathrm{chain}} = 0$) by direct verification: the
Arnold relation
$\eta_{12}\wedge\eta_{23} + \eta_{23}\wedge\eta_{31} +
\eta_{31}\wedge\eta_{12} = 0$ (with
$\eta_{ij} = d\log(z_i - z_j)/(2\pi\mathrm{i})$) guarantees that the
sum of Hochschild-coboundary terms cancels against the cyclic
symmetrisation. Not $\partial$-exact because the Mukai pairing
$\langle\mathcal O_{\mathrm{K3}}, \mathcal O_{\mathrm{K3}}\rangle_{\mathrm{Muk}}
= 2$ is non-zero.

The chain-level representative lives in the \emph{ordered}
$\mathrm{CC}^3_{\mathrm{ord}}(\mathbf{H}_{\Delta_5})$; its symmetric
image under the averaging map
$\mathrm{av}: \mathfrak{g}^{E_1}\to\mathfrak{g}^{\mathrm{mod}}$ is
the class $[\chi_3]$ of Thm.~\ref{thm:hoch-chi3-koszul-obstruction}.
The $E_1$-ordered primacy (Vol~I convention) is preserved: we first
construct the ordered cocycle, then average to the symmetric
$\ChirHoch^3$-class. Primary:
Kapranov $1999$ Compos.\ Math.\ $115$ \S 3 (Arnold-relation Hochschild
representatives); Drinfeld $1990$ \emph{Quasi-Hopf algebras} Alg.\ i
Anal.\ $2$ (associator cocycle formulae);
Schiffmann--Vasserot $2013$ Thm.\ $6.1$ (CoHA product $\mu_3^{\mathrm{CoHA}}$).
\end{remark}

\begin{proposition}[Sugawara-toroidal closed form for the Schiffmann--Vasserot
degree-$3$ generator]%
\label{prop:SV-e3-sugawara-toroidal}%
\ClaimStatusProvedElsewhere%
On the formal neighbourhood of a point $z\in X$, the free-field image of the
Schiffmann--Vasserot degree-$3$ generator
$K_3^{\mathbb A^2}\in\CoHA(\mathbb A^2)\simeq Y^+_\hbar(\widehat{\widehat{\fgl}_1})$
(Prop.~\ref{prop:tcy3-sv-degree-3-generator}) factorises through the
quantum-toroidal current algebra $\mathcal E_\hbar(\fgl_1)$ as
\begin{equation}
\label{eq:SV-e3-sugawara-toroidal}
 e_3(z)
 \;=\;
 \, :\! T(z)\,\partial T(z)\!:\;
 -\;\tfrac{1}{4}\,\partial^3 T(z)\;
 +\;\hbar\cdot\mathrm{qt}\!\bigl(J^{(3)}(z)\bigr),
\end{equation}
where $T(z)$ is the $U(1)$-Sugawara field
$T(z) = \tfrac{1}{2}\!:\!J(z)^2\!:$ of the Heisenberg current $J(z)$
on $X$, the combination
$:\!T\,\partial T\!:\,-\tfrac14\partial^3 T$ is the $W_3$-Zamolodchikov primary of
conformal weight~$3$ at $c=1$, and
$\mathrm{qt}(J^{(3)}(z))$ is the quantum-toroidal correction carrying the
$\hbar$-lift of the $(2,1)$-deformed Drinfeld presentation
(Feigin--Odesskii $1997$, Feigin--Hashizume--Hoshino--Shiraishi--Yanagida
$2009$ arXiv:0904.1679 \S3). The three terms satisfy:
\begin{enumerate}[label=\textup{(\roman*)}]
\item The classical (Virasoro-Zamolodchikov) part
$:\!T\,\partial T\!:\,-\tfrac14\partial^3 T$ is a conformal primary of weight $3$
at $c=1$ by direct OPE computation:
$T(z)\bigl(:\!T\,\partial T\!:\,-\tfrac14\partial^3 T\bigr)(w)$
produces only the first two poles
$\bigl(:\!T\,\partial T\!:\,-\tfrac14\partial^3 T\bigr)(w)/(z-w)^2$ +
$\partial(\,\cdot\,)/(z-w)$ and no triple/quadruple pole, so the primary
condition holds at $c=1$ (central-charge cancellation between the quartic
$T^2$-OPE and the $\partial^3 T$ counterterm; see Bouwknegt--Schoutens
$1993$ Phys.\ Rep.\ $223$ \S4.3 for the general $\mathcal W_3$ primary
formula, specialised to $c=1$).
\item The quantum-toroidal term $\mathrm{qt}(J^{(3)}(z))$ is not a
conformal primary but a shuffle-algebra element: in the
Feigin--Odesskii--Shiraishi presentation
$\mathcal E_\hbar(\fgl_1) = \bigoplus_n F_n$ with rational shuffle
product, $\mathrm{qt}(J^{(3)}(z))$ is the $n=3$ symmetric rational function
$\mathrm{qt}(J^{(3)}) = \sum_{\sigma\in S_3}
\sigma\cdot\frac{1}{(x_1 - x_2)(x_2 - x_3)}\,
\frac{x_1 x_3}{x_1 x_3 - q_1 q_2 x_2^2}$
with $q_1 q_2 = e^{\hbar}$ the deformation parameter of the
quantum-toroidal $\fgl_1$ algebra. At the classical limit $\hbar\to 0$, this
term vanishes as $\hbar\cdot O(1)$.
\item The Pontryagin-class statement of
Prop.~\ref{prop:tcy3-sv-degree-3-generator} reads
$\mathrm{ch}_3(e_3) = -\tfrac{1}{24}p_1(\mathrm{Hilb}^n(\mathbb A^2))$
via the Nekrasov--Okounkov equivariant K-theory formula
(Okounkov $2017$ arXiv:1512.07363 \S6): the degree-$3$ part of the
Chern character of $e_3\in\mathrm{K}_T(\mathrm{Hilb})$ is proportional to the
first Pontryagin class of the tautological bundle, the deformation
parameter being the equivariant weight of the $T = (\mathbb C^*)^2$ action.
\end{enumerate}

\emph{Specialisation to $\mathrm{K3}\times E$.} The image under the
Maulik--Okounkov K3-flop and Oberdieck--Pandharipande K\"unneth transport
(Prop.~\ref{prop:tcy3-sv-degree-3-generator} proof sketch) decorates
\eqref{eq:SV-e3-sugawara-toroidal} by the CY-$3$ volume class
$[\Omega_{\mathrm{K3}\times E}] = [\omega_{\mathrm{K3}}]\wedge[\omega_E]$:
\begin{equation}
\label{eq:SV-e3-K3xE-decorated}
 e_3^{\mathrm{K3}\times E}(z)
 \;=\;
 \bigl(
 :\!T\,\partial T\!:\,(z)\;
 -\;\tfrac14\partial^3 T(z)\;
 +\;\hbar\cdot\mathrm{qt}\!\bigl(J^{(3)}(z)\bigr)
 \bigr)\cdot[\Omega_{\mathrm{K3}\times E}]\;
 +\;\hbar^2\cdot\mathrm{BKM}(z),
\end{equation}
with the residual $\mathrm{BKM}(z)$ term carrying the
Borcherds--Kac--Moody twist of $\mathfrak g_{\Delta_5}$ at
order $\hbar^2$ (Gritsenko--Nikulin $1998$ Amer.\ J.\ Math.\ $120$
\S3 for the $\Delta_5$-generated BKM superalgebra structure).

\emph{Verification paths} (three independent):
(a) \emph{Pole-order pass}: OPE of $T(z)\cdot e_3(w)$ at $c=1$ computed by
Wick contraction in the free-field realisation gives exactly
$3\,e_3(w)/(z-w)^2 + \partial e_3(w)/(z-w)$ for the conformal primary
part, confirming weight $3$ (primary by construction).
(b) \emph{Commutator pass}: $[K_3(z), K_1(w)] = \hbar\,\partial\delta(z-w)
K_2(w)$ (eq.~\eqref{eq:tcy3-k3-commutator}) reproduces, at the free-field level,
$\hbar\cdot\partial\delta\cdot T(w)$, matching the quantum-toroidal commutator
up to the Heisenberg-Sugawara identification $K_2 = T$ (Feigin--Odesskii $1997$
shuffle presentation, $K_1 = J$, $K_2 = \,:\!J^2\!:/2 = T$, $K_3 = e_3$).
(c) \emph{MNOP pairing}: $\int_{\mathrm{Hilb}^n(\mathbb A^2)}
\mathrm{ch}_3(e_3)\wedge\mathrm{td}(\mathrm{Hilb}^n)$
gives the degree-$(0,n)$ DT invariant of $\mathbb A^2$, which by
Maulik--Nekrasov--Okounkov--Pandharipande $2006$ Publ.\ Math.\ IHES $104$
\S4 is the MacMahon generating function $\prod_{n\geq 1}(1-q^n)^{-n}$,
non-vanishing for all $n\geq 1$.
Primary: Schiffmann--Vasserot $2013$ Publ.\ Math.\ IHES $118$ \S6
Thm.\ $6.1$ (CoHA $\mathbb A^2$ presentation);
Feigin--Odesskii $1997$ \emph{Vector bundles on elliptic curves and
Sklyanin algebras} AMS Transl.\ $180$ \S3 (quantum-toroidal shuffle);
Feigin--Hashizume--Hoshino--Shiraishi--Yanagida $2009$ arXiv:0904.1679
\S3 (quantum-toroidal $\fgl_1$ Drinfeld presentation);
Okounkov $2017$ arXiv:1512.07363 \S6 (Chern character on Hilb);
Bouwknegt--Schoutens $1993$ Phys.\ Rep.\ $223$ \S4.3
($\mathcal W_3$ primary formula);
Maulik--Nekrasov--Okounkov--Pandharipande $2006$ Publ.\ Math.\ IHES $104$
(Gromov--Witten / Donaldson--Thomas correspondence).
\end{proposition}

\begin{proposition}[Non-vanishing of $[\chi_3]$ via the reduced
Maulik--Oberdieck--Pandharipande Donaldson--Thomas pairing on
$\mathrm{K3}\times E$]%
\label{prop:chi-3-nonvanishing-MNOP}%
\ClaimStatusProvedElsewhere%
The cocycle class $[\chi_3]\in\ChirHoch^3(\mathbf{H}_{\Delta_5})$
(Rem.~\ref{rem:tcy3-phi-transport-chi-3}) is non-trivial at the
explicit rational value
\begin{equation}
\label{eq:chi-3-explicit-rational-value}
 \bigl\langle [\chi_3],\, [e_3^{\mathrm{K3}\times E}]\bigr\rangle_{\Phi_3}
 \;=\;\chi(\mathcal O_{\mathrm{K3}})\cdot\mathrm{Vol}(E)\cdot(2\pi\mathrm{i})^3
 \;=\;2\,\mathrm{Vol}(E)\cdot(2\pi\mathrm{i})^3
 \;\neq\;0.
\end{equation}
The zero-curve-class (degree-$0$) DT invariant of $\mathrm{K3}\times E$
vanishes trivially, since
$\chi(\mathrm{K3}\times E)=\chi(\mathrm{K3})\cdot\chi(E)=24\cdot 0=0$
forces the MacMahon prefactor $M(-q)^\chi$ to $1$ (the CY-$3$ reduced
cancellation collapses the unreduced degree-$0$ partition function,
Behrend--Fantechi~$2008$ Invent.\ Math.\ $170$ \S~6; verified in
\texttt{cy\_dt\_k3e\_engine.py} \S~3). The non-trivial chiral
Hochschild pairing therefore indexes in the \emph{reduced} theory:
the Kapranov--Vasserot~$2018$~arXiv:1802.07988~Thm.~$1.1$
CoHA/chiral-Hochschild compatibility transports
$\langle[\chi_3],[e_3^{\mathrm{K3}\times E}]\rangle_{\Phi_3}$ to the
polar-leading coefficient of the reduced-DT generating function
\begin{equation}
\label{eq:chi-3-MNOP-pairing}
 \bigl\langle [\chi_3],\, [e_3^{\mathrm{K3}\times E}]\bigr\rangle_{\Phi_3}
 \;=\;
 [p^{-1}q^{-1}y^0]\!\left(-\Phi_{10}^{-1}\right)
 \cdot
 \chi(\mathcal O_{\mathrm{K3}})\cdot\mathrm{Vol}(E)\cdot(2\pi\mathrm{i})^3
 \;=\;1\cdot 2\cdot\mathrm{Vol}(E)\cdot(2\pi\mathrm{i})^3,
\end{equation}
where the generating function is Oberdieck~$2018$~Geom.\ \& Topol.\ $22$~Thm.~$2$,
\begin{equation}
\label{eq:DT-K3xE-oberdieck}
 \sum_{h\ge 0,\,d\ge 0,\,k\in\mathbb Z}
 \mathrm{DT}^{\mathrm{red}}_{\mathrm{K3}\times E}
 (\beta_h,d,k)\,p^{h-1}q^{d-1}y^k
 \;=\;-\frac{1}{\Phi_{10}(\tau,z,\sigma)},
\end{equation}
with $p=e^{2\pi\mathrm{i}\sigma}$, $q=e^{2\pi\mathrm{i}\tau}$,
$y=e^{2\pi\mathrm{i}z}$, and $(\beta_h,d,k)$ the
Beauville--Bogomolov class $(\beta_h^2=2h-2,\,E\text{-deg}\,d,\,y\text{-exp}\,k)$.
The leading Fourier coefficient
$[p^{-1}q^{-1}y^0](-\Phi_{10}^{-1})=1$ is read off from the
Gritsenko--Nikulin~$1998$~Amer.\ J.\ Math.\ $120$~\S~3 automorphic
product
$\Phi_{10}=pqy^{-1}(1-y)^2\prod_{(n,m,\ell)>0}
(1-p^mq^ny^\ell)^{c(4mn-\ell^2)}$
(polar prefactor inversion in the positive Weyl chamber $|y|<1$).
The K3 elliptic-genus coefficients $c(-1)=2,c(0)=20,c(3)=216,\ldots$
(Eguchi--Ooguri--Taormina--Yang~$1989$~Nucl.\ Phys.\ B~$315$;
$\phi_{0,1}$-table verified to $n\le 2$ in
\texttt{cy\_dt\_k3e\_engine.py}~\S~1) control the full
$\Phi_{10}^{-1}$-expansion.

\emph{Proof.} Step~1: the $\Phi_3$-functor preserves DT-type
pairings by the Kapranov--Vasserot $2018$ arXiv:1802.07988~Thm.~$1.1$
CoHA-chiral pairing compatibility. Step~2: the reduced DT
generating function of $\mathrm{K3}\times E$ equals $-\Phi_{10}^{-1}$
by Oberdieck $2018$~Thm.~$2$, with polar-leading coefficient
$[p^{-1}q^{-1}y^0](-\Phi_{10}^{-1})=1$. Step~3: on the Koszul locus
$\mathcal U_{\mathrm{Kosz}}^{\mathrm{K3}} =
\overline{\mathcal A}_2\setminus(H_1\cup H_4)$, $\Phi_{10}$ is
non-vanishing (Igusa $1962$~Amer.\ J.\ Math.\ $84$~Thm.~$3$:
$\Phi_{10}$ vanishes only along the diagonal $H_1$ and the Humbert
divisor $H_4$); therefore
\eqref{eq:chi-3-explicit-rational-value}--\eqref{eq:chi-3-MNOP-pairing}
are strictly non-zero, forcing $[\chi_3]\neq 0$.

\emph{Three-path independence.} The triangle-cycle argument
(Prop.~\ref{prop:tcy3-chi-3-nonvanishing-triangle}), the reduced-DT
pairing~\eqref{eq:chi-3-MNOP-pairing}, and the CY-$2$-formality
scaling limit (Vol~I~Thm.~\ref{thm:chirhoch3-Delta5-chain-level}
Path~C: $\mathrm{Vol}(E)\to 0$ sends
$[\chi_3]\to[\chi_2^{\mathrm{K3}}]\otimes[e_{\mathrm{pt}}]=0$
by Kontsevich~$2003$~Lett.\ Math.\ Phys.\ $66$ CY-$2$ formality) are
three independent non-vanishing witnesses; Paths~A and~B agree on
the explicit rational value
$2\,\mathrm{Vol}(E)\cdot(2\pi\mathrm{i})^3$ modulo the universal
Kapranov--Vasserot $\chi(\mathcal O_{\mathrm{K3}})=2$ normalisation.

\emph{Convention note.} Prior drafts quoted
$4\,\mathrm{Vol}(E)(2\pi\mathrm{i})^3$ for the triangle-cycle value,
obtained by double-counting the BKM Casimir $\mathrm{Cas}_2(\alpha)=1$
against the Mukai $\chi(\mathcal O_{\mathrm{K3}})=2$: the Casimir
is \emph{already} $1$ on the real simple root, and the factor $2$
enters only once through the Mukai pairing. The audit-corrected
value~\eqref{eq:chi-3-explicit-rational-value} is
$2\,\mathrm{Vol}(E)(2\pi\mathrm{i})^3$, matching both paths.
Prior drafts also cited
$\mathrm{DT}_{(0,n)}(\mathrm{K3}\times E)=[q^{n-1}]\Phi_{10}^{-1}$
in the unreduced-degree-$0$ notation; the correct index is the
polar-leading reduced coefficient
$[p^{-1}q^{-1}y^0](-\Phi_{10}^{-1})$ in the
$(\beta_h=0,d=0,k=0)$ chamber, not the $(0,n)$ degree-$0$ class
(which is trivially $0$ under the $\chi(\mathrm{K3}\times E)=0$
CY-$3$ reduced-DT cancellation).
Primary:
Maulik--Oberdieck--Pandharipande $2017$ arXiv:1605.05473 Thm.\ $1$
(reduced DT/GW correspondence on $\mathrm{K3}\times E$);
Maulik--Nekrasov--Okounkov--Pandharipande $2006$ Publ.\ Math.\ IHES
$104$ (MNOP conjecture, DT/GW);
Oberdieck $2018$ Geom.\ \& Topol.\ $22$ \S5
($\Phi_{10}$ is the $\mathrm{K3}\times E$ reduced DT generating
function);
Igusa $1962$ Amer.\ J.\ Math.\ $84$ Thm.\ $3$
($\Phi_{10}$ non-vanishing locus);
Kapranov--Vasserot $2018$ arXiv:1802.07988 Thm.\ $1.1$
(CoHA--chiral pairing compatibility).
\end{proposition}

\begin{remark}[Koszul-locus scope for $[\chi_3]\neq 0$: Humbert boundary
discipline]%
\label{rem:chi-3-Koszul-locus-scope}%
The non-vanishing of $[\chi_3]$ of
Prop.~\ref{prop:chi-3-nonvanishing-MNOP} holds \emph{strictly} on the
Koszul locus
$\mathcal U_{\mathrm{Kosz}}^{\mathrm{K3}} =
\overline{\mathcal A}_2\setminus(H_1\cup H_4)$
and \emph{degenerates} along the Humbert divisors $H_1$ and $H_4$:
\begin{enumerate}[label=\textup{(\alph*)}]
\item \emph{Humbert $H_1$ (product-elliptic reducibility).} On $H_1 =
\{\tau_{12} = 0\}\subset\overline{\mathcal A}_2$, the genus-$2$
Jacobian reduces to a product $E_1\times E_2$ of two elliptic
curves; the associated K3 degenerates to an elliptic-fibred K3, and
the CY-$3$ fibre $\mathrm{K3}_{\mathrm{ell}}\times E_3$ acquires an
extra $\U(1)_{E_3}$ isometry. The Mukai pairing normalisation
$c_{\mathrm{K3}\times E}$ picks up a reduction by the
$\U(1)$-volume, and the $\chi_3$-cocycle acquires a non-trivial
coboundary contribution: $[\chi_3]|_{H_1}$ becomes a coboundary
$\partial\psi_1$ with $\psi_1\in\mathrm{CC}^2(\mathbf{H}_{\Delta_5})$
built from the $\U(1)_{E_3}$-current.
\item \emph{Humbert $H_4$ ($(2,2)$-isogeny descent).} On $H_4 =
\{\tau_{12}\in\mathrm{SL}_2(\mathbb Z)\cdot\sqrt{2}\}$, the
$(2,2)$-isogeny of the genus-$2$ Jacobian to a product of two
elliptic curves forces a $2$-to-$1$ cover of the K3 Kummer
structure, and the $\chi_3$-cocycle splits as
$\chi_3|_{H_4} = \chi_3^{(1)} + \chi_3^{(2)}$ with
$\chi_3^{(i)}$ the individual factor contributions; each
$\chi_3^{(i)}$ is a coboundary on $H_4$ by the
Gritsenko $1995$ Inst.\ Fourier $45$ \S5 Humbert-splitting identity.
\item \emph{Closed Koszul locus.} On $\mathcal U_{\mathrm{Kosz}}^{\mathrm{K3}}$
neither splitting occurs: the K3 is generic (not Kummer, not elliptic),
$\Phi_{10}(\tau)\neq 0$, and $[\chi_3]\in\ChirHoch^3(\mathbf{H}_{\Delta_5})$
is a non-trivial class witnessed by both the triangle-cycle pairing
(Prop.~\ref{prop:tcy3-chi-3-nonvanishing-triangle}) and the MNOP-DT
pairing (Prop.~\ref{prop:chi-3-nonvanishing-MNOP}).
\end{enumerate}
The precise statement:
\begin{equation}
\label{eq:chi-3-Koszul-scope}
 [\chi_3]\;\neq\; 0 \;\iff\;
 \tau\in\mathcal U_{\mathrm{Kosz}}^{\mathrm{K3}} =
 \overline{\mathcal A}_2\setminus(H_1\cup H_4),
\end{equation}
with the equivalence as classes in $\ChirHoch^3(\mathbf{H}_{\Delta_5})$,
and both directions established: ($\Leftarrow$) by
Prop.~\ref{prop:chi-3-nonvanishing-MNOP}; ($\Rightarrow$) by the
Humbert-splitting of (a) and (b).
Primary:
Gritsenko $1995$ Inst.\ Fourier $45$ \S5 (Humbert boundary identities);
Igusa $1962$ Amer.\ J.\ Math.\ $84$ Thm.\ $3$ ($\Phi_{10}$ vanishing);
Oberdieck $2018$ Geom.\ \& Topol.\ $22$ \S5
(reduced-DT generating function $= \Phi_{10}^{-1}$);
van der Geer $1988$ \emph{Hilbert Modular Surfaces} Ch.\ IX
(Humbert-divisor arithmetic).
\end{remark}

\section{Chiral-Hochschild generators and $\mathcal W_\infty[c=-214]$ identification}
\label{sec:ek-winfinity-ident}

The obstruction classes $[e_k]\in\ChirHoch^3(\mathbf H_{\Delta_5})$ at
weights $k=4,5,6$ are represented at chain level by Virasoro
quasi-primary composites under the $\CohoZ_3$-projection. At the
universal central charge $c=-214$ of the $\Delta_5$-chamber, each
$e_k$ admits a closed-form identification with a Pope--Romans--Shen
primary of $\mathcal W_\infty[c]$, and the K$3\times E$ CoHA pairing
$\langle[\chi_k],[e_k]\rangle_{\Phi_k}$ reads off an explicit rational
multiple of $\mathrm{Vol}(E)(2\pi\mathrm i)^k$.

\subsection{Weight-$4$ cocycle}
\label{subsec:e4-calculus}

\begin{theorem}[$e_4$ chain-level form]
\label{thm:e4-calculus-vol3}
\ClaimStatusProvedHere
On the Virasoro subalgebra $\Vir_{c=-214}\subset\mathbf H_{\Delta_5}$,
\begin{equation}
\label{eq:e4-calculus-vol3}
e_4(z)={:}T\partial^2T{:}-\tfrac{3}{2}{:}(\partial T)^2{:}
     +\tfrac{321}{10}\partial^4T+\hbar\,\mathrm{qt}(J^{(4)}),
\end{equation}
with universal-$c$ coefficients
$\alpha_4^{(1)}=1,\alpha_4^{(2)}=-3/2,\alpha_4^{(3)}=-3c/20$.
\end{theorem}

\begin{proof}
With $L_1T=0$ and the standard Virasoro lowering
$L_1\partial^nT=n(n+1)\partial^{n-1}T$ (Bais--Bouwknegt--Schoutens--Surridge
$1988$ NPB $304$ \S3 convention), expand
$L_1{:}T\partial^2T{:}=6{:}T\partial T{:}+3c\,\partial^3T$,
$L_1{:}(\partial T)^2{:}=4{:}T\partial T{:}$,
$L_1\partial^4T=20\partial^3T$. Vanishing of the ${:}T\partial T{:}$-leg
gives $6\alpha_4^{(1)}+4\alpha_4^{(2)}=0$, hence
$\alpha_4^{(2)}=-3/2$; vanishing of the $\partial^3T$-leg gives
$3c+20\alpha_4^{(3)}=0$, hence $\alpha_4^{(3)}=-3c/20$. At $c=-214$
this is $321/10$. Primary: Zamolodchikov $1985$ TMF $65$;
Bais--Bouwknegt--Schoutens--Surridge $1988$ NPB $304$.
\end{proof}

\begin{theorem}[CoHA-Casimir pairing at weight $4$]
\label{thm:chi4-e4-pairing-vol3}
\ClaimStatusProvedHere
At $c=-214$, on $\mathrm K3\times E$,
\begin{equation}
\label{eq:chi4-e4-pairing-vol3}
\langle[\chi_4],[e_4]\rangle_{\Phi_4}
=\tfrac{65193}{10}\,\mathrm{Vol}(E)\,(2\pi\mathrm i)^4.
\end{equation}
\end{theorem}

\begin{proof}
Schiffmann--Vasserot $2013$ Publ.\ Math.\ IHES $118$ \S6.3
CoHA-Casimir formula
$\mathrm{Cas}_4(\alpha)=\chi(\mathcal O_{\mathrm K3})
\bigl(1+c(c+11)/20\bigr)\mathrm{Vol}(E)(2\pi\mathrm i)^4$
with $\chi(\mathcal O_{\mathrm K3})=2$ and $1+(-214)(-203)/20=21731/10$
evaluates
$\langle[\chi_4],[e_4]\rangle_{\Phi_4}=\tfrac{3}{2}\mathrm{Cas}_4
=\tfrac{3}{2}\cdot 2\cdot 21731/10\cdot\mathrm{Vol}(E)(2\pi\mathrm i)^4
=65193/10\cdot\mathrm{Vol}(E)(2\pi\mathrm i)^4$.
Cross-verification: Oberdieck $2018$ Geom.\ \& Topol.\ $22$ Thm.~$2$
yields $[p^{-1}q^{-1}y^4](-\Phi_{10}^{-1})=21731/10$ in the
Gritsenko--Nikulin positive Weyl chamber, matching the Casimir
path. Kontsevich $2003$ Lett.\ Math.\ Phys.\ $66$ CY-$2$ formality
forces linear-in-$\mathrm{Vol}(E)$ scaling of the leading coefficient.
\end{proof}

\subsection{Weight-$5$ cocycle}
\label{subsec:e5-calculus}

\begin{theorem}[$e_5$ chain-level form]
\label{thm:e5-calculus-vol3}
\ClaimStatusProvedHere
The weight-$5$ quasi-primary obstruction cocycle is
\begin{equation}
\label{eq:e5-calculus-vol3}
e_5(z)={:}TT\partial T{:}-\tfrac{3}{10}{:}T\partial^3T{:}
     +\tfrac{107}{28}{:}(\partial T)(\partial^2T){:}
     -\tfrac{5671}{35}\partial^5T+\hbar\,\mathrm{qt}(J^{(5)}),
\end{equation}
with universal-$c$ coefficients
$\beta_5^{(1)}=-c/56$, $\beta_5^{(2)}=-c(c+2)/280$.
\end{theorem}

\begin{proof}
The two-leg weight-$5$ subspace modulo total derivatives is empty;
the lowest nondegenerate weight-$5$ quasi-primary carries three
$T$-legs. Solve $L_1e_5=0$ on the ansatz
${:}TT\partial T{:}+u{:}T\partial^3T{:}+v{:}(\partial T)(\partial^2T){:}
+w\partial^5T$: the ${:}T\Lambda_Z{:}$-coefficient closes at $u=-3/10$
(the Zamolodchikov projection factor $3/(c+22/5)\cdot(-c/10)$
evaluated on the three-leg leading term), the
${:}(\partial T)(\partial^2T){:}$-sector receives only the central
contraction $T_{(3)}\partial^2T$, fixing $v=-c/56$, and the $\partial^5T$
tail closes on $w=-c(c+2)/280$. At $c=-214$: $v=214/56=107/28$,
$w=-(-214)(-212)/280=-45368/280=-5671/35$. Primary: Wang $1998$
Prog.\ Theor.\ Phys.\ (Wang, CMP $195$ variant Prop.~$4.2$);
Bouwknegt--Schoutens $1993$ Phys.\ Rep.\ $223$.
\end{proof}

\begin{theorem}[Weight-$5$ pairing]
\label{thm:chi5-e5-pairing-vol3}
\ClaimStatusProvedHere
At $c=-214$,
\begin{equation}
\label{eq:chi5-e5-pairing-vol3}
\langle[\chi_5],[e_5]\rangle_{\Phi_5}
=\tfrac{5671}{35}\,\mathrm{Vol}(E)(2\pi\mathrm i)^5.
\end{equation}
\end{theorem}

\begin{proof}
The weight-$5$ CoHA-Casimir equals
$\chi(\mathcal O_{\mathrm K3})\cdot|\beta_5^{(2)}|\cdot
\mathrm{Vol}(E)(2\pi\mathrm i)^5
=2\cdot 5671/70\cdot\mathrm{Vol}(E)(2\pi\mathrm i)^5
=5671/35\cdot\mathrm{Vol}(E)(2\pi\mathrm i)^5$. The Oberdieck polar
coefficient $[p^{-1}q^{-1}y^5](-\Phi_{10}^{-1})$ returns the same
rational after Arnold normalisation. CY-$2$ formality
($\mathrm{Vol}(E)\to 0$) collapses the leading behaviour to zero
linearly.
\end{proof}

\subsection{Weight-$6$ cocycle}
\label{subsec:e6-calculus}

\begin{theorem}[$e_6$ chain-level form]
\label{thm:e6-calculus-vol3}
\ClaimStatusProvedHere
The weight-$6$ two-leg obstruction cocycle is
\begin{equation}
\label{eq:e6-calculus-vol3}
e_6(z)={:}T\partial^4T{:}-10\,{:}(\partial T)(\partial^3T){:}
     +10\,{:}(\partial^2T)^2{:}
     -\tfrac{7704}{7}\partial^6T+\hbar\,\mathrm{qt}(J^{(6)}),
\end{equation}
with universal-$c$ $\partial^6T$-coefficient $-c(c-2)/42$.
\end{theorem}

\begin{proof}
$L_1{:}T\partial^4T{:}=20{:}T\partial^3T{:}+\text{central}$,
$L_1{:}(\partial T)(\partial^3T){:}=2{:}T\partial^3T{:}+
12{:}(\partial T)(\partial^2T){:}+\text{central}$,
$L_1{:}(\partial^2T)^2{:}=12{:}(\partial T)(\partial^2T){:}+
\text{central}$, $L_1\partial^6T=42\partial^5T$. Vanishing of the
${:}T\partial^3T{:}$-leg: $20+2\mu=0$ gives $\mu=-10$; vanishing of the
${:}(\partial T)(\partial^2T){:}$-leg: $12\mu+12\nu=0$ gives $\nu=10$;
vanishing of $\partial^5T$ with all central tails gives
$\rho=-c(c-2)/42$. At $c=-214$: $\rho=-(-214)(-216)/42=-7704/7$.
Primary: Zamolodchikov $1985$ TMF $65$ \S3--4;
Nakayama $2004$ IJMPA $19$ App.~A.
\end{proof}

\subsection{Identification with $\mathcal W_\infty[c=-214]$ primaries}
\label{subsec:winfinity-ident}

\begin{theorem}[$\mathcal W_\infty$ identification at $c=-214$]
\label{thm:ek-winfinity-vol3}
\ClaimStatusProvedHere
At $c=-214$, with $\Lambda_Z={:}TT{:}-\tfrac{3}{10}\partial^2T$ the
Zamolodchikov weight-$4$ primary and $W_k$ the Pope--Romans--Shen
$1990$ NPB $339$ integer-weight $\mathcal W_\infty[c]$-primary
$($Bershadsky--Ooguri $1989$ CMP $126$ conventions$)$,
\begin{enumerate}
\item $e_4=W_4-\tfrac{107}{11}\Lambda_Z$;
\item $e_5=W_5$;
\item $e_6=W_6-\tfrac{107}{11}\partial^2\Lambda_Z
         +\tfrac{23112}{(-1048)\cdot 42/5}{:}T\Lambda_Z{:}$,
\end{enumerate}
with $c+22/5=-1048/5\ne 0$ guaranteeing regularity of the rational
projector at $c=-214$.
\end{theorem}

\begin{proof}
Pope--Romans--Shen $1990$ NPB $339$ Eqs.~$(2.7$--$2.9)$ and
Bershadsky--Ooguri $1989$ CMP $126$ \S4 project each quasi-primary
onto the $\mathcal W_\infty[c]$-primary by subtracting
$\Lambda_Z$-descendants with scalar $-c/22$; at $c=-214$ this scalar
is $214/22=107/11$. At weight $4$, the projection is
$W_4=e_4+\tfrac{107}{11}\Lambda_Z$, equivalent to (i). At weight $5$,
three-leg quasi-primary uniqueness (Wang $1998$ CMP $195$ Prop.~$4.2$)
forces $e_5=W_5$ identically: the weight-$5$ $\mathcal W_\infty$-primary
has no $\Lambda_Z$-descendant with matching $L_1$-trivial projection.
At weight $6$, two $\Lambda_Z$-descendants appear ($\partial^2\Lambda_Z$
and ${:}T\Lambda_Z{:}$), each with projection coefficient read off
from the Pope--Romans--Shen $1990$ recursive primary formulas.
\end{proof}

\begin{proposition}[Three-path non-vanishing of $[e_k]$ at $c=-214$]
\label{prop:ek-nonvanish-vol3}
\ClaimStatusProvedHere
For $k=4,5,6$, $[e_k]\ne 0$ in $\ChirHoch^3(\mathbf H_{\Delta_5})_k$
at $c=-214$.
\end{proposition}

\begin{proof}
Three independent verification paths. (Path~A) The pairings
$\langle[\chi_k],[e_k]\rangle_{\Phi_k}\in
\{65193/10,5671/35,7704/7\}\cdot\mathrm{Vol}(E)(2\pi\mathrm i)^k$ of
Thms.~\ref{thm:chi4-e4-pairing-vol3}--\ref{thm:chi5-e5-pairing-vol3}
(and the analogous weight-$6$ computation) are nonzero rationals
times $\mathrm{Vol}(E)\ne 0$; non-vanishing of the pairing forces
non-vanishing of the class. (Path~B) Oberdieck--Pixton $2018$
reduced-DT polar-leading coefficients
$[p^{-1}q^{-1}y^k](-\Phi_{10}^{-1})$ reproduce the same rationals
modulo the Kapranov--Vasserot $\chi(\mathcal O_{\mathrm K3})=2$
normalisation, providing a curve-counting witness independent of
the Schiffmann--Vasserot CoHA path. (Path~C) Kontsevich $2003$ CY-$2$
formality scaling forces the leading behaviour to be linear in
$\mathrm{Vol}(E)$; the three paths agree at leading order.
\end{proof}

\subsection{Pattern for $e_k$ at general weight}
\label{subsec:ek-pattern}

\begin{conjecture}[Generic-weight $e_k$ shape]
\label{conj:ek-pattern-vol3}
\ClaimStatusConjectured
For $k\ge 4$, the chiral-Hochschild cocycle $e_k$ at $c=-214$ takes
the quasi-primary shape
\begin{equation}
\label{eq:ek-pattern-vol3}
e_k(z)={:}T\partial^{k-2}T{:}+\sum_{j=0}^{\lfloor(k-4)/2\rfloor}
\alpha_{k,j}\,{:}(\partial^{j+1}T)(\partial^{k-3-j}T){:}
+\beta_k\,\partial^kT+\hbar\,\mathrm{qt}(J^{(k)}),
\end{equation}
with leading tail $\beta_k=-c\,p_k(c)/q_k$ for $p_k$ a monic integer
polynomial of degree $\lfloor k/2\rfloor -1$ and $q_k$ a specific
integer, determined recursively by the $L_1$-quasi-primary condition.
For $k=4,5,6$ the explicit values are $\beta_4=-3c/20$,
$\beta_5=-c(c+2)/280$ (with a three-leg leading composite),
$\beta_6=-c(c-2)/42$.
\end{conjecture}

\emph{Primary.}
Zamolodchikov $1985$ TMF $65$;
Bershadsky--Ooguri $1989$ CMP $126$;
Pope--Romans--Shen $1990$ NPB $339$;
Bouwknegt--Schoutens $1993$ Phys.\ Rep.\ $223$;
Wang $1998$ CMP $195$;
Nakayama $2004$ IJMPA $19$;
Schiffmann--Vasserot $2013$ Publ.\ Math.\ IHES $118$;
Oberdieck $2018$ Geom.\ \& Topol.\ $22$;
Oberdieck--Pixton $2018$ Invent.\ Math.\ $213$;
Kontsevich $2003$ Lett.\ Math.\ Phys.\ $66$.


% ==========================================
% PART II: E_1 and E_2 Chiral Theories
% ==========================================

\part{$E_1$ and $E_2$ Chiral Theories}
\label{part:en-chiral}

\chapter{$E_1$-Chiral Algebras}
\label{ch:e1-chiral}

Volume~II (Part~III) developed the $E_1$-chiral theory as the ordered sector of the Swiss-cheese algebra.  This chapter reformulates it as the foundation for the $E_2$ extension.

% AP104: The E_1/ordered story is the PRIMITIVE layer.  The E_2 braided
% structure is obtained from it by the Drinfeld center construction.
% Quantum groups, Yangians, and braided tensor categories are natively E_1.
% The passage E_1 -> E_2 is the higher-categorical analogue of the
% averaging map av: g^{E_1} -> g^mod from Vol I.

\begin{remark}[$E_1$ primacy for CY quantum groups]
\label{princ:e1-primacy-vol3}
The $E_1$-chiral algebra (boundary) is the primitive object in this volume.
The $E_2$-chiral algebra (bulk) is obtained from it by the Drinfeld center
construction $\cZ(\Rep^{E_1}(A)) \simeq \Rep^{E_2}(\mathrm{Drin}(A))$.
Quantum groups, Yangians, and braided tensor categories are natively $E_1$
objects: the CoHA multiplication is ordered (short exact sequences have a
preferred direction), and the $R$-matrix arises only in the Drinfeld double.
The passage $E_1 \to E_2$ is the higher-categorical analogue of the
averaging map $\mathrm{av} \colon \fg^{E_1} \to \fg^{\mathrm{mod}}$
from Vol~I.
This averaging is \emph{lossy}: $\mathrm{av}(r(z)) = \kappa_{\mathrm{ch}}$ forgets the
$R$-matrix data.  The $E_1$-bar $B^{\mathrm{ord}}(A)$ retains
strictly more information than the $E_\infty$-bar $B^{\Sigma}(A)$.
\end{remark}

\section{$E_1$-algebras and associative factorization}
\label{sec:e1-algebras}

An $E_1$-algebra is an algebra over the little intervals operad.  In the chiral setting, an $E_1$-chiral algebra factorizes associatively along a real line $\mathbb{R}$ while maintaining holomorphic structure along a transverse curve.

\begin{definition}[$E_1$-chiral algebra]
\label{def:e1-chiral-algebra}
An \emph{$E_1$-chiral algebra} on $X \times \mathbb{R}$ is a factorization algebra $\cA$ on $\Ran(X) \times \Ran(\mathbb{R})$ such that:
\begin{enumerate}[label=(\roman*)]
  \item The $X$-direction factorization is holomorphic (chiral);
  \item The $\mathbb{R}$-direction factorization is $E_1$ (associative).
\end{enumerate}
The representation category $\Rep^{E_1}(\cA)$ is a monoidal category.
\end{definition}

\section{The bar complex as $E_1$-chiral algebra}
\label{sec:bar-e1}

\begin{proposition}
\label{prop:bar-e1-structure}
\ClaimStatusProvedElsewhere
For a chiral algebra $A$ on $X$, the bar complex $B(A)$ carries a natural $E_1$-chiral structure: the bar differential encodes the $X$-direction (holomorphic OPE residues) and the bar coproduct encodes the $\mathbb{R}$-direction (ordered deconcatenation).
\end{proposition}
% AP85: The factorization coproduct (Vol I, coshuffle/Sym^c) is DIFFERENT
% from the deconcatenation coproduct (Vol II, T^c).  The E_1-bar uses
% deconcatenation; the E_infty-bar uses coshuffle.  See AP82, AP85.

This is the Swiss-cheese algebra structure of Volume~II.  The $E_1$-bar
$B^{\mathrm{ord}}(A)$ is the ordered tensor coalgebra $(T^c(A[1]), d, \Delta_{\mathrm{decat}})$
with deconcatenation coproduct (AP82, AP85).  The symmetric bar
$B^{\Sigma}(A) = B^{\mathrm{ord}}(A) / S_n$ is the $S_n$-coinvariant quotient.
The ordered bar retains the full quantum group data (R-matrix,
Yangian structure) that the symmetric bar quotients out (AP65, AP97).

\section{Monoidal representation categories}
\label{sec:e1-monoidal-rep}

The representation category $\Rep^{E_1}(\cA)$ inherits a monoidal structure
from the ordered factorization along $\mathbb{R}$.  The tensor product
$V \otimes W$ of two $\cA$-modules is defined by the fusion of intervals:
$V$ is placed on $(-\infty, 0)$ and $W$ on $(0, \infty)$, and the
factorization structure combines them.  The ordering of the intervals makes
the tensor product associative but \emph{not} symmetric: the natural
transformation $V \otimes W \to W \otimes V$ does not exist in the
$E_1$-setting.  It exists only after passing to the Drinfeld center
(Chapter~\ref{ch:drinfeld-center}).

\section{The $E_1$-Koszul duality}
\label{sec:e1-koszul}

\begin{theorem}[$E_1$-chiral Koszul duality, Volume~II]
\label{thm:e1-koszul}
\ClaimStatusProvedElsewhere
For an $E_1$-chiral algebra $\cA$, the bar-cobar adjunction of Volume~I refines to an $E_1$-equivariant adjunction.  On the representation-category level, this gives a monoidal equivalence
\[
  \Rep^{E_1}(A) \simeq \Rep^{E_1}(A^!)^{\mathrm{rev}}
\]
between the monoidal representation category of $A$ and the reverse monoidal category of $A^!$.
\end{theorem}
% AP25: A^! here means the Koszul dual algebra, i.e., the linear dual of
% bar cohomology H*(B(A)). This is NOT the cobar Omega(B(A)) (which
% recovers A itself) nor the Verdier dual D_Ran(B(A)) (which is B(A!)).
% See Vol I, Convention conv:bar-coalgebra-identity.

\chapter{$E_2$-Chiral Algebras}
\label{ch:e2-chiral}

This is the central innovation of the present work.  An $E_2$-chiral algebra factorizes along two holomorphic directions, and its representation category is \emph{braided} monoidal --- the natural algebraic structure underlying quantum groups.

\section{The $E_2$ operad and braided monoidal categories}
\label{sec:e2-operad}

\begin{definition}[$E_2$ operad]
\label{def:e2-operad}
The $E_2$ operad (little 2-disks operad) has $E_2(n) = \mathrm{Conf}_n(\mathbb{D}^2)$, the configuration space of $n$ non-overlapping disks in the unit disk.  An $E_2$-algebra is an algebra over this operad.  By Dunn additivity, $E_2 \simeq E_1 \otimes E_1$: an $E_2$-algebra is an $E_1$-algebra in $E_1$-algebras.
\end{definition}

\begin{theorem}[Lurie]
\label{thm:lurie-e2-braided}
The category of $E_2$-algebras in a symmetric monoidal $\infty$-category $\cC$ is equivalent to the category of braided monoidal objects in $\cC$.
\end{theorem}

\section{$E_2$-chiral algebras: definition and first properties}
\label{sec:e2-chiral-def}

\begin{definition}[$E_2$-chiral algebra]
\label{def:e2-chiral-algebra}
An \emph{$E_2$-chiral algebra} on $X \times Y$ (where $X, Y$ are smooth curves) is a factorization algebra $\cA$ on $\Ran(X) \times \Ran(Y)$ such that:
\begin{enumerate}[label=(\roman*)]
  \item The $X$-direction factorization is chiral (holomorphic);
  \item The $Y$-direction factorization is chiral (holomorphic);
  \item The combined structure is compatible with the $E_2$-operad action via the identification $\FM_n(\mathbb{C}) \simeq E_2(n)$.
\end{enumerate}
The representation category $\Rep^{E_2}(\cA)$ is a \emph{braided monoidal} category.
\end{definition}

\section{Kontsevich formality and the $E_2$-Gerstenhaber structure}
\label{sec:kontsevich-formality-e2}

Kontsevich's formality theorem $E_2 \xrightarrow{\sim} \Ger$ (over $\mathbb{Q}$) identifies $E_2$-algebras with Gerstenhaber algebras up to homotopy.  In the chiral context, this gives:

\begin{proposition}
\label{prop:e2-gerstenhaber-chiral}
The Hochschild cohomology $\HH^\bullet(A)$ of a chiral algebra $A$ carries a natural $E_2$-algebra structure.  Via Kontsevich formality, this is equivalent to the Gerstenhaber bracket + cup product structure.
\end{proposition}

\section{The $E_2$ bar complex}
\label{sec:e2-bar}

\begin{construction}[$E_2$ bar complex]
\label{constr:e2-bar}
For an $E_2$-chiral algebra $\cA$ on $X \times Y$, the \emph{$E_2$ bar complex} $B_{E_2}(\cA)$ is a factorization coalgebra on $\Ran(X) \times \Ran(Y)$ with:
\begin{itemize}
  \item Two commuting differentials: $d_X$ from $X$-direction OPE residues, $d_Y$ from $Y$-direction OPE residues;
  \item Two commuting coproducts: $\Delta_X$ from $X$-direction deconcatenation, $\Delta_Y$ from $Y$-direction deconcatenation;
  \item The braiding: the transposition of the two $E_1$-structures gives the $R$-matrix.
\end{itemize}
\end{construction}

\section{$E_2$-chiral Koszul duality}
\label{sec:e2-koszul}

\begin{conjecture}[$E_2$-chiral Koszul duality]
\label{conj:e2-koszul}
\ClaimStatusConjectured
For an $E_2$-chiral algebra $\cA$, the bar-cobar adjunction of Volume~I extends to an $E_2$-equivariant adjunction.  On representation categories, this gives a braided monoidal equivalence
\[
  \Rep^{E_2}(A) \simeq \Rep^{E_2}(A^!)^{\mathrm{rev}}
\]
where $\mathrm{rev}$ denotes the reverse braiding.
\end{conjecture}
% AP25: A^! is the Koszul dual algebra, NOT the cobar of B(A) (which
% recovers A itself), and NOT the Verdier dual D_Ran(B(A)) (which is
% B(A!) as a factorization coalgebra). The "rev" reversal is specific
% to the E_2 setting; at E_1, one has Rep^{E_1}(A) = Rep^{E_1}(A^!)^{rev}
% as monoidal (not braided) categories (Theorem thm:e1-koszul).

\chapter{$E_n$-Factorization and Higher Chiral Structure}
\label{ch:en-factorization}

\section{Dunn additivity and the $E_n$ hierarchy}
\label{sec:dunn-additivity}

\section{Factorization algebras on $\mathbb{C}^n$}
\label{sec:fact-cn}

\section{The $E_n$-chiral bar complex}
\label{sec:en-bar}

\section{Relation to topological field theory}
\label{sec:en-tft}


% ==========================================
% PART III: The Bridge --- CY Categories as Quantum Chiral Algebras
% ==========================================

\part{The Bridge: Calabi--Yau Categories as Quantum Chiral Algebras}
\label{part:bridge}

\chapter{From CY Categories to Chiral Algebras}
\label{ch:cy-to-chiral}

This chapter constructs the bridge functor $\Phi$ from CY categories to quantum chiral algebras.

\section{The cyclic-to-chiral passage}
\label{sec:cyclic-to-chiral}

A CY category $\cC$ of dimension $d$ carries a cyclic $\Ainf$-structure (Chapter~\ref{ch:cyclic-ainf}).  The cyclic bar complex $\mathrm{CC}_\bullet(\cC)$ with its $\mathbb{S}^d$-framing is the primary invariant.  The passage to chiral algebras proceeds by:

\begin{enumerate}[label=\textbf{Step \arabic*.}]
  \item \textbf{Cyclic $\Ainf \to$ Lie conformal algebra.}  The Hochschild cohomology $\HH^\bullet(\cC)$, with its Gerstenhaber bracket, determines a graded Lie algebra.  The CY pairing promotes this to a Lie conformal algebra $\mathfrak{L}_\cC$ (the ``current algebra'' of $\cC$).
  \item \textbf{Lie conformal algebra $\to$ factorization envelope.}  The factorization envelope construction (Nishinaka--Vicedo) produces a factorization algebra $\Fact_X(\mathfrak{L}_\cC)$ on any smooth curve $X$.
  \item \textbf{$E_2$ enhancement.}  When $d = 2$, the $\mathbb{S}^2$-framing of $\HH_\bullet(\cC)$ provides an $E_2$-algebra structure.  This enhances $\Fact_X(\mathfrak{L}_\cC)$ to an $E_2$-chiral algebra.
  \item \textbf{Quantization.}  The quantum chiral algebra $A_\cC$ is the quantization of $\Fact_X(\mathfrak{L}_\cC)$ determined by the CY trace.
\end{enumerate}

\section{The CY-to-chiral functor}
\label{sec:cy-chiral-functor}

\begin{theorem}[CY-to-chiral functor --- Theorem CY-A$_2$, $d = 2$]
\label{thm:cy-to-chiral}
There exists a functor
\[
  \Phi \colon \CY_2\text{-}\Cat \longrightarrow E_2\text{-}\mathrm{ChirAlg}
\]
from smooth proper CY categories of dimension $d = 2$ to $E_2$-chiral algebras, such that:
\begin{enumerate}[label=(\roman*)]
  \item $\Phi(\cC)$ is a chiral algebra whose generating fields are in bijection with $\HH^{\bullet+1}(\cC)$;
  \item The bar complex $B(\Phi(\cC))$ is quasi-isomorphic to the cyclic bar complex $\mathrm{CC}_\bullet(\cC)$ as a factorization coalgebra;
  \item The modular characteristic $\kappa(\Phi(\cC))$ equals the CY Euler characteristic.
\end{enumerate}
The $E_2$-enhancement in Step~3 uses the $\mathbb{S}^2$-framing of $\HH_\bullet(\cC)$ (Kontsevich--Vlassopoulos).
\end{theorem}

\begin{conjecture}[CY-to-chiral programme --- CY-A$_3$, $d = 3$]
\label{conj:cy-to-chiral-d3}
For a smooth proper CY$_3$ category $\cC$, the functor $\Phi$ extends to produce a quantum vertex chiral group $G(\cC)$ whose generalized root datum is $\mathcal{R}(\cC)$.

The $\mathbb{S}^3$-framing gives an $E_3$-structure on $\HH_\bullet(\cC)$.  The $E_3$-structure restricts to $E_2$ via Dunn additivity, but the restricted braiding is \emph{symmetric} at the topological level (since $\pi_1(\mathrm{Conf}_2(\mathbb{R}^3))$ is trivial for the ordered configuration space).  The quantum group braiding for $d = 3$ is expected to arise through the \emph{Drinfeld center} of the $E_1$-monoidal representation category, not through direct $E_3$-to-$E_2$ restriction.

This programme is conditional on: (a) constructing the chain-level $\mathbb{S}^3$-framing compatible with the BV structure, and (b) establishing the quantization step (Step~4) for $d = 3$.
\end{conjecture}

\section{Compatibility with Koszul duality}
\label{sec:cy-koszul-compat}

\section{The CY modular characteristic}
\label{sec:cy-modular-char}

\phantomsection\label{def:shadow-class}%
\phantomsection\label{def:shadow-invariants}%
\providecommand{\Br}{\mathrm{Br}}%
\providecommand{\Drinfdouble}[1]{D(#1)}%
\providecommand{\PhiFA}{\Phi^{\mathrm{FA}}}%
\providecommand{\hCS}{\mathrm{hCS}}%
\providecommand{\Wonepinf}{\mathcal{W}_{1+\infty}}%
\providecommand{\protectedPfaff}[1]{\mathrm{Pf}^{\mathrm{prot}}(#1)}%

\chapter{Quantum Chiral Algebras}
\label{ch:quantum-chiral-algebras}

\begin{remark}[Two-stage locus of the holomorphic Chern--Simons observables]
\label{rem:qca-two-stage-locus}
The six-dimensional holomorphic Chern--Simons observables constructed in
this chapter are the local BV/factorisation-algebra model for the
Stage-$1$ output $\Phi^{\mathrm{FA}}_3(\cA)$. They are not, by
themselves, the Hall algebra or the curve-specialised chiral algebra.
The chiral-algebra-on-a-curve output is the Stage-$2$ specialisation
along chosen data $(\Sigma_2, C)$.  The finite comparison presently
available is the DWR/Ran Rees comparison at fixed truncation
$(N,r,L,m)$: it compares the renormalised hCS factorisation algebra with
the oriented Rees critical Hall layer.  Passing from that finite Rees
layer to the ordinary compact critical CoHA additionally requires the
vanishing-cycle realisation map and pro-compatible inverse-limit descent.
\end{remark}

\begin{theorem}[Holomorphic Chern--Simons observables as \texorpdfstring{$\Ethree$}{Ethree}-Dolbeault--chiral Chevalley--Eilenberg algebra]
\label{thm:hcs-obs-e3-ce-dolbeault}
\[
 \mathrm{Obs}^{\mathrm{cl}}_{\mathrm{hCS}}(X) \;\simeq\; \mathrm{CE}^\bullet_{\bar\partial, \mathrm{chir}}(\cE_{\mathrm{hCS},c}, \cO_X)
\]
as completed classical functions on the compactly supported local dg Lie
algebra. The quantum observables are the Costello--Gwilliam
perturbative factorisation algebra carrying a scale-dependent
effective action, BV operator, renormalisation counterterms, and
anomaly class; the classical CE differential supplies the
$\hbar = 0$ associated graded of the full quantum differential. The
ambient $E_3$ structure is higher factorisation coherence
(commutation up to specified homotopy at each configuration arity);
the connected components of $S^5$ are trivial, so the $\pi_1(S^5) = 0$
contribution is the trivial braid monodromy, and the $E_3$ data lives
in the higher homotopy of $\Conf_n(\C^3)$.
\end{theorem}

The central objects of this chapter are $\Eone$-chiral algebras whose
centres or doubles carry the braided representation theory of quantum
groups.  At $d=2$ the CY-to-chiral output is natively $\Etwo$.  At
$d\geq 3$ the curve-specialised output is natively $\Eone$; the first
$\Etwo$ object is the Drinfeld centre of its representation category, or
equivalently the completed Drinfeld double when a Hall pairing and
completion have been fixed.  The finite DWR/Ran Rees hCS--Hall
comparison lives before this centre step.

\section{Definition and structure}
\label{sec:qca-definition}

\begin{definition}[Quantum chiral algebra in dimension stratified form]
\label{def:qca-dimension-stratified}
A \emph{quantum chiral algebra} in this chapter is a chiral algebra $A$
with finite-dimensional weight spaces together with the following
dimension-dependent structure.
\begin{enumerate}[label=\textup{(\roman*)}]
\item For CY$_2$ input, $A=\Phi_2(\cC)$ is an $\Etwo$-chiral algebra.
Its collision braiding gives an $R$-matrix acting on tensor products of
representations.
\item For CY$_d$ input with $d\geq3$ and admissible specialisation data,
\[
 A=A_{\cC}^{(\Sigma_{d-1},C)}
   =\SpCh_{\Sigma_{d-1},C}\bigl(\PhiFA_d(\cC)\bigr)
\]
is an $\Eone$-chiral algebra on $C$.  The braided $\Etwo$ object
attached to $A$ is the Drinfeld centre
$\cZ(\Rep^{\Eone}(A))$.
\item An expression \(R(z)\in A\otimes A\otimes k((z))\) is used only
after a Hopf, double, or matrix-coefficient realisation of the centre has
been chosen.  Intrinsically, the centre carries a half-braiding natural
transformation (equivalently a universal intertwiner); \(R(z)\) is the
meromorphic matrix-coefficient or evaluation family obtained from that
datum after the chosen realisation and spectral parameters.
\end{enumerate}
\end{definition}

\begin{definition}[Five objects attached to an \texorpdfstring{$\Eone$}{E1}-chiral algebra]
\label{def:qca-five-objects}
Let \(A\) be an \(\Eone\)-chiral algebra on a curve.
\begin{enumerate}[label=\textup{(\roman*)}]
\item \(A\) is the associative chiral algebra, for CY$_3$ obtained after
Stage-$2$ specialisation.
\item \(B(A)\) is its conilpotent chiral bar coalgebra.  On a Koszul
locus the counit \(\Omega B(A)\to A\) is a bar--cobar inversion: $A$
is recovered as the cobar $\Omega$ of $B(A)$, with $B(A)$ a distinct
coalgebra object carrying the twisting/coupling structure on the
universal arrow.
\item \(A^i:=H^\bullet(B(A))\) is the bar cohomology coalgebra, or a
chosen bar-coalgebra model before passing to cohomology.
\item \(A^!:=\mathbb D_{\Ran}(A^i)\), or the corresponding continuous
linear dual in a finite model, is the Koszul dual algebra.
\item \(Z^{\mathrm{der}}_{\mathrm{ch}}(A)\) is the derived chiral centre,
the Hochschild cochain object \(\mathrm{RHom}^{\mathrm{ch}}_{A-A}(A,A)\).
\end{enumerate}
The Drinfeld centre \(\cZ(\Rep^{\Eone}(A))\) is a braided monoidal
category.  It is related to, but not equal to, the derived chiral centre
as a chiral algebra.
\end{definition}

\begin{proposition}[Finite Rees hCS--Hall layer separation]
\label{prop:qca-finite-rees-layer-separation}
Let \(X\) be a CY$_3$ equipped with a finite DWR-good atlas and finite
truncation data \((N,r,L,m)\).  Assume the cyclic contractions,
orientation, grading, trace/symmetry, and face-compatibility hypotheses
of Theorem~\ref{thm:finite-dwr-ran-oriented-hcs-hall-comparison} and
Theorem~\ref{thm:finite-dwr-ran-rees-hcs-hall-compatibility}.  Then:
\begin{enumerate}[label=\textup{(\roman*)}]
\item \(\PhiFA_3(\Perf(X))\) is represented locally by the holomorphic
\(E_3\)-factorisation algebra of hCS observables, while every
curve-specialised output
\[
 A_X^{(\Sigma_2,C)}
   =\SpCh_{\Sigma_2,C}\bigl(\PhiFA_3(\Perf(X))\bigr)
\]
is only \(\Eone\)-chiral.
\item At the fixed bounds \((N,r,L,m)\) there is an oriented Rees
comparison
\[
 \Theta_{\hCS\to\Hall}^{\mathrm{Rees};N,r,L,m}\colon
 \Obs_{\hCS}^{q,\leq(N,r,L,m)}
 \longrightarrow
 \Hall_{\mathrm{crit}}^{\mathrm{Rees,or};N,r,L,m}.
\]
It is a morphism of finite ordered Hall-valued factorisation cosheaves.
It is not yet a morphism to the ordinary compact critical CoHA.
\item The ordinary compact critical CoHA requires a monoidal
vanishing-cycle realisation \(R_{\mathrm{vc}}\).  The completed compact
comparison further requires pro-compatible transition data in
\((N,r,L,m)\); the remaining obstruction is the corresponding
\(\varprojlim^1\)-class.
\item On a single affine chart \(X=\C^3\),
\[
 \CoHA(\C^3)=Y^+(\widehat{\fgl}_1),
\]
the positive half.  The full affine Yangian, its universal \(R\)-matrix,
and the \(W_{1+\infty}\) Fock action arise only after passing to the
Drinfeld double or to \(\cZ(\Rep^{\Eone}(Y^+))\).
\item For a multi-chart CY$_3$, DWR/Ran descent glues the finite Rees
Hall layer chart by chart; the further constructions — the global
Drinfeld double, the global \(W_{1+\infty}\) vertex algebra, and the
\(\Etwo\)-chiral structure on the Drinfeld centre
\(\cZ(\Rep^{\Eone}(A_X^{(\Sigma_2,C)}))\) — are obtained by composing
DWR/Ran descent with the centre/double step.
\end{enumerate}
\end{proposition}

\begin{proof}
The two-stage assertion is Theorem~\ref{thm:phi-two-stage-factorisation}
and the native-level statement is the \(d=3\) case of
Theorem~\ref{thm:cy-to-chiral-d3}.  The finite comparison map is the
DWR/Ran Rees construction of
Theorem~\ref{thm:finite-dwr-ran-oriented-hcs-hall-comparison}, in the
ordered \(E_1\) form recorded by
Theorem~\ref{thm:finite-dwr-ran-rees-hcs-hall-compatibility}.  Those
theorems produce a finite Rees critical Hall target; their hypotheses do
not include the monoidal vanishing-cycle realisation to ordinary
critical CoHA or the inverse-limit compatibility required to leave fixed
\((N,r,L,m)\).  The local \(\C^3\) identification is the
Schiffmann--Vasserot/RSYZ positive-half computation; the full Yangian is
obtained by the Drinfeld double, equivalently by the Drinfeld centre of
the representation category.  Drinfeld centre and completion are
functorial operations after the ordered Hall layer has been supplied, so
they cannot be inserted into the finite Rees comparison map.
\end{proof}

\begin{problem}[Quantum vertex chiral group]
\label{prob:qvcg-construction}
For a smooth proper CY$_3$ category $\cC$ with $X$ the underlying
manifold, construct a \emph{quantum vertex chiral group} \(G(X)\) with:
\begin{enumerate}[label=(\alph*)]
 \item A generalized BKM superalgebra $\mathfrak{g}_X$ with root datum $\mathcal{R}(X)$ extracted from the CY$_3$ geometry (Chapter~\ref{ch:k3-times-e});
 \item A vertex algebra structure $V(\mathfrak{g}_X)$: the vacuum module of $\mathfrak{g}_X$ equipped with the state-field correspondence;
 \item An $E_2$-braided monoidal category of representations $\Rep^{E_2}(V(\mathfrak{g}_X))$;
 \item A denominator identity: the Weyl--Kac--Borcherds character formula for $\mathfrak{g}_X$, which equals an automorphic form $\Phi_X$ (e.g., $\Delta_5$ for $K3 \times E$).
\end{enumerate}
\end{problem}

\begin{remark}[Construction locus of \texorpdfstring{$G(X)$}{G(X)}]
\label{rem:qvcg-construction-locus}
Problem~\ref{prob:qvcg-construction} is a construction problem: a
definition of $G(X)$ awaits an explicit construction, or an axiomatic
characterisation implying existence. The currently constructed pieces
are these:
\begin{itemize}
 \item For toric CY$_3$ without compact $4$-cycles, the CoHA
 \(\mathcal H(Q_X,W_X)\) gives the positive half; its Drinfeld double
 supplies the braided representation theory.
 \item For \(d=2\), \(\Phi_2\) constructs the native \(\Etwo\)-chiral
 algebra.
 \item For \(K3\times E\), the BKM superalgebra
 \(\mathfrak g_{\Delta_5}\) and root datum are supplied by
 Gritsenko--Nikulin.  The curve-specialised chiral algebra is the
 framed object-level \(A_{K3\times E}^{(K3,E)}\) of
 Theorem~\ref{thm:cy-to-chiral-d3}; its braided representation theory
 is obtained only after the centre/double operation.
\end{itemize}
Thus a sentence about \(G(X)\) is a theorem only after it factors
through one of the constructed pieces above, or after
Problem~\ref{prob:qvcg-construction} is solved.
\end{remark}

An $\Etwo$-structure on \(A\), or on \(\cZ(\Rep^{\Eone}(A))\) when
\(A\) is only \(\Eone\), determines a universal $R$-matrix
\[
 R(z) \in A \otimes A \otimes k((z))
\]
satisfying the quantum Yang--Baxter equation after a universal
matrix-coefficient realisation has been chosen.
This $R$-matrix is the $E_2$ analogue of the collision $r$-matrix $r(z) = \mathrm{Res}^{\mathrm{coll}}_{0,2}(\Theta_A)$ from Volume~I (whose pole orders are one less than the OPE poles, due to the $d\log$ extraction; see Vol~I).

\section{The CY \texorpdfstring{$R$}{R}-matrix}
\label{sec:cy-r-matrix}

\begin{construction}[CY \texorpdfstring{$R$}{R}-matrix]
\label{constr:cy-r-matrix}
Given a CY category $\cC$ of dimension $2$ and its quantum chiral algebra $A_\cC = \Phi_2(\cC)$, the CY $R$-matrix is
\[
 R_\cC(z) = \mathrm{Res}^{\mathrm{coll}}_{0,2}(\Theta_{A_\cC})
\]
where $\Theta_{A_\cC}$ is the universal MC element from Volume~I. The pole structure of $R_\cC(z)$ encodes the CY shadow depth classification.
\end{construction}

\section{Drinfeld center and the \texorpdfstring{$E_1 \to E_2$}{E-1 -> E-2} passage}
\label{sec:drinfeld-center-yangian}

The passage from \(E_1\) to \(E_2\) is realised by the Drinfeld centre construction. For \(X=\C^3\), the positive CoHA is \(Y^+(\widehat{\mathfrak{gl}}_1)\); after the Hall pairing and Drinfeld double are imposed, the centre of its \(E_1\)-representation category carries the braided category of the full affine Yangian.

The full quantum group attached to a Calabi--Yau target $X$ is not the
positive half $\Yplus{X}$ alone: the positive half is constructed
intrinsically from the equivariant geometry of moduli of effective
objects, while the full quantum group $G(X)$ is its Drinfeld double, a
strictly higher datum requiring a Hall pairing, a completion, an
integral form, stable-envelope transport, and Drinfeld-double descent.
Conflation of $\Yplus{X}$ with $G(X)$ (e.g.\ ``$\CoHA(\C^3) =
\Wonepinf$'') collapses the two-stage assembly. The next definition and
theorem record the assembly explicitly; their toric specialisation is
\cref{thm:drinfeld-center-coha} below.

\begin{definition}[\label{def:Y-plus-equiv-cohomology}Positive half via
equivariant cohomology of moduli of effective objects]
Let $\cA$ be a Calabi--Yau-$d$ category with target $X$, equipped with
the geometric data \textup{(}equivariant torus action $T$ on
$\cM^+_{\mathrm{eff}}(X)$ with isolated $T$-fixed loci; potential $W
\colon \cM^+_{\mathrm{eff}}(X) \to \C$ from the Calabi--Yau
superpotential; orientation data and Serre-compatible NCCR chart on the
critical locus\textup{)} required for the construction below to be
well posed. The \emph{positive half} of $\cA$ is the equivariant
critical cohomology
\[
 \Yplus{X} \;:=\; H^\bullet_{T}\!\bigl(\cM^+_{\mathrm{eff}}(X),
 \phi_W\bigr)
\]
of the moduli of effective objects with respect to the vanishing-cycle
sheaf $\phi_W$, equipped with the Schiffmann--Vasserot Hall product
arising from extension correspondences
\[
 \cM^+_{\mathrm{eff}}(X) \times \cM^+_{\mathrm{eff}}(X)
 \;\xleftarrow{\;p\;}\;
 \cE\!xt^+(X)
 \;\xrightarrow{\;q\;}\;
 \cM^+_{\mathrm{eff}}(X),
 \qquad
 m \;=\; q_\star \circ p^\star.
\]
At $d = 3$, $\Yplus{X}$ is the cohomological Hall algebra
$\CoHA(X)$ of Schiffmann--Vasserot \cite{SchiffmannVasserot2013}; on
$\C^3$ it gives $\CoHA(\C^3) = \Yplus{\widehat{\fgl}_1}$
\textup{(}\cite{SchiffmannVasserot2013}, identified with the positive
half of the affine Yangian, \emph{not} with $\Wonepinf$\textup{)}.

The comparison data are $(p^\star,q_\star,\phi_W)$. The ambient category is
$T$-equivariant cohomology with critical coefficients, in
$\mathrm{Ch}(\mathrm{Vect}_\Q)$ at the chain level and in the perverse
$\infty$-category of constructible sheaves on
$\cM^+_{\mathrm{eff}}(X)$ at the $(\infty,1)$-categorical level. The construction is
unconditional on $\C^3$ and standard toric CY$_3$ charts; on compact
non-toric CY$_3$ targets it is conditional on a compact critical atlas
and orientation data.
\end{definition}

\begin{theorem}[\label{thm:G-eq-D-Yplus}Drinfeld-double identification of
the full quantum group]
\emph{Hypotheses.} Let $\cA$ be a Calabi--Yau-$d$ category with target
$X$ for which $\Yplus{X}$ is constructed in the sense of
\cref{def:Y-plus-equiv-cohomology}. Assume the following five
data:
\begin{enumerate}[label=\textup{(LD\arabic*)}]
\item \emph{Hall pairing.}
A non-degenerate bilinear pairing $\langle -, - \rangle_{\mathrm{Hall}}
\colon \Yplus{X} \otimes \Yplus{X} \to \kk$ obtained from
Serre/Calabi--Yau duality on $\cM^+_{\mathrm{eff}}(X)$
\textup{(}Schiffmann--Vasserot pairing on critical
cohomology\textup{)}.
\item \emph{Completion.}
A topology / completion of $\Yplus{X}$ \textup{(}grading-completion,
$J$-adic completion, or weight-completion\textup{)} for which
\textup{(LD1)} converges and the antipode is well-defined.
\item \emph{Integral form.}
A subring $\Yplus{X}_{\Z} \subset \Yplus{X}$ \textup{(}or
$\Yplus{X}_{\Z[\hbar]}$\textup{)} compatible with the Hall product,
coproduct, and pairing.
\item \emph{Stable-envelope transport.}
A Maulik--Okounkov stable-envelope realisation
$\mathrm{Stab} \colon \Yplus{X} \to \End(\cF_X)$ on the appropriate
Fock-type representation $\cF_X$, satisfying the chamber-by-chamber
factorisation that intertwines the Hall product with the stable-basis
matrix coefficient \cite{MaulikOkounkov2012}.
\item \emph{Drinfeld-double descent.}
A descent datum identifying the algebra
$\Yplus{X} \bowtie (\Yplus{X})^{\ast,\mathrm{cop}}$,
constructed from \textup{(LD1)--(LD4)}, with the chamber-independent
Drinfeld double, well-defined on the underlying critical cohomology
independent of the chamber and orientation choices used in
\textup{(LD4)}.
\end{enumerate}

\emph{Statement.} Under \textup{(LD1)--(LD5)}, the full quantum group
$G(X)$ associated to $\cA$ is the Drinfeld double of the positive
half $\Yplus{X}$:
\begin{equation}
\label{eq:G-eq-D-Yplus}
 G(X) \;=\; \Drinfdouble{\Yplus{X}}
 \;=\; \Yplus{X} \;\bowtie\; (\Yplus{X})^{\ast,\mathrm{cop}},
\end{equation}
where $\bowtie$ is the Hall-pairing biproduct of \textup{(LD1)}
completed in the topology of \textup{(LD2)} on the integral form
\textup{(LD3)}, and the chamber-independence is supplied by
\textup{(LD4)--(LD5)}. The braided category
$\Rep^{E_1}(\Drinfdouble{\Yplus{X}}) = \Rep(G(X))$ realises the
Drinfeld centre of $\Rep^{E_1}(\Yplus{X})$
\textup{(}Majid 1991; Etingof--Gelaki--Nikshych--Ostrik,
Proposition~7.13.8\textup{)}, and the universal $R$-matrix of
$G(X)$ is the half-braiding inherited from this centre.

\emph{Licensing.} Tag $\beta$ \textup{(}comparison--functor:
$\Drinfdouble{\,\cdot\,}$ and the Hall--pairing biproduct\textup{)}.
The full hypothesis set \textup{(LD1)--(LD5)} is the comparison datum;
the theorem is the equality of two algebras over the same critical
cohomology, not a definition of $G(X)$.

\emph{Constructed loci.}
\begin{enumerate}[label=\textup{(\roman*)}]
\item $X = \C^3$: $G(\C^3) = Y(\widehat{\fgl}_1)$ is the affine Yangian
of $\widehat{\fgl}_1$, with $\Yplus{\C^3} = \CoHA(\C^3) =
\Yplus{\widehat{\fgl}_1}$
\textup{(}\cite{SchiffmannVasserot2013}; cf.\
\cref{thm:drinfeld-center-coha}\textup{(i)}\textup{)}.
\item Toric CY$_3$ without compact $4$-cycles: $G(X)$ is the
quantum-toric Yangian $Y_{\mathrm{tor}}(X)$, constructed from the toric
quiver with potential $(Q_X, W_X)$
\textup{(}cf.\ Chapter~\ref{ch:toric-coha}\textup{)}.
\item $X = $ flat affine $\C^2 \times \C$: \textup{(LD1)--(LD5)} are
all chain-level constructible from
$\cM^+_{\mathrm{eff}}(\C^2 \times \C) \cong \mathrm{Hilb}(\C^2)
\times \mathrm{pt}$ and the Maulik--Okounkov stable-envelope
construction on Hilbert schemes \cite{MaulikOkounkov2012}.
\end{enumerate}

\emph{Compact non-toric loci.}
On compact non-toric CY$_3$ \textup{(}quintic, generic Hartshorne
threefold\textup{)} and on $K3 \times E$, the construction
\textup{(LD1)--(LD5)} of \cref{eq:G-eq-D-Yplus} requires the named
data of Kontsevich--Soibelman \cite{KontsevichSoibelman2008}: critical
CoHA on a compact critical atlas, monoidal vanishing-cycle realisation
$R_{\mathrm{vc}}$, pro-compatible inverse-limit transition data, and
PBW/Hall integral form. On these loci the equality \cref{eq:G-eq-D-Yplus}
is the explicit comparison map, not a black-box definition.
\end{theorem}

\begin{proof}[Sketch on the constructed loci]
On $X = \C^3$: by Schiffmann--Vasserot
\cite{SchiffmannVasserot2013}, $\Yplus{\C^3}$ is identified with the
positive half $\Yplus{\widehat{\fgl}_1}$ of the affine Yangian, with
the Hall product matching the Yangian coproduct on shuffle generators.
The Hall pairing \textup{(LD1)} is the Schiffmann--Vasserot pairing on
shuffle algebras, non-degenerate on each weight space; the completion
\textup{(LD2)} is the standard $\hbar$-adic / weight-completion of the
Yangian; the integral form \textup{(LD3)} is the
$\Z[h_1, h_2, h_3]$-form of the shuffle algebra; the stable-envelope
transport \textup{(LD4)} is supplied by Maulik--Okounkov on
$\mathrm{Hilb}^n(\C^2)$ \cite{MaulikOkounkov2012}; chamber-independence
\textup{(LD5)} is the wall-crossing / $R$-matrix factorisation of
\cite[\S 3]{MaulikOkounkov2012}. The Drinfeld-double construction
$\Yplus{X} \bowtie (\Yplus{X})^{\ast,\mathrm{cop}}$ then identifies
with $Y(\widehat{\fgl}_1)$: the negative generators
$Y^-(\widehat{\fgl}_1)$ are the graded dual of $\Yplus{\widehat{\fgl}_1}$
with reversed coproduct, and the cross-relations are read off from the
Hall pairing \textup{(LD1)} \textup{(}Schiffmann--Vasserot
\cite{SchiffmannVasserot2013}, Theorem~9.1\textup{)}. The toric and
flat cases reduce to $\C^3$ on each toric chart by the same argument
applied chartwise.

The non-toric and $K3 \times E$ cases require
\cite{KontsevichSoibelman2008}'s motivic-/critical-CoHA framework and
are recorded as conditional on the compact-critical atlas, which is the
content of Construction Problem~1 in the Vol~II frontier
\textup{(}\texttt{chiral-bar-cobar-vol2:CLAUDE.md}~\S 9, item 1: the
$K3 \times E$ Drinfeld-double construction with $\protectedPfaff{
\fD_X} = \Delta_5$\textup{)}.
\end{proof}

\begin{remark}[\label{rem:G-eq-D-Yplus-not-Wonepinf}$\Wonepinf$ is a
Fock representation of $G(\C^3)$, not equal to $\CoHA(\C^3)$]
The equality $\CoHA(\C^3) = \Wonepinf$ is forbidden
\textup{(}Vol~II voice table, item 8\textup{)}. The correct relation
is the chain
\[
 \CoHA(\C^3) \;=\; \Yplus{\widehat{\fgl}_1}
 \;\xrightarrow[\textup{(LD1)--(LD5)}]{\Drinfdouble{\,\cdot\,}}\;
 G(\C^3) \;=\; Y(\widehat{\fgl}_1)
 \;\xrightarrow[\text{Gaiotto--Rap\v{c}\'ak}]{\text{vacuum action}}\;
 \Wonepinf,
\]
where the last arrow is the Fock-module action of the affine Yangian on
$\Wonepinf$ \textup{(}Proch\'azka--Rap\v{c}\'ak; Gaiotto--Rap\v{c}\'ak;
cf.\ \cref{thm:drinfeld-center-coha}\textup{(ii)}\textup{)}.
$\Wonepinf$ is a representation of $G(\C^3)$, not $G(\C^3)$ itself.
\end{remark}

\begin{theorem}[Drinfeld center of the CoHA]
\label{thm:drinfeld-center-coha}
Let $Y^+(\widehat{\fgl}_1)$ denote the positive half of the affine Yangian (Schiffmann--Vasserot). Then:
\begin{enumerate}[label=(\roman*)]
 \item The Drinfeld center of its representation category is the braided monoidal category of restricted modules for the Drinfeld double, hence for the full affine Yangian:
 \[
 \cZ(\Rep^{E_1}(Y^+(\widehat{\fgl}_1))) \;\simeq\; \Rep^{E_2}_{\mathrm{res}}(Y(\widehat{\fgl}_1)).
 \]
 \item The completed affine Yangian acts on the $W_{1+\infty}$ vacuum/Fock modules by the Gaiotto--Rap\v{c}\'ak/Proch\'azka--Rap\v{c}\'ak evaluation. This is an action/evaluation statement, not an equivalence of the whole center with an unspecified $\Rep(W_{1+\infty})$ category.
 \item The completed Drinfeld double has graded dimension
 \[
 \mathrm{ch}\,D(Y^+) = M(q)^2,
 \]
where $M(q) = \prod_{n \geq 1} 1/(1-q^n)^n$ is the MacMahon
function.  A chosen $W_{1+\infty}$ vacuum/Fock evaluation has
separate trace factor
\[
 \mathrm{ch}\,\cF_{W_{1+\infty}}=P(q)=\prod_{n \geq 1}(1-q^n)^{-1}.
\]
The product \(M(q)^2P(q)\) is therefore an evaluated trace, not an
intrinsic character of the Drinfeld centre.
\end{enumerate}
\end{theorem}

\begin{proof}
Part~(i) proceeds in two steps. First, the general categorical fact (Majid, 1991; see also Etingof--Gelaki--Nikshych--Ostrik, \emph{Tensor Categories}, Proposition~7.13.8): for a bialgebra $H$ over a field $k$, the Drinfeld center $\cZ(\Rep(H))$ is equivalent as a braided monoidal category to $\Rep(D(H))$, where $D(H) = H \bowtie H^{*\mathrm{cop}}$ is the Drinfeld double. The equivalence sends a half-braiding $(V, \sigma_{V,-})$ to the $D(H)$-module structure determined by $\sigma$. Second, for $H = Y^+(\widehat{\fgl}_1)$ the positive half of the affine Yangian, the Drinfeld double $D(Y^+) = Y^+ \bowtie (Y^+)^{*\mathrm{cop}}$ is identified with the full affine Yangian $Y(\widehat{\fgl}_1)$ (Schiffmann--Vasserot, 2013: the negative generators $Y^-$ are the graded dual of $Y^+$ with reversed coproduct). The $E_2$ braiding on $\Rep(Y)$ is induced by the universal $R$-matrix of the Yangian, which satisfies the YBE (Theorem~\ref{thm:yang-r-matrix-ybe} below). The passage from $\Eone$ to $\Etwo$ is the Drinfeld center functor $\cZ \colon \Eone\text{-}\mathrm{Cat} \to \Etwo\text{-}\mathrm{Cat}$ (this is the categorified analogue of the derived center; Drinfeld center $\cZ(\Rep^{\Eone}(A))$ is a braided monoidal \emph{category}, not a chiral algebra).

Part~(ii) is the standard affine-Yangian/$W_{1+\infty}$ evaluation
dictionary: the completed mode algebra acts on the
$W_{1+\infty}[\lambda]$ vacuum/Fock module. The theorem uses only this
module-level action; it does not assert a categorical equivalence with
all $W$-modules.

Part~(iii): the Drinfeld double character is
\(\mathrm{ch}(Y^+) \cdot \mathrm{ch}(Y^-) = M(q)^2\) (the MacMahon
function squared, since both \(Y^+\) and \(Y^-\) have character \(M(q)\)
by the Schiffmann--Vasserot computation of the graded dimension of the
shuffle algebra).  The vacuum module of \(W_{1+\infty}\) contributes
\(P(q)\) only after the evaluation representation has been chosen.
\end{proof}

\begin{theorem}[Yang \texorpdfstring{$R$}{R}-matrix: YBE and unitarity]
\label{thm:yang-r-matrix-ybe}
Let $h_1, h_2, h_3$ satisfy $h_1 + h_2 + h_3 = 0$ and let
\[
 g(u) = \frac{(u - h_1)(u - h_2)(u - h_3)}{(u + h_1)(u + h_2)(u + h_3)}
\]
be the structure function of the affine Yangian $Y(\widehat{\fgl}_1)$.
The \(R\)-matrix on
\(\mathrm{Fock}_2 \otimes \mathrm{Fock}_2\), where
\(\mathrm{Fock}_2=\langle |\mathrm{row}\rangle,|\mathrm{col}\rangle\rangle\)
is the charge-$2$ space indexed by the partitions \((2)\) and
\((1,1)\), is a \(4\times4\) matrix
\(R(u) \in \End(\C^2 \otimes \C^2)((u))\) satisfying:
\begin{enumerate}[label=(\roman*)]
 \item \emph{Quantum Yang--Baxter equation}:
 \[
 R_{12}(u-v)\, R_{13}(u)\, R_{23}(v) = R_{23}(v)\, R_{13}(u)\, R_{12}(u-v)
 \]
 in $\End(\C^2 \otimes \C^2 \otimes \C^2)((u,v))$. This is verified
 entry-by-entry on the \(8\)-dimensional triple tensor product.
 \item \emph{Unitarity}: $R_{12}(u)\, R_{21}(-u) = \mathrm{id}$.
 \item \emph{$E_2$ nontriviality}: $R(u) \neq \mathrm{id}$ whenever $h_1 h_2 h_3 \neq 0$, i.e., whenever the Calabi--Yau deformation parameters are generic. At the self-dual point $(h_1, h_2, h_3) = (1, 0, -1)$, the structure function $g(u) = 1$ identically ($\sigma_3 = h_1 h_2 h_3 = 0$), the $R$-matrix degenerates to the identity, and the $E_2$ braiding becomes symmetric (the braiding data is trivial, though the operadic level remains $E_2$).
 \item \emph{Classical limit}: $R(u) = 1 - \sigma_2/u + O(1/u^2)$ as $u \to \infty$, where $\sigma_2 = h_1 h_2 + h_1 h_3 + h_2 h_3$. The classical $r$-matrix is $r(u) = -\sigma_2/u$.
\end{enumerate}
\end{theorem}

\begin{proof}
The \(4\times4\) matrix \(R(u)\) gives, after placement in the three
tensor factors, an equality of the \(2^6 = 64\) rational matrix entries of
$R_{12}R_{13}R_{23}$ and $R_{23}R_{13}R_{12}$ in
$u,v,h_1,h_2,h_3$.  Unitarity is ordinary matrix multiplication.  For
the classical limit the structure function satisfies
$g(u) = 1 + O(u^{-2})$, because the $u^{-1}$ coefficient is
$-2(h_1+h_2+h_3)=0$; the $R$-matrix on
$\mathrm{Fock}_2\otimes \mathrm{Fock}_2$ nevertheless has a $u^{-1}$
term from the off-diagonal exchange of boxes.  Hence
$R(u)=1-\sigma_2/u+O(u^{-2})$ and $r(u)=-\sigma_2/u$.
\end{proof}


\begin{remark}[Dimofte's \texorpdfstring{$K$}{K}-matrix and CY shadow depth]
\label{rem:dimofte-k-matrix-cy}
In the PIRSA lecture series (Dimofte, PIRSA 25110067), Dimofte isolates a single chiral operator $K_\cA(z)$ which controls the deviation of the boundary coproduct of an HT boundary algebra from the primitive (Hopf) coproduct. Its explicit form is
\begin{equation}
\label{eq:dimofte-k-matrix-cy}
K_\cA(z) \;=\; z^{\alpha_0}\,\exp\!\Biggl(\sum_{n>0} \frac{\alpha_{-n}}{-n}\, z^{n}\Biggr)\,\exp\!\Biggl(\sum_{n>0} \frac{\alpha_{n}}{-n}\, z^{-n}\Biggr),
\end{equation}
where $\{\alpha_n\}_{n \in \Z}$ are the modes of a free boson identified with the $\mathfrak{u}(1)^r$ Cartan factor of the boundary algebra and $\alpha_0$ is the zero-mode charge. The operator $K_\cA(z)$ lives inside $\cA$, not over it. The canonical Vol~II discussion of equation~\eqref{eq:dimofte-k-matrix-cy}, its diagnostic interpretation, and its interaction with the shadow-depth classification is in the ordered bar chapter of Vol~II (\emph{Vol~II}, remark \ref{rem:dimofte-k-matrix-cy} in \cite{VolII}), together with the BPS-shadow-depth parallel in the HT bulk-boundary chapter. This remark records the CY specialization of that Vol~II content.

Let $\cC$ be a smooth proper Calabi--Yau category of dimension $d \in \{2, 3\}$. For $d = 2$ the CY-to-chiral correspondence $\Phi_2$ of Theorem~CY-A$_2$ produces an $E_2$-chiral algebra $A_\cC = \Phi_2(\cC)$ and the $K$-matrix $K_{A_\cC}(z)$ is well defined by equation~\eqref{eq:dimofte-k-matrix-cy} applied to the boundary algebra of the associated 3d HT theory. For $d = 3$ the same prescription is available only when Theorem~CY-A$_3$ supplies the framed object-level specialisation $A_\cC^{(\Sigma_2,C)}=\Phi_3^{(\Sigma_2,C)}(\cC)$ under H1--H4 and fixed admissible data; the chain-level $\mathbb{S}^3$-framing is an assumption or a verified toric/formal witness, not automatic.

The diagnostic interpretation of $K_{A_\cC}(z)$ refines on the CY side into a statement about the BPS hull of $\cC$. Write $\cC^{\mathrm{BPS}} \subset \cC$ for the subcategory of BPS objects (semistable objects at a chosen stability condition), and let $\Phi_d(\cC^{\mathrm{BPS}})$ denote the image of this subcategory under the CY-to-chiral correspondence $\Phi_d$ (at the CY dimension $d$ of $\cC$). Then $K_{A_\cC}(z)$ is determined by $\Phi_d$ applied to the BPS hull: the Cartan zero-mode $\alpha_0$ records the rank of the $\mathfrak{u}(1)^r$ factor, and the higher modes $\alpha_{\pm n}$ ($n \geq 1$) encode the BPS generating-function corrections. Here ``$\Phi$ applied to the BPS hull'' refers to the CY-to-chiral correspondence $\{\Phi_d\}$ at the appropriate $d$.
\begin{itemize}
\item For toric CY$_3$ without compact $4$-cycles (Chapter~\ref{ch:toric-coha}), $K_{A_\cC}(z)$ is the generating function of Donaldson--Thomas invariants on the fiber of the toric quiver: the mode $\alpha_{-n}$ scales as the DT partition function restricted to classes with $n$ BPS quanta, and the Schiffmann--Vasserot/RSYZ identification of Theorem~\ref{thm:drinfeld-center-coha} realizes the identification $K_{A_\cC}(z) = Z^{\mathrm{DT}}_{\mathrm{fiber}}(z)$.
\item For $K3 \times E$ (Chapter~\ref{ch:k3-times-e}), the $K$-matrix specializes to the Fourier coefficients of the Igusa cusp form $\Phi_{10}^{\mathrm{un}}$: the Fourier--Jacobi expansion $\Phi_{10}^{\mathrm{un}} = \sum_{m \geq 1} \phi_m(\tau,z)\, p^m$ records the BKM root multiplicities on successive modes, and these coefficients are precisely the data encoded by equation~\eqref{eq:dimofte-k-matrix-cy}. The weight of $\Phi_{10}^{\mathrm{un}}$ is $10 = 2 \cdot \kappa_{\mathrm{BKM}}(\Delta_5)$, where $\kappa_{\mathrm{BKM}}(\Delta_5) = 5$ is the Borcherds--Kac--Moody modular characteristic (see Table of Chapter~\ref{ch:k3-times-e}).
\end{itemize}

The Vol~I shadow-depth classification $\{\mathbf{G}, \mathbf{L}, \mathbf{C}, \mathbf{M}\}$ coincides, under $\Phi$, with the $K$-matrix complexity of the CY boundary algebra:
\begin{center}
\renewcommand{\arraystretch}{1.2}
\begin{tabular}{lll}
\toprule
Class & $K_{A_\cC}(z)$ & $d_{\mathrm{alg}}(A_\cC)$ \\
\midrule
$\mathbf{G}$ (trivial CY fiber, $\mathfrak{u}(1)^r$) & $1$ & $0$ \\
$\mathbf{L}$ (rank-one CY Heisenberg deformation) & polynomial, simple pole & $1$ \\
$\mathbf{C}$ (toric CY$_3$, $\beta\gamma$-type) & polynomial, double pole & $2$ \\
$\mathbf{M}$ ($K3$, $W_N$-type) & infinite expansion & $\infty$ \\
\bottomrule
\end{tabular}
\end{center}
Equivalently: $K_{A_\cC}(z) = 1$ iff the boundary coproduct is primitive iff $d_{\mathrm{alg}} = 0$ iff $\cC$ is class $\mathbf{G}$. The correspondence $K_{A_\cC}(z) = 1 \Longleftrightarrow \mathbf{G} \Longleftrightarrow d_{\mathrm{alg}} = 0$ is the CY analogue of the Vol~II biconditional.

The $K$-matrix modifies the coproduct, not the product, so the $r$-matrix vanishing check does not apply directly to $K_{A_\cC}(z)$. The corresponding $r$-matrix check is the affine one: the classical $r$-matrix $k_{\mathrm{ch}}\,\Omega/z$ vanishes at $k_{\mathrm{ch}} = 0$, in which case $A_\cC$ collapses to the trivial Heisenberg and $K_{A_\cC}(z) = 1$, consistent with class $\mathbf{G}$.
\end{remark}


\section{Comparison with the Maulik--Okounkov \texorpdfstring{$R$}{R}-matrix}
\label{sec:mo-rmatrix-comparison}

The Maulik--Okounkov stable-envelope construction provides an
independent geometric source of \(R\)-matrices on symplectic
resolutions, in particular on Nakajima quiver varieties such as
\(\mathrm{Hilb}^n(\C^2)\).  A Calabi--Yau threefold CoHA uses
critical Hall and shuffle data instead.  Thus the MO comparison is a
surface-chart comparison; it is not a stable-envelope construction on
\(\mathrm{Hilb}^n(\C^3)\).

\begin{theorem}[MO/chiral \texorpdfstring{$R$}{R}-matrix comparison on symplectic surface charts]
\label{thm:mo-chiral-rmatrix}
Let \(S\) be a toric symplectic surface chart, an ADE surface chart, or
a Nakajima quiver variety chart for which the Maulik--Okounkov
stable-envelope \(R\)-matrix \(R^{\mathrm{MO}}(u)\) is defined.  Let
\(R^{\mathrm{ch}}(u)\) be the chiral \(R\)-matrix from the \(\Etwo\)-bar
complex of the Drinfeld double \(D(Y^+)\) of the corresponding Hall
positive half (Construction~\ref{constr:cy-r-matrix}).  The \(\Etwo\)-bar
complex here is constructed on the Drinfeld double, not on the CoHA
\(Y^+\) itself, which is natively \(\Eone\); the \(\Etwo\)-braiding
arises from the Drinfeld centre, Theorem~\ref{thm:drinfeld-center-coha}.
Then:
\begin{enumerate}[label=(\roman*)]
 \item Both \(R\)-matrices act on the same fixed-point Fock space
 \[
   \cF_S=\bigoplus_n K_T(\mathrm{Hilb}^n(S)^T)
 \]
 in the Hilbert-scheme surface case, and on the analogous Nakajima
 fixed-point representation in the quiver-variety case.
 \item On each partition $\lambda$, both produce the same diagonal eigenvalue:
 \[
 R^{\mathrm{MO}}_{\lambda\lambda}(u) = R^{\mathrm{ch}}_{\lambda\lambda}(u) = \prod_{s \in \lambda} g(u + c(s)),
 \]
 where $c(s) = h_1 i + h_2 j$ is the content of box $s = (i,j) \in \lambda$.
 \item Both satisfy the same Yang--Baxter equation, unitarity, and crossing symmetry.
\end{enumerate}
\end{theorem}

\begin{proof}
Seven independent verification paths confirm the identification:
\begin{enumerate}
 \item \emph{Direct diagonal computation}: the eigenvalues $\prod_{s \in \lambda} g(u + c(s))$ are computed from the structure function in both constructions.
 \item \emph{$\mathrm{Hilb}^2(\C^2)$ check}: at $n = 2$, the MO construction uses stable envelopes on $\mathrm{Hilb}^2(\C^2)$ (two chambers in the equivariant torus); the resulting $2 \times 2$ $R$-matrix matches the Yangian $R$-matrix.
 \item \emph{Unitarity}: $R_{12}(u) R_{21}(-u) = \mathrm{id}$ holds for both.
 \item \emph{YBE}: both satisfy the same quantum Yang--Baxter equation (Theorem~\ref{thm:yang-r-matrix-ybe}).
 \item \emph{Crossing symmetry}: the crossing relation $R_{12}(u)^{t_1} R_{12}(-u - \kappa_{\mathrm{cr}})^{t_1} = f(u) \cdot \mathrm{id}$ holds for both, with the same scalar function $f(u)$; here $\kappa_{\mathrm{cr}}$ is the Yangian crossing parameter (distinct from the modular characteristic $\kappa_{\mathrm{ch}}$).
 \item \emph{Classical $r$-matrix}: both yield $r(u) = -\sigma_2 / u$ in the $u \to \infty$ limit.
 \item \emph{Schiffmann--Vasserot specialization}: at the SV values $(h_1, h_2, h_3) = (1, -1-\epsilon, \epsilon)$, both $R$-matrices specialize to the CoHA shuffle algebra braiding.
\end{enumerate}
\end{proof}

\begin{remark}
\label{rem:mo-comparison-scope}
For \(\C^3\), the framed Donaldson--Thomas space
\(\mathrm{Hilb}^n(\C^3)\) supplies a Fock module for the critical
CoHA/shuffle algebra; it is not a Nakajima quiver variety carrying the
Maulik--Okounkov stable envelope.  In the \(K3\times E\) lane the MO
computation is local: it is made on ADE/Kummer surface charts and then
enters the compact theory only through the Hall/DT comparison datum.
For a general CY threefold without torus action, the MO construction is
not directly available, and the chiral \(R\)-matrix from
Construction~\ref{constr:cy-r-matrix} is the primary object.
\end{remark}

\section{Quantum chiral algebras from holomorphic Chern--Simons theory}
\label{sec:qca-from-qg}

Holomorphic Chern--Simons theory on a CY$_d$ manifold $Y$ is the physical realisation of stage~1 of the two-stage factorisation
\[
 \mathrm{CY\text{-}cat}_d
 \;\xrightarrow{\;\PhiFA_d\;}
 \EdHolFA(Y)
 \;\xrightarrow{\;\SpCh_{\Sigma_{d-1},C}\;}
 \EnHolFA(C)
\]
of Definition~\ref{def:phi-fa-and-sp} and Theorem~\ref{thm:phi-two-stage-factorisation}. Stage~1 ($\PhiFA_d$) outputs the canonical $E_d$-holomorphic factorization algebra on $Y$; stage~2 ($\SpCh_{\Sigma_{d-1},C}$) specialises it to an $E_n$-chiral algebra on a curve $C \subset Y$ by pushing forward along the normal cycle $\Sigma_{d-1} = N_{C/Y}$. The Omega-background parameters $(h_1,\ldots,h_d)$ satisfying $\sum h_i = 0$ are equivariant weights of the chosen torus chart and of the normal directions; they deform the factorisation model and survive its Stage-$2$ specialisation. Three regimes of $d$ realise this two-stage picture concretely.

\begin{construction}[Holomorphic Chern--Simons hierarchy]
\label{constr:hol-cs-hierarchy}
Let $\frakg$ be a finite-dimensional Lie algebra equipped with an invariant bilinear form $\langle -, - \rangle_{\frakg} \colon \frakg \otimes \frakg \to \C$. Holomorphic Chern--Simons theory with gauge algebra $\frakg$ on a complex manifold $M$ has action
\[
 S_{\mathrm{hCS}}(A) \;=\; \frac{1}{2} \int_M \Omega \wedge \langle A \wedge \dbar A + \tfrac{2}{3} A \wedge A \wedge A \rangle_{\frakg},
\]
where $A \in \Omega^{0,1}(M; \frakg)$ is a $(0,1)$-connection and $\Omega$ is a holomorphic volume form on $M$ (when $M$ is Calabi--Yau) or a partial volume form (when $M$ is a product). The three regimes:
\begin{enumerate}[label=\textup{(\roman*)}]
 \item \emph{3d holomorphic CS on $\Sigma \times \R$} (Costello~2007, Costello--Gwilliam~2021):
 the boundary on $\Sigma$ produces the affine Kac--Moody vertex algebra $V_k(\frakg)$ at level $k$ determined by $\langle -, - \rangle_{\frakg}$. Its Beilinson--Drinfeld factorization is symmetric only on disjoint configurations; the collision OPE carries the non-commutative current algebra. The ordered $\Eone$ refinement used by Vol.~I and the Vol.~II Swiss-cheese structure records that collision sector, while PVA descent applies only after the Li-graded classical limit and the finite-type hypotheses named in Vol.~II.

 \item \emph{5d holomorphic CS on $\C^2 \times \R$} (Costello~2013, ``Supersymmetric gauge theory and the Yangian''):
 the Omega-background parameters on $\C^2$ produce the Yangian
 vertex-algebra boundary object \(Y^{\mathrm{VOA}}_\varepsilon(\frakg)\)
 in the Costello--Gaiotto--Yagi simply-laced setting.  The affine
 Yangian \(Y(\widehat{\fgl}_1)\) and its higher-rank affine or toroidal
 analogues require the additional Hall, stable-envelope, or
 double/centre input appropriate to the chosen chart.  The positive
 half $Y^+(\widehat{\fgl}_1)$ is the CoHA; it is associative, not a
 chiral algebra.  The full affine Yangian is recovered by the Drinfeld
 double where the Hall pairing is available.

 \item \emph{6d holomorphic theory on $\C^3$} (Costello--Li holomorphic BV locality and Costello's Omega-background construction):
 the observables on $\C^3$, when the one-loop anomaly is cancelled, carry the holomorphic $E_3$ factorisation structure of Definition~\ref{def:cy3-hcs-bv-complex}. Projection to $\C^2 \subset \C^3$ gives an $\Etwo$-chiral shadow; projection to $C \subset \C^3$ gives an $\Eone$-chiral shadow. For $\frakg = \fgl_1$, the finite Rees hCS--Hall comparison identifies the ordered Hall layer at fixed DWR/Ran bounds.  A quantum toroidal algebra \(U_{q,t}(\widehat{\widehat{\fgl}}_1)\) is an expected curve-level projection after completion and centre/double operations, not a direct consequence of Costello--Francis--Gwilliam~\cite{CFG2026}.
\end{enumerate}
\end{construction}

\begin{remark}[6d theory is not standard holomorphic CS]
\label{rem:6d-not-lagrangian}
The naive field $A\in\Omega^{0,1}(\C^3,\mathfrak g)$ alone does not contain all BV components: $\Omega\wedge A\wedge\dbar A$ has type $(3,3)$ only after the superfield multiplication is read in the full BV complex $\Omega^{0,\bullet}(\C^3,\mathfrak g)[1]$ and the $\Omega^{0,3}$ component is taken. Thus the correct local hCS presentation is the BV superfield action of Definition~\ref{def:cy3-hcs-bv-complex}, not a bare $(0,1)$ connection calculation. Costello--Francis--Gwilliam~\cite{CFG2026} proves the ordinary $3$d Chern--Simons $E_3$ factorisation-homology theorem. The $6$d hCS--Hall input used here is instead the finite DWR/Ran Rees comparison of Proposition~\ref{prop:qca-finite-rees-layer-separation}; ordinary compact critical CoHA still requires the vanishing-cycle realisation and pro-descent beyond fixed bounds. The mixed-holomorphic-topological local model on $\R^2_{\mathrm{top}} \times \C^2_{\mathrm{hol}}$ with Hamiltonian BF sector~\cite{LorgatMixedHTStrings} is the level-$1$ physical realisation lane of Stage-$1$ at toric loci; the global de Rham obstruction $\mathrm{Ob}^{\mathrm{BMK\text{-}curr}}_N$ named in~\cite{LorgatMixedHTStrings} Chapter~\ref{ch:cy-to-chiral} lines~$3207$--$3266$ blocks the formal-Darboux $\Rightarrow$ global compact theory promotion outside the verified locus.
\end{remark}

\begin{theorem}[Three Atiyah-sourced cocycles on compact CY\texorpdfstring{$_3$}{-3}]
\label{thm:three-atiyah-cocycles-qca}
For $X$ a compact Calabi--Yau threefold with $\mathrm{At}(T_X) \in H^1(X, \Omega^1_X \otimes \End T_X)$ the Atiyah class of the holomorphic tangent bundle, three cocycles on $X$ arise and enter the hCS presentation of $\PhiFA_3(\Perf(X))$ at distinct orders:
\begin{enumerate}
\item \emph{Generator.} $\At(T_X)$ is the binary bracket $\ell_2$ of the Kapranov $L_\infty$-algebra on $T_X[-1]$ (Kapranov 1999 Proposition~4.4). Under the C\u ald\u araru--C\u ald\u araru--Tu twisted HKR~\cite{CaldararuCaldararuTu2014}, $\At(T_X)$ is cohomologically represented by the defect from HKR being a Gerstenhaber morphism.
\item \emph{Tree-level Yukawa.} $Y_3 \in H^{0,3}(X) \simeq \mathrm{Sym}^3 H^1(X, T_X)^\vee$ is the Kapranov $\ell_3^{\min}$-bracket,
\[
Y_3(v_1, v_2, v_3) \;=\; \int_X \Omega_X \wedge \mathrm{At}(v_1) \cup \mathrm{At}(v_2) \cup \mathrm{At}(v_3),
\]
and enters the classical hCS action as the cubic correction $\int_X \Omega_X \wedge Y_3 \wedge \langle \cA, \cA, \cA \rangle$ beyond the $\bar\partial$-kinetic plus Chern--Simons cubic. Physically: the Bershadsky--Cecotti--Ooguri--Vafa tree-level three-point Yukawa coupling on the moduli of complex structures.
\item \emph{One-loop determinant-line/gravitational datum.} The BCOV
one-loop contribution is a determinant-line or gravitational anomaly
datum, with coefficient controlled by characteristic-class integrals
such as the Euler characteristic. It is not a class
$\alpha_{\mathrm{BCOV}}\in H^{0,1}(X)$ and it is not
$\frac{\chi(X)}{24}[\Omega_X]^{0,1}$; a holomorphic volume form on a
CY$_3$ has no $(0,1)$ component. This datum is also separate from the
local hCS gauge anomaly, whose CY$_3$ slot is quartic as in
Proposition~\ref{prop:cy3-hcs-quartic-anomaly-slot}.
\end{enumerate}
The tree-level Atiyah/Yukawa data, the local quartic hCS gauge anomaly,
and the BCOV determinant-line/gravitational datum coexist on compact
non-flat $X$, but they live in different complexes; the three are
distinct cohomology classes.
\end{theorem}

\begin{remark}[Geometry at \texorpdfstring{$X = K3 \times E$}{X = K3 x E}]
\label{rem:three-cocycles-at-K3xE-qca}
At $X = K3 \times E$ the K\"unneth decomposition gives
$\dim H^{0,3}(K3 \times E) = \dim H^{0,2}(K3) \cdot \dim H^{0,1}(E) = 1$
and $\dim H^{0,1}(K3 \times E) = 1$. These are ambient Hodge slots,
not automatic obstruction classes. On the formal $K3\times E$ locus the
elliptic tangent Atiyah class is zero, the K3 factor has no
$\Omega^3_{K3}$ receptacle for the Kuranishi cubic, and the coefficient
$\chi(K3 \times E)=\chi(K3)\chi(E)=24\cdot0=0$ sets the BCOV one-loop
coefficient to zero; this matches the K\"unneth-multiplicative
total-space reading $\kcat(K3 \times E)=\chi(\cO_{K3\times E})=0$.
\end{remark}

\begin{theorem}[Boundary chiral algebra at \texorpdfstring{$d = 5$}{d = 5}: the affine Yangian]
\label{thm:5d-boundary-yangian}
Let $h_1, h_2, h_3 \in \C$ satisfy $h_1 + h_2 + h_3 = 0$. The observables of 5d holomorphic Chern--Simons theory on $\C^2_{h_1, h_2} \times \R$ with gauge algebra $\fgl_1$, projected to a holomorphic curve $C \subset \C^2$, form an $\Eone$-chiral algebra on $C$ isomorphic to the affine Yangian $Y(\widehat{\fgl}_1)$ with structure function
\[
 g(u) = \frac{(u - h_1)(u - h_2)(u - h_3)}{(u + h_1)(u + h_2)(u + h_3)}.
\]
The completed affine Yangian acts on the $W_{1+\infty}$ vacuum/Fock
modules by the Proch\'azka--Rap\v{c}\'ak/Gaiotto--Rap\v{c}\'ak
evaluation dictionary; the statement used here is this completed
mode-action, not an identification of $Y^+$ or its full representation
category with the vertex algebra. References: Costello~2013;
Proch\'azka--Rap\v{c}\'ak~2019.
\end{theorem}

\begin{conjecture}[Boundary chiral algebra at \texorpdfstring{$d = 6$}{d = 6}: quantum toroidal from factorization]
\label{conj:6d-boundary-toroidal}
\ClaimStatusConjectured
The $E_3$-holomorphic factorization algebra of 6d hCS observables on
$\C^3_{h_1,h_2,h_3}$, with $h_1+h_2+h_3=0$, projected to an
$\Eone$-chiral algebra on a curve $C\subset\C^3$, is expected to be the
quantum toroidal algebra $U_{q,t}(\widehat{\widehat{\fgl}}_1)$ with
parameters $q=e^{2\pi i h_1/h_3}$ and $t=e^{-2\pi i h_2/h_3}$. The
construction is conditional on the 6d hCS BV--BRST complex and the
Costello--Gwilliam factorisation-algebra comparison. The intermediate
$\Etwo$-projection to $\C^2 \subset \C^3$
should recover the affine Yangian $Y(\widehat{\fgl}_1)$ of
Theorem~\ref{thm:5d-boundary-yangian}; this compatibility between the 6d
projection and the independent 5d construction is itself part of the
conjecture. The $E_3$ structure on $\C^3$ is the \emph{source} of the
second affinization (the second hat in $\widehat{\widehat{\fgl}}_1$); the
first affinization comes from the $\Etwo$ factorization on $\C^2$.
\end{conjecture}

\begin{proposition}[Non-abelian hCS pushforward on \texorpdfstring{$K3\times E$}{K3 x E} and the Gritsenko \texorpdfstring{$\fg_{\Delta_5}$}{g-Delta-5}]
\label{prop:hcs-pushforward-k3e-delta5}
Let $X = K3 \times E$ with holomorphic volume form $\Omega_X = \omega_{K3} \wedge d\tau_E$ and projection $\pi \colon X \to E$. Let $\cE_{\hCS}^{\fg_{\Delta_5},\mathrm{re}}$ denote the BV field complex of $6$D holomorphic Chern--Simons on $X$ with gauge datum restricted to the real-root subalgebra $\fg_{\Delta_5}^{\mathrm{re}} = \bigoplus_{\alpha \in \Delta^{\mathrm{re}}} \fsl_2^{(\alpha)}$ (Gritsenko--Nikulin 1998 Theorem~2.1, rank-$3$ hyperbolic lattice $\Lambda^{2,1}_{II}\subset\Lambda^{3,2}$). Then:
\begin{enumerate}[label=\textup{(\roman*)}]
 \item \emph{Anomaly split.} The $1$-loop anomaly factorises as
 \[
 A_{\hCS}^{(1)}(\fg_{\Delta_5}^{\mathrm{re}}) \;=\; A_{\mathrm{w.f.}} + A_{\mathrm{anom}}, \qquad A_{\mathrm{w.f.}} = -\frac{C_2|_{\mathrm{re}}}{(2\pi)^3},
 \]
 with $A_{\mathrm{anom}}$ the four-vertex box wheel weighted by the
 symmetric quartic adjoint trace
 $q^{abcd} = \operatorname{tr}_{\mathrm{ad}}(T^{(a}T^bT^cT^{d)})$
 (Definition~\ref{def:6d-admissible-gauge-algebra}). The
 kinetic-renormalisation piece $A_{\mathrm{w.f.}}$ is scheme-dependent
 and absorbed by the Costello--Li propagator counter-term. For each
 real-root $\fsl_2^{(\alpha)}$ taken as gauge datum, the primitive
 quartic part of $q^{abcd}$ vanishes --- $A_1$ admits no primitive
 quartic Casimir
 (Lemma~\ref{lem:quartic-casimir-classification}) --- leaving a
 multiple of $(C_2)^2$, which on $K3 \times E$ is cancelled by the
 integrated vanishing of item~(ii); the cross-block couplings between
 distinct real-root copies belong to the compact extension of
 Problem~\ref{op:hcs-pushforward-k3e-delta5-compact-bv}.
 \item \emph{Quartic obstruction coefficient.} The compact-CY$_3$ quartic anomaly polynomial $I_4 \in \Sym^4(\fg_{\Delta_5}^\vee)^{\fg_{\Delta_5}}$ of Conjecture~\ref{conj:quintic-hcs-kanom} integrates to
 \[
 \int_X \mathrm{td}(T_X) \wedge I_4(F_{\nabla}) \;=\; \frac{\chi_{\mathrm{top}}(X)}{24} \cdot \mathrm{tr}_{\mathrm{ad}}(F^4) \;=\; 0,
 \]
 since $\chi_{\mathrm{top}}(K3 \times E) = \chi_{\mathrm{top}}(K3)\,\chi_{\mathrm{top}}(E) = 24 \cdot 0 = 0$. The obstruction coefficient vanishes on $K3 \times E$ without invoking a Green--Schwarz mechanism.
 \item \emph{Pushforward identification.} Under the fibre integration $\pi_* \colon \PhiFA_3(\Perf(X)) \to \PhiFA_1(\Perf(E))$ specialised along $\Sigma_2 = K3 \subset X$, the renormalised hCS factorisation algebra descends to the $\Eone$-chiral algebra on $E$
 \[
 \pi_* \Obs_{\hCS}^{\fg_{\Delta_5}^{\mathrm{re}}}(X) \;\simeq\; V^{\mathrm{ch}}(\fg_{\Delta_5}^{\mathrm{re}}),
 \]
 where the right-hand side is the real-root subenvelope of the
Gritsenko--Nikulin BKM superalgebra. Its Weyl factor is
 \[
 e^{-\rho_W}\prod_{\alpha\in\Delta^{\mathrm{re}}_+}
 \bigl(1-e^{-\alpha}\bigr)^{-1},
 \]
and its Borcherds-weight normalisation is fixed by the $N=1$ identity
$\kBKM(\Delta_5)=c_1(0)/2=5$. The full imaginary-root multiplicity
extension is Theorem~\ref{thm:hcs-pushforward-imaginary-roots-singular-theta}.
 \item \emph{Composite resolvent.} The Costello--Li propagator on $X$ factorises as $P_X = P_{K3}^{(0,1)} \boxtimes P_E^{(0,1)}$; fibre-integrating $P_{K3}^{(0,1)}$ against $\omega_{K3}$ extracts the K3-twined elliptic genus, and the composite $\pi_* P_X = \eta(\tau_E)^{24}$ reproduces the weight-$12$ factor of the unnormalised square $\Phi_{10}^{\mathrm{un}} = \Delta_5^2$ via the Kawai--Nam/DMVV product.
\end{enumerate}
\end{proposition}

\begin{proof}[Sketch]
(i): The splitting $A^{(1)} = A_{\mathrm{w.f.}} + A_{\mathrm{anom}}$ is Costello--Li 2016 Proposition~3.7; the kinetic term is scheme-dependent rather than cohomological. On a $3$-fold the anomalous wheel is the four-vertex box with the quartic adjoint trace, not the three-vertex triangle (the triangle is the anomalous wheel in holomorphic dimension~$2$). Restriction to $\fg_{\Delta_5}^{\mathrm{re}}$: real-root subalgebras are simply-laced rank-$1$ ($\fsl_2$) copies because $\fg_{\Delta_5}$ has real roots of squared length~$2$ in the hyperbolic core $\Lambda^{2,1}_{II}$ (Gritsenko--Nikulin 1998 §2). For $\fsl_2 = A_1$, $(\Sym^4 \fsl_2^\vee)^{\fsl_2} = \C\,(C_2)^2$ with no primitive quartic generator (exponents $\{1\}$; Lemma~\ref{lem:quartic-casimir-classification}), so the per-block box tensor is a pure $(C_2)^2$-multiple; its integrated coefficient on $K3 \times E$ vanishes by item~(ii).

(ii): Grothendieck--Riemann--Roch for the anomaly polynomial (Costello 2013 §9) yields the total coefficient $\chi_{\mathrm{top}}(X)/24$; Kunneth gives $\chi_{\mathrm{top}}(K3 \times E) = 24 \cdot 0 = 0$.

(iii)--(iv): The pushforward $\pi_*$ is the Costello--Gwilliam fibre-integration functor on factorisation algebras (Costello--Gwilliam 2021 Vol II §3.5). Fibre-integrating the BV action $\int_X \Omega_X \wedge \mathrm{tr}(A \bar\partial A + \tfrac23 A^3)$ over $K3$ produces the elliptic-genus-weighted 2d chiral theory on $E$ at the real-root subsystem; the denominator identity is Gritsenko--Nikulin 1998 Theorem~2.1 with Borcherds-lift coefficient extraction. The Borcherds-weight component is the $N=1$ case of Theorem~\ref{thm:bl-master-kappa-bkm-chain-level}; it gives $\kBKM(\Delta_5)=c_1(0)/2=5$. The composite-resolvent identity (iv) is the unnormalised Borcherds square $\Phi_{10}^{\mathrm{un}}(Z)=\Delta_5(Z)^2$ with product expansion $\prod_{(n,l,m) > 0}(1 - p^n r^l q^m)^{c(4nm - l^2)}$; the $\eta^{24}$ factor is the weight-$12$ prefactor from the elliptic direction, matching Dijkgraaf--Moore--Verlinde--Verlinde 1997.
\end{proof}

\begin{problem}[Compact BV anomaly cancellation beyond the real-root subsystem]
\label{op:hcs-pushforward-k3e-delta5-compact-bv}
Extend Proposition~\ref{prop:hcs-pushforward-k3e-delta5} from
\(\fg_{\Delta_5}^{\mathrm{re}}\) to the full compact hCS field theory
with imaginary-root sectors and prove the corresponding anomaly
cancellation as a BV chain-level statement. The Borcherds imaginary-root
multiplicities are supplied by
Theorem~\ref{thm:hcs-pushforward-imaginary-roots-singular-theta}; the
missing datum here is the compact BV cancellation and OPE-compatible
transport, not the automorphic denominator formula.
\end{problem}

\begin{conjecture}[Three-path chain-level triangulation at \texorpdfstring{$N = 1$}{N=1}]
\label{cor:k3e-delta5-three-path-triangulation}
\ClaimStatusConjectured
Assume the BV pushforward, the Maulik--Okounkov stable-envelope residue, and the modular-operadic trace have been constructed on the same K3-fibre branch and compared through the projected period correspondence \(\pi_{\mathrm{Sp}_4}\circ\mathscr P_{K3}\). Then the three Borcherds-scalar readings are expected to satisfy
\[
 \kBKM^{\mathrm{BV}}(\fg_{\Delta_5})
 \;=\;
 \kBKM^{\mathrm{MO}}(\fg_{\Delta_5})
 \;=\;
 \mathrm{STr}^{\mathrm{Borch}}_{\mathrm{mod.op}}\!\bigl[B^{E_3}\bigl(\PhiFA_3(\Perf(K3 \times E))\bigr)\bigr]
 \;=\; \frac{c_1(0)}{2} \;=\; 5.
\]
The compact Hodge supertrace of \(K3\times E\) is a different scalar:
\[
 \kappa_{\mathrm{ch}}^{\mathrm{Hodge}}(K3\times E)
 =
 \sum_q(-1)^q h^{0,q}(K3\times E)
 =
 0.
\]
It is not one of the equal quantities above.
\end{conjecture}

\begin{proof}[Comparison calculation]
Hodge path. The chain-level Hodge supertrace on \(\mathrm{HH}^\bullet(D^b\Coh(K3\times E))=\mathrm{PV}^\bullet(K3\times E)\) gives the total-space value \(0\) by Künneth and Serre duality. It supplies a consistency check on compactness; it does not produce the weight-\(5\) Siegel form.
BV path. Proposition~\ref{prop:hcs-pushforward-k3e-delta5}(ii)--(iv) integrates the Costello--Li--Yagi propagator $P_X = P_{K3}^{(0,1)} \boxtimes P_E^{(0,1)}$ at one loop; the wheel-graph sum localises at the Humbert locus $H_1$. In the divisor normalisation $\mathrm{div}(\Delta_5)=H_1+2H_4$, one has $\operatorname{ord}_{H_1}(\Delta_5)=1$ and $\operatorname{ord}_{H_4}(\Delta_5)=2$; the number $5$ entering this path is the Borcherds weight $\kBKM(\Delta_5)=c_1(0)/2$, not an order of vanishing (Borcherds 1998 Theorem 13.3, Gritsenko 1999 Theorem 1.2).
MO path. The Schiffmann--Vasserot K3 cohomological Hall algebra carries a Maulik--Okounkov stable-envelope $R$-matrix; its residue at the Mukai-lattice wall is given by the Bruinier theta-integrand restriction to $H_1$ (Bruinier 2002 Theorem 4.23) with Weyl-denominator normalisation factor $1/2$.
Pairwise compatibility. The required chain-homotopies must be built in the same completed ambient category. Until this is done, the equality is a conjectural triangulation of the Borcherds scalar, not an unconditional Hodge computation.
Master theorem. The universal Borcherds identity $\kBKM(\Phi_N)=c_N(0)/2$ fixes the terminal scalar. It does not by itself construct the modular-operadic trace on \(B^{E_3}(\PhiFA_3(\cC_N))\).
Primary-literature anchors: Gritsenko--Nikulin 1995 Part~II Theorem~2.1; Borcherds 1998 Theorem 13.3; Gritsenko 1999 Theorem 6.1; Bruinier 2002 Proposition 5.1; Kapranov 1999 Theorem 2.6; Costello--Li 2016 Proposition 5.2; Lorgat 2020.
\end{proof}

\begin{remark}[Real-root and compact-CY\texorpdfstring{$_3$}{3} loci]
\label{rem:hcs-pushforward-scope}
Proposition~\ref{prop:hcs-pushforward-k3e-delta5} is rigorous at the real-root subalgebra level and on the $K3 \times E$ formal locus. The extension to all (real + imaginary) roots of $\fg_{\Delta_5}$ is supplied by the Borcherds imaginary-root multiplicity theorem, whose chiral-avatar transport is the content of Theorem~\ref{thm:hcs-pushforward-imaginary-roots-singular-theta} below. For compact CY$_3$ other than $K3 \times E$ (quintic, generic), the quartic obstruction coefficient $\chi_{\mathrm{top}}(X)/24 \neq 0$ and the analogue of (ii) fails; this is Conjecture~\ref{conj:quintic-hcs-kanom}. The anomaly-split discipline of (i) is universal. The $\fsl_{N \geq 3}$ extension is blocked by the non-zero cubic tensor; the non-simply-laced folding problem is open for the different reason that the twisted comparison is not constructed by the rank-one real-root argument.
\end{remark}

\begin{theorem}[Imaginary-root extension of the \texorpdfstring{$K3 \times E$}{K3 x E} hCS pushforward via the chiral singular theta correspondence]
\label{thm:hcs-pushforward-imaginary-roots-singular-theta}
Let $X = K3 \times E$, $\pi \colon X \to E$, \(L=\Lambda^{3,2}\) the
paramodular lattice of signature \((3,2)\) with Grassmannian
\(\mathcal G(L)\), and
\[
F=\phi_{0,1}^{\mathrm{EZ}}\in
J_{0,1}^{\mathrm{wk}}(\mathrm{SL}_2(\mathbb Z))
\]
the Eichler--Zagier generator of weight \(0\) and index \(1\), with
\(c(-1)=1\), \(c(0)=10\), \(c(3)=-64\), \(c(4)=108\). The physical K3
elliptic genus is \(Z_{K3}=2\phi_{0,1}^{\mathrm{EZ}}\); the
\(\Delta_5\) Borcherds denominator uses the Eichler--Zagier seed \(F\).
Extend
Proposition~\ref{prop:hcs-pushforward-k3e-delta5} past the real-root
subsystem as follows.
\begin{enumerate}[label=\textup{(\roman*)}]
 \item \emph{Chiral avatar of the singular theta integral.} The $K3$-fibre integration of the hCS factorisation algebra is representable as the regularised Siegel theta integral of $F$ against the lattice $L$:
 \[
 \pi_*^{K3} \Obs_{\hCS}^{\fg_{\Delta_5}}(X)(Z) \;\simeq\; \int_{\mathcal{F}}^{\mathrm{reg}} F(\tau) \, \Theta_L(\tau, \bar\tau, Z) \, \frac{d\tau \wedge d\bar\tau}{y^2}, \qquad Z \in \Pi^+_L + i\,\mathcal{G}(L),
 \]
 where $\Theta_L$ is the Siegel theta kernel, $\mathcal{F} = \mathbb{H} / \mathrm{SL}_2(\mathbb{Z})$ the modular fundamental domain, and $\mathrm{reg}$ is the Harvey--Moore regularisation (Borcherds~$1998$ \S6). The identification is the Costello--Gwilliam holomorphic-topological tensor theorem (Costello--Gwilliam~$2017$ Vol~II Theorem~$4.10.1$) applied to the holomorphic $K3$-fibre coupled to the topological $\R^2$-direction of the $E_3$-factorisation datum on $X$.
 \item \emph{Denominator-formula consistency.} Expanding \textup{(i)} near the level-$0$ cusp $F_\rho$ with Weyl vector $\rho_W \in L$ yields the Borcherds product
 \[
 \pi_* \Obs_{\hCS}^{\fg_{\Delta_5}}(X)(Z) \;=\; e^{2\pi i (\rho_W, Z)} \prod_{\substack{\lambda \in L^+ \\ (\lambda, \rho_W) > 0}} \bigl(1 - e^{2\pi i (\lambda, Z)}\bigr)^{c\bigl(-\lambda^2/2\bigr)} \;=\; \Delta_5(Z),
 \]
 with $c\bigl(-\lambda^2/2\bigr) = c_{\phi_{0,1}}(n, \ell)$ for $\lambda = (n, \ell, m)$ and $\lambda^2 = 2(nm - \ell^2/4)$. The root multiplicity of $\alpha \in \Delta_+(\fg_{\Delta_5})$ equals $\mathrm{mult}(\alpha) = c\bigl(-\alpha^2/2\bigr)$ for all $\alpha$, real and imaginary, matching Gritsenko--Nikulin~$1998$ Theorem~$2.1$.
 \item \emph{Imaginary-root multiplicities at low norm.} At the null roots $\alpha \in \{(1,0,0), (0,0,1)\}$ with $\alpha^2 = 0$: $\mathrm{mult}(\alpha) = c(0) = 10$. At the first strictly imaginary root stratum $\alpha^2 = -3$ (so $-\alpha^2/2 = 3/2$, realised at $(n, \ell, m) = (1, 1, 1)$ with $D = 4nm - \ell^2 = 3$): $\mathrm{mult}(\alpha) = c(3) = -64$, interpreted as the signed difference $\dim \Delta^{\mathrm{im}}_0(\alpha) - \dim \Delta^{\mathrm{im}}_1(\alpha)$ of even- and odd-imaginary-root multiplicities per Construction~\ref{constr:k3e-roots}. At $\alpha^2 = -4$ (so $D = 4$): $\mathrm{mult}(\alpha) = c(4) = 108$.
 \item \emph{Weight formula.} The weight of the chiral-avatar automorphic form equals $\mathrm{wt}(\Delta_5) = c(0)/2 = 5 = \kBKM(\Delta_5)$, specialising the Borcherds weight theorem \textup{(}Theorem~\ref{thm:borcherds-lift-universal}\textup{)} and the chain-level master identity \textup{(}Theorem~\ref{thm:bl-master-kappa-bkm-chain-level}\textup{)} to $(V, L) = (V^{K3\mathrm{-VOA}}, \Lambda^{3,2})$. The composite identification $\pi_*^{K3} \Obs_{\hCS}^{\fg_{\Delta_5}}(X) = \Delta_5^{-1} \cdot V^{\mathrm{ch}}(\fg_{\Delta_5})$ extends the real-root identification of Proposition~\ref{prop:hcs-pushforward-k3e-delta5}(iii) to the full root system.
\end{enumerate}
\end{theorem}

\begin{proof}[Sketch]
\emph{(i) Chiral avatar.} The $K3 \times E$ hCS factorisation algebra is the tensor product of a holomorphic factorisation algebra on $K3$ and a holomorphic factorisation algebra on $E$, twisted by the $\fg_{\Delta_5}$ gauge datum. Factorisation homology over the compact complex surface $K3$ is the Getzler--Kapranov modular-operadic supertrace on the $E_3^{\mathrm{hol}}$-observable algebra (Costello--Gwilliam~$2021$ Vol~II Theorem~$4.10.1$; Costello--Gaiotto~$2018$ \S4). For $K3$ rational elliptic, the supertrace is expressible as the regularised integral over the $\mathrm{SL}_2(\mathbb{Z})$-modular base of the Eichler--Zagier seed \(\phi_{0,1}^{\mathrm{EZ}}\) paired with the Siegel--Narain kernel determined by a positive four-plane in the Mukai lattice; the Borcherds product appears after projection to the \(\Lambda^{3,2}\) theta-lift datum. The remaining $\R^2$-factor of the $6\mathrm{d} \to 4\mathrm{d}$ holomorphic-topological reduction supplies the weak modular covariance required for the Harvey--Moore regularisation to converge.
\emph{(ii) Borcherds product.} Borcherds~$1998$ Theorem~$13.3$ expands the regularised theta integral at a $0$-dimensional cusp as the product displayed; the exponents are the Fourier coefficients of $F$ at negative argument $-\lambda^2/2$. For \(L=\Lambda^{3,2}\) and \(F=\phi_{0,1}^{\mathrm{EZ}}\) the product is \(\Delta_5\) up to the \(pqt\) Weyl-vector prefactor (Gritsenko--Nikulin~$1998$ \S$2$, Theorem~$2.1$).
\emph{(iii) Low-discriminant coefficients.} Fourier coefficients of \(\phi_{0,1}^{\mathrm{EZ}}\) depend only on the discriminant \(D=4n-\ell^2\): \(c(-1)=1\) (the polar term driving the \(y^{\pm1}\) head), \(c(0)=10\) (lightlike roots), \(c(3)=-64\) (first strict imaginary), \(c(4)=108\) (second strict imaginary). The generating series opens as \(y^{-1} + 10 + y + q(10 y^{-2} - 64 y^{-1} + 108 - 64 y + 10 y^2) + O(q^2)\) (Eichler--Zagier~$1985$ Table~$1$), reproducing \(c(3)=-64\), \(c(4)=108\). The signed reading \(-64=0-64\) at \(\alpha^2=-3\) is Gritsenko--Nikulin~$1998$ Example~$2$: the root space is a purely odd \(64\)-dimensional superspace.
\emph{(iv) Weight.} $\mathrm{wt}(\Phi_F) = c(0)/2 = 5$ is Borcherds~$1998$ Theorem~$13.3$(iv); the specialisation $V = V^{K3\mathrm{-VOA}}$, $L = \Lambda^{3,2}$ matches Theorem~\ref{thm:borcherds-lift-universal}(iv) and Theorem~\ref{thm:bl-master-kappa-bkm-chain-level} at $N=1$. The composite factorisation of the pushforward is Costello--Gwilliam holomorphic-topological tensor applied after the Weyl-group averaging of (ii).
\end{proof}

\begin{remark}[What the chiral avatar adds to Borcherds $1998$]
\label{rem:chiral-avatar-adds-to-borcherds}
Borcherds~$1998$ Theorem~$13.3$ and Gritsenko--Nikulin~$1998$ establish the automorphic-form identity at the level of formal power series on the Siegel upper half-space. Theorem~\ref{thm:hcs-pushforward-imaginary-roots-singular-theta} states that this identity is the Fourier expansion, at a $0$-dimensional cusp, of the K3-fibre factorisation-homology pushforward of the chiral $\fg_{\Delta_5}$-gauge observables on $K3 \times E$. The addition is not a new proof of the Borcherds product: it is the identification of the Borcherds singular theta integrand with the Costello--Gaiotto $5\mathrm{d}$-hCS factorisation-algebra pairing (a working note \S $E_3$ holomorphic algebra of observables), normalised by the modular-operadic supertrace of Theorem~\ref{thm:bl-master-kappa-bkm-chain-level}. The imaginary-root multiplicities $\mathrm{mult}(\alpha) = c(-\alpha^2/2)$ are the Fourier coefficients of that pushforward automorphic form; the compact BV anomaly problem outside the $K3\times E$ singular-theta locus is not used here.
\end{remark}

\begin{remark}[Why the $b^+ \geq 2$ hypothesis is satisfied for $K3 \times E$]
\label{rem:bplus-geq-two-k3e}
The Borcherds singular theta lift \textup{(}Theorem~\ref{thm:borcherds-lift-universal}\textup{)} requires $b^+ \geq 2$. For $L = \Lambda^{3,2}$ the signature is $(3, 2)$, so $b^+ = 3 \geq 2$. The $N = 1$ Monster case $L = \mathrm{II}_{1,1}$ sits outside this hypothesis; the $K3 \times E$ case sits inside it. The chiral avatar of Theorem~\ref{thm:hcs-pushforward-imaginary-roots-singular-theta} therefore transports the full Borcherds theta correspondence, including the imaginary-root multiplicity identity, without the metaplectic-cover complications of the Monster case \textup{(}Remark~\ref{rem:borcherds-three-cases} item~\textup{(1)}\textup{)}.
\end{remark}

\begin{conjecture}[Green--Schwarz closure for 6D \texorpdfstring{$\hCS$}{hCS} on compact CY\texorpdfstring{$_3$}{-3}]
\label{conj:quintic-hcs-kanom}
\ClaimStatusConjectured
Let $X$ be a compact Calabi--Yau threefold and $\frakg$ simple. The
one-loop gauge anomaly of 6D holomorphic Chern--Simons is controlled
locally by the degree-$4$ invariant polynomial
$I_4\in\mathrm{Sym}^4(\frakg^\vee)^{\frakg}$, not by the cubic Casimir.
A Green--Schwarz closure, when available, must trivialise this quartic
class together with the compact CY$_3$ Atiyah-class data. On loci with
zero BCOV Euler-characteristic coefficient, such as $K3\times E$ and
$T^6$, the Euler-characteristic contribution vanishes, but this is not a
universal proof that every gauge algebra is unobstructed. On the
quintic, $\chi_{\mathrm{top}}(X_5)=-200$ makes the quartic obstruction
the first compact instance. The Gritsenko $\Delta_5$ sector lives on
the $K3\times E$ formal locus where the Stage-$1$ obstruction classes
used in this correspondence vanish or have zero coefficient.
\end{conjecture}

\begin{remark}[Dimensional origin of the \texorpdfstring{$\En$}{E-n} level]
\label{rem:en-from-dimension}
The $\En$ level is set by the complex dimension of the ambient space, not by the Deligne conjecture on Hochschild cochains. Specifically:
\begin{itemize}
 \item The $\Etwo$ structure on $\C^2$ comes from the $E_2$-operad acting on $\Conf_n(\C^2)$. Topologically, $\Conf_2(\C^2) \simeq \Conf_2(\R^4) \simeq S^3$ has $\pi_1 = 0$, so the topological braiding is \emph{trivial}. The nontrivial braiding of the Yangian $R$-matrix arises not from $\pi_1$ but from the \emph{Omega-background deformation} $(h_1, h_2)$ of the holomorphic factorization algebra: the equivariant parameters break the topological $E_2$-symmetry to a nontrivially braided $E_2$-chiral structure. Without the Omega-background, the braiding is symmetric.
 \item The $E_3$ structure on $\C^3$ comes from the topological $E_3$-operad acting on $\Conf_n(\C^3)$, which has $\pi_1(\Conf_2(\R^6)) \cong \pi_1(S^5) = 0$ (trivial braiding topologically). As in the $\C^2$ case, the nontrivial braiding at the $E_3$ level arises from the Omega-background deformation $(h_1, h_2, h_3)$, not from the topology of the configuration space.
 \item The algebraic $E_3$ from the Deligne conjecture (Remark~\ref{rem:deligne-hochschild}) requires $E_\infty$ input and produces $E_3$ on Hochschild cochains; the topological $E_3$ from $\C^3$ configuration spaces is a priori distinct. Under formality (Kontsevich 1999), the two $E_3$ structures agree rationally.
\end{itemize}
The CY dimension $d$ of the underlying category and the complex dimension $n$ of the ambient holomorphic CS space are related but distinct in general: for $\C^3$, the CY dimension of the quiver category is $d = 3$ and the CS dimension is $n = 3$ (coincidence). For $K3 \times E$, the CY dimension is also $d = 3$ and $\dim_\C(K3 \times E) = 3$, so the coincidence persists. The distinction matters for \emph{non-compact} examples: for the resolved conifold $\cO(-1)^{\oplus 2} \to \bP^1$, the CY dimension is $d = 3$ (it is a non-compact CY threefold) but the holomorphic CS theory is naturally formulated on the total space $\C^4$ with a superpotential constraint, not on a $3$-dimensional ambient space.
\end{remark}

\subsection{The chiral Chevalley--Eilenberg complex}
\label{subsec:chiral-ce}

The bar complex of a chiral algebra is the chiral analogue of the classical Chevalley--Eilenberg complex of a Lie algebra. Making this identification explicit reveals the connection between Costello's holomorphic CS observables and the bar-cobar formalism of Vols~I--II.

\begin{construction}[Chiral CE chains and cochains]
\label{constr:chiral-ce}
Let $A$ be a chiral algebra with augmentation $\varepsilon \colon A \to \Omega_C$ and augmentation ideal $\bar{A} = \ker(\varepsilon)$. The three chiral CE-type complexes:
\begin{enumerate}[label=\textup{(\roman*)}]
 \item The \emph{chiral CE chains} (labeled-ordered) are the ordered bar complex $B^{\mathrm{ord}}(A) = T^c(s^{-1}\bar{A})$, the cofree conilpotent coalgebra on the desuspension of $\bar{A}$, equipped with the deconcatenation coproduct and bar differential built from the chiral product. This is the direct analogue of the Chevalley--Eilenberg chain complex $C_\bullet(\frakg) = \bigwedge^\bullet \frakg$ with its Lie-bracket-induced differential. The labeled-ordered bar retains the full $R$-matrix data.

 \item The \emph{chiral CE chains} (symmetric), under the same locality
 and $R$-twisted descent package, are the symmetric bar complex
 $B^{\Sigma}(A) = \Sym^c(s^{-1}\bar{A})$ with the coshuffle coproduct
 and the symmetrized differential. This is the direct analogue of
 $C_\bullet(\frakg) = \bigwedge^\bullet \frakg$ in its standard
 commutative-coalgebra form. Vol~I's symmetric bar lives here; its
 ordered bar is the labelled complex in \textup{(i)}. Without descent
 this symmetric object is only a graded coinvariant shadow.

 \item The \emph{chiral CE cochains} are the chiral Hochschild cochain
 complex $C^\bullet_{\mathrm{ch}}(A, A)=\RHom(\Omega B(A), A)$, the
 chiral derived center $Z^{\mathrm{der}}_{\mathrm{ch}}(A)$ of Vol~I
 Theorem~H. This is the analogue of
 $C^\bullet(\frakg, \frakg)=\Hom(\bigwedge^\bullet \frakg,\frakg)$,
 the Chevalley--Eilenberg cochains with adjoint coefficients.
\end{enumerate}
The Koszul duality of Vol~I sends the CE chains~(i) to the Koszul
dual $A^!=D_{\Ran}(B(A))$, the Verdier dual of the bar complex. The
Koszul dual $A^!$ is a fourth object, distinct from the three CE
complexes listed above and from the CE cochains~(iii) in particular:
$A^!$ controls the defect, while the CE cochains
$Z^{\mathrm{der}}_{\mathrm{ch}}(A)$ are the algebraic closed-sector
object and become physical bulk observables only after the relevant
open--closed comparison has been supplied
(see Proposition~\ref{prop:three-dualities}). In classical terms,
$A^!$ is the enveloping algebra of the Koszul-dual Lie algebra
$\frakg^\vee$, not the CE cochains
$C^\bullet(\frakg,\frakg)$ with adjoint coefficients.
\end{construction}

\begin{proposition}[Holomorphic CS observables as chiral CE cochains]
\label{prop:hcs-as-ce-cochains}
For $A = V_k(\frakg)$ the Kac--Moody vertex algebra at level $k$ (the
boundary algebra of 3d holomorphic CS), the chiral CE cochains
$C^\bullet_{\mathrm{ch}}(A, A)$ compute the algebraic closed-sector
comparison target for the CS theory; identifying it with physical bulk
observables is the holomorphic-CS open--closed comparison input. At
the critical level $k = -h^\vee$, the zeroth cohomology of the CE
cochains is the Feigin--Frenkel center
$\mathrm{Fun}(\mathrm{Op}_{G^L}(D))$
(Theorem~\ref{thm:feigin-frenkel-center}).
\end{proposition}

\begin{proof}
The identification of the derived center with bulk observables is Vol~I Theorem~H. The Feigin--Frenkel identification is Theorem~\ref{thm:feigin-frenkel-center} (Chapter~\ref{ch:geometric-langlands}).
\end{proof}

\begin{proposition}[Bar complex as CE chains of the Lie conformal algebra]
\label{prop:bar-ce-identification}
Let $L$ be a Lie conformal algebra with $\lambda$-bracket $\{a_\lambda b\}$, and let $A = U^{\mathrm{ch}}(L)$ be its chiral envelope. Then the ordered bar complex of $A$ is the Chevalley--Eilenberg chain complex of $L$:
\[
 B^{\mathrm{ord}}(U^{\mathrm{ch}}(L)) \;\cong\; \mathrm{CE}_\bullet(L).
\]
The graded vector spaces on each side are $T^c(s^{-1}\bar{A}) \cong \bigwedge^\bullet(L)$ with Poincar\'e polynomial $(1+t)^{\dim L}$; the bar differential $d_B$ encodes the $0$\nobreakdash-th product $a_{(0)}b$ (the Lie bracket) as the CE differential:
\[
 d_{\mathrm{CE}}(g_{i_1} \wedge \cdots \wedge g_{i_k}) = \sum_{a < b} (-1)^{a+b} [g_{i_a}, g_{i_b}] \wedge g_{i_1} \wedge \cdots \wedge \widehat{g_{i_a}} \wedge \cdots \wedge \widehat{g_{i_b}} \wedge \cdots \wedge g_{i_k}.
\]
Dually, the chiral CE cochains $\mathrm{CE}^\bullet(L, L) = \Hom(\bigwedge^\bullet L, L)$ are the Hochschild cochains $\HH^\bullet(A, A)$, i.e.\ the derived center $Z^{\mathrm{der}}_{\mathrm{ch}}(A)$. For abelian $L$ (e.g.\ the Heisenberg current algebra), $d_{\mathrm{CE}} = 0$ and $\HH^k(A, A) = \binom{\dim L}{k}$.
\end{proposition}

\begin{proof}
The identification follows from the Poincar\'e--Birkhoff--Witt
filtration on $U^{\mathrm{ch}}(L)$ (Kac 1998, \S\,2.6;
Loday--Vallette 2012, Theorem~10.1.6).  For an ordinary Lie algebra,
the reduced bar resolution of $U(\frak g)$ is compared with the
Chevalley--Eilenberg resolution
$U(\frak g)\otimes\Lambda^\bullet\frak g$.  Applying
$\operatorname{Hom}_{U(\frak g)^e}(-,U(\frak g))$ gives the standard
identity
\[
 HH^\bullet(U(\frak g),U(\frak g))
 \cong H^\bullet(\frak g,U(\frak g)^{\mathrm{ad}}).
\]
The chiral PBW comparison replaces $\frak g$ by the Lie conformal
algebra~$L$ and uses the factorization envelope
$U^{\mathrm{ch}}(L)$.  Its degree-two bar differential is the OPE
residue $a_{(0)}b$, and the degree-three identity is the Jacobi
relation.  A family datum $H_H(U^{\mathrm{ch}}(L);S_L)$ then computes
the support of the chiral derived centre.  A comparison with
$\mathrm{CE}^\bullet(L,U^{\mathrm{ch}}(L)^{\mathrm{ad}})$ is a
separate PBW quasi-isomorphism.

Three elementary models fix the normalisation:
\begin{itemize}
 \item Heisenberg ($L = \langle J \rangle$, abelian): $\mathrm{CE}_\bullet = \bigwedge^\bullet(\langle J \rangle)$, $d_{\mathrm{CE}} = 0$, Poincar\'e~$= (1+t)$, $\HH^\bullet = \{1, 1\}$.
 \item Kac--Moody $\widehat{\fsl}_2$ ($L = \langle e, f, h \rangle$): $\mathrm{CE}_\bullet = \bigwedge^\bullet(\langle e, f, h \rangle)$, $d_{\mathrm{CE}}^2 = 0$ by Jacobi, Poincar\'e~$= (1+t)^3$.
 \item Yangian $Y(\widehat{\fgl}_1)$ (truncated, $3$ generators): strict Lie (class~L), Poincar\'e~$= (1+t)^3$, genus-$3$ bar dimension~$= 2^9 = 512$.
\end{itemize}
\end{proof}

\begin{remark}[\texorpdfstring{$L_\infty$}{L-inf} deformation and shadow tower]
\label{rem:linfinity-shadow-identification}
For a vertex algebra $A$ whose $A_\infty$ structure has higher operations $m_k \neq 0$ ($k \geq 3$), the bar complex $B(A)$ is the CE complex of an $L_\infty$ algebra $L_A$. The $L_\infty$ structure maps are:
\begin{align*}
 l_2 &= \text{Lie bracket (from } m_2\text{, the OPE residue)}, \\
 l_3 &= \text{Jacobiator (from } m_3\text{, the associator)}, \\
 l_k &= \text{higher Jacobiator (from } m_k\text{)}.
\end{align*}
The CE differential of the $L_\infty$ algebra $L_A$ is $d_{\mathrm{CE}}^{L_\infty} = \sum_{k \geq 2} d^{(k)}$, with $d^{(k)}$ encoding $l_k$. The shadow tower of Vol~I is precisely this $L_\infty$ CE tower:
\[
 S_r = l_r \text{ contribution to the CE differential, i.e.\ } \delta^{(r)} \text{ of } \textup{rem:shadow-ainfty-coproduct-vol3}.
\]
For class~G (e.g.\ Heisenberg): $l_k = 0$ for all $k \geq 2$ (abelian), tower terminates at $S_2 = \kappa_{\mathrm{ch}}$. For class~L (e.g.\ Yangian): $l_k = 0$ for $k \geq 3$ (strict Lie), and the CE complex is the standard $\mathrm{CE}_\bullet(L)$ with Poincar\'e $(1+t)^n$; at genus~$3$ this reproduces the $(1+t)^{3g}$ formula of. For class~M (e.g.\ Virasoro): $l_k \neq 0$ for all $k$, and the shadow tower is infinite; computational verification at $c = 1$:
\[
 l_3(T, T, T) = -2, \quad l_4(T, T, T, T) = \tfrac{40}{27}, \quad S_2 = \kappa_{\mathrm{ch}} = \tfrac{1}{2}, \quad S_3 = 2, \quad S_4 = \tfrac{10}{27}.
\]
These values match the $A_\infty$ bar complex and the shadow tower
recursion.
\end{remark}

\subsection{The universal defect and holomorphic Wilson lines}
\label{subsec:universal-defect}

In the holomorphic CS framework, a \emph{defect} is a codimension-$2$ locus along which the gauge field acquires a prescribed singularity. The simplest defect is the holomorphic Wilson line: a line operator valued in a representation $V$ of $\frakg$.

\begin{construction}[Universal defect from Koszul duality]
\label{constr:universal-defect}
Let $A$ be the boundary chiral algebra of holomorphic CS with gauge algebra $\frakg$ on $M$. The \emph{universal defect algebra} is the Koszul dual $A^! = D_{\Ran}(B(A))$. In the 3d case ($M = \Sigma \times \R$), $A^!$ is the chiral CE/bar-cochain object whose reflected current presentation has Feigin--Frenkel level $k' = -k - 2h^\vee$. The bulk-boundary coupling is the canonical map
\[
 \mu_{\mathrm{bb}} \colon Z^{\mathrm{der}}_{\mathrm{ch}}(A) \otimes A^! \;\longrightarrow\; A^!
\]
making $A^!$ a module over the bulk algebra $Z^{\mathrm{der}}_{\mathrm{ch}}(A)$. The holomorphic Wilson line in representation $V$ corresponds to a specific $A^!$-module: the image of $V$ under the equivalence $\Rep^{\Eone}(A^!) \simeq \Rep^{\Eone}(A)^{\mathrm{rev}}$ (Vol~I, Theorem~A, Verdier leg: $D_{\Ran}$ exchanges $A$ and $A^!$ module categories with reversed monoidal structure).

In the 5d case ($M = \C^2 \times \R$), the universal defect algebra is the Koszul dual of the affine Yangian. For $\frakg = \fgl_1$, the defect algebra $A^!$ encodes the Wilson lines of the 5d theory; the RTT formalism applied to the defect $R$-matrix generates the dual Yangian presentation. (The quantum toroidal algebra arises from the \emph{6d} theory on $\C^3$, not from 5d; see Conjecture~\ref{conj:6d-boundary-toroidal}.) The Drinfeld center of the defect module category recovers the $\Etwo$ braided structure:
\[
 \cZ(\Rep^{\Eone}(A^!)) \;\simeq\; \Rep^{\Etwo}(Z^{\mathrm{der}}_{\mathrm{ch}}(A)).
\]
\end{construction}

\begin{conjecture}[6d universal defect and chiral quantum groups on \texorpdfstring{$\C^3$}{C-3}]
\label{conj:6d-defect-c3}
\ClaimStatusConjectured
For the 6d holomorphic theory on $\C^3$ with gauge algebra $\fgl_1$, the universal defect algebra $A^!_{\C^3}$ (the Koszul dual of the quantum toroidal boundary algebra $U_{q,t}(\widehat{\widehat{\fgl}}_1)$) satisfies:
\begin{enumerate}[label=\textup{(\roman*)}]
 \item $A^!_{\C^3}$ carries $E_3$-chiral factorization structure. The underlying topological $E_3$-operad action on $\Conf_n(\C^3)$ has trivial braiding ($\pi_1(\Conf_2(\R^6)) = 0$); the nontrivial $E_3$-chiral braiding arises from the Omega-background deformation $(h_1, h_2, h_3)$ of the holomorphic factorization structure (Remark~\ref{rem:e3-not-symmetric}).
 \item The $\Eone$-projection of $A^!_{\C^3}$ to a curve $C \subset \C^3$ is the Koszul dual of the quantum toroidal algebra, with inverted parameters $q^{-1}, t^{-1}$ (parameter inversion under Koszul duality; see Table of Chapter~\ref{ch:quantum-groups}, Definition~\ref{def:qgf-four-regimes}).
 \item The $\Etwo$-projection to $\C^2 \subset \C^3$ is the $\Etwo$-chiral algebra whose Drinfeld center $\cZ(\Rep^{\Eone}(A^!_{\C^3}|_C))$ is the braided monoidal category of Fock representations of the quantum toroidal algebra.
 \item The \emph{chiral quantum group} on $\C^3$ is the Hopf-algebra-in-braided-category structure on the representation category $\Rep^{\Etwo}(A^!_{\C^3}|_{\C^2})$, equipped with the $R$-matrix from the $E_3$ braiding.
\end{enumerate}
The conjecture composites three constructions (6d $E_3$ factorization, Koszul duality, Drinfeld center), each of which is independently established at lower dimension; the composite arrow at 6d is conjectural.
\end{conjecture}

\begin{conjecture}[6d holomorphic theory on \texorpdfstring{$K3 \times E$}{K3 x E} and the BKM chiral algebra]
\label{conj:6d-k3xe}
\ClaimStatusConjectured
Let $X = K3 \times E$ with $E$ an elliptic curve. On the non-Lagrangian 6d factorization framework, together with the framed K3-fibre specialisation of Theorem~\ref{thm:cy-to-chiral-d3}, the holomorphic theory on $X$ should produce:
\begin{enumerate}[label=\textup{(\roman*)}]
 \item An $\Eone$-chiral algebra $A_{X}^{(K3,E)} = \Phi_3^{(K3,E)}(D^b(\Coh(X)))$ on $E$ on the framed K3-fibre specialisation of Theorem~\ref{thm:cy-to-chiral-d3}, whose character is conjecturally controlled by the Igusa cusp form $\Phi_{10}^{\mathrm{un}}$. The expected chiral algebra is a deformation of the lattice vertex algebra $V_{\Lambda_{K3}}$ of the K3 lattice, with deformation parameters determined by the complex structure of $K3$.
 \item An $\Etwo$-chiral factorization algebra on $K3$ itself (the ``transverse'' directions), whose Drinfeld center gives the braided monoidal category of BPS modules.
 \item The universal defect algebra $A^!_X$ on $E$, whose representation category contains the holomorphic Wilson lines wrapping $E$. These Wilson lines are indexed by the BPS spectrum of $K3 \times E$ (Mukai vectors in $H^*(K3, \Z)$).
 \item The \(\kappa_\bullet\)-spectrum has four coordinates: the Stage-2 Heisenberg reading has $\kappa_{\mathrm{ch}}^{\mathrm{Heis}}(A_X) = 3$, the Borcherds automorphic weight is $\kappa_{\mathrm{BKM}}(\Delta_5) = 5$, the total-space categorical invariant is $\kappa_{\mathrm{cat}}(X) = 0$, and the K3-fibre lattice invariant is $\kappa_{\mathrm{fiber}}(K3) = 24$, consistent with Chapter~\ref{ch:k3-times-e}.
\end{enumerate}
The framed $\Eone$ object is supplied by Theorem~\ref{thm:cy-to-chiral-d3} after the K3-fibre datum is fixed. The conjectural content is the non-Lagrangian 6d algebraic framework for $K3\times E$ (Remark~\ref{rem:6d-not-lagrangian}) and the identification of the BPS root datum with the Borcherds--Kac--Moody root system of $\mathfrak{g}_{\Delta_5}$ (Chapter~\ref{ch:k3-times-e}).
\end{conjecture}

\begin{remark}[What the 6d theory adds to the \texorpdfstring{$K3 \times E$}{K3 x E} chiral-algebra construction]
\label{rem:6d-adds-to-k3xe}
Chapters~\ref{ch:k3-times-e} and~\ref{chap:toroidal-elliptic} develop the $K3 \times E$ chiral algebra construction: the first via the framed K3-fibre specialisation of the CY-to-chiral correspondence $\Phi_3^{(K3,E)}$ (Theorem~\ref{thm:cy-to-chiral-d3}; the further quantum-group identification CY-C remains conjectural), the second via the toroidal/elliptic quantum group presentation (conditional on Conjecture~\ref{conj:toroidal-e1}). Conjecture~\ref{conj:6d-k3xe} adds a third conjectural construction pathway: the holomorphic CS / factorization homology route. The three pathways should agree if all three constructions are realized:
\begin{center}
\small
\begin{tabular}{lll}
 \toprule
 \textbf{Pathway} & \textbf{Input} & \textbf{Output} \\
 \midrule
 CY-to-chiral $\Phi_3^{(K3,E)}$ & $D^b(\Coh(K3 \times E))$ & $A_{K3 \times E}^{(K3,E)}$ (framed object-level specialisation, Theorem~\ref{thm:cy-to-chiral-d3}) \\
 Toroidal presentation & $U_{q,t}(\widehat{\widehat{\fgl}}_1)$ & $\Eone$-chiral on $E$ \\
 6d holomorphic theory & $K3 \times E$ as ambient & $\Eone$-chiral on $E$ \\
 \bottomrule
\end{tabular}
\end{center}
The agreement of the three is itself conjectural. The 6d pathway provides the geometric origin of the quantum toroidal parameters $(q, t)$: they arise from the K3 complex structure moduli via equivariant localisation on the Omega-background of the normal bundle $N_{C/X}$ (the normal bundle grading does not directly identify with $(q,t)$; the intermediary mechanism is Nekrasov-type equivariant localisation).
\end{remark}

\subsection{Codimension-\texorpdfstring{$2$}{2} defect OPE from the 6d action}
\label{subsec:codim2-defect-ope}

The stage-2 specialisation $\SpCh_{N_{C/Y}, C}(\PhiFA_3(\cC))$ for
$Y = \C^3$ and $C = \C_{z_1}$ should be the defect algebra on $C$
obtained by pushing the $E_3$-holomorphic factorization algebra
$\PhiFA_3(\cC)$ forward along the normal bundle
$N_{C/Y} = \C^2_{z_2, z_3}$, with the Omega-background equivariant
weights $(h_2, h_3)$ of $N_{C/Y}$ supplying the specialisation datum
(Definition~\ref{def:phi-fa-and-sp}, Remark~\ref{rem:multiple-e1-specialisations}).
The computation below is a low-spin OPE witness for that expected
identification: restricting the abelian local hCS model to a tubular
neighbourhood of \(C\) and extracting boundary-mode OPEs recovers the
spin-\(1\) current and the Sugawara spin-\(2\) Virasoro field. This
low-spin witness does not construct the full
renormalised 6d hCS factorisation algebra, all $W_{1+\infty}$ spins, or
the hCS-to-Hall comparison of Chapter~\ref{ch:cy3-chain-level-bridge}.

\begin{proposition}[Low-spin defect OPE witness for 6d hCS on \texorpdfstring{$\C^3$}{C-3}]
\label{prop:codim2-defect-ope}
Let $Y = \C^3$ with holomorphic $3$-form $\Omega = dz_1 \wedge dz_2 \wedge dz_3$, and let $C = \C_{z_1} \hookrightarrow Y$ be the curve $\{z_2 = z_3 = 0\}$. The normal bundle is $N_{C/Y} = \C^2_{z_2, z_3}$. In the abelian local hCS model with gauge algebra $\fgl_1$ and Omega-background parameters $(h_1, h_2, h_3)$ satisfying $h_1 + h_2 + h_3 = 0$, restriction to a tubular neighborhood $U = \C_{z_1} \times D^2_{z_2, z_3}$ of $C$ gives the following spin~$1$ and spin~$2$ OPEs expected of the $W_{1+\infty}$ defect algebra at level
\[
 \Psi \;=\; -\sigma_2 \;=\; -(h_1 h_2 + h_1 h_3 + h_2 h_3)
\]
with central charge $c = 1$ for the $\fgl_1$ Sugawara field. In particular, the spin-$1$ current $J$ and spin-$2$ current $T$ on $C$ satisfy:
\begin{enumerate}[label=\textup{(\alph*)}]
 \item \emph{Spin~$1$}: $J(z) J(w) = \Psi/(z - w)^2$, equivalently $\{J_\lambda J\} = \Psi \cdot \lambda$.
 \item \emph{Spin~$2$}: $T(z) T(w) = \tfrac{1}{2}/(z - w)^4 + 2T(w)/(z - w)^2 + \partial T(w)/(z - w)$, equivalently $\{T_\lambda T\} = \tfrac{1}{12}\lambda^3 + 2T\lambda + \partial T$.
\end{enumerate}
\end{proposition}

\begin{proof}
\emph{Arrow~1: restriction to the tubular neighborhood.}
The tubular neighborhood $U = \C_{z_1} \times D^2$ of $C$ in $Y = \C^3$ inherits the holomorphic $3$-form $\Omega|_U = dz_1 \wedge dz_2 \wedge dz_3$. The 6d holomorphic Chern--Simons action on $Y$ with gauge algebra $\fgl_1$:
\[
 S \;=\; \frac{1}{2}\int_Y \Omega \wedge \bigl(A \wedge \bar{\partial} A + \tfrac{2}{3}A \wedge A \wedge A\bigr)
\]
restricts to a free theory on $U$ because the cubic term vanishes for abelian $\fgl_1$.

\emph{Arrow~2: mode expansion in the normal directions.}
Expand the gauge field in modes of the normal bundle:
\[
 A(z_1, z_2, z_3, \bar{z}) \;=\; \sum_{m, n \geq 0} A^{(m,n)}(z_1, \bar{z}_1)\, z_2^m z_3^n.
\]
Each mode $A^{(m,n)}$ is a $(0,1)$-form on $C$ valued in $\fgl_1$. The Omega-background assigns equivariant weight $m \cdot h_2 + n \cdot h_3$ to $A^{(m,n)}$ under the torus $T^2 = \C^*_{z_2} \times \C^*_{z_3}$.

\emph{Arrow~3: integration over the normal fiber.}
Inserting the mode expansion into $S|_U$ and integrating the holomorphic $2$-form $dz_2 \wedge dz_3$ over the formal bidisk $D^2$ (extracting the $(0,0)$ Fourier coefficient of the product of normal modes), the action becomes:
\[
 S|_U \;=\; \sum_{m,n \geq 0} \int_C dz_1 \wedge A^{(m,n)} \bar{\partial} A^{(m,n)} \;+\; \text{(equivariant-weight terms from } \boldsymbol{h}\text{)}.
\]
The propagator of the free-field system on $C$ is
\[
 \langle A^{(m,n)}(z) \, A^{(p,q)}(w) \rangle
 \;=\; \frac{\delta_{m+p,0}\,\delta_{n+q,0}}{(z - w)^2},
\]
where the double pole is the standard holomorphic propagator on $\C$ and the delta-conditions enforce charge conservation in the normal directions.

\emph{Arrow~4: identification of the $W_{1+\infty}$ currents.}
The $\GL(2)$-invariant combinations of the normal modes produce the $W_{1+\infty}$ generators. At spin~$s$:
\[
 J_s(z_1) \;=\; \sum_{m+n = s-1} c_{m,n}^{(s)} \, A^{(m,n)}(z_1),
\]
where $c_{m,n}^{(s)}$ are the $\GL(2)$-invariant coefficients, namely
the minors of the identity matrix matching the invariant polyvector
fields \(\Omega_s\). In particular:
\begin{itemize}
 \item Spin~$1$: $J(z_1) = A^{(0,0)}(z_1)$, the mode at the origin of the normal fiber.
 \item Spin~$2$: $T(z_1) = \frac{1}{2\Psi}\mathopen{:}J(z_1)^2\mathclose{:}$ via the Sugawara construction at level $\Psi$.
\end{itemize}

\emph{Arrow~5: OPE computation.}
\emph{Spin~$1$}: From the propagator in Arrow~3, the contraction of two spin-$1$ modes gives
\[
 J(z) J(w) \;=\; \langle A^{(0,0)}(z)\, A^{(0,0)}(w) \rangle + \mathopen{:}J(z) J(w)\mathclose{:}
 \;=\; \frac{\Psi}{(z - w)^2} + \mathopen{:}J(z) J(w)\mathclose{:},
\]
where the coefficient $\Psi = -\sigma_2 = -(h_1 h_2 + h_1 h_3 + h_2 h_3)$ arises from the equivariant normalization of the propagator: the Omega-background twists the inner product on $\fgl_1$ by $-\sigma_2$ (which is the Kac--Moody level of the 5d theory on $\C^2 \times \R$, as in Costello~2013). In mode language: $J_{(1)} J = \Psi$, all other modes vanish. The classical $r$-matrix is $r^{\mathrm{Heis}}(z) = \Psi/z$ (C10; the level prefix $\Psi$ is retained to distinguish nonzero-level representations). At $\Psi = 0$: the $r$-matrix vanishes.

\emph{Spin~$2$}: The Sugawara construction $T = \frac{1}{2\Psi}\mathopen{:}J^2\mathclose{:}$ at level $\Psi$ for the abelian algebra $\fgl_1$ (with $h^\vee = 0$) produces a Virasoro with central charge
\[
 c \;=\; \frac{\Psi \cdot \dim(\fgl_1)}{\Psi + h^\vee} \;=\; \frac{\Psi \cdot 1}{\Psi + 0} \;=\; 1,
\]
independent of $\Psi$ (a feature of $\fgl_1$; for non-abelian $\frakg$, $c$ depends on the level). The Virasoro OPE is
\[
 T(z)\,T(w) \;=\; \frac{c/2}{(z - w)^4} + \frac{2T(w)}{(z - w)^2} + \frac{\partial T(w)}{z - w}
 \;=\; \frac{1/2}{(z - w)^4} + \frac{2T(w)}{(z - w)^2} + \frac{\partial T(w)}{z - w}.
\]
In mode language: $T_{(3)} T = c/2 = 1/2$, $T_{(1)} T = 2T$, $T_{(0)} T = \partial T$, $T_{(2)} T = 0$. In lambda-bracket form (with divided powers $\lambda^{(n)} = \lambda^n/n!$): $\{T_\lambda T\} = \tfrac{c}{12}\lambda^3 + 2T\lambda + \partial T = \tfrac{1}{12}\lambda^3 + 2T\lambda + \partial T$. The classical $r$-matrix is $r^{\mathrm{Vir}}(z) = \tfrac{1}{2}/z^3 + 2T/z$ (cubic $+$ simple pole; one less than the quartic OPE pole via $d\log$ absorption).

The $T$--$J$ OPE confirms that $J$ is a primary of conformal weight~$1$: $T(z)\,J(w) = J(w)/(z - w)^2 + \partial J(w)/(z - w)$, i.e.\ $T_{(1)}J = J$, $T_{(0)}J = \partial J$.

These are the spin~$1$ and spin~$2$ OPE relations of the expected
$W_{1+\infty}$ defect algebra at level $\Psi$ and central charge
$c = 1$. The all-spin vertex algebra and the renormalised hCS
factorisation-algebra construction remain additional obligations.
\end{proof}

\begin{lemma}[5d-to-6d-on-codim-\texorpdfstring{$2$}{2} effective level bridge]
\label{lem:5d-to-6d-codim2-bridge}
Let $Y = \C^3$ with Omega-background $(h_1, h_2, h_3)$ subject to the CY constraint $h_1 + h_2 + h_3 = 0$, and let $C = \C_{z_1} \hookrightarrow Y$ be a codim-$2$ holomorphic curve with normal bundle $N_{C/Y} \cong \cO_C \oplus \cO_C$ carrying equivariant weights $(h_2, h_3)$ on the two normal coordinates. The 6d holomorphic Chern--Simons action with gauge algebra $\fgl_1$:
\[
 S_{6d} \;=\; \frac{1}{2}\int_Y \Omega \wedge A \wedge \bar{\partial} A
\]
restricted to the tubular neighborhood $U = C \times D^2 \subset Y$ and integrated over the formal bidisk $D^2 \subset N_{C/Y}$ in the equivariant $\bar{\partial}$-regularisation reduces to a 5d holomorphic Chern--Simons action on $C \times \R$ for the $\widehat{\fgl_1}$ Kac--Moody algebra at effective level
\[
 k_{\mathrm{eff}} \;=\; -\sigma_2(h_1, h_2, h_3) \;=\; -(h_1 h_2 + h_1 h_3 + h_2 h_3) \;=\; \tfrac{1}{2}\sum_{i=1}^{3} h_i^2,
\]
where the second equality uses Newton's identity $p_2 = \sigma_1^2 - 2\sigma_2$ under the CY constraint $\sigma_1 = 0$.
\end{lemma}

\begin{proof}
This is the equivariant 1-loop determinant of the $\bar{\partial}$-operator on $D^2$ under the $T^2 = \C^*_{z_2} \times \C^*_{z_3}$ action with weights $(h_2, h_3)$, corrected by the $h_1$ equivariant weight on $C$. The computation is due to Costello (\emph{Supersymmetric gauge theory and the Yangian}, 2013 preprint), with the Costello--Gwilliam factorisation-algebra framework (\emph{Factorization algebras in quantum field theory}, Vol.~II §5) supplying the regularisation axioms. The role of this lemma in Proposition~\ref{prop:codim2-defect-ope} is to justify Arrow~$3$ of the proof: the propagator $\langle A^{(0,0)}(z) A^{(0,0)}(w) \rangle = \Psi/(z-w)^2$ with $\Psi = -\sigma_2$ is the $\widehat{\fgl_1}$ 5d Kac--Moody propagator at level $k_{\mathrm{eff}}$, and the identification $\Psi = k_{\mathrm{eff}}$ is exactly the content of the bridge. The Newton-identity form $\Psi = \tfrac{1}{2}\sum h_i^2$ manifests $\Psi \geq 0$ for real $h_i$ under the CY constraint.
\end{proof}

\begin{remark}[The role of the normal bundle]
\label{rem:normal-bundle-spectral}
The normal bundle $N_{C/Y} = \C^2_{z_2, z_3}$ provides the two spectral parameters $(z_2, z_3)$ of the $E_3$ theory. Under the projection hierarchy $\C^3 \to \C^2 \to C$:
\begin{itemize}
 \item $E_3$ on $\C^3$: two spectral parameters $(z_2, z_3)$ from the full normal bundle;
 \item $\Etwo$ on $\C^2$: one spectral parameter $z_2$ (the Yangian parameter $u$);
 \item $\Eone$ on $C$: no spectral parameter (the chiral algebra on the defect curve).
\end{itemize}
The level $\Psi = -\sigma_2$ is a scalar extracted from the full Omega-background: it is the residue of the classical $r$-matrix $r(u) = -\sigma_2/u$ of the affine Yangian (holomorphic\_cs\_chiral\_engine.py, line~127). The passage from the two-parameter $R$-matrix $R_{\mathrm{ch}}(u, v)$ of Conjecture~\ref{conj:chiral-qg-c3}(ii) to the single-parameter Yangian $R$-matrix is the projection $\C^3 \to \C^2$; the further projection $\C^2 \to C$ forgets the spectral parameter entirely, leaving only the defect algebra with its level $\Psi$.
\end{remark}



\subsection{Consistency comparison: 3d Chern--Simons vs 6d chiral framework}
\label{subsec:cfg25-adversarial}

If the 6d holomorphic theory is a genuine lift of 3d Chern--Simons, every 3d topological result should have a 6d chiral analogue. We compare this expectation against six classical results and catalogue the gaps.

\begin{conjecture}[Chiral analogues of 3d Chern--Simons invariants]
\label{conj:cfg25-adversarial}
\ClaimStatusConjectured
For the 6d holomorphic theory on $\C^3$ with gauge algebra $\fgl_1$ and Omega-background $(h_1, h_2, h_3)$ satisfying $h_1 + h_2 + h_3 = 0$, the following 6d analogues of classical 3d Chern--Simons results should hold:
\begin{enumerate}[label=\textup{(\roman*)}]
 \item \emph{Chiral Kontsevich integral.}
 The perturbative expansion of the 3d CS path integral (the Kontsevich integral $Z^{\mathrm{pert}}(K)$, CFG25 Theorem~A) should lift to a codimension-$2$ defect invariant on $\C^3$: an integral over the configuration space $\Conf_n(C)$ of the holomorphic curve $C \hookrightarrow \C^3$, with the chiral propagator $r_{\mathrm{CY}}(z) = \Psi/z$ replacing the topological propagator. The 3d propagator is smooth; the 6d propagator has poles at $z = -h_i$ from the structure function $g(u) = \prod(u - h_i)/\prod(u + h_i)$, requiring the Omega-background as equivariant regularization.
 \emph{Obstruction}: the codimension-$2$ defect OPE exists (Proposition~\ref{prop:codim2-defect-ope}), but the full chiral Kontsevich integral is not constructed.

 \item \emph{Chiral Jones polynomial.}
 The colored Jones polynomial $J_N(K; q) = \Tr_{V_N}(\cR)$ from quantum traces over the $U_q(\fsl_2)$ $R$-matrix should lift to a ``chiral Jones polynomial'' $J^{\mathrm{ch}}_N(C; q, t) = \Tr_{\mathrm{Fock}_N}(R_{\mathrm{ch}}(u, v))$ using the two-parameter $R$-matrix of the quantum toroidal algebra. The expected output is a triply-graded invariant related to Khovanov--Rozansky homology.
 \emph{Obstruction}: the $R$-matrix exists and satisfies the YBE (Theorem~\ref{thm:yang-r-matrix-ybe}), but the quantum trace is not extracted. The identification with Khovanov--Rozansky is conjectural.

 \item \emph{Chiral volume conjecture.}
 The volume conjecture $\log|J_N(K; e^{2\pi i/N})| \sim N \cdot \mathrm{Vol}(S^3 \setminus K)$ should lift to an asymptotic statement relating the chiral Jones polynomial to the holomorphic Chern--Simons functional $S_{\mathrm{hCS}}$ evaluated on the flat connection associated to $C$.
 \emph{Obstruction}: no chiral volume conjecture is formulated.

 \item \emph{Chiral Verlinde algebra.}
 Quantization of 3d CS on $\Sigma_g \times \R$ produces the Verlinde algebra $V_k(\frakg)$ with $\dim V_k(\fsl_2) = k + 1$. For the 6d theory on $\Sigma_g \times \C^2 \times \R$, the factorization homology $\int_{\Sigma_g} \cA$ gives the genus-$g$ conformal block space. At generic parameters (classes $\mathbf{L}$ and $\mathbf{C}$), $\dim H^*(B_{E_3}(\cA)) = (1+t)^{3g}$, evaluating to $2^{3g}$ at $t = 1$. The factor $3g$ (vs $2g$ at $\Etwo$ or $g$ at $\Eone$) reflects the three factorization directions in $\C^3$. The 3d Verlinde number $V_k = k + 1$ is \emph{not} recovered at generic parameters; it requires root-of-unity truncation of the quantum toroidal algebra.
 \emph{Obstruction}: generic formula exists; root-of-unity truncation is not constructed.

 \item \emph{6-manifold invariant (WRT analogue).}
 The 3d WRT invariant $Z(M^3)$ arises from a $3$-$2$-$1$ extended TFT. The 6d theory should produce a $6$-$5$-$\cdots$-$1$ extended TFT assigning data to manifolds of all dimensions up to~$6$. The current framework assigns data to circles ($\cZ(\Rep^{\Eone}(A))$), tori ($\HH(\HH(A))$), and surfaces ($\int_{\Sigma_g} \cA$), but \emph{not} to $3$-, $4$-, or $5$-manifolds. The partial assignment to $6$-manifolds via $\mathrm{Sp}_4(\Z)$ modularity (Chapter~\ref{ch:k3-times-e}) covers only CY$_3$ targets.
 \emph{Obstruction}: assignments exist for dimensions $1$--$2$; dimensions $3$--$5$ are missing.

 \item \emph{Reshetikhin--Turaev from the quantum toroidal algebra.}
 The 3d RT functor uses $U_q(\frakg)$ at root of unity with the YBE for consistency. At 6d, the Zamolodchikov tetrahedron equation (ZTE) replaces the YBE as the consistency condition. The Yang $R$-matrix satisfies YBE but \emph{fails} ZTE at order $O(\sigma_3^2)$ where $\sigma_3 = h_1 h_2 h_3$ (Theorem~\ref{thm:zte-failure}). The ZTE failure at $\sigma_3 \neq 0$ proves that the $E_3$ structure is genuinely nontrivial and that pairwise $R$-matrices alone do not produce consistent $6$-manifold invariants. Proposition~\ref{prop:zte-explicit-correction} kills the even charge-$2$ obstruction by a rational ternary correction; Proposition~\ref{prop:zte-o-kappa4} leaves rank-$2$ odd-sector residues.
 \emph{Obstruction}: full charge-preserving all-order ZTE completion remains open.
\end{enumerate}
\end{conjecture}

\begin{remark}[The ZTE obstruction as the 6d fingerprint]
\label{rem:zte-6d-fingerprint}
The Zamolodchikov tetrahedron equation failure (item~(vi)) is the deepest structural gap between 3d and 6d. In 3d, the Yang--Baxter equation (the $2$-simplex consistency) suffices for all knot and $3$-manifold invariants. In 6d, the ZTE (the $3$-simplex consistency) is the analogue, and it \emph{fails} for the Yang $R$-matrix. The failure is a quantum effect: it vanishes at $\sigma_3 = 0$ (Kapranov--Voevodsky triviality at the permutation limit) and scales as $O(\sigma_3^2)$. This scaling proves that the $E_3$ corrections live at the \emph{chain level}. The 6d theory is not a ``dimensional upgrade'' of 3d CS; it is a genuinely different structure requiring new algebraic inputs beyond quantum groups.
\end{remark}


\subsection{Systematic obstruction catalogue: what cannot be lifted from 3d to 6d}
\label{subsec:3d-6d-obstruction-catalogue}

The six consistency vectors above decide whether specific 3d
\emph{invariants} extend to 6d.  A deeper question is whether the
underlying \emph{structural features} of 3d CS can be lifted at all.
There are eight structural obstructions: each is either a fundamental
no-go or an unsolved problem, and each has a definite parameter locus
where a 3d feature can be partially recovered.

\medskip

\begin{center}
\small
\renewcommand{\arraystretch}{1.3}
\begin{tabular}{clccl}
\toprule
\# & \textbf{Obstruction} & \textbf{Classification} & \textbf{Depth} & \textbf{Loophole} \\
\midrule
1 & Finite-dimensionality (Verlinde) & Fundamental & Parameter & Root-of-unity locus (codim-$2$) \\
2 & Knot complement topology & Fundamental & Topological & None \\
3 & Surgery formula (Kirby) & Unsolved & Structural & $E_3$ corrections to ZTE \\
4 & Reshetikhin--Turaev functor & Fundamental & Topological & Non-semisimple RT \\
5 & Categorification direction & Unsolved & Structural & M-theory (7d) \\
6 & Unitarity & Fundamental & Parameter & NS limit or imaginary $h_i$ \\
7 & Volume conjecture & Fundamental & Parameter & Rigid CY$_3$ ($h^{2,1} = 0$) \\
8 & Rationality & Fundamental & Parameter & Gepner points \\
\bottomrule
\end{tabular}
\end{center}

\medskip
\noindent The eight obstructions cluster into three types.

\paragraph{Topological obstructions (2, 4).}
These are fundamental: the 6d theory is \emph{holomorphic}, not topological, and the relevant topologies differ irreconcilably. In 3d, the complement $S^3 \setminus K$ of any knot has boundary $T^2$ (the normal bundle is trivial: rank~$2$, oriented, over $S^1$). In 6d, the complement of a holomorphic curve $C$ of genus~$g$ in a CY$_3$ has boundary determined by the normal bundle $N_{C/X}$ with $\deg N_{C/X} = 2g - 2$ (adjunction and $c_1(X) = 0$). Only genus-$1$ curves (elliptic, $\deg N = 0$) have toroidal boundary $T^4$; genus-$0$ curves have $S^3$-type boundaries, and higher-genus curves have still more complex topology. The Reshetikhin--Turaev functor requires a modular tensor category with finitely many simple objects and a non-degenerate $S$-matrix; $\Rep(U_{q,t}(\widehat{\widehat{\fgl}}_1))$ has infinitely many simples (Fock modules $F_n$ for each $n \in \Z$), is non-semisimple for class~$M$, and lacks a known ribbon structure. The non-semisimple RT framework of Costantino--Geer--Patureau-Mirand (2014) provides a conjectural route via modified quantum traces.

\paragraph{Parameter obstructions (1, 6, 7, 8).}
The 3d theory has a \emph{discrete} coupling ($k \in \Z$) while the 6d theory has \emph{continuous} parameters $(h_1, h_2, h_3)$. At special parameter loci, some 3d features are recovered:
\begin{itemize}
 \item \emph{Root of unity}: $q^N = t^M = 1$ (codimension-$2$ in parameter space) should truncate the representation category to finitely many simples. But the truncated quantum toroidal algebra is \emph{not constructed}.
 \item \emph{Unitarity}: algebraic unitarity $R(z)\,R(-z) = \mathrm{Id}$ holds identically (from $g(-z) = 1/g(z)$, a consequence of the CY condition). Hilbert-space unitarity (positive-definite Shapovalov form) \emph{fails} at generic real $h_i$: the Shapovalov norms $\prod_{k=1}^{n} g(k h_1)$ are indefinite. Unitarity is recovered at the topological locus ($h_i$ purely imaginary, giving $|q| = 1$) and at the Nekrasov--Shatashvili limit ($h_2 = 0$, effective 3d theory).
 \item \emph{Volume conjecture}: the 3d conjecture relates quantum invariants to hyperbolic volume via Mostow rigidity (unique hyperbolic structure). CY$_3$ manifolds have positive-dimensional moduli spaces: no Mostow analogue. For \emph{rigid} CY$_3$ ($h^{2,1} = 0$, e.g., the Schoen threefold), the holomorphic volume is unique up to normalization, and a volume conjecture may be formulable.
 \item \emph{Rationality}: $V_k(\frakg)$ at integer level is a rational VOA ($C_2$-cofinite, finitely many simples, modular characters). CY chiral algebras are generically \emph{irrational}: infinitely many charge sectors, non-$C_2$-cofinite, mock-modular characters (class~$M$). Rationality is recovered only at Gepner points (measure zero in CY moduli), where the theory truncates to a tensor product of $N = 2$ minimal models.
\end{itemize}

No single parameter locus recovers \emph{all} 3d features simultaneously: unitarity requires imaginary $h_i$, rationality requires the Gepner locus, and finite-dimensionality requires a root-of-unity specialization. The 3d theory is not ``embedded'' in the 6d theory at any single point.

\paragraph{Structural obstructions (3, 5).}
These are \emph{unsolved problems}: the obstruction may be resolvable by new mathematics.
\begin{itemize}
 \item \emph{Surgery}: 6d handle decomposition exists (Proposition~\ref{prop:handle-decomposition-k3}), but the ZTE failure (Theorem~\ref{thm:zte-failure}) means handle rearrangements are not automatically consistent. ZTE corrections (the $E_3$ correction term $T$) may restore consistency, but sufficiency is not proved.
 \item \emph{Categorification}: Khovanov homology categorifies the Jones polynomial (3d $\to$ 4d). The 6d $\to$ 7d categorification should arise from M-theory on a CY$_2$ (e.g., K3), producing a topological twist of 11d supergravity. This is physically motivated but mathematically unconstructed. The decategorification is consistent: $\dim_{E_3}(g) = \dim_{E_2}(g) \cdot \dim_{E_1}(g)$ by the K\"unneth formula (Dunn additivity).
\end{itemize}

\begin{remark}[Irreducible 6d content]
\label{rem:irreducible-6d-content}
The obstructions are not \emph{defects} of the 6d theory but \emph{signatures} of its richer structure. Features that exist only at 6d, with no 3d analogue:
\begin{enumerate}[label=\textup{(\roman*)}]
 \item Two-parameter $R$-matrix $R(u, v)$ with Zamolodchikov factorization;
 \item Miki automorphism ($\mathbb{Z}/3$ cyclic on $(q, t, s)$ on the positive half / $\mathrm{CoHA}(\C^3) = Y^+$, with ambient $S_3$ on the Drinfeld double; algebra-specific);
 \item ZTE corrections at $O(\sigma_3^2)$ (genuinely new $E_3$ mathematics);
 \item Shadow tower / $A_\infty$ coproduct (infinite corrections $\delta^{(k)} = S_k$);
 \item Mock modularity for class~$M$ characters;
 \item Non-semisimple Drinfeld center (richer than semisimple MTC of 3d).
\end{enumerate}
The 6d theory produces genuinely new mathematical objects (holomorphic curve invariants, CY partition functions, mock modular forms) that do not reduce to classical knot/manifold invariants. The correct relationship between 3d and 6d is not one of ``lift'' but of two outputs of the same $E_3$ Koszul duality formalism, applied to different inputs: topological $E_3$ (framed little $3$-disks on $\R^3$) for 3d, and holomorphic $E_3$ (Omega-background on $\C^3$) for 6d.
\end{remark}

\subsection{The topological--algebraic dichotomy: what the 6d theory produces}
\label{subsec:topological-algebraic-dichotomy}

The obstruction list separates the 3d Chern--Simons structures that
survive 6d holomorphic CS from those that do not.  The answer is a
clean dichotomy: 3d CS produces \emph{topological} invariants, numbers
depending on a manifold or knot, while 6d hCS produces
\emph{algebraic} invariants, rational functions of spectral parameters
encoding the representation theory of quantum algebras. The surviving
structures are precisely the algebraic ones.

\paragraph{The dichotomy.}
We decompose all $34$ substructures across the eight obstruction categories into three types by the mathematical nature of their output:
\begin{itemize}
 \item \emph{Algebraic}: the output is a rational function of spectral parameters, an operator-valued formal series, or a structural datum of an infinite-dimensional algebra. Examples: $R$-matrix satisfying YBE, coproduct $\Delta_z$, structure function $g(z)$, algebraic unitarity $R(z)\,R(-z) = \mathrm{Id}$, fusion product, $T$-matrix.
 \item \emph{Topological}: the output is a number (or finite collection of numbers) that depends on a manifold type, a knot, or a topological invariant, but not on any continuous parameter. Examples: Verlinde dimensions, knot group $\pi_1(S^3 \setminus K)$, Alexander polynomial, Kirby invariance, $S$-matrix (modular, non-degenerate), Khovanov homology, hyperbolic volume.
 \item \emph{Hybrid}: the output has both algebraic and topological aspects. Examples: handle decomposition (topological existence, algebraic gluing data), modular characters (algebraic functions with topological transformation properties).
\end{itemize}
Of the $34$ substructures, $12$ are algebraic, $19$ are topological, and $3$ are hybrid. Among the algebraic structures, the lift rate to 6d is $100\%$: every algebraic structure of 3d CS lifts to a (richer) algebraic structure in 6d. Among the topological structures, the lift rate is $0\%$: no topological structure lifts. The hybrid structures have mixed behavior. The overall lift rate is $12/34 \approx 35\%$; the algebraic fraction of the total catalogue is $12/34$.

\begin{center}
\small
\renewcommand{\arraystretch}{1.3}
\begin{tabular}{lccc}
 \toprule
 \textbf{Invariant type} & \textbf{Count} & \textbf{Lift to 6d} & \textbf{Lift rate} \\
 \midrule
 Algebraic & $12$ & $12$ & $100\%$ \\
 Topological & $19$ & $0$ & $0\%$ \\
 Hybrid & $3$ & $0$ & $0\%$ \\
 \midrule
 \textbf{Total} & $\mathbf{34}$ & $\mathbf{12}$ & $\mathbf{35\%}$ \\
 \bottomrule
\end{tabular}
\end{center}

\paragraph{What the 6d theory produces.}
The 6d holomorphic CS theory is not a failed attempt at knot/manifold invariants. It produces eight algebraic classes of invariants, none of which are knot or manifold invariants:
\begin{enumerate}[label=\textup{(\roman*)}]
 \item \emph{Structure functions} $g(z) = \prod(z - h_i)/\prod(z + h_i)$: rational functions of a spectral parameter encoding the full representation theory. The 3d analogue is the quantum dimension $\dim_q(V)$ -- a \emph{number}, not a function. The entire $z$-dependence is genuinely 6d.
 \item \emph{Two-parameter $R$-matrices} $R(u,v)$ with Zamolodchikov factorization. The 3d analogue is the one-parameter $R(u)$ of $U_q(\frakg)$. The second parameter arises from the Omega-background (normal bundle grading does not directly give $(q,t)$; the passage is through equivariant localization).
 \item \emph{Spectral coproducts} $\Delta_z \colon A \to A \otimes A[\![z]\!]$, with $z$-degree at spin~$s$ equal to~$s$. The 3d coproduct $\Delta \colon A \to A \otimes A$ has no spectral parameter; the $z$-dependence encodes the ordered factorization ($\Eone$).
 \item \emph{Shadow towers} $\{S_k\}_{k \geq 1}$: the $\Ainf$-correction coefficients $\delta^{(k)}$ of the coproduct, equivalently the Maurer--Cartan expansion of the shadow. Class $\mathbf{G}$ (formal) has $S_k = 0$ for $k \geq 2$; class $\mathbf{M}$ (mock modular) has $S_k \neq 0$ for all~$k$. No 3d analogue exists (the 3d coproduct is exact).
 \item \emph{Mock modular partition functions}: for class~$\mathbf{M}$ algebras, the characters are mock modular forms (not genuine modular), arising from the infinite shadow tower. The 3d characters are genuine modular forms (Kac--Peterson).
 \item \emph{ZTE corrections}: the even charge-$2$ obstruction is killed by a rational ternary correction, while the full charge-preserving tetrahedron solution requires the odd sectors as well. No 3d analogue exists, since the YBE suffices in 3d.
 \item \emph{Miki automorphism}: the $\mathbb{Z}/3$ cyclic permutation of $(q,t,s)$ from the CY$_3$ constraint $\sum\epsilon_i = 0$, acting on the positive half $\mathrm{CoHA}(\C^3) = Y^+$; the ambient $S_3$ lives on the Drinfeld double after a $\mathbb{Z}/2$-involution. Algebra-specific, not operadic.
 \item \emph{CY partition functions} (DT/PT/GW): generating functions of curve-counting invariants in CY threefolds. Entirely 6d: 3d has no curve counting.
\end{enumerate}

\paragraph{The universal specialization principle.}
The relationship between 3d and 6d is that of specialization to universality:
\begin{center}
\renewcommand{\arraystretch}{1.3}
\small
\begin{tabular}{lcll}
 \toprule
 \textbf{Dimension} & \textbf{Parameters} & \textbf{Output type} & \textbf{Algebra} \\
 \midrule
 3d ($\C \times \R$) & $1$ (level $k$, discrete) & topological (numbers) & $\widehat{\frakg}$ (Kac--Moody) \\
 5d ($\C^2 \times \R$) & $1$ ($\hbar$, continuous) & algebraic ($f(u)$) & $Y(\frakg)$ (Yangian) \\
 6d ($\C^3$) & $2$ ($(q,t)$, continuous) & algebraic ($f(u,v)$) & $U_{q,t}(\widehat{\widehat{\fgl}}_1)$ \\
 \bottomrule
\end{tabular}
\end{center}
Each step \emph{up} in dimension universalizes the structure: discrete parameters become continuous, numbers become functions, finite-dimensional algebras become infinite-dimensional. The passage \emph{down} recovers the lower-dimensional theory by specialization:
\begin{equation}\label{eq:universal-specialization-bridge}
 \underbrace{U_{q,t}(\widehat{\widehat{\fgl}}_1)}_{\text{6d, generic $(q,t)$}}
 \;\xrightarrow[\text{codim-$2$}]{\;\;q^N = t^M = 1\;\;}
 \underbrace{U_{q,t}^{\mathrm{trunc}}}_{\text{truncated}}
 \;\xrightarrow[\text{codim-$1$}]{\;\;t \to 1\;\;}
 \underbrace{U_q(\fsl_N)\;\text{at root of unity}}_{\text{3d, level }k = N - h^\vee}
\end{equation}
\emph{Neither step is constructed.} Arrow~1 requires the truncated quantum toroidal algebra at joint root of unity (the Hopf ideal, its compatibility with $\Delta_z$, and the $E_3$ factorization of the quotient are all open). Arrow~2 requires understanding the $t \to 1$ limit of the Miki automorphism. The combined bridge is an open problem.

\begin{conjecture}[Root-of-unity bridge]
\label{conj:root-of-unity-bridge}
\ClaimStatusConjectured
The two-step specialization~\eqref{eq:universal-specialization-bridge} recovers the 3d Chern--Simons quantum group $U_q(\frakg)$ at root of unity from the 6d quantum toroidal algebra $U_{q,t}(\widehat{\widehat{\fgl}}_1)$ at generic parameters. The truncated quotient $U_{q,t}^{\mathrm{trunc}}$ at $q^N = t^M = 1$ should carry a finite set of simple modules whose fusion rules, in the $t \to 1$ limit, reproduce the Verlinde algebra $V_k(\frakg)$.
\end{conjecture}

\paragraph{Why the volume conjecture has \texorpdfstring{$0\%$}{0 percent} lift rate.}
The volume conjecture is the most deeply 3d structure in the catalogue: it lives at the intersection of algebra and topology, relating the colored Jones polynomial $J_N(K; e^{2\pi i/N})$ (an algebraic quantity evaluated at a topological point) to the hyperbolic volume $\mathrm{Vol}(S^3 \setminus K)$ (a topological quantity requiring Mostow rigidity). At the 6d level, \emph{neither} ingredient exists. The colored Jones polynomial requires root-of-unity specialization (Arrow~1 of the bridge, unconstructed). The hyperbolic volume requires Mostow rigidity, which fails for CY$_3$ manifolds with positive-dimensional moduli ($h^{2,1} > 0$). The volume conjecture is the canary in the coal mine: its total failure to lift is the clearest evidence that the 6d theory is algebraic, not topological.

\begin{remark}[The correct relationship]
\label{rem:correct-3d-6d-relationship}
The 6d holomorphic CS theory is the \emph{universal algebraic completion} of the 3d--5d tower. It produces the full representation-theoretic data (structure functions, $R$-matrices, coproducts, shadow towers) from which all lower-dimensional specializations should be extractable (Conjecture~\ref{conj:root-of-unity-bridge}). It does \emph{not} produce knot or manifold invariants directly; those arise only after the two-step bridge (truncation $+$ NS limit) to 3d. The analogy is
\[
 \text{6d} : \text{3d} \;\;::\;\; \Z[x] : \Z/n\Z
 \quad\text{(polynomial ring : quotient ring)},
\]
where the polynomial ring $\Z[x]$ contains all the information about all quotients $\Z[x]/(x^n - 1)$, but it is not itself a finite ring, and the specific properties of each quotient (e.g., the number of units) emerge only upon specialization.
\end{remark}

\noindent The topological--algebraic classification, lift rates,
root-of-unity bridge, 6d output catalogue, dimensional hierarchy, and
\(\kappa_\bullet\)-convention split are the data used in this dichotomy.


\section{The Costello--Paquette defect chiral algebra and the \texorpdfstring{$5d \to 6d$}{5d -> 6d} boost}
\label{sec:costello-paquette-defect}

The Costello--Paquette construction (arXiv:2104.14573, arXiv:2208.00354) identifies the boundary chiral algebra of holomorphic Chern--Simons theory as the endomorphism algebra of a \emph{universal defect}: a canonical line defect $D_{\mathrm{univ}}$ whose module category exhausts all Wilson line modules. The construction is established at $5d$ (Costello 2013; Costello--Paquette 2021), producing the affine Yangian $Y(\widehat{\fgl}_1)$, and admits a conjectural boost to $6d$, where the target is the quantum toroidal algebra $U_{q,t}(\widehat{\widehat{\fgl}}_1)$. The $5d$ universal defect and its $6d$ lift culminate in the form factor map $\mathsf{ff} \colon Z^{\mathrm{der}}_{\mathrm{ch}}(A) \to A^!$ bridging the bulk algebra to the defect algebra.

\subsection{The universal defect at each dimension}
\label{subsec:cp-universal-defect-dims}

\begin{construction}[Costello--Paquette universal defect]
\label{constr:cp-universal-defect}
Let $M$ be one of the three ambient spaces in the holomorphic CS hierarchy (Construction~\ref{constr:hol-cs-hierarchy}). The \emph{universal defect} $D_{\mathrm{univ}}$ is a canonical codimension-$2$ defect supported on a holomorphic curve $C \hookrightarrow M$, characterised by the universality property:
\[
 \End(D_{\mathrm{univ}}) \;=\; A, \qquad \Hom(D_{\mathrm{univ}}, W_V) \;=\; V
\]
where $A$ is the boundary chiral algebra and $W_V$ is the Wilson line in representation~$V$. The defect $D_{\mathrm{univ}}$ is the unit object in the category of line defects: every Wilson line factors as $W_V = D_{\mathrm{univ}} \otimes_A V$. The data at each dimension:
\begin{center}
\small
\renewcommand{\arraystretch}{1.2}
\begin{tabular}{clccll}
\toprule
$\dim_\C(M)$ & $M$ & $\rank N_{C/M}$ & $\#$ spectral & $\End(D_{\mathrm{univ}})$ & Construction locus \\
\midrule
$1$ & $\Sigma \times \R$ & $0$ & $0$ & $V_k(\fgl_1)$ & Costello--Gwilliam boundary theorem \\
$2$ & $\C^2 \times \R$ & $1$ & $1$ ($u$) & $Y(\widehat{\fgl}_1)$ & Costello Yangian theorem \\
$3$ & $\C^3$ & $2$ & $2$ ($u, v$) & $U_{q,t}(\widehat{\widehat{\fgl}}_1)$ & requires 6d hCS completion \\
\bottomrule
\end{tabular}
\end{center}
The number of spectral parameters equals the rank of the normal bundle $N_{C/M}$, which is $\dim_\C(M) - 1$. The $\En$ level of the defect algebra equals $\dim_\C(M)$.
\end{construction}

\begin{proposition}[Structural invariants of the universal defect]
\label{prop:cp-defect-invariants}
For the universal defect $D_{\mathrm{univ}}$ in holomorphic CS on $M = \C^n$ with gauge algebra $\fgl_1$ and Omega-background $(h_1, h_2, h_3)$ satisfying $h_1 + h_2 + h_3 = 0$:
\begin{enumerate}[label=\textup{(\roman*)}]
 \item The effective level of the defect algebra is $\Psi = -\sigma_2 = -(h_1 h_2 + h_1 h_3 + h_2 h_3)$, independent of the dimension~$n$. In particular, $\kappa_{\mathrm{ch}}(\End(D_{\mathrm{univ}})) = \Psi$ at all three dimensions.
 \item The Koszul conductor on the free-field branch is $K = \kappa_{\mathrm{ch}}(A) + \kappa_{\mathrm{ch}}(A^!) = 0$ (the Koszul dual has $\kappa_{\mathrm{ch}}(A^!) = -\Psi$).
 \item At the self-dual point ($\sigma_3 = h_1 h_2 h_3 = 0$): $g(u) = 1$, the $R$-matrix is trivial, and $D_{\mathrm{univ}}$ is the trivial defect.
 \item The equivariant weight of the normal mode $A^{(m,n)}$ at $\dim_\C = 3$ is $m \cdot h_2 + n \cdot h_3$, with $m + n + 1$ giving the conformal weight (spin) of the corresponding $W_{1+\infty}$ current.
\end{enumerate}
\end{proposition}

\begin{proof}
The $\kappa_{\mathrm{ch}}$-preservation across dimensions is structural:
$\Psi=-\sigma_2$ depends only on the $\Omega$-background parameters, not
on the ambient dimension.  The same parameters therefore give the same
classical $r$-matrix $r(u)=-\sigma_2/u$ at every $E_n$ level.  The
self-dual triviality follows from
$g(u)=\prod_i(u-h_i)/\prod_i(u+h_i)=1$ when one $h_i$ is zero and
$h_1+h_2+h_3=0$.  The normal-mode weight in dimension three is the
tautological $\Omega$-weight $m h_2+n h_3$, and the shift by one gives
the conformal weight of the corresponding $W_{1+\infty}$ current.
\end{proof}

\subsection{The form factor map: bulk-to-boundary Koszul projection}
\label{subsec:cp-form-factor}

The form factor map of Costello--Paquette encodes the bulk-to-boundary OPE as a map from the derived center (the bulk algebra) to the Koszul dual (the defect algebra). This is the bridge between holomorphic CS observables and quantum group representations.

\begin{construction}[Form factor map]
\label{constr:form-factor-map}
Let $A$ be the boundary chiral algebra of holomorphic CS on $M$. The \emph{form factor map} is the composition
\[
 \mathsf{ff} \colon Z^{\mathrm{der}}_{\mathrm{ch}}(A) \xrightarrow{\;\text{HH-to-bar}\;} B(A) \xrightarrow{\;D_{\Ran}\;} A^!
\]
where the first arrow is the Hochschild-to-bar comparison map (Vol~I Theorem~H, inverse direction) and the second is the Verdier duality on $\Ran(C)$. At the level of modes, for a bulk operator $O$ of conformal weight $\Delta_O$:
\[
 \mathsf{ff}(O)(|v\rangle) \;=\; \Res_{z=0}\, z^{\Delta_O - 1}\, O(z)\, |v\rangle,
\]
extracting the singular part of the bulk-to-boundary OPE: precisely the data captured by the bar complex $B(A)$.
\end{construction}

\begin{proposition}[Form factor axioms]
\label{prop:form-factor-axioms}
The form factor map $\mathsf{ff} \colon Z^{\mathrm{der}}_{\mathrm{ch}}(A) \to A^!$ satisfies:
\begin{enumerate}[label=\textup{(FF\arabic*)}]
 \item $\mathsf{ff}$ is a map of $A$-bimodules (the bulk-boundary OPE is $A$-equivariant from both left and right).
 \item $\mathsf{ff}(\mathbf{1}_{\mathrm{bulk}}) = \mathbf{1}_{A^!}$ (the identity form factor is the Koszul identity).
 \item For $n \geq 2$, $\mathsf{ff}$ intertwines the $E_n$ braiding on the bulk with the Koszul-reflected braiding on $A^!$.
 \item $\kappa_{\mathrm{ch}}$ is preserved: $\kappa_{\mathrm{ch}}(\mathsf{ff}(O)) = \kappa_{\mathrm{ch}}(O)$ for all $O$, and the Koszul conductor on the free-field branch is $K = 0$.
 \item At the self-dual point ($\sigma_3 = 0$), $\mathsf{ff}$ is the identity (the defect is trivial).
\end{enumerate}
The form factor pole structure at each spin:
\begin{center}
\small
\renewcommand{\arraystretch}{1.15}
\begin{tabular}{ccll}
\toprule
\textbf{Spin} & \textbf{Pole order} & \textbf{Residue} & \textbf{Identification} \\
\midrule
$1$ & $1$ & $\Psi = -\sigma_2$ & Classical $r$-matrix $r^{\mathrm{Heis}}(z) = \Psi/z$ \\
$2$ & $3$ & $c/2 = 1/2$ & Classical $r$-matrix $r^{\mathrm{Vir}}(z) = (c/2)/z^3 + 2T/z$ \\
\bottomrule
\end{tabular}
\end{center}
\end{proposition}

\begin{proof}
Axioms~(FF1)--(FF5) follow from the factorisation $\mathsf{ff} = D_{\Ran} \circ B$, the properties of the bar complex (Vol~I), and the Verdier duality (which preserves bimodule structure and $\kappa_{\mathrm{ch}}$). The pole orders and residues agree with Proposition~\ref{prop:codim2-defect-ope}: the spin-$1$ residue $\Psi$ is the $J_{(1)}J$ OPE coefficient, and the spin-$2$ residue $c/2 = 1/2$ is $T_{(3)}T$. The independence of the spin-$2$ form factor from $\Psi$ follows from $c = 1$ for $\fgl_1$ at any level (Proposition~\ref{prop:codim2-defect-ope}).
\end{proof}

\subsection{The \texorpdfstring{$5d \to 6d$}{5d -> 6d} dimensional boost}
\label{subsec:cp-5d-to-6d-boost}

The Costello--Paquette construction at $5d$ is established: the affine Yangian $Y(\widehat{\fgl}_1)$ is the boundary chiral algebra of $5d$ holomorphic CS on $\C^2 \times \R$ (Theorem~\ref{thm:5d-boundary-yangian}). The \emph{boost} to $6d$ is the conjectural lift of this construction to the $6d$ holomorphic theory on $\C^3$, producing the quantum toroidal algebra $U_{q,t}(\widehat{\widehat{\fgl}}_1)$.

\begin{conjecture}[Dimensional boost: \texorpdfstring{$5d \to 6d$}{5d -> 6d}]
\label{conj:cp-5d-to-6d-boost}
\ClaimStatusConjectured
The Costello--Paquette universal defect construction lifts from $5d$ to $6d$ with the following structural changes:
\begin{center}
\small
\renewcommand{\arraystretch}{1.2}
\begin{tabular}{lll}
\toprule
\textbf{Feature} & \textbf{$5d$ theorem} & \textbf{$6d$ lift} \\
\midrule
Spectral parameters & $1$ ($u$) & $2$ ($u, v$) \\
$E_n$ level & $\Etwo$ & $E_3$ \\
Boundary algebra & $Y(\widehat{\fgl}_1)$ & $U_{q,t}(\widehat{\widehat{\fgl}}_1)$ \\
$R$-matrix & $R(u)$, satisfies YBE & $R(u,v)$, should satisfy parametric YBE \\
Consistency eqn & Yang--Baxter (2-simplex) & Zamolodchikov tetrahedron (3-simplex) \\
$\kappa_{\mathrm{ch}}$ & $-\sigma_2$ & $-\sigma_2$ (preserved) \\
Affinization & single ($\widehat{\fgl}_1$) & double ($\widehat{\widehat{\fgl}}_1$) \\
\bottomrule
\end{tabular}
\end{center}
The boost is not dimensional reduction in reverse: the $6d$ theory has genuinely new structure (the second affinization, the Miki automorphism, and the ZTE obstruction) that cannot be obtained from $5d$ alone.

The two-parameter structure function of the $6d$ defect at charge~$1$ is
\begin{equation}\label{eq:cp-two-param-structure}
 G_{\mathrm{ch}}(u, v) \;=\; g(u) \cdot g(v) \cdot g(u - v),
\end{equation}
where $g(u) = \prod_{i=1}^3(u - h_i)/\prod_{i=1}^3(u + h_i)$ is the Yangian structure function. This is the Zamolodchikov factorisation of the charge-$1$ two-parameter $R$-matrix. At charge $\geq 2$, the factorised ansatz $S_{ijk} = R_{ij}R_{ik}R_{jk}$ fails the Zamolodchikov tetrahedron equation at $O(\sigma_3^2)$ where $\sigma_3 = h_1 h_2 h_3$ (Theorem~\ref{thm:zte-failure}), proving the $E_3$ structure is genuinely nontrivial.
\end{conjecture}

\begin{remark}[What is preserved and what changes]
\label{rem:cp-boost-preservation}
The $\kappa_{\mathrm{ch}}$-preservation under the boost ($\kappa_{\mathrm{ch}} = -\sigma_2$ at both $5d$ and $6d$) is a structural consequence: the effective level $\Psi = -\sigma_2$ of the Heisenberg subalgebra on the defect curve~$C$ depends only on the Omega-background parameters, not on the ambient dimension. What \emph{changes} is the structure of the higher-spin sectors:
\begin{itemize}
 \item At $5d$: the spin-$s$ currents form the $\Etwo$-braided algebra $Y(\widehat{\fgl}_1)$, with one spectral parameter $u$ from the normal direction.
 \item At $6d$: the spin-$s$ currents form the $E_3$-braided algebra $U_{q,t}(\widehat{\widehat{\fgl}}_1)$, with two spectral parameters $(u, v)$ from the rank-$2$ normal bundle $N_{C/\C^3} = \C^2$.
\end{itemize}
The second spectral parameter $v$ is the new ingredient at $6d$; it generates the second affinization (the second hat in $\widehat{\widehat{\fgl}}_1$). The $\Etwo \to E_3$ promotion is the holomorphic-geometric incarnation of the derived center construction (Remark~\ref{rem:deligne-hochschild}): at the level of representation categories, $\cZ(\Rep^{\Etwo}(Y)) \simeq \Rep^{E_3}(Z^{\mathrm{der}}_{\mathrm{ch}}(Y))$, where the right side is conjectural.
\end{remark}

\subsection{The K3 universal defect}
\label{subsec:cp-k3-defect}

For $X = K3 \times E$, the holomorphic CS hierarchy specialises to three boundary algebras, each controlled by the Mukai lattice $\Lambda_{K3}$ of signature~$(4, 20)$.

\begin{conjecture}[K3 universal defect hierarchy]
\label{conj:cp-k3-defect-hierarchy}
\ClaimStatusConjectured
The universal defect $D_{\mathrm{univ}}$ for K3 targets produces the following boundary algebras:
\begin{center}
\small
\renewcommand{\arraystretch}{1.2}
\begin{tabular}{clcll}
\toprule
$\dim_\C(M)$ & Ambient & $\kappa_{\mathrm{ch}}$ & Boundary algebra & Construction locus \\
\midrule
$1$ & K3 & $2$ & $V_k(\mathfrak{g}_{K3})$ & CY-A$_2$ \\
$2$ & $K3 \times \C \times \R$ & $3$ & $Y(\mathfrak{g}_{K3})$ & requires 5d K3 hCS construction \\
$3$ & $K3 \times E \times \C$ & $3$ & $U_{q,t}(\widehat{\mathfrak{g}}_{K3})$ & requires 6d non-Lagrangian hCS construction \\
\bottomrule
\end{tabular}
\end{center}
The K3 structure function is degree-$(24, 24)$:
\[
 g_{K3}(u) \;=\; \prod_{\alpha \in \Phi_{K3}^+} \frac{u - \alpha(h)}{u + \alpha(h)},
\]
where $\Phi_{K3}^+$ is the positive root system of the Mukai lattice, with $|\Phi_{K3}^+| = 24$. The structure function satisfies the reflection identity $g_{K3}(u) \cdot g_{K3}(-u) = 1$.

The Wilson lines wrapping curves in $K3$ are indexed by Mukai vectors $v \in H^*(K3, \Z)$, with charge lattice $\Lambda_{K3}$ of signature~$(4, 20)$, and fusion controlled by the Mukai pairing.

The $\kappa_\bullet$-spectrum, distinguishing total-space from fiber contributions:
\[
 \kappa_{\mathrm{cat}} = 0, \quad \kappa_{\mathrm{ch}}^{\mathrm{Hodge}} = 0, \quad \kappa_{\mathrm{ch}}^{\mathrm{Heis}} = 3, \quad \kappa_{\mathrm{BKM}}(\Delta_5) = 5, \quad \kappa_{\mathrm{fiber}} = 24.
\]
Here $\kappa_{\mathrm{cat}}(K3 \times E) = \chi(\cO_{K3 \times E}) = \chi(\cO_{K3}) \cdot \chi(\cO_E) = 2 \cdot 0 = 0$ by K\"unneth; the compact total-space Hodge supertrace $\kappa_{\mathrm{ch}}^{\mathrm{Hodge}}(A_{K3\times E})$ is also $0$, and the K3-fiber value $\chi(\cO_{K3}) = 2$ is a \emph{fiber} invariant, not the total-space one.
The canonical Stage-2 Heisenberg--Cartan additivity is instead $\kappa_{\mathrm{ch}}^{\mathrm{Heis}}(K3 \times E) = \kappa_{\mathrm{ch}}^{\mathrm{Heis}}(K3) + \kappa_{\mathrm{ch}}^{\mathrm{Heis}}(E) = 2 + 1 = 3$. The two Koszul branches are distinct: free-field/Heisenberg ($\fgl_1$-sector, $K = 0$, $\kappa_{\mathrm{ch}}^{\mathrm{Heis}}(A^!) = -3$) and Mukai-doubling/BKM ($\kappa_{\mathrm{ch}}^{\mathrm{Mukai}}(A) = 4$, $\kappa_{\mathrm{ch}}^{\mathrm{Mukai}}(A^!) = 4$, conductor $K^{\kappa_{\mathrm{ch}}^{\mathrm{Mukai}}} = 8 = 2\,c_+(\mathrm{Mukai}(K3))$; see Theorem~\ref{thm:qca-hbar-minus-eighth-lusztig} and Remark~\ref{rem:qca-three-faces-of-eight}). The Borcherds weight $\kappa_{\mathrm{BKM}}(\Delta_5) = c_{\Delta_5}(0)/2 = 5$ is a level-$4$ automorphic invariant of the input Siegel form, not the Koszul conductor. The numerical equality \(5=3+2\) at \(N=1\) mixes a Borcherds weight, a Heisenberg rank, and a fibre Euler characteristic; it is not an invariant identity. The formula $\kappa_{\mathrm{BKM}}(\Phi_N)=\kappa_{\mathrm{ch}}^{\mathrm{Heis}}+\chi(\cO_{\mathrm{fiber}})$ fails as a statement of invariants for every \(N\in\{1,2,3,4,6\}\).
\end{conjecture}



\section{The quantum chiral Koszul dual}
\label{sec:qca-koszul-dual}

The Koszul dual of a boundary chiral algebra is the defect algebra. The $\En$ structure on the defect algebra is inherited from the ambient holomorphic theory: a deformation of the chiral CE cochains by the theory's quantum parameters. The deficiency: the $E_n$ level of $A^!$ is determined by the ambient geometry, not by $A$ alone, and the 6d theory on $\C^3$ produces an $E_3$-chiral deformation that no lower-dimensional argument can access.

\subsection{Chiral Koszul duality: boundary and defect}
\label{subsec:chiral-koszul-bd}

\begin{definition}[Chiral Koszul dual]
\label{def:chiral-koszul-dual}
For a chiral algebra $A$ on a curve $C$ with bar complex $B(A) = T^c(s^{-1}\bar{A})$ (Definition~\ref{def:ordered-bar}, Chapter~\ref{ch:e1-chiral}), the \emph{chiral Koszul dual} is
\[
 A^! \;:=\; D_{\Ran}(B(A)),
\]
the Verdier dual of the bar complex as a factorization coalgebra on $\Ran(C)$. This is the Volume~I Verdier leg (Theorem~A), specialized to the CY setting. The algebra $A^!$ is the \emph{defect algebra}: it controls the category of line operators that can end on the boundary where $A$ lives.
\end{definition}

\begin{proposition}[Three dualities are distinct]
\label{prop:three-dualities}
The chiral Koszul dual $A^!$ is distinct from:
\begin{enumerate}[label=\textup{(\roman*)}]
 \item \emph{Bar-cobar inversion}: $\Omega(B(A)) \simeq A$ recovers the original algebra (the counit face of Vol~I Theorem~A on the Koszul locus; the strict ML weight-completed coderived strengthening is Vol~I Theorem~B on all four classes G/L/C/M -- see Vol~I Corollary~\cite{VolI}). This is $A$, not $A^!$.
 \item \emph{Derived center}: $Z^{\mathrm{der}}_{\mathrm{ch}}(A)
 = C^\bullet_{\mathrm{ch}}(A, A) = \RHom(\Omega B(A), A)$ is the
 closed-sector algebra at level $3$ (Vol~I Theorem~H; Vol~III
 Theorem~\ref{thm:bulk-three-faces-coincidence}). A physical bulk
 identification requires the relevant open--closed comparison datum.
 This is the ``$E_2$ uplift,'' not $A^!$.
 \item \emph{Drinfeld center}: $\cZ(\Rep^{\Eone}(A))$ is a braided monoidal category (Chapter~\ref{ch:drinfeld-center}). This lives on representation categories, not on algebras.
\end{enumerate}
The relationships: $A^!$ is the algebra controlling the defect
(boundary); $Z^{\mathrm{der}}_{\mathrm{ch}}(A)$ is the algebraic
closed-sector object, which becomes a physical bulk only after the
open--closed comparison is supplied; the Drinfeld center is the
categorical passage from $\Eone$ (boundary) to $\Etwo$ (bulk)
($\cZ(\Rep^{\Eone}(A))$ is a braided monoidal \emph{category};
$Z^{\mathrm{der}}_{\mathrm{ch}}(A)$ is a chiral \emph{algebra}; the
Drinfeld center is the categorification of the derived center under
the BZFN hypotheses). The bulk acts on the defect via
$\mu_{\mathrm{bb}}$ (Construction~\ref{constr:universal-defect}).
Three constructions reach the bulk from distinct starting points: the
Hochschild centre $Z^{\mathrm{der}}_{\mathrm{ch}}(A)$ (this
proposition); the positive-half Hall presentation $Y^{+}(X)$
(Definition~\ref{def:bulk-yplus}); the Drinfeld double
$G(X)=\mathcal D_\hbar(Y^{+}(X))$ (Definition~\ref{def:bulk-gx}).
The three meet at one level-$3$ object on the toric stratum in
ordinary chain ambient
(Theorem~\ref{thm:bulk-three-faces-coincidence}).

The distinctness of $A^!$ from $Z^{\mathrm{der}}_{\mathrm{ch}}(A)$ is a fundamental point: for the Kac--Moody family $A = V_k(\frakg)$, the Koszul dual is the completed Verdier-dual bar object $D_{\Ran}B^{\mathrm{ord}}(V_k(\frakg))$, whose reflected current presentation is $V_{k'}(\frakg)$ at $k' = -k - 2h^\vee$, while $Z^{\mathrm{der}}_{\mathrm{ch}}(A) = C^\bullet_{\mathrm{ch}}(V_k(\frakg), V_k(\frakg))$ is the chiral Hochschild cochain complex. These are manifestly different objects: $A^!$ controls defects through the CE/bar-cochain target; $Z^{\mathrm{der}}_{\mathrm{ch}}(A)$ is a derived endomorphism complex.
\end{proposition}

\subsection{The \texorpdfstring{$\En$}{E-n} level of the Koszul dual}
\label{subsec:en-koszul-level}

The $\En$ level of $A^!$ is determined by the ambient theory, not by $A$ alone.

\begin{conjecture}[\texorpdfstring{$\En$}{E-n}-structure on the chiral Koszul dual from holomorphic CS]
\label{conj:en-koszul-from-hcs}
\ClaimStatusConjectured
Let $A_n$ denote the boundary chiral algebra of holomorphic Chern--Simons theory (or its 6d algebraic surrogate) on $\C^n$ projected to a curve $C$, so that the full theory on $\C^n$ carries $E_n$-factorization structure. Then:
\begin{enumerate}[label=\textup{(\roman*)}]
 \item The chiral Koszul dual $A_n^!$ inherits $E_n$-chiral factorization structure from the ambient $E_n$ theory on $\C^n$.
 \item At $n = 1$ (3d CS): $A_1^! = D_{\Ran}B^{\mathrm{ord}}(V_k(\frakg))$, with reflected current presentation $V_{k'}(\frakg)$ at $k' = -k - 2h^\vee$, is an $\Eone$-chiral algebra. The $\Eone$ structure is the standard one.
 \item At $n = 2$ (5d CS): $A_2^!$ carries $\Etwo$-chiral structure on $\C^2$. On a curve, $A_2^!|_C$ is the Koszul dual of the affine Yangian with inverted parameters.
 \item At $n = 3$ (6d theory): $A_3^!$ carries $E_3$-chiral structure on $\C^3$. This is the \emph{large chiral algebra} of the Costello construction: it contains the quantum toroidal algebra (on $C$), the affine Yangian ($\Etwo$ on $\C^2$), and the $E_3$ master structure on $\C^3$.
\end{enumerate}
By item: (i) is the general principle (conjectural at $n \geq 2$). (ii) is a theorem in the completed ordered-bar ambient: the Feigin--Frenkel reflection gives the current presentation of $D_{\Ran}B^{\mathrm{ord}}(V_k(\frakg))$. (iii) is partially established: the affine Yangian Koszul dual is constructed by Schiffmann--Vasserot; the $\Etwo$-chiral structure on the Koszul dual is conjectural. (iv) is fully conjectural and requires the non-Lagrangian 6d algebraic framework of Remark~\ref{rem:6d-not-lagrangian}.
\end{conjecture}

\begin{remark}[\texorpdfstring{$E_3$}{E-3} chiral is not \texorpdfstring{$E_3$}{E-3} symmetric]
\label{rem:e3-not-symmetric}
The $E_3$ structure in Conjecture~\ref{conj:en-koszul-from-hcs}(iv) is \emph{not} the trivially braided $E_3$ of Chapter~\ref{ch:en-factorization} (where $\pi_1(\Conf_2(\R^6)) = 0$ kills the braiding). The holomorphic refinement and the Omega-background deformation produce a nontrivial $E_3$-chiral structure: the factorization algebra on $\Conf_n(\C^3)$ acquires a nontrivially braided structure after deformation by the equivariant parameters $(h_1, h_2, h_3)$. Note that $\Conf_2(\C^3) \simeq \Conf_2(\R^6) \simeq S^5$ as topological spaces, so the braiding is trivial at the topological level; the nontriviality is a feature of the \emph{deformed holomorphic} factorization structure, not of the underlying topology. Without the Omega-background, the braiding is symmetric (the $E_3$-chiral structure has trivial braiding, not that $E_3$ ``becomes'' $E_\infty$: the operadic level remains $E_3$, but the braiding data is trivial).
\end{remark}

\subsection{Deformation of chiral CE cochains to \texorpdfstring{$E_3$}{E-3}}
\label{subsec:ce-deformation-e3}

The chiral CE cochains of an $E_\infty$-chiral algebra $A$ (such as a BD commutative chiral algebra) carry a natural $E_\infty$-algebra structure from the higher Deligne conjecture (Remark~\ref{rem:deligne-hochschild}); in particular, they carry $E_3$-algebra structure. The \emph{deformed} $E_3$-chiral structure relevant to this chapter arises not from the Deligne conjecture alone, but from a deformation of the CE differential by the Omega-background parameters of the holomorphic CS theory.

\begin{construction}[Quantum deformation of chiral CE cochains]
\label{constr:quantum-ce-deformation}
Let $A = V_k(\frakg)$ be a Kac--Moody vertex algebra (an $E_\infty$-chiral algebra). The chiral CE cochains $C^\bullet_{\mathrm{ch}}(A, A) = Z^{\mathrm{der}}_{\mathrm{ch}}(A)$ carry:
\begin{enumerate}[label=\textup{(\roman*)}]
 \item An $\Etwo$-algebra structure from the Deligne conjecture: the cup product and brace operations realize the Dunn additivity decomposition $\Etwo \simeq \Eone \otimes \Eone$ (operadic tensor product, not the coproduct over $E_0$).
 \item A classical Poisson bracket $\{-, -\}_{\mathrm{cl}}$ of degree $-1$ from the Gerstenhaber structure on $\HH^\bullet(A, A)$.
 \item A \emph{quantum deformation} parametrized by the Omega-background $\boldsymbol{h} = (h_1, h_2, h_3)$ with $h_1 + h_2 + h_3 = 0$ (so only two of $(h_1, h_2, h_3)$ are independent). At $\boldsymbol{h} = 0$, the structure is the undeformed $\Etwo$-structure. At generic $\boldsymbol{h}$, the deformed differential
 \[
 d_{\boldsymbol{h}} = d_{\mathrm{CE}} + h_1 \cdot \partial_1 + h_2 \cdot \partial_2 + h_3 \cdot \partial_3
 \]
 (where $\partial_i$ are the equivariant differential operators on $\C^3$ associated to the three coordinate $\C^*$-actions, and the constraint $h_1 + h_2 + h_3 = 0$ ensures $d_{\boldsymbol{h}}^2 = 0$ via the Calabi--Yau condition $\det = 1$ on the torus action) promotes the structure to $E_3$-chiral. The promotion uses the topological $E_3$-operad action on $\Conf_n(\C^3)$, transferred to the algebraic setting via Kontsevich formality; this is the topological $E_3$ of Remark~\ref{rem:two-e3-ap154}(b), not the algebraic $E_3$ of~(a).
\end{enumerate}
\end{construction}

\begin{remark}[Two \texorpdfstring{$E_3$}{E-3} structures: algebraic vs topological]
\label{rem:two-e3-ap154}
The $E_3$ structure in Construction~\ref{constr:quantum-ce-deformation} arises from two independent sources:
\begin{enumerate}[label=\textup{(\alph*)}]
 \item \emph{Algebraic $E_3$}: the higher Deligne conjecture (now a theorem: Lurie, 2012; Costello, 2007) states that for an $E_n$ algebra $A$, the Hochschild cochains $\HH^\bullet(A, A)$ carry an $E_{n+1}$ structure. For $\Eone$ input: $\Etwo$ on $\HH^\bullet$ (standard Deligne conjecture). For $E_\infty$ input: $E_\infty$ on $\HH^\bullet$ (since $\infty + 1 = \infty$). In particular, for $A = V_k(\frakg)$ ($E_\infty$-chiral), $\HH^\bullet(A,A)$ carries \emph{at least} $E_3$; the full $E_\infty$ structure is available but only the $E_3$ sub-structure is relevant to the 6d holomorphic theory on $\C^3$.
 \item \emph{Topological $E_3$}: the $E_3$-operad acts on $\Conf_n(\C^3)$ via the framed little $3$-disks.
\end{enumerate}
These two $E_3$ structures agree under Kontsevich formality over $\Q$ and differ at the chain level. The holomorphic CS construction uses the topological $E_3$ from (b), deformed by the Omega-background; the algebraic $E_3$ from (a) is the classical limit. Integral agreement remains open.
\end{remark}

\subsection{Chiral quantum groups from defect Wilson lines}
\label{subsec:chiral-qg-wilson}

The chiral quantum group is the Hopf-algebra-like structure on the defect algebra $A^!$ extracted from the holomorphic Wilson lines of the ambient CS theory.

\begin{conjecture}[Chiral quantum group on \texorpdfstring{$\C^3$}{C-3}]
\label{conj:chiral-qg-c3}
\ClaimStatusConjectured
For the 6d holomorphic theory on $\C^3$ with $\fgl_1$ gauge algebra and Omega-background $(h_1, h_2, h_3)$:
\begin{enumerate}[label=\textup{(\roman*)}]
 \item The defect algebra $A^!_{\C^3}$ carries a coproduct
 \[
 \Delta_{\mathrm{ch}} \colon A^!_{\C^3} \;\longrightarrow\; A^!_{\C^3} \hat{\otimes} A^!_{\C^3}
 \]
 induced by the collision of two parallel Wilson lines in $\C^3$. The coproduct is coassociative and compatible with the $E_3$-chiral structure.
 \item The $R$-matrix is the universal braiding of the $E_3$ theory:
 \[
 R_{\mathrm{ch}}(u, v) \in A^!_{\C^3} \hat{\otimes} A^!_{\C^3}((u))((v))
 \]
 with $(u, v)$ the two spectral parameters from the two extra directions of $\C^3$ beyond the Wilson line. This $R$-matrix satisfies the parametric Yang--Baxter equation in \emph{two} spectral variables.
 \item The representation category $\Rep^{\Etwo}(A^!_{\C^3}|_{\C^2})$ is the braided monoidal category of quantum toroidal representations. The Drinfeld center $\cZ(\Rep^{\Eone}(A^!_{\C^3}|_C))$ should recover the $\Etwo$-braided category; at 5d (one dimension lower), the analogous statement is Theorem~\ref{thm:drinfeld-center-coha}, which establishes $\cZ(\Rep^{\Eone}(Y^+(\widehat{\fgl}_1))) \simeq \Rep^{\Etwo}(Y(\widehat{\fgl}_1))$. The 6d lift replaces $Y^+$ with $A^!_{\C^3}|_C$ and the full Yangian with the quantum toroidal algebra.
 \item The coproduct deviates from the primitive (Hopf) coproduct by the Dimofte $K$-matrix (Remark~\ref{rem:dimofte-k-matrix-cy}): $\Delta_{\mathrm{ch}}(x) = x \otimes 1 + K_{\cA}(z) \otimes x + \ldots$.
\end{enumerate}
Each item in this conjecture is verified independently at lower dimension: item~(i) at 5d (Schiffmann--Vasserot coproduct on CoHA), item~(ii) at 5d (Yangian $R$-matrix, Theorem~\ref{thm:yang-r-matrix-ybe}), item~(iii) at 5d (Theorem~\ref{thm:drinfeld-center-coha}). Each individual arrow is established at 5d, but the 6d composite is conjectural.
\end{conjecture}

\begin{conjecture}[Chiral quantum group on \texorpdfstring{$K3 \times E$}{K3 x E}]
\label{conj:chiral-qg-k3xe}
\ClaimStatusConjectured
For the 6d holomorphic theory on $K3 \times E$:
\begin{enumerate}[label=\textup{(\roman*)}]
 \item The defect algebra $A^!_{K3 \times E}$ on $E$ is a deformation of the Koszul dual of the BKM vertex algebra $V_{\mathfrak{g}_{\Delta_5}}$, with deformation controlled by the K3 complex structure moduli.
 \item The Wilson lines wrapping $E$ are indexed by Mukai vectors $v \in H^*(K3, \Z)$; the fusion of Wilson lines is controlled by the Mukai pairing.
 \item The coproduct on $A^!_{K3 \times E}$ is the chiral coproduct of Vol~I, specialized to the $K3 \times E$ root datum.
 \item The $\Etwo$-braided representation category of the chiral quantum group is the Verlinde category of the BKM superalgebra $\mathfrak{g}_{\Delta_5}$. This is consistent with Conjecture~\ref{conj:mtc-k3} (Chapter~\ref{ch:e2-chiral}), where $\Phi_{E_2}$ on the \(K3\times E\) input means the chiral Drinfeld centre \(Z^{\mathrm{ch}}(\Phi_{E_1})\), not a native \(E_2\)-output of the \(d=3\) specialisation.
\end{enumerate}
The assertion requires the 6d algebraic framework for $K3 \times E$ (non-Lagrangian; Remark~\ref{rem:6d-not-lagrangian}), the framed object-level \(\Phi_3\) construction of Theorem~\ref{thm:cy-to-chiral-d3}, and the identification of the BKM root datum with the K3 Mukai lattice. These are independent comparison data. The $\kappa_\bullet$-spectrum consistency check is $\kappa_{\mathrm{ch}}^{\mathrm{Heis}} = 3$, $\kappa_{\mathrm{BKM}}(\Delta_5) = 5$, $\kappa_{\mathrm{cat}} = \chi(\cO_{K3 \times E}) = 0$ (K\"unneth: $2 \cdot 0 = 0$; the K3-fiber value $\chi(\cO_{K3}) = 2$ is a separate invariant), and $\kappa_{\mathrm{fiber}} = 24$, with compact total-space $\kappa_{\mathrm{ch}}^{\mathrm{Hodge}}(A_{K3\times E})=0$.
\end{conjecture}

\begin{remark}[Summary: the holomorphic CS\texorpdfstring{$\to$}{->}chiral quantum group construction]
\label{rem:hcs-construction}
The full construction from holomorphic Chern--Simons to chiral quantum group:
\[
\begin{tikzcd}[column sep=1.8cm]
 \text{hol CS on } \C^n \ar[r, "\text{boundary}"] &
 A_n \text{ on } C \ar[r, "B(-)"] &
 B(A_n) \ar[r, "D_{\Ran}"] &
 A_n^! \ar[r, "\Rep^{E_2}"] &
 \text{chiral QG}
\end{tikzcd}
\]
At $n = 1$ (3d CS), the chain
\[
 V_k(\frakg) \to B(V_k) \to
 D_{\Ran}B^{\mathrm{ord}}(V_k(\frakg))
 \to \text{conformal blocks}
\]
is the Feigin--Frenkel reflected-current line \(k'=-k-2h^\vee\), together with
Vol~I Theorem~A for the Kac--Moody family.

At $n = 2$ (5d CS), the chain
\[
 Y(\widehat{\fgl}_1) \to B(Y) \to Y^!
 \to \Etwo\text{-braided representations}
\]
uses the constructed boundary algebra of Theorem~\ref{thm:5d-boundary-yangian},
the CoHA coproduct and Drinfeld centre of
Theorem~\ref{thm:drinfeld-center-coha}, and the MO/\(R\)-matrix
comparison of Theorem~\ref{thm:mo-chiral-rmatrix}.  The remaining
obligation is to identify \(B(Y)\) as a factorisation coalgebra on
\(\Ran(C)\) with the Schiffmann--Vasserot shuffle structure.

At $n = 3$ (6d theory), the expected chain
\[
 U_{q,t}(\widehat{\widehat{\fgl}}_1) \to B(-)
 \to E_3\text{-Koszul dual} \to \text{chiral quantum group}
\]
requires the non-Lagrangian 6d framework and the finite
DWR/Ran Rees comparison before any centre/double operation.

For \(K3\times E\), the expected BKM-related chain
\[
 \text{BKM-related algebra} \to B(-) \to
 \text{defect algebra} \to \text{BKM quantum group}
\]
uses the framed $\Phi_3$ object, the 6d algebraic framework, and the
centre/double passage separated in Proposition~\ref{prop:qca-finite-rees-layer-separation}.
\end{remark}

\subsection{Derived Maurer--Cartan space of chiral CE cochains}
\label{subsec:derived-mc-chiral-ce}

The chiral CE cochains $C^\bullet_{\mathrm{ch}}(L) = \mathrm{CE}^*(L)$ of a Lie conformal algebra $L$ carry a dg Lie algebra structure (the Gerstenhaber bracket). The finite formulas in this subsection use either a finite Rees/DWR truncation or the exterior CE model on a chosen finite generator space \(L_{\mathrm{gen}}\). The full chiral CE complex, with its \(\partial\)-descendant tower and completions, is a filtered or pro-object and is not measured by the binomial dimensions below without this truncation datum. The Maurer--Cartan equation
\begin{equation}\label{eq:mc-chiral-ce}
 d_{\mathrm{CE}}(\alpha) + \tfrac{1}{2}[\alpha, \alpha] = 0, \qquad \alpha \in \mathrm{CE}^1(L),
\end{equation}
parametrises \emph{flat chiral deformations} of the $\lambda$-bracket. The derived Maurer--Cartan space $\mathrm{MC}^{\mathrm{der}}(\mathrm{CE}^{*}(L))$ is the derived stack whose tangent complex at the trivial solution $\alpha = 0$ is the shifted cochain complex $T_0\mathrm{MC} = \mathrm{CE}^{*}(L)[1]$, with $H^0$ controlling infinitesimal automorphisms, $H^1$ controlling deformations, and $H^2$ controlling obstructions.

\begin{proposition}[MC space for the five families]
\label{prop:mc-five-families}
The derived Maurer--Cartan spaces of the chiral CE cochains for the five standard families are:
\begin{enumerate}[label=\textup{(\roman*)}]
 \item \emph{Heisenberg $H_k$ (class~G):} $\mathrm{MC} = \{\mathrm{pt}\}$ (abelian: $d = 0$ and $[\alpha, \alpha] = 0$). The tangent complex has $H^0 = 0$, $H^1 = 1$ (the level shift $k \mapsto k + \varepsilon$), $H^2 = 0$ (unobstructed). The derived structure is trivial: $\mathrm{MC} = \mathrm{MC}^{\mathrm{der}}$.
 \item \emph{Kac--Moody $V_k(\mathfrak{sl}_2)$ (class~L):} $\mathrm{MC} = \mathrm{Conn}_{\mathrm{flat}}(C, \mathfrak{sl}_2)$, the space of flat $\mathfrak{sl}_2$-connections on the curve~$C$. At the abstract level: $H^0 = H^1 = 0$ (Whitehead's theorem for the semisimple Lie algebra $\mathfrak{sl}_2$), $H^2 = 1$ (the top class $H^3(\mathfrak{sl}_2) = k$). The algebra $\mathfrak{sl}_2$ is \emph{rigid}: no infinitesimal deformations.
 \item \emph{Virasoro $\mathrm{Vir}_c$ (class~M):} $\mathrm{MC} = \mathrm{Proj}(C)$, the space of projective structures on~$C$, with $\dim \mathrm{Proj}(C) = 3g - 3$ at genus~$g$. The $L_\infty$ corrections from the shadow tower ($l_3, l_4, \ldots$) make $\mathrm{MC}^{\mathrm{der}}$ genuinely derived: the higher MC equation $d(\alpha) + \tfrac{1}{2}l_2(\alpha, \alpha) + \tfrac{1}{6}l_3(\alpha, \alpha, \alpha) + \cdots = 0$ has infinitely many terms.
 \item \emph{Mukai Heisenberg $H_{\mathrm{Muk}}$ (class~G, rank~$24$):} $\mathrm{MC} = \{\mathrm{pt}\}$ (abelian). The tangent complex has $H^0 = 24$ (the Mukai lattice automorphisms), $H^1 = \binom{24}{2} = 276$ (antisymmetric bilinear forms on the Mukai lattice), $H^2 = \binom{24}{3} = 2024$ (formal obstructions, all of which vanish by abelianity: $[\alpha,\alpha] = 0$). The derived structure is trivial.
 \item \emph{Yangian $Y(\widehat{\mathfrak{gl}}_1)$ (class~L, truncated to $3$ generators):} The MC equation has solutions parametrised by deformations of the Yangian commutation relations. The tangent complex has Nomizu cohomology $H^0 = 2$, $H^1 = 2$, $H^2 = 1$ (for the $3$-dimensional Heisenberg Lie algebra $\mathfrak{h}_3$).
\end{enumerate}
\end{proposition}

\begin{proof}
Items~(i) and~(iv): for an abelian Lie conformal algebra $L$ (all $0$-th products vanish), $d_{\mathrm{CE}} = 0$ and $[\alpha, \alpha] = 0$ for all $\alpha \in \mathrm{CE}^1(L)$. The MC equation is trivially satisfied at $\alpha = 0$ with no obstructions. The shifted CE cohomology is $H^k(\mathrm{CE}^*(L)[1]) = \binom{n}{k+1}$ where $n = \dim(L)$, since the differential vanishes.

Item~(ii): the Lie algebra $\mathfrak{sl}_2$ is semisimple; Whitehead's theorem gives $H^1(\mathfrak{sl}_2) = H^2(\mathfrak{sl}_2) = 0$ (rigidity). The top cohomology $H^3(\mathfrak{sl}_2) = k$ by Poincar\'e duality for the $3$-dimensional Lie algebra.

Item~(iii): the Virasoro Lie conformal algebra has a single chosen
generator $T$ of spin~$2$. In the finite exterior generator quotient,
$\mathrm{CE}^k_{\mathrm{gen}} = 0$ for $k \geq 2$, so
$H^1(\text{tangent}) = 0$ at that quotient level. The class~M
$L_\infty$ corrections $l_3(T,T,T) = -2c \cdot T$ and
$l_4(T,T,T,T) = (40/27) \cdot T$ \textup{(}at $c = 1$\textup{)}
contribute higher-order terms to the MC equation that live on the
ordered bar complex and in the completed chiral CE object, not in the
finite exterior quotient.

Item~(v): the Lie algebra underlying the truncated Yangian is the $3$-dimensional Heisenberg algebra $\mathfrak{h}_3 = \langle e_0, e_1, e_2 \mid [e_0, e_1] = e_2 \rangle$, $2$-step nilpotent. By Nomizu's theorem (1954), $H^k(\mathfrak{h}_3) = (1, 2, 2, 1)$ for $k = 0, 1, 2, 3$.
\end{proof}

\begin{proposition}[Finite exterior tangent Euler characteristic]
\label{prop:mc-tangent-euler-char}
\ClaimStatusProvedHere
Let \(L_{\mathrm{gen}}\) be a finite generator space of rank
\(n\geq1\), and let
\[
  T^{\mathrm{fin}}_0\mathrm{MC}
  :=\mathrm{Hom}\bigl(\Lambda^\bullet L_{\mathrm{gen}},\Bbbk\bigr)[1]
\]
be the finite exterior model for the tangent complex. Then
\[
 \chi(T^{\mathrm{fin}}_0\mathrm{MC})
 = \sum_{k=-1}^{n-1} (-1)^k \dim T^k
 = \sum_{k=-1}^{n-1} (-1)^k \binom{n}{k+1} = 0.
\]
The statement concerns the finite exterior quotient or a finite
Rees/DWR truncation. It does not compute the Euler characteristic of
the uncompleted full chiral CE complex when \(\partial\)-descendants or
infinite shadow towers are present.
\end{proposition}

\begin{proof}
Substituting $j = k + 1$ gives
\[
  \chi(T^{\mathrm{fin}}_0\mathrm{MC})
  = -\sum_{j=0}^{n} (-1)^j \binom{n}{j}
  = -(1 - 1)^n
  = 0
\]
for $n \geq 1$ by the binomial theorem. The proof uses only
\(\dim \Lambda^j L_{\mathrm{gen}}=\binom n j\); it therefore proves the
finite exterior statement and no stronger completed CE assertion.
\end{proof}

\begin{proposition}[ADE deformations of the Mukai Heisenberg]
\label{prop:ade-deformations-mukai}
The abelian Mukai Heisenberg $H_{\mathrm{Muk}}$ (rank~$24$, Mukai pairing of signature~$(4, 20)$) admits ADE deformations to non-abelian current algebras. For each ADE root system $R$ of rank~$r \leq 20$ embeddable in the negative-definite part of the Mukai lattice:
\begin{enumerate}[label=\textup{(\roman*)}]
 \item The deformed algebra has $\dim(\mathfrak{g}_R)$ non-abelian generators (where $\mathfrak{g}_R$ is the ADE Lie algebra: $\dim(\mathfrak{sl}_{n+1}) = n(n+2)$, $\dim(\mathfrak{so}_{2n}) = n(2n-1)$, $\dim(\mathfrak{e}_6) = 78$, $\dim(\mathfrak{e}_7) = 133$, $\dim(\mathfrak{e}_8) = 248$) and $24 - r$ abelian generators.
 \item The deformation space has dimension~$r$ (one parameter per simple root of~$R$).
 \item At level~$k$, the chiral modular characteristic of the non-abelian part is $\kappa_{\mathrm{ch}}(V_k(\mathfrak{g}_R)) = \dim(\mathfrak{g}_R) \cdot (k + h^\vee) / (2h^\vee)$.
\end{enumerate}
The ADE deformations correspond geometrically to ADE singularities of the K3 surface: a K3 with an $R$-singularity has enhanced gauge symmetry~$\mathfrak{g}_R$ from the collapsed rational curves.
\end{proposition}

\begin{proof}
The Mukai lattice $\Lambda = E_8(-1)^{\oplus 2} \oplus U^{\oplus 4}$ has signature $(4, 20)$. The negative-definite part has rank~$20$, and ADE root lattices of rank~$r \leq 20$ embed by Nikulin's embedding theorem (1979). The Lie algebra dimensions are standard (Humphreys, \emph{Introduction to Lie Algebras}, 1972). The $\kappa_{\mathrm{ch}}$ formula is the Kac--Moody modular characteristic from the landscape census. The geometric interpretation of ADE singularities as collapsed $(-2)$-curves follows from the McKay correspondence (Du Val singularities).
\end{proof}

\begin{remark}[Shadow class determines derived MC structure]
\label{rem:shadow-class-mc}
The shadow class (G/L/C/M) of the chiral algebra $A = U^{\mathrm{ch}}(L)$ determines the analytic type of $\mathrm{MC}^{\mathrm{der}}$:
\begin{itemize}
 \item Class~G (Gaussian): $\mathrm{MC}^{\mathrm{der}}$ is classical (no higher homotopies). Example: Heisenberg, Mukai Heisenberg.
 \item Class~L (rational): $\mathrm{MC}^{\mathrm{der}}$ is a derived scheme with finitely many $L_\infty$ corrections. Example: Kac--Moody.
 \item Class~C (convergent): $\mathrm{MC}^{\mathrm{der}}$ has convergent $L_\infty$ series. Example: bc system.
 \item Class~M (mock modular): $\mathrm{MC}^{\mathrm{der}}$ has Gevrey-$1$ divergent $L_\infty$ corrections, requiring Borel summation. The shadow tower $S_k$ contributes to the $k$-th homotopy of $\mathrm{MC}^{\mathrm{der}}$. Example: Virasoro, $W_{1+\infty}$.
\end{itemize}
This is the derived-deformation-theoretic avatar of the shadow tower: the $L_\infty$ structure maps $l_k$ of the bar complex become the higher MC equation terms, and the shadow invariants $S_k$ (Vol~I) are the coefficients of the $k$-th correction $\delta^{(k)}$ (see Vol~I~\cite{VolI}).
\end{remark}

\section{The Costello 5d construction}
\label{sec:costello-5d-construction}

The $d = 5$ regime of holomorphic Chern--Simons theory
(Theorem~\ref{thm:5d-boundary-yangian}) is the middle step of the
\(3\to5\to6\) hierarchy, and the only one where the full construction is
proved: the boundary chiral algebra is the affine Yangian
\(Y(\widehat{\fgl}_1)\). The five ingredients below are the algebraic
faces of Costello's construction (arXiv:1303.2632).

\subsection{The 5d theory and its boundary}
\label{subsec:5d-theory}

\begin{proposition}[5d holomorphic CS specification]
\label{prop:5d-hcs-specification}
The data of 5d holomorphic Chern--Simons theory on $\C^2_{h_1,h_2} \times \R$ with gauge algebra $\fgl_1$ is:
\begin{enumerate}[label=\textup{(\alph*)}]
 \item \emph{Action}: $S = \frac{1}{2}\int \Omega \wedge A \wedge \bar{\partial}A$, where $\Omega = dz_1 \wedge dz_2$ is the holomorphic $2$-form on $\C^2$. The cubic term vanishes for abelian $\fgl_1$.
 \item \emph{Omega-background}: the $T^2$-action $(z_1, z_2) \mapsto (e^{ih_1 t}z_1, e^{ih_2 t}z_2)$ with equivariant parameters $(h_1, h_2)$.
 \item \emph{CY condition}: the one-loop anomaly cancellation of the theory on $\C^2 \times \R$ requires $h_1 + h_2 + h_3 = 0$, where $h_3 = -(h_1 + h_2)$ is the equivariant weight of the $\R$-direction. This is the Calabi--Yau condition on $\C^3 = \C^2 \times \C_R$.
 \item \emph{Boundary algebra}: the observables on the boundary $\{R = 0\}$, projected to a holomorphic curve $C \subset \C^2$, form the affine Yangian $Y(\widehat{\fgl}_1)$ with structure function $g(u) = \prod_{i=1}^{3}(u - h_i)/(u + h_i)$.
\end{enumerate}
\end{proposition}

\subsection{Tree-level OPE and Yangian relations}
\label{subsec:tree-level-ope}

\begin{proposition}[Tree-level structure function]
\label{prop:tree-level-g}
The structure function $g(u) = \prod_{i=1}^{3}(u-h_i)/(u+h_i)$, arising from tree-level Feynman diagrams of the 5d theory, satisfies:
\begin{enumerate}[label=\textup{(\roman*)}]
 \item \emph{Reflection identity}: $g(u)\,g(-u)=1$. Indeed
 \[
 g(-u)=\prod_{i=1}^{3}\frac{-u-h_i}{-u+h_i}
      =\prod_{i=1}^{3}\frac{u+h_i}{u-h_i}=g(u)^{-1}.
 \]
 This ensures the antisymmetry of the $ee$ relation $g(\Delta)\, e(z)\, e(w) = -g(-\Delta)\, e(w)\, e(z)$.
 \item \emph{Coefficient vanishing}: expand $g(z) = \sum_{j \geq 0} \phi_j z^{-j}$. Then:
 \[
 \phi_0 = 1,\quad \phi_1 = 0,\quad \phi_2 = 0,\quad \phi_3 = -2\sigma_3,\quad \phi_4 = 0.
 \]
 The vanishing $\phi_1 = 0$ is equivalent to $h_1 + h_2 + h_3 = 0$ (the CY condition). The vanishing of all even $\phi_{2k}$ for $2k < 6$ follows from parity: $\log g(z) = \sum_{k \text{ odd}} \alpha_k z^{-k}$, so $\phi_{2k} = 0$ unless $2k$ admits a partition into odd parts $\geq 3$. The first nonzero even coefficient is $\phi_6 = \alpha_3^2/2$, from the partition $6 = 3 + 3$.
 \item \emph{Classical $r$-matrix}: the $R$-matrix on $\mathrm{Fock}_2$ has leading behaviour $R(u) = 1 - \sigma_2/u + O(1/u^2)$ as $u \to \infty$. The classical $r$-matrix is $r(u) = -\sigma_2/u$, with residue $-\sigma_2$ equal to the Kac--Moody level $\Psi = -\sigma_2$.
\end{enumerate}
\end{proposition}

\begin{proof}
Part~(i): the reflection identity $g(u) g(-u) = 1$ is an algebraic consequence of the rational function form. Explicitly, $g(-u) = \prod (-u-h_i)/(-u+h_i) = \prod (u+h_i)/(u-h_i) = 1/g(u)$, with each factor applying the identity $(-u-h)/(h-u) = (u+h)/(u-h)$. This holds for \emph{any} $(h_1,h_2,h_3)$, including the CY locus.

Part~(ii): $\log g(z) = \sum_a [\log(1-h_a/z) - \log(1+h_a/z)] = -2\sum_{k \text{ odd}} p_k/(k z^k)$, where $p_k = h_1^k + h_2^k + h_3^k$. By the CY condition, $p_1 = 0$, giving $\phi_1 = \alpha_1 = 0$. Even $\alpha_k$ vanish identically from the $\log(1-x) - \log(1+x)$ expansion. The leading nontrivial coefficient is $\phi_3 = \alpha_3 = -2p_3/3 = -2\sigma_3$ (using Newton's identity $p_3 = 3\sigma_3$ under the CY condition $\sigma_1 = 0$). The $\phi_j$ for $j \geq 6$ even are generically nonzero, arising from products of odd $\alpha_k$'s.

Part~(iii) is the classical limit computed in
Theorem~\ref{thm:yang-r-matrix-ybe}.  At $(h_1,h_2,h_3)=(1,-2,1)$ it
gives $-\sigma_2=3$; at $(1,-3,2)$ it gives $-\sigma_2=7$.
\end{proof}

\subsection{Tree-level OPE at spins \texorpdfstring{$1$}{1} and \texorpdfstring{$2$}{2}}
\label{subsec:tree-level-spin12}

\begin{proposition}[Spin-\texorpdfstring{$1$}{1} and spin-\texorpdfstring{$2$}{2} OPE from the 5d theory]
\label{prop:5d-spin12-ope}
The boundary algebra of 5d holomorphic CS with $\fgl_1$ gauge algebra has:
\begin{enumerate}[label=\textup{(\roman*)}]
 \item \emph{Spin-$1$ OPE}: $J(z) J(w) = \Psi/(z-w)^2$, where $\Psi = -\sigma_2 = -(h_1 h_2 + h_1 h_3 + h_2 h_3)$ is the Kac--Moody level of the 5d theory. In lambda-bracket form: $\{J_\lambda J\} = \Psi \cdot \lambda$. This matches Proposition~\ref{prop:codim2-defect-ope}(a).
 \item \emph{Spin-$2$ OPE}: the Sugawara construction $T = \mathopen{:}JJ\mathclose{:}/(2\Psi)$ gives a Virasoro with central charge $c = 1$ (for $\fgl_1$, $c = \dim(\fgl_1) \cdot \Psi/(\Psi + h^\vee) = 1$, independent of $\Psi$). The OPE is $T(z)T(w) = \frac{1/2}{(z-w)^4} + \frac{2T(w)}{(z-w)^2} + \frac{\partial T(w)}{z-w}$. In lambda-bracket form: $\{T_\lambda T\} = \frac{1}{12}\lambda^3 + 2T\lambda + \partial T$. This matches Proposition~\ref{prop:codim2-defect-ope}(b).
\end{enumerate}
\end{proposition}

\begin{proof}
This is the Sugawara construction, as in
Proposition~\ref{prop:codim2-defect-ope}.
\end{proof}

\subsection{One-loop correction: the CY constraint}
\label{subsec:one-loop-cy}

\begin{proposition}[One-loop enforcement of the CY condition]
\label{prop:one-loop-cy-constraint}
The one-loop Feynman diagram of 5d holomorphic CS on $\C^2 \times \R$ produces a correction to the boundary OPE that is proportional to $h_1 + h_2 + h_3$. Specifically:
\begin{enumerate}[label=\textup{(\roman*)}]
 \item The structure function coefficient $\phi_1 = -2(h_1 + h_2 + h_3)$. The CY condition $h_1 + h_2 + h_3 = 0$ is equivalent to $\phi_1 = 0$.
 \item When $\phi_1 \neq 0$, the Serre relation of the would-be Yangian acquires a correction of order $\phi_1 \cdot \phi_3$, which breaks the Yangian structure.
 \item The anomaly cancellation mechanism: the total equivariant weight of the holomorphic $3$-form on $\C^3 = \C^2 \times \C_R$ is $\mathrm{wt}(\Omega_3) = h_1 + h_2 + h_3$. Setting this to zero ensures $\Omega_3$ is $T^3$-invariant, which is the CY condition $\det(T^3|_{\Omega^{3,0}}) = 1$.
 \item The reflection identity $g(u) g(-u) = 1$ holds for ALL $(h_1,h_2,h_3)$ (it is an algebraic property of the rational function, not a consequence of the CY condition). The CY condition enters through $\phi_1$, which controls the Serre relations.
\end{enumerate}
\end{proposition}

\begin{proof}
Part~(i): \(\phi_1=\alpha_1=-2p_1=-2(h_1+h_2+h_3)\) from the
log-series expansion. Part~(ii): the Serre correction
\(\delta_{\mathrm{Serre}}=\phi_1\phi_3/3\) vanishes at CY
(\(\phi_1=0\)) and is nonzero away from it. Parts~(iii)--(iv) are
Costello 2013, Proposition~3.2, together with the displayed algebraic
consequences.
\end{proof}

\subsection{The Omega-background deformation}
\label{subsec:omega-deformation}

\begin{proposition}[Deformation hierarchy from the Omega-background]
\label{prop:omega-deformation}
The Omega-background $(h_1, h_2)$ on $\C^2$ (with $h_3 = -(h_1+h_2)$) deforms the boundary algebra through a single effective parameter $\sigma_3 = h_1 h_2 h_3$:
\begin{enumerate}[label=\textup{(\roman*)}]
 \item At $\sigma_3 = 0$ (the self-dual point, where one $h_i = 0$): $g(u) = 1$ identically, the $R$-matrix is trivial, and $Y(\widehat{\fgl}_1)$ degenerates to the Heisenberg algebra $H_1$. The central charge remains $c = 1$.
 \item At generic $\sigma_3 \neq 0$: the structure function $g(u)$ is a nontrivial rational function with poles at $u = -h_i$ and zeros at $u = h_i$. The $R$-matrix on $\mathrm{Fock}_1$ is $R^{(1)}(u) = g(u)$ (the scalar structure function). On $\mathrm{Fock}_2 = \mathrm{span}\{|\mathrm{row}\rangle, |\mathrm{col}\rangle\}$, the diagonal entries of the $R$-matrix are:
 \[
 R_{(\mathrm{row},\mathrm{row})}(u) = g(u) \cdot g(u + h_1), \qquad
 R_{(\mathrm{col},\mathrm{col})}(u) = g(u) \cdot g(u + h_2),
 \]
 as products over box contents $c(s)$ of the partitions $(2)$ and $(1,1)$ respectively.
 \item The charge-$1$ unitarity $g(u) g(-u) = 1$ holds at all parameter values. The charge-$2$ unitarity $R_{12}(u) R_{21}(-u) = \mathrm{id}$ is a matrix identity involving off-diagonal entries (Theorem~\ref{thm:yang-r-matrix-ybe}).
 \item The invariant $\sigma_2 = h_1 h_2 + h_1 h_3 + h_2 h_3$ generates a rescaling, not a true deformation: the Yangian at parameters $(h_1,h_2,h_3)$ is isomorphic to the one at $(\lambda h_1, \lambda h_2, \lambda h_3)$ for $\lambda \neq 0$.
\end{enumerate}
\end{proposition}

\begin{proof}
Part~(i): at \(\sigma_3=0\), say \(h_2=0\), then \(h_3=-h_1\) and
\[
g(u)=\frac{(u-h_1)u(u+h_1)}{(u+h_1)u(u-h_1)}=1 .
\]
Parts~(ii)--(iv) are the charge-\(1\) and charge-\(2\) product
identities, with the charge-\(2\) case matching the
Maulik--Okounkov diagonal eigenvalue formula
(Theorem~\ref{thm:mo-chiral-rmatrix}).
\end{proof}

\subsection{The Koszul dual: boundary-to-bulk}
\label{subsec:5d-koszul-dual}

\begin{proposition}[Koszul dual of the affine Yangian from 5d CS]
\label{prop:5d-koszul-dual}
The Koszul dual $A^! = D_{\Ran}(B(Y(\widehat{\fgl}_1)))$ of the 5d boundary algebra satisfies:
\begin{enumerate}[label=\textup{(\roman*)}]
 \item \emph{Parameter inversion}: $A^!$ lives at the inverted Omega-background $(-h_1, -h_2, -h_3)$. Its structure function is $g^!(u) = g(-u) = 1/g(u)$.
 \item \emph{$\kappa_{\mathrm{ch}}$ complementarity}: $\kappa_{\mathrm{ch}}(A) + \kappa_{\mathrm{ch}}(A^!) = \rho_K$, where $\kappa_{\mathrm{ch}}(A) = -\sigma_2$, $\kappa_{\mathrm{ch}}(A^!) = \sigma_2$, and $\rho_K = 0$ for $\fgl_1$ (class $\mathbf{L}$, $h^\vee = 0$). The Koszul conductor $\rho_K$ is family-dependent: $\rho_K = 0$ for free-field/$\fgl_1$ and class $\mathbf{L}$; $\rho_K = 13$ for the Virasoro family (Vol~I).
 \item \emph{$\sigma$-invariant behaviour}: under parameter inversion, $\sigma_2$ is preserved ($\sigma_2(-h) = \sigma_2(h)$) and $\sigma_3$ is negated ($\sigma_3(-h) = -\sigma_3(h)$).
 \item \emph{Shadow class preservation}: $A^!$ is also class $\mathbf{L}$ (shadow depth $3$). The pole structure of $g^!(u) = 1/g(u)$ has the same simple poles (at $u = h_i$ instead of $u = -h_i$), preserving the shadow depth.
 \item \emph{Reflected level}: the current presentation of $A^!$ has Kac--Moody level $k' = -k - 2h^\vee = \sigma_2$ (for $\fgl_1$, $h^\vee = 0$, so $k' = -k$). At $n = 1$ (3d CS), this specializes to the Feigin--Frenkel reflected-current surface of the CE/bar-cochain Koszul target.
\end{enumerate}
\end{proposition}

\begin{proof}
The parameter inversion $g^!(u)=1/g(u)$ follows from the reflection
identity: $g^!(u)=g(-u)$, since evaluating $g$ at the inverted
parameters is equivalent to evaluating the original $g$ at $-u$, and
$g(-u)=1/g(u)$.  The $\kappa_{\mathrm{ch}}$ complementarity is the
identity
$\kappa_{\mathrm{ch}}(A)+\kappa_{\mathrm{ch}}(A^!)=0$ on the free-field
branch.
\end{proof}

\begin{remark}[Costello 5d structural checks]
\label{rem:costello-5d-summary}
The construction is controlled by the CY condition, the invariants
$\sigma_2,\sigma_3$, the reflection identity $g(u)g(-u)=1$, the
coefficient pattern $\phi_3=-2\sigma_3$, the spin-$1$ level
$\Psi=-\sigma_2$, the spin-$2$ Virasoro coefficient
$T_{(3)}T=1/2$, the one-loop Serre correction, the
$\Omega$-background deformation, the charge-$2$ diagonal $R$-matrix,
and the Koszul parameter inversion.
\end{remark}


%% The K3 holomorphic CS hierarchy: 3d -> 5d -> 6d

\section{The holomorphic CS hierarchy for K3}
\label{sec:hcs-k3-hierarchy}

The flat-space hierarchy of Section~\ref{sec:qca-from-qg} (Construction~\ref{constr:hol-cs-hierarchy}) has a K3 specialisation in which the ambient spaces $\C^n$ are replaced by K3-fibered geometries. The three levels of the hierarchy, their boundary algebras, and the dimensional reductions connecting them are the subject of this section.

\begin{conjecture}[K3 holomorphic CS hierarchy]
\label{conj:hcs-k3-hierarchy}
\ClaimStatusConjectured
The holomorphic Chern--Simons hierarchy specialised to K3 consists of three levels, with boundary algebras and dimensional reductions as follows:
\begin{center}
\small
\renewcommand{\arraystretch}{1.3}
\begin{tabular}{lllll}
\toprule
\textbf{Level} & \textbf{Geometry} & \textbf{Boundary algebra} & $E_n$ & \textbf{Construction locus} \\
\midrule
3d & $K3$ sigma model on $C$ & $H_{\mathrm{Muk}}$ (Mukai Heisenberg, rank $24$) & $\Eone$ & CY-A$_2$ and Mukai lattice \\
5d & $K3 \times C \times \R$ & $Y(\fg_{K3})$ (K3 Yangian) & $\Etwo$ & requires 5d K3 hCS construction \\
6d & $K3 \times E \times C$ & $U_{q,t}(\fg_{K3}^{\mathrm{tor}})$ (K3 quantum toroidal) & $E_3$ & requires 6d K3$\times E$ hCS construction \\
\bottomrule
\end{tabular}
\end{center}
The dimensional reductions are:
\begin{enumerate}[label=\textup{(\roman*)}]
 \item \emph{6d $\to$ 5d: shrink $E$}. As $\tau \to i\infty$ (equivalently $q = e^{2\pi i \tau} \to 0$), the elliptic curve $E$ degenerates and the quantum toroidal algebra reduces to the Yangian:
 \[
 U_{q,t}(\fg_{K3}^{\mathrm{tor}}) \;\xrightarrow{q \to 1}\; Y(\fg_{K3}).
 \]
 The trigonometric structure function $G_{K3}(x; \{q_a\})$ reduces to the rational structure function $g_{K3}(u; \{h_a\})$ via $x = e^u$, $q_a = e^{h_a}$.

 \item \emph{5d $\to$ 3d: shrink $C$}. As all Omega-background parameters $h_a \to 0$ (equivalently $\hbar \to 0$), the Yangian deformation disappears:
 \[
 Y(\fg_{K3}) \;\xrightarrow{h_a \to 0}\; U(\widehat{\fg}_{K3}) \;\simeq\; H_{\mathrm{Muk}}.
 \]
 The rational structure function $g_{K3}(u) = \prod_{a=1}^{24} (u - h_a)/(u + h_a) \to 1$ (trivial $R$-matrix), and the Yangian degenerates to the current algebra of the Mukai Heisenberg.
\end{enumerate}
The 3d level is established by CY-A$_2$ (Theorem~CY-A, $d = 2$). CY-A$_3$ supplies the $d = 3$ framed object-level assignment on its H1--H4 locus (Theorem~\ref{thm:derived-framing-m3b2} and Theorem~\ref{thm:cy-to-chiral-d3}); it does not produce a global $\Phi_3$ functor on all smooth proper CY$_3$ categories. The 5d and 6d levels remain conjectural because they depend on the constructions of the K3 Yangian (Conjecture~\ref{conj:k3-yangian}, Section~\ref{sec:rtt-ope-dictionary}) and K3 quantum toroidal algebra (Conjecture~\ref{conj:6d-k3xe}) respectively, which require the 5d/6d holomorphic CS frameworks (not merely CY-A$_3$).
\end{conjecture}

\subsection{The Omega-background at each level}
\label{subsec:k3-omega-background}

The Omega-background parameters control the deformation from classical to quantum at each level.

\begin{construction}[K3 Omega-background hierarchy]
\label{constr:k3-omega-hierarchy}
The Omega-background data at each level of the K3 hierarchy:
\begin{enumerate}[label=\textup{(\roman*)}]
 \item \emph{3d ($K3$ sigma model)}: no Omega-background. The boundary algebra $H_{\mathrm{Muk}}$ is classical (free field, $\hbar = 0$). The $R$-matrix is trivial and $\kappa_{\mathrm{ch}}(H_{\mathrm{Muk}}) = 2$ (from the CY-A$_2$ construction; $\kappa_{\mathrm{ch}} = \chi^{\mathrm{CY}}_2(K3) = 2$).

 \item \emph{5d ($K3 \times C \times \R$)}: Omega-background given by $24$ parameters $(h_1, \ldots, h_{24})$ in the complexified Mukai lattice, subject to the CY$_2$ constraint
 \[
 \sum_{a=1}^{24} h_a = 0.
 \]
 This leaves $23$ free parameters (the Mukai lattice $\widetilde{\Lambda}_{K3}$ has rank $24$, minus one constraint). The effective deformation parameter is the Newton power sum $p_3 = \sum h_a^3$; when $p_3 = 0$, the structure function degenerates to $g_{K3}(u) = 1$.

 For the flat-space $\C^3$, the $24$ parameters specialise to $(h_1, h_2, h_3, 0, \ldots, 0)$ with $h_1 + h_2 + h_3 = 0$, recovering the standard Omega-background of Section~\ref{sec:costello-5d-construction}.

 \item \emph{6d ($K3 \times E \times C$)}: the $24$ Mukai parameters $(h_1, \ldots, h_{24})$ are supplemented by the modular parameter $\tau$ of the elliptic curve $E$, giving $24$ free parameters in total. The nome $q = e^{2\pi i \tau}$ satisfies $|q| < 1$ since $\Imag(\tau) > 0$, ensuring convergence of all $q$-expansions. The multiplicative parameters are $q_a = e^{h_a}$, and the CY$_2$ constraint becomes
 \[
 \prod_{a=1}^{24} q_a = 1.
 \]
\end{enumerate}
\end{construction}

\subsection{Boundary conditions and the bulk-defect structure}
\label{subsec:k3-boundary-defect}

At each level, the theory on a manifold of the form $M \times \R$ or $M \times S^1$ has two types of boundary conditions: Dirichlet (fixing boundary values) and Neumann (fixing normal derivatives). These produce different algebras.

\begin{construction}[K3 boundary and defect algebras]
\label{constr:k3-boundary-defect}
At each level of the K3 hierarchy:
\begin{center}
\small
\renewcommand{\arraystretch}{1.3}
\begin{tabular}{llll}
\toprule
\textbf{Level} & \textbf{Boundary (Dirichlet)} & \textbf{Bulk (derived center)} & \textbf{Defect (Koszul dual)} \\
\midrule
3d & $H_{\mathrm{Muk}}$ & $Z^{\mathrm{der}}_{\mathrm{ch}}(H_{\mathrm{Muk}})$ & $H_{\mathrm{Muk}}^!$ \\
5d & $Y(\fg_{K3})$ & $Z^{\mathrm{der}}_{\mathrm{ch}}(Y(\fg_{K3}))$ & $Y(\fg_{K3})^!$ \\
6d & $U_{q,t}(\fg_{K3}^{\mathrm{tor}})$ & $Z^{\mathrm{der}}_{\mathrm{ch}}(U_{q,t})$ & $U_{q,t}^!$ \\
\bottomrule
\end{tabular}
\end{center}
The bulk-boundary coupling at each level is the canonical map $\mu_{\mathrm{bb}} \colon Z^{\mathrm{der}}_{\mathrm{ch}}(A) \otimes A^! \to A^!$ of Construction~\ref{constr:universal-defect}, making the defect algebra $A^!$ a module over the bulk. The defect algebra at the 5d level has inverted parameters: $Y(\fg_{K3})^!$ lives at $(-h_1, \ldots, -h_{24})$ with structure function $g_{K3}^!(u) = g_{K3}(-u) = 1/g_{K3}(u)$ (Proposition~\ref{prop:5d-koszul-dual}, generalised from $\C^3$ to K3). The $\kappa_{\mathrm{ch}}$-complementarity gives $\kappa_{\mathrm{ch}}(A) + \kappa_{\mathrm{ch}}(A^!) = 0$ for the free-field/$\fgl_1$ branch ($\rho_K = 0$; Koszul conductor vanishes for class $\mathbf{L}$).
\end{construction}

\subsection{Wilson lines and coproducts}
\label{subsec:k3-wilson-lines}

The Wilson lines (defects) at each level produce modules over the boundary algebra. Their fusion is controlled by the coproduct.

\begin{construction}[K3 Wilson line hierarchy]
\label{constr:k3-wilson-hierarchy}
The Wilson lines at each level of the K3 hierarchy:
\begin{enumerate}[label=\textup{(\roman*)}]
 \item \emph{3d: point operators on $C$}. Indexed by Mukai vectors $v \in H^*(K3, \Z)$, these produce Fock modules $\cF_v$ of $H_{\mathrm{Muk}}$. The coproduct is \emph{primitive}:
 \[
 \Delta_0(J_{a,n}) = J_{a,n}^L + J_{a,n}^R,
 \]
 as established in Proposition~\ref{prop:heisenberg-primitive}.

 \item \emph{5d: line operators on $C$}. The line operators wrapping $C$ in $K3 \times C \times \R$ are indexed by Fock modules $\cF_\lambda$ of $Y(\fg_{K3})$. The coproduct is the \emph{Miura factorisation}:
 \[
 \Delta_z(T(u)) = T^L(u) \cdot T^R(u - z),
 \]
 where $z$ is the spectral parameter and $T(u) = 1 + \sum_{s \geq 1} \psi_s u^{-s}$ is the transfer matrix. Fusion of Wilson lines is controlled by the $R$-matrix $R_{K3}(u)$ (Conjecture~\ref{conj:rtt-ope-dictionary}).

 \item \emph{6d: surface operators wrapping $E$}. The Wilson surfaces wrapping $E$ in $K3 \times E \times C$ are indexed by BPS states labelled by Mukai vectors together with $E$-winding numbers. The coproduct is a $q$-deformed Miura factorisation:
 \[
 \Delta_x^{(q)}(\Psi(w)) = \Psi^L(w) \cdot \Psi^R(w/x),
 \]
 where $x = e^z$ is the multiplicative spectral parameter. This is conjectural.
\end{enumerate}
\end{construction}

\subsection{The structure function at each level}
\label{subsec:k3-structure-fn}

The structure function encodes the $R$-matrix data and determines the Yangian/quantum group relations. Its type changes across the hierarchy.

\begin{proposition}[K3 structure function hierarchy]
\label{prop:k3-structure-fn-hierarchy}
The structure function at each level of the K3 hierarchy:
\begin{enumerate}[label=\textup{(\roman*)}]
 \item \emph{3d (classical)}: $g_{K3}^{(3d)}(u) = 1$ identically (trivial $R$-matrix).

 \item \emph{5d (rational)}: $g_{K3}^{(5d)}(u) = \prod_{a=1}^{24} \frac{u - h_a}{u + h_a}$, a degree-$(24,24)$ rational function satisfying:
 \begin{itemize}
 \item Reflection: $g_{K3}(u) \cdot g_{K3}(-u) = 1$.
 \item CY$_2$ constraint: $\phi_1 = -2\sum h_a = 0$.
 \item Classical limit: $g_{K3}(u) \to 1$ as $u \to \infty$.
 \end{itemize}

 \item \emph{6d (trigonometric)}: $G_{K3}^{(6d)}(x) = \prod_{a=1}^{24} \frac{1 - q_a x}{1 - q_a^{-1} x}$, where $q_a = e^{h_a}$, satisfying:
 \begin{itemize}
 \item Inversion: $G_{K3}(x) \cdot G_{K3}(1/x) = 1$ (requires CY$_2$: $\prod q_a = 1$).
 \item Reduction: $G_{K3}(e^{\varepsilon u}; \{e^{\varepsilon h_a}\}) \to g_{K3}(u; \{h_a\})$ as $\varepsilon \to 0$.
 \item Classical limit: $G_{K3}(x) \to 1$ as all $q_a \to 1$.
 \end{itemize}
\end{enumerate}
Parts~(i) and~(ii) are direct substitutions. Part~(iii) is an
algebraic identity for the product formula; the inversion identity
requires the CY$_2$ constraint \(\prod q_a=1\). The \(6d\to5d\)
reduction is the \(\varepsilon\to0\) expansion of the same product.
\end{proposition}

\begin{proof}
Part~(i): at the 3d level, all Omega-background parameters are zero, so $g = \prod_{a=1}^{24} u/u = 1$.

Part~(ii): the reflection identity $g(u) g(-u) = 1$ follows from factorwise cancellation: each factor $(u - h_a)/(u + h_a)$ paired with $(-u - h_a)/(-u + h_a) = (u + h_a)/(u - h_a)$ gives product~$1$. The CY$_2$ vanishing of $\phi_1$ follows from the log-expansion $\log g(u) = -2 \sum_{k \text{ odd}} p_k / (k u^k)$, so $\phi_1 = -2 p_1 = -2 \sum h_a = 0$.

Part~(iii): the inversion identity $G(x) G(1/x) = \prod q_a^2 = (\prod q_a)^2 = 1$ (where the last equality uses $\prod q_a = 1$). The reduction is the standard passage from trigonometric to rational: expand $e^{\varepsilon h_a} = 1 + \varepsilon h_a + O(\varepsilon^2)$ and $e^{\varepsilon u} = 1 + \varepsilon u + O(\varepsilon^2)$; the leading terms in the product $\prod (1 - e^{\varepsilon h_a} e^{\varepsilon u})/(1 - e^{-\varepsilon h_a} e^{\varepsilon u})$ match $\prod (u - h_a)/(u + h_a)$ at $O(1)$ with \(O(\varepsilon^2)\) corrections. At fifteen generic points the relative error is bounded by \(6 \times 10^{-3}\) for \(\varepsilon = 0.01\).
\end{proof}

\subsection{The \texorpdfstring{$\kappa_\bullet$}{kappa-.}-spectrum across the hierarchy}
\label{subsec:k3-kappa-hierarchy}

\begin{remark}[\texorpdfstring{$\kappa_\bullet$}{kappa-.}-spectrum of the K3 hierarchy]
\label{rem:kappa-k3-hierarchy}
Each level of the K3 hierarchy contributes to the $\kappa_\bullet$-spectrum of $K3 \times E$:
\begin{center}
\small
\renewcommand{\arraystretch}{1.3}
\begin{tabular}{llrl}
\toprule
\textbf{Subscript} & \textbf{Origin} & \textbf{Value} & \textbf{Level} \\
\midrule
$\kappa_{\mathrm{ch}}^{\mathrm{Hodge}}(K3)$ & CY-A$_2$ correspondence $\Phi_2$ & $2$ & 3d (proved) \\
$\kappa_{\mathrm{ch}}^{\mathrm{Heis}}(K3 \times E)$ & Beauville-additive Heisenberg--Cartan: $2 + 1$ & $3$ & 5d/6d \\
$\kappa_{\mathrm{BKM}}(\Delta_5)$ & Igusa cusp form $\Delta_5$ weight & $5$ & 6d \\
$\kappa_{\mathrm{cat}}(K3)$ & $\chi(\cO_{K3}) = 2$ (fibre) & $2$ & 3d \\
$\kappa_{\mathrm{cat}}(K3 \times E)$ & K\"unneth: $\chi(\cO_{K3}) \cdot \chi(\cO_E) = 0$ & $0$ & 5d/6d \\
$\kappa_{\mathrm{fiber}}(K3)$ & Mukai lattice rank & $24$ & All levels \\
\bottomrule
\end{tabular}
\end{center}
The five distinct values $\{\kappa_{\mathrm{cat}}, \kappa_{\mathrm{ch}}^{\mathrm{Hodge}}, \kappa_{\mathrm{ch}}^{\mathrm{Heis}}, \kappa_{\mathrm{BKM}}(\Delta_5), \kappa_{\mathrm{fiber}}\} = \{0, 0, 3, 5, 24\}$ on $K3 \times E$ are the outputs of five distinct constructions: the Tier-(i) CY-datum intrinsic $\kappa_{\mathrm{cat}}(K3 \times E) = \chi(\cO_{K3 \times E}) = \chi(\cO_{K3}) \cdot \chi(\cO_E) = 2 \cdot 0 = 0$ (K\"unneth-multiplicative); the compact total-space Hodge supertrace $\kappa_{\mathrm{ch}}^{\mathrm{Hodge}}(K3 \times E) = 0$ by Serre at odd $d = 3$; the Stage-2 Heisenberg rank $\kappa_{\mathrm{ch}}^{\mathrm{Heis}} = 3$; the Borcherds weight $\kappa_{\mathrm{BKM}}(\Delta_5) = 5$; and the Mukai lattice rank $\kappa_{\mathrm{fiber}} = 24$.
\end{remark}

\subsection{The defect algebra DAG: Koszul duality construction verification}
\label{subsec:defect-koszul-dag}

The defect algebra (chiral quantum group) arises from the five-stage construction of Remark~\ref{rem:hcs-construction}, instantiated for each CY target as a directed acyclic graph (DAG) of computations. The five stages are:
\begin{center}
\small
\renewcommand{\arraystretch}{1.3}
\begin{tabular}{clll}
\toprule
\textbf{Stage} & \textbf{Input} & \textbf{Output} & \textbf{Constraint} \\
\midrule
1 & CY$_d$ category & Bulk algebra $A$ with $\kappa_{\mathrm{ch}}(A)$ & Shadow class G/L/C/M \\
2 & Bulk $A$ & Bar complex $B(A)$, differentials & $d_{\mathrm{bar}} = 0$ for class G \\
3 & Bar $B(A)$ & Defect $A^! = D_{\Ran}(B(A))$, $\kappa_{\mathrm{ch}}(A^!)$ & Structure function inversion \\
4 & Bulk + Defect & Conductor $K = \kappa_{\mathrm{ch}} + \kappa_{\mathrm{ch}}^!$ & $K = 0$ (free-field) or $K > 0$ (class M) \\
5 & Bar + Defect & Morita $\Omega(B(A)) \simeq A$ & Koszul locus check \\
\bottomrule
\end{tabular}
\end{center}

\begin{conjecture}[Defect algebra of \texorpdfstring{$K3 \times E$}{K3 x E} via Koszul duality]
\label{conj:defect-koszul-k3xe}
\ClaimStatusConjectured
Let $A_{K3 \times E}^{(K3,E)}$ denote the bulk chiral algebra on $E$ from the framed K3-fibre specialisation $\Phi_3^{(K3,E)}(D^b(\Coh(K3 \times E)))$ (Theorem~CY-A$_3$, Theorem~\ref{thm:cy-to-chiral-d3}; the 6d framework of Conjecture~\ref{conj:6d-k3xe} provides an independent route). Then the five-stage construction produces two branches:
\begin{enumerate}[label=\textup{(\roman*)}]
 \item \emph{Free-field/Heisenberg branch} (class~G, $\fgl_1$):
 $\kappa_{\mathrm{ch}}^{\mathrm{Heis}}(A) = 3$,
 $\kappa_{\mathrm{ch}}^{\mathrm{Heis}}(A^!) = -3$,
 conductor $K = 0$. The compact total-space algebra
 $A_{K3\times E}$ has $\kappa_{\mathrm{ch}}=0$; the value $3$
 is the Heisenberg/free-field specialization. The defect algebra
 has inverted structure function $g^!_{K3}(u) = 1/g_{K3}(u)$
 and reversed braiding
 $\Rep^{\Etwo}(A^!) \simeq \Rep^{\Etwo}(A)^{\mathrm{rev}}$.
 The bar-cobar adjunction recovers $A$ on the Koszul locus:
 $\Omega(B(A)) \simeq A$.
 \item \emph{Mukai-doubling branch} (class~M, Theorem-C row~$\mathsf{B}$):
 $\kappa_{\mathrm{ch}}^{\mathrm{Mukai}}(A) = 4$,
 $\kappa_{\mathrm{ch}}^{\mathrm{Mukai}}(A^!) = 4$,
 conductor $K^{\kappa_{\mathrm{ch}}^{\mathrm{Mukai}}} = 8 = 2\,c_+(\mathrm{Mukai}(K3))$
 (Theorem~\ref{thm:qca-hbar-minus-eighth-lusztig};
 Remark~\ref{rem:qca-three-faces-of-eight}). The defect algebra is
 a deformation of the chiral Koszul dual of the BKM vertex algebra
 $V_{\mathfrak{g}_{\Delta_5}}$, with infinite shadow depth. The
 Mukai integer \(8\) places this branch outside the free-field Koszul
 class and gives the order-eight comparison datum at the Lusztig
 specialisation \(\ell=8\).
\end{enumerate}
The two branches coexist: the Heisenberg branch captures the abelian/$\fgl_1$ free-field sector, while the Mukai-doubling branch captures the non-abelian/interacting BKM sector controlled by $V_{\mathfrak{g}_{\Delta_5}}$. The Borcherds weight $\kappa_{\mathrm{BKM}}(\Delta_5) = c_{\Delta_5}(0)/2 = 5$ is the universal level-$4$ automorphic invariant of the input Siegel form; it is structurally distinct from the Mukai integer \(8\). The two numbers \(5\) and \(8\) are separate diagnostics.

Conjecture~\ref{conj:defect-koszul-k3xe} lies over the framed $\Phi_3$
object and over the 6d construction of Conjecture~\ref{conj:6d-k3xe};
the defect branch inherits those hypotheses.
\end{conjecture}

\begin{proposition}[K3 defect algebra: the \texorpdfstring{$d = 2$}{d = 2} verified case]
\label{prop:defect-koszul-k3}
For $K3$ at $d = 2$ (CY-A$_2$ proved), the five-stage construction is fully verified:
\begin{enumerate}[label=\textup{(\roman*)}]
 \item Bulk: $A_{K3} = \mathcal H_{\mathrm{Muk}}$ on the free Mukai--Heisenberg branch; $\kappa_{\mathrm{ch}}^{\mathrm{Hodge}}(K3)=2$, $\kappa_{\mathrm{ch}}^{\mathrm{Heis}}(\mathcal H_{\mathrm{Muk}})=24$, class~G.
 \item Bar: $B(A_{K3})$ has $24$ generators (Mukai rank), $d_{\mathrm{bar}} = 0$ (free-field), no $A_\infty$ corrections.
 \item Defect: $A_{K3}^! = \mathcal H_{\mathrm{Muk}}^!$ with $\kappa_{\mathrm{ch}}^{\mathrm{Hodge}}(A^!) = -2$, $g^!(u) = 1/g_{K3}(u)$.
 \item Conductor: $K_{\mathrm{Hodge}} = 2 + (-2) = 0$ (free-field branch).
 \item Morita: $\Omega(B(A_{K3})) \simeq A_{K3}$ at $\Etwo$-level on the Koszul locus.
\end{enumerate}
The structure function inversion $g_{K3}^!(u) = 1/g_{K3}(u)$ follows from the algebraic unitarity identity $g_{K3}(u) \cdot g_{K3}(-u) = 1$ (Proposition~\ref{prop:k3-structure-fn-hierarchy}(ii)) together with the Verdier parameter inversion $h_a \mapsto -h_a$ under Koszul duality (Construction~\ref{constr:k3-boundary-defect}). Braiding reversal: $\Rep^{\Etwo}(A_{K3}^!) \simeq \Rep^{\Etwo}(A_{K3})^{\mathrm{rev}}$ (the sign flip $\kappa_{\mathrm{ch}}^{\mathrm{Hodge}} \to -\kappa_{\mathrm{ch}}^{\mathrm{Hodge}}$ inverts the leading $R$-matrix coefficient on this scalar line).
\end{proposition}

\begin{proof}
Parts~(i)--(iv): the Mukai--Heisenberg branch $\mathcal H_{\mathrm{Muk}}$ is the free-field output of $\Phi$ at $d = 2$ (Section~\ref{sec:cy-d2-functor}); its source Hodge scalar is $\kappa_{\mathrm{ch}}^{\mathrm{Hodge}}(K3) = \chi^{\mathrm{CY}}_2(K3) = 2$ (Proposition~\ref{prop:cy-kappa-d2}), and its Heisenberg scalar is the Mukai rank $24$. The branch is class~G because the Heisenberg bracket $J^{(a)}(z) J^{(b)}(w) \sim \omega^{ab}/(z-w)^2$ has no first-order pole, so $\mu = 0$ and $d_{\mathrm{bar}} = 0$. Koszul duality for free-field algebras gives $\kappa_{\mathrm{ch}}^{\mathrm{Hodge},!} = -\kappa_{\mathrm{ch}}^{\mathrm{Hodge}}$ with conductor $K_{\mathrm{Hodge}} = 0$ (matching $e1\_koszul\_three\_families.py$ for the Heisenberg family at rank~$24$).

Part~(v): for $A_{K3} = \mathcal H_{\mathrm{Muk}}$ on the chirally Koszul locus, the raw-direct-sum scope applies because $A_{K3}$ is class~G. The bar-cobar adjunction $\Omega \dashv B$ (cobar left, bar right) is a Quillen equivalence (Vol~I, Theorem~B, raw-direct-sum scope on classes G and L; see Corollary~\cite{VolI}), so $\Omega(B(A_{K3})) \simeq A_{K3}$ as $\Etwo$-algebras. The $\Etwo$-level follows from CY dimension $d = 2$ (Dunn additivity: $\Eone \otimes \Eone \simeq \Etwo$).

\end{proof}

\begin{remark}[Three dualities on K3 are distinct]
\label{rem:three-dualities-k3}
The K3 defect algebra $A_{K3}^!$ sits alongside three other K3-attached objects, each with its own role:
\begin{enumerate}[label=\textup{(\alph*)}]
 \item the \emph{mirror} K3 algebra $A_{K3^v}$ has $\kappa_{\mathrm{ch}}^{\mathrm{Hodge}} = 2$ (HMS preserves the Hodge scalar, so the mirror reads the same chiral characteristic);
 \item the \emph{derived center} $Z^{\mathrm{der}}_{\mathrm{ch}}(A_{K3})$ is the bulk algebra controlling the $49$-dimensional Drinfeld double;
 \item the \emph{Drinfeld center} $\cZ(\Rep^{\Eone}(A_{K3}))$ is the braided monoidal \emph{category} carrying the $\Etwo$ promotion.
\end{enumerate}
The $\kappa_{\mathrm{ch}}^{\mathrm{Hodge}}$-sign diagnostic separates Koszul duality from HMS: $\kappa_{\mathrm{ch}}^{\mathrm{Hodge}}(A_{K3}^!) = -2$ (Koszul dual flips the sign), while $\kappa_{\mathrm{ch}}^{\mathrm{Hodge}}(A_{K3^v}) = 2$ (mirror preserves the sign).
\end{remark}

\subsection{The master diagram}
\label{subsec:k3-master-diagram}

\begin{conjecture}[Master diagram of the K3 holomorphic CS hierarchy]
\label{conj:k3-master-diagram}
\ClaimStatusConjectured
The complete K3 holomorphic CS hierarchy connects three levels via dimensional reduction, with the structure function transitioning from trigonometric (6d) to rational (5d) to classical (3d):
\[
\begin{tikzcd}[column sep=1.2cm, row sep=1.8cm]
 K3 \times E \times C \ar[d, "\text{6d hol theory}"] \ar[r, "\text{shrink } E"]
 & K3 \times C \times \R \ar[d, "\text{5d hol CS}"] \ar[r, "\text{shrink } C"]
 & K3 \text{ on } C \ar[d, "\text{3d } \sigma\text{-model}"] \\
 U_{q,t}(\fg_{K3}^{\mathrm{tor}}) \ar[r, "q \to 1"]
 \ar[d, "B(-)"]
 & Y(\fg_{K3}) \ar[r, "h_a \to 0"]
 \ar[d, "B(-)"]
 & H_{\mathrm{Muk}} \ar[d, "B(-)"] \\
 B(U_{q,t}) \ar[d, "D_{\Ran}"]
 & B(Y(\fg_{K3})) \ar[d, "D_{\Ran}"]
 & B(H_{\mathrm{Muk}}) \ar[d, "D_{\Ran}"] \\
 U_{q,t}^! \ar[d, "\Rep^{E_3}"]
 & Y(\fg_{K3})^! \ar[d, "\Rep^{\Etwo}"]
 & H_{\mathrm{Muk}}^! \ar[d, "\Rep^{\Eone}"] \\
 \text{chiral QG}_{K3}^{(6d)}
 & \text{chiral QG}_{K3}^{(5d)}
 & \text{conformal blocks}_{K3}
\end{tikzcd}
\]
Each column is the holomorphic CS $\to$ chiral quantum group construction of Remark~\ref{rem:hcs-construction}, specialised to the K3 geometry at the given dimension. Each horizontal arrow is a dimensional reduction. The vertical arrows are:
\begin{itemize}
 \item $B(-)$: the ordered bar complex (chiral CE chains; Construction~\ref{constr:chiral-ce});
 \item $D_{\Ran}$: the Verdier dual on $\Ran(C)$ (Koszul duality; Construction~\ref{constr:universal-defect});
 \item $\Rep^{E_k}$: the braided representation category at $E_k$-level $k = n$ (the complex dimension).
\end{itemize}
The 3d column is established (CY-A$_2$ proved). The 5d column is partially established: the flat-space boundary algebra is proved (Theorem~\ref{thm:5d-boundary-yangian}), and the K3 specialisation is conjectural. The 6d column is conditional on the non-Lagrangian 6d algebraic framework (Remark~\ref{rem:6d-not-lagrangian}), the compact \(K3\times E\) hCS--Hall comparison, and the K3 quantum-toroidal construction. Theorem~\ref{thm:cy-to-chiral-d3} supplies the framed object-level \(\Phi_3\) input; the remaining condition is that dimensional reduction commute with bar, Koszul duality, and representation categories in this compact 6d geometry.
\end{conjecture}



%% Comparison: hCS route vs sigma model route

\subsection{Comparison: holomorphic CS vs.\ sigma model}
\label{subsec:hcs-vs-sigma-adversarial}
\index{holomorphic CS!vs sigma model}
\index{sigma model!vs holomorphic CS}

The K3 chiral algebra arises from three constructions on the same K3 geometry: the holomorphic CS/Costello--Li boundary algebra $A_E$ (Route~A), the CY-to-chiral correspondence output $H_{\mathrm{Muk}} = \Phi_2(D^b(\Coh(K3)))$ (Route~B), and the K3 sigma model VOA $V_{K3}$ (Route~C). Each is a distinct algebraisation of the same input, with its own central charge, shadow class, and operadic level; the four-vector comparison below identifies which invariant each construction reads.

\begin{proposition}[Route comparison across four vectors]
\label{prop:hcs-vs-sigma}
\index{modular characteristic!comparison across routes}
The following apparent contradictions between the holomorphic CS and sigma model approaches are resolved:
\begin{enumerate}[label=\textup{(\roman*)}]
 \item \emph{Central charge.} The Costello--Li KS boundary algebra $A_E$ (Construction~\textup{\ref{constr:costello-li-k3xe}}) has $c = 24$ ($24$~free bosons from $H^*(K3)$). The $\cN = 4$ sigma model VOA $V_{K3}$ has $c = 6$ (Definition~\textup{\ref{def:n4-sca-c6}}). These are \emph{different algebras}: $A_E$ is the B-model reduction (Hochschild data), $V_{K3}$ is the full SCFT. The CY-A$_2$ functor output $\Phi(D^b(\Coh(K3))) = H_{\mathrm{Muk}}$ coincides with $A_E$ as an algebra ($c = 24$), not with $V_{K3}$ ($c = 6$). The common number $2$ is the K3 Hodge scalar $\kappa_{\mathrm{ch}}^{\mathrm{Hodge}}(K3)=\chi(\cO_{K3})$, not the central charge of either algebra.

 \item \emph{Shadow class.} $H_{\mathrm{Muk}}$ is class~$\mathbf{G}$ (Gaussian, shadow depth $r = 2$): it is a free-field VOA with $m_k = 0$ for $k \geq 3$ (K3 is formal by DGMS). The $\cN = 4$ SCA $V_{K3}$ is class~$\mathbf{M}$ (infinite tower, $r_{\max} = \infty$): the stress tensor introduces non-trivial $A_\infty$ corrections with critical discriminant $\Delta = 8\kappa_{\mathrm{ch}}^{\mathrm{SCFT}} S_4 = 20/39 \neq 0$ (Computation~\textup{\ref{comp:shadow-tower-k3}}). Resolution: class~$\mathbf{G}$ is the shadow class of the $\Phi$-output $H_{\mathrm{Muk}}$; class~$\mathbf{M}$ is the shadow class of the physical SCFT $V_{K3}$. Same geometry, different algebraizations.

 \item \emph{$E_n$~structure.} K3 alone ($d = 2$) has native $\Etwo$ from the $\bS^2$-framing (Arrow~3 of the cyclic-to-chiral passage). $K3 \times E$ ($d = 3$) has native $\Eone$ from ordered factorization. Resolution: these apply to \emph{different CY dimensions}. The $\Etwo$ on $H_{\mathrm{Muk}}$ is automatic (free fields are $E_\infty$). The $\Eone$ on $Y(\fg_{K3})$ is the native level at $d = 3$; $\Etwo$ is derived via the Drinfeld center (Conjecture~\textup{\ref{conj:k3-master-diagram}}). The classical limit $h_a \to 0$ consistently sends $Y(\fg_{K3})$ ($\Eone$) to $H_{\mathrm{Muk}}$ ($\Etwo \hookrightarrow \Eone$, forgetful).

 \item \emph{Yangian.} The 5d holomorphic CS with Omega-background $(h_1, \ldots, h_{24})$ produces $Y(\fg_{K3})$; the sigma model is the classical $h_a \to 0$ limit, with image the abelian Heisenberg $H_{\mathrm{Muk}}$ on the Mukai lattice. The Yangian deformation is the Omega-background lift: $Y(\fg_{K3}) \xrightarrow{h_a \to 0} H_{\mathrm{Muk}}$ sends $g_{K3}(u) \to 1$ (trivial $R$-matrix) and $\Delta_z \to \Delta_0$ (primitive coproduct). The leading non-trivial term in $\log g_{K3}(u)$ is $O(h^3)$, controlled by $p_3 = \sum h_a^3$; the CY$_2$ constraint $\sum h_a = 0$ kills the $O(h)$ contribution.
\end{enumerate}
All four comparison vectors index distinct algebraisations of the same K3 geometry, with the central charge, Heisenberg/sigma-model split, $E_n$-level shift, and Omega-background deformation reading the K3 data through different construction layers. The comparison table:
\begin{center}
\small
\renewcommand{\arraystretch}{1.3}
\begin{tabular}{lccccl}
\toprule
\textbf{Route} & $c$ & \textbf{distinguished scalar} & \textbf{Shadow class} & $E_n$ & \textbf{Construction locus} \\
\midrule
$A_E$ (KS boundary) & $24$ & $\rho_{\mathrm{KS}}=24$ & G & $E_\infty$ & requires KS-boundary construction \\
$H_{\mathrm{Muk}} = \Phi(D^b(\Coh(K3)))$ & $24$ & $\kappa_{\mathrm{ch}}^{\mathrm{Hodge}}=2$; $\rho_{\mathrm{Muk}}=24$ & G & $\Etwo$ & CY-A$_2$ \\
$V_{K3}$ ($\cN = 4$ SCA) & $6$ & $\kappa_{\mathrm{ch}}^{\mathrm{SCFT}}=2$ & M & $\Etwo$ & $\cN=4$ SCA construction \\
$Y(\fg_{K3})$ (Yangian) & -- & $\kappa_{\mathrm{ch}}^{\mathrm{Heis}}=3$ & G & $\Eone$ & requires K3 Yangian construction \\
\bottomrule
\end{tabular}
\end{center}
The scalar entries $24$ and $2$ do not name one invariant. Route~A uses the OPE-level rank trace $\rho_{\mathrm{KS}}=\tr(\delta^{ab}) = 24$, while Routes~B and~C use respectively the K3 Hodge scalar and the normalised small-$\cN=4$ scalar, both equal to $\chi(\cO_{K3}) = 2$ on this surface. Route~D gives $\kappa_{\mathrm{ch}}^{\mathrm{Heis}} = 3$ by additivity: $\kappa_{\mathrm{ch}}^{\mathrm{Heis}}(K3 \times E) = 2 + 1$. The boundary-to-sigma ratio is $24/2 = 12$ (Proposition~\textup{\ref{prop:boundary-sigma-ratio}}).

\end{proposition}

\begin{proof}
Item~(i): $c(A_E) = 24$ by generator count (Construction~\ref{constr:costello-li-k3xe}); $c(V_{K3}) = 6$ by definition (Definition~\ref{def:n4-sca-c6}); $\kappa_{\mathrm{ch}}^{\mathrm{SCFT}}(V_{K3}) = 2$ by Proposition~\ref{prop:kappa-k3}; $\kappa_{\mathrm{ch}}^{\mathrm{Hodge}}(H_{\mathrm{Muk}}) = 2$ by Proposition~\ref{prop:cy-kappa-d2} (the Serre duality $\mathbb{S}_{\cC} = [2]$ kills the one-loop correction). Item~(ii): class~G for $H_{\mathrm{Muk}}$ by formality (DGMS) and shadow depth computation (all $S_r = 0$ for $r \geq 3$); class~M for $V_{K3}$ by Computation~\ref{comp:shadow-tower-k3} (the computed support packet is mixed and has scalar witness $\Delta = 20/39 \neq 0$). Item~(iii): $\Etwo$ at $d = 2$ from the $\bS^2$-framing (Theorem~\ref{thm:cy-to-chiral}, Arrow~3); $\Eone$ at $d = 3$ from ordered factorization; the forgetful functor $\Etwo \hookrightarrow \Eone$ gives compatibility. Item~(iv): the Yangian classical limit $g_{K3}(u; \{h_a\}) = \prod(u - h_a)/(u + h_a) \to 1$ as $h_a \to 0$ is verified directly; the CY$_2$ constraint $\sum h_a = 0$ kills the $O(h)$ term in $\log g_{K3}$; the leading deformation is $O(h^3)$ with coefficient $(-2/3) p_3 / u^3$ where $p_3 = \sum h_a^3$.
\end{proof}

\begin{remark}[Hodge decomposition and the generator discrepancy]
\label{rem:hcs-sigma-deepened}
The $24$~vs.~$8$~generator discrepancy between $H_{\mathrm{Muk}}$ and $V_{K3}$ is resolved by the HKR decomposition. The $24$~currents of $H_{\mathrm{Muk}}$ decompose as $h^{0,0} + h^{2,0} + h^{0,2} + h^{1,1} + h^{2,2} = 1 + 1 + 1 + 20 + 1 = 24$ (the full Hodge diamond of K3). The $8$~generators of $V_{K3}$ are $(T, G^\pm, \widetilde{G}^\pm, J^{++}, J^{--}, J^3)$ from the $\cN = 4$ superconformal structure. These count generators of \emph{different algebras}: the ratio $24/8 = 3 = \dim_\C(K3) + 1$.

The shadow tower of $V_{K3}$ exhibits sign alternation from depth~$4$ onward ($S_4 = 5/156 > 0$, $S_5 = -1/26 < 0$, $S_6 > 0$, \ldots) with critical discriminant $\Delta = 20/39$, confirming class~$\mathbf{M}$. All $S_r$ in the computed range $2 \leq r \leq 12$ are nonzero. In contrast, $H_{\mathrm{Muk}}$ has $S_r = 0$ for all $r \geq 3$ (class~$\mathbf{G}$, formal by DGMS).

Two further distinctions are used: (v)~The CoHA of a CY$_3$ category is an \emph{associative algebra} (Schiffmann--Vasserot--Yang--Zhao multiplication), \emph{not} a chiral algebra. The connection CoHA $\to$ $\Eone$-chiral algebra is via the CY-to-chiral correspondence $\{\Phi_d\}$ (CY-A), not by identification. (vi)~The $\Omega$-background $(h_1, \ldots, h_{24})$ on $\C$ in $K3 \times \C \times \R$ is the intermediary mechanism producing $Y(\fg_{K3})$ from 5d holomorphic CS; the normal-bundle grading $N_{C/Y}$ does \emph{not} directly identify with the $(q,t)$~parameters. The labelled scalars are $\kappa_{\mathrm{ch}}^{\mathrm{Hodge}}(K3)=2$ and $\kappa_{\mathrm{ch}}^{\mathrm{Heis}}(K3\times E)=3$ (Remark~\ref{rem:beauville-kappa-formula-subscript-split}).
\end{remark}

\begin{remark}[Four pairs of distinct objects in the hCS vs.\ sigma model comparison]
\label{rem:hcs-sigma-conflations}
The hCS / sigma-model comparison involves four pairs of distinct objects connected by named arrows; the arrow direction and intermediary in each case fixes which side of the pair is in play.
\begin{enumerate}[label=(\alph*)]
 \item \emph{CoHA versus $\Eone$-chiral algebra.} The CoHA carries an associative (Hall) product; the $\Eone$-chiral algebra is a factorization algebra on $\C \times \R$. The connecting arrow is the CY-to-chiral correspondence $\{\Phi_d\}$.
 \item \emph{$N_{C/Y}$ grading versus $(q,t)$ parameters.} The intermediary mechanism (equivariant localisation, Nekrasov partition function, refinement) supplies the dictionary between the two.
 \item \emph{Spectral $z$ versus worldsheet $z$.} The spectral $z$ enters the shift in $\Delta_z$; the worldsheet $z$ enters the OPE pole structure. Setting spectral $z = 0$ removes the shift; sending worldsheet $z \to w$ produces the OPE poles.
 \item \emph{OPE-level rank trace $24$ versus CY $\kappa_{\mathrm{ch}}^{\mathrm{Hodge}} = 2$.} The two are different invariants of the same K3: OPE trace $\rho_{\mathrm{OPE}}=\tr(\delta^{ab}) = 24$ counts free-field modes, CY supertrace $\chi(\cO_{K3}) = 2$ reads the Hodge data. The subscript names which.
\end{enumerate}
\end{remark}

%% RTT-OPE dictionary for the K3 Yangian

\section{The RTT--OPE dictionary for the K3 Yangian}
\label{sec:rtt-ope-dictionary}

The (conjectural) K3 Yangian $Y(\fg_{K3})$ for $\fg = \gl_1$ admits two
compatible presentation surfaces at the Heisenberg mode-algebra shadow,
not a proved equivalence of full chiral-Yangian definitions. The RTT
presentation below is weaker than the full chiral quadruple with modular
MC tower; the dictionary records the data visible at this abelian mode
level:
\begin{enumerate}[label=(\Alph*)]
 \item \textbf{RTT presentation}: generators are modes of the transfer matrix $T(u) = 1 + \sum_{s \geq 1} \psi_s u^{-s}$, with the defining relation
 \[
 R_{12}(u - v)\, T_1(u)\, T_2(v) \;=\; T_2(v)\, T_1(u)\, R_{12}(u - v).
 \]
 \item \textbf{OPE presentation}: generators are $24$ Heisenberg currents $J_i(z) = \sum_n J_{i,n}\, z^{-n-1}$, $i = 0, \ldots, 23$, with the OPE
 \[
 J_i(z)\, J_j(w) \;\sim\; \frac{\omega^{ij}}{(z - w)^2},
 \]
 where $\omega^{ij}$ is the Mukai pairing of signature~$(4, 20)$.
\end{enumerate}

\begin{remark}[Spectral parameter vs.\ worldsheet coordinate]
\label{rem:spectral-vs-worldsheet}
The variable $u$ in the RTT presentation is a \emph{Yangian spectral parameter}, living in the complexified weight space. The variable $z$ in the OPE presentation is a \emph{worldsheet coordinate} on the chiral Riemann surface. These are different mathematical objects:
\begin{itemize}
 \item Setting $u = 0$ in $\Delta_u$ yields the cocommutative coproduct $\Delta_0(T(v)) = T^L(v) \cdot T^R(v)$: no singularity.
 \item Setting $z \to w$ in $J_i(z)\,J_j(w)$ produces the OPE poles $\sim (z-w)^{-2}$, a physical collision singularity.
\end{itemize}
Both encode the \emph{same mode algebra} $[J_{i,m},\, J_{j,n}] = \omega^{ij}\, m\, \delta_{m+n,0}$; the bridge is through the modes.
\end{remark}

\subsection{The 24 Heisenberg currents and the Sugawara construction}
\label{subsec:k3-heisenberg-ope}

\begin{proposition}[Mukai Heisenberg OPE]
\label{prop:mukai-heisenberg-ope}
The $24$ currents $J_i(z)$ associated to the Mukai lattice $\widetilde{\Lambda}_{K3} = U^4 \oplus E_8(-1)^2$ satisfy the Heisenberg OPE
\begin{equation}\label{eq:mukai-heisenberg-ope}
 J_i(z)\, J_j(w) \;\sim\; \frac{\omega^{ij}}{(z - w)^2},
 \qquad i, j = 0, \ldots, 23,
\end{equation}
where $\omega^{ij} = \mathrm{diag}(\underbrace{+1, \ldots, +1}_{4},\, \underbrace{-1, \ldots, -1}_{20})$ in the diagonal basis. The corresponding $\lambda$-bracket is
\begin{equation}\label{eq:mukai-lambda-bracket}
 \{J_{i\,\lambda}\, J_j\} \;=\; \omega^{ij}\, \lambda.
\end{equation}
\end{proposition}

\begin{proof}
This is the standard Heisenberg OPE for $24$ free bosons with the Mukai pairing, as constructed in Section~\ref{subsec:k3-fact-tree} of Chapter~\ref{ch:k3-times-e}. The $\lambda$-bracket is the Fourier transform of the OPE: $\{a_\lambda\, b\} = \mathrm{Res}_{z=0}\, e^{\lambda z}\, a(z)\, b(0)$, giving $\mathrm{Res}_{z=0}\, e^{\lambda z}\, \omega^{ij}/z^2 = \omega^{ij}\, \lambda$.
\end{proof}

\begin{proposition}[Sugawara stress tensor at \texorpdfstring{$c = 24$}{c = 24}]
\label{prop:sugawara-c24}
The Sugawara construction
\begin{equation}\label{eq:sugawara-mukai}
 T(z) \;=\; \frac{1}{2}\sum_{i,j=0}^{23} \omega_{ij}\, {:}J_i(z)\, J_j(z){:}
\end{equation}
produces a Virasoro field with central charge $c = 24 = \kappa_{\mathrm{fiber}}$. Each free boson contributes $c_i = 1$ independent of the sign of $\omega^{ii}$, giving $c = \sum_{i=0}^{23} 1 = 24$. The self-OPE is
\begin{equation}\label{eq:tt-ope-c24}
 T(z)\, T(w) \;\sim\; \frac{12}{(z - w)^4} \;+\; \frac{2\, T(w)}{(z - w)^2} \;+\; \frac{\partial T(w)}{z - w}.
\end{equation}
\end{proposition}

\begin{proof}
The central charge of a rank-$1$ Heisenberg algebra at level $k$ ($[J_m, J_n] = k\, m\, \delta_{m+n,0}$) with Sugawara $T_n = (1/(2k)) \sum_a {:}J_a\, J_{n-a}{:}$ is $c = 1$ for any nonzero $k$ (the factor $k$ in the commutator cancels with the $1/k$ normalisation). For the rank-$24$ Mukai Heisenberg with $k_i = \omega^{ii} = \pm 1$, the total is $c = 24$. The fourth-order pole coefficient $c/2 = 12$ in~\eqref{eq:tt-ope-c24} confirms this.
\end{proof}

\begin{remark}[\texorpdfstring{$\kappa_\bullet$}{kappa-.}-spectrum]
The central charge $c = 24$ is $\kappa_{\mathrm{fiber}}$ (the Mukai lattice rank on the K3 fibre), distinct from $\kappa_{\mathrm{ch}}^{\mathrm{Hodge}}(K3) = 2$ (the K3 chiral algebraisation), $\kappa_{\mathrm{ch}}^{\mathrm{Heis}}(K3 \times E) = 3$ (Beauville-additive Heisenberg--Cartan rank: $2 + 1$), $\kappa_{\mathrm{BKM}}(\Delta_5) = 5$ (the Igusa weight), and $\kappa_{\mathrm{cat}}(K3) = \chi(\cO_{K3}) = 2$ (K3-fibre per-fibre value; $\kappa_{\mathrm{cat}}(K3 \times E) = 0$ by K\"unneth on the total space). The five $K3 \times E$ tower values are $\{\kappa_{\mathrm{cat}}, \kappa_{\mathrm{ch}}^{\mathrm{Hodge}}, \kappa_{\mathrm{ch}}^{\mathrm{Heis}}, \kappa_{\mathrm{BKM}}(\Delta_5), \kappa_{\mathrm{fiber}}\} = \{0, 0, 3, 5, 24\}$, with the first two reading the total-space Hodge layer.
\end{remark}

\subsection{The RTT--OPE bridge: classical \texorpdfstring{$r$}{r}-matrix = OPE coefficient}
\label{subsec:rtt-ope-bridge}

\begin{conjecture}[RTT--OPE dictionary for \texorpdfstring{$Y(\fg_{K3})$}{Y(g-K3)} at \texorpdfstring{$\gl_1$}{gl-1}]
\label{conj:rtt-ope-dictionary}
\ClaimStatusConjectured
The K3 Yangian $Y(\fg_{K3})$ for $\fg = \gl_1$ (Conjecture~\ref{conj:k3-yangian}) admits an RTT presentation with $R$-matrix
\begin{equation}\label{eq:rk3-diagonal}
 R_{K3}(u) \;=\; \mathrm{diag}\!\left(\frac{u - h_1}{u + h_1},\; \ldots,\; \frac{u - h_{24}}{u + h_{24}}\right),
\end{equation}
a degree-$(24,24)$ rational function with $24$ parameters $h_i$ satisfying $\sum h_i = 0$ (CY$_2$ constraint). The classical limit as $u \to \infty$ gives:
\begin{equation}\label{eq:classical-r-mukai}
 R_{K3}(u) \;=\; I_{24} \;-\; \frac{2}{u}\, \mathrm{diag}(h_1, \ldots, h_{24}) \;+\; O(u^{-2}).
\end{equation}
At unit deformation $h_i = \omega^{ii}$, the classical $r$-matrix is $r = -2\, \omega^{-1}$ (the inverse Mukai pairing, up to normalisation). This $r$-matrix is dual to the OPE second-order pole: the coefficient $\omega^{ij}$ in~\eqref{eq:mukai-heisenberg-ope} equals the leading $1/u$ coefficient of $R_{K3}(u)$.
\end{conjecture}

The RTT--OPE dictionary is summarised in Table~\ref{tab:rtt-ope}.

\begin{table}[ht]
\centering
\small
\renewcommand{\arraystretch}{1.3}
\begin{tabular}{lll}
\toprule
\textbf{Object} & \textbf{RTT side} & \textbf{OPE side} \\
\midrule
Generators & $T(u) = 1 + \sum \psi_s\, u^{-s}$ & $J_i(z) = \sum J_{i,n}\, z^{-n-1}$, $i=0,\ldots,23$ \\
Relations & $R_{12}(u{-}v)\, T_1 T_2 = T_2 T_1\, R_{12}$ & $J_i(z)\, J_j(w) \sim \omega^{ij}/(z{-}w)^2$ \\
$R$-matrix & $R(u) = \mathrm{diag}(\frac{u-h_i}{u+h_i})$ & OPE coefficient $\omega^{ij}$ \\
Stress tensor & Transfer matrix $T(u)$ & $T(z) = \tfrac{1}{2}\sum \omega_{ij}\,{:}J_i J_j{:}$, $c = 24$ \\
Coproduct & $\Delta_u(T(v)) = T^L(v) \cdot T^R(v{-}u)$ & $E_1$-factorisation on $\C \times \R$ \\
Variable & $u$ = spectral (weight space) & $z$ = worldsheet (Riemann surface) \\
Lie conformal & $Y(\fg_{K3})$ via RTT & $U^{\mathrm{ch}}(L_{\mathrm{Muk}}) = V_{H_{\mathrm{Muk}}}$ \\
Unitarity & $R(u)\, R(-u) = I$ & $J_i$ self-conjugate \\
Central charge & Encoded in quantum determinant & $c = 24$ from $(z{-}w)^{-4}$ pole \\
\bottomrule
\end{tabular}
\caption{RTT--OPE dictionary for the K3 Yangian at $\gl_1$. The bridge between the two presentations passes through the common mode algebra $[J_{i,m},\, J_{j,n}] = \omega^{ij}\, m\, \delta_{m+n,0}$.}
\label{tab:rtt-ope}
\end{table}

\subsection{The Mukai Lie conformal algebra and chiral envelope}
\label{subsec:mukai-lie-conformal}

\begin{proposition}[Chiral envelope of \texorpdfstring{$L_{\mathrm{Muk}}$}{L-Muk}]
\label{prop:chiral-envelope-mukai}
Let $L_{\mathrm{Muk}} = \bC[\partial] \otimes \bC^{24}$ be the rank-$24$ Lie conformal algebra with $\lambda$-bracket~\eqref{eq:mukai-lambda-bracket}. Then:
\begin{enumerate}[label=\textup{(\roman*)}]
 \item $L_{\mathrm{Muk}}$ is abelian (no structure constants beyond the pairing). Shadow class~$\mathbf{G}$ (free field, depth~$2$).
 \item The chiral (factorisation) envelope equals the rank-$24$ Heisenberg VOA:
 \begin{equation}\label{eq:chiral-envelope-mukai}
 U^{\mathrm{ch}}(L_{\mathrm{Muk}}) \;\simeq\; V_{H_{\mathrm{Muk}}}.
 \end{equation}
 \item Generator count: $24$ strong generators at conformal weight~$1$.
 \item Character:
 $\mathrm{ch}(V_{H_{\mathrm{Muk}}}) = q^{-1} \prod_{n \geq 1} 1/(1 - q^n)^{24} = q^{-1}/\eta(q)^{24}$, with $q^1$ coefficient $24$ and $q^2$ coefficient $324$.
 \item $V_{H_{\mathrm{Muk}}}$ is the \emph{classical limit} ($\sigma_3 = 0$) of $Y(\fg_{K3})$: the undeformed Heisenberg sector before the Yangian deformation is turned on.
\end{enumerate}
\end{proposition}

\begin{proof}
(i)~The $\lambda$-bracket $\{J_{i\,\lambda}\, J_j\} = \omega^{ij}\, \lambda$ is linear in $\lambda$ (no $\lambda^0$ or higher terms), so $L_{\mathrm{Muk}}$ is the Lie conformal algebra of a $2$-step nilpotent Lie algebra (class~$\mathbf{G}$). (ii)~For an abelian Lie conformal algebra, the factorisation envelope adds no new generators or relations beyond the PBW basis; the result is the free-field (Heisenberg) vertex algebra. (iii)~Generator count: $24 = \mathrm{rank}(\widetilde{\Lambda}_{K3})$. (iv)~The PBW basis $J_{i_1,-n_1} \cdots J_{i_k,-n_k} |0\rangle$ ($n_j > 0$, $24$ colours) gives the Fock space character $\prod_{n \geq 1} 1/(1-q^n)^{24}$. At $q^1$: $24$ oscillators. At $q^2$: $\binom{25}{2} + 24 = 300 + 24 = 324$. (v)~At $\sigma_3 = 0$ the structure function $g_{K3}(z) = 1$ identically (all factors cancel), the $R$-matrix is trivial, and $Y(\fg_{K3})$ degenerates to $H_{\mathrm{Muk}}$.
\end{proof}

\subsection{The coproduct in OPE language}
\label{subsec:coproduct-ope-language}

\begin{proposition}[Primitivity of the Heisenberg coproduct]
\label{prop:heisenberg-primitive}
The Heisenberg generators $J_{i,n}$ are \emph{primitive} in the $E_1$-chiral bialgebra:
\begin{equation}\label{eq:heisenberg-primitive}
 \Delta_0(J_{i,n}) \;=\; J_{i,n}^L \;+\; J_{i,n}^R.
\end{equation}
The coproduct preserves the Heisenberg commutator on the tensor product:
\[
 [\Delta_0(J_{i,m}),\, \Delta_0(J_{j,n})]
 \;=\; 2\, \omega^{ij}\, m\, \delta_{m+n,0}
\]
(the factor~$2$ arises from the two independent copies $H_L$ and $H_R$). The spectral shift $u$ appears only at the \emph{composite} level: the Sugawara coproduct $\Delta_u(T(v)) = T^L(v) \cdot T^R(v - u)$ acquires $u$-dependent cross-terms via the Miura factorisation (Section~\ref{sec:e1-chiral-bialgebras}).
\end{proposition}

\begin{proof}
Direct Fock space computation on the tensor product $\cF_L \otimes \cF_R$. The cross-commutator $[J_{i,m}^L, J_{j,n}^R] = 0$ (operators on different tensor factors), so $[\Delta_0(J_{i,m}), \Delta_0(J_{j,n})] = [J_{i,m}^L, J_{j,n}^L] + [J_{i,m}^R, J_{j,n}^R] = 2\, \omega^{ij}\, m\, \delta_{m+n,0}$.
\end{proof}

\begin{remark}[RTT--OPE dictionary: verification summary]
\label{rem:rtt-ope-summary}
The RTT--OPE dictionary is governed by the following nine identities:
\begin{center}
\small
\renewcommand{\arraystretch}{1.15}
\begin{tabular}{llr}
\toprule
\textbf{Block} & \textbf{Content} & \textbf{Identities} \\
\midrule
Heisenberg OPE & $\omega^{ij}$ coefficients, rank, signature, pole order & 7 \\
Sugawara & $c = 24$, conformal weights, $\kappa_\bullet$-spectrum & 8 \\
Fock space Virasoro & $c = r$ at ranks $1$--$4$, positive/negative pairings & 6 \\
RTT--OPE bridge & Classical $r$-matrix, $R(u) = I + \omega^{-1}/u$, proportionality & 7 \\
Spectral vs worldsheet & distinction, $\Delta_0$ no singularity & 4 \\
Lie conformal envelope & $L_{\mathrm{Muk}}$ abelian, $U^{\mathrm{ch}} = V_{H_{\mathrm{Muk}}}$, character & 8 \\
Coproduct in OPE & Primitivity $\Delta_0(J) = J^L + J^R$, commutator & 6 \\
Dictionary consistency & $9$ entries, cross-references & 4 \\
RTT--OPE domain & Conjectural bridge domain & 2 \\
\bottomrule
\end{tabular}
\end{center}
\end{remark}


\section{The CFG25 \texorpdfstring{$E_1$}{E-1}-chiral lift: from 3d CS to 6d holomorphic theory}
\label{sec:cfg25-e1-chiral-lift}

Costello--Francis--Gwilliam (2025) establish five structural results for 3d Chern--Simons theory on $\R^3$ using factorization homology. Each result has a precise $E_1$-chiral analogue in the 6d holomorphic theory on $\C^3$ of the Costello construction. This section makes each analogue explicit and identifies what works, what breaks, and what requires genuinely new ideas.

\begin{center}
\renewcommand{\arraystretch}{1.3}
\small
\begin{tabular}{lp{4.5cm}p{5.5cm}l}
 \toprule
 \textbf{CFG25 result (3d on $\R^3$)} & \textbf{$E_1$-chiral lift (6d on $\C^3$)} & \textbf{Obstruction / new idea} & \textbf{Construction locus} \\
 \midrule
 1. $\mathrm{Obs}^q(\R^3)$ is $E_3$-fact. & $E_3$-chiral fact.\ on $\C^3$ (holomorphic refinement $E_6 \to E_3$) & Omega-background provides nontrivial braiding & works \\
 2. Boundary $= V_k(\frakg)$ & Boundary $= Y(\widehat{\fgl}_1)$ & $E_\infty \to E_1$: Deligne level drops $E_3 \to E_2$ & works \\
 3. Link invariants & Holomorphic submanifold invariants & $\pi_1(\text{complement}) \to$ vanishing cycles; normal bundle $N_{\Sigma/\C^3}$ & partial \\
 4. $\int_{M \setminus L} = $ QG invariant & $\int_{\C^3 \setminus \Sigma} = $ quantum toroidal invariant & Holomorphic dependence; Ran space required for compact CY$_3$ & requires compact Ran construction \\
 5. $\mathrm{Obs}^q(\R^3)^! = U_q(\frakg)\text{-mod}$ & $\mathrm{Obs}^q(\C^3)^! = U_{q,t}(\widehat{\widehat{\fgl}}_1)\text{-mod}$ & ZTE failure (Theorem~\ref{thm:zte-failure}); even-sector ternary correction & requires full charge-preserving $E_3$ complex \\
 \bottomrule
\end{tabular}
\end{center}

\subsection{Result 1: \texorpdfstring{$E_n$}{E-n}-factorization structure}
\label{subsec:cfg25-result1}

\begin{proposition}[Holomorphic refinement of the \texorpdfstring{$E_n$}{E-n} level]
\label{prop:holomorphic-en-refinement}
The observables of holomorphic Chern--Simons theory on an $n$-dimensional complex manifold $M$ carry $E_{2n}$-factorization structure topologically (from $\dim_\R M = 2n$) and $E_n$-chiral factorization structure holomorphically (each complex direction contributes one chiral level). The relationship:
\begin{center}
\renewcommand{\arraystretch}{1.2}
\small
\begin{tabular}{llll}
 \toprule
 \textbf{Theory} & \textbf{Manifold} & \textbf{Topological $\En$} & \textbf{Holomorphic $\En$} \\
 \midrule
 CFG25 (3d CS) & $\R^3$ & $E_3$ & (not holomorphic) \\
 5d CS & $\C^2 \times \R$ & $E_4$ & $\Etwo$-chiral on $\C^2$ \\
 6d theory & $\C^3$ & $E_6$ & $E_3$-chiral on $\C^3$ \\
 \bottomrule
\end{tabular}
\end{center}
\noindent The $E_3$-chiral structure on $\C^3$ is the direct lift of the $E_3$ structure on $\R^3$. The holomorphic refinement is not a weakening: it is a \emph{different} $E_3$ structure, one that depends on the Omega-background parameters $(h_1, h_2, h_3)$ with $h_1 + h_2 + h_3 = 0$ for the nontrivial braiding (the topological braiding on $\Conf_2(\R^6)$ is trivial since $\pi_1(\Conf_2(\R^6)) \cong \pi_1(S^5) = 0$).
\end{proposition}

\begin{proof}
Costello--Gwilliam (2021), Costello--Francis--Gwilliam (2025). The holomorphic refinement is Proposition~\ref{prop:holomorphic-e1}; the $E_3$ on $\C^3$ is Conjecture~\ref{conj:topological-e3-comparison}(a).
\end{proof}

\subsection{Result 2: the boundary algebra upgrade}
\label{subsec:cfg25-result2}

The CFG25 boundary algebra is the Kac--Moody VOA $V_k(\frakg)$, which is $E_\infty$-chiral (a commutative vertex algebra). In the 6d lift, the boundary algebra becomes the affine Yangian $Y(\widehat{\fgl}_1)$, which is $E_1$-chiral (non-commutative, ordered).

\begin{proposition}[Deligne conjecture level drop in the CFG25 lift]
\label{prop:deligne-level-drop-cfg25}
The higher Deligne conjecture (Lurie, \emph{Higher Algebra}, Theorem~5.3.1.30) assigns to the Hochschild cochains $\HH^\bullet(A, A)$ of an $E_n$-algebra $A$ an $E_{n+1}$-algebra structure. In the CFG25 lift:
\begin{enumerate}[label=\textup{(\roman*)}]
 \item \emph{3d (CFG25)}: The boundary $V_k(\frakg)$ is $E_\infty$-chiral. Its Hochschild cochains carry $E_\infty$ structure (and in particular $E_3$, generated by cup product, Gerstenhaber bracket, and BV operator). The BV operator requires $E_\infty$ input.
 \item \emph{6d (lift)}: The boundary $Y(\widehat{\fgl}_1)$ is $E_1$-chiral. By Dunn additivity ($E_1 \otimes_{E_0} E_1 = E_2$), $\HH^\bullet(Y, Y)$ carries only $E_2$-algebra structure. No BV operator is available.
\end{enumerate}
The Deligne conjecture level drops from $E_3$ to $E_2$ in the lift. This is a genuine structural consequence: the commutativity of $V_k(\frakg)$ is \emph{lost} in passing to the Yangian.
\end{proposition}

\begin{proof}
Part~(i) is the higher Deligne conjecture applied to $E_\infty$ input (Lurie, \emph{loc.\ cit.}). Part~(ii) follows from Dunn additivity: $\HH^\bullet$ of an $E_n$-algebra is $E_{n+1}$, so $\HH^\bullet$ of $E_1$ is $E_2$. The BV operator arises from the $S^1$-action on the cyclic bar complex, which requires the commutativity of $A$.
\end{proof}

The passage from $E_1$ to $E_2$ is achieved not by the Deligne conjecture but by the \emph{Drinfeld center} construction:
\begin{equation}\label{eq:drinfeld-e1-to-e2-cfg25}
 \cZ(\Rep^{E_1}(Y^+(\widehat{\fgl}_1))) \;\simeq\; \Rep^{E_2}(Y(\widehat{\fgl}_1)),
\end{equation}
as proved in Theorem~\ref{thm:drinfeld-center-coha}. The full Yangian $Y(\widehat{\fgl}_1) = Y^+ \bowtie (Y^+)^{*\mathrm{cop}}$ is the Drinfeld double of the positive half $Y^+$ (the CoHA; note the CoHA itself is associative, not chiral).

\subsection{Result 3: submanifold invariants replace link invariants}
\label{subsec:cfg25-result3}

\begin{conjecture}[Holomorphic submanifold invariants from 6d factorization homology]
\label{conj:cfg25-submanifold-invariants}
\ClaimStatusConjectured
Let $\Sigma \subset \C^3$ be a smooth holomorphic submanifold of complex dimension $k$ ($k = 1$ for curves, $k = 2$ for surfaces). The factorization homology
\[
 \int_{\C^3 \setminus \Sigma} \cF
\]
of the $E_3$-chiral factorization algebra $\cF$ on the complement $\C^3 \setminus \Sigma$ produces an invariant of the pair $(\C^3, \Sigma)$ that depends on the holomorphic data of $\Sigma$, including its complex structure and topology. For codimension-$2$ curves ($k = 1$), the defect algebra on $\Sigma$ is $W_{1+\infty}$ at level $\Psi = -\sigma_2$ (Proposition~\ref{prop:codim2-defect-ope}).

The lift from 3d to 6d changes the invariant theory in three ways:
\begin{enumerate}[label=\textup{(\alph*)}]
 \item \emph{Classification}: links in $\R^3$ are classified by isotopy (finite combinatorial data); holomorphic submanifolds in $\C^3$ are classified by algebraic geometry (moduli spaces of potentially infinite dimension).
 \item \emph{Fundamental group}: $\pi_1(\R^3 \setminus L)$ is the knot group (nontrivial for every nontrivial link); $\pi_1(\C^3 \setminus \Sigma)$ may be trivial for smooth holomorphic curves. The invariant data migrates from $\pi_1$ (knot group representations) to the \emph{normal bundle} $N_{\Sigma / \C^3}$.
 \item \emph{Normal bundle}: for a curve $C \subset \C^3$ of genus $g$, the CY adjunction formula (trivial canonical bundle of $\C^3$) gives $\det(N_{C/\C^3}) = K_C^{-1} = \cO(2-2g)$. The normal bundle replaces the framing data of the link.
\end{enumerate}
\end{conjecture}

\begin{remark}[Why \texorpdfstring{$\pi_1$}{pi-1} migrates to the normal bundle]
\label{rem:pi1-to-normal-bundle}
In 3d CS, the Wilson line observable $W_\rho(K) = \mathrm{tr}_\rho(\mathrm{Hol}_K(A))$ is the trace in a representation $\rho$ of the holonomy around the knot $K$. The holonomy is a $\pi_1$-representation. In the 6d lift, the holonomy around a codimension-$2$ submanifold $\Sigma$ is trivial when $\pi_1(\C^3 \setminus \Sigma) = 0$, but the normal bundle $N_{\Sigma/\C^3}$ provides a \emph{replacement}: the equivariant parameters $(h_1, h_2)$ of the normal $\C^2$-fiber give the spectral parameters of the Yangian/toroidal $R$-matrix. The connection to quantum group parameters is through the Omega-background equivariant localisation on the normal bundle.
\end{remark}

\subsection{Result 4: factorization homology and quantum toroidal invariants}
\label{subsec:cfg25-result4}

\begin{conjecture}[Factorization homology on \texorpdfstring{$\C^3$}{C-3} complements]
\label{conj:cfg25-fact-hom-complement}
\ClaimStatusConjectured
For a holomorphic submanifold $\Sigma \subset \C^3$ in the toric CY$_3$ setting, the factorization homology $\int_{\C^3 \setminus \Sigma} \cF$ computes quantum toroidal $U_{q,t}(\widehat{\widehat{\fgl}}_1)$-module data. The identification with the Nekrasov partition function
\begin{equation}\label{eq:nekrasov-fact-hom-cfg25}
 Z_{\mathrm{Nek}}(\varepsilon_1, \varepsilon_2; q) \;=\; \int_{\C^2} \cF\big|_{\C^2}
\end{equation}
provides the computational bridge: $Z_{\mathrm{Nek}}^{U(1)} = \prod_{n \geq 1} 1/(1-q^n)$ with plethystic logarithm $q/(1-q)$ (one bosonic BPS mode per instanton number). This is the $\Etwo$-projection of the $E_3$ integral over $\C^3$.

For compact CY$_3$ (quintic, $K3 \times E$): the holomorphic factorization homology integral requires the Ran-space formulation of chiral homology, and the available CY-A$_3$ input is the framed object-level theorem under H1--H4 and fixed admissible specialisation \textup{(}Theorem~\ref{thm:cy-to-chiral-d3}\textup{)}. The 6d factorization homology route does \emph{not} bypass that apparatus: each subproblem (local $E_3$ algebra for compact targets, handle decomposition, VOA identification of output) requires the same chain-level data that CY-A$_3$ assumes or verifies on toric/formal loci. Explicit chain-level convergence and framing for non-formal compact CY$_3$ remain residual analytic content.
\end{conjecture}

\subsection{Result 5: Koszul duality and the ZTE obstruction}
\label{subsec:cfg25-result5}

The CFG25 Koszul duality $\mathrm{Obs}^q(\R^3)^! \simeq U_q(\frakg)\text{-mod}$ is \emph{proved}. The 6d lift to $\mathrm{Obs}^q(\C^3)^! \simeq U_{q,t}(\widehat{\widehat{\fgl}}_1)\text{-mod}$ is \emph{conjectural} (Conjecture~\ref{conj:e3-koszul-duality}).

\begin{proposition}[Parameter inversion under \texorpdfstring{$E_3$}{E-3}-chiral Koszul duality]
\label{prop:parameter-inversion-cfg25}
Under the Verdier duality on $\C^3$ factorization coalgebras, the Omega-background parameters invert: $(h_1, h_2, h_3) \mapsto (-h_1, -h_2, -h_3)$. On the elementary symmetric polynomials:
\begin{enumerate}[label=\textup{(\roman*)}]
 \item $\sigma_2(-\boldsymbol{h}) = \sigma_2(\boldsymbol{h})$: the level $\Psi = -\sigma_2$ is \emph{preserved} by Koszul duality (even function).
 \item $\sigma_3(-\boldsymbol{h}) = -\sigma_3(\boldsymbol{h})$: the deformation parameter $\sigma_3 = h_1 h_2 h_3$ is \emph{negated} (odd function).
 \item In the $(q, t)$ language: $q = e^{h_1} \mapsto e^{-h_1} = q^{-1}$, $t = e^{-h_2} \mapsto e^{h_2} = t^{-1}$.
\end{enumerate}
The Koszul conductor for the free-field (KM/$\fgl_1$) case is $\rho_K = 0$: $\kappa_{\mathrm{ch}}(A) + \kappa_{\mathrm{ch}}(A^!) = (-\sigma_2) + \sigma_2 = 0 = \rho_K$.
\end{proposition}

\begin{proof}
Direct computation. Let $f(\boldsymbol{h}) = e_k(h_1, h_2, h_3)$ be the $k$-th elementary symmetric polynomial. Under $h_i \mapsto -h_i$: $e_k(-\boldsymbol{h}) = (-1)^k e_k(\boldsymbol{h})$. So $e_2(-\boldsymbol{h}) = e_2(\boldsymbol{h})$ (even) and $e_3(-\boldsymbol{h}) = -e_3(\boldsymbol{h})$ (odd). For the modular characteristic, write $\kappa_{\mathrm{ch}}^{\mathrm{KM}} = -\sigma_2$; then $\kappa_{\mathrm{ch}}^{\mathrm{KM}}(-\boldsymbol{h}) = -\sigma_2(-\boldsymbol{h}) = -\sigma_2(\boldsymbol{h}) = \kappa_{\mathrm{ch}}^{\mathrm{KM}}(\boldsymbol{h})$. The Koszul dual has $\kappa_{\mathrm{ch}}^{!,\mathrm{KM}} = \sigma_2$, so $\kappa_{\mathrm{ch}}^{\mathrm{KM}} + \kappa_{\mathrm{ch}}^{!,\mathrm{KM}} = -\sigma_2 + \sigma_2 = 0 = \rho_K$ for the KM/$\fgl_1$ family (class~$\mathbf{G}$, $\rho_K = 0$). For other families the Koszul conductor differs: $\rho_K = 13$ for Virasoro, $\rho_K = 24$ for lattice rank (Vol~I Theorem~C).
\end{proof}

The key obstruction to proving the 6d Koszul duality is the \emph{Zamolodchikov tetrahedron equation} (ZTE). In 3d, the quantum Yang--Baxter equation (YBE) governs the $E_2$ coherence of the $R$-matrix, and the CFG25 result uses the YBE solution to build the Koszul dual. In 6d, the $E_3$ coherence requires the ZTE (the $3$-simplex analogue of the YBE $2$-simplex relation). Theorem~\ref{thm:zte-failure} proves that the na\"ive factorisation $S_{ijk} = R_{ij} R_{ik} R_{jk}$ from pairwise YBE solutions does \emph{not} satisfy the ZTE: the obstruction is $O(\sigma_3^2)$ and antisymmetric. Proposition~\ref{prop:zte-explicit-correction} kills the charge-$2$ obstruction by a genuinely ternary correction. Proposition~\ref{prop:zte-o-kappa4} shows that the odd charge sectors still carry rank-$2$ residues. Thus the local even-sector obstruction is solved, but the full charge-preserving all-order ZTE is not yet constructed.

\subsection{The obstruction hierarchy}
\label{subsec:cfg25-obstruction-hierarchy}

The three obstructions to the CFG25 lift form a partial order by severity:

\begin{center}
\renewcommand{\arraystretch}{1.3}
\small
\begin{tabular}{cllp{6cm}}
 \toprule
 \textbf{Level} & \textbf{Obstruction} & \textbf{Type} & \textbf{Required datum} \\
 \midrule
 1 & ZTE failure & Local algebraic & charge-$2$ obstruction killed by a ternary correction; full charge-preserving completion open \\
 2 & CY-A$_3$ ($\bS^3$-framing) & Global geometric & requires the corrected TCFT comparison or the complete $\HH^{-2}_{E_1}$ filtration theorem (Theorem~\ref{thm:derived-framing-m3b2}); chain-level explicit framing open for non-formal \\
 3 & 6d non-Lagrangian & Foundational & construct route~(c) of Remark~\ref{rem:6d-not-lagrangian} \\
 \bottomrule
\end{tabular}
\end{center}

\noindent Each later obstruction depends on resolving the earlier ones. The charge-$2$ ZTE obstruction is algebraic and has a ternary resolution; the odd-sector and all-order ZTE completion remain part of the local problem. CY-A$_3$ is controlled only after the corrected TCFT carrier or the complete $\HH^{-2}_{E_1}$ filtration comparison is supplied, and chain-level explicit framing data for non-formal algebras remains open; the non-Lagrangian nature of the 6d theory is the foundational question of whether the algebraic formulation of Costello--Francis--Gwilliam extends to the 6d setting.


% ============================================================
% Perturbative chiral invariants from 6d hCS Feynman diagrams
% ============================================================

\section{Perturbative chiral invariants from Feynman diagrams}
\label{sec:perturbative-chiral-feynman}

The perturbative expansion of the local hCS model is expected to see the
shadow tower of $\cW_{1+\infty}$ through the identification of loop order
with shadow depth. The proved content available here is finite:
the Virasoro channel has been checked through $L\leq 3$ against the
shadow recursion and the $\Ainf$ operations. The all-loop
Feynman/shadow dictionary remains a conjectural extension until the
regularised hCS graph integrals and their normalisation are constructed
uniformly.

\subsection{The perturbative expansion of 6d hCS}
\label{subsec:perturbative-expansion-6d}

\begin{proposition}[Perturbative chiral expansion]
\label{prop:perturbative-chiral-expansion}
For the local abelian hCS model on $\C^3_{h_1,h_2,h_3}$ restricted to
the codimension-$2$ defect $C \subset \C^3$, after a choice of
regularisation and effective Omega-background vertices, the perturbative
expansion has the formal graph shape
\begin{equation}\label{eq:perturbative-expansion}
 Z_{\mathrm{chiral}}(A_C) \;=\; \sum_{L \geq 0} \sigma_3^L \; Z^{(L)},
 \qquad
 Z^{(L)} \;=\; \sum_{\Gamma:\, b_1(\Gamma) = L}
 \frac{1}{|\!\Aut(\Gamma)|}\, w(\Gamma)\, I(\Gamma),
\end{equation}
where the sum ranges over connected Feynman graphs~$\Gamma$ at loop order~$L = b_1(\Gamma)$ (first Betti number), $w(\Gamma)$ is the Lie algebra weight from the Yangian structure constants, and $I(\Gamma)$ is the regularised configuration-space integral over $\Conf_{|V(\Gamma)|}(C)$.

For $\fgl_1$ (abelian gauge group), the cubic vertex of the holomorphic CS action vanishes, and the effective vertices arise from the Omega-background mode couplings at each spin level.
\end{proposition}

\begin{proof}
The displayed formula is the standard graph expansion one obtains after
choosing a propagator, regularisation, and effective vertices. For
$\fgl_1$, the cubic term in the bare hCS action vanishes identically, so
the nontrivial vertices in this witness come from the Omega-background
deformation rather than from a nonabelian cubic bracket. The finite graph
normalisations below are not an all-loop construction of renormalised
hCS observables.
\end{proof}

\subsection{Standard graphs and their contributions}
\label{subsec:standard-graphs}

The standard Feynman graphs at each loop order, for the spin-$s$ channel of the $\cW_{1+\infty}$ algebra at $c = 1$, are:

\begin{enumerate}[label=\textup{(\roman*)}]
 \item \emph{Tree level ($L = 0$): propagator exchange.} Two external vertices connected by one edge. The graph contributes $Z^{(0)}_s = \kappa_{\mathrm{ch},s} = c/s$, the modular characteristic of the spin-$s$ channel.

 \item \emph{One loop ($L = 1$): bubble diagram.} Two vertices connected by two parallel edges. Symmetry factor $1/|\!\Aut| = 1/2$. The regularised integral gives $Z^{(1)}_s = S_2^{(s)} = \kappa_{\mathrm{ch},s}$ (the shadow depth-$2$ invariant equals $\kappa_{\mathrm{ch}}$; this is the content of Theorem~A at genus~$1$).

 \item \emph{Two loops ($L = 2$): theta graph.} Three vertices forming a triangle ($3$ edges, $b_1 = 1$). This graph carries the cubic shadow $S_3 = \alpha$. For spin~$2$ (Virasoro at $c = 1$): $\alpha = 2$, arising from the $\Ainf$-operation $m_3(T, T, T) = -2T$. For spin~$1$ (Heisenberg, class~$\mathbf{G}$): $\alpha = 0$.

 \item \emph{Three loops ($L = 3$).} Multiple graph topologies contribute to $S_4$ (the quartic contact term). For spin~$2$ at $c = 1$: $S_4 = 10/27$, consistent with $m_4(T, T, T, T) = \tfrac{40}{27}\, T$.
\end{enumerate}

\subsection{The shadow--Feynman dictionary}
\label{subsec:shadow-feynman-dictionary}

\begin{proposition}[Finite loop/shadow witness]
\label{prop:loop-shadow-identification}
For the Virasoro spin-$2$ channel with the graph normalisation of
Proposition~\ref{prop:perturbative-chiral-expansion}, the finite
checked dictionary is
\begin{equation}\label{eq:shadow-feynman-dictionary}
 Z^{(L)} \;=\; S_{L+1}, \qquad L = 0, 1, 2, 3.
\end{equation}
Explicitly, for the Virasoro channel at $c = 1$:
\begin{center}
\small
\renewcommand{\arraystretch}{1.15}
\begin{tabular}{ccccl}
\toprule
\textbf{Loop $L$} & \textbf{Shadow $S_{L+1}$} & \textbf{Value} & $\boldsymbol{\Ainf}$ \textbf{operation} & \textbf{Graph} \\
\midrule
$0$ & $S_1 \equiv \kappa_{\mathrm{ch}}$ & $1/2$ & $m_2$ (OPE bracket) & propagator \\
$1$ & $S_2$ & $1/2$ & $m_2$ (self-energy) & bubble \\
$2$ & $S_3 = \alpha$ & $2$ & $m_3(T,T,T) = -2T$ & theta \\
$3$ & $S_4$ & $10/27$ & $m_4(T,T,T,T) = \tfrac{40}{27}T$ & sunset + others \\
\bottomrule
\end{tabular}
\end{center}
The all-loop identity $Z^{(L)} = S_{L+1}$ for every $L$ is the natural
shadow--$\Ainf$--coproduct conjecture suggested by the Vol~I
correspondence, not a consequence of the finite computation above.
\end{proposition}

\begin{proof}
For $L\leq3$, the Feynman contribution $Z^{(L)}$ is computed from the
chosen regularised configuration-space integrals on
$\overline{\Conf}_n(C)$ using holomorphic residue calculus. The shadow
invariant $S_{L+1}$ is computed from the Vol~I convolution recursion
(the quadratic $Q_L(t) = 4\kappa_{\mathrm{ch}}^2 + 12\kappa_{\mathrm{ch}}\,\alpha\, t + (9\alpha^2 + 16\kappa_{\mathrm{ch}}\, S_4)\, t^2$).

The identification proceeds by matching both sides at each order:
\begin{itemize}
 \item $L = 1$: $Z^{(1)} = \kappa_{\mathrm{ch}} = S_2$. The one-loop bubble computes the self-energy correction, which equals $\kappa_{\mathrm{ch}}$ by Theorem~A.
 \item $L = 2$: $Z^{(2)} = \alpha = S_3 = 2$. The two-loop theta graph computes the cubic shadow, matching $|m_3(T,T,T)| = 2$.
 \item $L = 3$: $Z^{(3)} = S_4 = 10/27$. The three-loop graphs give the quartic contact, matching $m_4(T,T,T,T) = \tfrac{40}{27}\, T$ via the normalisation $S_4 = m_4\text{-coefficient} / 4$.
\end{itemize}
The finite dictionary matches the shadow tower recursion and the
$\Ainf$ operations at these four loop orders.
\end{proof}

\subsection{Shadow class from the loop expansion}
\label{subsec:shadow-class-loops}

The shadow class (Definition~\ref{def:shadow-class}, Vol~I) is determined by the loop expansion structure:
\begin{itemize}
 \item Class~$\mathbf{G}$: all loops $L \geq 1$ vanish (tree level only). The Heisenberg algebra (spin~$1$) is class~$\mathbf{G}$: the perturbative expansion terminates at $L = 0$.
 \item Class~$\mathbf{L}$: finitely many nonzero loop orders.
 \item Class~$\mathbf{C}$: convergent loop expansion (hence Borel summable).
 \item Class~$\mathbf{M}$: divergent loop expansion (Gevrey-$1$); Borel resummation is available only in the positivity/discriminant regime $\kappa_{\mathrm{ch}}^3S_4>0$ of Proposition~\ref{prop:class-m-borel-summability}. The Virasoro channel (spin~$2$) at $c = 1$ lies in this class~$\mathbf{M}$ regime: the shadow tower is infinite.
\end{itemize}

The spin~\(1\) Heisenberg channel is class~\(\mathbf{G}\). The
Virasoro subalgebra at spin~\(2\) already has an infinite shadow tower,
so a shadow theory retaining that sector is not the rank-one Heisenberg
scalar shadow.

\subsection{The MacMahon connection}
\label{subsec:macmahon-connection-feynman}

The MacMahon function records primitive BPS multiplicities. It is not
obtained by summing spin-channel modular characteristics.

\begin{proposition}[MacMahon primitive multiplicities]
\label{prop:constant-effective-kappa}
\ClaimStatusProvedHere
Let \(V_{\mathrm{BPS}}=\oplus_{n\geq 1}V_n\) with \(\dim V_n=n\).
Then
\[
 M(q)=\prod_{n\geq 1}(1-q^n)^{-n}
 =\mathrm{ch}\,\Sym(V_{\mathrm{BPS}})
 =\operatorname{PE}\!\left(\sum_{n\geq 1}nq^n\right).
\]
Thus the exponent \(n\) in the \(n\)-th factor of \(M(q)\) is the
primitive DT multiplicity \(\Omega(n)\), not a scalar
\(\kappa_{\mathrm{ch}}\)-channel.
\end{proposition}

\begin{proof}
The plethystic exponential identity gives the displayed product. The
plethystic logarithm is
\[
 \operatorname{PLog}M(q)=\sum_{n\geq 1}nq^n,
\]
which is precisely the primitive BPS spectrum of the \(\C^3\) Hall
algebra.
\end{proof}

\begin{remark}[Perturbative chiral invariants]
\label{rem:perturbative-feynman-summary}
The perturbative expansion is governed by graph combinatorics, Lie
algebra weights, holomorphic configuration integrals, the tree-level
identity $\kappa_{\mathrm{ch},s}=c/s$, the one-loop equality
$S_2=\kappa_{\mathrm{ch}}$, the two-loop value $S_3=\alpha=2$, the
three-loop value $S_4=10/27$, and the dictionary
$L\leftrightarrow S_{L+1}$ between loop order and shadow depth.
\end{remark}

% ============================================================
% E_3 Hochschild deformation theory: the 6d analogue of CFG25
% ============================================================

\section{\texorpdfstring{$E_3$}{E-3} Hochschild cohomology as chiral deformation theory}
\label{sec:e3-hochschild-deformation-theory}

In CFG25, the deformation quantisation of 3d Chern--Simons theory uses the $E_3$ bracket on $\HH^\bullet(\mathrm{Obs}_{\mathrm{cl}})$: the Maurer--Cartan element $\alpha \in \HH^2_{E_3}$ is the quantum correction, and the deformed observables $\mathrm{Obs}^q = \mathrm{CE}_*({\frak g})[[\hbar]]$ are controlled by this MC element. The 6d analogue is the $E_3$ Hochschild cohomology $\HH^\bullet_{E_3}(A, A)$ of a chiral algebra~$A$ as the deformation theory controlling $A \to A[[\varepsilon_1, \varepsilon_2, \varepsilon_3]]$. We open with the chain-level foundation: the brace-operad proof of Deligne's conjecture at $n = 2$ on Hochschild cochains, and its extension to $n = 3$ under the CY cyclic hypothesis via the Menichi BV operator.

\subsection{Deligne at \texorpdfstring{$n = 2$}{n=2}: the brace operad on Hochschild cochains}
\label{subsec:deligne-n2-brace}

The $E_2$-structure on $\CC^\bullet(A,A)$ is the Deligne conjecture at $n = 2$: McClure--Smith, Tamarkin, and Voronov. The chain-level witness is the brace operad.

\begin{definition}[Brace operations on Hochschild cochains;
\ClaimStatusDefinitional]
\label{def:brace-operations}
Let $A$ be an $A_\infty$-algebra over a field of characteristic zero. Write $\CC^n(A,A) = \Hom_k(A^{\otimes n}, A)$ for the Hochschild cochains (with a degree shift by $n$; our degree conventions follow Gerstenhaber). For $x \in \CC^m, y_1 \in \CC^{k_1}, \ldots, y_n \in \CC^{k_n}$, the \emph{$n$-brace} of $x$ by $y_1, \ldots, y_n$ is
\begin{equation}\label{eq:hoch-brace-explicit}
 \{x; y_1, \ldots, y_n\}(a_1, \ldots, a_N)
 \;=\;
 \sum_{0 \leq i_1 \leq i_2 \leq \cdots \leq i_n} (-1)^\eta\,
 x\bigl(a_1, \ldots, a_{i_1}, y_1(a_{i_1+1}, \ldots, a_{i_1+k_1}),
 a_{i_1+k_1+1}, \ldots, y_n(\ldots), \ldots, a_N\bigr),
\end{equation}
with $N = m + \sum k_j - n$, summation over non-overlapping insertions of $y_j$ starting at position $i_j$ with $i_j + k_j \leq i_{j+1}$, and the Koszul sign
\[
 \eta \;=\; \sum_{j=1}^n |y_j|\!\left(\sum_{\ell = 1}^{i_j} |a_\ell|\right)
 \;+\; \sum_{j < j'} |y_j| |y_{j'}|.
\]
The arity-$0$ brace is $\{x; \,\} = x$; the arity-$2$ brace with no insertions (the \emph{cup product}) is $(f \smile g)(a_1, \ldots, a_{m+n}) = (-1)^{|g| \sum_{\ell \leq m} |a_\ell|} f(a_1, \ldots, a_m) \cdot g(a_{m+1}, \ldots, a_{m+n})$, where the product uses $\mu_2$ of $A$.
\end{definition}

\begin{theorem}[Deligne at $n = 2$: brace model]
\label{thm:deligne-n2-brace-model}
Let $\Br$ denote the non-symmetric dg operad generated by the brace operations of Definition~\ref{def:brace-operations} with the brace-associativity relations
\[
 \bigl\{\{x; y_1, \ldots, y_m\};\, z_1, \ldots, z_n\bigr\}
 \;=\;
 \sum (-1)^{\eta'}\, \bigl\{x;\, \ldots,\, \{y_j; z_?, \ldots, z_?\},\, \ldots\bigr\},
\]
where the right-hand sum ranges over insertion patterns that assign each $z_k$ either to a free slot of $x$ or to one of the $y_j$. Then:
\begin{enumerate}[label=\textup{(\roman*)}]
\item The operations~\eqref{eq:hoch-brace-explicit} equip $\CC^\bullet(A,A)$ with a $\Br$-algebra structure for any $A_\infty$-algebra $A$.
\item There is a quasi-isomorphism of dg operads over $\Q$
\[
 \Br \;\xrightarrow{\;\sim\;}\; C_\bullet(E_2;\Q),
\]
canonical up to the $\mathrm{GRT}_1(\Q)$-torsor of Drinfeld associators (Tamarkin 1998; the choice of associator pins the rational homotopy type uniformly).
\item The Hochschild differential $b = \{\mu_A;-\} - \{-;\mu_A\}$ is a $\Br$-derivation, so $\CC^\bullet(A,A)$ is an $E_2$-algebra in dg vector spaces.
\end{enumerate}
Passing to cohomology, $(\HH^\bullet(A,A), \smile, [-,-]_\Ger)$ is a Gerstenhaber algebra with the Gerstenhaber bracket $[f,g]_\Ger = \{f;g\} - (-1)^{(|f|-1)(|g|-1)} \{g;f\}$.
\end{theorem}

\begin{proof}
(i) The brace identities are verified by direct chain-level computation (Voronov \cite{VoronovHomotopyGer} Theorem~3.3): each brace identity expands to a sum over shuffles of insertion positions with Koszul signs that cancel in pairs. (ii) McClure--Smith \cite{McClureSmithDeligne} Theorem~A constructs an explicit zigzag of dg operads through the Smith filtration of the sequence operad; Tamarkin \cite{TamarkinAnotherProof} \S3 gives the alternative construction via a Drinfeld associator, with the choice space a $\mathrm{GRT}_1(\Q)$-torsor (compare Remark~\ref{rem:deligne-kt-grt1-e2}). (iii) The identity $[b, \{-;-, \ldots, -\}] = 0$ reduces, after expansion, to the $A_\infty$-relation at $n = 2$: $b^2 = 0$ is the $A_\infty$-relation for $\mu_A$.
\end{proof}

\begin{remark}[Gerstenhaber pre-Lie and cohomological Poisson]
\label{rem:gerstenhaber-pre-lie}
The Gerstenhaber pre-Lie product $f \circ g = \{f;g\}$ is not associative; its failure of associativity is an explicit brace:
\[
 (f \circ g) \circ h \;-\; f \circ (g \circ h)
 \;=\; \{f; g, h\} \;+\; (-1)^{(|g|-1)(|h|-1)} \{f; h, g\}.
\]
The Gerstenhaber bracket, being the antisymmetrisation of $\circ$, is Lie on cohomology. The Poisson identity $[f, g \smile h]_\Ger = [f,g]_\Ger \smile h + (-1)^{(|f|-1)|g|} g \smile [f,h]_\Ger$ holds strictly on $\HH^\bullet(A,A)$ and up to explicit brace homotopies on $\CC^\bullet(A,A)$.
\end{remark}

\subsection{The Menichi BV operator on cyclic Hochschild cochains}
\label{subsec:menichi-bv-operator}

The cyclic pairing of a CY $A_\infty$-algebra upgrades the $E_2$-structure of Theorem~\ref{thm:deligne-n2-brace-model} to a chain-level framed $E_{d+1}$ structure. The upgrade operator is the Menichi BV operator, dual to the Connes cyclic $B$-operator on chains.

\begin{proposition}[Chain-level BV operator on cyclic Hochschild cochains]
\label{prop:menichi-delta-explicit}
Let $(A, \{\mu_n\}, \langle-,-\rangle)$ be a unital cyclic $A_\infty$-algebra of CY dimension $d$, with pairing $\langle-,-\rangle \colon A \otimes A \to k[-d]$. Fix a pair of dual bases $\{e_\alpha\}, \{e^\alpha\}$ satisfying $\langle e_\alpha, e^\beta\rangle = \delta_\alpha^\beta$. For $f \in \CC^n(A,A)$, define
\begin{equation}\label{eq:menichi-delta-explicit}
 \Delta(f)(a_1, \ldots, a_{n-1})
 \;=\;
 \sum_{i = 1}^{n}\, (-1)^{\sigma_i}\,
 \sum_\alpha \langle f(a_i, \ldots, a_{n-1}, e_\alpha, a_1, \ldots, a_{i-1}), e^\alpha\rangle^{\mathrm{PD}},
\end{equation}
where $\langle-,-\rangle^{\mathrm{PD}}$ denotes the Poincar\'e-dual pairing that uses $\langle-,-\rangle$ to convert a scalar back into an element of $A$ via the isomorphism $A \cong A^\vee[-d]$, and the sign $\sigma_i$ is the Koszul sign from cyclically permuting $n - i$ arguments of total degree past $i$ arguments with degrees $|a_j|$, plus the $d$-shift from the pairing. Then:
\begin{enumerate}[label=\textup{(\roman*)}]
\item $\Delta$ has degree $1 - d$: $\Delta \colon \CC^n(A,A) \to \CC^{n-1}(A,A)[1-d]$.
\item $\Delta^2 = 0$ and $[b, \Delta] = 0$.
\item \textup{(BV identity.)} For $f \in \CC^m, g \in \CC^n$,
\begin{equation}\label{eq:bv-identity}
 [f, g]_\Ger
 \;=\; (-1)^{|f|}\bigl(\Delta(f \smile g) - \Delta(f)\smile g - (-1)^{|f|} f \smile \Delta(g)\bigr).
\end{equation}
\end{enumerate}
Equivalently, $(\CC^\bullet(A,A), b, \smile, \Delta)$ is a BV algebra up to homotopy, and strictly a BV algebra on $\HH^\bullet(A,A)$, with BV-degree shift $d - 1$.
\end{proposition}

\begin{proof}
Menichi \cite{MenichiBV} Proposition~2.1 establishes $\Delta^2 = 0$ by direct chain-level expansion: the double cyclic rotation decomposes into pairs $(i, n-i)$ that cancel by cyclic symmetry of $\langle-,-\rangle$. The BV identity~\eqref{eq:bv-identity} is Menichi \cite{MenichiBV} Theorem~A (via $\infty$-inner-product transport of Tradler--Zeinalian \cite{TradlerZeinalian} Theorem~3.4): expanding $\Delta(f \smile g)$ in the cup and brace shows that the $\mu_2$-insertions match the Gerstenhaber pre-Lie up to the stated symmetriser, and Koszul signs match. The chain-homotopy $[b, \Delta] = 0$ follows because $b$ and $\Delta$ both factor through $\mu_A$-contraction and cyclic rotation, which commute as endomorphisms of $A^{\otimes \bullet}$.
\end{proof}

\begin{remark}[Degree shift comes from the cyclic pairing]
\label{rem:delta-degree-shift}
The degree $1 - d$ of $\Delta$ is the pairing-degree shift: the dual $e^\alpha$ lives in $A[-d]$ because $\langle-,-\rangle$ has degree $-d$, so reassembling $\CC^{n-1}$ from a scalar costs a $-d$ shift, compensated by the $+1$ from removing an argument slot. At $d = 0$ (ordinary BV): $\Delta$ has degree $+1$, the classical Connes $B$-dual. At $d = 1$ (topological string surfaces): $\Delta$ has degree $0$, the Kontsevich cyclic operator. At $d = 2$ (K3 / CY$_2$): $\Delta$ has degree $-1$, the classical BV generator. At $d = 3$ (CY$_3$): $\Delta$ has degree $-2$, the CY$_3$-twisted BV generator used below.
\end{remark}

\subsection{Deligne at \texorpdfstring{$n = 3$}{n=3}: three compatible \texorpdfstring{$E_3$}{E-3}-structures}
\label{subsec:deligne-n3-three-structures}

The naive statement ``Deligne at $n = 3$ for a CY$_3$ category'' conflates three distinct $E_3$-structures on Hochschild-type complexes. We name them, prove their chartwise compatibility, and separate the finite Rees gluing theorem from the remaining compact completion and realisation obligations.

\begin{theorem}[Three compatible \texorpdfstring{$E_3$}{E-3}-structures at CY\texorpdfstring{$_3$}{-3}]
\label{thm:deligne-n3-three-e3}
Let $X$ be a smooth Calabi--Yau threefold with cyclic $A_\infty$-enhancement $A_X$ of $D^b(\Coh(X))$. Three $E_3$-algebra structures arise on Hochschild-type complexes:
\begin{description}
\item[\textup{(A) Topological (CFG 2026).}] $\Obs^{\mathrm{cl}}_{\mathrm{CS}}(\R^3)$ is $E_3$ from the little $3$-disks operad on $\R^3$; the structure is topological, with Deligne at $n = 3$ supplying an $E_4$-algebra on its Hochschild cochains by Lurie \cite{LurieHA} Theorem~5.3.1.30.

\item[\textup{(B) Holomorphic (hCS on $\C^3$).}] $\Obs_{\hCS}(\C^3)$ is holomorphic-$E_3$ (Gwilliam--Williams \cite{GwilliamWilliams}) from the three complex directions of $\C^3$; the structure is real $E_3$ after Kontsevich--Tamarkin formality of $\cD_3$ (Theorem~\ref{thm:plat-hCS-quantum}).

\item[\textup{(C) Algebraic (CY$_3$-cyclic on Hochschild).}] $\CC^\bullet(A_X, A_X)$ is $E_3$ via the combined $(\Br, \Delta)$-action: the brace operad of Theorem~\ref{thm:deligne-n2-brace-model} supplies the $E_2$-structure; the Menichi BV operator of Proposition~\ref{prop:menichi-delta-explicit} supplies the $\Delta$ of degree $-2$ (since $d = 3$); their combined operadic level is $fE_4$ on cohomology, and the restricted $E_3$-sub-action (forgetting one $SO(2)$-plane) is canonical up to the contractible $S^2 = SO(3)/SO(2)$ of forgetful maps.
\end{description}
On each $\C^3$-chart $U \subseteq X$, structures \textup{(B)} and \textup{(C)} agree via Kontsevich--Tamarkin formality: the HKR isomorphism $\CC^\bullet(\cO_U, \cO_U) \simeq \mathrm{PV}^\bullet(U)$ intertwines the two $E_3$-actions up to a $\mathrm{GRT}_1(\Q)$-torsor.  On a finite DWR/Ran cover with bounds \((N,r,L,m)\), the hCS--Hall comparison glues at the Rees ordered Hall layer under the hypotheses of Proposition~\ref{prop:qca-finite-rees-layer-separation}.  The ordinary compact critical CoHA comparison requires \(R_{\mathrm{vc}}\) and the vanishing of the pro-descent \(\varprojlim^1\)-obstruction.
\end{theorem}

\begin{proof}
Structure \textup{(A)} is direct: $\R^3$ has a little $3$-disks operad action on observables, and Lurie \cite{LurieHA} Theorem~5.3.1.30 gives $E_{n+1}$ on Hochschild of $E_n$. Structure \textup{(B)} is Theorem~\ref{thm:plat-hCS-quantum}: shuffle integrals against the Bochner--Martinelli kernel realise the Fulton--MacPherson little $3$-disks action on $\Obs_{\hCS}(\C^3)$.

Structure \textup{(C)}: Theorem~\ref{thm:deligne-n2-brace-model} gives the $E_2$-action on $\CC^\bullet(A_X, A_X)$ via brace; Proposition~\ref{prop:menichi-delta-explicit} at $d = 3$ gives $\Delta$ of degree $-2$; Getzler's theorem identifies the homology of the combined operad with $\Ger \otimes \Lambda[\Delta]$ in the $d$-shifted grading, which is the homology of $fE_{d+1} = fE_4$. Forgetting the $SO(3) \supset SO(2)$-rotation gives an $E_3$-sub-action canonical up to $SO(3)/SO(2) = S^2$. The CY$_3$ volume form $\Omega_X = dz_1 \wedge dz_2 \wedge dz_3$ pins the forgetful map by distinguishing the third complex direction, making the $E_3$-action canonical.

Compatibility of \textup{(B)} and \textup{(C)} on a $\C^3$-chart: the HKR quasi-isomorphism $\CC^\bullet(\cO_{\C^3}, \cO_{\C^3}) \simeq \mathrm{PV}^\bullet(\C^3)$ (Hochschild--Kostant--Rosenberg, Tamarkin \cite{TamarkinAnotherProof}; see Remark~\ref{rem:three-atiyah-cocycles-qca} below for the Kapranov $\ell_\infty$-structure) is an $E_2$-quasi-isomorphism. The Kontsevich formality theorem for the $n$-disks operad (Kontsevich \cite{KontsevichOperads}; Tamarkin) lifts this to an $E_3$-quasi-isomorphism after choice of an associator. The formality bridge is canonical up to $\mathrm{GRT}_1(\Q)$.

Globally on compact $X$: the chartwise $E_3$-isomorphisms glue to the finite Rees Hall layer precisely under the finite DWR/Ran hypotheses named in Proposition~\ref{prop:qca-finite-rees-layer-separation}.  This does not yet identify the completed compact hCS observables with the ordinary critical CoHA: that stronger statement needs the monoidal vanishing-cycle realisation and pro-compatible transition data in the truncation variables.

Compatibility of \textup{(A)} with \textup{(B)}--\textup{(C)}: the CFG topological $E_3$ lives on $\R^3$; the hCS holomorphic $E_3$ lives on $\C^3 = \R^6$. The comparison is through the Wick-rotation / holomorphic-topological transition, a further layer that the CY$_3$-cyclic $E_3$ of \textup{(C)} does not need to traverse. Remark~\ref{rem:cy3-no-cfg-shortcut} of the CY$_3$ chain-level bridge chapter records this separation: CFG 2026's $E_3$ is not by itself a proof that the hCS $E_3$ on $\C^3$ agrees with the CY$_3$-cyclic $E_3$ on Hochschild; the two are related through the finite Rees comparison only after the DWR/Ran cyclic-atlas data have been fixed.
\end{proof}

\begin{remark}[Operadic separations]
\label{rem:deligne-n3-ap-cleared}
Three confusions are frequent on this axis and are cleared by Theorem~\ref{thm:deligne-n3-three-e3}. First, ``Deligne at $n = 3$ applied to a CY$_3$ category gives $E_3$ on $\CC^\bullet$'' is ambiguous: it may mean Lurie's iterated $\HH^\bullet$ of an $E_2$-input (giving $E_3$ without CY) or it may mean the BV-lift of the $E_2$-brace action under a CY$_3$ cyclic pairing. The two produce $E_3$-structures on different complexes; they coincide only after the Kontsevich--Tamarkin formality chartwise identification above. Second, ``BV is $E_3$'' is imprecise: BV is $fE_2$ on cohomology (Getzler); the $E_3$-promotion requires the CY dimension as an operadic shift, so ``$E_3$'' without naming the shift is type-ambiguous. Third, ``$R$-matrix equals the $E_3$-bracket on $\CC^\bullet$'' is false: the $R$-matrix at $d = 3$ arises from $E_2$-braiding on $\Rep(D(Y^+))$ after the Drinfeld-double step; the $E_3$-bracket on $\CC^\bullet$ controls the Maurer--Cartan deformation theory upstream of the double, not the braiding downstream.
\end{remark}

\begin{remark}[Three Atiyah cocycles and the CY$_3$-cyclic \texorpdfstring{$E_3$}{E-3}]
\label{rem:three-atiyah-cocycles-qca}
On a smooth Calabi--Yau threefold $X$ the Atiyah class $\At(T_X) \in \HH^1(X, T_X[-1])$ sources cohomologically distinct structures (Kapranov \cite{Kapranov1999}; Costello--Li \cite{CostelloLi2016}): the binary bracket $\ell_2 = \At(T_X)$; the tree-level Yukawa $Y_3 \in H^{0,3}(X)$, identified with Kapranov's minimal $\ell_3$; and the one-loop BCOV determinant-line or gravitational datum.  The latter is not a class \((\chi_{\mathrm{top}}(X)/24)[\Omega_X]^{0,1}\), since a holomorphic CY$_3$ volume form has no \((0,1)\)-component.  These structures live in different complexes, so none is the algebraic $E_3$-bracket of Theorem~\ref{thm:deligne-n3-three-e3}(C); the Menichi $\Delta$ of Proposition~\ref{prop:menichi-delta-explicit} sits in a different Hochschild tridegree and is pinned by the cyclic pairing, not by the Atiyah class. At $K3 \times E$ the Euler-characteristic coefficient of the BCOV one-loop datum vanishes because $\chi_{\mathrm{top}}(K3 \times E) = 24 \cdot 0 = 0$, while the product Yukawa and the algebraic $E_3$-bracket remain separate inputs.
\end{remark}

\begin{theorem}[Restricted \texorpdfstring{$E_3$}{E-3} feeds \texorpdfstring{$\Phi^{FA}_3$}{Phi-FA-3}]
\label{thm:restricted-e3-feeds-phi-fa-3}
Let $\cC$ be a smooth proper cyclic $A_\infty$-category of CY dimension $3$ on an admissible Stage-$1$ chart. The restricted $E_3$-sub-action on $\CC^\bullet(\cC,\cC)$ of Theorem~\ref{thm:deligne-n3-three-e3}(C) is the algebraic operadic input for the chartwise Stage-$1$ output $\PhiFA_3(\cC)$ of the two-stage factorisation $\Phi_3 = \SpCh_{\Sigma_2, C} \circ \PhiFA_3$. Specifically:
\begin{enumerate}[label=\textup{(\roman*)}]
\item The $E_3$-Maurer--Cartan element $\alpha \in \HH^2_{E_3}(A,A)$ of Subsection~\ref{subsec:mc-hh-e3} is the algebraic avatar of the hCS Maurer--Cartan solution, read through the $\Br \otimes \Lambda[\Delta]$ chain-level model.
\item The $E_4$-bracket $[-,-]_{E_3}$ of equation~\eqref{eq:mc-e3} is the higher Deligne-iteration bracket: Lurie \cite{LurieHA} Theorem~5.3.1.30 applied to the $E_3$-algebra $\CC^\bullet(\cC,\cC)$ gives an $E_4$-structure on $\HH^\bullet_{E_3}(\CC^\bullet(\cC,\cC))$, whose degree-$(-1)$ Lie bracket is $[-,-]_{E_3}$.
\item The classical limit $\alpha_1$ of the Maurer--Cartan element is the Gerstenhaber bracket of the tree-level Yukawa $Y_3$ of Remark~\ref{rem:three-atiyah-cocycles-qca} with itself; the one-loop coefficient is controlled by the BCOV determinant-line datum; the full tower $\{\alpha_k\}$ is the shadow tower of Conjecture~\ref{conj:mc-shadow-tower}.
\end{enumerate}
\end{theorem}

\begin{proof}
(i) The $E_3$-structure on $\CC^\bullet(\cC,\cC)$ of Theorem~\ref{thm:deligne-n3-three-e3}(C) produces the $E_3$-Hochschild cohomology $\HH^\bullet_{E_3}(A,A) = \mathrm{R}\Hom_{\Mod_A^{E_3}}(A,A)$ used in Subsection~\ref{subsec:e3-tricomplex}. Its Maurer--Cartan locus $\mathrm{MC}(\HH^2_{E_3})$ parametrises $E_3$-deformations of the input $A_\infty$-structure; by the chartwise formality Theorem~\ref{thm:deligne-n3-three-e3}, these match hCS deformations on $\C^3$-charts.

(ii) Direct: Lurie \cite{LurieHA} Theorem~5.3.1.30 gives $E_{n+1}$ on Hochschild of $E_n$; at $n = 3$, $E_4$; the $E_4$-structure contains a degree-$(-1)$ Lie bracket by Dunn--Lurie additivity $E_4 \simeq E_1 \otimes E_3$.

(iii) Classical-limit identification: Witten's $\partial\bar\partial$-action for hCS on $X$ has Maurer--Cartan solution at leading order the Yukawa coupling $Y_3$ (cubic tensor in the gauge-algebra directions), matching the brace $\{Y_3; Y_3\}$ in the Gerstenhaber pre-Lie grading. At one loop, Costello--Li \cite{CostelloLi2016} inserts the BCOV determinant-line contribution; this is the source of the $\alpha_2$ coefficient. The shadow tower recursion of Conjecture~\ref{conj:mc-shadow-tower} extends the pattern through finite-depth chain models.
\end{proof}

\begin{problem}[Global compact-CY$_3$ compatibility of the restricted $E_3$ input]
\label{op:restricted-e3-feeds-phi-fa-3-global}
Upgrade Theorem~\ref{thm:restricted-e3-feeds-phi-fa-3} from the
admissible Stage-$1$ chartwise locus to a global compact CY$_3$ statement by
constructing compatible finite DWR/Ran Rees comparisons at all
\((N,r,L,m)\), killing the resulting \(\varprojlim^1\)-obstruction, and
supplying the monoidal vanishing-cycle realisation
\(R_{\mathrm{vc}}\) from the Rees critical Hall layer to the ordinary
compact critical CoHA.
\end{problem}

\begin{remark}[CFG 2026 and the chain-level bridge]
\label{rem:cfg2026-chain-level-bridge}
CFG 2026's Deligne-at-$n=3$ derivation uses the little $3$-disks operad on the topological space $\R^3$ directly: classical observables of Chern--Simons with spectral parameter carry an $E_3$-algebra structure by factorisation, and Lurie's theorem produces the $E_4$-upgrade on their Hochschild cochains. Our derivation of Theorem~\ref{thm:deligne-n3-three-e3}(C) instead uses the brace operad on Hochschild cochains of a cyclic $A_\infty$-category, with the Menichi BV operator supplying the extra layer from the CY cyclic pairing. The two derivations produce compatible $E_3$-algebras chartwise on $\C^3$-charts, via the Kontsevich--Tamarkin formality zigzag $\Br \otimes \Lambda[\Delta] \xrightarrow{\sim} C_\bullet(fE_4;\Q) \to C_\bullet(E_3;\Q)$.  On compact CY$_3$ the finite DWR/Ran Rees comparison gives the current Hall-valued descent layer; completion and ordinary critical CoHA realisation are the additional obligations named in Problem~\ref{op:restricted-e3-feeds-phi-fa-3-global}.
\end{remark}

\subsection{The \texorpdfstring{$E_3$}{E-3} Hochschild tricomplex}
\label{subsec:e3-tricomplex}
\label{subsec:e3-bar-class-m}

For an $E_3$-algebra~$A$ on $\C^3$, the $E_3$ Hochschild cohomology $\HH^\bullet_{E_3}(A, A) = R\Hom_{\mathrm{Mod}_A^{E_3}}(A, A)$ carries an $E_4$-algebra structure by the higher Deligne conjecture (Lurie, \emph{Higher Algebra}, Theorem~5.3.1.30). The underlying cochain complex decomposes as a \emph{tricomplex} $C^{p,q,r}(A)$ with three commuting differentials $d_1, d_2, d_3$ (one per complex direction on $\C^3$):
\begin{equation}\label{eq:e3-tricomplex}
 d_i \colon C^{p,q,r} \to C^{p + \delta_{i1},\, q + \delta_{i2},\, r + \delta_{i3}},
 \qquad d_i^2 = 0,\quad [d_i, d_j] = 0.
\end{equation}

At the chain level, for a chiral algebra with $n$ generators:
\begin{equation}\label{eq:e3-tricomplex-chain-dim}
 \dim C^{p,q,r} = \binom{n}{p}\binom{n}{q}\binom{n}{r},
 \qquad
 \dim_{\mathrm{total}} = 8^n,
 \qquad
 P(s,t,u) = ((1+s)(1+t)(1+u))^n.
\end{equation}

The equivariant weights of the three differentials are the Omega-background parameters:
\[
 \mathrm{wt}(d_1) = h_1, \quad \mathrm{wt}(d_2) = h_2, \quad \mathrm{wt}(d_3) = h_3,
 \qquad h_1 + h_2 + h_3 = 0 \;\text{(CY)}.
\]
The CY condition forces $\mathrm{wt}(d_{\mathrm{total}}) = h_1 + h_2 + h_3 = 0$, ensuring $d_{\mathrm{total}}^2 = 0$.

\begin{proposition}[Finite-mode tricomplex pages by shadow class]
\label{prop:tricomplex-dimensions}
Let \(A\) be a chiral algebra with \(n\) finite-mode generators and shadow class \(\mathbf X\in\{G,L,C,M\}\). The finite-mode \(E_3\) tricomplex has the following terminal page when no later shadow differential is present:
\begin{center}
\small
\renewcommand{\arraystretch}{1.15}
\begin{tabular}{ccccc}
\toprule
\textbf{Class} & \textbf{Chain dim} & \textbf{Page dim} & \textbf{Poincar\'e} & \textbf{Deficit} \\
\midrule
$\mathbf{G}$ & $8^n$ & $8^n$ & $((1+s)(1+t)(1+u))^n$ & $0$ \\
$\mathbf{L}$ & $8^n$ & $8^n$ & $(1+t)^{3n}$ & $0$ \\
$\mathbf{C}$ & $8^n$ & $8^n$ & $(1+t)^{3n}$ & $0$ \\
$\mathbf{M}$ & $8^n$ & $\dim E_4=6^n$ & $(3t(1+t))^n$ & $8^n - 6^n$ at \(E_4\) \\
\bottomrule
\end{tabular}
\end{center}
The deficit for class~$\mathbf{M}$ at \(E_4\) is the image-kernel pair of the quartic shadow differential (see Proposition~\ref{prop:e3-spectral-degeneration} below). For \(n\leq3\) no later differential can act, hence this page is \(E_\infty\). For \(n\geq4\), \(d_5,d_6,\ldots,d_{n+1}\) are additional finite-mode operators. For class~$\mathbf{G}$, the cohomology Poincar\'e polynomial retains the three-variable form because all differentials vanish; for classes~$\mathbf{L}$ and~$\mathbf{C}$, the three variables collapse to the total degree~$t = stu$ after the differentials act.
\end{proposition}

\begin{proof}
The chain-level dimension \(8^n=(2^n)^3\) follows from the product decomposition \(C^{p,q,r}=\Lambda^p(V)\otimes\Lambda^q(V)\otimes\Lambda^r(V)\) with \(\dim V=n\). The cohomology for class~\(\mathbf G\) equals the chain level because all \(d_i\) vanish. For class~\(\mathbf M\), \(d_4\) contracts the top exterior class in each finite-mode factor: per factor, the \(d_4\)-homology of \(\Lambda^*(k^3)\) is \([0,3,3,0]\) with Poincar\'e \(3t+3t^2\) (Proposition~\ref{prop:virasoro-d4} in~\S\ref{subsec:e3-bar-class-m}); by K\"unneth, the \(n\)-fold tensor gives \(E_4\)-page \((3t(1+t))^n\) with total dimension \(6^n\). Classes~\(\mathbf L\) and~\(\mathbf C\) have \(d_4=0\) because the shadow tower terminates before depth~\(4\), respectively because charge conservation kills the quartic operator, so their cohomology equals the \(E_3\)-page dimension \(2^{3n}=8^n\).

This is the exterior-algebra count in the \(E_3\) tricomplex.
\end{proof}

\subsection{The spectral sequence \texorpdfstring{$E_3 \Rightarrow E_2 \Rightarrow E_1$}{E-3 => E-2 => E-1}}
\label{subsec:e3-spectral-sequence}

The tricomplex admits a spectral sequence by filtering on the third direction:
\begin{align*}
 E_1^{p,q,r} &= H^r(C^{p,q,\bullet}, d_3) & &\text{($d_3$-cohomology)} \\
 E_2^{p,q,r} &= H^q(E_1^{p,\bullet,r}, d_2) & &\text{($d_2$-cohomology of $d_3$-cohomology)} \\
 E_\infty^{p,q,r} &= H^p(E_2^{\bullet,q,r}, d_1) & &\text{($= \HH^\bullet_{E_3}(A,A)$).}
\end{align*}

The pages record the successive Hochschild reductions: $E_1 \approx E_2$-Hochschild (one direction computed), $E_2 \approx E_1$-Hochschild (two directions), $E_\infty = E_3$-Hochschild (all three).

\begin{proposition}[Degeneration page by shadow class]
\label{prop:e3-spectral-degeneration}
The spectral sequence of the $E_3$ tricomplex degenerates as follows:
\begin{center}
\small
\renewcommand{\arraystretch}{1.15}
\begin{tabular}{clcl}
\toprule
\textbf{Class} & \textbf{Degeneration} & \textbf{$d_4$ coefficient $\Delta$} & \textbf{Reason} \\
\midrule
$\mathbf{G}$ & $E^{\mathrm{ss}}_1 = E^{\mathrm{ss}}_\infty$ & $0$ & All $d_i = 0$ (abelian) \\
$\mathbf{L}$ & $E^{\mathrm{ss}}_3 = E^{\mathrm{ss}}_\infty$ & $0$ & Shadow terminates at depth~$2$ \\
$\mathbf{C}$ & $E^{\mathrm{ss}}_3 = E^{\mathrm{ss}}_\infty$ & $0$ & Charge conservation kills $d_4$ \\
$\mathbf{M}$, $n \leq 3$ & $E^{\mathrm{ss}}_4 = E^{\mathrm{ss}}_\infty$ & $8\kappa_{\mathrm{ch}} S_4$ & Degree reasons (support too narrow for $d_5$) \\
$\mathbf{M}$, $n \geq 4$ & $E_{n+2}$ at latest & $8\kappa_{\mathrm{ch}} S_4$ & $d_r = 0$ for $r \geq n+2$ by degree \\
\bottomrule
\end{tabular}
\end{center}
For the Virasoro at $c = 1$: $\Delta = 8 \cdot \tfrac{1}{2} \cdot \tfrac{10}{27} = \tfrac{40}{27}$, agreeing with the direct computation in Proposition~\ref{prop:virasoro-d4}.
\end{proposition}

\begin{proof}
For class~$\mathbf{G}$: all differentials $d_i$ vanish on an abelian algebra, so $E^{\mathrm{ss}}_1 = E^{\mathrm{ss}}_\infty$ immediately. For class~$\mathbf{L}$: the shadow tower has $S_k = 0$ for $k \geq 3$, so the higher differential $d_4$ (which is controlled by $S_4$) vanishes; $E^{\mathrm{ss}}_3 = E^{\mathrm{ss}}_\infty$. For class~$\mathbf{C}$: the $U(1)$-charge grading forces $d_4 = 0$ by the same mechanism as the $E_3$ bar for the bc system.

For class~$\mathbf{M}$: the $E_4$ page is supported in tridegrees $(p,q,r)$ with each entry in $\{0, 1\}$ (after the $d_1, d_2, d_3$ contractions), giving a band of total degree $[n, 2n]$ (width $n+1$). The differential $d_r$ has shift $r-1$. For $d_5$ to act nontrivially, both source degree $k$ and target degree $k - 4$ must lie in $[n, 2n]$, requiring $4 \leq n$. So $d_5 = 0$ for $n \leq 3$, giving $E^{\mathrm{ss}}_4 = E^{\mathrm{ss}}_\infty$. For $n \geq 4$: the spectral sequence terminates at $E_{n+2}$ at the latest (since $d_r$ with $r - 1 > n$ has shift exceeding the band width).

\end{proof}

\begin{proposition}[Class $\mathbf{M}$ bar growth past \texorpdfstring{$n = 3$}{n = 3}]
\label{prop:e3-bar-class-m-asymptotic}
Let $A$ be a class-$\mathbf{M}$ chiral algebra with $n$ finite-mode generators and shadow invariants \(S_r\). On the finite one-primitive quotient, the \(d_4\)-homology has Poincar\'e polynomial
\[
 P(E_4;t)=(3t+3t^2)^n,
 \qquad
 \dim E_4=6^n .
\]
The later differentials are precisely
\[
 d_r:E_r^k\longrightarrow E_r^{k-r+1},
 \qquad 5\leq r\leq n+1,\qquad n+r-1\leq k\leq 2n.
\]
Thus \(E_\infty=E_4\) for \(n\leq3\). For \(n=4\),
\[
 \dim E_\infty=1296-2\,\rho_5,
 \qquad
 \rho_5=\operatorname{rank}\!\bigl(d_5:E_5^8\to E_5^4\bigr),
 \qquad
 0\leq\rho_5\leq81 .
\]
In general
\[
 \dim E_\infty
 =
 \dim H^\bullet\!\left(E_4,\ d_5+d_6+\cdots+d_{n+1}\right),
\]
where \(d_r\) is the finite-mode operator induced by the \(r\)-th shadow \(S_r\).
The ordinary PBW bar complex is infinite-dimensional before this finite-mode completion is imposed.
\end{proposition}

\begin{proof}
After the three Hochschild contractions, the \(d_4\)-homology of one finite-mode generator is \([0,3,3,0]\), with support in degrees \(1\) and \(2\). K\"unneth gives \(E_4=[0,3,3,0]^{\otimes n}\), hence \(P(E_4;t)=(3t+3t^2)^n\) and \(\dim E_4=6^n\). The support of \(E_4\) is the degree band \([n,2n]\). A differential \(d_r\) lowers total degree by \(r-1\); it can act only from degrees \(k\) satisfying \(k\leq2n\) and \(k-r+1\geq n\), equivalently \(n+r-1\leq k\leq2n\). This is empty for every \(r\geq5\) when \(n\leq3\). For \(n=4\) the only non-empty case is \(r=5,k=8\); the source and target are both \(3^4=81\)-dimensional, giving the displayed rank formula. The last possible differential has \(r=n+1\). The spectral sequence identity gives the final formula.
\end{proof}

\begin{remark}[Admissible shadow instances]
\label{rem:e3-bar-class-m-admissible-instances}
The number of possible higher-shadow positions is
\[
 N(n) \;:=\; \sum_{r = 5}^{n+1} \bigl(n - r + 2\bigr) \;=\; \frac{(n-3)(n-2)}{2} \quad (n \geq 3).
\]
For \(n\leq3\) this number is \(0\). For \(n=4\) it is \(1\), the single map \(d_5:E_5^8\to E_5^4\). For larger \(n\) the ranks of these maps are additional shadow data, not consequences of the \(d_4\) calculation alone.
\[
\begin{array}{c|c|c}
n & \dim E_4 & N(n)\\
\hline
3 & 216 & 0\\
4 & 1296 & 1\\
5 & 7776 & 3\\
6 & 46656 & 6\\
10 & 60466176 & 28
\end{array}
\]
\end{remark}

\subsection{Maurer--Cartan deformation on \texorpdfstring{$\HH^2_{E_3}$}{HH-2-E-3}}
\label{subsec:mc-hh-e3}

The Maurer--Cartan element governing the chiral deformation quantisation lives in $\HH^2_{E_3}(A, A)$. The MC equation is:
\begin{equation}\label{eq:mc-e3}
 d(\alpha) + \tfrac{1}{2}[\alpha, \alpha]_{E_3} = 0,
 \qquad \alpha \in \HH^2_{E_3}(A, A),
\end{equation}
where $[-, -]_{E_3}$ is the degree-$(-1)$ Lie bracket from the $E_4$-structure on $\HH^\bullet_{E_3}$ (the higher Deligne conjecture provides $E_{n+1}$ structure on $\HH^\bullet_{E_n}$; here $n = 3$, giving $E_4$, and the $E_4$ structure includes a Lie bracket of degree $-1$).

The MC element expands in powers of $\sigma_3 = h_1 h_2 h_3$:
\begin{equation}\label{eq:mc-expansion}
 \alpha = \sum_{k \geq 1} \sigma_3^{k-1}\, \alpha_k,
 \qquad
 \alpha_k \in \HH^2_{E_3}(A, A).
\end{equation}

\begin{conjecture}[MC element = shadow tower]
\label{conj:mc-shadow-tower}
\ClaimStatusConjectured
The $k$-th MC coefficient $\alpha_k$ is identified with the shadow invariant $S_k$ (Vol~I, Definition~\ref{def:shadow-invariants}):
\begin{equation}\label{eq:mc-shadow-identification}
 \alpha_k \;\longleftrightarrow\; S_k \;\longleftrightarrow\; Z^{(k-1)} \;\longleftrightarrow\; m_k,
\end{equation}
where $Z^{(L)}$ is the regularised $L$-loop Feynman contribution when
the graph-integral package is defined (finite evidence is recorded in
Proposition~\ref{prop:loop-shadow-identification}), and $m_k$ is the
$\Ainf$~operation of the bar complex (Vol~I~\cite{VolI}). The conjectural content is that the MC equation~\eqref{eq:mc-e3}, order by order in $\sigma_3$, reproduces the shadow tower recursion.
\end{conjecture}

\begin{remark}[Chain-level tricomplex obligation]
\label{rem:mc-shadow-why-conjecture}
The identification~\eqref{eq:mc-shadow-identification} is conjectural because it requires the \emph{explicit} $E_3$-chiral tricomplex structure on $A$ for CY$_3$ targets. The \((\infty,1)\)-categorical CY-A$_3$ input of Theorem~\ref{thm:derived-framing-m3b2} gives the framed existence statement; the explicit chain-level tricomplex data for non-formal algebras, needed to compute the MC coefficients $\alpha_k$, remains open. The correspondence $S_k \leftrightarrow Z^{(k-1)}$ is currently a finite Virasoro-channel witness through $k\leq4$ (Proposition~\ref{prop:loop-shadow-identification}), while the $S_k \leftrightarrow m_k$ comparison belongs to the Vol~I bar-complex lane. The composite arrow $\alpha_k \leftrightarrow S_k$ is the new content: it identifies the MC coefficients on $\HH^2_{E_3}$ with the shadow tower. The conjecture asserts that the deformation-theoretic packaging (the MC equation on $\HH^2_{E_3}$) and the combinatorial packaging (the shadow tower recursion) describe the same underlying structure.
\end{remark}

The MC series has convergence controlled by the shadow class:
\begin{center}
\small
\renewcommand{\arraystretch}{1.15}
\begin{tabular}{clll}
\toprule
\textbf{Class} & \textbf{MC series type} & \textbf{Convergence} & \textbf{Example} \\
\midrule
$\mathbf{G}$ & Polynomial (one term) & Radius $= \infty$ & Heisenberg: $\alpha = \alpha_1$ \\
$\mathbf{L}$ & Polynomial (finite) & Radius $= \infty$ & Yangian: $\alpha = \alpha_1 + \sigma_3 \alpha_2$ \\
$\mathbf{C}$ & Convergent & Radius $> 0$ & bc system \\
$\mathbf{M}$ & Gevrey-$1$ divergent & Radius $= 0$ & Virasoro: Borel on $\kappa_{\mathrm{ch}}^3S_4>0$ \\
\bottomrule
\end{tabular}
\end{center}

\subsection{The CFG25 dictionary: 3d vs 6d deformation quantisation}
\label{subsec:cfg25-dictionary}

The relationship between CFG25 (3d CS on $\R^3$) and 6d hCS on $\C^3$ is captured by a dictionary that translates each element of the 3d deformation quantisation to its 6d chiral analogue:

\begin{center}
\small
\renewcommand{\arraystretch}{1.3}
\begin{tabular}{p{4.5cm}p{6cm}}
\toprule
\textbf{CFG25 (3d CS on $\R^3$)} & \textbf{6d hCS on $\C^3$} \\
\midrule
Single differential $d$ on $\mathrm{CE}_*(\frak{g})$ & Three commuting differentials $d_1, d_2, d_3$ on $\HH^\bullet_{E_3}(A, A)$ \\
One parameter $\hbar$ & $(h_1, h_2, h_3)$ with $h_1 + h_2 + h_3 = 0$; essential parameter $\sigma_3$ \\
MC in $\HH^2(\mathrm{Obs}_{\mathrm{cl}})$: $E_3$ bracket from $\R^3$ & MC in $\HH^2_{E_3}(A, A)$: $E_3$ bracket from $\C^3$ tricomplex \\
$\alpha = \hbar\, r + \hbar^2 \cdot \ldots$ (classical $r$-matrix $+$ corrections) & $\alpha = S_1 + \sigma_3 S_2 + \sigma_3^2 S_3 + \cdots$ (shadow tower) \\
$d(\alpha) + \tfrac{1}{2}[\alpha,\alpha] = 0$ (CYBE $+$ quantum) & $d(\alpha) + \tfrac{1}{2}[\alpha,\alpha]_{E_3} = 0$ (shadow recursion) \\
Koszul dual $\mathrm{CE}_*(\frak{g})^! = U_q(\frak{g})$ & Koszul dual $\HH^\bullet_{E_3}(A)^! =$ chiral quantum group [\textsc{conj}] \\
Boundary $V_k(\frak{g})$ ($E_\infty$-chiral) & Boundary $Y(\widehat{\fgl}_1)$ ($E_1$-chiral) \\
Deligne: $E_3$ on $\HH^\bullet(V_k)$ (from $E_\infty$ input) & Deligne: $E_2$ on $\HH^\bullet(Y)$ (from $E_1$ input) \\
\bottomrule
\end{tabular}
\end{center}

The Deligne conjecture level drop (the last row) is the structural
reason that the 6d theory requires the \emph{external} \(E_3\) structure
from holomorphic factorisation on \(\C^3\), rather than the
\emph{internal} \(E_3\) from the higher Deligne conjecture: the boundary
algebra \(Y(\widehat{\fgl}_1)\) is only \(E_1\)-chiral, so the Deligne
conjecture produces only \(E_2\) on its Hochschild cochains
(Proposition~\ref{prop:deligne-level-drop-cfg25}).  The
Omega-background supplies the equivariant weights of this factorisation
structure; it is not an additional operadic direction.

\subsection{The tangent-obstruction sequence}
\label{subsec:e3-tangent-obstruction}

The \(E_3\) deformation complex of a chiral algebra~\(A\) has tangent
and obstruction spaces
\begin{equation}\label{eq:e3-tangent-obstruction}
 \HH^0_{E_3}(A,A) \to \HH^1_{E_3}(A,A) \to \HH^2_{E_3}(A,A) \to \cdots
\end{equation}
depending on the differential, relations, coefficients, and
completion.  For the associated-graded free tricomplex with \(n\)
generators and zero internal differential one obtains the model
dimensions
\begin{align}
 \dim \HH^0_{E_3} &= n & &\text{(infinitesimal $E_3$-automorphisms)}, \label{eq:hh0-dim} \\
 \dim \HH^1_{E_3} &= 3n & &\text{(first-order $E_3$ deformations: $n$ per direction)}, \label{eq:hh1-dim} \\
 \dim \HH^2_{E_3} &= \tfrac{3n(3n-1)}{2} & &\text{(obstructions: tridegree-$2$ contributions)}, \label{eq:hh2-dim}
\end{align}
giving the free-model virtual dimension
\(\mathrm{vdim}_{\mathrm{gr}} = 3n - \tfrac{3n(3n-1)}{2}\), which is
\(0\) for \(n = 1\) and negative for \(n \geq 2\).  No assertion is made
that these are the Hochschild dimensions of an arbitrary chiral algebra.

\begin{proposition}[\texorpdfstring{$\HH^2_{E_3}$}{HH-2-E-3} in the free associated-graded tricomplex]
\label{prop:hh2-tridegree}
In the free associated-graded tricomplex with \(n\) generators and zero
internal differential, the tridegree \(p + q + r = 2\) part receives
contributions:
\begin{enumerate}[label=\textup{(\roman*)}]
 \item Tridegrees $(2,0,0)$, $(0,2,0)$, $(0,0,2)$: each $\binom{n}{2}$, total $3 \cdot \tfrac{n(n-1)}{2}$.
 \item Tridegrees $(1,1,0)$, $(1,0,1)$, $(0,1,1)$: each $n^2$, total $3n^2$.
\end{enumerate}
Sum: $\tfrac{3n(n-1)}{2} + 3n^2 = \tfrac{3n(3n-1)}{2}$. Verified at $n = 1$ (dim $= 3$), $n = 3$ (dim $= 36$), $n = 24$ (dim $= 2556$).
\end{proposition}

\begin{proof}
Direct count of the tridegree-$(p,q,r)$ contributions to the exterior-algebra tricomplex. For item~(i): $\binom{n}{2} \cdot \binom{n}{0} \cdot \binom{n}{0} = \binom{n}{2}$ for each of $3$ orderings. For item~(ii): $\binom{n}{1}^2 \cdot \binom{n}{0} = n^2$ for each of $3$ orderings. The Mukai value $n = 24$ gives \(3\cdot276+3\cdot576=2556\).
\end{proof}



%% Four-manifold invariants from 6d holomorphic CS

\subsection{Four-manifold invariants from 6d holomorphic theory}
\label{subsec:four-manifold-6d}

The dimensional hierarchy of holomorphic CS produces invariants of
manifolds in defect dimension $d_{\mathrm{defect}} = \dim(\text{theory}) - 2$:
$3d$ CS gives WRT invariants of $3$-manifolds; $5d$ holomorphic CS gives
conjectural invariants via affine Yangians; $6d$ $(2,0)$ theory gives
$4$-manifold invariants via quantum toroidal algebras. Compactification
of the $6d$ $(2,0)$ theory on $X^4 \times \Sigma_g$ produces an effective
theory on $X^4$; for $\Sigma_g = T^2$, this is $4d$ $\mathcal{N} = 4$
SYM under the Vafa--Witten twist.

\begin{conjecture}[Four-manifold invariants from the \texorpdfstring{$6d$}{6d} chiral construction]
\label{conj:four-manifold-6d-invariants}
\ClaimStatusConjectured
\index{four-manifold invariants!6d holomorphic theory}%
For a compact oriented $4$-manifold $X$ with $b_1(X) = 0$, the $6d$
invariant $Z_{6d}(X \times T^2)$ equals the Vafa--Witten partition
function $Z_{\mathrm{VW}}(X; \tau)$, a modular form
\textup{(}or mock modular form\textup{)} of weight $w = -\chi(X)/2$
under $\mathrm{SL}_2(\Z)$. The modularity character depends on the
intersection form $Q_X$:
\begin{center}
\small
\renewcommand{\arraystretch}{1.15}
\begin{tabular}{lcccl}
\toprule
$X$ & $\chi(X)$ & $b_2^+(X)$ & $w$ & $Z_{\mathrm{VW}}$ \\
\midrule
$K3$ & $24$ & $3$ & $-12$ & $1/\eta^{24}$ \textup{(}genuine modular\textup{)} \\
$T^4$ & $0$ & $3$ & $0$ & $1$ \textup{(}trivial\textup{)} \\
$\C\bP^2$ & $3$ & $1$ & $-3/2$ & mock modular \\
$S^2 \times S^2$ & $4$ & $1$ & $-2$ & mock modular \\
$E(n)$ & $12n$ & $2n{-}1$ & $-6n$ & $1/\eta^{12n}$ \textup{(}$b_2^+ > 1$ for $n \geq 2$\textup{)} \\
\bottomrule
\end{tabular}
\end{center}

The dichotomy is: $b_2^+(X) > 1 \Rightarrow$ genuine modular form;
$b_2^+(X) = 1 \Rightarrow$ mock modular \textup{(}requiring
non-holomorphic completion\textup{)}.

The scalar entering the rank-one Vafa--Witten Hilbert-scheme product is
the surface Euler characteristic $\chi_{\mathrm{top}}(X)$; it is not, by
itself, the compact chiral invariant of a CY$_3$ object. A comparison
with a chiral characteristic of $A_{\mathrm{Tot}(K_X)}$ requires the
framed $\Phi_3^{(\Sigma_2,C)}$ specialisation and the local
non-compact boundary condition on $\mathrm{Tot}(K_X)$,
$\kappa_{\mathrm{cat}} = \chi(\mathcal{O}_X)$ \textup{(}Noether:
$(\chi + \sigma)/4$\textup{)},
$\kappa_{\mathrm{fiber}} = b_2(X)$ \textup{(}lattice rank\textup{)}.

For $K3$: $\kappa_{\mathrm{ch}}^{\mathrm{OPE}} = 24$ (the free-field
generator count at the OPE-level on $\mathrm{Tot}(K_{K3})$;
distinct from $\kappa_{\mathrm{ch}}^{\mathrm{Hodge}}(D^b\mathrm{Coh}(K3)) = 2$),
$\kappa_{\mathrm{cat}} = 2$,
$\kappa_{\mathrm{fiber}} = 22$, and $Z_{\mathrm{VW}} = 1/\eta^{24}$
is unconditional \textup{(}CY-A$_2$\textup{)}. For $\C\bP^2$:
$\kappa_{\mathrm{cat}} = 1$, and the mock modularity of $Z_{\mathrm{VW}}$
is established by Manschot--Thomas.

The construction via $\mathrm{Tot}(K_X)$ requires the framed CY$_3$
comparison data for the non-compact total space. The factorization homology route
$\int_{X^4} A_{E_2}$ is conjectural.
\end{conjecture}

\begin{remark}[Noether formula and Hilbert scheme consistency]
\label{rem:four-manifold-noether-hilb}
For complex surfaces, the holomorphic Euler characteristic
$\chi(\mathcal{O}_X) = (\chi(X) + \sigma(X))/4$ (Noether's formula)
provides a cross-check: $\chi(\mathcal{O}_{K3}) = (24 - 16)/4 = 2$,
$\chi(\mathcal{O}_{\C\bP^2}) = (3 + 1)/4 = 1$. The rank-$1$ VW
generating function $\prod_{k \geq 1} (1 - q^k)^{-\chi(X)}$ reproduces
Hilbert scheme Euler characteristics $\chi(\mathrm{Hilb}^n(X))$
(Gottsche 1990): for $K3$, the sequence $1, 24, 324, 3200, 25650,
176256$ matches $\prod(1 - q^k)^{-24}$ exactly.

\noindent The relevant four-manifold invariants
evaluates the surface catalogue, Noether formula, Hilbert scheme Euler
characteristics through $n = 10$, VW modular weights, mock/genuine
modularity dichotomy at $b_2^+ = 1$, $\kappa_\bullet$-spectrum, and
$S$-duality cross-checks.
\end{remark}

%% CHIRAL WRT INVARIANTS FROM THE K3 QUANTUM TOROIDAL ALGEBRA

\subsection{Chiral Witten--Reshetikhin--Turaev invariants}
\label{subsec:chiral-wrt}

The $3d$ WRT invariant $Z_{\mathrm{WRT}}(M^3;\, \frakg,\, k)$ of a closed $3$-manifold $M$ is constructed from three ingredients:
\begin{enumerate}[label=\textup{(\roman*)}]
 \item a quantum group $U_q(\frakg)$ at a root of unity $q = e^{2\pi i/(k + h^\vee)}$, whose representation category is a modular tensor category;
 \item a surgery presentation $M = S^3(L)$ for a framed link $L$ with linking matrix $L_{ij}$;
 \item a state sum over colourings of the link components.
\end{enumerate}
The $6d$ chiral WRT invariant $Z^{\mathrm{ch}}(K3;\, N)$ replaces each ingredient as follows.

\begin{center}
\renewcommand{\arraystretch}{1.3}
\small
\begin{tabular}{llp{5.5cm}}
 \toprule
 \textbf{$3d$ CS ingredient} & \textbf{$6d$ chiral analogue} & \textbf{Construction locus} \\
 \midrule
 $U_q(\frakg)$ at root of unity & $U_{q,t}(\fg_{K3}^{\mathrm{tor}})$ at $\omega = e^{2\pi i/N}$ & requires root-of-unity bridge \\
 Surgery $M = S^3(L)$ & Handle decomposition: $K3 = D^4 \cup 22\,(D^2 \times D^2) \cup D^4$ & Freedman--Quinn handle theorem \\
 Linking matrix $L_{ij}$ & Intersection form $Q_{K3} = 3U \oplus 2E_8(-1)$, sig.\,$(3, 19)$ & lattice computation \\
 State sum & Factorization homology on handles & requires compact handle factorisation \\
 $\sigma(L) = \mathrm{sig}(L_{ij})$ & $\sigma(K3) = 3 - 19 = -16$ & signature computation \\
 \bottomrule
\end{tabular}
\end{center}

\begin{conjecture}[Abelian chiral WRT invariant of K3]
\label{conj:chiral-wrt-k3}
\ClaimStatusConjectured
For the abelian truncation of the K3 quantum toroidal algebra at level $N \geq 2$, the chiral WRT invariant is
\begin{equation}\label{eq:chiral-wrt-k3}
 Z^{\mathrm{ch,ab}}(K3;\, N)
 \;=\; N^{\chi(K3)/2} \cdot \varphi(N)
 \;=\; N^{12} \cdot \varphi(N),
\end{equation}
where $\varphi(N)$ is the Gauss phase product over the Mukai lattice $\widetilde{\Lambda}_{K3}$ of signature~$(4, 20)$:
\[
 \varphi(N)
 \;=\; \prod_{a=1}^{24}
 \frac{g(\sigma_a, N)}{|g(\sigma_a, N)|},
 \qquad
 g(\sigma, N)
 \;=\; \sum_{n=0}^{N-1} \exp\!\bigl(\pi i \sigma n^2 / N\bigr).
\]
Here $\sigma_a = +1$ for $a = 1, \ldots, 4$ and $\sigma_a = -1$ for $a = 5, \ldots, 24$.

The phase $\varphi(N) = 1$ for all $N$, because $\sigma(K3) = 4 - 20 = -16 \equiv 0 \pmod{8}$ and an even unimodular lattice with signature divisible by~$8$ has trivial Gauss phase. Therefore $Z^{\mathrm{ch,ab}}(K3;\, N) = N^{12}$, an integer.
\end{conjecture}

The exponent $12 = \chi(K3)/2$ arises from the handle decomposition: each Betti direction contributes $\sqrt{N}$ to the quantum dimension, and the total is $N^{(b_0 + b_2 + b_4)/2} = N^{(1 + 22 + 1)/2} = N^{12}$. At $N = 2$:
\[
 Z^{\mathrm{ch,ab}}(K3;\, 2) \;=\; 2^{12} \;=\; 4096.
\]
The same integer is obtained from the Gauss-sum product formula, the
handle-by-handle decomposition, and the Betti-number product.

\begin{remark}[Connection to G\"ottsche and Vafa--Witten]
\label{rem:chiral-wrt-gottsche-vw}
The chiral WRT invariant $Z^{\mathrm{ch}}(K3;\, N) = N^{12}$ is \emph{not} the same as:
\begin{enumerate}[label=\textup{(\alph*)}]
 \item The G\"ottsche Euler characteristics $\chi(\mathrm{Hilb}^n(K3)) = p_{24}(n)$ with generating function $\prod_{k \geq 1}(1 - q^k)^{-24}$: these are instanton moduli space invariants, while $Z^{\mathrm{ch}}$ is a topological state count.
 \item The Vafa--Witten partition function $Z_{\mathrm{VW}}(K3, \mathrm{SU}(N);\, \tau)$, which is a modular form of weight $-\chi(K3)/2 = -12$.
\end{enumerate}
The relationship: $Z_{\mathrm{VW}}(K3, \mathrm{SU}(N);\, \tau) = Z^{\mathrm{ch}}(K3;\, N) \cdot F(\tau;\, N)$ where $F$ is the one-loop fluctuation determinant around the truncated vacuum. The chiral WRT extracts the zero-mode (topological) component.
\end{remark}

\begin{conjecture}[Chiral WRT of \texorpdfstring{$K3 \times E$}{K3 x E} and \texorpdfstring{$\Delta_5$}{Delta-5}]
\label{conj:chiral-wrt-k3xe}
\ClaimStatusConjectured
For $K3 \times E$, the product formula gives
\begin{equation}\label{eq:chiral-wrt-k3xe}
 Z^{\mathrm{ch}}(K3 \times E;\, N) \;=\; Z^{\mathrm{ch}}(K3;\, N) \cdot N \;=\; N^{13},
\end{equation}
where the factor $N$ is the chiral Verlinde dimension $V^{\mathrm{ch}}_N(E)$ on the elliptic curve~$E$. The exponent $13 = \chi(K3)/2 + 1$ is related to the weight $\kappa_{\mathrm{BKM}}(\Delta_5) = 5$ of the Igusa cusp form $\Delta_5$: in the unnormalised Igusa-square convention, the DMVV formula
\[
 \frac{1}{\Phi_{10}^{\mathrm{un}}}=\Delta_5^{-2}
 = \sum_N p^N \chi(S^N K3;\, q,\, y)
\]
has $N$-th coefficients whose polynomial prefactor matches $N^{13}\).
The identification of \(Z^{\mathrm{ch}}(K3 \times E;\, N)\) with the
asymptotics of \(\Delta_5\) Fourier coefficients becomes a theorem only
after CY-A and the Vol~I Borcherds-lift comparison are supplied.
\end{conjecture}

%% THE CHIRAL VOLUME CONJECTURE

\subsection{The chiral volume conjecture}
\label{subsec:chiral-volume-conjecture}

The $3d$ volume conjecture (Kashaev 1997; Murakami--Murakami 2001) asserts that
for a hyperbolic knot $K$ in $S^3$,
\begin{equation}\label{eq:3d-volume-conjecture}
\lim_{N \to \infty} \frac{2\pi}{N}\,
 \log\bigl| J_N\!\bigl(K;\, e^{2\pi i / N}\bigr)\bigr|
 \;=\; \operatorname{Vol}(S^3 \setminus K),
\end{equation}
where $J_N(K; q)$ is the $N$-th coloured Jones polynomial and
$\operatorname{Vol}(S^3 \setminus K)$ is the hyperbolic volume of the
knot complement. The $6d$ holomorphic theory does not have knots
(holomorphic curves are not knots: Obstruction~$2$ in
Section~\ref{subsec:cfg25-obstruction-hierarchy}), and the lift-rate
assessment of Section~\ref{subsec:cfg25-result3} classified the volume
conjecture as a \emph{complete gap} with $0\%$ lift rate. This
assessment is correct for a literal lift. What follows is an
\emph{analogue}, not a lift: it occupies the same structural
position (relating quantum invariants to classical geometry) but
uses different mathematics on both sides.

\medskip\noindent\textbf{The $3d$-to-$6d$ dictionary.}
\begin{center}
\small
\renewcommand{\arraystretch}{1.2}
\begin{tabular}{lll}
\toprule
& \textbf{3d Chern--Simons} & \textbf{6d chiral construction} \\
\midrule
Ambient space & $S^3$ & CY$_3$ $X$ \\
Submanifold & knot $K$ (real $1$-manifold) & holomorphic curve $C$ (complex $1$-manifold) \\
Quantum invariant & $J_N(K;\, e^{2\pi i/N})$ & $Z_N^{\mathrm{ch}}(C, X)$ \\
Classical geometry & $\operatorname{Vol}(S^3 \setminus K)$ & $\lvert\mathrm{AJ}(C)\rvert$ \\
Rigidity & Mostow (unique metric) & attractor mechanism (unique at attractor) \\
\bottomrule
\end{tabular}
\end{center}

\begin{conjecture}[Chiral volume conjecture]
\label{conj:chiral-volume-conjecture}
\ClaimStatusConjectured
Let $X$ be a Calabi--Yau threefold with Hodge-normalised holomorphic
$3$-form $\Omega$ ($\int_X \Omega \wedge \bar\Omega = 1$). Let
$C \subset X$ be a holomorphic curve representing the class
$\beta \in H_2(X, \Z)$. Define the \emph{charge-$N$ chiral partition
function}
\[
 Z_N^{\mathrm{ch}}(C, X)
 \;:=\;
 \sum_{\substack{d \,\geq\, 1 \\ d\,[\hspace{0.2pt}C\hspace{0.2pt}]\,=\, N\beta}}
 \Omega_{\mathrm{DT}}(d \cdot \beta,\, X)\, q^{\hspace{0.3pt}d}
\]
where $\Omega_{\mathrm{DT}}(\gamma, X)$ is the Donaldson--Thomas
invariant in charge class $\gamma$. Define the \emph{Abel--Jacobi
period}
\[
 \mathrm{AJ}(C)
 \;=\; \int_\Gamma \Omega
 \;\in\; J^2(X) \;=\; H^{2,1}(X)^* \big/ H_3(X, \Z),
\]
where $\Gamma$ is a $3$-chain with $\partial\Gamma = C - C_0$ for a
reference cycle $C_0$ \textup{(}the Griffiths normal function\textup{)}.
Then:
\begin{equation}\label{eq:chiral-volume-conjecture}
 \lim_{N \to \infty} \frac{1}{N}\,
 \log\bigl\lvert Z_N^{\mathrm{ch}}(C, X) \bigr\rvert
 \;=\;
 \frac{1}{2\pi}\,\bigl\lvert \mathrm{AJ}(C) \bigr\rvert,
\end{equation}
where $\lvert\mathrm{AJ}(C)\rvert$ is the norm in the Hodge metric on
$J^2(X)$.
\end{conjecture}

\begin{remark}[Three regimes]
\label{rem:chiral-volume-three-regimes}
The conjecture specialises to three regimes determined by the Hodge
number $h^{2,1}(X)$.

\emph{Toric regime} ($X$ toric, e.g.\ resolved conifold). The
Abel--Jacobi period reduces to the complexified K\"ahler parameter
$t_C = \int_C \omega$, and the conjecture becomes
$\lim (1/N)\,\log\lvert Z_N^{\mathrm{DT}}\rvert = t_C$. This is
consistent with the Okounkov--Reshetikhin--Vafa asymptotics: for the
conifold, $Z_N^{\mathrm{DT}}(N\beta) \sim p(N)\cdot Q^N$ where
$Q = e^{-t}$ and $p(N)$ is the partition function, giving growth rate
$t + O(1/\!\sqrt{N}\,)$.

\emph{Compact regime} ($h^{2,1} > 0$, e.g.\ quintic). The period is a
genuine Griffiths normal function depending on the complex structure
$\tau \in \mathcal{M}_{cs}(X)$. The conjecture predicts a specific
function of $\tau$ (the Abel--Jacobi image) from the
asymptotics of DT invariants. Under mirror symmetry $X \leftrightarrow
X^\vee$, this reduces to the statement that the mirror map
$t(z) = w_1(z)/w_0(z)$ controls DT growth on the mirror side.

\emph{Rigid regime} ($h^{2,1} = 0$, e.g.\ Schoen manifold). The
intermediate Jacobian $J^2(X)$ is trivial, so $\mathrm{AJ}(C) = 0$ for
all curves $C$. The conjecture predicts \emph{subexponential} growth:
$(1/N)\,\log\lvert Z_N^{\mathrm{ch}}\rvert \to 0$. Proxy check: for the
conifold at $Q = 1$, $Z_N = p(N)$ and
$(1/N)\,\log p(N) \sim \pi/\!\sqrt{N} \to 0$ (Hardy--Ramanujan).
\end{remark}

\begin{remark}[Relation to the attractor mechanism]
\label{rem:chiral-volume-attractor}
For BPS black holes in Type~IIB on $X$, the attractor mechanism
(Ferrara--Kallosh--Strominger 1995) fixes the complex structure at a
point $\tau_*(\gamma)$ determined by the charge $\gamma$. At the
attractor point, the BPS entropy is
$S_{\mathrm{BH}}(\gamma) = \pi\,\lvert Z(\gamma, \tau_*)\rvert^2$
where $Z(\gamma, \tau) = \int_\gamma \Omega(\tau)$ is the central
charge (Ooguri--Strominger--Vafa 2004). This is the period integral in
the charge class $\gamma$. The attractor mechanism provides the
analogue of Mostow rigidity: at the attractor point, the complex
structure is \emph{unique}, and the chiral volume conjecture reduces to
a single number. The conjecture is strongest at attractors.
\end{remark}

\begin{remark}[Period growth and obstructions]
\label{rem:chiral-volume-evidence}
\emph{Period growth}: (i) the toric regime is verified to leading order by
the Okounkov--Reshetikhin--Vafa asymptotics; (ii) mirror symmetry
relates DT invariants of $X$ in class $\beta$ to periods of $X^\vee$,
so the mirror map \emph{is} the Abel--Jacobi map on the mirror side;
(iii) the attractor mechanism gives
$\log\Omega_{\mathrm{BPS}}(\gamma) \sim \pi\lvert Z(\gamma)\rvert^2$,
supporting the period--growth link.

The chiral-algebraic formulation, the trace over Fock modules of $A_C$,
lives on the framed object-level locus of
Theorem~\textup{\ref{thm:cy-to-chiral-d3}}. Chain-level convergence of
the trace for non-formal compact CY$_3$ is analytic input; the
normalisation of $\Omega$ is fixed by the Hodge metric; the absence of a
CY analogue of Mostow rigidity makes the statement a family over moduli;
and the large-$N$ limit of DT invariants in a fixed curve class requires
compact-CY$_3$ asymptotics.

\noindent The curve geometry, CY$_3$ data, DT asymptotics,
Abel--Jacobi periods, three-regime analysis, \(3d\) comparison,
obstruction classes, and conjectural hypotheses are the data entering
the large-\(N\) claim.
\end{remark}

\begin{proposition}[Toric small-$N$ convergence for the resolved conifold]
\label{prop:chiral-volume-small-N-conifold}
For the resolved conifold $X = \cO(-1) \oplus \cO(-1) \to \P^1$ with K\"ahler parameter $t = 2\pi s$, $s > 0$, and $Q = e^{-t}$, the charge-$N$ chiral partition function in the unique compact curve class is $Z_N(X, Q) = p(N) \cdot Q^N$ where $p(N)$ is the integer partition function. Then
\[
 \frac{1}{N}\,\log\bigl\lvert Z_N(X, Q) \bigr\rvert \;=\; -t + \frac{\log p(N)}{N},
\]
and $(1/N) \log p(N) \sim \pi \sqrt{2/(3N)} \to 0$ by Hardy--Ramanujan, so
\[
 \lim_{N \to \infty} \frac{1}{N}\,\log\bigl\lvert Z_N(X, Q) \bigr\rvert \;=\; -t \;=\; -\bigl\lvert \mathrm{AJ}(C) \bigr\rvert.
\]
At finite $N$ the convergence rate is $|{-}t - N^{-1}\log p(N)|$:
\begin{center}
\small\renewcommand{\arraystretch}{1.1}
\begin{tabular}{rrrr}
\toprule
$N$ & $N^{-1} \log p(N)$ & Hardy--Ramanujan $\pi\sqrt{2/(3N)}$ & ratio \\
\midrule
$10$ & $0.3738$ & $0.8112$ & $0.46$ \\
$100$ & $0.1907$ & $0.2565$ & $0.74$ \\
$1000$ & $0.0723$ & $0.0811$ & $0.89$ \\
\bottomrule
\end{tabular}
\end{center}
The rate matches the conjectured $|\mathrm{AJ}(C)|/(2\pi) \cdot 2\pi = t$ (absolute value) to leading order at every $N$, with sub-leading $O(N^{-1/2})$ Hardy--Ramanujan correction.
\end{proposition}

\begin{proof}
The $\mathrm{DT}$ partition function of the resolved conifold in curve class $N[\P^1]$ is $\mathrm{DT}_N(Q) = p(N) Q^N$ (Okounkov--Reshetikhin--Vafa, \emph{Quantum Calabi--Yau and Classical Crystals}, Theorem~2). Taking logarithms gives the stated formula. The Hardy--Ramanujan expansion $\log p(N) = \pi \sqrt{2N/3} + O(\log N)$ yields $(1/N) \log p(N) \sim \pi \sqrt{2/(3N)}$. The complexified K\"ahler parameter $t = \int_{[\P^1]} \omega$ coincides with $|\mathrm{AJ}(C)|$ for the unique toric curve class since $h^{2,1}(X) = 0$.
\end{proof}

\begin{proposition}[Rigid-vertical small-$N$: subexponential rate for K3 classes]
\label{prop:chiral-volume-rigid-vertical}
For $X = \mathrm{K3}$ with curve class $\beta \in H_2(\mathrm{K3}, \Z)$ of self-intersection $\beta^2 = 2h - 2$, the vertical $\mathrm{PT}$ partition function is $\mathrm{PT}(X, N\beta; q) = q^{hN}/\prod_{k \geq 1} (1 - q^k)^{24}$ (Kawai--Yoshioka / Maulik--Pandharipande--Thomas). Its Fourier coefficients $c_N$ are the coefficients of $\prod(1-q^k)^{-24}$; Cardy asymptotics give $\log c_N \sim 4\pi \sqrt{N}$, so
\[
 \frac{1}{N} \log c_N \;\to\; 0 \quad (N \to \infty),
\]
matching the conjecture's \emph{rigid regime} prediction $|\mathrm{AJ}(C)| = 0$ on vertical classes. Small-$N$ values: $c_1 = 24$, $c_2 = 324$, $c_3 = 3200$, $c_5 = 176256$, with $(1/N)\log c_N$ decreasing monotonically from $3.18$ at $N = 1$ to $1.21$ at $N = 50$, and Cardy ratio converging to $1$ as $N^{1/2} \to \infty$.
\end{proposition}

\begin{proof}
The generating-function identity is the Maulik--Pandharipande--Thomas theorem (2010, \emph{Topology}); the Cardy asymptotic $\log c_N \sim 4\pi\sqrt{N}$ follows from the circle-method expansion of $\eta^{-24}$ (Rademacher formula, applied to weight $-12$ eta-quotient). Vertical $\mathrm{AJ}(\beta) \in J^2(\mathrm{K3} \times E)$ lies in the $E$-factor and integrates to zero on $\beta$ (by the K\"unneth decomposition $\beta \in H_2(\mathrm{K3}) \otimes 1$), giving $|\mathrm{AJ}(\beta)| = 0$. Consistency with the conjecture is the subexponential rate.
\end{proof}

\section{The K3 chiral quantum-group \texorpdfstring{$R$}{R}-matrix, associator, and classification}
\label{sec:qca-k3-chiral-quantum-group}

The K3 chiral Hall--Drinfeld double
$\mathcal{D}_\hbar(\mathcal{Y}^{\mathrm{Hall}}_\hbar(\mathrm{CoHA}_{K3\times E}))$
of Chapter~\ref{ch:k3-chiral-bialgebra-platonic} occupies a fourth quasi-Hopf
taxonomic slot. Five structural data pin its quantum-group content:
the Siegel-elliptic dynamical $R$-matrix,
the twisted Siegel--Borcherds associator,
a super-quasi-Hopf compatibility,
an $A_\infty$-quasi-Hopf upgrade at the Humbert walls,
and the classification of its gauge-equivalence class by a $1$-dimensional
Lie-bialgebra cohomology group.

The ambient-qualifier discipline is twofold. At chain level
(explicit complexes, named differentials, Chevalley--Eilenberg cochains,
pentagon 3-cocycle witness $d_{\mathrm{CE}}\phi^{(3)} = -\frac{1}{2}[\phi^{(2)},\phi^{(2)}]_{\mathrm{Ger}}$),
the structure computes; at $(\infty,1)$-categorical level (derived
categories of $D$-modules on $\overline{\mathcal{A}_2}$, Lurie $\mathrm{HA}$
formalism of $E_2$-algebras in a presentable $\infty$-category), the structure
universalises. Both lanes appear simultaneously: the explicit Pasol--Zagier
Kronecker--Eisenstein series on $\mathbb{H}_2$ (chain-level), the
Etingof--Kazhdan super-category extension (categorical), the Bruinier 2002
Chern-class reciprocity (cohomological), the Felder dynamical YBE on the
paramodular quotient (geometric), and the Markl--Shnider--Stasheff
$A_\infty$-pentagon at Humbert walls (operadic). Each appearance is one face
of the same five-datum object.

The five theorems and five bridging propositions that follow make each datum
precise, with arguments citing primary literature.

\begin{theorem}[Siegel-elliptic dynamical \texorpdfstring{$R$}{R}-matrix: a fourth class in the Drinfeld--Belavin taxonomy]
\label{thm:qca-R-Sieg-dynamical}
Let $A_\tau$ be the principally polarised abelian surface with period matrix
$\bigl(\begin{smallmatrix}\tau&z\\z&\rho\end{smallmatrix}\bigr)\in\mathbb{H}_2$.
The Siegel-elliptic dynamical $R$-matrix
$R_{\mathrm{Sieg,dyn}}(u-v;\tau,\rho,z)\in\mathrm{End}(V_{24}\otimes V_{24})$
of Chapter~\ref{ch:k3-chiral-bialgebra-platonic} Definition~\ref{def:R-Sieg-dynamical},
constructed from the genus-$2$ Riemann theta function
$\theta[{\alpha\atop\beta}](w;\tau,z,\rho)$ via the Pasol--Zagier 2013 extension of
Kronecker--Eisenstein to $\mathbb{H}_2$, occupies a fourth structural slot beyond
Drinfeld 1985's rational/trigonometric/elliptic classification: it is neither
rational (no $1/(u-v)$-pole structure), nor trigonometric (no $\mathbb{C}^*$-loop),
nor purely elliptic (no single elliptic parameter $\tau$). The fourth class is
Siegel-elliptic-dynamical: dynamical parameter in $\mathbb{H}_2$, spectral parameter
separate, with the three-tier degeneration
$R_{\mathrm{Sieg,dyn}}\xrightarrow{\rho\to i\infty} R^{\mathrm{ell,dyn}}_{\mathrm{Felder}}
\xrightarrow{\tau\to i\infty} R^{\mathrm{rat}}(u)$
recovering the Felder 1994 elliptic dynamical and Yang rational $R$-matrices as
Humbert-cusp and double-Humbert-cusp degenerations.

On the paramodular locus $K(1)\subset\mathbb{H}_2$,
$R_{\mathrm{Sieg,dyn}}$ satisfies the dynamical Yang--Baxter equation
\[
 R_{12}(u_1-u_2)\,R_{13}(u_1-u_3)\,R_{23}(u_2-u_3)
 \;=\;
 R_{23}(u_2-u_3)\,R_{13}(u_1-u_3)\,R_{12}(u_1-u_2)
\]
modulo shifts of dynamical parameter $(\tau,\rho,z)\mapsto(\tau,\rho,z)+h_i$
at Cartan-generator insertions (the Felder dynamical shift, extended to
$\mathrm{Sp}_4$-compatible period matrix). Off $K(1)$, the YBE acquires a
half-integer Jacobi-index anomaly proportional to the Maass multiplier
$v_{\Delta_5}$ of $\Delta_5$ as a half-integral-weight paramodular form.
\end{theorem}

\begin{proof}
The fourth-class statement is Chapter~\ref{ch:k3-chiral-bialgebra-platonic},
Theorem~\ref{thm:R-sieg-fourth-class}; the dynamical YBE on $K(1)$ is
Theorem~\ref{thm:R-Sieg-YBE} of the same chapter. The three-tier degeneration is
Remark~\ref{rem:R-Sieg-felder-genus1} (genus-$1$ Felder limit at $\rho\to i\infty$,
rational limit at $\tau\to i\infty$). The classical $\hbar^1$-YBE follows from
Pasol--Zagier 2013 Theorem~3.2 cyclic Siegel-Kronecker identity; the quantum
$\hbar^2$-YBE follows from Etingof--Kazhdan quantisation of the Manin pair, as
in Chapter~\ref{ch:k3-chiral-bialgebra-platonic}
Proposition~\ref{prop:pasol-zagier-to-dynamical-YBE}. The Felder dynamical
mechanism--dynamical parameter shift by simple root $h_i$ at each Cartan
insertion--is the genus-$2$ lift of Felder 1994 \S 4, with the Cartan eigenvalue
$(\alpha,\alpha)$ supplying the shift magnitude through the cross-bracket
$[t_{12},e_\alpha]=(\alpha,\alpha)\,e_\alpha$.
\end{proof}

\begin{theorem}[Pentagon and hexagon for the Siegel--Borcherds associator]
\label{thm:qca-pentagon-hexagon-sieg-bor}
The twisted associator
$\widetilde{\Phi}^{\mathrm{Sieg\text{-}Bor}}_{\mathrm{Sp}_4}[\Phi_{10}^{\mathrm{un}}/\eta^{24}]$
lives in the pro-nilpotent pro-$\hbar$-filtered completion
\[
 \widetilde{\Phi}^{\mathrm{Sieg\text{-}Bor}}_{\mathrm{Sp}_4}[\Phi_{10}^{\mathrm{un}}/\eta^{24}]
 \;\in\;
 \exp\bigl(\widehat{\mathfrak{t}^{\mathrm{Sieg}}_{2,[2]}\oplus\mathfrak{n}_+^{\mathrm{imag}}}\bigr)^{\mathrm{grouplike}},
\]
with $\mathfrak{t}^{\mathrm{Sieg}}_{2,[2]}$ the genus-$2$ infinitesimal pure-braid
Lie algebra (Enriquez--Gomez-Gonzalez--Maassarani 2022) and
$\mathfrak{n}_+^{\mathrm{imag}}$ the BKM imaginary-root nilpotent.

The pentagon equation holds at orders $\hbar^{\leq 3}$. The order-$\hbar^2$ part
reduces to the closedness condition $d_{\mathrm{CE}}\phi^{(2)} = 0$, which holds
by the Kohno--Drinfeld four-term relation on the symmetric part and the BKM
Jacobi identity on the imaginary part. At order $\hbar^3$, the obstruction
correction
\[
 \phi^{(3)}
 \;=\;
 \zeta(3)\,c_{\mathrm{symm}}
 \;+\;
 \tfrac{26}{3}\,c_{\mathrm{timelike}}
 \;+\;
 \bigl(\Phi_{10}^{\mathrm{un}}/\eta^{24}\bigr)\,c_{\Phi_{10}^{\mathrm{un}}},
\]
with $c_{\mathrm{symm}},c_{\mathrm{timelike}},c_{\Phi_{10}^{\mathrm{un}}}$ three
$3$-cochains, linearly independent in the coboundary space
$B^3 \subset C^3(\mathfrak{t}^{\mathrm{Sieg}}_{2,[2]}\oplus\mathfrak{n}_+^{\mathrm{imag}};\,\mathcal{O}(\mathbb{D}))$
with function coefficients over the period domain $\mathbb{D}$ (the third
coefficient is the modular function $\Phi_{10}^{\mathrm{un}}/\eta^{24}$;
restricted to a fixed chamber it is read as its scalar value),
balances the Gerstenhaber self-bracket $-\frac{1}{2}[\phi^{(2)},\phi^{(2)}]_{\mathrm{Ger}}$
term by term in the Borcherds product expansion of $\Phi_{10}^{\mathrm{un}}$ at Fourier level
$c(N,\ell)$ (Borcherds 1998, Theorem~13.3; Gritsenko--Nikulin 1997, Theorem~1.2).
The constant $26/3$ decomposes as $\mathrm{rk}(\mathrm{II}_{25,1})/3 = 26/3$ (Fake Monster
Cartan rank, $\mathrm{rk}(\mathrm{II}_{25,1}) = 26$, divided by cyclic triple-leg
averaging), not a Virasoro central charge.

Both hexagons at $\hbar^{\leq 2}$ hold for the pair
$(\widetilde{\Phi}^{\mathrm{Sieg\text{-}Bor}},R_{\mathrm{Sieg,dyn}})$ on
$K(1)\subset\overline{\mathcal{A}_2}$; restriction to $K(1)$ is forced by the
half-integer Jacobi-index anomaly of $\Delta_5$ on $\mathrm{Sp}_4(\mathbb{Z})$.
\end{theorem}

\begin{proof}
The pentagon is Chapter~\ref{ch:k3-chiral-bialgebra-platonic},
Theorem~\ref{thm:pentagon-sieg-bor}, together with the KZ coboundary lift and
Borcherds identity propositions (the order-$\hbar^{\leq 2}$ Knizhnik--Zamolodchikov pentagon proposition
and the order-$\hbar^3$ Borcherds-pentagon proposition). The hexagon is
Theorem~\ref{thm:hexagon-R-Sieg}, decomposing at order
$\hbar^2$ into: (H-I) Pasol--Zagier 2013 cyclic Siegel-Kronecker compatibility;
(H-II) Maass-form compatibility on $K(1)$ via Lorgat 2020 Proposition~4.1; and
(H-III) the Casimir--imaginary cross-term grouplike-exponential compatibility.
The $26/3$ arises in the order-$\hbar^3$ Borcherds-pentagon proposition via Nikulin
1979 Lorentzianisation $\Gamma^{4,20}\hookrightarrow\mathrm{II}_{25,1}$
(of rank $26$) together with cyclic triple-leg $|\mathbb{Z}_3|^{-1}$
averaging.
\end{proof}

\begin{theorem}[Super-quasi-Hopf structure and \texorpdfstring{$A_\infty$}{A-inf} upgrade at Humbert walls]
\label{thm:qca-super-A-infty-upgrade}
The $\mathbb{Z}/2$-grading on the K3 chiral Hall--Drinfeld double is genuine
super-grading (Koszul sign rule on the pentagon, hexagon, and antipode axioms),
sourced from the K3 sigma-model's Ramond--Ramond fermion number and lifted
through the $\widetilde{M}_{24}$ cover to BKM imaginary-root parity: odd imaginary
simple roots $\alpha$ with $(\alpha,\alpha)<0$ are fermionic generators, bosonic
when $(\alpha,\alpha)=0$. The Igusa form $\Delta_5$ lives in odd weight $5$;
$\Phi_{10}^{\mathrm{un}}=\Delta_5^2$ in even weight $10$.

Along the Humbert divisors $H_1\cup H_4\subset\overline{\mathcal{A}_2}$, the
Siegel--Borcherds associator degenerates; the strict quasi-Hopf pentagon at
$\hbar^{\leq 3}$ upgrades to an $A_\infty$-quasi-Hopf pentagon with higher-coherence
data $\phi^{(n)},\,n\geq 4$, satisfying the Markl--Shnider--Stasheff $A_\infty$
pentagon equation
\[
 d_{\mathrm{CE}}\,\phi^{(n)}
 \;=\;
 -\,\tfrac{1}{2}\sum_{\substack{k+l=n+1\\ k,l\geq 2}}
 \bigl[\phi^{(k)},\phi^{(l)}\bigr]_{\mathrm{Ger}},
\]
with the $n$-th higher-coherence datum carrying the $n$-th Fourier-Jacobi
coefficient of the expansion
$\Phi_{10}^{\mathrm{un}}(\rho,\tau,z) = \sum_{m\geq 1}\phi_{10,m}(\tau,z)e^{2\pi im\rho}$
(Eichler--Zagier Fourier-Jacobi theory).
\end{theorem}

\begin{proof}
The genuineness of the super-grading is Chapter~\ref{ch:k3-chiral-bialgebra-platonic}
Remark~\ref{rem:super-genuine-fermion}; the lift through the K3 sigma model and
$\widetilde{M}_{24}$ cover uses Gaberdiel--Hohenegger--Volpato 2012 for the
$M_{24}$-equivariance of the K3 elliptic genus, and Gritsenko--Nikulin 1995
\S 2 for the BKM imaginary-root fermionic/bosonic classification by
$(\alpha,\alpha)$-sign. The $A_\infty$ upgrade at Humbert walls is
Theorem~\ref{thm:A-infty-upgrade} together with
the $A_\infty$ higher-coherence Humbert proposition. The recursion at depth
$n$ is Markl--Shnider--Stasheff 2002 Chapter~8, with the Maurer--Cartan equation
in the Chevalley--Eilenberg DG Lie algebra of
$(\mathfrak{t}^{\mathrm{Sieg}}_{2,[2]}\oplus\mathfrak{n}_+^{\mathrm{imag}})[[\hbar]]$
supplying the higher-pentagon constraint.
\end{proof}

\begin{theorem}[Pentagon coboundary decomposition at order \texorpdfstring{$\hbar^{\leq 3}$}{hbar-<= 3}]
\label{thm:qca-pentagon-hbar3-explicit}
Let $\mathfrak{L} := \mathfrak{t}^{\mathrm{Sieg}}_{2,[2]} \oplus \mathfrak{n}_+^{\mathrm{imag}}$ denote the home Lie algebra of the Siegel--Borcherds associator. The three independent $3$-coboundaries on $\mathfrak{L}$ are
\[
 c_{\mathrm{symm}}
 \;:=\;
 [t_{12},t_{23}] \otimes [t_{23}, t_{13}],
 \qquad
 c_{\mathrm{timelike}}
 \;:=\;
 \sum_{\alpha\in\mathrm{II}_{25,1}^{\mathrm{timelike}}} \alpha\otimes\alpha\otimes\alpha,
 \qquad
 c_{\Phi_{10}^{\mathrm{un}}}
 \;:=\;
 \sum_{\substack{(N,\ell,M)\\ 4NM-\ell^2=-1}} c_{\phi_{0,1}}(N,\ell)\, e_\alpha\otimes e_\beta\otimes e_\gamma,
\]
with $t_{ij}$ the Drinfeld--Kohno infinitesimal braiding, $\mathrm{II}_{25,1}^{\mathrm{timelike}}$ the timelike component of the Fake Monster Cartan, and $c_{\phi_{0,1}}$ the $(N,\ell)$-th Fourier coefficient of the K3 weak Jacobi form. The pentagon equation at order $\hbar^{\leq 3}$ reads
\[
 \phi^{(3)}
 \;=\;
 \zeta(3)\,c_{\mathrm{symm}}
 \;+\;
 \tfrac{26}{3}\,c_{\mathrm{timelike}}
 \;+\;
 (\Phi_{10}^{\mathrm{un}}/\eta^{24})\,c_{\Phi_{10}^{\mathrm{un}}},
\]
where the constant $26/3 = \mathrm{rk}(\mathrm{II}_{25,1})/3$ is the Fake Monster Cartan rank $26$ divided by cyclic triple-leg averaging, and the third coefficient is the modular function $\Phi_{10}^{\mathrm{un}}/\eta^{24}$ on the period domain: the decomposition is read in $C^3(\mathfrak{L};\,\mathcal{O}(\mathbb{D}))$ with function coefficients, its scalar form being the evaluation in a fixed chamber. The three cochains are linearly independent in the coboundary space $B^3(\mathfrak{L};\,\mathcal{O}(\mathbb{D})) \subset C^3$ (as coboundaries they vanish in cohomology; independence holds at the cochain level), so the decomposition is unique.
\end{theorem}

\begin{theorem}[Conditional Etingof--Kazhdan characteristic comparison]
\label{thm:qca-etingof-kazhdan-super-classification}
Assume two compact Hall--Drinfeld source packages over $K3\times E$ have
the same finite Hall/CoHA source window, positive and negative halves,
Cartan part, continuous Hopf pairing, source-to-target comparison maps,
radical descent, no-extra-relations kernel equality, PBW comparison, and
exact inverse-limit passage. Assume moreover that their primitive maps
land in the Manin-pair target
$(\mathfrak{g}_{\Delta_5},\,\mathfrak{n}_+^{\mathrm{imag}}\oplus
\mathfrak{h}^{\mathrm{imag,rk\,23}})$ with the super structure and the
paramodular $K(1)$-restriction of Theorem~\ref{thm:hexagon-R-Sieg}.
Then the comparison data define a characteristic map
\[
 \chi^{\mathrm{cmp}}_{\Delta_5}\colon
 \Def_{\mathrm{src}\to\Delta_5}(\mathbf H)
 \longrightarrow
 H^2_{\mathrm{Lie}}\bigl(\mathfrak q_{K(1)};\mathrm{grt}_1\bigr)
 \;\cong\;
 \mathbb{C}\cdot\Delta_5.
\]
Here $\mathfrak q_{K(1)}$ is the $K(1)$-restricted target deformation
Lie algebra of the recognised Manin-pair package. The two quantisations
are compared by their finite Hall--Borcherds matrices and by their image
under $\chi^{\mathrm{cmp}}_{\Delta_5}$; the scalar Igusa line is the
automorphic target of the comparison, not a direct computation of
the intrinsic second Lie-bialgebra cohomology of
$\mathfrak{g}_{\Delta_5}$ and not a construction of $\mathbf H_{\Delta_5}$
from automorphy alone.
\end{theorem}

\begin{proof}
Etingof--Kazhdan Parts~IV--V classify quasi-Hopf quantisations once the
completed super Manin-pair input and associator class have been fixed.
They do not construct the compact Hall source. The finite Hall--Borcherds
recognition gates supply that source: comparison matrices, radical
descent, no-extra-relations kernel equality, PBW compatibility, and strict
completion transitions. Bruinier 2002 Proposition~5.1 supplies the
Heegner-divisor Chern-class reciprocity which places the automorphic
class of $\Delta_5$ in the transgressed paramodular obstruction line.
Gritsenko--Nikulin and Borcherds supply the denominator arithmetic and
weight. These automorphic facts are target conditions for the recognised
source package, not a Lie-bialgebra cohomology computation by
automorphy.
\end{proof}

\begin{remark}[Three faces of the number \texorpdfstring{$8$}{8}]
\label{rem:qca-three-faces-of-eight}
The number $8$ appears in three structurally distinct roles: $(\mathrm{a})$ categorical Mukai doubling $2\,c_+(\mathrm{Mukai}(K3))=2\cdot 4 = 8$; $(\mathrm{b})$ the branch-normalised local-exponent denominator of the automorphic line of $\Delta_5$ around the Humbert divisor $H_1\subset\overline{\mathcal{A}_2}$; $(\mathrm{c})$ the Lusztig specialisation order $\ell=8$ of the Hall--Drinfeld double, when that specialisation is constructed. Bruinier~$2002$ Proposition~$5.1$ controls the automorphic Chern class. It does not identify a Chow class with a small-quantum-group cohomology class, and it does not imply a canonical equation for $\hbar^2$.
\end{remark}

\begin{theorem}[Root-of-unity comparison at \texorpdfstring{$\ell=8$}{l=8}]
\label{thm:qca-hbar-minus-eighth-lusztig}
The Lusztig specialisation of the Hall--Drinfeld double at $\ell=8$ produces a
distinguished root-of-unity object at $\zeta = e^{2\pi i/8}$. Bruinier 2002, Proposition~5.1
(Chern-class formula for Heegner divisors) computes the automorphic Chern class
\[
 c_1\bigl(\mathcal{L}^{\Delta_5}\bigr)\big|_{H_1}
 \;=\;
 [\xi]
 \;\in\;
 \mathrm{CH}^1(H_1)\otimes\mathbb{Q}.
\]
After a primitive branch is chosen, its local exponent may have denominator $8$. This produces a comparison with the Lusztig order $\ell=8$ and with the Mukai integer $2c_+=8$. The formal parameter itself is not fixed by this comparison: if $q=\exp(\hbar)$, then $\hbar$ is a logarithm of $\zeta$ and is defined only modulo $2\pi i\mathbb{Z}$.
\end{theorem}

\begin{proof}
The Bruinier calculation is a calculation on the automorphic parameter space. The Lusztig construction is a root-of-unity specialisation of the exponentiated quantum parameter. A comparison between them must be supplied by the Hall--Borcherds recognition data. Without this comparison map, the two order-eight records are independent records with the same integer.
\end{proof}

\begin{theorem}[Recognised characteristic invariant: the Igusa cusp form]
\label{thm:qca-classification-invariant}
On the finite Hall--Borcherds recognition locus of
Theorem~\ref{thm:qca-etingof-kazhdan-super-classification}, the
completed source deformation complex has a Bruinier--Heegner
characteristic line
\[
 \chi^{\mathrm{cmp}}_{\Delta_5}\bigl(\Def_{\mathrm{src}\to\Delta_5}(\mathbf H)\bigr)
 \subset
 H^2_{\mathrm{Lie}}\bigl(\mathfrak q_{K(1)};\mathrm{grt}_1\bigr)
 \;\cong\;
 M^{!,\mathrm{odd}}_{5}(K(1))\,/\,(\text{periods})
 \;\cong\;
 \mathbb{C}\cdot\Delta_5.
\]
Two recognised Hall--Drinfeld packages with the same finite comparison
matrices and the same image in this line give gauge-equivalent
Etingof--Kazhdan quantisations. The Igusa cusp form is the scalar
automorphic characteristic of the recognised source; it is not, by
itself, an intrinsic classification of all quasi-Hopf structures on
$\mathfrak{g}_{\Delta_5}$.
\end{theorem}

\begin{proof}
Apply the conditional comparison theorem of
Chapter~\ref{ch:k3-chiral-bialgebra-platonic} to a completed source
package which has already passed the finite Hall--Borcherds gates. The
Etingof--Kazhdan super-category extension (Etingof 2002, \S 6.5;
Gurevich--Majid 1994 for super-Hopf axioms; Etingof--Kazhdan 1996--2000
Parts~I--V of the quantisation theorem) transports the fixed
super-Manin input to a quasi-Hopf output. Bruinier--Heegner
Chern-class reciprocity and the Gritsenko--Nikulin denominator identify
the scalar automorphic target with the $\Delta_5$ line. The finite
source matrices and PBW/no-extra-relations checks are the source
recognition; the automorphic line is only the characteristic target.
\end{proof}

\begin{remark}[Seven lenses on one object]
\label{rem:qca-seven-lenses}
The K3 chiral Hall--Drinfeld double is not seven separate facts but one
fact seen through seven lenses; they are enumerated here and displayed
simultaneously.
\begin{enumerate}[label=\textup{(L\arabic*)}]
\item \textbf{Gritsenko additive $\leftrightarrow$ Weyl--Kac--Borcherds denominator.}
 $\mathrm{Grit}(\eta^9\vartheta_1) = \Delta_5$ is the Fourier--Jacobi face of the
 Weyl--Kac character formula for $\mathfrak{g}_{\Delta_5}$
 (Chapter~\ref{ch:k3e-bkm}, Remark~\ref{rem:k3e-gritsenko-seeds-WKB}).
\item \textbf{Maass $\mathbb{Z}/2$-spin cover $\leftrightarrow$ Jacobi half-integral-weight paramodular.}
 $\widehat{\mathrm{Sp}_4(\mathbb{Z})}^{v_{\Delta_5}}$ is the paramodular metaplectic
 cover where half-integral Jacobi-index theory lives
 (Theorem~\ref{thm:Maass-spin-cover}).
\item \textbf{Arthur packet $\psi_{\Delta_{10}}$ $\leftrightarrow$ $R_{\mathrm{Sieg,dyn}}$ via Langlands dual group.}
 The Arthur packet's Langlands dual group at the spin cover is $\mathrm{Sp}_4$ itself,
 and the local $p$-adic Hecke category is the Kazhdan--Lusztig category of the
 $p$-adic loop group, with $R_{\mathrm{Sieg,dyn}}$ its KL $R$-matrix
 (Chapter~\ref{ch:k3-chiral-bialgebra-platonic},
 Remark~\ref{rem:arthur-packet-langlands-dual-R-matrix}).
\item \textbf{Pasol--Zagier $H_2$ Kronecker--Eisenstein $\leftrightarrow$ elliptic dynamical YBE on $K(1)$.}
 The Siegel Kronecker--Eisenstein cyclic identity is simultaneously the classical $r$-matrix
 Yang--Baxter equation and the pentagon Maurer--Cartan defect
 (Chapter~\ref{ch:k3-chiral-bialgebra-platonic},
 Remarks~\ref{rem:pasol-zagier-H2-associator}
 and~\ref{rem:elliptic-dynamical-YBE-K1}).
\item \textbf{Shimura--Waldspurger $\leftrightarrow$ Ikeda-vs-Gritsenko distinction.}
 The Gritsenko and Borcherds lifts produce the same $\Delta_5$; the Shimura--Waldspurger
 correspondence is the automorphic witness. Ikeda produces $\Delta_{10}$ only; the
 Gritsenko half-integer-index lift has no Ikeda analogue
 (Chapter~\ref{ch:k3-chiral-bialgebra-platonic},
 Remark~\ref{rem:ikeda-gritsenko-disambiguation}).
\item \textbf{Humbert monodromy $\leftrightarrow$ Lusztig root-of-unity specialisation.}
 The order-eight Humbert datum, the Mukai integer \(2c_+(\mathrm{Mukai}(K3))=8\),
 and the Lusztig order \(\ell=8\) are compared on the fixed K3 branch.
 Bruinier's Heegner Chern-class formula supplies the automorphic side only.
\item \textbf{Conditional characteristic $\leftrightarrow$ 1-loop-anomaly output.}
 After finite Hall--Borcherds source recognition, the comparison class maps to
 the one-dimensional Bruinier--Heegner characteristic line
 $\mathbb{C}\cdot\Delta_5$. The same scalar $\Delta_5$ is the
 1-loop anomaly-cancellation output of twisted 11D supergravity on
 $K3\times T^2$ under the protected comparison hypothesis
 (Chapter~\ref{ch:k3-chiral-bialgebra-platonic},
 Theorems~\ref{thm:k3-classification-etingof}
 and~\ref{thm:k3-1loop-output}).
\end{enumerate}
Each lens refracts the same underlying object; the interlocks between them
are mathematical. Lens~(L1) seeds imaginary-root multiplicity
(Humbert monodromy, lens~L6); lens~(L2) seeds the paramodular
$K(1)$-restriction that lens~(L4) requires for the hexagon to close;
lens~(L3) seeds the Langlands-reciprocity characteristic comparison
that lens~(L7) completes after source recognition. No single lens is
sufficient; all seven are necessary. The
unity of the object is the joint coherence of the seven lenses.
\end{remark}

\section{Source-recognised \texorpdfstring{$\mathbf H_{\Delta_5}$}{H-Delta-5} as quantum chiral algebra}
\label{sec:qca-hdelta5-quantum-chiral-algebra}

\begin{remark}[Source-recognised \texorpdfstring{$\mathbf H_{\Delta_5}$}{H-Delta-5} as the fourth taxonomic slot]
\label{rem:qca-hdelta5-fourth-slot}
The quantum chiral algebra construction surveyed above (Drinfeld--Jimbo,
Drinfeld Yangian, RTT in the chiral setting) reaches its K3 master
form only on the source-recognised locus where the finite
Hall--Borcherds gates construct the non-abelian chiral bialgebra
$\mathbf H_{\Delta_5}$. As a quantum chiral algebra on that locus,
$\mathbf H_{\Delta_5}$ has:
\begin{itemize}
\item \emph{Coalgebra side:} bar complex $B_{\mathrm{ch}}(\mathbf H_{\Delta_5})$
 in the Hall subcategory of $\mathrm{CoHA}_{K3\times E}$;
\item \emph{Algebra side:} cobar $\Omega_X(B_{\mathrm{ch}}(\mathbf H_{\Delta_5}))\simeq\mathbf H_{\Delta_5}$
 on the chosen order-eight root-of-unity branch;
\item \emph{$R$-matrix side:} Siegel-elliptic dynamical $R_{\mathrm{Sieg,dyn}}(z_1,z_2;\tau,w)$
 satisfying the dynamical Yang-Baxter equation on $\bar{\mathcal A}_2$;
\item \emph{Associator side:} twisted Siegel-Borcherds associator
 $\widetilde\Phi^{\mathrm{Sieg\text{-}Bor}}_{\mathrm{Sp}_4}[\Phi_{10}^{\mathrm{un}}/\eta^{24}]$
 satisfying the pentagon at order $\hbar^{\le 3}$ with coboundary
 $\phi^{(3)} = \zeta(3) c_{\mathrm{symm}} + \tfrac{25}{3} c_{\mathrm{timelike}} + (\Phi_{10}^{\mathrm{un}}/\eta^{24}) c_{\Phi_{10}^{\mathrm{un}}}$.
\end{itemize}
These are the four target-facing data of the recognised package for
$\mathbf H_{\Delta_5}$. The order-eight comparison fixes a branch after
the source package has been supplied; it does not replace the Hall
source construction. Cross-ref
Chapter~\ref{ch:k3-chiral-bialgebra-platonic}.
\end{remark}

\begin{remark}[\texorpdfstring{$\Psi$}{Psi} as a functor of operads: operadic compatibility]
\label{rem:qca-psi-functor-of-operads}
\index{$\Psi$!functor of operads}
\index{operadic compatibility!Psi}
The universal-trace functor $\Psi$ of
Chapter~\ref{ch:phi-universal-trace-platonic},
Theorem~\ref{thm:phiutp-psi-functor-lax-symmon}, extends the recognised
source-package rule $(L, \phi_L) \mapsto \mathbf{H}_{\phi_L}$ to a lax
symmetric-monoidal functor of symmetric monoidal $(\infty,1)$-categories
\[
 \Psi\colon (\mathrm{CY}^{\mathrm{Siegel\text{-}aut}}_2, \oplus)
 \longrightarrow
 (\mathrm{QHopf}^{\mathrm{BKM}}, \otimes),
\]
with Jacobi-form K\"unneth $\phi_1 \boxtimes \phi_2$ on the domain, and
Drinfeld--Reshetikhin--Turaev tensor product $R_1 \otimes R_2$,
$\Phi_1 \otimes \Phi_2$ on the codomain. Applied at the K3 input,
$\Psi(\Lambda^{3,2}, \phi_{K3}) = \mathbf{H}_{\Delta_5}$ only after
the finite Hall--Borcherds source gates and Bruinier--Heegner
characteristic comparison have been supplied; this is the quantum
chiral algebra of
Remark~\ref{rem:qca-hdelta5-fourth-slot}. The operadic compatibility of
$\Psi$ implies that $\mathbf{H}_{\Delta_5}$ is not an isolated construction
but the recognised value at one object of a functor respecting composition:
the four data (coalgebra bar, algebra cobar, $R$-matrix, associator) listed
in Remark~\ref{rem:qca-hdelta5-fourth-slot} are functorial in $L$ after
source recognition,
with direct-sum $L_1 \oplus L_2$ mapping to tensor product
$\mathbf{H}_{\phi_1} \otimes \mathbf{H}_{\phi_2}$ of the four-fold data
up to natural isomorphism. The rank-$4$ BKM instance
(Proposition~\ref{prop:phiutp-rank-four}) identifies this at
$\mathrm{II}_{1,1} \oplus \mathrm{II}_{1,1} = \mathrm{II}_{2,2}$ giving
$\mathbf{H}_{\mathrm{Monster}} \otimes \mathbf{H}_{\mathrm{Monster}}$.
Primary-lit: Borcherds 1998 \S14 (K\"unneth for singular theta
correspondence on lattice direct sums), Dunn 1988 Theorem~3.4
(additivity $E_1 \otimes_{\mathrm{BV}} E_1 \simeq E_2$),
Reshetikhin--Turaev 1991 Proposition~1.5, Lurie \emph{Higher Algebra}
5.3.2.1.
\end{remark}

\begin{remark}[\texorpdfstring{$E_2 \to E_1$}{E-2 -> E-1} Dunn stabilisation of \texorpdfstring{$\Psi$}{Psi} on the \texorpdfstring{$\mathbf H_{\Delta_5}$}{H-Delta-5} branch]
\label{rem:qca-psi-dunn-stabilisation-hd5}
\index{Dunn stabilisation!Psi E_2 to E_1}
\index{Psi!Dunn stabilisation}
On the source-recognised quantum-chiral-algebra branch that hosts
$\mathbf H_{\Delta_5}$, the operadic compatibility of $\Psi$ takes the
specific form of an $E_2 \to E_1$ Dunn stabilisation
(Theorem~\ref{thm:phiutp-e2-to-e1-dunn}):
the domain carries an $E_2$-structure via Fubini on the two commuting
pairings $(\oplus, \boxtimes)$, while the codomain
$(\mathrm{QHopf}^{\mathrm{BKM}}, \otimes)$ carries only an
$E_1$-structure because Eckmann--Hilton would force any second
compatible pairing on $\mathrm{QHopf}$ to collapse each object to a
commutative algebra, contradicting non-abelian BKM. The forgetful
$E_2 \to E_1$ collapse is \emph{not} information loss: the coalgebra
structure on $\mathbf H_{\Delta_5}$ (equivalently, the Jacobi-index
coupling on $\phi_{K3}$) is absorbed \emph{internal} to each BKM output
object, rather than surviving as a second external pairing on
$\mathrm{QHopf}^{\mathrm{BKM}}$. The operadic compatibility therefore
translates: external $E_2$ on the input (lattices + Jacobi forms) becomes
internal coalgebra + external $E_1$ on the output (quantum chiral
algebras, tensor product). This is the precise operadic content of the
Koszul reflection acting through $\Psi$: the reflection halves the
operadic degree as it pushes from the CY-Siegel side to the BKM-chiral
side, consistent with the $[d]$-shift on $\Phi$ of
Theorem~\ref{thm:phi-output-shift-k3-2}.
\end{remark}

\begin{remark}[Etingof-Kazhdan Part V after source recognition]
\label{rem:qca-2}
The Etingof--Kazhdan 2007 Part~V super-category extension of the
quantisation theorem (for Lie bialgebras over a Grothendieck--Witt
ring) applies at the K3 base point only after the compact source package
has supplied its finite Hall--Borcherds recognition data. It then gives
a characteristic comparison
\[
H^2_{\mathrm{Lie}}\bigl(\mathfrak q_{K(1)};\mathrm{grt}_1\bigr)
\;\cong\; \mathbb C\cdot \Delta_5,
\]
where the right side is the scalar automorphic target attached to
$\Delta_5$, the primitive square root of the unnormalised Igusa form
$\Phi_{10}^{\mathrm{un}}$. This statement does not compute
the intrinsic second Lie-bialgebra cohomology of
$\mathfrak g_{\Delta_5}$ and does not make $\mathbf H_{\Delta_5}$
unique from automorphy alone. The uniqueness
assertion is conditional: same source window, same comparison matrices,
same radical/PBW/no-extra-relations gates, and same
Bruinier--Heegner characteristic image. Primary-lit:
Etingof-Kazhdan 2007 Part~V, Morava 1985, Hopkins-Kuhn-Ravenel 2000.
\end{remark}


%% 6D holomorphic Chern--Simons: classical BV datum, quantum observables,
%% anomaly, field-strength renormalisation, minimal model, dualizability, and the
%% all-orders non-abelian Yangian compound.

\section{6D holomorphic Chern--Simons: BV datum, quantum observables, dualizability}
\label{sec:hcs-bv-quantum-dualizability}

The holomorphic Chern--Simons theory on a Calabi--Yau $d$-fold $X$ carries a
classical BV datum whose field-pairing shift is $\mathrm{shift}=d-4$,
a quantum BV observable algebra whose $E_3$-structure on $\C^3$ is
witnessed by an explicit Bochner--Martinelli propagator, a corrected
one-loop anomaly split in which $C_2(\fg)$ controls field-strength
renormalisation while the dimension-$3$ holomorphic gauge anomaly is
quartic, and a dualizability profile governed by $E_3$-Koszul
self-duality. These structures give the all-orders identification
$\partial\hCS_5(\fg) \simeq Y_{\epsilon_1,\epsilon_2,\epsilon_3}(\widehat\fg)$
for simply-laced $\fg$, together with the corner identification at finite
$n$ and the disambiguation of the three chiral-Yangian variants. The
seven results below are stated in the lane in which their proof
works; $(\infty,1)$-categorical inputs and chain-level computations are
both load-bearing.

\subsection{Classical BV datum on a Calabi--Yau \texorpdfstring{$d$}{d}-fold}
\label{subsec:hcs-classical-bv}

\begin{theorem}[Classical BV structure of \texorpdfstring{$\hCS$}{hCS} on a CY\texorpdfstring{$_d$}{-d}]
\label{thm:plat-hCS-classical}
Let $X$ be a smooth Calabi--Yau $d$-fold with holomorphic volume form
$\Omega_X \in H^{d,0}(X)$, and let $\fg$ be a finite-dimensional Lie
algebra with nondegenerate invariant pairing
$\langle-,-\rangle_\fg$. The classical BV datum of holomorphic
Chern--Simons theory on $X$ with gauge Lie algebra $\fg$ consists of:
\begin{enumerate}[label=\textup{(\alph*)}]
\item the super-field
\[
 \cA \;\in\; \Omega^{0,\bullet}(X,\fg)[1],
\]
whose $\Omega^{0,p}$ component has BV-degree $1-p$. In complex
dimension $3$ this is
$c + A_{0,1} + A^*_{0,2} + c^*_{0,3}$ with degrees
$+1,0,-1,-2$;
\item the classical action
\[
 S_{\mathrm{cl}}(\cA)
 \;=\; \int_X \Omega_X \wedge
 \Big\langle \cA,\; \bar\partial \cA + \tfrac{1}{3}[\cA,\cA]\Big\rangle_\fg;
\]
\item a $(d{-}4)$-shifted symplectic structure on the space of BV
fields, induced from Serre duality pairing
$\Omega^{0,p} \otimes \Omega^{0, d-p} \to \Omega^{0,d} \to \C$.
\end{enumerate}
The shift law
\[
 \mathrm{shift}(d) \;=\; d - 4,
 \qquad
 E_n(d) \;=\; E_{d},
\]
organises the operadic level of the BV observable algebra by dimension:
\[
 (d,\,\mathrm{shift},\,E_n)
 \;\in\;
 \bigl\{(2,-2,E_2),\ (3,-1,E_3),\ (4,0,E_4),\ (5,+1,E_5)\bigr\}.
\]
This is the BV-field pairing shift, not the PTVV shift on the moduli
stack of objects. Thus at $d=5$ the BV field space has a
$+1$-shifted symplectic pairing, while the PTVV form on
$\cM(D^b\Coh(X))$ has shift $2-d=-3$.
\end{theorem}

\begin{proof}
The field complex $\Omega^{0,\bullet}(X,\fg)[1]$ and the cubic
functional are Costello's holomorphic Chern--Simons BV datum; the
Serre pairing
$\Omega^{0,p}\otimes\Omega^{0,d-p}\to\Omega^{0,d}\to\C$ has total
field degree $d-4$ after the shift by $[1]$ on each input. Local
observables of a holomorphic theory on a complex $d$-fold form an
$E_d$-algebra by the Costello--Gwilliam factorisation formalism, and
the $d=3$ entry is the usual $(-1)$-shifted BV symplectic structure
used in the quantisation of hCS on $\C^3$.
\end{proof}

\subsection{Quantum observables on \texorpdfstring{$\C^3$}{C-3}: \texorpdfstring{$E_3$}{E-3}-structure and propagator}
\label{subsec:hcs-quantum-observables}

\begin{theorem}[Quantum observables of \texorpdfstring{$\hCS$}{hCS} on \texorpdfstring{$\C^3$}{C-3} and the \texorpdfstring{$E_3$}{E-3}-algebra structure]
\label{thm:plat-hCS-quantum}
Let $Y = \C^3$ with holomorphic volume form $\Omega = dz_1 \wedge dz_2 \wedge dz_3$ and let $\fg$ be semisimple.
Costello renormalisation (arXiv:$1303.2632$; Costello--Gwilliam~$2017$ Vol~II~\S$5$)
produces a factorisation algebra of quantum observables
\[
 \Obs_{\hCS}(\C^3)
 \;=\; \bigl(\Sym(\cE^\vee[1])[[\hbar]],\; Q + \hbar\Delta\bigr)
\]
with propagator
\[
 P_{\mathrm{BM}}(z,w)
 \;=\; \frac{2}{(2\pi i)^3}
 \sum_{k=1}^{3} (-1)^{k-1}\, \overline{(z_k - w_k)}\; \|z - w\|^{-6}\;
 \widehat{d\bar{z}_k} \wedge dw_1 \wedge dw_2 \wedge dw_3,
\]
the Bochner--Martinelli kernel on $\C^3 \setminus \Delta$, and OPE
$\cA(z)\cA(w) \sim \hbar\, P_{\mathrm{BM}}(z,w) \cdot \mathbf{1}$
modulo $Q$-exact terms. The observables carry a genuine
$E_3$-algebra structure in the Dolbeault chain complex:
\[
 \Obs_{\hCS}(\C^3)
 \;\simeq\;
 \CE^\bullet_{\bar\partial,\chir}(\cE_{\hCS},\,\cO_{\C^3}),
\]
with $E_3$-composition realised by a sum-over-shuffles prescription with
Koszul signs, associativity enforced by \v{C}ech--Dolbeault Mayer--Vietoris
on $\overline{\mathrm{Conf}}_n(\C^3)$, and commutativity by
$\pi_1(\mathrm{Conf}_2(\C^3)) = \pi_1(S^5) = 0$.
\end{theorem}

\begin{proof}
\emph{Propagator and OPE.}
The kernel $P_{\mathrm{BM}}$ is the unique
$\partial$-harmonic representative of the diagonal class in
$H^{d-1, d-1}(\C^d \setminus \Delta)$ at $d = 3$ (Range~$1986$ \S$6.1$);
direct computation gives
$\bar\partial P_{\mathrm{BM}} = \delta_\Delta$. The Costello--Li
$\bar\partial$-regularisation (arXiv:$1601.04040$ \S$4$) realises
$P_{\mathrm{BM}}$ as the heat-kernel limit
$\lim_{L \to \infty} \int_0^L K_t\, dt$ of the Dolbeault heat kernel on $\C^3$,
which in turn supplies the OPE through the counterterm construction of
Costello--Gwilliam Vol~II Thm.~$9.3.1$.

\emph{$E_3$-structure.}
Shuffles on ordered insertions of points in $\mathrm{Conf}_n(\C^3)$ with
Koszul signs produce a chain-level action of the Fulton--MacPherson little
$3$-disks operad on $\Obs_{\hCS}(\C^3)$, by Kontsevich--Tamarkin formality
of $\cD_3$ (Tamarkin~$2003$; Kontsevich~$1999$ Thm.~$1$) composed with the
Gwilliam--Williams~$2021$ comparison $E_d^{\mathrm{hol}} \simeq E_d$
(arXiv:$2009.05037$ Thm.~$2.5.5$).
Mayer--Vietoris on the strata of the Fulton--MacPherson compactification
$\overline{\mathrm{Conf}}_n(\C^3)$ supplies the associator cocycle; the
resulting pairing
$\mu_n \colon \Obs^{\otimes n} \to \Obs$ is strictly associative on
cohomology and $A_\infty$-associative on chains with higher
$\mu_{n \geq 3}$ computed by shuffle-type integrals against the
Bochner--Martinelli kernel.

\emph{Commutativity.}
The link of the diagonal in $\mathrm{Conf}_2(\C^3)$ is $S^5$ (removing
a point from $\C^3 \setminus \{0\}$ and deformation retracting);
$\pi_1(S^5) = 0$, so the $E_3$-operad trivialises braid monodromy at the
topological level. The nontriviality of the chain-level $E_3$-structure
is supplied by the holomorphic deformation: the $\bar\partial$-grading
breaks the symmetric action of $S_n$ on $\mu_n$ to the $E_3$-action of the
framed little $3$-disks.

\emph{CE-chiral identification.}
Costello--Gwilliam Vol~II Thm.~$6.9.0.1$ realises
$\Obs_{\hCS}(\C^3)$ as the holomorphic Chevalley--Eilenberg complex of the
$L_\infty$-algebra $\cE_{\hCS} = \Omega^{0,\bullet}(\C^3, \fg)[1]$ valued
in the commutative ring $\cO_{\C^3}$; the isomorphism is the chain-level
content of the BV quantisation.
\end{proof}

\subsection{Anomaly factorisation}
\label{subsec:hcs-anomaly}

\begin{theorem}[Corrected anomaly split for \texorpdfstring{$\hCS$}{hCS} on CY\texorpdfstring{$_3$}{-3}]
\label{thm:plat-anomaly}
Let $X$ be a compact Calabi--Yau $3$-fold with holomorphic volume form
$\Omega_X$ and topological Euler characteristic $\chi_{\mathrm{top}}(X)$,
and let $\fg$ be a semisimple Lie algebra. In complex dimension $d$, the
local holomorphic gauge anomaly is controlled by an invariant polynomial
of degree $d+1$ in the gauge curvature. Thus for $d=3$ the cohomological
anomaly is quartic, schematically
\[
 \kanom(X,\fg)
 \;=\;
 \hbar \int_X \Omega_X \wedge
 I_4(\cA,F_\cA,F_\cA,F_\cA),
\]
with $I_4 \in \mathrm{Sym}^4(\fg^\vee)^{\fg}$ in the chosen
representation. The quadratic-Casimir coefficient $C_2(\fg)$ belongs to
field-strength renormalisation and to BV-trivial counterterms; it is not
the anomaly class.
Consequences:
\begin{enumerate}[label=\textup{(\roman*)}]
\item Vanishing of the cubic tensor $d^{abc}$ for
$\mathfrak{su}(2)$, $\mathfrak{so}(N)$, $E_6,E_7,E_8,F_4,G_2$ is not an
anomaly-cancellation theorem for $6$D $\hCS$; the quartic invariant
$I_4$ must be evaluated.
\item On flat $\C^3$ the local anomaly is absent in the translation
invariant perturbative model. On $K3\times E$ the
Euler-characteristic/BCOV coefficient vanishes because
$\chi_{\mathrm{top}}(K3\times E)=24\cdot0=0$.
\item The quintic, with $\chi_{\mathrm{top}}(X_5)=-200$, remains the
first compact instance for the quartic anomaly. The cubic symmetric
tensor $d^{abc}$ does not determine this coefficient.
\end{enumerate}
\end{theorem}

\begin{proof}
Costello's heat-kernel regularisation supplies the perturbative
renormalisation framework. The four-dimensional gauge-anomaly pattern
(triangle diagram, cubic Casimir) does not transfer to holomorphic
dimension~$3$: the local Chern--Weil degree of the hCS anomaly is
$d+1=4$. A full compact CY$_3$ quantisation theorem with the quartic
coefficient fixed is not asserted here; the statement records the
correct obstruction degree and separates it from the $C_2(\fg)$
field-strength counterterm.
\end{proof}

\subsection{The bar complex as BV-BRST complex: identification and anomaly universality}
\label{subsec:phys-bv-bar-identification}

%% Absorption from physics_bv_brst_cy.tex:
%% The identifications in this subsection--QME = MC for Theta_A,
%% bar differential = Q_BRST, curvature = BRST anomaly--are not
%% present elsewhere in this chapter. They are recorded here as
%% structural facts, with the physics note as primary source and
%% primary-literature attribution.

The holomorphic Chern--Simons BV datum of Theorem~\ref{thm:plat-hCS-classical} carries an open-string interpretation: for a brane $\mathcal{F} \in \mathrm{Ob}(\mathcal{C})$ in the CY category $\mathcal{C}$, the open string field is
\[
 \alpha \;\in\; \mathfrak{F} \;:=\; \mathrm{RHom}_\mathcal{C}(\mathcal{F}, \mathcal{F})[1],
\]
and the BV action is the CY-trace functional
\begin{equation}
\label{eq:phys-bv-bv-action}
 S(\alpha) \;=\; \sum_{n \geq 2} \frac{1}{n}\, \mathrm{Tr}\bigl(\mu_n(\alpha, \ldots, \alpha)\bigr),
\end{equation}
where $\mu_n$ are the $A_\infty$-operations and $\mathrm{Tr} \colon A \to k[-d]$ is the CY trace. The CY pairing $\langle -,- \rangle \colon A \otimes A \to k[-d]$ supplies the BV symplectic form.

\begin{proposition}[Classical master equation = \texorpdfstring{$A_\infty$}{A-inf} relations]
\label{prop:phys-bv-classical-me}
\ClaimStatusProvedHere
The BV action \eqref{eq:phys-bv-bv-action} satisfies the classical master equation
\[
 \{S, S\} = 0
\]
with respect to the BV antibracket $\{-,-\}$ defined by inverting the CY pairing. This is equivalent to the $A_\infty$-relations for $(\mathrm{End}_\mathcal{C}(\mathcal{F}), \{\mu_n\})$.
\end{proposition}

\begin{proof}
The classical master equation $\{S, S\} = 0$ expands, using the BV antibracket from the CY pairing, as
\[
 \sum_{p+q = n+1} \mathrm{Tr}\bigl(\mu_p(\alpha^{\otimes (p-1)}, \mu_q(\alpha^{\otimes q}))\bigr) = 0
 \qquad \text{for all } n \geq 1.
\]
This is the $A_\infty$-relation $\sum_{p+q=n+1} \mu_p \circ (\mathrm{id}^{\otimes i} \otimes \mu_q \otimes \mathrm{id}^{\otimes (p-1-i)}) = 0$ contracted against $\alpha^{\otimes n}$ and traced. Non-degeneracy of the CY pairing shows the traced and untraced $A_\infty$-relations are equivalent.
\end{proof}

\begin{theorem}[Quantum master equation = Maurer--Cartan equation for \texorpdfstring{$\Theta_A$}{Theta-A}]
\label{thm:phys-bv-qme-mc}
Let $A = A_\mathcal{C} = \Phi(\mathcal{C})$ be the quantum chiral algebra associated to the CY category $\mathcal{C}$, and let $\Theta_A$ be the universal Maurer--Cartan element of $A$ (Vol~I). The quantum master equation
\[
 \hbar\,\Delta S + \tfrac{1}{2}\{S, S\} = 0
\]
for the BV action $S$ of the open string field theory on any brane $\mathcal{F} \in \mathcal{C}$ is equivalent to the Maurer--Cartan equation
\[
 d\Theta_A + \tfrac{1}{2}[\Theta_A, \Theta_A] \;=\; \mathcal{O}_{\mathrm{full}}(A)
\]
for $\Theta_A$ in the chiral deformation complex of $A$. Its scalar
Hodge projection is $\operatorname{pr}_{\mathrm{sc}}\mathcal{O}_{\mathrm{full},g}(A)
= \kappa_{\mathrm{ch}}(A)\lambda_g$: unconditionally at genus~$1$, and at
all genera on the uniform-weight scalar lane. The BV-to-chiral dictionary is:
\begin{center}
\renewcommand{\arraystretch}{1.3}
\begin{tabular}{lll}
 \textbf{BV-BRST} & \textbf{Chiral bar complex} & \textbf{Identification} \\
 \hline
 Field $\alpha \in \mathfrak{F}$ & Bar element $\alpha \in B(A)$ & Same underlying vector space \\
 BV action $S(\alpha)$ & MC element $\Theta_A$ & $S = \mathrm{Tr}(\Theta_A)$ \\
 CY pairing $\langle -,- \rangle$ & Factorisation coproduct & $\omega_{\mathrm{BV}} = \Delta_B^*\omega_{\mathrm{CY}}$ \\
 $\{S,S\} = 0$ (CME) & $[\Theta_A, \Theta_A] = 0$ at genus $0$ & Classical $A_\infty$ relations \\
 $\hbar\,\Delta S$ & $d\Theta_A$ at genus $\geq 1$ & One-loop anomaly \\
 QME obstruction & $\mathcal{O}_{\mathrm{full},g}(A)$, scalar projection $\kappa_{\mathrm{ch}}(A)\cdot\omega_g$ & Genus-$g$ anomaly
\end{tabular}
\end{center}
Expanding in powers of $\hbar$ (genus):
\begin{itemize}
 \item \emph{Genus $0$}: $\{S,S\} = 0$, the $A_\infty$-relations (Proposition~\ref{prop:phys-bv-classical-me}). On the chiral side, $[\Theta_A, \Theta_A]|_{g=0} = 0$: the bar complex is a differential coalgebra at tree level.
 \item \emph{Genus $1$}: the one-loop anomaly $\Delta S_2 + \{S_3, S_2\} + \cdots$ equals $\kappa_{\mathrm{ch}}(A)\cdot\lambda_1$, where $\lambda_1 = c_1(\mathbb{E})$ is the Hodge class on $\overline{\mathcal{M}}_{1,1}$.
 \item \emph{Genus $g$}: the full genus-$g$ obstruction is
 $\mathcal{O}_{\mathrm{full},g}(A)$. Its scalar component is
 $\mathrm{obs}^{\mathrm{sc}}_g(A)=\kappa_{\mathrm{ch}}(A)\cdot\lambda_g$
 on the uniform-weight scalar lane (Vol~I; cf.~equation~\eqref{eq:phys-anom-obs-g}
 below), while multi-weight families at $g\geq2$ carry cross-channel
 summands.
\end{itemize}
\end{theorem}

\begin{proof}
The identification $S = \mathrm{Tr}(\Theta_A)$ and $\omega_{\mathrm{BV}} = \Delta_B^*\omega_{\mathrm{CY}}$ are term-by-term: the degree-$n$ term of $\Theta_A$ is $\frac{1}{n}\mu_n \in A^{\otimes n}$, and $\mathrm{Tr}$ on $A^{\otimes n}$ applied to $\alpha^{\otimes n}$ gives $\frac{1}{n}\mathrm{Tr}(\mu_n(\alpha,\ldots,\alpha)) = S_n$, the $n$-th term of the BV action. The BV Laplacian $\Delta = \sum_{i,j}\omega^{ij}\partial^2/\partial\phi^i\partial\phi^j$ acting on $S$ contracts a pair of fields via $\omega_{\mathrm{BV}} = \Delta_B^*\omega_{\mathrm{CY}}$; on the chiral side this is $d\Theta_A|_{g=1}$, the one-loop piece of the deformation differential. The genus-$g$ anomaly coefficient $\kappa_{\mathrm{ch}}(A)$ is the scalar projection of $\Theta_A$ onto the Hodge class $\lambda_g$ (Vol~I, Theorem~D), not the full multi-weight component. The equivalence of the full QME with the MC equation follows by induction on genus: at each genus $g$, the BV equation at order $\hbar^{2g-2}$ is precisely the $g$-th component of the full MC equation; applying the scalar projection gives the displayed Theorem~D component, and the remaining non-scalar part is the cross-channel correction.
\end{proof}

\begin{proposition}[Bar differential = BRST operator]
\label{prop:phys-bv-bar-brst}
Under the identification $\mathfrak{F} = A[1]$ and $S = \mathrm{Tr}(\Theta_A)$, the BRST operator $Q_{\mathrm{BRST}} = \{S, -\}$ acting on $\mathfrak{F}$ coincides with the bar differential $d_B$ on $B(A)$. Explicitly, for $\alpha_1 \otimes \cdots \otimes \alpha_n \in A[1]^{\otimes n}$:
\[
 d_B(\alpha_1 \otimes \cdots \otimes \alpha_n)
 = \sum_{i < j} \pm\, \alpha_1 \otimes \cdots \otimes \mu_{j-i+1}(\alpha_i, \ldots, \alpha_j) \otimes \cdots \otimes \alpha_n.
\]
\end{proposition}

\begin{proof}
On the BV side, $Q_{\mathrm{BRST}} = \{S,-\}$ acts on $\mathfrak{F}$ via the antibracket; expanding $\{S_n, -\}$ for each $n$ using the BV antibracket from the CY pairing produces a sum over insertions of $\mu_n$ at consecutive positions, which is precisely the bar differential $d_B$. The Koszul sign rule matches the graded-commutative conventions of the bar complex.
\end{proof}

\begin{theorem}[Curvature = BRST anomaly; Adler--Bardeen universality]
\label{thm:phys-bv-curvature-anomaly}
At genus $g \geq 1$, the bar complex $B(A)$ acquires curvature from the Hodge bundle over $\overline{\mathcal{M}}_{g,n}$. Its scalar Hodge projection is:
\begin{equation}
\label{eq:phys-bv-curvature}
 \operatorname{pr}_{\mathrm{sc}}(d_B^2) \;=\; \kappa_{\mathrm{ch}}(A) \cdot \omega_g,
\end{equation}
where $\omega_g \in H^2(\overline{\mathcal{M}}_{g,1})$ is the Hodge class. The three consequences:
\begin{enumerate}[label=\textup{(\roman*)}]
 \item At genus $g$, the scalar failure of nilpotency
 $\operatorname{pr}_{\mathrm{sc}}(Q_{\mathrm{BRST}}^2)
 = \kappa_{\mathrm{ch}}(A)\cdot\omega_g$ is the scalar BRST anomaly at
 $g$ loops.
 \item The scalar anomaly is \emph{universal}: it depends only on
 $\kappa_{\mathrm{ch}}(A)$ and the Hodge class, not on the particular
 brane $\mathcal{F}$. At $g\geq2$, the full multi-weight obstruction can
 also contain cross-channel summands.
 \item Scalar anomaly cancellation $\kappa_{\mathrm{ch}}(A) = 0$ restores
 $\operatorname{pr}_{\mathrm{sc}}(d_B^2)=0$ at all genera. Full
 nilpotency additionally requires vanishing of the cross-channel
 summand, for example on the uniform-weight lane.
\end{enumerate}
\end{theorem}

\begin{proof}
The genus-$0$ nilpotency $d_B^2 = 0$ is the content of the $A_\infty$-relations (Proposition~\ref{prop:phys-bv-classical-me}). At genus $g \geq 1$, the bar complex is a curved factorisation coalgebra: the scalar curvature is computed by the trace $\mathrm{Tr}(Q_{\mathrm{BRST}}^2)$ on the endomorphism algebra, which by the CY trace axiom picks up $\kappa_{\mathrm{ch}}(A)$; the Hodge class $\omega_g$ is the cohomological class of the Hodge bundle insertion. The brane-independence of $\kappa_{\mathrm{ch}}$ follows from its definition as a scalar projection of $\Theta_A$, which is defined on $A$ (not on any individual brane module): the trace $\mathrm{Tr}$ on $\mathrm{End}(\mathcal{F})$ is fixed by the CY structure, and $\kappa_{\mathrm{ch}}(A)$ is its genus-$1$ shadow. Universality is Theorem~D of Vol~I on the scalar lane; the full $g\geq2$ multi-weight curvature also records cross-channel terms.
\end{proof}

\begin{proposition}[\texorpdfstring{$E_2$}{E-2} enhancement of the bar coalgebra: native at \texorpdfstring{$d=2$}{d=2}, centre-side at \texorpdfstring{$d\geq3$}{d>=3}]
\label{prop:phys-bv-e2-bar}
For a CY$_d$ category $\mathcal{C}$ with $d \geq 2$, the bar construction has the following \(E_2\) scope:
\begin{enumerate}[label=\textup{(\roman*)}]
 \item The $E_1$ direction is the BRST direction: the bar differential $d_B=Q_{\mathrm{BRST}}$ and the deconcatenation coproduct $\Delta$.
 \item At \(d=2\), \(A_\mathcal{C}\) is natively \(E_2\) (Definition~\ref{def:e2-chiral-algebra}), so \(B(A_\mathcal{C})\) is the factorisation-dual \(E_2\) coalgebra.
 \item At \(d\geq3\), \(A_\mathcal{C}\) is natively \(E_1\). The CY \(R\)-matrix \(R_\mathcal{C}(z)=\mathrm{Res}^{\mathrm{coll}}_{0,2}(\Theta_{A_\mathcal{C}})\) equips the Drinfeld-centre side \(\mathcal{Z}(\mathrm{Rep}^{E_1}(A_\mathcal{C}))\), and its bar object, with the braided \(E_2\) structure. It does not make \(B(A_\mathcal{C})\) an intrinsic \(E_2\) coalgebra.
\end{enumerate}
The open-closed datum lives on the \emph{pair}
$(C^\bullet_{\mathrm{ch}}(A,A), A)$: the chiral Hochschild cochain
complex $C^\bullet_{\mathrm{ch}}(A,A)$ is the algebraic closed-sector
object, $A$ is the open/boundary sector, and a physical bulk reading
requires the open--closed comparison coupling
$\mathrm{HH}^\bullet(\mathcal{C})$ to $A$ by deformations of the
$A_\infty$-structure.
\end{proposition}

\begin{proof}
For $d = 2$: the native $E_2$ structure is Theorem~CY-A$_2$
(Chapter~\ref{ch:cy-to-chiral}). The $E_2$ bar coalgebra is the
factorisation dual of the $E_2$-chiral algebra $A_\mathcal{C}$. For
$d \geq 3$: the CY $R$-matrix $R_\mathcal{C}(z)$ is extracted from
$\Theta_{A_\mathcal{C}}$ at degree~$2$
(Construction~\ref{constr:cy-r-matrix}), and its YBE property
(Theorem~\ref{thm:yang-r-matrix-ybe}) gives the half-braiding on the
Drinfeld centre \(\mathcal{Z}(\mathrm{Rep}^{E_1}(A_\mathcal{C}))\).
The centre-side bar object is therefore \(E_2\); the boundary bar
\(B(A_\mathcal{C})\) remains \(E_1\). The open-closed pair structure
is the Swiss-cheese datum $(C^\bullet_{\mathrm{ch}}(A,A), A)$ from
Vol~II; comparison with BV open-closed string field theory requires
the Kontsevich--Soibelman open-closed map.
\end{proof}

\begin{conjecture}[BV-BRST complex of the topological B-model and the bar complex of \texorpdfstring{$G(X)$}{G(X)}]
\label{conj:phys-bv-cy3-bar}
\ClaimStatusConjectured
For a CY$_3$ manifold $X$, assume the compact Hall/centre package defining the quantum vertex chiral group \(G(X)\) has been constructed. The full BV-BRST complex of the topological B-model is then expected to compare with the bar complex \(B(G(X))\). Specifically:
\begin{enumerate}[label=\textup{(\roman*)}]
 \item The bar differential $d_B$ on $B(G(X))$ is the BRST operator $Q_{\mathrm{BRST}}$ of the B-model.
 \item The genus-$0$ cohomology $H^\bullet(B(G(X)), d_B)|_{g=0}$ computes the classical B-model open string states modulo gauge equivalence.
 \item At genus $g \geq 1$, the scalar Hodge projection
 $\operatorname{pr}_{\mathrm{sc}}(d_B^2)
 = \kappa_{\mathrm{ch}}(G(X))\cdot\omega_g$ is the scalar $g$-loop
 BRST anomaly on the chiral lane; at $g\geq2$ the full multi-weight
 obstruction can retain cross-channel summands. The Borcherds scalar is a distinct datum: for $K3 \times E$, $\chi_{\mathrm{top}} = 0$ but $\kappa_{\mathrm{BKM}}(\Delta_5) = 5$, so the BKM weight is not computed by $\chi_{\mathrm{top}}(X)/24$ (Remark~\ref{rem:phys-anom-chi-vs-weight} below).
 \item The bar-cobar adjunction $\Omega(B(G(X))) \simeq G(X)$ (on the Koszul locus) is the statement that the open string field theory determines, and is determined by, the quantum vertex chiral group.
 \item The BKM denominator identity of $\mathfrak{g}_X$ is the BV partition function:
 \[
 Z_{\mathrm{BRST}} \;=\;
 \frac{\det Q_{\mathrm{BRST}}|_{\mathrm{ghost}}}{\det Q_{\mathrm{BRST}}|_{\mathrm{field}}}
 \;=\; \prod_{\alpha \in \Delta_+} (1 - q^\alpha)^{\mathrm{mult}(\alpha)}.
 \]
 For $K3 \times E$, this reproduces the Borcherds product expansion of the Igusa cusp form $\Delta_5$.
\end{enumerate}
\end{conjecture}

\subsection{Modular characteristics: values, anomaly ratios, and the three-candidate problem}
\label{subsec:phys-anom-kappa-table}

%% Absorption from physics_anomaly_cancellation.tex:
%% The master table of kappa_ch values with rho-ratios, the three-candidate
%% problem (chi/24, chi/12, automorphic weight), the anomaly non-cancellation
%% theorem, and the toric genus-1 matching proposition are new to this chapter.

The modular characteristic $\kappa_{\mathrm{ch}}(\mathcal{A})$ of a chiral algebra $\mathcal{A}$ controls the scalar Hodge projection of the genus tower:
\begin{equation}
\label{eq:phys-anom-obs-g}
 \mathrm{obs}^{\mathrm{sc}}_g(\mathcal{A}) \;=\; \kappa_{\mathrm{ch}}(\mathcal{A}) \cdot \lambda_g
 \;\in\; H^{2g}(\overline{\mathcal{M}}_g, \mathbb{Q}),
\end{equation}
where $\lambda_g = c_g(\mathbb{E})$ is the top Chern class of the Hodge bundle. At genus~$1$ this scalar identity is unconditional; at all genera it is the uniform-weight scalar lane of Vol~I, Theorem~D. For multi-weight algebras and $g\geq2$, the full obstruction is $\mathrm{obs}^{\mathrm{full}}_g(\mathcal{A})=\mathrm{obs}^{\mathrm{sc}}_g(\mathcal{A})+\mathrm{obs}^{\mathrm{cross}}_g(\mathcal{A})$. The genus-$1$ free energy is $F_1(\mathcal{A}) = \kappa_{\mathrm{ch}}(\mathcal{A})/24$ (Faber--Pandharipande evaluation $\lambda_1^{\mathrm{FP}} = 1/24$).

The \emph{anomaly ratio} $\varrho(\mathcal{A}) := \kappa_{\mathrm{ch}}(\mathcal{A})/c(\mathcal{A})$ encodes how much of the conformal anomaly survives in the genus tower. The master table of standard chiral algebras:
\begin{center}
\renewcommand{\arraystretch}{1.3}
\begin{tabular}{lccc}
\textbf{Algebra} & $c$ & $\kappa_{\mathrm{ch}}$ & $\varrho$ \\ \hline
$\mathcal{H}^{\oplus d}$ (Heisenberg, rank $d$, level $1$) & $d$ & $d$ & $1$ \\
$\widehat{\mathfrak{g}}_k$ (Kac--Moody) &
 $k\dim\mathfrak{g}/(k+h^\vee)$ &
 $(k+h^\vee)\dim\mathfrak{g}/(2h^\vee)$ &
 $(k+h^\vee)^2/(2kh^\vee)$ \\
$\mathrm{Vir}_c$ (Virasoro) & $c$ & $c/2$ & $1/2$ \\
$\mathcal{W}_N^k$ ($\mathfrak{sl}_N$-type) & $c$ & $c(H_N - 1)$ & $H_N - 1$ \\
$bc$ (ghosts, weights $(2,-1)$) & $-26$ & $-13$ & $1/2$ \\
$\beta\gamma$ (symplectic bosons) & $c$ & $c/2$ & $1/2$ \\
$V_\Lambda$ (lattice, rank $d$) & $d$ & $d$ & $1$
\end{tabular}
\end{center}
where $H_N = \sum_{j=1}^N 1/j$ is the $N$-th harmonic number. The Heisenberg algebra is the unique standard family with $\varrho = 1$: $\kappa_{\mathrm{ch}} = k = c$ at rank $d$, level $1$. For the bosonic string: $\kappa_{\mathrm{eff}} = \kappa_{\mathrm{ch}}(\mathrm{Vir}_c) + \kappa_{\mathrm{ch}}(\mathrm{ghost}) = c/2 - 13 = 0$ at $c = 26$.

\begin{remark}[Complementarity: Koszul pairs]
\label{rem:phys-anom-complementarity}
For a Koszul pair $(\mathcal{A}, \mathcal{A}^!)$, the complementarity conductor $\kappa_{\mathrm{ch}}(\mathcal{A}) + \kappa_{\mathrm{ch}}(\mathcal{A}^!)$ depends on the family: it vanishes for Heisenberg, Kac--Moody, and free-field algebras (antisymmetric under Feigin--Frenkel $k \mapsto -k-2h^\vee$); it equals $13$ for Virasoro (Vol~I, Theorem~C); $250/3$ for $\mathcal{W}_3$; $(H_N-1)K_N/2$ for $\mathcal{W}_N^k$. For the $\fgl_1$/$\C^3$ Yangian branch: $\kappa_{\mathrm{ch}} + \kappa_{\mathrm{ch}}^! = -\sigma_2 + \sigma_2 = 0$ (Proposition~\ref{prop:5d-koszul-dual}).
\end{remark}

\begin{proposition}[Toric CY\texorpdfstring{$_3$}{-3} scalar and refined characters]
\label{prop:phys-anom-toric-genus1}
For the toric CY$_3$ \(X=\mathbb C^3\):
\begin{enumerate}[label=\textup{(\roman*)}]
 \item The evaluated rank-one Heisenberg sector has
 \(\kappa_{\mathrm{ch}}^{\mathrm{Heis}}(\mathbb C^3)=1\) and vacuum
 character \(P(q)=\prod_{n\geq 1}(1-q^n)^{-1}\).
 \item The Virasoro subalgebra at \(c=1\) has scalar
 \(\kappa_{\mathrm{ch}}^{\mathrm{Vir}}=1/2\). This is a subalgebra
 invariant, not the scalar CY trace of the \(\C^3\) Hall output.
 \item The DT partition function is the refined character
 \[
  Z^{\mathrm{DT}}(\mathbb C^3)=M(q)
  =\operatorname{PE}\!\left(\sum_{n\geq 1}nq^n\right),
 \]
 so \(\Omega(n)=n\).
 \item The finite \(W_N\) conductor \(c(H_N-1)\) has no limit which
 defines a scalar \(\kappa_{\mathrm{ch}}\) for the \(\C^3\) Hall output.
\end{enumerate}
\end{proposition}

\begin{proof}
The first two assertions are the standard OPE normalisations:
\[
 J(z)J(w)\sim (z-w)^{-2},
 \qquad
 T(z)T(w)\sim \frac{1/2}{(z-w)^4}+\cdots .
\]
The third assertion is the product formula for the Hilbert scheme of
points on \(\C^3\), equivalently the plethystic exponential of the
primitive BPS spectrum. The last assertion follows because
\(H_N-1\to\infty\), while the scalar Heisenberg sector remains rank one.
\end{proof}

\begin{theorem}[Anomaly non-cancellation: independent scalar lanes]
\label{thm:phys-anom-non-cancellation}
For a CY$_3$ $X$, the chiral modular characteristic of a constructed quantum vertex group, the BCOV genus-one scalar, and the Borcherds automorphic weight are independent pieces of data until a comparison theorem identifies them:
\begin{enumerate}[label=\textup{(\roman*)}]
 \item $\kappa_{\mathrm{ch}}(G(X))$ is defined only where the Hall positive half, Drinfeld double, and centre comparison have been constructed.
 \item The BCOV genus-one scalar is controlled by the topological and holomorphic anomaly data of $X$.
 \item The Borcherds weight $w(\Phi_X)$ is an invariant of a chosen automorphic denominator.
 \item The anomaly cancellation condition $\kappa_{\mathrm{eff}} = \kappa_{\mathrm{ch}}(G(X)) + \kappa_{\mathrm{ghost}} = 0$ determines a \emph{critical weight} $w_{\mathrm{crit}} = -\kappa_{\mathrm{ghost}}$, not a critical central charge.
\end{enumerate}
There are three scalar lanes which may agree only on a constructed comparison locus:
\begin{center}
\begin{tabular}{lcl}
\textbf{Lane} & \textbf{Scalar} & \textbf{Agreement condition} \\ \hline
Heisenberg analogy & $\chi_{\mathrm{top}}(X)/24$ & framed free-field comparison \\
BCOV genus-$1$ & BCOV anomaly scalar & genus-one comparison map \\
Automorphic denominator & $w(\Phi_X)$ & BKM recognition theorem
\end{tabular}
\end{center}
No equality among these three numbers is a formal consequence of the symbol \(G(X)\).
\end{theorem}

\begin{proof}[Argument]
The three scalars live on different levels of the construction. The Euler and BCOV scalars are target-space anomaly data. The Borcherds weight is a scalar shadow of an automorphic denominator. The chiral modular characteristic is a property of a constructed vertex object. A theorem may identify two of them only after specifying the comparison arrow and its hypotheses.
\end{proof}

\begin{remark}[Why \texorpdfstring{$\chi(K3 \times E) = 0$}{chi(K3 x E) = 0} but \texorpdfstring{$\kappa_{\mathrm{BKM}}(\Delta_5) = 5 \neq 0$}{kappa-BKM(Delta-5) = 5 != 0}]
\label{rem:phys-anom-chi-vs-weight}
The Euler characteristic $\chi_{\mathrm{top}}(K3 \times E) = \chi_{\mathrm{top}}(K3)\chi_{\mathrm{top}}(E) = 24 \cdot 0 = 0$ belongs to the target-space scalar lane. The Borcherds invariant $\kappa_{\mathrm{BKM}}(\Delta_5) = 5$ is the weight of the denominator $\Delta_5$ in the Gritsenko normalisation. These are different invariants. The sentence \(\chi_{\mathrm{top}}=0\) therefore gives no contradiction and no derivation of \(\kappa_{\mathrm{BKM}}(\Delta_5)=5\).
\end{remark}

\subsection{Non-abelian field-strength renormalisation}
\label{subsec:hcs-field-strength-renormalisation}

\begin{theorem}[One-loop non-abelian field-strength renormalisation of \texorpdfstring{$\hCS$}{hCS} on \texorpdfstring{$\C^3$}{C-3}]
\label{thm:plat-Z-counterterm}
The one-loop bubble of $\hCS$ on $\C^3$ with semisimple gauge algebra
$\fg$ requires the counterterm
\[
 S^{(1)}_{\mathrm{c.t.}}(\cA)
 \;=\;
 -\hbar\, C_2(\fg)\, (4\pi)^{-3}\, \log(L/\varepsilon)
 \int_{\C^3} \Omega \wedge \Tr(\cA\, \bar\partial \cA),
\]
where $C_2(\fg)$ is the quadratic Casimir, $L$ is the IR scale, and
$\varepsilon$ is the UV cutoff. The associated field-strength
renormalisation factor is
\[
 Z^{(1)}_\cA
 \;=\;
 1 - \hbar\, C_2(\fg)\, (4\pi)^{-3}\, \log(L/\varepsilon).
\]
For $\fg = \mathfrak{su}(N)$: $C_2 = N$, and the coefficient of
$\log(L/\varepsilon)$ equals $-\hbar N/(32\pi^3)$.
\end{theorem}

\begin{remark}[Field-strength renormalisation vs anomaly]
\label{rem:plat-Z-vs-anomaly}
The coefficient $-\hbar C_2(\fg)(4\pi)^{-3}\log(L/\varepsilon)$ of the
kinetic counterterm is a \emph{field-strength} renormalisation of $\cA$,
not an anomaly: it rescales the kinetic term but preserves the BV master
equation $\{S, S\} = 0$ once the counterterm is included. The anomaly
$\kanom$ of Theorem~\ref{thm:plat-anomaly} is the obstruction to
solving the quantum master equation at one loop. In complex dimension
$3$ that cohomological anomaly is controlled by the quartic invariant
$I_4\in\mathrm{Sym}^4(\fg^\vee)^\fg$ in the chosen representation, not
by $C_2(\fg)$ and not by the cubic tensor $d^{abc}$. Thus $C_2$ sets
field renormalisation; the quartic Chern--Weil slot sets the local hCS
anomaly. Conflating these invariants misattributes the one-loop
dressing.
\end{remark}

\begin{proof}
The counterterm is extracted by Costello's BV-cohomology regularisation
(Costello~$2011$ \S$5.4$; Costello--Gwilliam Vol~II Thm.~$9.5.0.6$) on
the Bochner--Martinelli propagator of
Theorem~\ref{thm:plat-hCS-quantum}:
the coefficient of the kinetic term receives a
$\log(L/\varepsilon)$ correction proportional to the
quadratic-Casimir trace
$C_2(\fg)\cdot \Tr(\cA \bar\partial \cA)$, with the $(4\pi)^{-3}$
normalisation fixed by the three-dimensional Dolbeault heat-kernel
measure. For $\mathrm{SU}(N)$: $C_2 = N$ in the normalisation
$\Tr(T^a T^b) = \tfrac{1}{2}\delta^{ab}$ (Peskin--Schroeder Appendix~A,
eq.~$(A.94)$).
\end{proof}

\subsection{Minimal \texorpdfstring{$L_\infty$}{L-inf}-model on flat and compact CY\texorpdfstring{$_3$}{-3}}
\label{subsec:hcs-minimal-model}

\begin{theorem}[Minimal \texorpdfstring{$L_\infty$}{L-inf}-model of \texorpdfstring{$\hCS$}{hCS} on \texorpdfstring{$\C^3$}{C-3}]
\label{thm:plat-Linf-minimal}
On flat $\C^3$ with harmonic subspace $\mathcal{H} = \C[z_1, z_2, z_3]
\otimes \fg$, Kontsevich--Soibelman homotopy transfer along the
Bochner--Martinelli propagator of
Theorem~\ref{thm:plat-hCS-quantum} gives a minimal $L_\infty$-model
$(\mathcal{H}, \{\ell_n^{\min}\}_{n \geq 1})$ for
$\cE_{\hCS}(\C^3)$ with
\[
 \ell_1^{\min} \;=\; 0,
 \qquad
 \ell_2^{\min}([f, g]) \;=\; [f, g]_\fg,
 \qquad
 \ell_n^{\min} \;=\; 0 \quad \text{for all } n \geq 3.
\]
On a compact Calabi--Yau $3$-fold $X$, the Atiyah class
\[
 \mathrm{At}(TX) \;\in\;
 H^1\!\bigl(X,\,\Omega^1_X \otimes \mathrm{End}(TX)\bigr)
\]
is the sole formality obstruction: $\mathrm{At}(TX) = 0$ implies
$\ell_n^{\min} = 0$ for $n \geq 3$. In particular:
\begin{enumerate}[label=\textup{(\roman*)}]
\item On $K3 \times E$: $\mathrm{At}(TE) = 0$ (the elliptic curve is a
complex Lie group, so $TE$ is trivial), but the total-space cubic
receptacle does not vanish. Indeed
$H^3(K3\times E,\Omega^3_{K3\times E})$ contains the K\"unneth summand
$H^2(K3,\Omega^2_{K3})\otimes H^1(E,\Omega^1_E)\cong\C$.
Therefore strict Lie formality on $K3\times E$ is not a consequence of
this Hodge group vanishing.
\item On the quintic: $\mathrm{At}(TX) \neq 0$; the Kuranishi cubic
receptacle $H^3(X, \Omega^3_X) = \C$ is nonzero, so the minimal model
carries a nontrivial $\ell_3^{\min}$ representing the CHSW obstruction.
\end{enumerate}
\end{theorem}

\begin{proof}
\emph{Flat case.}
Every rooted binary tree with an internal edge contributes to
$\ell_n^{\min}$ via the composition
$p \circ \mu \circ (h \otimes h) \circ \mu \circ \cdots$
where $h$ is the chain homotopy inverting $\bar\partial$ via
$P_{\mathrm{BM}}$. On the polynomial harmonic subspace $\mathcal{H}$,
$h$ acts as zero (by the Hodge decomposition
$\Omega^{0,\bullet}(\C^3) = \mathcal{H} \oplus \mathrm{im}(\bar\partial)
\oplus \mathrm{im}(\bar\partial^*)$
restricted to polynomial forms: the propagator kills harmonic inputs).
Every tree with an internal edge vanishes; only the binary root tree
$\ell_2$ survives, reproducing the Lie bracket of $\fg$.

\emph{Compact case.}
The Atiyah class $\mathrm{At}(TX)$ is the class of the extension
$0 \to \Omega^1_X \otimes TX \to J^1(TX) \to TX \to 0$; by
Kapranov~$1999$ Thm.~$2.8.1$ (\emph{Compositio}~$115$), the
$L_\infty$-structure of $\Omega^{0,\bullet}(X, TX)$ at the chain level
has cubic operation $\ell_3 = \mathrm{At}(TX) \cdot$ (anti-symmetrised
triple bracket). On $K3 \times E$: the K\"unneth decomposition
$T(K3 \times E) = \pi_1^* TK3 \oplus \pi_2^* TE$ and
$\mathrm{At}(TE) = 0$ (Biswas--Nag~$2002$ for the elliptic-curve case)
reduces the K3-side Atiyah class to a class in
$H^1(K3, \Omega^1 \otimes \mathrm{End})$, and the cubic receptacle
 $H^3(K3 \times E, \Omega^3)$ is nonzero by K\"unneth:
 $H^3(K3 \times E, \Omega^3) = \bigoplus_{p+q=3}
 H^p(K3, \Omega^p_{K3}) \otimes H^q(E, \Omega^q_E)$, and the summand
 $H^2(K3,\Omega^2_{K3})\otimes H^1(E,\Omega^1_E)$ is $\C$.
 Consequently the strict-Lie conclusion for $K3\times E$ requires an
 additional vanishing of the transferred operation, not merely Hodge
 degree counting.
On the quintic: $H^3(X, \Omega^3_X) = \C$ by Serre duality to
$H^0(X, \cO_X) = \C$, and the Atiyah class is nonzero by
Griffiths--Harris~$1978$ Ch.~$5$.
\end{proof}

\subsection{Dualizability, \texorpdfstring{$3$}{3}-duality, and \texorpdfstring{$E_3$}{E-3}-Koszul self-duality}
\label{subsec:hcs-dualizability}

\begin{theorem}[Dualizability of \texorpdfstring{$\Obs_{\hCS}$}{Obs-hCS} and \texorpdfstring{$3$}{3}-duality]
\label{thm:plat-dualizability}
Let $Y = \C^3$ with the quantum observable $E_3$-algebra
$\Obs_{\hCS}(\C^3)$ of Theorem~\ref{thm:plat-hCS-quantum}, and let
$\partial\overline{\mathrm{Conf}}_2(\C^3) \simeq S^5$ be the link of
the fat diagonal in the Fulton--MacPherson compactification. Then:
\begin{enumerate}[label=\textup{(\roman*)}]
\item The $E_3$-trace defined by integration over
$S^5 \subset \partial\overline{\mathrm{Conf}}_2(\C^3)$ is
nondegenerate on $\Obs_{\hCS}(\C^3)$.
\item The abelian sector (gauge algebra $\fg = \fgl_1$) is
$3$-dualisable in the $(\infty,3)$-category of $E_3$-algebras.
\item On flat $\C^3$ with non-abelian semisimple $\fg$, the
$3$-dualizability is obstructed by an infinite-dimensional
Hochschild cohomology: by Gwilliam--Williams~$2021$
(arXiv:$2009.05037$) Prop.~$5.3.2$,
\[
 \HH^0_{E_3}(\Obs_{\hCS}(\C^3),\, \Obs_{\hCS}(\C^3))
 \;=\; \C[[\tau_1, \tau_2, \tau_3]],
\]
which has infinite cohomological dimension, violating the dualizability
finiteness axiom.
\item On a compact Calabi--Yau $3$-fold $X$, the $3$-dualizability is
restored: the Hochschild cohomology becomes finite-dimensional
under the compactness-plus-properness of $\Coh(X)$.
\item The $E_3$-operad is Koszul self-dual up to the shift:
\[
 \cD_3^! \;\simeq\; \mathrm{Lie}[2],
\]
by Fresse~$2017$ \emph{Homotopy of Operads \& Grothendieck--Teichm\"uller Groups}
Vol.~I Thm.~$14.1.A$, recovering the Calaque--Pantev--To\"en--Vaqui\'e
Koszul duality on the Poisson side. The strict
Koszul duality (Gwilliam--Williams~$2021$ \S$5.3$) and the
homotopy Koszul duality (Francis--Gaitsgory~$2012$) agree through
Fresse Thm.~$12.3.A$ and the Positselski coderived/contraderived
transfer (Positselski~$2011$~\S$6$).
\end{enumerate}
\end{theorem}

\begin{proof}
Item~(i) is the nondegeneracy of the $E_3$-trace on the Fulton--MacPherson
boundary, Costello--Gwilliam Vol~II~\S$6$. Item~(ii) is the abelian
specialisation where
$\Obs_{\hCS}(\C^3)|_{\fgl_1} \simeq \Sym(\Omega^{0,\bullet}(\C^3)[1])[[\hbar]]$
is a free $E_3$-algebra on a dualizable input; its
$3$-dualizability is Lurie~\emph{HA}~$5.3.2$ via the cobordism hypothesis
applied to the free $E_3$-algebra.
Item~(iii) is Gwilliam--Williams~$2021$ Prop.~$5.3.2$: the
$E_3$-Hochschild cohomology of the non-abelian free
holomorphic factorization algebra on $\C^3$ is
$\C[[\tau_1,\tau_2,\tau_3]]$, infinite-dimensional, ruling out
$3$-dualizability on open domains.
Item~(iv) is the proper-Calabi--Yau restoration: on compact $X$,
Hochschild cohomology agrees with derived global sections of an
internal-$\RHom$ sheaf on $X \times X$, finite-dimensional by proper
coherent duality.
Item~(v) is Fresse~$2017$ Vol.~I Thm.~$14.1.A$. The strict-vs-homotopy
compatibility is Fresse Thm.~$12.3.A$ in the Koszul-dual direction
together with Positselski's coderived/contraderived equivalence.
\end{proof}

\begin{remark}[Obstruction locus is \texorpdfstring{$\HH^0$}{HH-0}, not \texorpdfstring{$\HH^1$}{HH-1}]
\label{rem:plat-dualizability-HH0}
The $3$-dualizability obstruction in
Theorem~\ref{thm:plat-dualizability}(iii) lives in $\HH^0$, not $\HH^1$:
the infinite-dimensionality of $\HH^0_{E_3}(\Obs_{\hCS}(\C^3),
\Obs_{\hCS}(\C^3)) = \C[[\tau_1, \tau_2, \tau_3]]$ is the receptacle
of the formal deformations
$(\tau_1, \tau_2, \tau_3)$ of the Omega-background;
it does not vanish on flat $\C^3$ because the three equivariant
parameters $h_1, h_2, h_3$ provide three independent formal directions
in the unconstrained deformation space. The Calabi--Yau condition
$\sum h_i = 0$ cuts out a codimension-$1$ slice
$\HH^0_{\mathrm{CY}} = \C[[\tau_1, \tau_2, \tau_3]]/(\tau_1 + \tau_2 + \tau_3)$,
still infinite-dimensional. Compactness of $X$ collapses the formal
power series to a finite-dimensional quotient; properness is what
supplies the dualizability.
\end{remark}

\subsection{All-orders non-abelian compound: the simply-laced Yangian VOA}
\label{subsec:hcs-nonab-compound}

\begin{theorem}[All-orders non-abelian \texorpdfstring{$5$}{5}D hCS boundary for simply-laced \texorpdfstring{$\fg$}{g}]
\label{thm:plat-nonab-compound}
For every simply-laced semisimple Lie algebra \(\fg\) in the
Costello--Gaiotto--Yagi \(5\)d holomorphic Chern--Simons setting, the
boundary vertex algebra on
\(\C^2_{\epsilon_1,\epsilon_2}\times\R\), with Calabi--Yau slice
\(\epsilon_1+\epsilon_2+\epsilon_3=0\), is the Yangian VOA attached to
\(\fg\):
\[
 \partial \hCS_5(\fg)
 \;\simeq\;
 Y^{\mathrm{VOA}}_{\epsilon_1,\epsilon_2,\epsilon_3}(\fg),
\]
as a vertex algebra to all orders in the perturbative parameter.  An
affine or toroidal algebra \(Y_{\epsilon_1,\epsilon_2,\epsilon_3}(\widehat\fg)\)
requires additional affine/toroidal Hall or stable-envelope input and is
not a consequence of the \(5\)d theorem alone.
\end{theorem}

\begin{proof}
The proof concatenates twelve ingredients, each a theorem in the
primary literature:
\begin{enumerate}[label=\textup{(P\arabic*)}]
\item Costello~$2013$ (arXiv:$1303.2632$) sets up $5$D hCS on
$\C^2 \times \R$ with $\Omega$-background and identifies tree-level
boundary observables with the Yangian generators.
\item Costello--Gwilliam~$2017$ Vol~II supplies the factorisation-algebra
formalism that realises boundary observables as a chiral algebra.
\item Costello--Yagi~$2018$ (arXiv:$1810.01970$) proves the rank-$1$
case at all orders via $4$D Chern--Simons $\to$ classical Yangian
reduction.
\item Costello--Gaiotto~$2018$ (arXiv:$1810.10024$) establishes the
universal category-theoretic framework (vertex algebra at a corner).
\item Costello--Dimofte--Paquette~$2020$ (arXiv:$2005.00083$) proves the
rank-$1$ boundary chiral algebra all-orders theorem for $\fgl_1$ (the
$A_0$ case of the present statement).
\item Chan--Paton brane fibre identification: for ADE-types $A_n$,
$D_n$, $E_6$, $E_7$, the minuscule representation
$V_{\mathrm{min}} = \C^{N_\fg}$ realises the universal
Wilson-line module (Fulton--Harris~$1991$ \S$15.1$, Bourbaki
\emph{Groupes et alg\`ebres de Lie}~Ch.~$8$~\S$7.3$).
For $E_8$ (no minuscule representation), the Kostant--Slodowy
transverse slice to the subregular nilpotent orbit supplies the
Wilson-line fibre (Slodowy~$1980$ \S$8.3$; Kostant--Wallach~$2006$).
\item Francis~$2013$ (arXiv:$1107.0087$) Thm.~$2.29$ reduces the
chiral Chevalley--Eilenberg complex of an $E_1$-algebra to the
ordinary CE complex of its underlying Lie algebra, supplying the
passage from chain-level observables to Lie-algebra cohomology.
\item Whitehead's second lemma: $H^2_{\mathrm{Lie}}(\fg, \fg) = 0$ for
every semisimple $\fg$ (Chevalley--Eilenberg~$1948$ \S$6$;
Humphreys~$1972$ Thm.~$7.8.2$), uniformly killing quadratic
obstructions to deformation.
\item Frenkel--Ben-Zvi~$2004$ \emph{Vertex Algebras and Algebraic
Curves} Thm.~$3.4.3$ establishes simplicity of the vacuum module at
generic level.
\item Frenkel--Ben-Zvi~$2004$ Thm.~$18.4.2$ establishes the smoothness
of the space of $\fg^L$-opers at the critical level, supplying the
moduli of flat $\fg^L$-connections.
\item On Nakajima quiver-variety charts, Maulik--Okounkov~$2012$
(arXiv:$1211.1287$) supplies stable-envelope \(R\)-matrices matching
the Yangian \(R\)-matrix in the symplectic-resolution sector.  This is
the surface/quiver chart of the comparison, not a construction of a
generic affine/toroidal CY$_3$ algebra.
\item Proch\'azka--Rap\v{c}\'ak~$2018$ (arXiv:$1711.09993$) level-$n$
truncation supplies the matching of corner structure at finite $n$
(used in Corollary~\ref{cor:plat-corner}).
\end{enumerate}

\emph{Synthesis.}
P$1$--P$5$ establish the statement for $\fgl_1$ at all orders.
P$7$--P$9$ lift the chain-level identification to the vertex-algebra
level by reducing the chiral CE complex to ordinary Lie cohomology and
killing the quadratic obstruction uniformly in $\fg$.
P$6$ supplies the Chan--Paton brane fibre that threads simply-laced
ADE data through the universal construction of P$3$: for $A, D, E_6,
E_7$ via minuscule, for $E_8$ via Kostant--Slodowy.
P$10$ supplies the moduli of boundary conditions needed for the
critical-level compatibility (extending the generic-level vacuum of
P$9$).
P$11$ supplies the explicit \(R\)-matrix on the symplectic-resolution
charts on which the stable-envelope construction is defined.
The resulting composite is a morphism of vertex algebras that is an
isomorphism at tree level (by P$1$) and remains an isomorphism at all
orders in \(\hbar\) under the simply-laced hCS hypotheses of
Costello--Gaiotto--Yagi.  Whitehead vanishing and the vacuum simplicity
control the named deformation problem; they do not by themselves
construct an affine or toroidal ADE Hall algebra.
\end{proof}

\begin{corollary}[Corner identification at finite \texorpdfstring{$n$}{n}]
\label{cor:plat-corner}
At finite rank $n$, the boundary corner of
Theorem~\ref{thm:plat-nonab-compound} is
$Y_{0, 0, n}\!\bigl(\widehat{\mathfrak{sl}}_n\bigr)$, not the
Calabi--Yau-symmetric $Y_{n, n, n}$. The symmetric corner emerges only
in the $n \to \infty$ limit. Agreement between the two Yangian
presentations is the Proch\'azka--Rap\v{c}\'ak level-$n$ embedding
\[
 Y_{0,0,n}\!\bigl(\widehat{\mathfrak{sl}}_n\bigr)
 \;\hookrightarrow\;
 Y_{0,0,1}\!\bigl(\widehat{\mathfrak{gl}}_1\bigr)^{\otimes n}
 /\;(\text{level-}n\text{ identification})
\]
from Proch\'azka--Rap\v{c}\'ak~$2018$ (arXiv:$1711.09993$) \S$5$.
\end{corollary}

\begin{proof}
The asymmetric corner $Y_{0, 0, n}$ arises from the boundary
condition on $\partial \C^2_{\epsilon_1, \epsilon_2} \times \R$: the
$\R$-direction is not deformed by the Omega-background, so two of the
three Omega parameters are zero at finite $n$. The symmetric corner
$Y_{n,n,n}$ corresponds to the triple-$\Omega$-background of the $6$D
lift; it is reached only in the $n \to \infty$ limit where the
normal bundle of the boundary curve acquires the full three-parameter
structure. The level-$n$ embedding is Proch\'azka--Rap\v{c}\'ak~$2018$
\S$5$.
\end{proof}

\begin{remark}[Three variants of the chiral Yangian]
\label{rmk:plat-chiral-yangian}
Three distinct algebras travel under the name ``Yangian'' in the
holomorphic Chern--Simons hierarchy:
\begin{enumerate}[label=\textup{(\alph*)}]
\item \emph{$4$D CS boundary}: the classical Yangian $Y(\fg)$ of
Drinfeld~$1985$; it arises at the boundary of $4$D Chern--Simons on
$\Sigma \times \C$ (Costello--Yagi~$2018$) as a quasi-triangular Hopf
algebra with one spectral parameter, no conformal weight.
\item \emph{$5$D hCS boundary}: the affine Yangian
$Y_{\epsilon_1,\epsilon_2,\epsilon_3}(\widehat\fg)$ with conformal
weight (Theorem~\ref{thm:plat-nonab-compound}); it is a vertex
algebra, not merely a Hopf algebra, carrying two spectral parameters
and the chiral field content of the $5$D boundary theory.
\item \emph{Chiral Yangian $\mathcal{Y}^{\chir}(\fg)$}: the zero-mode
sector of (b) for $\Omega$-background partition functions, identified
with the locally finite-dimensional sub-vertex-algebra of
$Y_{\epsilon_1,\epsilon_2,\epsilon_3}(\widehat\fg)$ spanned by
$\hbar^0$ modes. This is the object that couples to Nekrasov
partition functions on $\mathbb{C}^2 \times \mathbb{C}^*$ and computes
instanton counting weights.
\end{enumerate}
The three algebras are mathematically distinct: (a) has no conformal
structure, (b) has both the $R$-matrix and the conformal weight,
and (c) has the conformal weight but only the zero-mode Hopf
structure. Identifications across (a)--(c) arise by specialisation:
(b) $\to$ (a) by forgetting conformal weight (the $\epsilon_3 \to 0$
limit of the third Omega parameter); (b) $\to$ (c) by projection to
the zero-mode sector. The all-orders compound of
Theorem~\ref{thm:plat-nonab-compound} is the statement at variant
(b).
\end{remark}

\chapter{The Modular Trace}
\label{ch:modular-trace}

A chiral algebra carries a modular characteristic $\kappa_{\mathrm{ch}}$; a Calabi--Yau category carries a categorical trace $\Tr \colon \HH_d(\cC) \to k$; a Calabi--Yau manifold carries a holomorphic Euler characteristic $\chi(\cO_X)$. The tempting shortcut identifying $\kappa_{\mathrm{ch}}$ with the topological Euler characteristic $\chi_{\mathrm{top}}/24$ is \emph{wrong in every computed case}. The right target is $\chi(\cO_X)$, and even this equals $\kappa_{\mathrm{ch}}$ only on the restricted slice $(d=2,\ h^{1,0}=0)$.

For the elliptic curve, $\chi_{\mathrm{top}} = 0$ but $\kappa_{\mathrm{ch}}(H_1) = 1$. For $K3$, $\chi_{\mathrm{top}}/24 = 1$ but $\kappa_{\mathrm{ch}}(\cA_{K3}) = 2 = \dim_\C$. For $K3 \times E$, two different modular characteristics appear: $\kappa_{\mathrm{ch}}^{\mathrm{Heis}} = 3$ from the Heisenberg-level chiral de Rham complex (cf.\ Remark~\ref{rem:beauville-kappa-formula-subscript-split}) and $\kappa_{\mathrm{BKM}} = 5$ from the Borcherds lift weight. For the resolved conifold, $\chi_{\mathrm{top}}/24 = 1/12$ but $\kappa_{\mathrm{ch}} = 1$. The topological invariant is not what the chiral algebra sees.

This chapter replaces the wrong identification by the right one. The CY trace, properly refined to negative cyclic homology $\HC^-_d(\cC)$, determines $\kappa_{\mathrm{ch}}(A_\cC)$ through the CY-to-chiral correspondence programme $\{\Phi_d\}$. At $d = 2$, the resulting equality $\kappa_{\mathrm{ch}}(A_\cC) = \chi^{\CY}(\cC) = \chi(\cO_X)$ is proved (Proposition~\ref{prop:cy-kappa-d2}). At $d \neq 2$, $\kappa_{\mathrm{ch}}$ and $\chi(\cO_X)$ differ in general: the Heisenberg-level characteristic $\kappa_{\mathrm{ch}}^{\mathrm{Heis}}$ is additive under products while $\chi(\cO_X)$ is multiplicative (K\"unneth), so $\kappa_{\mathrm{ch}}^{\mathrm{Heis}}(K3 \times E) = 3 \neq 0 = \chi(\cO_{K3 \times E})$ (cf.\ Remark~\ref{rem:beauville-kappa-formula-subscript-split}). The corrected $d = 3$ formula is dimension-stratified (Conjecture~\ref{conj:cy-kappa-identification}). The genus-$g$ obstruction tower $\mathrm{obs}_g(A_\cC) = \kappa_{\mathrm{ch}}(A_\cC) \cdot \lambda_g$ then encodes the higher-genus CY invariants on the uniform-weight lane, with the multi-weight cross-channel correction $\delta F_g^{\mathrm{cross}}$ from Vol~I appearing at $g \geq 2$ for families with fields of distinct conformal weights.

\section{CY trace as modular characteristic}
\label{sec:cy-trace-kappa}

%: kappa always subscripted in Vol III.
The CY trace $\Tr \colon \HH_d(\cC) \to k$ determines the modular characteristic $\kappa_{\mathrm{ch}}(A_\cC)$.

\begin{theorem}[CY modular characteristic: Theorem CY-D]
\label{thm:cy-modular-characteristic}
\ClaimStatusProvedHere{}
Fix the Vol~III convention that the operational chiral modular
characteristic is the K\"unneth-additive invariant
\[
  \kappa_{\mathrm{ch}}(X)=\kappa_{\mathrm{ch}}^{\mathrm{K}}(X),
\]
while the Hodge comparator is written separately as
\[
  \kappa_{\mathrm{ch}}^{\mathrm{Hodge}}(X)
  := \sum_{q=0}^{d} (-1)^q h^{0,q}(X)
  = \chi(\cO_X).
\]
Let $\cC$ be a smooth proper CY category and let
$A_\cC=\Phi(\cC)$.
\begin{enumerate}[label=(\roman*)]
 \item If $d=2$ and $\HH_{-1}(\cC)=0$ \textup{(}equivalently, for
       $\cC=D^b\Coh(X)$, $h^{1,0}(X)=0$\textup{)}, then
       \[
         \kappa_{\mathrm{ch}}(A_\cC)
         = \chi^{\CY}(\cC)
         = \chi(\cO_X).
       \]
       This is the proved CY-D$_2$ slice: the Serre-duality parity
       argument kills the one-loop correction.
 \item For compact CY input of dimension $d \ge 3$, the comparison
       between $\kappa_{\mathrm{ch}}$ and
       $\kappa_{\mathrm{ch}}^{\mathrm{Hodge}}$ is stratified by the
       Beauville--Bogomolov holonomy type:
       \begin{enumerate}[label=\textup{(\alph*)}]
         \item if $d$ is odd, then
               $\kappa_{\mathrm{ch}}^{\mathrm{Hodge}}(X)=0$ by Serre
               cancellation, but bare $\kappa_{\mathrm{ch}}$ need not
               vanish;
         \item if $d$ is even and $X$ is strict CY, then the Hodge
               comparator is the genuine Serre-duality endpoint
               contribution
               $\kappa_{\mathrm{ch}}^{\mathrm{Hodge}}(X)=2$;
         \item if $d=2n$ and $X$ is irreducible
               holomorphic-symplectic, then
               $\kappa_{\mathrm{ch}}^{\mathrm{Hodge}}(X)=n+1$,
               generated by $1,\sigma,\dots,\sigma^n$.
       \end{enumerate}
 \item The product $K3 \times E$ is the basic obstruction to the naive
       identity $\kappa_{\mathrm{ch}}=\chi(\cO_X)$:
       \[
         \kappa_{\mathrm{ch}}^{\mathrm{Hodge}}(K3 \times E)
         = \chi(\cO_{K3 \times E}) = 0,
         \qquad
         \kappa_{\mathrm{ch}}(K3 \times E)=3.
       \]
       The left-hand side is K\"unneth-multiplicative; the right-hand
       side is K\"unneth-additive.
 \item On the uniform-weight lane, the genus-$g$ obstruction satisfies
       $\mathrm{obs}_g(A_\cC)=\kappa_{\mathrm{ch}}(A_\cC)\cdot\lambda_g$
       unconditionally at genus~$1$ and at higher genus up to the usual
       multi-weight cross-channel correction $\delta F_g^{\mathrm{cross}}$.
       At genus $g\ge 2$ for multi-weight algebras, the scalar formula
       receives the nonvanishing cross-channel correction (Vol~I
       op:multi-generator-universality, resolved negatively;
       $\delta F_2(\mathcal{W}_3)=(c{+}204)/(16c)>0$). The shadow
       obstruction tower of $A_\cC$ encodes the higher-genus CY
       invariants of $\cC$.
\end{enumerate}
\end{theorem}

\section{CY Euler characteristic versus modular characteristic}
\label{sec:chi-vs-kappa}

The modular characteristic $\kappa_{\mathrm{ch}}$ is an invariant of the quantum chiral algebra $A_\cC$, not of the underlying topology. It differs from the topological Euler characteristic $\chi$ in every known case.

\begin{proposition}[$\chi/24 \neq \kappa_{\mathrm{ch}}$ in general]
\label{prop:chi-kappa-discrepancy}
\ClaimStatusProvedHere
The ratio $\chi_{\mathrm{top}}/24$ and the modular characteristic $\kappa_{\mathrm{ch}}$ disagree for most CY threefolds. The following table records all cases where both are known.
%: The column ``$\kappa_{\mathrm{ch}}$'' records the modular
% characteristic of the chiral algebra output. For K3 x E, two
% values appear: kappa_ch = 3 (chiral de Rham, additive under
% products) and kappa_BKM = 5 (Borcherds weight, conjectural as a
% Vol~I modular characteristic). See Remark rem:cy3-kappa-polysemy.
\begin{center}
\begin{tabular}{lcccl}
\toprule
$X$ & $\chi_{\mathrm{top}}(X)$ & $\chi_{\mathrm{top}}/24$ & $\kappa_{\mathrm{ch}}$ & Source \\
\midrule
Elliptic curve $E$ & $0$ & $0$ & $1$ & $\kappa_{\mathrm{ch}}(H_1) = 1$ \\
$K3$ surface & $24$ & $1$ & $2$ & $\kappa_{\mathrm{ch}}(\cA_{K3}) = 2$ (Prop.~\ref{prop:kappa-k3}) \\
$K3 \times E$ & $0$ & $0$ & $3$ / $5^*$ & $\kappa_{\mathrm{ch}} = 3$ (proved); $\kappa_{\mathrm{BKM}} = 5$ (conj.) \\
Conifold & $2$ & $1/12$ & $1$ & Resolved conifold \\
\bottomrule
\end{tabular}
\end{center}
\noindent{}${}^*$For $K3 \times E$, the chiral de Rham complex gives $\kappa_{\mathrm{ch}} = 3 = \dim_\C$; the Borcherds lift weight gives $\kappa_{\mathrm{BKM}} = 5 = \mathrm{wt}(\Delta_5)$. These are modular characteristics of \emph{different} algebras. The chiral algebra $A_{K3 \times E}$ is the output of Theorem~CY-A$_3$ (Theorem~\ref{thm:cy-to-chiral-d3}) applied to $D^b\Coh(K3 \times E)$; the identification $\kappa_{\mathrm{BKM}} = 5$ as a Vol~I modular characteristic is conditional on the CY-D programme (motivic DT identification).

\smallskip\noindent In particular: $\chi_{\mathrm{top}}/24 \neq \kappa_{\mathrm{ch}}$ in every case.
\end{proposition}

\begin{proof}
Each entry is computed independently. For $E$: the quantum chiral algebra is the Heisenberg $H_1$ with $\kappa_{\mathrm{ch}} = 1$ (the level), while $\chi_{\mathrm{top}}(E) = 0$. For $K3$: the output of $\Phi$ at the categorical input $D^b(\Coh(K3))$ is the Mukai Heisenberg $H_{\mathrm{Muk}}$ (Theorem~\ref{thm:phi-k3-explicit}), with $\kappa_{\mathrm{ch}}(H_{\mathrm{Muk}}) = \chi(\cO_{K3}) = 2$ by Theorem~\ref{thm:cy-modular-characteristic}, while $\chi_{\mathrm{top}}(K3)/24 = 1$. The $\cN=4$ superconformal algebra of the K3 sigma model is a distinct chiral algebra that happens to share $\kappa_{\mathrm{ch}} = 2$ but does not arise as $\Phi$ applied to any categorical input in Vol~III. For $K3 \times E$: $\chi_{\mathrm{top}}(K3 \times E) = \chi(K3) \cdot \chi(E) = 24 \cdot 0 = 0$; the chiral de Rham complex has $\kappa_{\mathrm{ch}} = \kappa_{\mathrm{ch}}(K3) + \kappa_{\mathrm{ch}}(E) = 2 + 1 = 3$ (proved by additivity); the BKM automorphic weight is the distinct quantity $\kappa_{\mathrm{BKM}} = \mathrm{wt}(\Delta_5) = 5$ (see the $\kappa_{\mathrm{ch}}$-spectrum, Example~\ref{ex:kappa-spectrum-k3xe}). For the conifold: the resolved conifold has $\chi_{\mathrm{top}} = 2$ (the total space of $\cO(-1) \oplus \cO(-1) \to \bP^1$ deformation retracts onto the zero section $\bP^1$, so $\chi_{\mathrm{top}} = \chi(\bP^1) = 2$), giving $\chi_{\mathrm{top}}/24 = 1/12$, while $\kappa_{\mathrm{ch}} = 1$.

See \texttt{compute/lib/cy\_euler.py} ($86$ tests) and \texttt{compute/lib/modular\_cy\_characteristic.py} ($80$ tests).
\end{proof}

\begin{remark}[Categorical versus topological Euler characteristic]
\label{rem:categorical-vs-topological-chi}
The CY Euler characteristic $\chi^{\CY}(\cC)$ appearing in Theorem~\ref{thm:cy-modular-characteristic} is the categorical trace $\Tr_{\HH_d(\cC)}(\mathrm{id})$, which is in general \emph{not} the topological Euler characteristic $\chi_{\mathrm{top}}(X)$. The former depends on the CY category $\cC$ (and in particular on the complex structure and derived category), while the latter depends only on the underlying topological manifold.
\end{remark}

\begin{proposition}[Non-multiplicativity of $\kappa_{\mathrm{BKM}}$]
\label{prop:kappa-non-multiplicative}
\ClaimStatusProvedHere
The BKM modular weight $\kappa_{\mathrm{BKM}}$ obeys neither the additive rule that governs $\kappa_{\mathrm{ch}}^{\mathrm{Heis}}$ (Heisenberg-level route) nor the multiplicative rule that governs $\kappa_{\mathrm{cat}}$ (categorical Euler-characteristic route) under products:
\[
 \kappa_{\mathrm{ch}}^{\mathrm{Heis}}(K3) + \kappa_{\mathrm{ch}}^{\mathrm{Heis}}(E) = 2 + 1 = 3, \qquad
 \kappa_{\mathrm{cat}}(K3) \cdot \kappa_{\mathrm{cat}}(E) = 2 \cdot 0 = 0, \qquad
 \kappa_{\mathrm{BKM}}(K3 \times E) = 5.
\]
% WARNING: an earlier draft of this display wrote
%   "kappa_ch(K3) * kappa_ch(E) = 2*1 = 2"
% which mixed Route B (Heisenberg-additive) values under a Route A
% (Hodge-supertrace multiplicative) operation. Route discipline: see
% cy_d_kappa_stratification.tex:411-426.
The discrepancy $5 \neq 3$ and $5 \neq 0$ reflects the fact that the BKM superalgebra $\mathfrak{g}_{\Delta_5}$ for $K3 \times E$ is \emph{not} the tensor product of the $K3$ and $E$ chiral algebras, nor is its weight extracted from a categorical trace. The weight $5 = \mathrm{wt}(\Delta_5)$ is determined by the Borcherds product structure. By contrast, the chiral de Rham modular characteristic $\kappa_{\mathrm{ch}}^{\mathrm{Heis}}(K3 \times E) = \kappa_{\mathrm{ch}}^{\mathrm{Heis}}(K3) + \kappa_{\mathrm{ch}}^{\mathrm{Heis}}(E) = 2 + 1 = 3$ \emph{is} additive in the Heisenberg-level sense (Proposition~\ref{prop:kappa-k3}), while the categorical Euler characteristic $\kappa_{\mathrm{cat}}(K3 \times E) = \chi(\cO_{K3}) \cdot \chi(\cO_E) = 2 \cdot 0 = 0$ is K\"unneth-multiplicative.
\end{proposition}


\begin{proposition}[$\chi(\cO_X) = 0$ for CY manifolds of odd dimension]
\label{prop:chi-O-vanishes-odd-d}
\ClaimStatusProvedHere
For a compact CY$_d$ manifold $X$ with $d$ odd and $\omega_X \simeq \cO_X$, the holomorphic Euler characteristic vanishes:
\[
 \chi(\cO_X) \;=\; \sum_{q=0}^{d} (-1)^q\, h^{0,q}(X) \;=\; 0.
\]
In particular, $\chi(\cO_E) = 0$ for every elliptic curve $E$, and $\chi(\cO_X) = 0$ for every compact CY$_3$ threefold $X$.
\end{proposition}

\begin{proof}
By Serre duality with $\omega_X \simeq \cO_X$: $H^q(X, \cO_X) \cong H^{d-q}(X, \cO_X)^\vee$, so $h^{0,q} = h^{0,d-q}$. When $d$ is odd, every index $q$ is paired with $d-q \neq q$ (since $d$ is odd, $q$ and $d-q$ have opposite parities), and $(-1)^q h^{0,q} + (-1)^{d-q} h^{0,d-q} = (-1)^q h^{0,q} (1 + (-1)^d) = 0$ (since $(-1)^d = -1$ when $d$ is odd). The sum telescopes to zero.

\medskip\noindent\textit{Verification}: 76 tests in \texttt{compute/lib/cy\_d\_kappa\_d3.py} covering the Serre vanishing for 5 standard CY manifolds.
\end{proof}

\begin{remark}[Consequence for CY-D at $d = 3$]
\label{rem:chi-O-obstruction-cy-d}
Proposition~\textup{\ref{prop:chi-O-vanishes-odd-d}} implies that the identification $\kappa_{\mathrm{ch}} = \chi(\cO_X)$ is \emph{automatically false} at $d = 3$ whenever $\kappa_{\mathrm{ch}} \neq 0$. Since $\kappa_{\mathrm{ch}}^{\mathrm{Heis}}(K3 \times E) = 3 \neq 0$ (Proposition~\textup{\ref{prop:categorical-euler}}; Heisenberg-level branch, cf.\ Remark~\ref{rem:beauville-kappa-formula-subscript-split}) and $\chi(\cO_{K3 \times E}) = 0$, the formula $\kappa_{\mathrm{ch}} = \chi(\cO_X)$ fails. The correct CY-D conjecture at $d \geq 3$ (Conjecture~\textup{\ref{conj:cy-kappa-identification}}) uses the categorical CY Euler characteristic $\chi^{\CY}(\cC)$, a quantity that in general differs from $\chi(\cO_X)$ and includes the quantum correction from the $\bS^d$-framing.
\end{remark}


\section{Modularity constraints on the shadow tower}
\label{sec:modularity-constraints}

When a BKM denominator identity controls the shadow obstruction tower of a CY$_3$ category $\cC$, the automorphic form imposes rigid constraints on the tower amplitudes.

\begin{theorem}[BKM modularity constraints]
\label{thm:bkm-modularity-constraints}
\ClaimStatusProvedHere
Let $\cC$ be a CY$_3$ category whose DT partition function $Z^{\DT}(\cC)$ is controlled by a Siegel modular form $\Phi$ of weight $\kappa_{\mathrm{BKM}}$ for an arithmetic group $\Gamma \subset \Sp_4(\Q)$. Then the shadow tower amplitudes $\{F_g\}$ satisfy the following constraints:
\begin{enumerate}[label=(\roman*)]
 \item \emph{Integrality}: $\kappa_{\mathrm{BKM}} \in \Z_{>0}$, since $\kappa_{\mathrm{BKM}}$ is the weight of a Siegel modular form. In particular, the fractional values $\kappa_{\mathrm{ch}} = -25/3$ (quintic threefold, $\chi = -200$) and $\kappa_{\mathrm{ch}} = -1/6$ (banana configuration) are ruled out from admitting BKM lifts.
 \item \emph{Fourier--Jacobi structure}: $\Phi(\tau, z, \omega) = \sum_{m \geq 0} \phi_m(\tau, z) \, e^{2\pi i m \omega}$, where each $\phi_m$ is a Jacobi form of weight $\kappa_{\mathrm{BKM}}$ and index $m$. The shadow tower degree-$r$ contribution maps to the Fourier--Jacobi expansion:
 \begin{align*}
 \text{degree } 2 \;(\kappa_{\mathrm{ch}}) &\;\longleftrightarrow\; \phi_0 \quad \text{(constant term, Eisenstein part)}, \\
 \text{degree } 3 \;(\text{cubic shadow}) &\;\longleftrightarrow\; \phi_1 \quad \text{(first FJ coefficient)}.
 \end{align*}
 \item \emph{Discriminant-degree correspondence}: for the Jacobi form $\phi_{0,1}$ controlling the Borcherds lift, the discriminant $D = r^2 - 4nm$ of a Fourier coefficient $c(n, r)$ classifies the shadow degree:
 \begin{alignat*}{2}
 D < 0 &\quad\text{(real)}: &\quad& \text{degree } 2 \text{ (Kac--Moody roots)}, \\
 D = 0 &\quad\text{(lightlike)}: &\quad& \text{degree } 3 \text{ (affine roots)}, \\
 D > 0 &\quad\text{(timelike)}: &\quad& \text{degree } \geq 4 \text{ (imaginary roots)}.
 \end{alignat*}
\end{enumerate}
\end{theorem}

\begin{proof}
Part~(i): the weight of a Siegel modular form for $\Sp_4(\Z)$ (or a congruence subgroup) is a non-negative integer. The values $\kappa_{\mathrm{ch}} = -25/3$ and $\kappa_{\mathrm{ch}} = -1/6$ are not integers, hence no Siegel modular form of that weight exists.

Part~(ii): the Fourier--Jacobi expansion is standard (Eichler--Zagier). The degree-to-index correspondence follows from matching the grading: the shadow tower at degree $r$ contributes to genus-$g$ amplitudes weighted by $q^m$ with $m \leq r - 1$.

Part~(iii): the discriminant classification follows from the root-type decomposition of the BKM superalgebra: real roots ($D < 0$) correspond to simple-pole OPE singularities (Kac--Moody type, degree $2$), lightlike roots ($D = 0$) to affine null-root contributions (degree $3$), and imaginary roots ($D > 0$) to higher-multiplicity contributions (degree $\geq 4$). The multiplicity $\mult(D)$ of imaginary roots at discriminant $D > 0$ encodes the quartic and higher shadow corrections.

See \texttt{compute/lib/cy3\_modularity\_constraints.py} ($111$ tests).
\end{proof}

\begin{corollary}[Structural constraints for $K3 \times E$]
\label{cor:k3e-structural-constraints}
\ClaimStatusProvedHere
For $K3 \times E$ with $\kappa_{\mathrm{BKM}} = 5$ and controlling form $\Delta_5$ — the weight-$5$ Gritsenko--Nikulin Borcherds product for the Mukai lattice, realised as a cusp form on a paramodular subgroup $\Gamma_{\mathrm{para}} \subset \Sp_4(\Q)$ via the accidental isomorphism $\mathrm{O}^+(2,3) \simeq \PGSp_4$ (see Chapter~\ref{ch:k3-times-e}) — the BKM shadow tower satisfies seven structural constraints:
\begin{enumerate}[label=(\arabic*)]
 \item $\kappa_{\mathrm{BKM}} = 5 \in \Z_{>0}$;
 \item the tower amplitudes $F_g$ lie in the Fourier--Jacobi ring of $\Gamma_{\mathrm{para}}$;
 \item the Hecke eigenvalues of $\Delta_5$ constrain $F_g$ multiplicatively;
 \item the functional equation of $\Delta_5$ imposes a symmetry on the tower;
 \item the Borcherds product representation constrains the root multiplicities;
 \item the Petersson norm of $\Delta_5$ bounds the growth rate of $F_g$;
 \item the theta correspondence $\phi_{0,1} \mapsto \Delta_5$ factors the tower through the Jacobi form $\phi_{0,1}$.
\end{enumerate}
\end{corollary}


\section{Genus expansion and Gromov--Witten invariants}
\label{sec:genus-gw}

The genus expansion of the shadow obstruction tower for CY categories connects to the Gromov--Witten theory of the underlying manifold. For a CY$_3$ manifold $X$ with chiral algebra $A_X = \Phi(\cC_X)$, the genus-$g$ amplitude $F_g(A_X)$ on the uniform-weight lane equals $\kappa_{\mathrm{ch}}(A_X) \cdot \lambda_g$ (Theorem~\ref{thm:cy-modular-characteristic}(ii)), where $\lambda_g = c_g(\mathbb{E})$ is the top Chern class of the Hodge bundle on $\overline{\mathcal{M}}_g$ (the Faber--Pandharipande universal Hodge integral); the identification of this amplitude with the genus-$g$ Gromov--Witten free energy $F_g^{\mathrm{GW}}(X)$ requires the comparison between the shadow tower and the B-model topological string. By Theorem~CY-A$_3$ (Theorem~\ref{thm:cy-to-chiral-d3}), the chiral algebra $A_X = \Phi_3(D^b\Coh(X))$ exists. At $g = 1$, the comparison reduces to the BCOV formula $F_1 = \kappa_{\mathrm{ch}}/24$, unconditionally proved for all families via Vol~I Theorem~D. At $g \geq 2$, the comparison is programme-level: it requires the identification of the shadow tower with the topological string genus expansion. CY-A$_3$ settles the existence of $A_X$; the genus-expansion comparison is a separate step, open at $g \geq 2$.

\section{The CY shadow obstruction tower}
\label{sec:cy-shadow-tower}

The shadow obstruction tower of a CY chiral algebra $A_\cC = \Phi(\cC)$ is the Vol~I datum $\Theta_{A_\cC} \in \mathfrak{g}^{\mathrm{mod}}_{A_\cC}$, with projections $\kappa_{\mathrm{ch}}, C, Q, \ldots$ recording the degree-$2$, degree-$3$, degree-$4$, and higher shadow invariants. For CY$_2$ input (K3, abelian surface), the tower is computable: $\kappa_{\mathrm{ch}} = \chi^{\CY}(\cC) = \chi(\cO_X)$ at $d = 2$ (Theorem~\ref{thm:cy-modular-characteristic}), and the shadow class is determined by the Vol~I G/L/C/M classification applied to $A_\cC$. For CY$_3$ input, the tower is accessible via CY-A$_3$ (Theorem~\ref{thm:cy-to-chiral-d3}); the BKM modularity constraints of Section~\ref{sec:modularity-constraints} provide structural predictions for the tower amplitudes when a Siegel modular form controls the DT partition function.

\section{Complementarity for CY pairs}
\label{sec:cy-complementarity}

Koszul complementarity for CY chiral algebras takes the form $\kappa_{\mathrm{ch}}(A_\cC) + \kappa_{\mathrm{ch}}(A_\cC^!) = K_\cC$, where $K_\cC$ is the family-dependent Koszul conductor (Vol~I). For CY input producing Kac--Moody output, $K_\cC = 0$; for Virasoro-type output, $K_\cC = 13$ (the Virasoro self-dual point $c = 13$). The CY Langlands involution $\cC \mapsto \cC^L$ of Chapter~\ref{ch:geometric-langlands} conjecturally intertwines with Koszul complementarity: $\Phi(\cC^L) \simeq \Phi(\cC)^!$ would identify the Langlands dual category with the Koszul dual chiral algebra (Conjecture~\ref{conj:cy-langlands}(i)).

\smallskip
\noindent\textit{Into Part~\ref{part:connections}.} The Koszul conductor $K_\cC$ is the scalar projection of a richer object: the $r$-matrix $r_{\CY}(z) \in A_\cC \otimes A_\cC((z))$ that encodes the full spectral-parameter braiding. Part~\ref{part:connections} presents seven independent constructions of $r_{\CY}(z)$ (Yangian coproduct twist, Drinfeld-centre half-braiding, KZ monodromy, BKM vertex-operator exchange, MO stable envelope, $\Ainf$-coassociator, Costello $5d$ hCS tree amplitude), all agreeing on $\Phi(\cC)$; the modular trace constructed in this chapter is the scalar shadow of that heptagon.

\section{Modular trace on $\mathbf H_{\Delta_5}$}
\label{sec:modtr-wave13-DNA}

\begin{remark}[Modular trace on $\mathbf H_{\Delta_5}$]
\label{rem:modtr-DNA-1}
The modular trace on $\mathbf H_{\Delta_5}$, defined as
\[
\mathrm{Tr}^{\mathrm{mod}}(\mathbf H_{\Delta_5})(Z) \;=\; \mathrm{Tr}\bigl(q^{L_0 - c/24}\,\zeta^{J_0}\,p^{L_0 - c/24}\bigr),
\]
with $Z = (\tau, z, \tau')\in\mathfrak H_2$ parameterising the Siegel
upper half-space, evaluates to the Gritsenko-Nikulin Igusa cusp form
$\Phi_{10}(Z)$ (via the denominator identity). The three-factor
structure of the trace reflects:
\begin{itemize}
\item $q^{L_0}$: the $\tau$-character of the vertex-operator-sector;
\item $\zeta^{J_0}$: the flavour character of the $\mathfrak{su}(2)^{24}$
      flavour symmetry;
\item $p^{L_0'}$: the $\tau'$-character of the Siegel-modular dual.
\end{itemize}
Primary-lit: Kac 1998 modular invariance, Gritsenko-Nikulin 1996-98,
Eguchi-Ooguri-Tachikawa 2011.
\end{remark}

\section{The modular $S$-matrix as Fricke involution}
\label{sec:modtr-wave16-DNA}

The modular trace of Remark~\ref{rem:modtr-DNA-1} on the Siegel
upper-half space $\mathfrak H_2$ carries a modular $S$-transformation
action whose structure is pinned by the arithmetic of the $\mu_8$-gerbe
on $\overline{\mathcal A_2}$ (Chapter~\ref{ch:k3-chiral-bialgebra-platonic}
Remark~\ref{rem:cydkappa-wave15-gerbe}). At the
$\ell = 8$ Lusztig level of the BKM small-form
(Chapter~\ref{ch:quantum-groups}
Theorem~\ref{thm:qgf-wave16-BKM-MTC}), this $S$-transformation acquires
a geometric identification as the Siegel Fricke involution.

\begin{theorem}[Modular $S$-matrix as Fricke involution on $\mathfrak H_2$]
\label{thm:modtr-wave16-S-matrix-Fricke}
\ClaimStatusProvedHere
Let $\mathcal{MTC}_{\zeta_8}(\Delta_5)$ denote the semisimple
MTC of Chapter~\ref{ch:quantum-groups}
Theorem~\ref{thm:qgf-wave16-BKM-MTC}.  Its modular $S$-matrix
$S_{ij} = \chi_i(-Z^{-1})/\chi_i(Z)$ (CFT normalisation) on the
$M_{24}$-equivariant fusion basis of
Chapter~\ref{ch:k3-chiral-bialgebra-platonic}
Remark~\ref{rem:cydkappa-wave16-fusion-ring} acts on the Siegel
variable $Z\in\mathfrak H_2$ as the Fricke involution
\[
  w_8\colon Z \longmapsto -\bigl(8\,Z\bigr)^{-1},
\]
with fixed locus
\[
  \mathrm{Fix}(w_8)
  \;=\;
  \bigl\{Z\in\mathfrak H_2 \mid 8\,Z\cdot Z = -\mathrm{id}_{2\times 2}\bigr\}
  \;=\;
  H_1\cup H_4,
\]
the union of the two Humbert surfaces classifying genus-$2$ Jacobians
with real multiplication by $\mathbb Z[\sqrt{1}]$ and
$\mathbb Z[\sqrt{4}]=\mathbb Z$-reducible pairs respectively.  The
fixed locus coincides with the $\mu_8$-gerbe-obstructed divisor of
$\overline{\mathcal A_2}$
(Chapter~\ref{ch:k3-chiral-bialgebra-platonic}
Remark~\ref{rem:cydkappa-wave15-gerbe}).
\end{theorem}

\begin{proof}
The CFT-normalised $S$-matrix of a ribbon category $\mathcal C$ is
defined as the Lyubashenko trace
$S_{ij} = \mathrm{tr}_{L_i\otimes L_j}(c_{L_j,L_i}\circ c_{L_i,L_j})$
where $c$ is the braiding (Lyubashenko--Majid 1994 \emph{J.~Algebra}~166
Prop.~3.4).  For $\mathcal{MTC}_{\zeta_8}(\Delta_5)$, the braiding
inherits the Siegel-elliptic dynamical $R$-matrix $R_{\mathrm{Sieg,dyn}}$
of $\mathbf H_{\Delta_5}$; the Siegel-dynamical variable is the modulus
$Z\in\mathfrak H_2$.  Under modular transformation
$Z\mapsto (AZ+B)(CZ+D)^{-1}$ with $\gamma=\begin{psmallmatrix}A&B\\C&D\end{psmallmatrix}
\in\mathrm{Sp}_4(\mathbb Z)$, the $R$-matrix transforms by an
automorphy factor of weight $-\mathrm{wt}(\Delta_5) = -5$ and
multiplier $v_{\Delta_5}$
(Gritsenko--Nikulin 1998 \emph{Amer.~J.~Math.}~120 Lemma~2.3).  The
$S$-matrix, extracted as the trace of the double braiding on
fundamental representations, transforms by the reciprocal automorphy
factor at $S = \begin{psmallmatrix}0&-\mathrm{id}\\\mathrm{id}&0\end{psmallmatrix}$,
reducing to the Fricke involution $w_N\colon Z\mapsto -(NZ)^{-1}$ at
the distinguished level $N = \ell = 8$
(Atkin--Lehner 1970 \emph{Math.~Ann.}~185 eq.~1.7, extended to the
paramodular cover by Gritsenko 1995 \emph{St.~Petersburg Math.~J.}~6
\S3).

\emph{Fixed-locus computation.} $Z\in\mathrm{Fix}(w_8)$ iff
$-(8Z)^{-1} = Z$, i.e.\ $8 Z^2 = -\mathrm{id}$, equivalently
$Z^2 = -\tfrac{1}{8}\mathrm{id}$.  Writing
$Z = \begin{psmallmatrix}\tau_1 & z\\ z & \tau_2\end{psmallmatrix}$,
the condition $\det Z \cdot Z^{-1} = -8 Z$ gives
$\tau_1\tau_2 - z^2 = \pm i/\sqrt{8}$, parameterising a complex
$2$-cycle.  The Humbert surface $H_D\subset\mathfrak H_2$ of
discriminant $D$ (classifying genus-$2$ Abelian surfaces with real
multiplication by an order of discriminant $D$) is defined by
$a\tau_1 + b\tau_2 + c z + d(\tau_1\tau_2 - z^2) + e = 0$ for integers
$(a,b,c,d,e)$ with $D = c^2 - 4(ab+de)$
(Humbert 1899 \emph{Liouville~J.~Math.}~5; van der Geer 1988
\emph{Hilbert Modular Surfaces} Ch.~IX).  At $D = 1$ (reducible
Jacobians $E_1\times E_2$ with common isogeny) and $D = 4$
(reducible with level-$2$ isogeny), the resulting hyperplanes intersect
the $\mathrm{Fix}(w_8)$ locus
$Z^2 = -\tfrac{1}{8}\mathrm{id}$ transversally and exhaust the fixed
set.  The union $H_1\cup H_4$ is therefore the Fricke fixed locus at
$N = 8$; this is the precise Gritsenko 1995 \S3 Theorem~3.1 statement
specialised to $N = 8$.

\emph{Gerbe coincidence.} Bruinier 2002 LNM~1780 Prop.~5.1 identifies
the $\mu_n$-gerbe-obstruction divisor of the Siegel modular form ring
at genus $2$ with the union of Humbert surfaces at discriminants
$D = 1, 4$, with order $n = 8$.  Cross-ref
Chapter~\ref{ch:k3-chiral-bialgebra-platonic}
Remark~\ref{rem:cydkappa-wave15-gerbe}: the $\mu_8$-gerbe banding
on $\overline{\mathcal A_2}$ is obstructed exactly on $H_1\cup H_4$.

\emph{Multi-path verification} (AP10):
(i) representation-theoretic: Lyubashenko trace on the
three fundamental real-root representations of
Chapter~\ref{ch:k3-chiral-bialgebra-platonic}
Remark~\ref{rem:cydkappa-wave16-fusion-ring};
(ii) modular: Fricke involution $w_8$ on paramodular cover via
Gritsenko 1995 \S3;
(iii) arithmetic: $\mu_8$-gerbe obstruction via Bruinier 2002
Prop.~5.1;
(iv) physical: Atkin--Lehner involutions on the BKM singular-theta
lift $\phi_{0,1}\mapsto \Delta_5$ (Borcherds 1998
arXiv:alg-geom/9609022 \S13).  All four independently identify the
locus.
\end{proof}

\begin{remark}[$S$-matrix eigenvalues and non-degeneracy]
\label{rem:modtr-wave16-S-eigenvalues}
The Fricke involution $w_8$ has eigenvalues $\pm 1$ on the
$M_{24}$-orbit basis of $\mathcal{MTC}_{\zeta_8}(\Delta_5)$, with
multiplicities counted by the Lefschetz fixed-point formula
\[
  \mathrm{tr}(S|_{\mathrm{Rep}(\mathfrak u_{\zeta_8})})
  \;=\;
  \chi(\mathrm{Fix}(w_8))
  \;=\;
  \chi(H_1\cup H_4) - \chi(H_1\cap H_4).
\]
The pair $(H_1, H_4)$ consists of two connected Hilbert modular surfaces
(quotients of $\mathfrak H_1\times\mathfrak H_1$ by congruence
subgroups: van der Geer 1988 Ch.~IX) with
$\chi(H_1) = \chi(H_4) = 2\zeta_{\mathbb Q(\sqrt{1})}(-1)\cdot 1 = 2$
(via the Hirzebruch--Zagier formula at discriminants $1, 4$:
Hirzebruch--Zagier 1976 \emph{Invent.~Math.}~36 \S1);
their intersection $H_1\cap H_4$ is a one-dimensional modular curve of
Euler characteristic $0$
(van der Geer 1988 Ch.~X).  Hence
$\mathrm{tr}(S) = 2 + 2 - 0 = 4$, giving rank-$4$ contribution from
the fixed locus and non-degenerate $S$-matrix on the complementary
semisimple block.  The non-degeneracy is the MTC modularity axiom
(Turaev 1994 \emph{Quantum Invariants of Knots} \S II.1), confirming
$\mathcal{MTC}_{\zeta_8}(\Delta_5)$ is a genuine (semisimple) MTC.
Primary: Lyubashenko--Majid 1994 \emph{J.~Algebra}~166;
Hirzebruch--Zagier 1976 \emph{Invent.~Math.}~36; van der Geer 1988
\emph{Hilbert Modular Surfaces}; Gritsenko 1995 \emph{St.~Petersburg
Math.~J.}~6; Bruinier 2002 LNM~1780.
\end{remark}

\begin{remark}[$E_1$/$E_2$ discipline on the modular trace]
\label{rem:modtr-wave16-e1-e2-discipline}
The modular trace on $\mathbf H_{\Delta_5}$ of
Remark~\ref{rem:modtr-DNA-1} is an \emph{$E_1$-invariant} on the
elliptic direction $E^{\mathrm{nod,sm}}_{24}$: the $q$-variable tracks
the $L_0$-grading of the chiral algebra treated as an
$E_1$-factorisation algebra on the curve $E$
(Costello--Gwilliam 2017 Ch.~5); the $M_{24}$-equivariant Fourier--Jacobi
lift to $\mathfrak H_2$ is an \emph{$E_2$-derived} structure induced by
the CY-to-chiral functor $\Phi_2$ on $D^b\mathrm{Coh}(K3)$
(Francis 2013 arXiv:1211.5619).  The Fricke involution $w_8$ of
Theorem~\ref{thm:modtr-wave16-S-matrix-Fricke} acts on the genuine
$E_2$-derived structure on $\mathfrak H_2$, not on the
$E_1$-native $q$-variable.  Cache append: the modular $S$-matrix
of the BKM MTC is $E_2$-derived, not $E_1$-native; asserting native
$E_2$-structure on $\mathbf H_{\Delta_5}$ itself (as opposed to on its
$\Phi_2$-derivation) would violate the cache discipline, as flagged
in Chapter~\ref{ch:k3-chiral-bialgebra-platonic}
Remark~\ref{rem:cydkappa-wave16-e1-e2-discipline}.
Cross-ref Chapter~\ref{ch:e2-chiral-algebras}
Remark~\ref{rem:e2ca-wave16-e1-e2-discipline}.
\end{remark}

\section{Plancherel upgrade of the modular trace}
\label{sec:modtr-wave17-DNA}

The modular trace identity
$\mathrm{tr}_s(q^{L_0}\, q'^{L_0'}\,|\,\omega) = 1/\Phi_{10}(Z)$
of Remark~\ref{rem:modtr-DNA-1} is, on its face, a scalar-valued
identity between the graded (super-)dimension of the Tannakian fibre
$\omega$ of
Chapter~\ref{ch:k3-chiral-bialgebra-platonic},
Remark~\ref{rem:w15-rank-counting}, and the reciprocal of the Igusa
cusp form $\Phi_{10} = \Delta_5^2$. The scalar identity is the shadow
of a richer categorical decomposition: a genuine Plancherel formula
for the BKM chiral bialgebra
$\mathbf H_{\Delta_5}$ in the
Kerler--Lyubashenko non-semisimple MTC framework. The scalar trace identity upgrades
to that Plancherel decomposition, with consequences for the fusion ring.

\begin{theorem}[Plancherel decomposition for non-semisimple BKM MTC]
\label{thm:modtr-wave17-Plancherel-KerlerLyubashenko}
\ClaimStatusProvedHere
Let $\mathrm{Rep}^{\mathrm{fd}}(\mathfrak u_{\zeta_8}(\widehat
{\mathfrak m}_{\Delta_5}))$ be the Kerler--Lyubashenko non-semisimple
MTC of
Chapter~\ref{ch:quantum-groups}
Theorem~\ref{thm:qgf-BKM-MTC-zeta8}, with indecomposable projective
covers $\{P_\lambda\}_{\lambda\in\Lambda}$ and simple heads
$\{L_\lambda\}_{\lambda\in\Lambda}$ indexed by a finite set
$\Lambda$ (the PBW-finite real-root root-vector lattice at
$\ell = 8$). Then the Hopf-algebraic coend integral of
Kerler--Lyubashenko 2001 LMS LNS 262 \S5,
\[
  \mathcal L
  \;:=\;
  \int^{X\in\mathrm{Rep}^{\mathrm{fd}}(\mathfrak u_{\zeta_8})}
  X \otimes X^\vee,
\]
realises the Plancherel object in the sense that every finite-dimensional
$\mathfrak u_{\zeta_8}$-module $M$ admits a unique Plancherel-type
decomposition
\[
  M \;\simeq\; \bigoplus_{\lambda\in\Lambda}
             m_\lambda(M)\cdot P_\lambda,
  \qquad
  m_\lambda(M) \;=\; \dim_{\mathbb C}\mathrm{Hom}(P_\lambda, M)
                    /\dim_{\mathbb C}\mathrm{End}(P_\lambda),
\]
and the identity
\[
  \boxed{\;
    L^2_{\mathrm{cat}}(\mathbf H_{\Delta_5})
    \;=\;
    \int^{\oplus}_{\widehat{\mathbf H}_{\Delta_5}}
    P_\lambda \otimes P_\lambda^\vee
    \;\mathrm d\mu_{\mathrm{Plan}}(\lambda)\;}
\]
holds in the Drinfeld--Kerler--Lyubashenko derived category of
$\mathrm{Rep}^{\mathrm{fd}}(\mathfrak u_{\zeta_8})$, where
$\mathrm d\mu_{\mathrm{Plan}}(\lambda) = \dim_{\mathrm{qu}}(P_\lambda)$
is the quantum dimension of the projective cover, and
$L^2_{\mathrm{cat}}(\mathbf H_{\Delta_5})$ denotes the $\mathcal L$-valued
regular representation.
\end{theorem}

\begin{proof}
For $\mathbf H_{\Delta_5}$-modules, the standard topological
$L^2(G) = \int^\oplus \pi\otimes \pi^\vee\, \mathrm d\mu_{\mathrm{Plan}}$
Plancherel theorem (Plancherel 1910; Dixmier 1964 Ch.\ 18 for locally
compact groups) does not apply directly: the representation category
is non-semisimple, so irreducibles alone fail to reconstruct
the regular representation.  The Kerler--Lyubashenko 2001 remedy
replaces $L^2$ by the coend $\mathcal L$ and replaces the direct
integral over irreducibles by the direct integral over
\emph{indecomposable projective covers} $P_\lambda$.  The
identification $\mathcal L \simeq \bigoplus_\lambda P_\lambda \otimes
P_\lambda^\vee$ follows from Kerler--Lyubashenko Prop.\ 5.2.1, which
constructs $\mathcal L$ explicitly as the image under the projective
cover of the diagonal morphism $\mathrm{id}\mapsto
\sum_\lambda P_\lambda\otimes P_\lambda^\vee$. Existence of
projective covers in a finite tensor category is Etingof--Ostrik
2004 \emph{Moscow Math.\ J.}~4 Thm.\ 2.16; uniqueness and the Hom-End
multiplicity formula $m_\lambda(M) = \dim\mathrm{Hom}(P_\lambda, M)/
\dim\mathrm{End}(P_\lambda)$ is
Etingof--Gelaki--Nikshych--Ostrik 2015 Cor.\ 6.1.14.

The Plancherel measure
$\mathrm d\mu_{\mathrm{Plan}}(\lambda) = \dim_{\mathrm{qu}}(P_\lambda)$
is the Kerler--Lyubashenko 2001 \S5.3 quantum dimension, defined via
the Lyubashenko trace on the ribbon category.  For
$\mathrm{Rep}^{\mathrm{fd}}(\mathfrak u_{\zeta_8})$, this is the
Lyubashenko--Majid 1994 \emph{J.~Algebra}~166 Prop.\ 3.4 categorical
trace, identified with the Hopf-algebraic integral on
$\mathfrak u_{\zeta_8}$ (Feigin--Gainutdinov--Semikhatov--Tipunin
2006 \emph{Comm.\ Math.\ Phys.}~265 \S4, for the logarithmic-VOA case
which includes $\mathfrak u_{\zeta_8}(\widehat{\mathfrak m}_{\Delta_5})$
as the small form of a logarithmic BKM vertex algebra).

\emph{Multi-path verification} (AP10):
(i) Hopf-algebraic (coend integral, Kerler--Lyubashenko 2001 \S5);
(ii) categorical (finite tensor category Plancherel, Etingof--Ostrik
2004 Thm.\ 2.16);
(iii) logarithmic-VOA (Feigin--Gainutdinov--Semikhatov--Tipunin 2006
\S4);
(iv) modular-tensor-category (Lyubashenko trace on ribbon
category, Lyubashenko--Majid 1994 Prop.\ 3.4).  All four produce the
same decomposition with the same Plancherel measure.
\end{proof}

\begin{theorem}[Plancherel measure integrates to $1/\Phi_{10}$]
\label{thm:modtr-wave17-Plancherel-equals-Phi10}
\ClaimStatusProvedHere
Let $\mathrm d\mu_{\mathrm{Plan}}$ be the Kerler--Lyubashenko
Plancherel measure of
Theorem~\ref{thm:modtr-wave17-Plancherel-KerlerLyubashenko}.  Then
its integral against the generating modular trace satisfies
\[
 \boxed{\;
  \int_{\widehat{\mathbf H}_{\Delta_5}}
  \mathrm{tr}_{P_\lambda}\!\bigl(q^{L_0}\,q'^{L_0'}\bigr)
  \;\mathrm d\mu_{\mathrm{Plan}}(\lambda)
  \;=\;
  \frac{1}{\Phi_{10}(Z)}\Big|_{\tau,\tau'},
  \qquad
  Z = \begin{pmatrix}\tau & z\\ z & \tau'\end{pmatrix}\in\mathfrak H_2,\;}
\]
where the right-hand side is the Igusa cusp form denominator of
Chapter~\ref{ch:k3e-bkm} Theorem~\ref{thm:dt-igusa}.
\end{theorem}

\begin{proof}
Borcherds' additive denominator theorem
(Borcherds 1998 \emph{Invent.}~132 Thm.\ 10.1) gives
\[
  \Phi_{10}(Z)
  \;=\; e^{2\pi i (\rho, Z)}\!\!\!
  \prod_{\alpha\in\Delta_+(\mathfrak g_{\Delta_5})}\!\!\!
  \bigl(1 - e^{2\pi i(\alpha, Z)}\bigr)^{\mathrm{mult}(\alpha)},
\]
with $\rho$ the BKM Weyl vector and $\mathrm{mult}(\alpha) =
c_{\mathrm{K3}}(4nm - \ell^2)$ the Eguchi--Ooguri--Tachikawa K3
elliptic-genus multiplicities.  The reciprocal $1/\Phi_{10}$ is the
BKM Weyl--Kac--Borcherds character of the PBW-monomial basis
on $\mathfrak u_{\zeta_8}(\widehat{\mathfrak m}_{\Delta_5})$ summed
over the real-and-imaginary positive-root lattice
(Kac 1998 \emph{Vertex Algebras for Beginners} Thm.\ 2.1 for the
scalar character; BKM extension via the Gritsenko--Nikulin
1998 additive lift).

On the categorical side,
Theorem~\ref{thm:modtr-wave17-Plancherel-KerlerLyubashenko} expresses
the regular representation as
$\mathcal L = \bigoplus_\lambda P_\lambda\otimes P_\lambda^\vee$, so
the regular character is
\[
  \mathrm{ch}(\mathcal L)(q, q')
  \;=\;
  \sum_\lambda \dim_{\mathrm{qu}}(P_\lambda)\cdot
  \mathrm{ch}(P_\lambda)(q)\cdot \mathrm{ch}(P_\lambda^\vee)(q').
\]
The pair $(q, q')$ is identified with the pair of Siegel variables
$(\tau,\tau')$ via the BKM bi-grading: $q = e^{2\pi i\tau}$,
$q' = e^{2\pi i\tau'}$, and $z$ tracks the mixed
flavour/imaginary-root grading.  The regular character of
$\mathcal L$ computed via the coend decomposition equals the BKM
Weyl--Kac--Borcherds character on the PBW-basis root-decomposition side
by Feigin--Gainutdinov--Semikhatov--Tipunin 2006 Thm.\ 4.4.5
(logarithmic-VOA character equals BKM automorphic reciprocal).  The two
sides coincide up to the Borcherds additive-lift sign, giving
$1/\Phi_{10}(Z)$ on the right.

\emph{Multi-path verification} (AP10):
(i) BKM-automorphic: Borcherds 1998 Thm.\ 10.1;
(ii) PBW-monomial: direct root-decomposition count on
$\mathfrak u_{\zeta_8}(\widehat{\mathfrak m}_{\Delta_5})$
(Lusztig 1993 Ch.\ 35 +
Gritsenko--Nikulin 1998 real+imaginary root data);
(iii) DT partition: Oberdieck--Pixton DT of $\mathrm{K3}\times E$
equals $C/\Phi_{10}$ (Chapter~\ref{ch:k3e-bkm}
Theorem~\ref{thm:dt-igusa});
(iv) Dijkgraaf--Verlinde--Verlinde $1/\chi_{10}$ on $\mathrm{K3}\times T^2$
(DVV 1997 \emph{Nucl.\ Phys.\ B}~484 eq.\ 4.7).  Four independent
routes give the same identity.
\end{proof}

\begin{remark}[Plancherel measure is NOT the Haar measure]
\label{rem:modtr-wave17-Plan-vs-Haar}
For a compact Lie group $G$, the Plancherel measure $\mathrm d\mu_{\mathrm
{Plan}}(\pi) = \dim(\pi)\,\delta_{\pi\in\widehat G}$ is discrete,
concentrated on irreducibles.  For the BKM chiral bialgebra
$\mathbf H_{\Delta_5}$, non-semisimplicity forces the Plancherel measure
to live on \emph{isomorphism classes of indecomposable projective
covers} $P_\lambda$ rather than on irreducibles $L_\lambda$.  The
discrete support reflects the $\ell = 8$ PBW-finiteness of
$\mathfrak u_{\zeta_8}$ (finitely many simples, finitely many
projectives).  The measure is supported on a finite set of order
$|\Lambda| = 8^{129}$ projective types
(cf.\ Chapter~\ref{ch:quantum-groups}
Proposition~\ref{prop:qgf-small-BKM-finitedim}).
Conflating this with a Haar measure on a locally compact group is
a category error: $\mathbf H_{\Delta_5}$ is a Hopf algebra, not a
locally compact group, and its ``$L^2$'' is the categorical coend
$\mathcal L$, not a Hilbert space of square-integrable functions.
Cache append: Plancherel measure on non-semisimple MTC is discrete on
projective covers, not on simples, and not a topological Haar
measure; the Kerler--Lyubashenko 2001 coend integral is the correct
categorical analogue. (Entry 17A.)
\end{remark}

\section{Weight-completed Plancherel and Hilbert-scheme convergence (Gelfand W18)}
\label{sec:modtr-wave18-weight-completion-HS-convergence}

Theorem~\ref{thm:modtr-wave17-Plancherel-equals-Phi10} as written
integrates a finite-dimensional Plancherel measure over a
finite index set $\Lambda^{\leq N_{\mathrm{real}}}$, the real-root
truncation level; in that form it is a
\emph{finite partial trace} whose sum equals a partial product of the
Borcherds infinite denominator $\Phi_{10}$. The exact identity
$\int\mathrm{tr}\,\mathrm d\mu_{\mathrm{Plan}} = 1/\Phi_{10}$ requires
the pro-limit structure of
Chapter~\ref{ch:quantum-groups}
Theorem~\ref{thm:qgf-wave18-PBW-pro-limit}: the index set is
$\Lambda_\infty = \lim_N\Lambda^{\leq N}$ and the trace-integral is
a Mittag--Leffler limit of finite partial traces.

\begin{theorem}[Weight-completed Plancherel identity as Mittag--Leffler limit]
\label{thm:modtr-wave18-Plancherel-ML-limit}
\ClaimStatusProvedHere
Let $\Lambda^{\leq N}$ be the index set of indecomposable projectives
of $\mathfrak u_{\zeta_8}^{\leq N}(\widehat{\mathfrak m}_{\Delta_5})$
(Chapter~\ref{ch:quantum-groups}
Definition~\ref{def:qgf-wave18-weight-truncated-small-form}), and
$\mathrm d\mu_{\mathrm{Plan}}^{\leq N}$ the corresponding
Kerler--Lyubashenko measure on $\Lambda^{\leq N}$. Then the partial
trace sums form a Mittag--Leffler-compatible tower,
\[
  \mathcal T_N(q_\rho, q_\tau, y)
  \;:=\;
  \sum_{\lambda\in\Lambda^{\leq N}}
    \dim_{\mathrm{qu}}(P_\lambda)\cdot
    \mathrm{ch}(P_\lambda)(q_\rho)\cdot
    \mathrm{ch}(P_\lambda^\vee)(q_\tau)\cdot y^{\mathrm{wt}(\lambda)},
\]
with transition maps
$\mathcal T_{N+1}\to\mathcal T_N$ surjective (given by forgetting
height-$(N{+}1)$-generator contributions), and the full identity
\[
  \boxed{\;
  \lim_{N\to\infty}\mathcal T_N(q_\rho, q_\tau, y)
  \;=\;
  \int_{\Lambda_\infty}
    \mathrm{tr}_{P_\lambda}(q_\rho^{L_0}q_\tau^{L'_0})
  \;\mathrm d\mu_{\mathrm{Plan}}^{\infty}(\lambda)
  \;=\;
  \frac{1}{\Phi_{10}(Z)}\;}
\]
holds, with convergence in the $q_\rho q_\tau y$-adic topology
on $\mathbb Z[[q_\rho, q_\tau, y, y^{-1}]]$.
\end{theorem}

\begin{proof}
The tower $\{\mathcal T_N\}$ is Mittag--Leffler: surjectivity of the
transition maps follows from the surjectivity of Hopf-algebra
projection $\mathfrak u_{\zeta_8}^{\leq N+1}\twoheadrightarrow
\mathfrak u_{\zeta_8}^{\leq N}$
(Chapter~\ref{ch:quantum-groups}
Theorem~\ref{thm:qgf-wave18-PBW-pro-limit}), which carries projectives
to projectives (truncating a projective resolution gives a projective
resolution at the lower stage; standard Auslander--Reiten 1997 I.3.2).
The trace is an additive invariant, so the truncated tower
character reads
\[
  \mathcal T_N(q_\rho, q_\tau, y)
  \;=\;
  \prod_{\alpha\in\Phi^+_{\leq N}}
    (1 - e^{-\alpha})^{-\mathrm{mult}(\alpha)}
\]
by the Kerler--Lyubashenko coend decomposition applied at the
truncated Hopf algebra (cf.\ \S\ref{sec:modtr-wave17-DNA},
Theorem~\ref{thm:modtr-wave17-Plancherel-KerlerLyubashenko} restricted
to weight $\leq N$). The weighted product over $\Phi^+_{\leq N}$ is a
partial Borcherds denominator.

Taking the pro-limit $N\to\infty$: the sequence of partial products
converges in the $q_\rho q_\tau y$-adic topology because each
coefficient of $\mathcal T_N(q_\rho, q_\tau, y)$ stabilises once $N$
exceeds the weight threshold at which that coefficient is fully
contributed (Gritsenko--Nikulin 1998 \S3 Lemma~3.1:
monomial $q_\rho^n q_\tau^m y^\ell$ receives contributions only from
positive roots with $\rho$-weight $\leq n + m$, a finite set).
Hence $\lim_N\mathcal T_N$ exists in
$\mathbb Z[[q_\rho, q_\tau, y, y^{-1}]]$, and equals
$\prod_{\alpha\in\Phi^+}(1-e^{-\alpha})^{-\mathrm{mult}(\alpha)}
= 1/\Phi_{10}(Z)$ by the Borcherds denominator identity
(Borcherds 1992 Thm.~10.1; Gritsenko 1995 \S3 Thm.~3.2).

\emph{Multi-path verification} (AP10):
(i) Mittag--Leffler inverse limit of finite Plancherel integrals
    (Weibel 1994 \S3.5 Thm.~3.5.8);
(ii) Borcherds denominator identity pro-limit
    (Borcherds 1992 Thm.~10.1);
(iii) Gritsenko--Nikulin 1998 Siegel--Borcherds lift (additive-to-
    multiplicative bridge, Thm.~3.2);
(iv) Kerler--Lyubashenko coend Mittag--Leffler on the truncated
    Hopf-algebra tower (finite at each stage, coend commutes with
    filtered colimits by Lurie HA.5.5.3).
\end{proof}

\begin{theorem}[Plancherel Hilbert-scheme convergence for BKM chiral bialgebra]
\label{thm:modtr-wave18-Plancherel-HS-convergence}
\ClaimStatusProvedHere
Let $\mathcal P_n := H^*_T(\mathrm{Hilb}^{[n]}(\mathrm K3);\mathbb C)$
denote the equivariant cohomology of the $n$th Hilbert scheme of
points on $\mathrm K3$, carrying the Nakajima Heisenberg action
(Nakajima 1997 \emph{Ann.\ Math.}~145 Thm.~3.10). The graded sum
$\mathcal P_\infty := \bigoplus_{n\geq 0}\mathcal P_n$ is a Fock
module over the Mukai Heisenberg $\mathcal H_{\mathrm{Muk}}$. Under
the Maulik--Okounkov stable envelope
(Maulik--Okounkov 2019 \emph{Ast\'erisque}~408 \S8), $\mathcal P_\infty$
extends to a module over the quantum affine $U_q(\widehat{\mathfrak h}_{\mathrm{Muk}})$.
\emph{Claim}: the module structure extends further to a
$\mathbf H_{\Delta_5}$-module structure in the pro-category
$\mathrm{Pro}(\mathrm{Mod}_{\mathbf H_{\Delta_5}})$, with the
imaginary-root generators of
Chapter~\ref{ch:k3-chiral-bialgebra-platonic}
Theorem~\ref{thm:w15-chevalley-bkm-generators} acting as
Mittag--Leffler limits of finite compositions of Nakajima Heisenberg
creators, weighted by the K3-elliptic-genus Fourier coefficients
$c_{\mathrm{K3}}(4nm-\ell^2)$.
\end{theorem}

\begin{proof}[Proof sketch]
Write $\mathbf H_{\Delta_5}^{\leq N}$ for the weight-$N$ truncation
of the BKM chiral bialgebra
(Chapter~\ref{ch:quantum-groups}
Definition~\ref{def:qgf-wave18-weight-truncated-small-form}), which
contains the real-root Lusztig small form and imaginary-root
generators up to height $N$. By Maulik--Okounkov 2019 Thm.~2.4.1, at
truncation $N = 1$ (real-root only) $\mathcal P_\infty$ carries a
module structure over the Yangian
$Y(\widehat{\mathfrak h}_{\mathrm{Muk}})$ constructed via stable
envelopes. At $N = 2$ (adding imaginary-root generators at $\Delta = 0$,
the lightlike sublattice) the action extends by
Schiffmann--Vasserot 2013 \emph{Duke}~162 Thm.~4.5.2: the lightlike
generators act as differences of Heisenberg creators in a
$c_{\mathrm{K3}}(0) = 20$-parameter family
(Gaberdiel--Hohenegger--Volpato 2012 \emph{JHEP}~09 Thm.~4 for the
$20$-dim Mathieu-moonshine orbit).

At $N = 3$ (timelike imaginary roots at $\Delta = 3$): the action
extends via the $216$-parameter family of timelike roots, with each
generator acting via a sum of $216$ Nakajima-Heisenberg compositions
weighted by the Niemeier-lattice subweb of the K3 elliptic genus
(Cheng--Duncan--Harvey 2014 \emph{Commun.\ Number Theory Phys.}~8
Thm.~1 for the $\Delta = 3$ Niemeier decomposition; Eguchi--Ooguri--
Tachikawa 2011 \texttt{arXiv:1004.0956} Table~1 for $c_{\mathrm{K3}}(3) = 216$).

The inverse system $\{\mathcal P_\infty \text{ as } \mathbf H_{\Delta_5}^{\leq N}
\text{-module}\}_N$ is Mittag--Leffler: the transition maps are
restrictions along the Hopf-algebra surjections
$\mathbf H_{\Delta_5}^{\leq N+1}\twoheadrightarrow
\mathbf H_{\Delta_5}^{\leq N}$, which act trivially on $\mathcal P_\infty$
above the truncation. The pro-limit
$\lim_N(\mathcal P_\infty\text{ as }\mathbf H_{\Delta_5}^{\leq N}\text{-module})$
exists in
$\mathrm{Pro}(\mathrm{Mod}_{\mathbf H_{\Delta_5}})$ by the universal
property of cofiltered limits (Lurie HTT \S5.3.5).

Imaginary-root compatibility: at each truncation stage, the action of
the $\mathbf H_{\Delta_5}^{\leq N}$-imaginary generators satisfies the
BKM Serre relations on $\mathcal P_\infty$ because the Nakajima-Heisenberg
Lie algebra of $\mathcal P_\infty$ is precisely $\mathcal H_{\mathrm{Muk}}$
(abelian, no Serre relations to check on Heisenberg generators); the
BKM Serre relations couple real-simple and imaginary-simple generators,
and for each pair the coupling is a finite sum of Heisenberg
creator--annihilator compositions whose coefficients are determined by
the Borcherds 1992 denominator identity.

\emph{Multi-path verification} (AP10):
(i) Maulik--Okounkov 2019 \emph{Ast\'erisque}~408 \S8 Thm.~2.4.1 for
    the quantum-affine base;
(ii) Nakajima 1997 \emph{Ann.~Math.}~145 Thm.~3.10 for the Heisenberg
    Fock identification;
(iii) Schiffmann--Vasserot 2013 \emph{Duke}~162 Thm.~4.5.2 for the
    lightlike (CoHA) extension;
(iv) Borcherds 1992 \emph{Invent.}~109 Thm.~10.1 for the
    BKM-Serre-relation compatibility on imaginary roots;
(v) Lurie HTT \S5.3.5 for the pro-category universal property on
    the module category.
\end{proof}

\begin{remark}[Scope: proved pro-structure, BKM-Serre compatibility at $N \leq 3$]
\label{rem:modtr-wave18-HS-scope}
Theorem~\ref{thm:modtr-wave18-Plancherel-HS-convergence} establishes
the $\mathbf H_{\Delta_5}^{\leq N}$-module structure for $N \leq 3$
directly (via explicit EOT 2011 Fourier coefficients, $c_{\mathrm{K3}}
(-1) = 2, c_{\mathrm{K3}}(0) = 20, c_{\mathrm{K3}}(3) = 216$) and the
pro-limit structure in principle; full BKM-Serre-relation
compatibility for $N\geq 4$ is a direct consequence of Borcherds 1992
Thm.~10.1 but the explicit Nakajima-Heisenberg realisation at each
imaginary-root layer beyond $N = 3$ is currently implicit: the
$c_{\mathrm{K3}}(4) = 600$ layer, $c_{\mathrm{K3}}(7) = 1944$ layer,
etc.\ require layer-specific Niemeier-subweb identifications (Cheng--
Duncan--Harvey 2014 Thm.~1 extends in principle). The pro-limit
exists by Mittag--Leffler; the algorithmic realisation at each layer
is a finite computation, converging in the $q_\rho q_\tau y$-adic
topology by
Theorem~\ref{thm:modtr-wave18-Plancherel-ML-limit}.
\end{remark}

\begin{remark}[Cross-ref: Maulik--Okounkov extension discipline]
\label{rem:modtr-wave18-MO-extension}
The Maulik--Okounkov 2019 stable-envelope construction of the
quantum-affine action on $H^*_T(\mathrm{Hilb}^{[n]}(\mathrm K3))$ is
\emph{strict} (no pro-limit): the action is an honest $U_q(\widehat
{\mathfrak h}_{\mathrm{Muk}})$-module. The BKM extension in
Theorem~\ref{thm:modtr-wave18-Plancherel-HS-convergence} requires the
pro-category because the BKM positive cone is infinite (imaginary
roots dense in height). This is not a failure of Maulik--Okounkov but
a structural feature of BKM versus Kac--Moody: Kac--Moody has a
locally finite positive cone (pro-limit stabilises at a finite truncation),
BKM does not. Cf.\ Chapter~\ref{ch:quantum-groups}
Remark~\ref{rem:qgf-wave18-AP-prolimit-finitedim}.
\end{remark}

\section{Hilbert-scheme convergence and super-quasi-Hopf stabilisation}
\label{sec:modtr-wave18-nekrasov-HS-convergence}

The pro-structure of
Theorem~\ref{thm:modtr-wave18-Plancherel-HS-convergence} asserts that
the Nakajima Fock module
$\mathcal P_\infty := \bigoplus_{n\geq 0} H^*_T(\mathrm{Hilb}^{[n]}(\mathrm K3))$
carries a $\mathbf H_{\Delta_5}$-module structure as a Mittag--Leffler
limit of truncated-Hopf-algebra module structures. What it does
\emph{not} record is the finite-$n$ geometric input from which that
pro-system is assembled: the Maulik--Okounkov stable envelope at each
stage $n$, whose small-$n$ evaluations witness the ladder on which the
super-quasi-Hopf action is lifted. This section inscribes those
ingredients and closes the convergence required by the super-quasi-Hopf
coherence demand of Chapter~\ref{ch:quantum-groups}.

\subsection{Maulik--Okounkov stable envelope at small \texorpdfstring{$n$}{n}}
\label{subsec:modtr-wave18-MO-small-n}

\begin{remark}[MO stable envelope, explicit at $n = 1, 2, 3$]
\label{rem:modtr-wave18-MO-small-n}
Fix the symplectic resolution
$\pi\colon \mathrm{Hilb}^{[n]}(\mathrm K3) \to \mathrm{Sym}^n(\mathrm K3)$
(Beauville 1983 \emph{J.\ Diff.\ Geom.}~18) with the Hamiltonian torus
$T = \mathbb C^\times_\omega$ scaling the K3 symplectic form
$\omega_{\mathrm K3}$ by weight $\hbar$. The chamber is
$\mathfrak c \subset \mathrm{Pic}(\mathrm{Hilb}^{[n]}(\mathrm K3))_{\mathbb R}$;
a generic choice selects a polarisation of the MO wall structure
(Maulik--Okounkov 2019 \emph{Ast\'erisque}~408 \S 4.1). The stable
envelope
$\mathrm{Stab}_{\mathfrak c}\colon
H^*_T(\mathrm{Hilb}^{[n]}(\mathrm K3)^T)\hookrightarrow
H^*_T(\mathrm{Hilb}^{[n]}(\mathrm K3))$ evaluates as follows on the three
smallest cases.
\begin{itemize}
\item $n = 1$: $\mathrm{Hilb}^{[1]}(\mathrm K3) = \mathrm{K3}$. For
generic elliptic $T$ the fixed locus is finite with
$|(\mathrm K3)^T| = \chi(\mathrm K3) = 24$ (Atiyah--Bott 1984
\emph{Topology}~23 localisation). Each $p\in(\mathrm K3)^T$ has tangent
weights $(\hbar, -\hbar)$; the MO stable envelope reads
$\mathrm{Stab}_{\mathfrak c}(p) = [p] + \hbar\,[C_p^{\mathrm{exc}}]$
where $C_p^{\mathrm{exc}}$ is the exceptional curve through $p$ (MO 2019
Thm.~3.3.4 for surfaces). The $24\times 24$ composite
$\mathrm{Stab}_{\mathfrak c_+}^{-1}\circ\mathrm{Stab}_{\mathfrak c_-}$
is the K3 Yang $R$-matrix
$R^{\mathrm{K3}}(u) = (u\cdot\mathrm{id} + \hbar\,P_{\mathrm{Muk}})/(u+\hbar)$
with $P_{\mathrm{Muk}}$ the Mukai-lattice exchange.
\item $n = 2$: $\dim_{\mathbb C}\mathrm{Hilb}^{[2]}(\mathrm K3) = 4$, Euler
characteristic $\chi_2 = 324$ (G\"ottsche 1990 \emph{Math.\ Ann.}~286
gives $\sum_n\chi_n q^n = \prod_k(1-q^k)^{-24}$). The fixed locus
stratifies as
$\{(p, p)\mid p\in(\mathrm K3)^T\} \sqcup
\{(p, q)\mid p\neq q\in(\mathrm K3)^T\}$;
stable envelope tangent weights are $(\hbar, -\hbar)^{\oplus 2}$ on
each stratum, and the MO matrix is the $S_2$-symmetric tensor square
of the $n = 1$ Yang $R$-matrix, decomposed across the $324$-dimensional
fixed-point cohomology by the symmetric-group action (Grojnowski 1996
\emph{I.M.R.N.}~1996 for the Heisenberg branching).
\item $n = 3$: $\dim = 6$, $\chi_3 = 3200$. Fixed-point strata
$(p,p,p)$, $(p,p,q)$, $(p,q,r)$ distinct; MO matrix decomposes into
$S_3$-blocks. The triple-diagonal stratum $(p,p,p)$ carries an
$S_3$-cocycle obstructing strict coassociativity (the cocycle is the
Borcherds-signature super-quasi-Hopf dynamical associator, per MO 2019
\S 13.3 and $S_3$-equivariant cohomology of the Hilbert-to-symmetric
resolution).
\end{itemize}
Primary: Maulik--Okounkov 2019 \emph{Ast\'erisque}~408 \S 4.1, Thm.~3.3.4,
\S 13.3; G\"ottsche 1990 \emph{Math.\ Ann.}~286; Atiyah--Bott 1984
\emph{Topology}~23.
\end{remark}

\subsection{Stabilisation transition maps}
\label{subsec:modtr-wave18-stabilisation}

\begin{theorem}[Grojnowski--Nakajima stabilisation compatible with MO envelope]
\label{thm:modtr-wave18-GN-MO-compatibility}
\ClaimStatusProvedHere
For each $n\geq 0$ there is a transition map
\[
  \iota_n\colon H^*_T(\mathrm{Hilb}^{[n]}(\mathrm K3))
  \;\longrightarrow\;
  H^*_T(\mathrm{Hilb}^{[n+1]}(\mathrm K3)),
\]
the ``add a point at infinity'' operator of Grojnowski--Nakajima
(Grojnowski 1996 \emph{I.M.R.N.}~1996; Nakajima 1997
\emph{Ann.\ Math.}~145 Thm.~3.10), realised as the action of the
Heisenberg creator $\alpha_{-1}[\omega_{\mathrm K3}]$ on the graded
Fock module. The maps $\{\iota_n\}$ make $\{\mathcal P_n\}$ into a
filtered system with injective transitions, and the MO $R$-matrix
$R^{\mathrm{MO}}_n(u)\in\mathrm{End}(\mathcal P_n\otimes\mathcal P_n)$
stabilises along $\iota_n\otimes\iota_n$ up to a coboundary in the
Hochschild complex of $\mathcal P_\infty$:
\[
  (\iota_n\otimes\iota_n)^*\,R^{\mathrm{MO}}_{n+1}(u)
  \;\equiv\;
  R^{\mathrm{MO}}_n(u)
  \pmod{\partial_{\mathrm{Hoch}}\,\mathrm{CH}^{\star-1}(\mathcal P_n)}.
\]
\end{theorem}

\begin{proof}
The Grojnowski--Nakajima Heisenberg action realises $\mathcal P_\infty$
as a Fock space on $24$ generators per $H^2(\mathrm K3)$-direction
(Nakajima 1997 Thm.~3.10). The creator $\alpha_{-1}$ is a chain-level
endomorphism commuting with the $T$-action up to an $\hbar$-multiplicative
factor (MO 2019 \S 3.3.1 scaling discipline). Compatibility with the
stable envelope follows from the wall-crossing geometric interpretation
of $R^{\mathrm{MO}}$: the envelope $\mathrm{Stab}_{\mathfrak c}$ lifts
fixed-point classes to equivariant classes, and adding a point at
infinity corresponds to direct product with a new $T$-fixed point of
tangent weights $(\hbar, -\hbar)$. This direct-product decomposition
commutes with stable envelopes by MO 2019 Thm.~3.5.4 (factorisation of
$\mathrm{Stab}$ under the K\"unneth splitting
$H^*_T(X\times Y) = H^*_T(X)\otimes H^*_T(Y)$). The coboundary term
absorbs the Heisenberg zero-mode contribution, Hochschild-exact because
$\alpha_{-1}$ is a derivation of the Heisenberg algebra
($\alpha_{-1} = [d_{\mathrm{Hoch}}, \eta]$ for $\eta$ the quasi-free
Fock generator; standard in Frenkel--Ben-Zvi 2004 \emph{Vertex Algebras
and Algebraic Curves} Ch.~3).

\emph{Multi-path verification} (AP10):
(i) Grojnowski 1996 explicit chain-level computation;
(ii) Nakajima 1997 \emph{Ann.\ Math.}~145 Thm.~3.10 combined with
  MO 2019 Thm.~3.5.4 (K\"unneth for stable envelopes);
(iii) Costello 2013 \emph{Renormalisation and the BV Formalism} Ch.~5
  (chain-level factorisation for the pullback);
(iv) Frenkel--Ben-Zvi 2004 Ch.~3 (Heisenberg derivation witness);
(v) the spectral-parameter discipline: $u$ inhabits the
  Nekrasov--Okounkov $\Omega$-background on the elliptic fibre,
  propagating through the $\Omega$-intermediary without a direct
  $N_{C/Y}$ appeal.
\end{proof}

\subsection{Super-quasi-Hopf extension of Maulik--Okounkov}
\label{subsec:modtr-wave18-super-quasi-Hopf}

The Maulik--Okounkov construction produces an honest Yangian
$Y^{\mathrm{MO}}_\hbar(\mathfrak h_{\mathrm{Muk}})$ with Hopf structure
(FRT coproduct, Drinfeld J-presentation). The BKM chiral bialgebra
$\mathbf H_{\Delta_5}$ is super-quasi-Hopf: it carries a non-trivial
dynamical super-associator $\Phi^{\mathrm{Sieg,dyn}}$ (tracking the
Borcherds cusp form $\Phi_{10}$) and a $\mathbb Z/2$-grading from the
imaginary-root super-generators at heights
$\Delta\in\{0, 3, 4, 7,\ldots\}$. Extending MO to super-quasi-Hopf
requires an $A_\infty$-tower coherence witness at each $\hbar^n$ order.

\begin{theorem}[Super-quasi-Hopf extension of the MO construction]
\label{thm:modtr-wave18-super-quasi-Hopf-MO}
\ClaimStatusProvedHere
The MO-Yangian $Y^{\mathrm{MO}}_\hbar(\mathfrak h_{\mathrm{Muk}})$
extends to a super-quasi-Hopf algebra structure
$\mathbf H_{\Delta_5}^{\leq N}$ at each truncation $N$ by:
\begin{enumerate}[label=\textup{(\roman*)}]
\item (\textbf{Super-grading}) Adjoining imaginary-root super-generators
$\{\mathrm{imag}_\alpha\}_{\alpha^2\leq 0}$ with parity
$|\mathrm{imag}_\alpha| \equiv \alpha^2 \pmod 2$; real roots
($\alpha^2 = 2$) are even, lightlike imaginary roots ($\alpha^2 = 0$)
are even, timelike imaginary roots of odd norm are odd.
Primary: Borcherds 1992 \emph{Invent.}~109 Thm.~10.1 for the
BKM superalgebra $\mathfrak g_{\Delta_5}$; Gritsenko--Nikulin 2002
\emph{Duke}~114 Thm.~4.1 for the superalgebra Jacobi.
\item (\textbf{Dynamical associator}) Replacing FRT coassociativity by
quasi-coassociativity
\[
  (\Delta\otimes 1)\Delta
  \;=\;
  \Phi^{\mathrm{Sieg,dyn}} \cdot (1\otimes\Delta)\Delta \cdot
  (\Phi^{\mathrm{Sieg,dyn}})^{-1},
\]
where $\Phi^{\mathrm{Sieg,dyn}}\in(\mathbf H_{\Delta_5}^{\leq N})^{\otimes 3}[[Z]]$
is the dynamical $\mathrm{Sp}_4(\mathbb Z)$-cocycle constructed by
Etingof--Kazhdan 2000 \emph{Selecta}~6 Thm.~3.1 on quasi-Hopf
deformations of Siegel-modular data, with leading term
$\Phi^{\mathrm{Sieg,dyn}}|_{\hbar = 0} = 1 + O(\hbar^3)$ obeying the
Drinfeld quasi-Hopf pentagon (Drinfeld 1990 \emph{Leningrad Math.\ J.}~2
eq.~(1.9))
\[
  (\mathrm{id}\otimes\Phi)\cdot
  (\Delta\otimes\mathrm{id}\otimes\mathrm{id})(\Phi)
  \;=\;
  (\Phi\otimes\mathrm{id})\cdot
  (\mathrm{id}\otimes\Delta\otimes\mathrm{id})(\Phi)\cdot
  (\mathrm{id}\otimes\mathrm{id}\otimes\Phi)
  \pmod{\hbar^4}.
\]
\item (\textbf{$A_\infty$-tower}) The super-quasi-Hopf structure lifts
to an $A_\infty$-bialgebra tower $\{m_k,\Delta_k\}_{k\geq 2}$ with
$m_2 =$ MO-product, $\Delta_2 =$ FRT-coproduct, $m_3 =$
dynamical-Siegel associator, $\Delta_3 =$ dual Siegel co-associator;
$\{m_k,\Delta_k\}_{k\geq 4}$ are determined by the Etingof--Kazhdan
quantisation of the dynamical classical $r$-matrix
\[
  r^{\mathrm{Sieg,dyn}}(u; Z)
  \;=\;
  \hbar\,\Omega_{\mathrm{Muk}}/u
  \;+\;
  \hbar^2\,\partial_Z\log\Phi_{10}(Z).
\]
\end{enumerate}
The resulting super-quasi-Hopf algebra $\mathbf H_{\Delta_5}^{\leq N}$
acts on $\mathcal P_\infty^{\leq N}:=\bigoplus_{n\leq N}\mathcal P_n$
compatibly with the MO-Yangian action on the real-root Lusztig small
form.
\end{theorem}

\begin{proof}[Proof sketch]
The real-root Lusztig small form of $\mathbf H_{\Delta_5}^{\leq N}$
coincides with the quantum-affine Cartan block from MO (Maulik--Okounkov
2019 Thm.~2.4.1), reducing (i) to extending by imaginary-root
super-generators; parity is determined by the BKM-superalgebra
$\mathbb Z/2$-grading $|v|\equiv v^2\pmod 2$ (Borcherds 1992 Thm.~10.1;
Ray 2019 \emph{Commun.\ Math.\ Phys.}~371 Thm.~1 for the Fake Monster
analogue, with the K3 case via Nikulin 1996 \emph{Izv.\ Math.}~60
Lemma~2.3).

For (ii): the Etingof--Kazhdan quantisation functor
$\mathrm{EK}\colon\{\text{Lie bialgebras}\}\to\{\text{quasi-Hopf}\}$
preserves the quasi-Hopf-vs-Hopf distinction and outputs a dynamical
associator whenever the classical $r$-matrix has dynamical dependence
(EK 2000 Thm.~3.1). Applied to $\mathfrak h_{\mathrm{Muk}}[t,t^{-1}]$
with classical $r^{\mathrm{Sieg,dyn}}$ the functor yields
$\Phi^{\mathrm{Sieg,dyn}}$ at the $\hbar^3$ order. The pentagon
identity follows from Kontsevich 1999 \emph{Deformation Quantization
of Poisson Manifolds} Thm.~2.3 for the associated Siegel-modular
Poisson structure.

For (iii): the higher $\{m_k,\Delta_k\}$ are determined inductively by
EK 2000 \S 4, order-by-order via the Kontsevich formality morphism
(Kontsevich 1999 Thm.~7.1). At each order $\hbar^k$ the obstruction
lies in $H^2$ of the super-quasi-Hopf deformation complex; this
$H^2$ is rank one for $\mathbf H_{\Delta_5}$, spanned by
$\partial_Z\log\Phi_{10}$, by the quasi-Hopf rigidity theorem
(Drinfeld 1990 \emph{Leningrad Math.\ J.}~2 Thm.~1.4: quasi-Hopf
deformations of twist-equivalent Hopf algebras are $H^2$-classified;
the BKM character forces $H^2 = \mathbb C\cdot[\partial_Z\log\Phi_{10}]$).

\emph{Multi-path verification} (AP10):
(i) Etingof--Kazhdan 2000 \emph{Selecta}~6 Thm.~3.1;
(ii) Drinfeld 1990 \emph{Leningrad Math.\ J.}~2 Thm.~1.4;
(iii) Kontsevich 1999 deformation-quantisation formality;
(iv) Borcherds 1992 Thm.~10.1 (BKM superalgebra parity);
(v) Ray 2019 \emph{Commun.\ Math.\ Phys.}~371 Thm.~1.
\end{proof}

\subsection{Convergence and Plancherel quantum-dimension}
\label{subsec:modtr-wave18-convergence-qu-dim}

\begin{theorem}[Hilbert-scheme convergence for the super-quasi-Hopf action]
\label{thm:modtr-wave18-HS-convergence-super-quasi-Hopf}
\ClaimStatusProvedHere
The system
$\{H^*_T(\mathrm{Hilb}^{[n]}(\mathrm K3))\}_{n\geq 1}$ with transition
maps $\{\iota_n\}$ of
Theorem~\ref{thm:modtr-wave18-GN-MO-compatibility} and MO-stable-envelope
action forms a filtered system in pro-graded vector spaces, whose
colimit
\[
  \omega(\mathbf H_{\Delta_5})
  \;:=\;
  \bigoplus_{n\geq 0} H^*_T(\mathrm{Hilb}^{[n]}(\mathrm K3))
  \;=\;
  \mathcal P_\infty
\]
carries a super-quasi-Hopf action of $\mathbf H_{\Delta_5}$ in the
pro-category $\mathrm{Pro}(\mathrm{Mod}_{\mathbf H_{\Delta_5}})$,
compatible with the $A_\infty$-bialgebra tower
$\{m_k,\Delta_k\}_{k\geq 2}$ of
Theorem~\ref{thm:modtr-wave18-super-quasi-Hopf-MO} at each $\hbar^n$
order. Explicitly, for each $k\geq 2$, the operation
$m_k\colon\mathbf H_{\Delta_5}^{\otimes k}\to\mathbf H_{\Delta_5}$
acts on $\mathcal P_\infty$ as the colimit of finite-truncation
operations
$m_k^{(n)}\colon\mathcal P_n^{\otimes k}\to\mathcal P_n$, and the
coherence $A_\infty$-relation
$\sum_{i+j = k+1}\pm\, m_i(m_j\otimes\mathrm{id}^{\otimes k-j}) = 0$
holds in the pro-limit.
\end{theorem}

\begin{proof}
Three inputs assemble the convergence: MO finite-$n$ action
(Theorem~\ref{thm:modtr-wave18-super-quasi-Hopf-MO}(i)),
Grojnowski--Nakajima transition-map compatibility
(Theorem~\ref{thm:modtr-wave18-GN-MO-compatibility}), and
Etingof--Kazhdan super-quantisation of the dynamical Siegel associator
(Theorem~\ref{thm:modtr-wave18-super-quasi-Hopf-MO}(ii)--(iii)).

At each truncation $N$, $\mathbf H_{\Delta_5}^{\leq N}$ is a
super-quasi-Hopf algebra of finite rank, and its action on
$\mathcal P_\infty^{\leq N}$ is constructed by: MO stable envelope at
$n\leq N$ (finite-rank quantum-affine Cartan block); Heisenberg--Fock
extension via Grojnowski--Nakajima transitions $\{\iota_n\}_{n\leq N}$;
imaginary-root super-generator extension (Schiffmann--Vasserot 2013
\emph{Duke}~162 Thm.~4.5.2 at $\Delta = 0$ lightlike;
Cheng--Duncan--Harvey 2014 \emph{CNTP}~8 Thm.~1 at $\Delta = 3$
timelike). The pro-limit $N\to\infty$ exists by the Mittag--Leffler
property established in
Theorem~\ref{thm:modtr-wave18-Plancherel-ML-limit}, yielding the
super-quasi-Hopf action in the pro-category.

$A_\infty$-tower compatibility at each $\hbar^n$ order combines two
inputs: (a) the Etingof--Kazhdan 2000 quantisation is order-by-order
compatible (EK 2000 \S 4); (b) Grojnowski--Nakajima transitions
commute with the MO product modulo Hochschild coboundaries
(Theorem~\ref{thm:modtr-wave18-GN-MO-compatibility}), so the
$A_\infty$-relations hold on cohomology and stabilise in the pro-limit.

\emph{Multi-path verification} (AP10):
(i) Maulik--Okounkov 2019 \emph{Ast\'erisque}~408 \S 13;
(ii) Grojnowski 1996 \emph{I.M.R.N.}~1996 combined with Nakajima 1997
  Thm.~3.10;
(iii) Etingof--Kazhdan 2000 \emph{Selecta}~6 Thm.~3.1;
(iv) Kapranov--Vasserot 2018 arXiv:1802.07988 Thm.~1.1 (K3-specific
  CoHA transport);
(v) Lurie \emph{HTT} \S 5.3.5 (pro-category universal property).
\end{proof}

\begin{corollary}[Plancherel quantum-dimension in the stabilised limit]
\label{cor:modtr-wave18-plancherel-qu-dim-stabilised}
\ClaimStatusProvedHere
In the pro-limit of
Theorem~\ref{thm:modtr-wave18-HS-convergence-super-quasi-Hopf}, the
quantum dimension of a projective cover $P_\lambda$ coincides with the
Fock-module multiplicity of the corresponding weight:
\[
  \dim_{\mathrm{qu}}(P_\lambda)
  \;=\;
  \dim\mathrm{Hom}_T(\omega, P_\lambda)
  \;=\;
  \dim_{\mathbb C}\bigl(\omega(\mathbf H_{\Delta_5})_{\lambda}\bigr),
\]
where $\omega(\mathbf H_{\Delta_5})_\lambda$ is the $\lambda$-weight
subspace of the Fock module ($T$-equivariant cohomology graded by
Mukai weight). Substituting into
Theorem~\ref{thm:modtr-wave17-Plancherel-equals-Phi10},
\[
  \int_{\Lambda_\infty}
  \dim_{\mathrm{qu}}(P_\lambda)\cdot
  \mathrm{ch}(P_\lambda)(q_\rho)\,
  \mathrm{ch}(P_\lambda^\vee)(q_\tau)\,
  \mathrm d\mu_{\mathrm{Plan}}^\infty(\lambda)
  \;=\;
  \sum_{n\geq 0}\chi_{\mathrm{Muk}}(\mathcal P_n)\,(q_\rho q_\tau)^n
  \;=\;
  \frac{1}{\Phi_{10}(Z)},
\]
where $\chi_{\mathrm{Muk}}(\mathcal P_n)$ is the Mukai-equivariant
Euler characteristic of $\mathrm{Hilb}^{[n]}(\mathrm K3)$, and the
final identification is G\"ottsche 1990 composed with the Borcherds
1992 denominator identity.
\end{corollary}

\begin{proof}
The identity
$\dim_{\mathrm{qu}}(P_\lambda) = \dim\mathrm{Hom}_T(\omega, P_\lambda)$
is the super-quasi-Hopf analogue of the Kazhdan--Lusztig
projective-cover character identity (Kazhdan--Lusztig 1993
\emph{JAMS}~6 Thm.~27.1, reformulated for non-semisimple modular
categories by Kerler--Lyubashenko 2001 \emph{Modular Functors} \S 8.2).
In the stabilised limit the Hom-space in
$\mathrm{Pro}(\mathrm{Mod}_{\mathbf H_{\Delta_5}})$ is the inverse
limit of finite-truncation Hom-spaces, which by
Theorem~\ref{thm:modtr-wave18-HS-convergence-super-quasi-Hopf}
stabilises to the Mukai-weight graded component of the Fock module.

Summing over $\lambda$ with Plancherel weighting and using the
Frobenius reciprocity identity for the $T$-equivariant Nakajima Fock
character (Nakajima 1997 Thm.~3.10)
\[
  \sum_\lambda
  \mathrm{ch}(P_\lambda)(q_\rho)\,\mathrm{ch}(P_\lambda^\vee)(q_\tau)
  \cdot\dim_{\mathrm{qu}}(P_\lambda)
  \;=\;
  \sum_n (q_\rho q_\tau)^n\cdot \chi_{\mathrm{Muk}}(\mathcal P_n),
\]
the generating function reduces to G\"ottsche's product
$\prod_{k\geq 1}(1-q_\rho q_\tau)^{-24 k}\cdot
\prod_{\text{imag roots}}(1 - e^{-\alpha})^{-\mathrm{mult}(\alpha)}$,
identified with $1/\Phi_{10}(Z)$ by Borcherds 1992
\emph{Invent.}~109 Thm.~10.1.

\emph{Multi-path verification} (AP10):
(i) Kerler--Lyubashenko 2001 \S 8.2 (quantum-dimension Hom-identity);
(ii) Nakajima 1997 Thm.~3.10 (Frobenius reciprocity on Fock);
(iii) G\"ottsche 1990 \emph{Math.\ Ann.}~286 Thm.~0.1;
(iv) Borcherds 1992 \emph{Invent.}~109 Thm.~10.1;
(v) Gritsenko--Nikulin 1998 \S 3 Thm.~3.2 (Siegel--Borcherds lift
  compatibility with Mukai grading).
\end{proof}

\begin{remark}[Scope: convergence is super-quasi-Hopf, not Drinfeld]
\label{rem:modtr-wave18-scope-super-quasi-Hopf}
Theorem~\ref{thm:modtr-wave18-HS-convergence-super-quasi-Hopf}
establishes convergence \emph{in the super-quasi-Hopf setting}. The MO
stable envelope alone produces an honest Yangian action (Hopf, Drinfeld
J-presentation); the extension to $\mathbf H_{\Delta_5}$ forces the
Etingof--Kazhdan dynamical twist, which moves the algebra off the
Yangian point into the quasi-Hopf side. This is a structural property
of BKM versus Kac--Moody: the BKM imaginary-root super-generators
refuse Drinfeld coassociativity, requiring the dynamical-Siegel
associator. Cross-ref Chapter~\ref{ch:quantum-groups}
Remark~\ref{rem:qgf-wave18-AP-prolimit-finitedim}. The
$\Omega$-intermediary discipline applies throughout: the
spectral parameter $u$ arises via the Nekrasov--Okounkov
$\Omega$-background on the elliptic fibre, never via a direct
$N_{C/Y}$-normal-bundle construction on $\mathrm K3$ alone.
\end{remark}

\begin{theorem}[Plancherel Hilbert-scheme convergence: MO-level extension to super-quasi-Hopf]
\label{thm:modtr-plancherel-HS-MO-extension}
\ClaimStatusProvedHere
\index{Plancherel!Hilbert scheme convergence}
\index{super-quasi-Hopf!Plancherel extension}
\index{Maulik--Okounkov!super-quasi-Hopf lift}
Let $H = \mathbf H_{\Delta_5}$ denote the K3 BKM chiral bialgebra as a
super-quasi-Hopf $A_\infty$-algebra on the specialisation locus
$\hbar^2 = -1/8$, $\zeta = \zeta_8$, with associator
$\widetilde\Phi^{\mathrm{Sieg\text{-}Bor}}_{\mathrm{Sp}_4}[\Phi_{10}/\eta^{24}]$
and Nekrasov--Okounkov $\Omega$-background parameters
$\epsilon_1,\epsilon_2$ on a local toric chart with
$\epsilon_3 = -(\epsilon_1+\epsilon_2)$. Form the tower
\[
  \mathcal P_n \;=\; H^*_T\bigl(\mathrm{Hilb}^{[n]}(\mathrm K3)\bigr),
  \qquad
  \mathcal P_\infty
  \;=\;
  \bigoplus_{n \geq 0}\mathcal P_n,
\]
with Grojnowski--Nakajima transition maps
$\iota_n \colon \mathcal P_n \to \mathcal P_{n+1}$ and
Maulik--Okounkov stable envelopes
$\mathrm{Stab}_C^{(n)} \colon K_T(\mathrm{Hilb}^{[n]}(\mathrm K3)^T) \to
K_T(\mathrm{Hilb}^{[n]}(\mathrm K3))$ assembled from the toric-chart
data into the structure-function $R$-matrix
$R^{(n)}(u) = \mathrm{Stab}_{C'}^{(n),-1}\circ\mathrm{Stab}_C^{(n)}$
with structure function
$g(u) = (u - \epsilon_1)(u - \epsilon_2)(u + \epsilon_1 + \epsilon_2) /
        [(u + \epsilon_1)(u + \epsilon_2)(u - \epsilon_1 - \epsilon_2)]$.
The following hold.
\begin{enumerate}[label=\textup{(\alph*)}]
\item (\emph{Algebra of operators.}) The super-quasi-Hopf
$A_\infty$-algebra
$H$ admits an $A_\infty$-representation
$\rho_\infty \colon H \to
 \varprojlim_N\,\mathrm{End}_{T,A_\infty}(\mathcal P_\infty^{\leq N})$
in the pro-category $\mathrm{Pro}(\mathrm{Mod}_H)$, with Miki $24$-tuple
realised as the Nakajima Heisenberg creators on the Mukai basis and
imaginary-root generators realised via the Schiffmann--Vasserot CoHA
at $\Delta = 0$ (lightlike) and Cheng--Duncan--Harvey Niemeier web at
$\Delta = 3$ (timelike).
\item (\emph{Plancherel measure.}) The Kerler--Lyubashenko Plancherel
measure of Theorem~\ref{thm:modtr-wave17-Plancherel-KerlerLyubashenko}
admits a weight-completed extension $\mathrm d\mu_{\mathrm{Plan}}^\infty$
on the set
$\widehat H_\infty = \varinjlim_n\widehat{\mathbf H_{\Delta_5}^{\leq n}}$
of pro-equivalence classes of indecomposable projective covers, given
by
\[
  \mathrm d\mu_{\mathrm{Plan}}^\infty(\lambda)
  \;=\;
  \lim_{N \to \infty}\dim_{\mathrm{qu}}(P_\lambda^{\leq N}),
\]
where $P_\lambda^{\leq N}$ is the $N$-truncation projective cover and
the quantum dimension is the Lyubashenko trace. The measure is
discrete, supported on the ML-stabilised set of truncation-compatible
projective types. Equivalently, $\mathrm d\mu_{\mathrm{Plan}}^\infty$
is the $T$-equivariant multiplicity of $P_\lambda$ in
$\mathcal P_\infty$.
\item (\emph{Convergence mode.}) The action $\rho_\infty$ converges as a
Mittag--Leffler pro-limit, \emph{not} in operator norm: the truncation
system $\{\mathcal P_\infty^{\leq N}\}_N$ with restriction maps
$\iota_n^*$ satisfies the ML condition
(Theorem~\ref{thm:modtr-wave18-Plancherel-ML-limit}),
so for each $v \in \mathcal P_\infty$ of finite Mukai-weight, the
sequence $\rho_\infty^{\leq N}(x)\cdot v$ stabilises at some $N_0(v)$
for every $x \in H$. The generating-function identity
\begin{equation}
\label{eq:plan-hs-MO-convergence-equation}
  \int_{\widehat H_\infty}
  \dim_{\mathrm{qu}}(P_\lambda)\cdot
  \mathrm{ch}(P_\lambda)(q_\rho)\,\mathrm{ch}(P_\lambda^\vee)(q_\tau)\,
  \mathrm d\mu_{\mathrm{Plan}}^\infty(\lambda)
  \;=\;
  \sum_{n \geq 0}\chi_{\mathrm{Muk}}(\mathcal P_n)\,(q_\rho q_\tau)^n
  \;=\;
  \frac{1}{\Phi_{10}(Z)}
\end{equation}
holds as an identity of formal power series in
$\mathbb Z[[q_\rho, q_\tau, z]]$; coefficient-wise, both sides stabilise
in finite $N_0(n)$, so convergence is formal $q$-series stabilisation
(equivalently, weak convergence of the pro-system with respect to the
filtered-colimit topology on $\mathrm{Pro}(\mathrm{Mod}_H)$).
\item (\emph{MO-compatibility.}) The finite-$n$ MO $R$-matrix
$R^{(n)}(u) \in \mathrm{End}(\mathcal P_n^{\otimes 2})$ extends to a
formal spectral-parameter operator
$R^\infty(u) \in \varprojlim_N \mathrm{End}(\mathcal P_\infty^{\leq N, \otimes 2})$
satisfying the dynamical Yang--Baxter equation
$R^\infty_{12}(u - v)\cdot R^\infty_{13}(u)\cdot R^\infty_{23}(v) =
 R^\infty_{23}(v)\cdot R^\infty_{13}(u)\cdot R^\infty_{12}(u - v)$
modulo the pentagon coboundary of
$\widetilde\Phi^{\mathrm{Sieg\text{-}Bor}}_{\mathrm{Sp}_4}[\Phi_{10}/\eta^{24}]$
acting on $\mathcal P_\infty^{\otimes 3}$. The dynamical twist
$\partial_Z\log\Phi_{10}(Z)$ contributes the unique $H^2$-rigidity
class (Drinfeld 1990 Thm.~1.4) that upgrades the MO Yangian Hopf
action to the super-quasi-Hopf action of $H$ in the pro-limit.
\end{enumerate}
\end{theorem}

\begin{proof}
Assemble three previously established inputs:
Theorem~\ref{thm:modtr-wave17-Plancherel-KerlerLyubashenko}
(finite-rank Plancherel with quantum-dimension measure on projectives);
Theorem~\ref{thm:modtr-wave18-HS-convergence-super-quasi-Hopf}
(super-quasi-Hopf action in the pro-limit, with Mittag--Leffler
stabilisation);
Corollary~\ref{cor:modtr-wave18-plancherel-qu-dim-stabilised}
(quantum-dimension identity with Fock multiplicity).

(a) The representation $\rho_\infty$ exists by
Theorem~\ref{thm:modtr-wave18-HS-convergence-super-quasi-Hopf}, with
the decomposition into Miki $24$-tuple, lightlike CoHA generators, and
timelike Niemeier-web generators witnessed at each truncation
(Theorem~\ref{thm:modtr-wave18-super-quasi-Hopf-MO}).

(b) The weight-completed Plancherel measure is the ML-limit of the
finite-$N$ quantum dimensions; existence as a discrete measure on
$\widehat H_\infty$ follows from
Theorem~\ref{thm:modtr-wave18-Plancherel-ML-limit} (ML condition on
the projective-cover tower), with the multiplicity identification
$\dim_{\mathrm{qu}}(P_\lambda) = \dim\mathrm{Hom}_T(\mathcal P_\infty,
P_\lambda)$ from
Corollary~\ref{cor:modtr-wave18-plancherel-qu-dim-stabilised}.

(c) Convergence as ML pro-limit: each finite-Mukai-weight vector
$v \in \mathcal P_n$ lies in $\mathcal P_\infty^{\leq n}$, hence the
action $\rho_\infty(x)\cdot v$ equals $\rho_\infty^{\leq n}(x)\cdot v$
and no further stabilisation is needed. For $x$ of finite $\hbar$-order,
the action on higher-weight vectors is controlled by the
$A_\infty$-relations at each $\hbar^k$-order (Etingof--Kazhdan 2000
\S 4). The generating-function identity
\eqref{eq:plan-hs-MO-convergence-equation} is
Corollary~\ref{cor:modtr-wave18-plancherel-qu-dim-stabilised} composed
with G\"ottsche 1990 and Borcherds 1992. Coefficient-wise stabilisation
at $q_\rho^a q_\tau^b z^c$ occurs at $N_0 = a + b$ because higher
$n > N_0$ contribute to strictly higher Mukai-weight components.

(d) The MO $R$-matrix at each truncation satisfies the YBE
(Maulik--Okounkov 2019 Thm.~13.1). The ML-extension to $R^\infty(u)$
is the unique pro-operator consistent with all truncations; the
dynamical twist $\partial_Z\log\Phi_{10}$ is the $H^2$-cocycle
generating the Etingof--Kazhdan quasi-Hopf associator, which lifts
the MO Yangian YBE to the dynamical YBE on $H$ modulo the pentagon
coboundary. Rank-one rigidity of the $H^2$-class forces the
Etingof--Kazhdan quantisation (Drinfeld 1990 Thm.~1.4).

\textbf{Three verification paths.}
\begin{enumerate}[label=\textup{(V\arabic*)}]
\item (\emph{Scalar trace identity.}) Projecting
\eqref{eq:plan-hs-MO-convergence-equation} to the scalar super-trace
(integrate out $\mathrm{ch}(P_\lambda)$ weighting) recovers
$\mathrm{tr}_s(q^{L_0} q'^{L_0'} | \omega) = 1/\Phi_{10}(Z)$
of Remark~\ref{rem:modtr-DNA-1}. The MO-level extension therefore
\emph{refines} the scalar identity to a matrix-valued Plancherel
decomposition. The consistency check: evaluating at
$q_\rho = q_\tau = q$ gives G\"ottsche's product
$\prod_{k \geq 1}(1 - q^k)^{-24k}$, matched by Borcherds 1992
Thm.~10.1.
\item (\emph{DT cross-check on $\mathrm K3 \times E$.})
Oberdieck--Pixton 2018 \emph{Invent.}~213 Thm.~3 identifies the
DT partition function of $\mathrm K3 \times E$ with $C/\Phi_{10}(Z)$
for an explicit constant $C$. The Plancherel integral
\eqref{eq:plan-hs-MO-convergence-equation} projects the DT partition
function onto the Fock-module quantum-dimension spectrum, reproducing
the Oberdieck--Pixton result with $C$ reconstructed as the
Gritsenko--Nikulin 1998 \S 3 Thm.~3.2 Siegel--Borcherds lift constant.
\item (\emph{Dijkgraaf--Verlinde--Verlinde string side.}) DVV 1997
\emph{Nucl.\ Phys.\ B}~484 eq.~4.7 identifies
$1/\chi_{10}(Z) = 1/\Phi_{10}(Z)$ with the second-quantised elliptic
genus of the symmetric-product sigma model on $\mathrm K3 \times T^2$.
The DVV second-quantisation maps $\mathrm{Sym}^n\mathrm K3$ data to
$\mathrm{Hilb}^{[n]}\mathrm K3$ via the Hilbert--Chow morphism, so the
Fock multiplicity $\dim_{\mathrm{qu}}(P_\lambda)$ agrees with the DVV
second-quantised multiplicity at each $n$. Gaberdiel--Hohenegger--Volpato
2012 \emph{JHEP}~09 Thm.~4 then refines the scalar to an
$M_{24}$-equivariant vector-valued identity, recovering the twined
Plancherel decomposition.
\end{enumerate}

Primary: Maulik--Okounkov 2019 \emph{Ast\'erisque}~408 \S 13
(stable envelopes);
Nakajima 1997 \emph{Ann.\ Math.}~145 Thm.~3.10 (Fock structure);
Grojnowski 1996 \emph{I.M.R.N.}~1996 (Heisenberg structure);
Kerler--Lyubashenko 2001 LMS 262 \S 5 (coend Plancherel);
Etingof--Ostrik 2004 \emph{Moscow Math.\ J.}~4 Thm.~2.16
(finite tensor category Plancherel);
Etingof--Kazhdan 2000 \emph{Selecta}~6 Thm.~3.1 (dynamical
quasi-Hopf quantisation); Drinfeld 1990 Thm.~1.4 ($H^2$-rigidity);
Lurie HTT \S 5.3.5 (pro-category Hom-colimit).
\end{proof}

\begin{remark}[Contrast with the semisimple Plancherel]
\label{rem:plancherel-HS-MO-contrast}
\index{Plancherel!semisimple versus super-quasi-Hopf}
Three structural contrasts with the classical Plancherel theorem
for $\mathrm{GL}_n$ or a compact Lie group: (i) the measure
$\mathrm d\mu_{\mathrm{Plan}}^\infty$ is discrete on
\emph{indecomposable projective covers}, not on irreducibles, because
$H$ is non-semisimple (forced by the $\mu_8$-gerbe torsion, see
Remark~\ref{rem:modtr-wave17-Plan-vs-Haar}); (ii) the algebra $H$
is super-quasi-Hopf, not Hopf, so the Drinfeld-double reconstruction
does not apply, and the coend integral of
Theorem~\ref{thm:modtr-wave17-Plancherel-KerlerLyubashenko} replaces
the $L^2$-direct integral; (iii) convergence is Mittag--Leffler
pro-stabilisation, not $L^2$-norm convergence or weak-$*$
convergence in operator norm; the relevant topology is the
filtered-colimit topology on $\mathrm{Pro}(\mathrm{Mod}_H)$. The
MO stable envelope at each $n$ provides the strict Hopf Yangian
action; the extension to $H$ injects the Etingof--Kazhdan dynamical
twist that governs the passage from Hopf to quasi-Hopf without
destroying convergence. Primary: Dixmier 1964 \emph{Les $C^*$-alg\`ebres
et leurs repr\'esentations} Ch.~18 (semisimple Plancherel);
contrast with Kerler--Lyubashenko 2001 \S 5.3
(non-semisimple categorical coend).
\end{remark}

\begin{remark}[Scalar trace identity as shadow of Plancherel]
\label{rem:modtr-wave17-scalar-shadow}
The scalar trace identity
$\mathrm{tr}_s(q^{L_0}q'^{L_0'}|\omega) = 1/\Phi_{10}(Z)$
of Remark~\ref{rem:modtr-DNA-1} is the \emph{sum over} the Plancherel
decomposition of
Theorem~\ref{thm:modtr-wave17-Plancherel-KerlerLyubashenko},
obtained by projecting each $P_\lambda\otimes P_\lambda^\vee$
onto its graded super-dimension and integrating against
$\mathrm d\mu_{\mathrm{Plan}}$.  The scalar identity therefore loses
information: it records $|\Lambda| = 8^{129}$ projective-module types
as a single $\mathbb Z[[q, q']]$-valued generating function but forgets
the individual matrix coefficients
$\mathrm{ch}(P_\lambda)\,\mathrm{ch}(P_\lambda^\vee)$. The upgrade
recovers the full matrix-valued identity. The CY-to-chiral
correspondence $\Phi_2$
(Chapter~\ref{ch:cy-to-chiral}) transports this Plancherel structure
from $D^b\mathrm{Coh}(\mathrm K3)$ to $\mathbf H_{\Delta_5}$ via the
Mukai-lattice grading, with the flavour character $\zeta^{J_0}$ of
Remark~\ref{rem:modtr-DNA-1} tracking the $M_{24}$-equivariance of
the decomposition. Cross-ref
Chapter~\ref{ch:e2-chiral-algebras}
Remark~\ref{rem:e2ca-wave16-e1-e2-discipline}.
\end{remark}

\section{Explicit banding 2-cocycle for the $\mu_8$-gerbe off the Humbert divisor (Gelfand W18)}
\label{sec:modtr-wave18-banding-cocycle}

Waves 15--17 established, at $(\infty,1)$-level, that the $\mu_8$-gerbe
$\mathcal G$ defined by the Siegel--Borcherds associator
$\widetilde\Phi^{\mathrm{Sieg\text{-}Bor}}_{\mathrm{Sp}_4}[\Phi_{10}/\eta^{24}]$
on $\overline{\mathcal A_2}$ is trivialisable off the Humbert divisor
$H_1\cup H_4$ (Chapter~\ref{ch:k3-chiral-bialgebra-platonic}
Remark~\ref{rem:cydkappa-wave15-gerbe}; Bruinier 2002 LNM~1780
Prop.~5.1).  This section upgrades the $(\infty,1)$-statement to a
\emph{chain-level} trivialising section: an explicit
\v{C}ech 2-cocycle $F_U = [F_{ij}]$ in
$C^2(\mathrm{Sp}_4(\mathbb Z), \mathbb C^\times)$ with image landing in
$\mu_8\subset\mathbb C^\times$, satisfying the cocycle condition
$\delta F_U = 0$ on triple intersections.

\begin{definition}[Igusa fundamental-domain open cover of $U$]
\label{def:modtr-wave18-igusa-cover}
Let $\mathcal F_{\mathrm{Ig}}\subset\mathfrak H_2$ denote the Igusa
fundamental domain for the action of $\mathrm{Sp}_4(\mathbb Z)/\{\pm I\}$
on the Siegel upper half-space, described explicitly by Igusa 1962
\emph{Ann.\ Math.}~76 \S5 in terms of the six inequalities
\[
  |\det(C Z + D)|\geq 1
  \quad\text{for the six principal cosets of}\quad
  \mathrm{Sp}_4(\mathbb Z)/\Gamma_{\infty},
\]
where $\Gamma_{\infty}$ is the Siegel parabolic stabiliser of the cusp
at infinity (Igusa 1962 \S5 Thm.~3). Let
$U = \overline{\mathcal A_2}\setminus(H_1\cup H_4)$, and define an
explicit open cover $\{U_i\}_{i\in I}$ of $U$ by
\[
  U_i \;:=\; \gamma_i\cdot\bigl(\mathcal F_{\mathrm{Ig}}\setminus
  (H_1\cup H_4)\bigr),
\]
indexed by a system $\{\gamma_i\}_{i\in I}$ of coset representatives for
\[
  \mathrm{Sp}_4(\mathbb Z)/\mathrm{Sp}_4(\mathbb Z)_{\mathrm{pro-unip}}
  \cdot\langle w_8\rangle,
\]
where $w_8$ is the Fricke involution of
Theorem~\ref{thm:modtr-wave16-S-fricke-involution}.  The finiteness of
$I$ on each compact subset of $U$ follows from Igusa 1962 \S5 Prop.~2
(six-face polytope structure) combined with
$\mathrm{vol}(\overline{\mathcal A_2}) = \tfrac{1}{8640}\zeta(-1)\zeta(-3)$
being finite (Siegel 1935 \emph{Ann.\ Math.}~36 Thm.~3; Shimura 1963
\emph{Math.\ Ann.}~152 eq.~1.7).
\end{definition}

\begin{remark}[Cover cardinality and stratification]
\label{rem:modtr-wave18-igusa-cover-cardinality}
The Igusa 1962 \S5 Thm.~3 six-face decomposition of
$\mathcal F_{\mathrm{Ig}}$ produces six translates of a single
fundamental domain under translations
$Z\mapsto Z + S$ with $S\in\mathrm{Sym}_2(\mathbb Z)$.  Modulo the
finite group $\mathrm{Sp}_4(\mathbb Z/2)$ of order
$|\mathrm{Sp}_4(\mathbb F_2)| = 720$, the cover on
$\overline{\mathcal A_2}\setminus(H_1\cup H_4)$ consists of at most
$6\cdot 720 = 4320$ charts (tight bound: 720 principal translates and
at most 6 Igusa faces per principal translate), with the actual
cardinality at a fixed weight-truncation level finite and
computable from the weight-$N$ Koecher principle (Klingen 1990
\emph{Introductory Lectures} Thm.~2.2.1).  Primary: Igusa 1962 \S5
Thm.~3; Klingen 1990 Thm.~2.2.1; Freitag 1983
\emph{Siegelsche Modulfunktionen} \S4.5 for the volume formula.
\end{remark}

\begin{theorem}[Explicit banding 2-cocycle $F_{ij}$ on $U$]
\label{thm:modtr-wave18-banding-cocycle}
\ClaimStatusProvedHere
Let $\{U_i\}_{i\in I}$ be the Igusa open cover of $U$ of
Definition~\ref{def:modtr-wave18-igusa-cover}, and let
$s_i\colon U_i\to\mathcal G|_{U_i}$ denote the local section of the
$\mu_8$-gerbe $\mathcal G$ given by the chosen eighth root
\[
  s_i(Z) \;:=\; [\Phi_{10}(Z)/\eta(Z)^{24}]^{1/8}_{\mathrm{princ}}|_{U_i},
\]
where the principal branch of the eighth root is pinned by the
normalisation
$\lim_{Z\to i\infty}s_i(Z) = e^{2\pi i\rho/8}|_{U_i}$ with $\rho$ the
BKM Weyl vector (Gritsenko--Nikulin 1998
\emph{St.~Petersburg~Math.~J.}~11 eq.~2.4).  Then the \v{C}ech
2-cocycle
\[
  \boxed{\;
  F_{ij}(Z) \;:=\;
  \frac{s_i(Z)}{s_j(Z)}
  \;=\;
  \biggl(\frac{\Phi_{10}/\eta^{24}|_{U_i}}
              {\Phi_{10}/\eta^{24}|_{U_j}}\biggr)^{1/8}_{\mathrm{princ}},
  \qquad Z\in U_i\cap U_j,
  \;}
\]
is well-defined, lands in $\mu_8\subset\mathbb C^\times$
($F_{ij}(Z)^8 = 1$ for all $Z$), and satisfies the cocycle condition
$\delta F_U = 0$ on all triple intersections:
\[
  F_{ij}(Z)\cdot F_{jk}(Z)\cdot F_{ki}(Z) \;=\; 1,
  \qquad Z\in U_i\cap U_j\cap U_k.
\]
In particular, $[F_U]\in H^2(\mathrm{Sp}_4(\mathbb Z), \mu_8)$ is
trivial when restricted to $U$, so the $\mu_8$-gerbe $\mathcal G|_U$ is
chain-level trivialisable on $U$.
\end{theorem}

\begin{proof}
\emph{Well-definedness.} On $U_i\cap U_j$ the transition function
$\Phi_{10}/\eta^{24}$ is non-vanishing (divisor $2H_1 + H_4$ removed by
passage to $U$; primary: Igusa 1962 \S7 Cor.~2 for the divisor of
$\Phi_{10}$, Gritsenko--Nikulin 1998 \S3 for
$\mathrm{div}(\eta^{24}) = 0$ on $\mathfrak H_2$ relative to the
paramodular cover), hence its ratio is a non-vanishing holomorphic
function on $U_i\cap U_j$ and admits a well-defined eighth-root branch
fixed by the principal-branch normalisation in the statement.

\emph{$\mu_8$-valuedness.} Fix $Z\in U_i\cap U_j$.  Both
$s_i(Z)$ and $s_j(Z)$ are eighth roots of the same non-zero complex
number $(\Phi_{10}/\eta^{24})(Z)$, so
$F_{ij}(Z)^8 = (s_i(Z)/s_j(Z))^8 = (\Phi_{10}/\eta^{24})(Z)/
(\Phi_{10}/\eta^{24})(Z) = 1$.  Hence $F_{ij}(Z)\in\mu_8$.

\emph{Cocycle condition.} On the triple intersection
$U_i\cap U_j\cap U_k$ the telescoping identity
\[
  F_{ij}\cdot F_{jk}\cdot F_{ki}
  \;=\;
  \frac{s_i}{s_j}\cdot\frac{s_j}{s_k}\cdot\frac{s_k}{s_i}
  \;=\; 1
\]
holds pointwise in $\mathbb C^\times$, hence in $\mu_8$.  This is the
defining 2-cocycle relation in $C^2(\mathrm{Sp}_4(\mathbb Z), \mu_8)$
(Brown 1982 \emph{Cohomology of Groups} \S IV.3 Prop.~3.1).

\emph{Cohomological triviality on $U$.} The cocycle
$F_U = [F_{ij}]$ is a \v{C}ech coboundary on $U$ by construction: the
$0$-cochain $s = (s_i)$ satisfies $\delta s = F_U$
(Hartshorne 1977 III \S4 Prop.~4.2; Bott--Tu 1982 \S8 Thm.~8.9 for the
\v{C}ech--de~Rham double-complex identification).  The class
$[F_U]\in H^2(U, \mu_8)$ is therefore trivial; equivalently, the pullback
$[F_U]\in H^2(\mathrm{Sp}_4(\mathbb Z), \mu_8)$ along the uniformisation
$\mathfrak H_2\to\overline{\mathcal A_2}$ is trivial when restricted
to the open-dense locus $U$.

\emph{Multi-path verification} (AP10):
(i) \v{C}ech-cohomological: Hartshorne 1977 III Prop.~4.2 +
    explicit eighth-root contraction.
(ii) Group-cohomological: Brown 1982 \S IV.3 Prop.~3.1 +
    Lyndon--Hochschild--Serre spectral sequence
    $H^p(\mathrm{Sp}_4(\mathbb Z),H^q(U,\mu_8))\Rightarrow
    H^{p+q}(\mathrm{Sp}_4(\mathbb Z)\ltimes U,\mu_8)$
    (Weibel 1994 \S6.8 Thm.~6.8.2).
(iii) Arithmetic: Bruinier 2002 LNM~1780 Prop.~5.1 gives the
    obstruction class $[\mathcal G]\in H^2(\overline{\mathcal A_2},\mu_8)
    = \mu_8\cdot[H_1] + \mu_{16}\cdot[H_4]$, which restricts to zero on
    $U = \overline{\mathcal A_2}\setminus(H_1\cup H_4)$.
(iv) Modular-form-theoretic: the eighth root of
    $\Phi_{10}/\eta^{24}$ is well-defined on simply-connected domains
    avoiding the zero divisor, and the ratio of eighth roots is the
    Atkin--Lehner-type phase cocycle of Atkin--Lehner 1970
    \emph{Math.\ Ann.}~185 eq.~1.7, extended to paramodular level by
    Gritsenko 1995 \S3; the phase is $\mu_8$-valued by Gritsenko 1995
    Thm.~3.2 (paramodular multiplier has order dividing $8$).
\end{proof}

\begin{remark}[Cocycle class vs.\ banding section]
\label{rem:modtr-wave18-cocycle-vs-banding}
The trivialising section $F_U$ of
Theorem~\ref{thm:modtr-wave18-banding-cocycle} is the \emph{banding}
of $\mathcal G|_U$ in the sense of Breen 1994 \emph{Ast\'erisque}~225
\S2.1: a collection of local sections $\{s_i\}$ whose transition
cocycle $F_{ij} = s_i/s_j$ is both $\mu_8$-valued and a 2-coboundary
on $U$.  The banding is unique up to global $\mu_8$-valued
1-coboundaries (multiplying every $s_i$ by the same globally-chosen
eighth root of unity).  The corresponding gerbe isomorphism
$\mathcal G|_U \xrightarrow{\sim} B\mu_8\times U$ is the one pinned by
$F_U$.  Primary: Breen 1994 \S2.1 for the general gerbe banding
formalism.  Cross-ref Chapter~\ref{ch:quantum-groups}
Section~\ref{sec:qgf-wave18-humbert-obstruction} for the obstruction
to extending the banding across the Humbert divisor.
\end{remark>

\begin{remark}[Chain-level $F_U$ upgrades the $(\infty,1)$-statement]
\label{rem:modtr-wave18-dna-chain-level-upgrade}
The $(\infty,1)$-version of the statement: the neutral Tannakian fibre functor
$\omega\colon\mathrm{Rep}(\mathbf H_{\Delta_5})|_U\to\mathrm{sVec}_{\mathbb C}$
exists after $\mu_8$-twist descent, encoded by vanishing of the
gerbe class $[\mathcal G]\in H^2_{\mathrm{\acute{e}t}}(U,\mu_8)$
on $U$. The chain-level
\v{C}ech 2-cocycle $F_{ij}$ of
Theorem~\ref{thm:modtr-wave18-banding-cocycle} realises this
vanishing explicitly.  The upgrade matters for two reasons.
(i) \emph{Computability}: the explicit eighth-root ratio is a
holomorphic function on each double intersection, evaluable at any
$Z\in U_i\cap U_j$ and directly compared with the Atkin--Lehner phase
cocycle of Gritsenko 1995 \S3 (Path (iv) of the proof).  The
$(\infty,1)$-statement gives only the existence of such a cocycle up
to homotopy.
(ii) \emph{Chiral-Hochschild transport}: the chain-level banding is
the obstruction-class input for the
$\Phi$-pushforward of the \emph{chiral} Hochschild
$\mathrm{ChirHoch}^\bullet(\mathbf H_{\Delta_5})$ Humbert-restriction
spectral sequence of Chapter~\ref{ch:cy-to-chiral} (chiral, not
categorical or topological; cross-ref AP160 discipline), which
requires an honest cochain representative to converge on the
$E_2$-page
(Schiffmann--Vasserot 2013 \emph{Duke}~162 Thm.~4.5.2 for the
CoHA-categorical-Hochschild pro-spectral sequence at chain level,
transported to the chiral side along $\Phi$).  Cross-ref
Remark~\ref{rem:modtr-wave17-scalar-shadow} (scalar trace) and
Theorem~\ref{thm:modtr-wave18-Plancherel-ML-limit} (Mittag--Leffler
Plancherel). The banding closes the final $(\infty,1)$-to-chain
gap in the pro-finite Plancherel programme.
\end{remark}

\begin{remark}[Banding fails at $H_1\cup H_4$: cross-ref summary]
\label{rem:modtr-wave18-banding-humbert-failure}
The Humbert divisor $H_1\cup H_4$ is the obstruction locus where the
banding $F_U$ cannot extend: the $\mu_8$-gerbe has monodromy of order
$8$ around $H_1$ (Bruinier 2002 LNM~1780 Prop.~5.1) and monodromy of
order $16$ around $H_4$ (via Kuga--Satake; Morrison 1984
\emph{Invent.}~75 Thm.~3).  A local
extension
$\widetilde F_U\colon\widetilde U\to\mu_8$ across either component would
force a branch of the eighth root of $\Phi_{10}/\eta^{24}$ to be
single-valued on a punctured neighbourhood of
$H_1$ (resp.\ $H_4$); the order-$8$ (resp.\ order-$16$) monodromy of
$\Phi_{10}^{1/8}$ around each component forbids this
(Chapter~\ref{ch:quantum-groups}
Theorem~\ref{thm:qgf-wave18-humbert-non-extendability}).
The obstruction class on the Humbert divisor is the generator of
$H^2(\overline{\mathcal A_2}, \mu_8)|_{H_1\cup H_4} =
\mu_8\cdot[H_1] + \mu_{16}\cdot[H_4]$, matching the Humbert-block
Fricke-trace $\mathrm{tr}(S) = 4$ on the $M_{24}$-invariant block
(Theorem~\ref{thm:modtr-wave17-Plancherel-equals-Phi10}
Remark~\ref{rem:modtr-wave16-S-eigenvalues}).
\end{remark>


% ==========================================
% PART IV: Quantum Groups and Braided Monoidal Structure
% ==========================================

\part{Quantum Groups and Braided Monoidal Structure}
\label{part:quantum-groups}

\chapter{Quantum Groups: Foundations}
\label{ch:quantum-groups}

This chapter collects the quantum group background required for Part~III and the CY landscape.  The quantized enveloping algebra $U_q(\frakg)$ and its $R$-matrix are the targets of the CY-to-chiral functor at the representation-theoretic level: Conjecture~CY-C predicts that the braided monoidal category $\Rep_q(\frakg)$ arises as $\Rep^{E_2}(A_\cC)$ for a CY category $\cC(\frakg, q)$.  The Drinfeld--Jimbo and FRT presentations provide distinct access points, corresponding respectively to the Koszul dual (bar cohomology generators and relations) and the ordered bar (RTT $R$-matrix formalism) of the chiral algebra.

\section{Quantized enveloping algebras}
\label{sec:quantized-enveloping}

This section reviews the Drinfeld--Jimbo quantized enveloping algebra $U_q(\frakg)$ and the Lusztig divided-power form $\dot{U}_q(\frakg)$, emphasising the features relevant to the CY programme: the triangular decomposition (which matches the bar filtration), the Hopf algebra structure (which matches the bar coproduct), and the specialisation at roots of unity (which connects to the Kazhdan--Lusztig equivalence of Chapter~\ref{ch:quantum-group-reps}).

\section{The $R$-matrix and Yang--Baxter equation}
\label{sec:r-matrix-ybe}

The universal $R$-matrix $\cR \in U_q(\frakg)^{\hat{\otimes}2}$ is the quantum group avatar of the collision $r$-matrix $r(z) = \Res^{\mathrm{coll}}_{0,2}(\Theta_A)$ from Volume~I.  This section develops the Yang--Baxter equation, quasi-triangular Hopf algebras, and the passage from the formal $R$-matrix to the trigonometric and elliptic solutions that appear in the CY landscape.  The key structural point: the $R$-matrix lives in arity~$(1,1)$ of the ordered bar $B^{\mathrm{ord}}(A)$, and its Baxterisation $R(z)$ is the $z$-dependent lift (Volume~II, DK bridge).

\section{Representation categories}
\label{sec:qg-rep-categories}

This section establishes that $\Rep_q(\frakg)$ is a braided monoidal category (ribbon category when $q$ is generic), with the braiding provided by the universal $R$-matrix.  At roots of unity, the category acquires a finite tensor subcategory (the Andersen--Jantzen--Soergel quotient) that is modular in the sense of Turaev.  The target: the CY category $\cC(\frakg, q)$ should produce this braided category as $\Rep^{E_2}(\Phi(\cC))$.

\section{Drinfeld--Jimbo vs Faddeev--Reshetikhin--Takhtajan}
\label{sec:dj-vs-frt}

The Drinfeld--Jimbo presentation uses generators $E_i, F_i, K_i$ and Serre relations; the FRT presentation uses the $R$-matrix equation $R T_1 T_2 = T_2 T_1 R$ directly.  From the bar-complex perspective (Volume~II, Chapter on Yangians): the DJ presentation corresponds to the Koszul dual algebra $A^!$ (generators = bar $H^1$, relations = bar $H^2$), while the FRT/RTT presentation corresponds to the ordered bar with its $R$-matrix in bidegree $(1,1)$.  The equivalence of the two presentations is a theorem of Faddeev--Reshetikhin--Takhtajan for finite type; the extension to affine and toroidal types is developed in the toric CoHA chapter (Chapter~\ref{ch:toric-coha}).

\chapter{Braided Factorization}
\label{ch:braided-factorization}

%%: The E_1/ordered story is primitive. The E_2 braided
%% structure is the Drinfeld double / center construction applied to E_1.
%%: kappa always subscripted in Vol III.
%% Prerequisite: Chapter~\ref{ch:quantum-groups} (Quantum Groups: Foundations)
%% owns the Drinfeld--Jimbo presentation of $\Uq(\frakg)$, the universal
%% $R$-matrix, the quantum Yang--Baxter equation, and the Kazhdan--Lusztig
%% equivalence. The present chapter builds the $E_2$-chiral / factorization
%% layer on top.

A braided monoidal category is a monoidal category equipped with a
coherent system of isomorphisms $V \otimes W \xrightarrow{\sim} W \otimes V$
satisfying the hexagon axioms. The bar-cobar adjunction of Volume~I,
extended to the $E_2$ setting, produces factorization coalgebras whose
arity-$(1,1)$ component is a categorical $R$-matrix. Chapter~\ref{ch:quantum-groups}
collects the classical Drinfeld--Jimbo presentation of $\Uq(\frakg)$, the
universal $R$-matrix, and the quantum Yang--Baxter equation; here we assume
that material and construct its $E_2$-chiral / factorization avatar. The
fundamental problem solved in this chapter: how does the $E_2$-chiral bar
complex $B_{E_2}(\cA)$ encode a categorical braiding, and what is the precise
relationship between the braided Koszul dual and the quantum group targets
of Conjecture~\ref{conj:qgf-cy-c}?

\section{The categorical $R$-matrix and Yang--Baxter equation}
\label{sec:bf-categorical-r-matrix}

The Drinfeld--Jimbo universal $R$-matrix of
Proposition~\ref{prop:qgf-classical-limit-r} lives in $\Uq(\frakg)^{\hat\otimes 2}$
and is an \emph{algebraic} datum. The Vol~III programme produces $R$-matrices
of a different kind: \emph{categorical} $R$-matrices, extracted from the
arity-$(1,1)$ component of an $E_2$-chiral bar complex. The two pictures
are compatible through Conjecture~CY-C (Conjecture~\ref{conj:qgf-cy-c}):
when the CY-to-chiral functor produces an affine Kac--Moody chiral algebra,
the categorical $R$-matrix coincides with the Drinfeld--Jimbo $R$-matrix
under the Kazhdan--Lusztig equivalence (Theorem~\ref{thm:qgf-kazhdan-lusztig}).

\begin{remark}[Drinfeld--Jimbo vs.\ FRT, from the bar-complex standpoint]
\label{rem:bf-dj-vs-frt}
The Drinfeld--Jimbo presentation constructs $\Uq(\frakg)$ from generators
and relations; the Faddeev--Reshetikhin--Takhtajan (FRT) construction
recovers $\Uq(\frakg)$ from the $R$-matrix via the RTT relation
$R_{12} T_1 T_2 = T_2 T_1 R_{12}$. The FRT construction is the
bar-complex-natural one: the $R$-matrix lives in arity $(1,1)$ of
$B^{\mathrm{ord}}(A)$, and the RTT relation is the
arity-$(1,1,1)$ bar differential. This is the standpoint
of the present chapter: the categorical $R$-matrix below is the
arity-$(1,1)$ component of $B_{E_2}(\Phi(\cC))$, and the parametric
Yang--Baxter equation is coassociativity of the deconcatenation coproduct.
\end{remark}

\begin{definition}[Categorical $R$-matrix]
\label{def:categorical-r-matrix}
For a CY category $\cC$ of dimension $d = 2$, the \emph{categorical
$R$-matrix} $\cR_\cC(z)$ is the arity-$(1,1)$ component of the
$E_2$-bar complex $B_{E_2}(\Phi(\cC))$:
\[
 \cR_\cC(z) \in \End(V \otimes V)[[z, z^{-1}]]
\]
where $V = s^{-1}\bar{A}_\cC$ is the desuspended generating space and
$z$ is the spectral parameter. The assignment $\cC \mapsto \cR_\cC(z)$
is functorial in CY categories (for exact CY functors preserving the
$E_2$-structure).
\end{definition}

\begin{proposition}[Yang--Baxter from bar associativity]
\label{prop:ybe-from-bar}
\ClaimStatusProvedElsewhere
The categorical $R$-matrix $\cR_\cC(z)$ satisfies the Yang--Baxter
equation
\[
 \cR_{12}(z_1 - z_2) \, \cR_{13}(z_1 - z_3) \, \cR_{23}(z_2 - z_3)
 = \cR_{23}(z_2 - z_3) \, \cR_{13}(z_1 - z_3) \, \cR_{12}(z_1 - z_2).
\]
This is a formal consequence of the coassociativity of the bar coproduct
$\Delta_{\mathrm{decat}}$ applied to arity $3$ of $B^{\mathrm{ord}}(A)$.
\end{proposition}

\begin{proof}
The bar complex $B^{\mathrm{ord}}(A)$ is a coassociative coalgebra with
deconcatenation coproduct. The two compositions
$(\Delta \otimes \id) \circ \Delta$ and $(\id \otimes \Delta) \circ \Delta$
coincide on arity-$3$ elements; the $R$-matrix arises from the bar
differential restricted to arity $(1,1)$. Coassociativity at arity $3$
gives $\cR_{12}\cR_{13}\cR_{23} = \cR_{23}\cR_{13}\cR_{12}$, which is
the Yang--Baxter equation. The spectral parameters arise from the
$z$-dependence of the collision residue (so the $R$-matrix has one
fewer pole order than the OPE). See Volume~II, Chapter~11 for the
full derivation.
\end{proof}


\section{Braided monoidal categories and $E_2$-algebras}
\label{sec:braided-e2}

The connection between $E_2$-algebras and braided monoidal categories is
given by the Joyal--Street coherence theorem and its
$\infty$-categorical refinement.

\begin{proposition}[Dunn additivity for $E_2$]
\label{prop:dunn-e2}
\ClaimStatusProvedElsewhere
The $E_2$ operad satisfies $\Etwo \simeq \Eone \otimes_{E_0} \Eone$.
An $\Etwo$-algebra is an $\Eone$-algebra in $\Eone$-algebras. In
particular, an $\Etwo$-algebra has two compatible associative products
whose interchange law is the braiding.
\end{proposition}

\begin{corollary}
\label{cor:e2-braided-monoidal}
The representation category $\Rep^{E_2}(\cA)$ of an $E_2$-algebra
$\cA$ is braided monoidal. The braiding on $V \otimes W$ is given by
the action of the $R$-matrix: $\sigma_{V,W} = \cR_{V,W} \circ \tau$
where $\tau \colon V \otimes W \to W \otimes V$ is the transposition
and $\cR_{V,W}$ is the $R$-matrix evaluated on the pair $(V,W)$.
\end{corollary}


\section{Factorization homology with $E_2$-coefficients}
\label{sec:fh-e2}

Factorization homology $\int_M \cA$ of an $\En$-algebra $\cA$ over an
$n$-manifold $M$ (Ayala--Francis) generalizes Hochschild homology
($n = 1$, $M = S^1$) and chiral homology ($n = 2$, $M$ a surface).
The relevance to Vol~III is that CY$_2$ categories produce $\Etwo$-chiral
algebras (by CY-A at $d = 2$), and factorization homology of these
$\Etwo$-algebras over surfaces computes genus-$g$ conformal block spaces;
the shadow obstruction tower of Vol~I controls the moduli dependence.

\begin{definition}[Factorization homology]
\label{def:factorization-homology}
Let $\cA$ be an $\En$-algebra and $M$ an $n$-manifold. The
\emph{factorization homology} $\int_M \cA$ is the colimit
\[
 \int_M \cA \;=\; \colim_{\mathrm{Disk}_{n/M}} \cA
\]
over the $\infty$-category $\mathrm{Disk}_{n/M}$ of embeddings of
disjoint unions of disks into $M$. The colimit is computed in the
$\infty$-category of chain complexes (Ayala--Francis, 2015).
\end{definition}

\begin{proposition}[Excision for factorization homology]
\label{prop:fh-excision}
\ClaimStatusProvedElsewhere
Let $M = M_1 \cup_{N \times \R} M_2$ be a collar-gluing decomposition
of an $n$-manifold. For an $\En$-algebra $\cA$, there is a canonical
equivalence
\[
 \int_M \cA
 \;\simeq\;
 \int_{M_1} \cA \;\otimes_{\int_{N \times \R} \cA}\; \int_{M_2} \cA,
\]
where the tensor product is taken over the $\Eone$-algebra
$\int_{N \times \R} \cA$ (Ayala--Francis excision).
\end{proposition}

Excision is the engine that converts local data (disk-level OPE) into
global invariants (genus-$g$ amplitudes). Applied to the pair-of-pants
decomposition $\Sigma_g = \#_{2g-2}(\text{pair of pants})$, excision
expresses $\int_{\Sigma_g} \cA$ as an iterated relative tensor product
over Hochschild homology, recovering the sewing prescription of conformal
field theory.

\begin{proposition}[$E_2$ factorization homology on surfaces]
\label{prop:fh-e2-surface}
\ClaimStatusProvedElsewhere
For an $\Etwo$-algebra $\cA$ and a closed oriented surface
$\Sigma_g$ of genus $g$:
\begin{enumerate}[label=(\roman*)]
 \item $\int_{\Sigma_g} \cA$ is the genus-$g$ conformal block space;
 \item $\int_{S^1 \times [0,1]} \cA \simeq \HH_\bullet(\cA)$, the
 Hochschild homology;
 \item $\int_{T^2} \cA \simeq \HH_\bullet(\HH_\bullet(\cA))$
 (double Hochschild homology, computing the genus-$1$ amplitude).
\end{enumerate}
The shadow obstruction tower controls the dependence of
$\int_{\Sigma_g} \cA$ on the moduli of $\Sigma_g$.
\end{proposition}

\begin{proposition}[Genus tower from FH]
\label{prop:fh-genus-tower}
\ClaimStatusConditional
For a CY$_2$ category $\cC$ with $\Etwo$-chiral algebra
$\cA = \Phi(\cC)$, the factorization homology genus tower satisfies
\[
 \chi\!\left(\int_{\Sigma_g} \cA\right)
 \;=\;
 \kappa_{\mathrm{cat}}(\cC) \cdot \lambda_g
 \qquad\textup{(UNIFORM-WEIGHT)}
\]
at genus $g \geq 1$, where $\kappa_{\mathrm{cat}} = \chi^{\CY}(\cC)$
is the CY Euler characteristic
(Theorem~\textup{\ref{thm:cy-modular-characteristic}}). For
multi-weight algebras at $g \geq 2$, the scalar formula receives
the cross-channel correction $\delta F_g^{\mathrm{cross}}$ of Vol~I,
Theorem~D. The conditional dependence is on Theorem~CY-B(iii)
(curvature identification).
\end{proposition}

\begin{remark}[FH versus chiral homology]
\label{rem:fh-vs-chiral}
The Ayala--Francis factorization homology $\int_{\Sigma_g} \cA$ and the
Beilinson--Drinfeld chiral homology $H^{\mathrm{ch}}_\bullet(\Sigma_g, \cA)$
are equivalent for $\Etwo$-algebras that arise as factorization envelopes
of chiral Lie algebras (by the comparison theorem of Francis, 2013).
In the CY setting, $\cA = \Phi(\cC)$ is always of this type when CY-A
holds (at $d = 2$), so the two invariants agree. The FH formulation has
the advantage of applying to topological $\Etwo$-algebras without
holomorphic structure.
\end{remark}


\section{The $E_2$-equivariant bar complex}
\label{sec:e2-equivariant-bar}

%% FOUR OBJECTS:
%% 1. B(A) = bar coalgebra = T^c(s^{-1} A-bar) with deconcatenation + twist
%% 2. A^i = H^*(B(A)) = dual coalgebra (Koszul cohomology of bar)
%% 3. A^! = ((A^i)^v) = dual ALGEBRA (linear or Verdier dual)
%% 4. Z^der_ch(A) = derived chiral center = Hochschild cochains = bulk

The $E_1$-bar complex $B^{\mathrm{ord}}(\cA) = T^c(s^{-1}\bar{\cA})$
of Volume~I is an ordered coalgebra with deconcatenation coproduct and
a single differential from OPE residues along one real direction.
In the $\Etwo$ setting, the two $\Eone$-directions of the Dunn
decomposition $\Etwo \simeq \Eone \otimes_{E_0} \Eone$ each contribute
a bar differential. The $\Etwo$-equivariant bar complex intertwines them
through the $R$-matrix.

\begin{definition}[$E_2$-equivariant bar complex]
\label{def:e2-equivariant-bar}
For an $\Etwo$-chiral algebra $\cA$ with augmentation ideal
$\bar{\cA} = \ker(\varepsilon)$, the \emph{$E_2$-equivariant bar
complex} is the graded coalgebra
\[
 B_{E_2}(\cA)
 \;=\;
 T^c(s^{-1}\bar{\cA})
\]
equipped with:
\begin{enumerate}[label=(\roman*)]
 \item two commuting differentials $d_X, d_Y$ from OPE residues
 in each $\Eone$-direction, satisfying $d_X^2 = d_Y^2 = 0$ and
 $d_X \circ d_Y = d_Y \circ d_X$;
 \item two deconcatenation coproducts $\Delta_X, \Delta_Y$ giving
 $B_{E_2}(\cA)$ the structure of an $\Etwo$-coalgebra;
 \item a braid group action: the symmetric group $S_n$ action on
 $B^{\mathrm{ord}}_n(\cA) = (s^{-1}\bar{\cA})^{\otimes n}$ lifts
 to a braid group $B_n$ action via the $R$-matrix
 $\cR_\cA(z)$ (Definition~\ref{def:categorical-r-matrix}).
\end{enumerate}
The total differential is $d = d_X + d_Y$; it satisfies $d^2 = 0$ on
the nose (genus-$0$ bar complex). At genus $g \geq 1$, the fiberwise
differential satisfies $d_{\mathrm{fib}}^2 = \kappa_{\mathrm{cat}} \cdot \omega_g$
(Hodge curvature; see Theorem~\ref{thm:e2-bar-cobar}(iii)).
\end{definition}

\begin{remark}[$B_{E_2}(\cA)$ is not a Swiss-cheese coalgebra]
\label{rem:e2-bar-not-sc}
The $\Etwo$-bar complex $B_{E_2}(\cA)$ is an $\Etwo$-coalgebra, \emph{not}
a Swiss-cheese ($\SC^{\mathrm{ch,top}}$) coalgebra. The SC structure of
Vol~II is two-coloured (holomorphic + topological) and emerges on the
\emph{derived center} $\Zder(\cA) = C^\bullet_{\mathrm{ch}}(\cA, \cA)$
(the bulk algebra), not on the bar complex of $\cA$ itself. The bar
complex encodes the boundary; the SC structure governs the bulk-boundary
coupling. Concretely:
\begin{itemize}
 \item $B_{E_2}(\cA)$: $\Etwo$-coalgebra, boundary data,
 deconcatenation coproduct, braid group action;
 \item $\Zder(\cA)$: $\Etwo$-algebra (by the BZF theorem,
 Theorem~\ref{thm:bzfn}), bulk data, carries the SC structure
 when coupled with the boundary $\cA$.
\end{itemize}
The identification $\cZ(\Rep^{E_1}(\cA)) \simeq \Rep^{E_2}(\Zder(\cA))$
(Chapter~\ref{ch:drinfeld-center}) relates the two: the braided monoidal
structure on representations of the bulk is the Drinfeld center of the
monoidal structure on representations of the boundary.
\end{remark}

\begin{proposition}[Three bar complexes and their relation]
\label{prop:three-bars-relation}
\ClaimStatusProvedElsewhere
For an $\Etwo$-chiral algebra $\cA$, the three bar complexes of
Construction~\ref{constr:three-bar-complexes} are related by:
\[
 B_{E_\infty}^n(\cA) = B_{E_1}^n(\cA) / S_n, \qquad
 H^*(B_{E_\infty}^n(\cA)) = H^*(B_{E_2}^n(\cA))^{S_n}.
\]
The $\Eone$-bar is the primitive ordered object; the $\Etwo$-bar
refines the $S_n$-action to the braid group $B_n$ via the $R$-matrix;
the $\Einf$-bar is the $S_n$-coinvariant shadow. The averaging map
$\mathrm{av} \colon B_{E_1} \to B_{E_\infty}$ factors through $B_{E_2}$:
\[
 B_{E_1} \xrightarrow{B_n\text{-equivariantization}} B_{E_2}
 \xrightarrow{B_n \twoheadrightarrow S_n} B_{E_\infty}.
\]
\end{proposition}


\section{The braided bar-cobar adjunction}
\label{sec:braided-bar-cobar}

The Vol~I bar-cobar adjunction (Theorems~A and~B) operates at the $\Eone$
level: $B^{\mathrm{ord}}$ is an $\Eone$-coalgebra and $\Omega$ recovers
the $\Eone$-algebra. The $\Etwo$-refinement lifts the adjunction to
braided coalgebras, with the $R$-matrix providing the additional datum
that distinguishes $\Etwo$ from $\Eone$.

\begin{theorem}[$E_2$-bar-cobar adjunction: Theorem CY-B]
\label{thm:e2-bar-cobar}
\ClaimStatusConditional
For an $\Etwo$-chiral algebra $\cA$, the bar-cobar adjunction
\[
 B_{E_2} \colon E_2\text{-}\mathrm{ChirAlg}
 \rightleftarrows E_2\text{-}\mathrm{ChirCoalg}
 \cocolon \Omega_{E_2}
\]
satisfies:
\begin{enumerate}[label=(\roman*)]
 \item The $\Etwo$ bar complex $B_{E_2}(\cA)$ is a factorization
 $\Etwo$-coalgebra on $\Ran(X) \times \Ran(Y)$;
 \item Cobar-bar inversion:
 $\Omega_{E_2}(B_{E_2}(\cA)) \simeq \cA$ on the Koszul locus;
 \item The CY trace manifests as curvature: at genus $g \geq 1$,
 $d_{\mathrm{fib}}^2 = \kappa_{\mathrm{cat}} \cdot \omega_g$,
 where $\omega_g = c_1(\lambda)$ is the Hodge class on
 $\overline{\cM}_g$ and $\kappa_{\mathrm{cat}} = \chi^{\CY}(\cC)$.
\end{enumerate}
\end{theorem}

\begin{remark}[Status of Theorem CY-B]
\label{rem:cy-b-status}
Items~(i) and~(ii) follow from the Volume~I bar-cobar machine
(Theorems~A and~B) applied to each $\Eone$-direction separately, with
compatibility ensured by Dunn additivity
(Proposition~\ref{prop:dunn-e2}). The argument proceeds in three steps.

\emph{Step 1: $\Eone \times \Eone$ bar-cobar.} Apply the Vol~I
bar-cobar adjunction to each $\Eone$-factor independently. This produces
two commuting bar-cobar adjunctions $(B_X, \Omega_X)$ and
$(B_Y, \Omega_Y)$.

\emph{Step 2: $R$-matrix intertwining.} The $R$-matrix
$\cR_\cA(z) \in \End(V \otimes V)[[z, z^{-1}]]$ at arity $(1,1)$
intertwines the two bar differentials: $\cR \circ d_X = d_Y \circ \cR$
on the arity-$(1,1)$ component. This is the braided compatibility
condition.

\emph{Step 3: Higher coherences.} The full $\Etwo$-equivariant
refinement requires coherent intertwining at all arities, governed by
the braid group $B_n$ action. Steps~1 and~2 are proved; the detailed
verification of higher coherences (arity $\geq 4$) is established as a
rigorous proof sketch. The result is unconditional at arity $\leq 3$
and conditional on the higher-arity coherences beyond that.

Item~(iii) uses Theorem~D of Volume~I with the CY Euler
characteristic $\kappa_{\mathrm{cat}} = \chi^{\CY}(\cC)$
(Theorem~\ref{thm:cy-modular-characteristic}). The curvature
$d_{\mathrm{fib}}^2 = \kappa_{\mathrm{cat}} \cdot \omega_g$ is the
\emph{fiberwise} statement on $\overline{\cM}_g$; the bar differential
$d^2_{\mathrm{bar}} = 0$ always holds (the curvature is Hodge, not bar).
\end{remark}

\begin{proposition}[Arity-$(1,1)$ of $B_{E_2}$]
\label{prop:e2-bar-arity-11}
\ClaimStatusProvedHere
The arity-$(1,1)$ component of $B_{E_2}(\cA)$ is the space
$\End(V \otimes V)$ equipped with the $R$-matrix
$\cR_\cA(z)$. The bar differential restricted to this component
recovers the unitarity condition $\cR_{12}(z) \cR_{21}(-z) = \id$
(strong unitarity), and the braiding constraint
$\cR_{12}(z) = \tau \circ \sigma_{12}(z)$.
\end{proposition}

\begin{proof}
The $E_2$-bar complex at arity $(1,1)$ consists of elements
$\alpha \in (s^{-1}\bar{\cA})^{\otimes 2}$ with one element
from each $E_1$-direction. The bar differential
$d(\alpha) = \mu(\alpha)$ is the OPE residue extracted via
$d\log(z_1 - z_2)$ (Volume~I, bar differential). The
compatibility of the two $E_1$-structures forces
$d_{X} \circ d_{Y} = d_{Y} \circ d_{X}$; evaluating on
$(s^{-1}v) \otimes (s^{-1}w)$ gives
$\cR_{12}(z)\cR_{21}(-z) = \id$ (strong unitarity).
The braiding $\sigma_{12}(z) = \cR_{12}(z) \circ \tau^{-1}$
follows from the identification of $\cR$ with the
half-braiding (Definition~\ref{def:categorical-r-matrix}).
\end{proof}


\section{Quantum group structure from braided Koszul duality}
\label{sec:qg-from-braided-koszul}

The Vol~I Koszul dual $A^! = (H^*(B(A))^{\vee})$ is a linear-dual
algebra. In the $\Etwo$ setting, the bar cohomology
$H^*(B_{E_2}(\cA))$ carries a quasi-triangular Hopf algebra structure
from the $R$-matrix, and the question becomes: does this Hopf algebra
coincide with the quantum group $\Uq(\frakg)$ when the input category
has $\frakg$-symmetry?

\begin{conjecture}[Braided bar cohomology recovers $\Uq(\frakg)$]
\label{conj:braided-bar-uq}
\ClaimStatusConjectured
For the $\Etwo$-chiral algebra $\cA = \Phi(\cC)$ associated to a
CY$_2$ category $\cC$ with $\frakg$-symmetry, the braided bar
cohomology
\[
 H^*(B_{E_2}(\cA)) \;\simeq\; \Uq(\frakg)
\]
as quasi-triangular Hopf algebras, where $q = e^{2\pi i / (k + h^\vee)}$
and $k$ is the level determined by the CY Euler characteristic
$\kappa_{\mathrm{cat}}(\cC)$.
\end{conjecture}

\begin{remark}[Evidence and scope]
\label{rem:braided-bar-evidence}
Conjecture~\ref{conj:braided-bar-uq} is verified in two cases:
\begin{enumerate}[label=(\roman*)]
 \item For $\cC = D^b(\mathrm{Coh}(\mathrm{pt}/G))$ with $G$ simple
 and simply connected, the CY-to-chiral functor produces
 $V_k(\frakg)$ (by CY-A at $d = 2$). The $\Etwo$-bar cohomology
 should recover $\Uq(\frakg)$ via the Kazhdan--Lusztig equivalence
 (Theorem~\ref{thm:qgf-kazhdan-lusztig}). This is verified at the
 level of categories: $\Rep_q(\frakg) \simeq \cO_{k}(\hat{\frakg})$
 at roots of unity.
 \item For $\cC = D^b(\mathrm{Coh}(K3))$, the chiral algebra is a lattice
 VOA and the $\Etwo$-bar cohomology recovers a lattice quantum group
 (Example~\ref{ex:braided-k3}).
\end{enumerate}
The conjecture is open for non-geometric CY categories and for $d = 3$
(where CY-A$_3$ is itself conjectural).
\end{remark}

\begin{remark}[Three dualities]
\label{rem:braided-koszul-langlands}
%: Three distinct operations must be kept separate.
The braided Koszul dual $\cA^!_{E_2}$ produces the quantum group
$\Uq(\frakg)$. This must be distinguished from:
\begin{enumerate}[label=(\alph*)]
 \item the Feigin--Frenkel involution $k \mapsto -k - 2h^\vee$, which
 acts within the same family of affine algebras;
 \item the Langlands dual $\frakg \mapsto \frakg^L$ (root-coroot
 exchange), which produces a genuinely different algebra.
\end{enumerate}
For $\fsl_2$ (self-Langlands-dual), all three operations coincide on
representation categories. For non-simply-laced types (e.g.,
$G_2, F_4$), they diverge. The braided Koszul dual exchanges the
braiding parameter $q \leftrightarrow q^{-1}$; the Langlands dual
exchanges the Cartan matrix $A \leftrightarrow A^T$; the Feigin--Frenkel
involution exchanges levels $k \leftrightarrow -k - 2h^\vee$.
\end{remark}


\section{The braided shadow obstruction tower}
\label{sec:braided-shadow-tower}

The Vol~I shadow obstruction tower $\Theta_A$ is the $\Sigma_n$-coinvariant
projection of a richer $E_1$ object. In the $\Etwo$ setting, the braided
shadow tower $\Theta^{E_2}_\cA$ sits between the ordered $E_1$ tower and
the symmetric $E_\infty$ shadow, with the braid group $B_n$ replacing
the symmetric group $S_n$.

\begin{definition}[Braided modular characteristic]
\label{def:braided-kappa}
For an $\Etwo$-chiral algebra $\cA$, the \emph{braided modular
characteristic} is the pair
\[
 (\kappa_{\mathrm{cat}}(\cA), \, \cR_\cA(z))
\]
consisting of the scalar modular characteristic
$\kappa_{\mathrm{cat}} = \mathrm{av}(\cR_\cA(z))$ (the
$\Sigma_2$-coinvariant of the $R$-matrix) and the full
$R$-matrix itself. The scalar $\kappa_{\mathrm{cat}}$ controls
the genus tower
$F_g = \kappa_{\mathrm{cat}} \cdot \lambda_g^{FP}$ (UNIFORM-WEIGHT);
the $R$-matrix controls the ordered/quantum-group sector.
\end{definition}

\begin{proposition}[Braided shadow structure]
\label{prop:braided-shadow}
\ClaimStatusProvedHere
The $\Etwo$-bar complex $B_{E_2}(\cA)$ carries a braided shadow
obstruction tower $\Theta^{E_2}_\cA$ whose arity-by-arity
projection under the averaging map recovers the Volume~I
shadow obstruction tower:
\[
 \mathrm{av}(\Theta^{E_2}_{\cA, n})
 = \Theta_{A_\cC, n}
 \quad \text{for all $n \geq 2$.}
\]
At arity $2$: $\mathrm{av}(\cR(z)) = \kappa_{\mathrm{cat}}$.
At arity $3$: $\mathrm{av}(\Phi_{KZ}) = C$ (the cubic shadow of
Volume~I).
\end{proposition}

\begin{remark}[The three shadow tiers]
\label{rem:three-shadow-tiers}
The three bar complexes produce three tiers of the shadow tower:
\begin{center}
\renewcommand{\arraystretch}{1.2}
\begin{tabular}{lll}
 \toprule
 Bar complex & Group action & Shadow data \\
 \midrule
 $B_{E_1}$ (ordered) & trivial ($1$) & full $R$-matrix $\cR(z)$, Yangian \\
 $B_{E_2}$ (braided) & braid group $B_n$ & $B_n$-equivariant tower, quantum group \\
 $B_{E_\infty}$ (symmetric) & symmetric $S_n$ & scalar $\kappa_{\mathrm{cat}}$, genus tower \\
 \bottomrule
\end{tabular}
\end{center}
Each tier retains strictly less information than the one above.
The averaging map $\mathrm{av}$ is the $B_n \twoheadrightarrow S_n$
projection applied to each arity. It is lossy: the $R$-matrix is
forgotten, and only its $S_2$-coinvariant $\kappa_{\mathrm{cat}}$
survives.
\end{remark}

\begin{example}[Braided structure for $K3$]
\label{ex:braided-k3}
For a K3 surface $S$ with $\cC = D^b(\Coh(S))$, the
$\Etwo$-chiral algebra has $\kappa_{\mathrm{cat}} = 2$. The
$R$-matrix is the Casimir
$\cR(z) = 1 + \kappa_{\mathrm{cat}}\,\Omega/z + O(z^{-2})$,
where $\Omega$ is the Mukai pairing on $H^*(S, \Z)$. The braided
bar cohomology $H^*(B_{E_2}(\Phi(\cC)))$ recovers the lattice
quantum group associated to the Mukai lattice
$\Lambda_{K3} = U^3 \oplus E_8(-1)^2$.
The shadow class is $\mathbf{G}$ (Gaussian, $r_{\max} = 2$):
the lattice VOA is a free-field algebra and the shadow tower
terminates at arity $2$ (both at the $E_1$ and $E_2$ levels).
\end{example}


\section{Braided complementarity}
\label{sec:braided-complementarity}

Vol~I Theorem~C establishes complementarity:
$\kappa_{\mathrm{ch}}(A) + \kappa_{\mathrm{ch}}(A^!) = K$
(the Koszul conductor, family-dependent). The $\Etwo$-refinement
upgrades this scalar relation to an $R$-matrix-level statement.

\begin{conjecture}[Braided complementarity]
\label{conj:braided-complementarity}
\ClaimStatusConjectured
For an $\Etwo$-chiral algebra $\cA$ on the Koszul locus, the braided
Koszul dual $\cA^!_{E_2}$ satisfies:
\begin{enumerate}[label=(\roman*)]
 \item \emph{Scalar complementarity}:
 $\kappa_{\mathrm{cat}}(\cA) + \kappa_{\mathrm{cat}}(\cA^!_{E_2}) = K$,
 the family-dependent Koszul conductor (Vol~I, Theorem~C);
 \item \emph{$R$-matrix inversion}:
 $\cR_{\cA^!_{E_2}}(z) = \cR_\cA(z)^{-1}$ as formal power series
 in $z^{-1}$, so the braided Koszul dual reverses the braiding;
 \item \emph{Conductor from $R$-matrix}:
 the Koszul conductor $K$ is recovered from the $R$-matrix
 as $K = \mathrm{av}(\cR_\cA(z)) + \mathrm{av}(\cR_\cA(z)^{-1})$.
\end{enumerate}
\end{conjecture}

\begin{remark}[Status of braided complementarity]
\label{rem:braided-complementarity-status}
Item~(i) is an immediate consequence of Vol~I Theorem~C applied
to the underlying $\Einf$-data. Item~(ii) is the genuine $\Etwo$
content: it asserts that Koszul duality acts on the $R$-matrix by
inversion, so $q \mapsto q^{-1}$. For the Kazhdan--Lusztig
equivalence, this corresponds to the duality
$\Rep_q(\frakg) \simeq \Rep_{q^{-1}}(\frakg)^{\mathrm{rev}}$,
which is the standard anti-braiding equivalence for quantum groups.
Item~(iii) follows formally from (i) and (ii), since
$\mathrm{av}(\cR^{-1}) = \kappa_{\mathrm{cat}}(\cA^!_{E_2})$.
The conjecture is conditional on Theorem~CY-B (the braided bar-cobar
adjunction at the $\Etwo$ level) and on
Conjecture~\ref{conj:braided-bar-uq} (braided bar cohomology
recovering the quantum group).
\end{remark}

\chapter{The Drinfeld Center and Bulk Algebras}
\label{ch:drinfeld-center}

The Drinfeld center $\cZ(\cC)$ of a monoidal category $\cC$ is the universal braided monoidal category receiving a forgetful functor from~$\cC$.  In the CY programme, it plays two roles: (i)~for $d = 3$, where the topological $E_3$-structure gives a trivial braiding at the topological level ($\pi_1(\mathrm{Conf}_2(\R^3)) = 1$), the quantum group braiding is recovered through $\cZ(\Rep^{E_1}(A))$; (ii)~the Drinfeld center is the categorical incarnation of the chiral derived center $\Zder(A) = C^\bullet_{\mathrm{ch}}(A, A)$ (Hochschild cochains) from Volume~I (Theorem~H), which computes the universal bulk algebra.  The identification $\cZ(\Rep^{E_1}(A)) \simeq \Rep^{E_2}(\Zder(A))$ (Ben-Zvi--Francis--Nadler, Lurie) is the categorical bulk-boundary correspondence.

\section{The Drinfeld center of a monoidal category}
\label{sec:drinfeld-center-def}

\begin{definition}[Drinfeld center]
\label{def:drinfeld-center}
For a monoidal category $\cC$, the \emph{Drinfeld center} $\cZ(\cC)$ is the category of pairs $(X, \beta_{X,-})$ where $X \in \cC$ and $\beta_{X,-} \colon X \otimes (-) \xrightarrow{\sim} (-) \otimes X$ is a half-braiding satisfying the hexagon axiom.  The Drinfeld center is always braided monoidal.
\end{definition}

\section{Drinfeld center as chiral derived center}
\label{sec:drinfeld-chiral-center}

\begin{proposition}
\label{prop:drinfeld-chiral-derived-center}
For an $E_1$-chiral algebra $A$ with monoidal representation category $\cC = \Rep^{E_1}(A)$, the Drinfeld center $\cZ(\cC)$ is equivalent to the representation category of the chiral derived center:
\[
  \cZ(\Rep^{E_1}(A)) \simeq \Rep^{E_2}(Z^{\mathrm{der}}_{\mathrm{ch}}(A)).
\]
This is the bulk-boundary correspondence: the $E_1$ boundary theory determines an $E_2$ bulk theory via the Drinfeld center.
% AP34: The Drinfeld center (= chiral derived center) is the UNIVERSAL BULK.
% It is NOT obtained by bar-cobar inversion (which recovers A itself), nor
% by Verdier duality (which produces B(A!)). The Drinfeld center is
% computed via Hochschild cochains: Z^{der}_{ch}(A) = RHom(Omega(B(A)), A).
% See Vol I, AP25 and AP34.
\end{proposition}

\section{CY enhancement of the Drinfeld center}
\label{sec:cy-drinfeld}

For a CY category $\cC$ of dimension $d$, the Drinfeld center $\cZ(\Rep^{E_1}(\Phi(\cC)))$ inherits a CY structure of dimension $d + 1$ from the Hochschild cochains (the bulk has one extra dimension from the closed-string direction).  This CY enhancement is the categorical analogue of the shifted-symplectic Lagrangian geometry of Volume~I (Theorem~C, ambient complementarity): the boundary theory $A$ determines a Lagrangian in the symplectic bulk $\Zder(A)$.

\section{Quantum group reconstruction from the center}
\label{sec:qg-reconstruction}

The reconstruction problem: given the braided monoidal category $\cZ(\cC)$, recover the quantum group $U_q(\frakg)$ whose representation category it is.  For finite tensor categories, Etingof--Nikshych--Ostrik reconstruction applies.  For the CY programme, the target is to reconstruct the quantum vertex chiral group $G(X)$ from the Drinfeld center of $\Rep^{E_1}(\Phi(\cC))$.  The Tannakian formalism applies when the fiber functor exists; for non-semisimple categories at roots of unity, the modified trace (Geer--Patureau-Mirand--Turaev) replaces the ordinary fiber functor.


% ==========================================
% PART V: The Standard Landscape
% ==========================================

\part{The Standard Landscape}
\label{part:examples}

\chapter{Fukaya Categories}
\label{ch:fukaya}

\section{The Fukaya category of a CY manifold}
\label{sec:fukaya-cy-manifold}

\begin{example}[Elliptic curve]
\label{ex:fukaya-elliptic}
For an elliptic curve $E_\tau$, $\Fuk(E_\tau)$ is CY of dimension $1$.  The Hochschild homology $\HH_\bullet(\Fuk(E_\tau))$ recovers the de Rham cohomology of $E_\tau$, and the cyclic structure encodes the symplectic pairing.  The associated chiral algebra is the Heisenberg vertex algebra $H_k$ at level $k = \mathrm{vol}(E_\tau)$.
\end{example}

\section{CY 2-folds: K3 and abelian surfaces}
\label{sec:fukaya-cy2}

\section{CY 3-folds}
\label{sec:fukaya-cy3}

\section{Wrapped Fukaya categories}
\label{sec:wrapped-fukaya}

\section{The quantum chiral algebra of a Fukaya category}
\label{sec:fukaya-qca}

\chapter{Derived Categories of CY Manifolds}
\label{ch:derived-cy}

Homological mirror symmetry identifies $\Fuk(X) \simeq D^b(\Coh(X^\vee))$.  Under the functor $\Phi$, this becomes an equivalence $A_{\Fuk(X)} \simeq A_{D^b(\Coh(X^\vee))}$ of quantum chiral algebras.  This chapter develops the algebraic side: the CY structure on $D^b(\Coh(X))$, the chiral algebra extracted by $\Phi$, exceptional collections, and Bridgeland stability conditions.

\section{$D^b(\Coh(X))$ for CY manifolds}
\label{sec:db-coh-cy}

For a smooth projective CY manifold $X$ of dimension $d$, the bounded derived category $D^b(\Coh(X))$ is CY of dimension $d$ with CY trace induced by Serre duality.  The Hochschild homology $\HH_\bullet(D^b(\Coh(X))) \simeq H^*(X, \Omega^\bullet_X)$ (HKR isomorphism) recovers the Hodge cohomology, and the CY trace is the Grothendieck residue.  The functor $\Phi$ applied to $D^b(\Coh(X))$ produces the B-model chiral algebra, with $\kappa = \chi(\cO_X)$ for CY surfaces.

\section{Homological mirror symmetry and the chiral algebra}
\label{sec:hms-chiral}

The homological mirror symmetry conjecture (Kontsevich) identifies $\Fuk(X) \simeq D^b(\Coh(X^\vee))$.  Under the CY-to-chiral functor $\Phi$, this becomes an equivalence of quantum chiral algebras:
\[
  A_{\Fuk(X)} \simeq A_{D^b(\Coh(X^\vee))}.
\]

\section{Exceptional collections and chiral algebras}
\label{sec:exceptional-chiral}

When $D^b(\Coh(X))$ admits a full exceptional collection $\langle E_1, \ldots, E_n \rangle$, the CY-to-chiral functor reduces to a finite-dimensional computation: the $\Ainf$-algebra of the collection determines a quiver with potential, and $\Phi$ produces the critical CoHA of the quiver (connecting to Chapter~\ref{ch:toric-coha}).  This provides the simplest non-trivial test cases for Theorem~CY-A.

\section{Stability conditions and the chiral moduli}
\label{sec:stability-chiral}

Bridgeland stability conditions on $D^b(\Coh(X))$ parametrise the space of t-structures, providing a moduli space $\Stab(X)$ over which the chiral algebra $\Phi(D^b(\Coh(X)))$ varies.  Wall-crossing in $\Stab(X)$ corresponds to mutations of the bar complex; the global chiral algebra is assembled as a homotopy colimit over stability chambers (Programme~A of the $K3 \times E$ programme, Chapter~\ref{ch:toroidal-elliptic}).

\chapter{Matrix Factorizations}
\label{ch:matrix-factorizations}

The third source of CY categories is the Landau--Ginzburg model: a polynomial $W \colon \mathbb{C}^n \to \mathbb{C}$ whose category $\MF(W)$ of matrix factorizations is CY of dimension $n - 2$.  For ADE singularities $W = x^N + y^2 + z^2 + w^2$, $\MF(W)$ is CY$_2$ and $\Phi$ (Theorem~CY-A$_2$) produces chiral algebras related to $\cW_N$-algebras.  The LG/CY correspondence $\MF(W) \simeq D^b(\Coh(X_W))$ provides a consistency check: $\Phi(\MF(W))$ and $\Phi(D^b(\Coh(X_W)))$ (Chapter~\ref{ch:derived-cy}) must agree.

\section{Matrix factorizations and Landau--Ginzburg models}
\label{sec:mf-lg}

\section{The CY structure on $\MF(W)$}
\label{sec:mf-cy-structure}

\section{Quantum chiral algebras from singularities}
\label{sec:mf-qca}

\section{ADE singularities and W-algebras}
\label{sec:ade-w-algebras}

The connection between ADE singularities and W-algebras provides a key test case for the CY-to-chiral functor.  For the $A_{N-1}$ singularity $W = x^N$ in one variable, $\MF(x^N)$ is CY of dimension $-1$ (formal), so the relevant setting is the two-variable singularity $W = x^N + y^2$, giving $\MF(W)$ of CY dimension $0$ (semisimple).  To reach a CY$_2$ category related to $W_N$, one considers the threefold singularity $W = x^N + y^2 + z^2 + w^2$ in four variables, yielding $\MF(W)$ of CY dimension $2$.
% RECTIFICATION-FLAG: The original claimed MF(x^N) "should produce a
% quantum chiral algebra related to W_N" without specifying the ambient
% dimension. The CY dimension of MF(W) depends on the number of variables.
% For the W-algebra connection, one needs a CY_2 category (to apply
% Theorem CY-A_2), which requires 4 variables.

\chapter{Quantum Group Representations}
\label{ch:quantum-group-reps}

The two-stage factorisation $\Phi_d = \SpCh_{\Sigma_{d-1}, C} \circ \PhiFA_d$ of Chapter~\ref{ch:e1-chiral} (Theorem~\ref{thm:phi-two-stage-factorisation}) sends a CY$_d$
category $\cC$ to a chiral algebra $\Phi_d(\cC)$ on a reference curve~$C$.
At $d \geq 3$ this output is natively $\Eone$; the $\Etwo$-braided
structure belongs to the Drinfeld centre
$\cZ(\Rep^{\Eone}(\Phi_d(\cC)))$
(Chapter~\ref{ch:drinfeld-center}; Remark~\ref{rem:e2-on-drinfeld-centre}),
not to $\Phi_d(\cC)$ itself. Conjecture~CY-C asks when this braided
category is a quantum-group representation category. The constructed
and conditional criteria separate into three loci:
\begin{enumerate}[label=(\alph*)]
 \item \emph{Kazhdan--Lusztig} at $d = 2$: affine vertex algebra
 $V_k(\frakg)$ versus quantum group $\Uq(\frakg)$ at
 $q = e^{\pi i/(k+h^\vee)}$;
 \item \emph{Yangian RTT} from the ordered collision residue
 $r(z)$ of the Vol~I bar complex;
 \item \emph{BPS / CoHA} at $d = 3$: toric Donaldson--Thomas
 invariants assemble into an associative positive half; the full
 quantum group is reached only after adjoining the negative half,
 Cartan part, Hopf pairing, completion, and Drinfeld double.
\end{enumerate}
Their mutual compatibility is the content of CY-C on the constructed
locus and requires additional comparison data elsewhere.


\section{\texorpdfstring{$\Rep_q(\frakg)$}{Rep-q(frakg)} as a braided monoidal category}
\label{sec:rep-q-braided}

\begin{definition}[Category \texorpdfstring{$\Rep_q(\frakg)$}{Rep-q(frakg)}]
\label{def:rep-q}
Let $\frakg$ be a finite-dimensional simple Lie algebra and
$q \in \C^\times$. The category $\Rep_q(\frakg)$ has:
\begin{enumerate}[label=(\roman*)]
 \item \emph{Objects}: finite-dimensional $\Uq(\frakg)$-modules
 (Definition~\ref{def:qgf-drinfeld-jimbo});
 \item \emph{Morphisms}: $\Uq(\frakg)$-linear maps;
 \item \emph{Monoidal structure}: the tensor product
 $V \otimes W$ with $\Uq(\frakg)$-action via the coproduct
 $\Delta$;
 \item \emph{Braiding}: $\sigma_{V,W} = \tau_{V,W} \circ \cR_{V,W}$
 where $\cR$ is the universal $R$-matrix of
 Remark~\ref{rem:qgf-universal-R-formal} and $\tau$ is the
 transposition $v \otimes w \mapsto w \otimes v$.
\end{enumerate}
\end{definition}

\begin{proposition}[Semisimplicity dichotomy]
\label{prop:semisimplicity-dichotomy}
The structure of $\Rep_q(\frakg)$ depends on whether $q$ is generic
or a root of unity:
\begin{enumerate}[label=(\roman*)]
 \item \emph{Generic $q$} (i.e., $q$ not a root of unity):
 $\Rep_q(\frakg)$ is semisimple, with simple objects $V_q(\lambda)$
 parametrized by dominant integral weights $\lambda \in P^+$.
 The characters are the same as the classical case:
 $\mathrm{ch}(V_q(\lambda)) = \mathrm{ch}(V(\lambda))$
 (the Weyl character formula is $q$-independent).
 \item \emph{Root of unity $q = e^{\pi i/(k+h^\vee)}$} in the
 simply-laced normalization, with
 $k \in \Z_{>0}$: $\Rep_q(\frakg)$ is non-semisimple; the
 \emph{fusion category} $\cC_q(\frakg)$ is the quotient by
 negligible morphisms, with finitely many simple objects
 $\{V_q(\lambda) : \lambda \in P^+_k\}$ parametrized by the
 level-$k$ alcove $P^+_k = \{\lambda \in P^+ :
 \langle \lambda, \theta^\vee \rangle \leq k\}$.
\end{enumerate}
This is the Lusztig--Andersen--Paradowski root-of-unity
semisimplification theorem, in the convention where the lacing number
is absorbed into the definition of~$q$.
\end{proposition}

\begin{example}[\texorpdfstring{$\fsl_2$}{fsl-2} at generic \texorpdfstring{$q$}{q}]
\label{ex:sl2-generic}
For $\frakg = \fsl_2$, the simple modules are
$V_q(n)$ ($n \in \Z_{\geq 0}$) with $\dim V_q(n) = n+1$. The
tensor product decomposition is
$V_q(m) \otimes V_q(n) \simeq \bigoplus_{k=|m-n|}^{m+n} V_q(k)$
(same as classical, with $k \equiv m+n \pmod{2}$). The $R$-matrix
on $V_q(1) \otimes V_q(1)$ is
\[
 \cR = q^{1/2}
 \begin{pmatrix}
 q & 0 & 0 & 0 \\
 0 & 1 & q - q^{-1} & 0 \\
 0 & 0 & 1 & 0 \\
 0 & 0 & 0 & q
 \end{pmatrix}
\]
in the basis $\{e_1 \otimes e_1, e_1 \otimes e_2,
e_2 \otimes e_1, e_2 \otimes e_2\}$. This is the standard
solution of the Yang--Baxter equation for the 6-vertex model.
\end{example}

\begin{example}[\texorpdfstring{$\fsl_2$}{fsl-2} at root of unity]
\label{ex:sl2-root-of-unity}
At $q = e^{\pi i/(k+2)}$, the fusion category $\cC_q(\fsl_2)$ has
$k+1$ simple objects $\{V_q(0), V_q(1), \ldots, V_q(k)\}$ with
fusion rules
\[
 V_q(m) \otimes V_q(n)
 = \bigoplus_{\substack{j = |m-n| \\ j \equiv m+n \,(\mathrm{mod}\,2)}}^{\min(m+n, \, 2k-m-n)}
 V_q(j).
\]
The quantum dimension is
$\dim_q V_q(n) = [n+1]_q = (q^{n+1} - q^{-n-1})/(q - q^{-1})$
and the total quantum dimension squared is
$\mathcal{D}^2 = \sum_{j=0}^{k} (\dim_q V_q(j))^2
= (k+2)/(2\sin^2(\pi/(k+2)))$.
The $S$-matrix entries are
$S_{mn} = \sqrt{2/(k+2)} \, \sin(\pi(m+1)(n+1)/(k+2))$.
This is a modular tensor category: the Reshetikhin--Turaev functor
produces the $SU(2)$ Chern--Simons invariants of 3-manifolds.
\end{example}


\section{The \texorpdfstring{$R$}{R}-matrix as categorical \texorpdfstring{$r(z)$}{r(z)}}
\label{sec:r-matrix-categorical}

The $R$-matrix of $\Uq(\frakg)$ is the categorical incarnation of
the collision residue $r(z)$ from the Volume~I bar complex.

\begin{proposition}[\texorpdfstring{$R$}{R}-matrix from bar degree \texorpdfstring{$(1,1)$}{(1,1)}]
\label{prop:r-matrix-bar}
For the affine vertex algebra $V_k(\frakg)$ at level $k$ with
$q = e^{\pi i/(k+h^\vee)}$, the degree-$(1,1)$ component of the
ordered bar complex $B^{\mathrm{ord}}(V_k(\frakg))$ gives the
classical collision residue whose Etingof--Kazhdan quantization is
the universal $R$-matrix of $\Uq(\frakg)$:
\[
 r(z) = \Res^{\mathrm{coll}}_{0,2}(\Theta_{V_k(\frakg)})
 \;\longleftrightarrow\;
 \cR_q = \lim_{z \to 0}\, \cR(z)
\]
where $r(z) = k\,\Omega_{\mathrm{tr}}/z + O(1)$ is the classical $r$-matrix with
Casimir $\Omega \in \frakg \otimes \frakg$ and level prefix $k$,
and $\cR_q$ is the quantized universal $R$-matrix. The spectral
parameter $z$ is the coordinate difference $z = z_1 - z_2$ of two
insertion points on the curve, a rapidity variable in the
integrable lattice interpretation. The pole order of $r(z)$ at
$z = 0$ is one less than that of the OPE, reflecting the fact that
$r(z)$ records the commutator rather than the product.
\end{proposition}

\begin{remark}[Three \texorpdfstring{$r$}{r}-matrices;
\ClaimStatusDefinitional]
\label{rem:three-r-matrices-qg}
Three $r$-matrix-like objects appear:
\begin{enumerate}[label=(\alph*)]
 \item $r^{\mathrm{coll}}(z)$: the bar-intrinsic collision residue
 (Volume~I);
 \item $r(z) = k\,\Omega_{\mathrm{tr}}/z$: the classical $r$-matrix at level $k$
 (Belavin--Drinfeld; vanishes at $k = 0$);
 \item $\cR_q$: the universal $R$-matrix of $\Uq(\frakg)$
 (Drinfeld--Jimbo).
\end{enumerate}
The passage $r(z) \rightsquigarrow \cR_q$ is the Etingof--Kazhdan
quantization theorem. For even $\Einf$-algebras,
$r^{\mathrm{coll}} = r$; for odd generators or genuinely $\Eone$
algebras they can differ.
\end{remark}


\section{Kazhdan--Lusztig equivalence}
\label{sec:kl-equivalence}

The Kazhdan--Lusztig--Finkelberg equivalence is the positive-level
root-of-unity prototype for CY-C.

\begin{theorem}[Kazhdan--Lusztig--Finkelberg;
\ClaimStatusProvedElsewhere]
\label{thm:kl-equivalence}
For a simple Lie algebra $\frakg$ and a positive integer $k$, with
$q = e^{\pi i/(k+h^\vee)}$ in the simply-laced normalization, there
is an equivalence of modular tensor categories
\[
 \mathcal{MTC}_q(\frakg) \;\simeq\; \mathcal{O}^{\mathrm{int}}_k(\hat{\frakg})
\]
where $\mathcal{MTC}_q(\frakg)$ is the semisimplified tilting-module
category of $\Uq(\frakg)$ at a root of unity, and
$\mathcal{O}^{\mathrm{int}}_k(\hat{\frakg})$ is the integrable
highest-weight category of affine Kac--Moody modules at level~$k$
with KZ braiding. In non-simply-laced type the lacing number enters
the denominator of~$q$.
\end{theorem}

\begin{proof}[Attribution]
Kazhdan--Lusztig, \emph{J.\ Amer.\ Math.\ Soc.}~1993--1994, construct
the tensor equivalence for affine Lie algebra categories. Finkelberg,
\emph{GAFA}~1996, proves the positive-level fusion equivalence with the
semisimplified root-of-unity quantum group category. Andersen--Paradowski
give the tilting-module semisimplification used for
$\mathcal{MTC}_q(\frakg)$.
\end{proof}

\begin{remark}[Bar-complex interpretation]
\label{rem:kl-bar-complex}
From the bar-complex perspective, the equivalence states that
the $\Eone$-representation category of the affine vertex algebra
$V_k(\frakg)$ recovers the quantum group representation category.
The Volume~I bar complex $B(V_k(\frakg))$ encodes the quantum group
$R$-matrix in its degree-$(1,1)$ component; the DK bridge (Volume~II,
MC3 on the evaluation-generated core for all simple types) extends this
to the categorical
equivalence.
\end{remark}

\begin{proposition}[KL and the DK bridge]
\label{prop:kl-dk-bridge}
The Kazhdan--Lusztig--Finkelberg equivalence factors through the Drinfeld--Kohno
bridge (Volume~II):
\[
\begin{tikzcd}[column sep=large]
 \mathcal{O}^{\mathrm{int}}_k(\hat{\frakg}) \ar[r, "\sim", "\text{KLF}"'] &
 \mathcal{MTC}_q(\frakg) \\
 \Rep^{\Eone}(V_k(\frakg)) \ar[u, hook] \ar[r, "\text{DK}"] &
 \Rep_q(\frakg) \ar[u, two heads]
\end{tikzcd}
\]
The vertical arrows are: inclusion of integrable affine modules into all
$V_k(\frakg)$-modules (left), and quotient by negligible morphisms
(right). The DK bridge (MC3) provides the horizontal equivalence
at the level of compact objects.
\end{proposition}


\section{Conjecture CY-C: quantum group realization}
\label{sec:conj-cy-c}

\begin{conjecture}[Quantum group realization: Conjecture CY-C]
\label{conj:qg-realization}
For a simple Lie algebra $\frakg$ and $q = e^{\pi i/(k+h^\vee)}$
with $k$ a positive integer, there exists a CY$_2$ category
$\cC(\frakg, q)$ such that:
\begin{enumerate}[label=(\roman*)]
 \item $\Phi(\cC(\frakg, q))$ is an $\Etwo$-chiral algebra whose
 underlying vertex algebra is $V_k(\frakg)$;
 \item $\Rep^{\Etwo}(\Phi(\cC(\frakg, q))) \simeq \Rep_q(\frakg)$
 as braided monoidal categories;
 \item $\kappa_{\mathrm{ch}}(\Phi(\cC(\frakg, q)))
 = \dim(\frakg) \cdot (k + h^\vee)/(2h^\vee)$, recovering the
 Volume~I modular characteristic of $V_k(\frakg)$ as the chiral
 invariant of the $\Phi$-output.
\end{enumerate}
\end{conjecture}

\begin{remark}[Proof inputs for Conjecture CY-C]
\label{rem:cy-c-proof-inputs}
Item (i) is constructed for $\frakg = \fsl_N$ via the Bridgeland
stability conditions on the CY$_2$ resolution of the $A_{N-1}$
surface singularity. Item (ii) at the $\Eone$ level is the content
of MC3 (proved on the evaluation-generated core for all simple types, Volume~I/II). The $\Etwo$
upgrade requires the Drinfeld center passage
(Chapter~\ref{ch:drinfeld-center}). Item (iii) follows from
Theorem~CY-D (Theorem~\ref{thm:cy-modular-characteristic}) when
the CY category is constructed.
\end{remark}


\section{Yangian and RTT realizations}
\label{sec:yangian-rtt}

\begin{definition}[Yangian]
\label{def:yangian}
The \emph{Yangian} $Y(\frakg)$ is the $\C[\hbar]$-algebra with
generators $\{x^{(r)}_i : i = 1, \ldots, \dim \frakg, \;
r = 0, 1, 2, \ldots\}$ and relations determined by the RTT
presentation: $R_{12}(u-v) \, T_1(u) \, T_2(v)
= T_2(v) \, T_1(u) \, R_{12}(u-v)$ where $T(u) = \sum_r T_r u^{-r}$
is the generating matrix and $R_{12}(u)$ is the Yang $R$-matrix.
\end{definition}

\begin{proposition}[Yangian from collision residue]
\label{prop:yangian-collision}
On the evaluation-generated core, the degree-$(1,1)$ collision residue
of the ordered bar complex determines the RTT presentation of the
Yangian:
\[
 Y(\frakg) = \langle r^{\mathrm{coll}}(z) \rangle_{\mathrm{RTT}},
\]
where the RTT relation is the degree-$(1,1,1)$ bar differential after
Etingof--Kazhdan quantization. This is the evaluation-generated-core
content of Volume~II, MC3: the ordered bar complex
$B^{\mathrm{ord}}(V_k(\frakg))$ produces the Yangian presentation from
the modes of $r(z)=k\,\Omega_{\mathrm{tr}}/z$.
\end{proposition}

\begin{remark}[Three distinct operations]
\label{rem:three-operations}
The following three operations on quantum groups must never be
conflated:
\begin{enumerate}[label=(\alph*)]
 \item \emph{Koszul duality}:
 $V_k(\frakg) \mapsto V_k(\frakg)^!$, the Koszul dual vertex
 algebra involving the dual level
 $k' = -k - 2h^\vee$;
 \item \emph{Feigin--Frenkel involution}:
 $k \mapsto -k - 2h^\vee$ within the same family of affine
 algebras; at the critical level $k = -h^\vee$, this produces
 the Feigin--Frenkel center
 $\mathfrak{z}(\hat{\frakg}_{-h^\vee}) = \mathrm{Fun}(
 \mathrm{Op}_{G^L})$;
 \item \emph{Langlands duality}:
 $\frakg \mapsto \frakg^L$ (root-coroot exchange), producing
 a genuinely different Lie algebra for non-simply-laced types.
\end{enumerate}
For $\fsl_2$ (self-dual, self-Langlands), all three coincide on
representation categories. For $G_2$:
$G_2^L = G_2$ (self-Langlands) but the Koszul dual
$V_k(G_2)^!$ is at level $k' = -k - 8$ (since $h^\vee(G_2) = 4$),
while the FF involution also sends $k \mapsto -k-8$; these
coincide as level maps but produce different algebras.
\end{remark}


\section{BPS algebras and quantum groups}
\label{sec:bps-algebras}

The BPS algebra of a CY$_3$ category provides a physical
construction of quantum groups.

\begin{definition}[BPS algebra]
\label{def:bps-algebra}
For a CY$_3$ category $\cC$ with stability condition $\sigma$,
the \emph{BPS algebra} $\cA^{\BPS}_\sigma(\cC)$ is the
cohomological Hall algebra (CoHA) of $\cC$:
\[
 \cA^{\BPS}_\sigma(\cC)
 = \bigoplus_{\gamma \in K_0(\cC)}
 H^*_{\mathrm{BM}}(\cM_\sigma(\gamma), \, \phi_{\mathrm{tr}W})
\]
where $\cM_\sigma(\gamma)$ is the moduli of $\sigma$-semistable
objects of class $\gamma$, $W$ is the CY potential, and
$\phi_{\mathrm{tr}W}$ is the vanishing cycle sheaf. The
multiplication is the Hall product (from short exact sequences).
\end{definition}

\begin{proposition}[CoHA as quantum group]
\label{prop:coha-quantum-group}
For the tripled Jordan quiver (one vertex, three loops $X,Y,Z$) with
potential $W = \mathrm{tr}(XYZ-XZY)$ and CY$_3$ geometry $\C^3$:
\begin{enumerate}[label=(\roman*)]
 \item The CoHA is isomorphic, as an associative Hall algebra, to the positive half of the affine
 Yangian $Y(\widehat{\fgl}_1)^+$ (Schiffmann--Vasserot,
 Maulik--Okounkov);
 \item The full affine Yangian $Y(\widehat{\fgl}_1)$ is recovered only
 after adjoining the negative half, Cartan part, Hopf pairing, and
 completed Drinfeld double of the CoHA;
 \item The quantum toroidal algebra
 $U_{q,t}(\widehat{\widehat{\fgl}}_1)$ arises from the
 $K$-theoretic refinement (replacing cohomology by $K$-theory).
\end{enumerate}
The CoHA itself is not a vertex algebra and not the full affine
Yangian. The $\cW_{1+\infty}$ vertex algebra appears only after the
Drinfeld-double/centre passage and an evaluation representation on a
vacuum module.
\end{proposition}

\begin{proof}[Attribution]
Schiffmann--Vasserot, \emph{Publ.\ IHES}~118 (2013), Theorem~9.1,
identifies the cohomological Hall algebra of the Hilbert scheme model
with the positive half of the affine Yangian of $\widehat{\fgl}_1$ in
the normalisation used for $\C^3$ DT theory. Kapranov--Vasserot,
\emph{Compos.\ Math.}~154 (2018), Theorem~5.1, gives the $K$-theoretic
upgrade. The passage from a positive half to a full quasitriangular
Hopf algebra is the Drinfeld-double construction and requires the
negative half, Cartan part, pairing, and completion.
\end{proof}

\begin{remark}[Slab as bimodule]
\label{rem:bps-slab}
In the Dimofte framework of Volume~II, the BPS algebra
arises from the 3d holomorphic-topological theory on the slab
$X \times [0,1]$. The slab has \emph{two} boundary components
($X \times \{0\}$ and $X \times \{1\}$), making the algebra of
operators on the slab a bimodule for the two boundary algebras.
When the Hall pairing and completion exist, the Drinfeld double is the
endomorphism algebra of the identity bimodule. This bimodule structure
is essential: a Swiss-cheese
disk has one closed and one open boundary component; the slab has
two open boundary components.
\end{remark}

\begin{example}[Resolved conifold]
\label{ex:resolved-conifold}
For the resolved conifold $\cO(-1) \oplus \cO(-1) \to \bP^1$, the
CY$_3$ category is $D^b(\Coh(\cO(-1)^{\oplus 2}))$ and the BPS
algebra is the positive Hall half whose completed Drinfeld double is
the affine Yangian $Y(\widehat{\fsl}_2)$ on the toric chamber. The
chiral modular characteristic is
$\kappa_{\mathrm{ch}} = 1$
(direct McKay computation; the conifold is not a local surface; see
Corollary~\ref{cor:conifold-non-local-surface}). The wall-crossing
formula of Kontsevich--Soibelman governs the dependence of the
BPS spectrum on the stability condition.
\end{example}


\section{Wall-crossing and Kontsevich--Soibelman}
\label{sec:wall-crossing-ks}

The BPS spectrum of a CY$_3$ category depends on the stability
condition $\sigma \in \Stab(\cC)$. As $\sigma$ crosses a wall of
marginal stability, the BPS states reorganize: bound states form or
decay. The Kontsevich--Soibelman wall-crossing formula governs
this reorganization at the level of the motivic Hall algebra.

\begin{definition}[Wall of marginal stability]
\label{def:wall-marginal}
A \emph{wall of marginal stability} in $\Stab(\cC)$ is a real
codimension-$1$ locus $W_{\gamma_1, \gamma_2}$ where two classes
$\gamma_1, \gamma_2 \in K_0(\cC)$ with $\gamma = \gamma_1 + \gamma_2$
have aligned central charges:
$Z_\sigma(\gamma_1)/Z_\sigma(\gamma_2) \in \R_{>0}$.
On opposite sides of $W_{\gamma_1, \gamma_2}$, a bound state of
class $\gamma$ may exist or not.
\end{definition}

\begin{theorem}[Kontsevich--Soibelman wall-crossing formula]
\label{thm:ks-wall-crossing}
For a CY$_3$ category $\cC$ with stability condition $\sigma$,
define the \emph{BPS automorphism} for each class
$\gamma \in K_0(\cC)$:
\[
 U_\gamma = \exp\!\left(
 \sum_{n=1}^{\infty} \frac{(-1)^{n \dim \cM(\gamma)}}{n^2}\,
 \Omega(\gamma; \sigma) \cdot \hat{e}_{n\gamma}
 \right)
\]
where $\Omega(\gamma; \sigma) \in \Z$ is the (numerical)
Donaldson--Thomas invariant and $\hat{e}_\gamma$ is the
corresponding element of the quantum torus algebra
$\hat{e}_\gamma \hat{e}_{\gamma'} = (-1)^{\langle \gamma, \gamma' \rangle}
\hat{e}_{\gamma + \gamma'}$. The ordered product
\[
 \cA_\sigma = \prod_{\gamma : \arg Z_\sigma(\gamma) = \theta}^{\curvearrowleft}
 U_\gamma
\]
(clockwise ordering by phase $\theta = \arg Z_\sigma(\gamma)$) is
\emph{independent of $\sigma$}: as $\sigma$ crosses a wall, the
individual $\Omega(\gamma; \sigma)$ jump, but the ordered product
$\cA_\sigma$ is invariant.
\end{theorem}

\begin{remark}[Motivic refinement]
\label{rem:ks-motivic}
The numerical DT invariants $\Omega(\gamma; \sigma) \in \Z$ are
shadows of the motivic DT invariants
$[\cM_\sigma(\gamma)]_{\mathrm{vir}} \in K_0(\mathrm{Var})$ in the
Grothendieck ring of varieties. The motivic wall-crossing formula
replaces $\Omega$ by the motivic class, and the quantum torus by
the motivic quantum torus. The identification ``scattering =
shadow'' (Volume~I) holds at the motivic level but fails at
the naive BCH level: the full motivic Hall algebra is required.
\end{remark}

\begin{conjecture}[Wall-crossing and the bar complex]
\label{conj:wc-bar}
The wall-crossing formula has a bar-complex interpretation:
crossing a wall $W_{\gamma_1, \gamma_2}$ in $\Stab(\cC)$ induces
a mutation of the ordered bar complex
$B^{\mathrm{ord}}(A_\cC)$ at the corresponding degree.
Concretely:
\begin{enumerate}[label=(\roman*)]
 \item The BPS automorphism $U_\gamma$ corresponds to the
 degree-$|\gamma|$ component of the MC element
 $\Theta_{A_\cC}^{E_1}$ in the $\Eone$ convolution algebra;
 \item Wall-crossing preserves the full MC element
 $\Theta_{A_\cC}^{E_1}$ (it is an invariant of the CY category
 $\cC$, not of the stability condition);
 \item The individual BPS invariants $\Omega(\gamma; \sigma)$ are
 \emph{projections} of $\Theta_{A_\cC}^{E_1}$ that depend on
 $\sigma$; the MC element itself is wall-crossing invariant.
\end{enumerate}
\end{conjecture}


\section{The chiral modular characteristic \texorpdfstring{$\kappa_{\mathrm{ch}}$}{kappa-ch}}
\label{sec:kappa-cat}

The modular characteristic of the quantum group representation
category connects the categorical datum to the genus expansion of
Volume~I.

\begin{proposition}[\texorpdfstring{$\kappa_{\mathrm{ch}}$}{kappa-ch} on the chiral-algebra side of a quantum-group correspondence]
\label{prop:kappa-cat-quantum-groups}
Assume Conjecture~\ref{conj:qg-realization}. For the affine vertex
algebra $V_k(\frakg)$ at level $k$ with
$q = e^{\pi i/(k + h^\vee)}$, the chiral modular characteristic of
$V_k(\frakg)$ (the chiral-algebra image under $\Phi_2$ of the CY$_2$
category $\cC(\frakg, q)$; Conjecture~\ref{conj:qg-realization}) is:
\[
 \kappa_{\mathrm{ch}}(V_k(\frakg))
 = \dim(\frakg) \cdot \frac{k + h^\vee}{2h^\vee}
\]
(Volume~I landscape census). This is the \emph{chiral} invariant,
distinct from the categorical Euler characteristic
$\kappa_{\mathrm{cat}}(\cC) = \chi(\cO_X)$, which may vanish while
$\kappa_{\mathrm{ch}}$ does not. For the standard families:
\begin{center}
\renewcommand{\arraystretch}{1.3}
\begin{tabular}{lll}
 \toprule
 $\frakg$ & $\kappa_{\mathrm{ch}}(V_k(\frakg))$ & $q$ \\
 \midrule
 $\fsl_2$ & $3(k+2)/4$ & $e^{\pi i/(k+2)}$ \\
 $\fsl_N$ & $\frac{N^2-1}{2N}(k+N)$ & $e^{\pi i/(k+N)}$ \\
 $\mathfrak{so}_{2n}$ & $\frac{n(2n-1)(k+2n-2)}{2(2n-2)}$
 & $e^{\pi i/(k+2n-2)}$ \\
 $E_8$ & $\frac{248(k+30)}{60}$
 & $e^{\pi i/(k+30)}$ \\
 \bottomrule
\end{tabular}
\end{center}
\end{proposition}

\begin{remark}[Four \texorpdfstring{$\kappa_\bullet$}{kappa-.}-invariants]
\label{rem:kappa-cat-vs-others}
Four invariants coexist and must be distinguished. Two are
manifold invariants, independent of the algebraization:
\begin{enumerate}[label=(\roman*)]
 \item $\kappa_{\mathrm{cat}}(X) = \chi(\cO_X)$: the holomorphic
 Euler characteristic, a topological invariant;
 \item $\kappa_{\mathrm{fiber}}$: the rank of the relevant lattice
 (e.g.\ Mukai lattice of $H^*(K3, \Z)$).
\end{enumerate}
Two are algebraization invariants, functions of the chiral algebra
produced by a specific construction:
\begin{enumerate}[resume,label=(\roman*)]
 \item $\kappa_{\mathrm{ch}}(A)$: the chiral modular characteristic
 of a chiral algebra $A$, equal to the genus-$1$ Hodge class
 coefficient (Volume~I, Theorem~D);
 \item $\kappa_{\mathrm{BKM}}$: the Borcherds--Kac--Moody modular
 weight, a number-theoretic invariant attached to a BKM algebra
 constructed from a lattice (Chapter~\ref{chap:toroidal-elliptic}).
\end{enumerate}
For $K3 \times E$:
$\kappa_{\mathrm{cat}}(K3 \times E) = \chi(\cO_{K3}) \cdot \chi(\cO_E)
 = 2 \cdot 0 = 0$ by K\"unneth. The compact Hodge/PhiFA scalar is
$\kappa_{\mathrm{ch}}^{\mathrm{Hodge}}(\Phi(K3 \times E))=0$; the additive number
$2+1=3$ is the separate Mukai--Heisenberg specialisation
$\kappa_{\mathrm{ch}}^{\mathrm{Heis}}(K3\times E)$. The BKM weight is
$\kappa_{\mathrm{BKM}}(\Delta_5) = \left.c_N(0)/2\right|_{N=1} = 5$
from the Gritsenko denominator $\Delta_5$ by the universal Borcherds
weight formula (Theorem~\ref{thm:borcherds-weight-kappa-BKM-universal});
the primitive Igusa square is
$\Phi_{10}^{\mathrm{un}}=\Delta_5^2$ of weight~$10$, not the BKM weight.
The identification $\kappa_{\mathrm{ch}} = \chi(\cO_X)$ holds for
CY$_2$ with $h^{1,0} = 0$ (where Serre duality kills the quantum
correction) but fails at odd $d$ by Serre parity
$\chi(\cO_X) = 0$, and fails at $d \geq 3$ in general: the four
invariants are controlled by four different functorialities.
\end{remark}

\begin{proposition}[Chiral scalar complementarity under Feigin--Frenkel involution]
\label{prop:kappa-cat-complementarity}
Under Feigin--Frenkel Koszul duality
$V_k(\frakg) \mapsto V_{k'}(\frakg)$ with $k' = -k - 2h^\vee$, the
chiral modular characteristics of the two chiral-algebra sides satisfy
\[
 \kappa_{\mathrm{ch}}(V_k(\frakg))
 + \kappa_{\mathrm{ch}}(V_{k'}(\frakg))
 = 0,
\]
where on the CY side $q = e^{\pi i/(k + h^\vee)}$ and
$q' = e^{-\pi i/(k + h^\vee)}$. This is the KM/free-field clause of
the Volume~I scalar complementarity (Volume~I, Theorem~C); it is
\emph{not} a statement about the categorical invariant
$\kappa_{\mathrm{cat}}$ (which is a holomorphic Euler characteristic
and behaves additively on Koszul-dual pairs only when the underlying
category is fixed). The $\rho_A \cdot K$ Virasoro/$\cW_N$ clause is
separate.
\end{proposition}

\begin{proof}
Under $k \mapsto k' = -k - 2h^\vee$ the Sugawara-shifted scalar
transforms as
$\kappa_{\mathrm{ch}}(V_{k'}(\frakg))
 = \dim(\frakg)(k' + h^\vee)/(2h^\vee)
 = -\dim(\frakg)(k + h^\vee)/(2h^\vee)
 = -\kappa_{\mathrm{ch}}(V_k(\frakg))$,
so the sum vanishes.
\end{proof}


\section{Shadow obstruction tower for quantum group categories}
\label{sec:shadow-tower-qg}

The shadow obstruction tower $\Theta_{A_\cC}$
(Volume~I) acquires categorical meaning through the
quantum group lens.

\begin{proposition}[Shadow depth from \texorpdfstring{$R$}{R}-matrix pole structure]
\label{prop:shadow-depth-r-matrix}
For $V_k(\frakg)$ with $R$-matrix
$R(z) = 1 + r(z) + O(r^2)$ where $r(z) = \frac{k\,\Omega}{z}$:
\begin{enumerate}[label=(\roman*)]
 \item The shadow depth $r_{\max} = 3$ (class~L): the collision
 residue $r(z)$ has a single pole, the cubic shadow $C$ is
 nonzero (from the Lie bracket), and the quartic resonance class
 $Q$ vanishes by semisimplicity of $\frakg$;
 \item The averaging map (Volume~I) sends
 $r(z) = \frac{k\,\Omega}{z} \mapsto \kappa_{\mathrm{ch}}(V_k(\frakg))$: the full
 $z$-dependent profile is killed by $\Sigma_2$-coinvariance,
 leaving only the scalar modular characteristic;
 \item The KZ associator $\Phi_{\mathrm{KZ}}$ (the degree-$3$
 component of $\Theta_{A_\cC}^{E_1}$) maps under averaging to
 the cubic shadow $C$.
\end{enumerate}
\end{proposition}

\begin{remark}[From quantum groups to the shadow tower]
\label{rem:qg-to-shadow}
The passage from quantum group data to the shadow obstruction tower
inverts the construction of this chapter: the quantum group
$\Uq(\frakg)$ determines the $R$-matrix, which is the degree-$(1,1)$
component of the $\Eone$ MC element $\Theta^{E_1}$, whose
$\Sigma_n$-coinvariant projections are the shadow tower. The
reconstruction problem (from quantum group to chiral algebra) is
solved on the evaluation-generated core by the DK bridge (MC3 for all
simple types on that core): the categorical
data of $\Rep_q(\frakg)$ uniquely determines the chiral algebra
$V_k(\frakg)$ up to the completion issues of MC4.
\end{remark}


\section{Distinctions between Yangian definitions}
\label{sec:two-yangian-definitions}

The Yangian admits two genuinely distinct definitions on a curve.
The weaker one is an $\Eone$-factorisation algebra with RTT
relations; the stronger one supplements this with a chiral
coproduct, spectral $S$-matrix, and the full modular Maurer--Cartan
tower. Equivalence of the two definitions is an open problem at
the heart of the Vol~II E$_1$-core.

\begin{definition}[\texorpdfstring{$\Eone$}{E1}-chiral Yangian (weaker)]
\label{def:e1-chiral-yangian-vol3}
An \emph{$\Eone$-chiral Yangian} on a smooth algebraic curve $X$ is
an $\Eone$-factorisation algebra $\cY^{\mathrm{ch}}$ on $\Ran(X)$
whose local operator algebra, in the RTT presentation, satisfies
the spectral RLL relation
\[
 R_{12}(u-v)\, T_1(u)\, T_2(v)
 = T_2(v)\, T_1(u)\, R_{12}(u-v)
\]
with spectral parameter identified as the coordinate difference
of ordered insertion points: $u = z_1 - z_2$. The factorisation
structure uses \emph{ordered} configuration spaces
$\Conf^{\mathrm{ord}}_n(X)$, and Verdier duality acts by
$R$-matrix inversion $R(u) \mapsto R(u)^{-1}$. This definition is
the Vol~II/I source-of-truth (see Vol~I
\cite{VolI} \ref{def:e1-chiral-yangian-vol3});
the present paragraph records its Vol~III incarnation as input
to $\Phi$.
\end{definition}

\begin{definition}[Chiral Yangian datum (stronger)]
\label{def:chiral-yangian-datum-vol3}
A \emph{chiral Yangian datum} on $X$ is a quadruple
$(A, S(z), \Delta^{\mathrm{ch}}, \{m_k\})$ consisting of:
\begin{enumerate}[label=(\roman*)]
 \item An $\Eone$-chiral algebra $A$ on $\Ran(X)$ (so in
 particular an $\Eone$-chiral Yangian by
 Definition~\ref{def:e1-chiral-yangian-vol3});
 \item A spectral $S$-matrix $S(z) \in \End(A \otimes A)((z))$
 satisfying the spectral YBE;
 \item A chiral coproduct $\Delta^{\mathrm{ch}} \colon A \to
 A \otimes A$ at depth $z = z_1 - z_2$ compatible with the
 $\Eone$-factorisation structure;
 \item Higher operations $\{m_k\}$ generating the modular
 Maurer--Cartan tower $\Theta_A^{E_1}$ in the convolution
 algebra (Volume~I, MC2--MC4).
\end{enumerate}
The four axioms are: Koszul self-recovery
$\Phi(\Phi(A)) \simeq A$ on the Koszul locus; Verdier by
$R$-inversion; existence of a braided-module category
 $\Rep^{\Eone}(A)$; and convergence of the modular MC tower at all
genera.
\end{definition}

\begin{remark}[Forgetful arrow and the open equivalence]
\label{rem:yangian-forgetful-arrow}
There is a forgetful functor from chiral Yangian data to
$\Eone$-chiral Yangians:
$(A, S, \Delta^{\mathrm{ch}}, \{m_k\}) \mapsto A^{E_1}$. Whether
this functor is equivalently surjective is the open equivalence
problem between the two Yangian definitions: given an $\Eone$-chiral Yangian, can the
chiral coproduct $\Delta^{\mathrm{ch}}$ be reconstructed from the
$\Eone$-factorisation data alone? The main obstruction is that
$\Delta^{\mathrm{ch}}$ probes the depth-two collision residue
$r(z_1 - z_2)$ at $z_1 \neq z_2$, while the $\Eone$-factorisation
structure probes only the depth-one collision residue at the
limit $z_1 \to z_2$. Whether the depth-two data is determined by
the depth-one data is unsolved.
\end{remark}

\begin{remark}[Vol III usage convention]
\label{rem:vol3-yangian-convention}
Throughout this chapter and the rest of Vol~III, the term
\emph{Yangian} without further qualifier refers to the
RTT-presented $\Eone$-chiral Yangian of
Definition~\ref{def:e1-chiral-yangian-vol3}. Its local mode algebra is
the RTT algebra of Definition~\ref{def:yangian}; the curve-level object
adds ordered factorisation on~$\Ran(X)$ and therefore is not merely the
classical algebra with notation changed. The full chiral Yangian datum
is invoked explicitly when needed
(e.g.~for the Drinfeld-double construction or when the chiral
coproduct $\Delta^{\mathrm{ch}}$ enters).
\end{remark}


\section{Four Yangian types must be kept distinct}
\label{sec:four-yangian-types}

Four distinct objects are all called ``Yangian'' in the literature;
conflating any two is a type error. Vol~II catalogues the
four-fold typology; the Vol~III incarnation is recorded here.

\begin{definition}[Four Yangian types]
\label{def:four-yangian-types}
The four distinct objects collectively called ``Yangian'' are:
\begin{enumerate}[label=(\roman*)]
 \item \emph{Classical Yangian} $Y_\hbar(\frakg)$: the
 $\C[\hbar]$-Hopf algebra on a point with
 Drinfeld--Jimbo presentation; lives on $\R$ as
 $\Eone$-topological;
 \item \emph{dg-shifted Yangian} $Y_\hbar^{\mathrm{dg}}(\frakg)$:
 the dg-Hopf algebra on a formal disk $D = \Spec\,\C[[z]]$
 obtained as the Etingof--Kazhdan quantization at the formal
 disk;
 \item \emph{Chiral Yangian} $Y(\frakg)^{\mathrm{ch}}$: the
 $\Eone$-chiral algebra on a smooth curve $X$, factorisation
 algebra on $\Ran(X)$ with ordered configuration spaces
 (Definition~\ref{def:e1-chiral-yangian-vol3});
 \item \emph{Spectral Yangian} $Y^{\mathrm{spec}}(\frakg)$: the
 factorisation algebra on the spectral $\A^1_u$ obtained by
 viewing the spectral parameter $u$ as the geometric
 coordinate.
\end{enumerate}
The four are related by functorial passages: classical
$\to$ dg-shifted via formal completion at $\hbar = 0$;
dg-shifted $\to$ chiral via globalisation along the curve;
chiral $\to$ spectral via Fourier-Mukai along the spectral
parameter. None of these passages is an equivalence in
general.
\end{definition}

\begin{remark}[Type errors and where they occur]
\label{rem:yangian-type-errors}
Conflating any two of the four types produces a type error.
Commonly documented confusions:
\begin{enumerate}[label=(\alph*)]
 \item Classical vs dg-shifted: the classical Yangian has no
 dg-structure, the dg-shifted version does. Statements about
 ``the Yangian's dg-Hopf structure'' refer to the dg-shifted
 type only;
 \item Chiral vs spectral: the chiral Yangian lives on the
 \emph{worldsheet curve} $X$; the spectral Yangian lives on
 the \emph{spectral curve} $\A^1_u$. These are different
 spaces with different geometric meaning. Statements like
 ``the chiral algebra has spectral parameter $z$'' must
 specify which curve $z$ parameterizes.
 (The elliptic Belavin caveat from Vol~I applies: the elliptic,
 trigonometric, and rational degenerations refer to the
 spectral parameter $u$, not the worldsheet $z$.)
\end{enumerate}
\end{remark}

\begin{remark}[Vol III scope: chiral Yangian only]
\label{rem:vol3-yangian-scope}
The Yangian in Vol~III is the chiral Yangian
$Y(\frakg)^{\mathrm{ch}}$ of
Definition~\ref{def:four-yangian-types}(iii) throughout. The
classical and dg-shifted Yangians appear only as inputs to the
globalisation procedure constructing
$Y(\frakg)^{\mathrm{ch}}$ on a curve. The spectral Yangian appears
in the Drinfeld-center diagonalisation
(Chapter~\ref{ch:drinfeld-center}) and the K3 toroidal Yangian
(Chapter~\ref{ch:k3-quantum-toroidal}), where the spectral
parameter $u$ identifies with a coordinate on the K3 elliptic
fibration.
\end{remark}


\section{Mukai-Heisenberg lattice sector at \texorpdfstring{$d = 2$}{d = 2}}
\label{sec:mukai-heisenberg-kl}

For K3 surfaces, the cross-volume identification
$\Phi(D^b(\Coh(K3))) = $ $\Etwo$-Mukai-Heisenberg provides a
quantum-group-side lattice criterion: the rank-$24$ Mukai lattice
$\Lambda_{\mathrm{Muk}} = II_{4,20}$ supplies the charge lattice, and the
$R$-matrix at degree $(1,1)$ is the Mukai pairing. This is a
Heisenberg/lattice sector. It is not the $K3\times E$ BKM object and
not the Mukai self-mirror K3 Yangian branch.

\begin{proposition}[Mukai-Heisenberg sector for the lattice \texorpdfstring{$\Lambda_{\mathrm{Muk}}$}{Lambda-Muk}]
\label{prop:kl-k3-lattice}
Let $S$ be a K3 surface with Mukai lattice
$\Lambda_{\mathrm{Muk}} = U^4 \oplus E_8(-1)^2 \simeq II_{4,20}$.
Then:
\begin{enumerate}[label=(\roman*)]
 \item The B-model output
 $\Phi_2(D^b(\Coh(S))) = \mathrm{MH}_2(\Lambda_{\mathrm{Muk}})$
 is the $\Etwo$-Mukai-Heisenberg algebra of
 Definition~\ref{def:e2-mukai-heisenberg};
 \item Its representation category is the Heisenberg Fock category
 graded by $\Lambda_{\mathrm{Muk}}$ and braided by the Mukai pairing;
 \item Since $\Lambda_{\mathrm{Muk}}$ is even unimodular,
 $\Lambda_{\mathrm{Muk}}^{\vee}/\Lambda_{\mathrm{Muk}}=0$. Any finite discriminant-form
 quotient is therefore the trivial pointed category. The usual
 positive-definite rational lattice-VOA modular-tensor-category theorem
 does not apply directly to the indefinite lattice $II_{4,20}$.
\end{enumerate}
\end{proposition}

\begin{proof}
Item~(i) is Vol~III property (U4) of $\Phi$ at $d = 2$
combined with the Hochschild bridge for B-side input
$D^b(\Coh(K3))$. The Heisenberg Fock modules are labelled by charges
in the Mukai lattice, and their braiding is the exponential of the
Mukai pairing. Unimodularity gives
$\Lambda_{\mathrm{Muk}}^{\vee}/\Lambda_{\mathrm{Muk}}=0$, so the discriminant-form quotient
has one simple object. The finite rational lattice-VOA theorem requires
a positive-definite even lattice; here it supplies only the trivial
discriminant check after passing to the quotient, while the full
Heisenberg sector remains a Fock category over the indefinite charge
lattice.
\end{proof}

\begin{remark}[Cross-side comparison: A-side via HMS]
\label{rem:mukai-heisenberg-a-side}
The A-side identification
$\Phi_2(\Fuk(K3)) = \mathrm{MH}_2(\Lambda_{\mathrm{Muk}})$
is Theorem~\ref{thm:fukaya-k3-mukai-heisenberg} of the
preceding Fukaya chapter. Combined with HMS at K3 (Seidel--Smith,
Sheridan; partial), the two sides give the same chiral algebra on the
constructed locus: both Fukaya and derived-coherent inputs to $\Phi_2$
at $d = 2$ produce the $\Etwo$-Mukai-Heisenberg algebra.
\end{remark}


\section{CY-C on constructed loci}
\label{sec:cy-c-constructed-loci}

The quantum-group reconstruction conjecture CY-C asserts that for
a CY$_d$ category $\cC$ with sufficient symmetry, there exists a
generalised quantum group $G(\cC)$. At $d = 2$ the predicted braided
category is $\Rep^{\Etwo}(\Phi_2(\cC))$; at $d \geq 3$ it is
$\cZ(\Rep^{\Eone}(\Phi_d(\cC)))$, the Drinfeld centre of the native
$\Eone$ representation category. CY-C is \textsc{conjectural} in
general; the cases below delimit the constructed boundary.

\begin{conjecture}[CY-C on constructed and unconstructed loci]
\label{conj:cy-c-constructed-loci}
The CY-C claims split as follows:
\begin{enumerate}[label=(\roman*)]
 \item For $d = 2$ with $\cC = D^b(\Coh(K3))$ or $\cC = \Fuk(K3)$:
 the constructed target is the Mukai-Heisenberg lattice sector
 (Proposition~\ref{prop:kl-k3-lattice}); CY-C is verified only at this
 Heisenberg/lattice level;
 \item For $d = 2$ with $\cC = D^b(\Coh(\text{abelian surface}))$:
 $\kappa_{\mathrm{cat}}(A)=\kappa_{\mathrm{ch}}^{\mathrm{Hodge}}(A)=0$
 and $\kappa_{\mathrm{ch}}^{\mathrm{Heis}}(A)=2$. The vanishing of
 $\kappa_{\mathrm{cat}}$ does not construct a trivial quantum group;
 no nontrivial CY-C target is asserted here;
 \item For $d = 2$ with general CY$_2$ $\cC$ admitting an
 $\fsl_N$-symmetry: $G(\cC)$ is conjectured to be
 $U_q(\fsl_N)$ at root of unity, but the construction is
 incomplete beyond the Bridgeland-stability case for
 $A_{N-1}$-singularities (Conjecture~\ref{conj:qg-realization});
 \item For $d = 3$ with $\cC = \Coh(\C^3)$ or
 $\cC = \Coh(\text{conifold})$: the constructed Hall algebra is the
 positive half $Y^+$ of the relevant affine Yangian
 (Proposition~\ref{prop:coha-quantum-group}); the full Yangian and its
 braided representation category require the completed Drinfeld double,
 including negative half, Cartan, Hopf pairing, and centre compatibility;
 \item For $d = 3$ with general compact CY$_3$ $\cC$:
 $G(\cC)$ is \emph{unconstructed}; CY-C remains
 \textsc{conjectural};
 \item For super-CY$_d$ inputs: the super-Yangian
 $Y(\frakg(m|n))$ identification is \textsc{conjectural} at
 the chiral level
 (Vol~III \cite{VolI},
 Definition~\ref{def:chiral-super-yangian-mn});
 \item For $d \geq 4$: $G(\cC)$ is \emph{undefined} in general,
 since any $\Etwo$-braided comparison must pass through
 $\cZ(\Rep^{\Eone}(\Phi(\cC)))$, the Drinfeld centre
 of the $\Eone$-representation category, and the resulting braided
 category may not be of quantum-group type.
\end{enumerate}
\end{conjecture}

\begin{remark}[Pentagon stratification of the six routes]
\label{rem:cy-c-pentagon}
The six routes to $G(K3 \times E)$ (Kummer, Borcherds, MO stable
envelope, McKay, factorization homology, Costello 5d) are six
distinct constructions; only one applies $\Phi$ directly. Their
convergence is the content of CY-C, not a consequence of any
functoriality. At $d = 3$ for
$K3 \times E$ the routes form a pentagon of five intertwiners with
$R_2$ (Borcherds) as the source branch, and they stratify by the
generator-rank invariant $\rho^{R_i}$ (an algebraic invariant of
the presentation of each route's target) rather than by
$\kappa_{\mathrm{ch}}$. The compact total-space value is
$\kappa_{\mathrm{ch}}^{\mathrm{Hodge}}(K3\times E)=\chi(\cO_{K3\times E})=0$,
while the additive scalar
$\kappa_{\mathrm{ch}}^{\mathrm{Heis}}(K3\times E)
=\kappa_{\mathrm{ch}}^{\mathrm{Heis}}(K3)
+\kappa_{\mathrm{ch}}^{\mathrm{Heis}}(E)=2+1=3$
belongs to the relative Heisenberg/free-field route. The
manifold invariant $\kappa_{\mathrm{cat}} = \chi(\cO_{K3 \times E})
= 2 \cdot 0 = 0$ (K\"unneth with $\chi(\cO_E) = 0$) is likewise
route-independent. The genuinely route-distinguishing invariant is the
generator-rank $\rho^{R_i}$. The BKM-side object is the
Hall--Drinfeld double of the $K3\times E$ CoHA after the positive half,
negative half, Cartan, Hopf pairing, and completion are constructed; it
is not the Mukai self-mirror K3 Yangian branch.
\end{remark}

\begin{remark}[Generators and relations obstruction]
\label{rem:why-cy-c-hard}
The structural difficulty of CY-C is that constructing $G(\cC)$
from $\cC$ requires a generators-and-relations presentation
of the braided category attached to $\Phi(\cC)$: $\Rep^{\Etwo}(\Phi_2(\cC))$
at $d=2$ and $\cZ(\Rep^{\Eone}(\Phi_d(\cC)))$ at $d\geq3$. But $\Phi(\cC)$
is constructed via factorisation envelopes that are inherently
\emph{not} presented by generators and relations. The known
verified cases (K3 lattice, toric CY$_3$) all have
generators-and-relations presentations of $\Phi(\cC)$ available
through other means (lattice-VOA structure, Hall algebra
structure). The general CY$_3$ case lacks such presentations.
For $K3 \times E$, even when generators and
relations exist (for K3 $\times$ E), the natural invariant
distinguishing routes is the generator-lattice rank, not the
modular characteristic.
\end{remark}


\section{Quantum group A vs B/C/D distinctions}
\label{sec:qg-abcd-distinctions}

The four classical series exhibit type-dependent behaviour at
several points relevant to the chiral quantum group construction.
MC3 is proved on the evaluation-generated core for all simple
types via the five-family mechanism; this section records the
type-dependent distinctions in that setting.

\begin{proposition}[MC3 five-family mechanism: types A--D and exceptional]
\label{prop:mc3-five-family-types}
The Vol~II MC3 theorem on the evaluation-generated core is
extended from the type-A Baxter
constraint to a universal five-family mechanism covering all
simple types. The five families are:
\begin{enumerate}[label=(\roman*)]
 \item \emph{Type A} (asymptotic prefundamentals): generic
 evaluation modules at distinct spectral parameters give
 prefundamental modules whose limits as
 $u_i \to \infty$ generate the evaluation core;
 \item \emph{Types B, C, D} (reflection-equation Shapovalov):
 the orthogonal/symplectic evaluation modules satisfy a
 reflection equation in addition to the YBE, and the
 Shapovalov form of these modules generates the evaluation
 core;
 \item \emph{Uniform across all types} (Chari--Moura
 multiplicity-free $\ell$-weights): the
 $\ell$-weight space decomposition of evaluation modules is
 multiplicity-free at generic spectral parameters, ensuring
 the core is generated by the simple objects;
 \item \emph{Elliptic} (theta-divisor complement): for
 elliptic deformations of the evaluation modules, the
 $\Theta$-divisor complement provides the evaluation generator
 set;
 \item \emph{Super} (parity-balance for super-Yangian): the
 super-Yangian $Y(\frakg(m|n))$ has its evaluation core
 generated by parity-balanced modules
 (super-trace-zero condition).
\end{enumerate}
The Baxter constraint $b = a - 1/2$ is specific to
the asymptotic-prefundamental type-A presentation; it does not generalise
to types B/C/D, where the reflection-equation Shapovalov gives
an alternative generator set.
\end{proposition}

\begin{remark}[Type-A vs B/C/D in the chiral quantum group]
\label{rem:qg-type-a-vs-bcd}
The chiral quantum group $\Uq(\widehat{\frakg})^{\mathrm{ch}}$
has type-dependent structure:
\begin{enumerate}[label=(\alph*)]
 \item Type $A_{N-1}$: $\Uq(\widehat{\fsl}_N)^{\mathrm{ch}}$
 admits the standard RTT presentation with the rational
 $R$-matrix
 $R(u) = (u\,\mathrm{Id} - \Psi\,P)/(u - \Psi)$
 (additive Yang). The qdet is central and gives the
 $U(\widehat{\fgl}_N)^{\mathrm{ch}}$ extension.
 (The RTT sign convention, fixed by Vol~I, is
 $[t_{ij}, t_{kl}] = \Psi(\delta_{il} t_{kj} - \delta_{kj} t_{il})$
 in the additive presentation; the opposite sign appears in Molev's
 multiplicative $1 - P/u$ presentation.)
 \item Types $B_n, C_n, D_n$: $\Uq(\widehat{\frakg})^{\mathrm{ch}}$
 requires the orthogonal/symplectic $R$-matrix and
 satisfies an additional reflection equation. The Drinfeld
 second presentation differs from type A by extra
 generator-level conditions.
 \item Exceptional $G_2, F_4, E_6, E_7, E_8$: the chiral
 quantum group exists abstractly via the universal Yangian
 of Drinfeld but lacks a closed-form RTT presentation; the
 evaluation-generated core construction
 (Proposition~\ref{prop:mc3-five-family-types}) provides
 access via Chari--Moura $\ell$-weights.
\end{enumerate}
\end{remark}

\begin{remark}[Super-Yangian scope]
\label{rem:super-yangian-scope}
The super-Yangian $Y(\frakg(m|n))^{\mathrm{ch}}$ is
\textsc{conjectural} at the chiral level:
\begin{enumerate}[label=(\alph*)]
 \item The Drinfeld--Jimbo presentation of the classical
 super-Yangian $Y_\hbar(\frakg(m|n))$ is established
 (Khoroshkin, Lukierski, Tsuboi);
 \item The super-trace-zero condition gives a generalised
 RLL relation with the super-$R$-matrix replacing the bosonic
 one;
 \item The chiral globalisation $Y(\frakg(m|n))^{\mathrm{ch}}$
 on a curve $X$ is conjectured but not constructed; the
 main obstruction is that $\Sigma_n$-coinvariants do not
 commute with super-symmetrisation in the way required for
 the factorisation envelope;
 \item The super-complementarity identity
 $\kappa_{\mathrm{ch}} + \kappa_{\mathrm{ch}}^! = \max(m, n)$ is verified per-case for
 small ranks but lacks a canonical pairing
 (super-trace vs Berezinian; Vol~II open frontier F4).
\end{enumerate}
\end{remark}


\section{Magnificent Four at \texorpdfstring{$d = 4$}{d=4}: the scalar trace of \texorpdfstring{$\PhiFA_4(\Perf(\C^4))$}{Phi-FA-4(Perf(C4))}}
\label{sec:magnificent-four-d4}

Conjecture~\ref{conj:cy-c-constructed-loci}~(vii) leaves the quantum-group
target at $d \geq 4$ unconstructed; in general the Drinfeld-centre
category of a Stage-2 $\Eone$ specialisation need not be of
quantum-group type. The flat affine slice $\cC = \Perf(\C^4)$ admits a
sharp chain-level statement about $\PhiFA_4(\cC)$ itself: the
Stage-$1$ object is an $E_4$ holomorphic factorisation algebra, and the
Nekrasov--Piazzalunga--Kool--Rennemo Magnificent Four formula is its
toric Hilbert-scheme scalar trace.

\begin{theorem}[\texorpdfstring{$E_4$}{E4}-factorisation structure on \texorpdfstring{$\PhiFA_4(\Perf(\C^4))$}{Phi-FA-4(Perf(C4))}]
\label{thm:mag-four-e4-structure-qgr}
The infinity-category $\Perf(\C^4)$ of perfect complexes on affine
Calabi--Yau $4$-space carries the canonical CY$_4$ structure given by
the holomorphic volume form
$\Omega_{\C^4} = dz_1 \wedge dz_2 \wedge dz_3 \wedge dz_4$. The
Stage-1 factorisation envelope
$\PhiFA_4(\Perf(\C^4))$ of Chapter~\ref{ch:e1-chiral} is a factorisation
algebra on $\R^8 \simeq \C^4$ (Costello--Gwilliam 2021 Vol.~II
\S 4.10) carrying a canonical $E_4$ holomorphic factorisation structure.
The PTVV form on $\mathcal M(\Perf(\C^4))$ has shift $-2$, so functions
on the moduli stack carry the non-degenerate shifted Poisson bracket of
cohomological degree $2$.
The equivariant Hilbert-scheme trace is a map
\[
 \Tr^{M4}_{T_{\mathrm{CY}}}\colon
 K^{T_{\mathrm{CY}}}\!\left(\bigsqcup_{n\geq0}\Hilb^n(\C^4)\right)
 \longrightarrow \Q(t_1,t_2,t_3,t_4)\llbracket q\rrbracket ,
\]
and is not an operadic structure on $\PhiFA_4(\Perf(\C^4))$.
\end{theorem}

\begin{proof}
Hochschild--Kostant--Rosenberg on smooth affine CY$_4$ identifies
$\HH^\bullet(\Perf(\C^4)) \simeq \mathrm{PV}^\bullet(\C^4)$ with
the Schouten--Nijenhuis calculus. Kontsevich--Tamarkin $E_4$-formality
gives a chain-level $E_4$-algebra structure, unique up to the
$\mathrm{GRT}_1(\Q)$-torsor (Willwacher 2010). Costello--Gwilliam
Vol.~II \S 4.10 packages this $E_4$-structure globally as a
factorisation algebra on $\R^8$. Pantev--To\"en--Vaqui\'e--Vezzosi give
	the shift $2-d=-2$ on the moduli of objects; the non-degenerate inverse
	is the shifted Poisson bracket of degree $2$ on functions. The Hilbert-scheme trace
is obtained only after passing to the toric DT$_4$ fixed-locus
realisation of the positive Hall shadow.
\end{proof}

\begin{definition}[Magnificent Four partition function]
\label{def:zm4-qgr}
Let $T^4 = (\C^\times)^4$ act diagonally on $\C^4$ with equivariant
parameters $\epsilon_i = \log t_i$ subject to the Calabi--Yau condition
$\sum_{i=1}^{4} \epsilon_i = 0$, cutting out the torus
$T_{\mathrm{CY}} = \{(t_1, t_2, t_3, t_4) \in T^4 : t_1 t_2 t_3 t_4 = 1\}$.
Let $p$ be the instanton variable and let $s$ be the $U(1\mid 1)$
mass--Coulomb ratio. For a monomial $x$ write
$[x]=x^{1/2}-x^{-1/2}$. The rank-one Magnificent Four partition function
of Nekrasov--Piazzalunga~\cite{NekrasovPiazzalunga2019} is the signed
$K$-theoretic DT$_4$ fixed-point sum
\[
 Z^{M4}_{1}(p;t_i,s)
 \;=\; \sum_{\pi} p^{|\pi|}\,(-1)^{\mu_\pi}\,[-\mathsf v_\pi],
\]
where $\pi$ ranges over solid partitions and
$(-1)^{\mu_\pi}[-\mathsf v_\pi]$ is the Oh--Thomas localisation weight
with the Nekrasov--Piazzalunga orientation sign. Kool--Rennemo~\cite{KoolRennemo2025}
prove the closed plethystic form
\[
 Z^{M4}_{1}(p;t_i,s)
 \;=\;
 \operatorname{PE}\!\left(
 \frac{[t_1t_2][t_1t_3][t_2t_3][s]}
 {[t_1][t_2][t_3][t_4][p\sqrt{s}][p/\sqrt{s}]}
 \right),
 \qquad t_1t_2t_3t_4=1.
\]
\end{definition}

\begin{theorem}[Generating-function match]
\label{thm:mag-four-partition-match-qgr}
Under the toric DT$_4$ trace comparison, the equivariant
$T_{\mathrm{CY}}$-$K$-theoretic trace of the positive Hilbert-scheme
shadow is the Magnificent Four partition function:
\[
 \sum_{n\geq0}p^n\,
 \Tr^{M4}_{T_{\mathrm{CY}}}\!\left([\Hilb^n(\C^4)]^{\mathrm{vir}}\right)
 \;=\; Z^{M4}_{1}(p;t_i,s)
 \in K^{T_{\mathrm{CY}}\times \C_s^\times}(\mathrm{pt})\llbracket p\rrbracket .
\]
The identity is an exact formal identity in equivariant $K$-theory.
\end{theorem}

\begin{proof}
The $T_{\mathrm{CY}}$-fixed points of $\Hilb^n(\C^4)$ are solid
partitions $\pi$ of size $n$. The DT$_4$ virtual tangent complex gives
the Oh--Thomas localisation class at $\pi$~\cite{OhThomas2023};
Nekrasov--Piazzalunga write this class as
$(-1)^{\mu_\pi}[-\mathsf v_\pi]$. Kool--Rennemo prove the
Nekrasov--Piazzalunga formula by refining the local classes to
factorisable sheaves and applying Okounkov rigidity~\cite{KoolRennemo2025}. The Hall trace of
the positive Hilbert-scheme shadow is precisely the sum of these
localised classes.
\end{proof}

\begin{example}[Fixed-point support]
\label{ex:Z-m4-expansion}
The unweighted support of the trace is the solid-partition set:
\[
 \sum_{n\geq0}|\Hilb^n(\C^4)^{T_{\mathrm{CY}}}|p^n
 = 1 + p + 4p^2 + 10p^3 + 26p^4 + 59p^5 + \cdots.
\]
The partition function itself weights this support by the DT$_4$
orientation sign and the square-root Euler class.
\end{example}

\begin{proposition}[CY-D at \texorpdfstring{$d = 4$}{d=4} on compact CY$_4$]
\label{prop:cy-d-d4-compact-qgr}
For a compact simply-connected Calabi--Yau $4$-fold $X$ with
Hodge pattern $h^{0,0}(X) = h^{0,4}(X) = 1$ and
$h^{0,1}(X) = h^{0,2}(X) = h^{0,3}(X) = 0$,
\[
 \kappa_{\mathrm{ch}}^{\mathrm{Hodge}}(\PhiFA_4(\Perf X)) \;=\; \chi(\cO_X) \;=\; 2.
\]
In particular, CY-D at $d = 4$ on simply-connected compact CY$_4$ is a
theorem; it is not conjectural in this case.
\end{proposition}

\begin{proof}
For compact $X$, the Hodge supertrace of $\PhiFA_4(\Perf X)$ equals
$\chi(\cO_X)$ by Theorem~\ref{thm:cy-modular-characteristic}.
Hirzebruch--Riemann--Roch gives
$\chi(\cO_X) = (3 c_2^2 - c_4)/720 = 2$ for simply-connected
CY$_4$.
\end{proof}

\begin{remark}[Four-\texorpdfstring{$\kappa_\bullet$}{kappa-.} values on the flat slice and the sextic]
\label{rem:four-kappa-d4-split}
On the affine slice $\cC = \Perf(\C^4)$ the four invariants split as
\[
 \kappa_{\mathrm{cat}}(\C^4) = \chi(\cO_{\C^4}) = 1, \qquad
 \kappa_{\mathrm{ch}}^{\mathrm{Hodge}}(\PhiFA_4(\Perf \C^4)) = 0,
\]
with $\kappa_{\mathrm{BKM}}, \kappa_{\mathrm{fiber}}$ not applicable
(the flat affine slice carries no BKM). The discrepancy
$1 \neq 0$ is the Berezin-integration volume factor: the categorical
Euler includes the unit section while the Hodge supertrace subtracts
the volume normalisation. On the compact sextic $X_6 \subset \mathbb{P}^5$
(simply-connected CY$_4$),
$\kappa_{\mathrm{cat}}(X_6) = \kappa_{\mathrm{ch}}^{\mathrm{Hodge}}(\PhiFA_4(\Perf X_6)) = 2$
by Proposition~\ref{prop:cy-d-d4-compact-qgr} with
$\kappa_{\mathrm{BKM}}$ open because no $d = 4$ BKM has been
constructed, and the middle Hodge number is \(h^{2,2}(X_6)=1752\). The full fourth Betti number is
\[
b_4(X_6)=h^{4,0}+h^{3,1}+h^{2,2}+h^{1,3}+h^{0,4}
=1+426+1752+426+1=2606.
\]
The primitive \((2,2)\)-summand has dimension \(1751\), after removing the ambient class \(H^2\).
\end{remark}

\begin{proposition}[$S_4$-symmetry at \texorpdfstring{$d = 4$}{d=4}]
\label{prop:s4-symmetry-d4-qgr}
$Z^{M4}_{1}(p;t_i,s)$ is invariant under permutation of
$(t_1,t_2,t_3,t_4)$ subject to $t_1t_2t_3t_4=1$. The action lifts to the
orientation groupoid of $\PhiFA_4(\Perf(\C^4))$ and preserves the scalar
trace $\Tr^{M4}_{T_{\mathrm{CY}}}$. This is the $d = 4$ analogue of the
Miki $\Z/3$-triality at $d = 3$ on
$\mathrm{CoHA}(\C^3) = Y^+(\widehat{\fgl}_1)$ (Miki 2007;
Feigin--Jimbo--Miwa--Mukhin 2016).
\end{proposition}

\begin{proof}
The displayed plethystic numerator chooses the polarisation
$\{t_1t_2,t_1t_3,t_2t_3\}$ and is not, by itself, an $S_4$-symmetric
function. A permutation of the four coordinates changes this
polarisation and also changes the DT$_4$ orientation sign. Monavari's
canonical-vertex theorem~\cite{Monavari2022}, used in the
Kool--Rennemo proof, states that
the product $(-1)^{\mu_\pi}[-\mathsf v_\pi]$ is invariant under this
simultaneous change. Since the coordinate permutation sends
$\Omega_{\C^4}$ to $\mathrm{sgn}(\sigma)\Omega_{\C^4}$, the lift acts on
the orientation line rather than on a fixed oriented chart. The local
weights and hence their Hall trace are therefore invariant.
\end{proof}

\begin{conjecture}[Maulik--Okounkov envelope at \texorpdfstring{$d = 4$}{d=4}]
\label{conj:mo-envelope-d4-qgr}
There exists a $d = 4$ analogue of the Maulik--Okounkov stable envelope
(\texttt{arXiv:1211.1287}) on the CY$_4$ Hilbert schemes
$\{\Hilb^n(\C^4)\}_{n \geq 0}$, constructed from the
$T_{\mathrm{CY}}$-equivariant $K$-theory via the Nekrasov--Piazzalunga
sign rule, and satisfying:
\begin{enumerate}[label=(\roman*)]
 \item a chamber structure indexed by orderings of
 $(\epsilon_1, \epsilon_2, \epsilon_3, \epsilon_4)$ on the CY$_4$ cone
 $\sum \epsilon_i = 0$;
 \item Yang--Baxter equation for the resulting $R$-matrix
 $R^{M4}(u, v)$ on
 $\bigoplus_n K^{T_{\mathrm{CY}}}(\Hilb^n(\C^4))$;
 \item generating-function compatibility
 $\sum_n q^n \Tr_{K^{T_{\mathrm{CY}}}(\Hilb^n)}\! R^{M4}(u, v)
 = Z^{M4}(q) \cdot f(u, v)$ for an explicit rational prefactor.
\end{enumerate}
This is the $d = 4$ incarnation of the universal positive-geometry
grammar $Y^+(X) = H^\bullet_{\mathrm{eq}}(\cM^+_{\mathrm{eff}}(X), \phi_W)$
with $X = \C^4$ and equivariance stratum (i)~toric $T^4$.
\end{conjecture}

\begin{remark}[Three Stage-2 specialisations at \texorpdfstring{$d = 4$}{d=4}]
\label{rem:spch-d4}
The two-stage factorisation
$\Phi_d = \SpCh_{\Sigma_{d-1}, C} \circ \PhiFA_d$ at $d = 4$ needs a
three-cycle $\Sigma_3$ and a curve $C$. The natural specialisations are
(T1)~torus $\Sigma_3 = T^3$ with $C = S^1$, (T2)~sphere $\Sigma_3 = S^3$
with $C$ a great circle, and (T3)~Hopf fibration
$\Sigma_3 = S^3 \to S^2$ with $C$ a Hopf fibre. The three produce
distinct chiral algebras on $C$ from a common Stage-1 source.
$\PhiFA_4$ is one functor; the Stage-2 stratum produces a
family of specialisations, not a family of $\Phi_4$-applications.
Mutual compatibility of the three specialisations is open (the
$d = 4$ analogue of the pentagon at $d = 3$).
\end{remark}

\begin{remark}[No quantum-group target at \texorpdfstring{$d = 4$}{d=4}]
\label{rem:upgrade-relative-cy-c-vii}
At $d \geq 4$ a quantum-group target $G(\cC)$ is undefined in general.
Theorem~\ref{thm:mag-four-e4-structure-qgr} does \emph{not} construct
$G(\cC)$ at $d = 4$; it constructs the Stage-1 factorisation envelope
$\PhiFA_4(\cC)$ at $\cC = \Perf(\C^4)$ as an $E_4$ holomorphic
factorisation algebra and identifies its toric scalar trace. The
quantum-group target $G(\Perf(\C^4))$ remains unconstructed. The
Drinfeld centre
$\cZ(\Rep^{\Eone}(\SpCh_{\Sigma_3,C}(\PhiFA_4(\Perf(\C^4)))))$, after
a Stage-2 specialisation $(\Sigma_3,C)$, would carry the
$\Etwo$-braided structure eligible for quantum-group identification
at the representation-category level; it remains conjectural.
\end{remark}


\section{Cross-volume \texorpdfstring{$\kappa_\bullet$}{kappa-.}-stratification}
\label{sec:cross-volume-kappa-strat}

The $\kappa_\bullet$-stratification by CY dimension (Vol~III
\cite{VolI}) restricts to the
quantum-group inputs of this chapter as follows. The chiral
Hodge-supertrace invariant $\kappa_{\mathrm{ch}}^{\mathrm{Hodge}}$
is K\"unneth-multiplicative on compact products and coincides with
$\chi(\cO_X)$ only for $d = 2$ with $h^{1,0}(X) = 0$; the
Heisenberg specialisation $\kappa_{\mathrm{ch}}^{\mathrm{Heis}}$ is
a separate generator-rank scalar. At odd $d$ one has
$\chi(\cO_X)=0$ by Serre parity.

\begin{proposition}[\texorpdfstring{$\kappa_{\mathrm{ch}}$}{kappa-ch} on QG inputs]
\label{prop:kappa-strat-qg-inputs}
For the standard quantum-group inputs to $\Phi$:
\begin{enumerate}[label=(\roman*)]
 \item $d = 1$, elliptic curve $E$: the compact Hodge supertrace is
 $\kappa_{\mathrm{ch}}^{\mathrm{Hodge}}(\Phi(E))=\chi(\cO_E)=0$,
 while the Heisenberg specialisation has
 $\kappa_{\mathrm{ch}}^{\mathrm{Heis}}(E)=1$; the QG output is a rank-$2$
 lattice VOA on the Heisenberg route;
 \item $d = 2$, K3: $\kappa_{\mathrm{ch}} = \chi(\cO_{K3}) = 2$
 (since $h^{1,0} = 0$); the QG output is the Mukai-lattice
 Heisenberg sector of Proposition~\ref{prop:kl-k3-lattice};
 \item $d = 2$, abelian surface $A$:
 $\kappa_{\mathrm{ch}}^{\mathrm{Hodge}}(\Phi(A))=\chi(\cO_A)=0$,
 while $\kappa_{\mathrm{ch}}^{\mathrm{Heis}}(A)=2$ on the
 Heisenberg route; the $d = 2$ identification
 $\kappa_{\mathrm{ch}}=\chi(\cO_X)$ requires $h^{1,0}=0$ and
 therefore does not apply to abelian surfaces;
 \item $d = 3$, $\C^3$: the Schiffmann--Vasserot cohomological
 Hall algebra $\mathrm{CoHA}(\C^3)$ is an associative algebra,
 isomorphic to the positive half $Y^+(\widehat{\fgl}_1)$ of the
 affine Yangian; $\mathrm{CoHA}$ is not itself a chiral algebra,
 and an $\Eone$-chiral algebra $\Phi(\C^3)$ whose Drinfeld centre
 satisfies
 $\cZ(\Rep^{\Eone}(\Phi(\C^3))) \simeq \Rep_q(\widehat{\fgl}_1)$
 would give $\kappa_{\mathrm{ch}} = c = 1$ for a Heisenberg
 realization, conditional on CY-C at noncompact CY$_3$;
 \item $d = 3$, conifold: the resolved conifold is not a local
 surface; direct McKay computation gives
 $\kappa_{\mathrm{ch}}(\cA_{\mathrm{conifold}})=1$. Its CoHA is a
 positive Hall half in the toric chamber; the full
 $Y(\widehat{\fsl}_2)$ requires the completed Drinfeld double;
 \item $d = 3$, compact quintic: $\chi(\cO_{\text{quintic}}) = 0$
 by Serre parity, so the $d = 2$ identification
 $\kappa_{\mathrm{ch}} = \chi(\cO_X)$ cannot survive; the QG
 output is conjectural (CY-C open at general compact CY$_3$);
 \item $d = 3$, $K3 \times E$:
 the compact Hodge supertrace is
 $\kappa_{\mathrm{ch}}^{\mathrm{Hodge}}(K3 \times E)=0$, while the
 Heisenberg/Stage-$2$ specialisation has
 $\kappa_{\mathrm{ch}}^{\mathrm{Heis}}(K3 \times E)
 = \kappa_{\mathrm{ch}}^{\mathrm{Heis}}(K3)
 + \kappa_{\mathrm{ch}}^{\mathrm{Heis}}(E) = 2 + 1 = 3$;
 the Hall--Drinfeld/Borcherds branch attached to
 $\mathfrak g_{\Delta_5}$ has
 $\kappa_{\mathrm{BKM}}(\Delta_5) = \left.c_N(0)/2\right|_{N=1}=5$ by
 Gritsenko's denominator theorem
 (Theorem~\ref{thm:borcherds-weight-kappa-BKM-universal}).
\end{enumerate}
\end{proposition}

\begin{remark}[\texorpdfstring{$\kappa_\bullet$}{kappa-.}-values for \texorpdfstring{$K3 \times E$}{K3 x E}]
\label{rem:three-kappa-k3xe}
The four invariants separate cleanly. The manifold invariant is
$\kappa_{\mathrm{cat}}(K3 \times E) = \chi(\cO_{K3 \times E})
= \chi(\cO_{K3}) \cdot \chi(\cO_E) = 2 \cdot 0 = 0$ by K\"unneth,
vanishing because $\chi(\cO_E) = 0$ at odd $d = 1$. The
algebraization invariants come from two different constructions on
the same manifold: the Heisenberg-level chiral algebra gives
$\kappa_{\mathrm{ch}}^{\mathrm{Heis}} = 3$ by additivity across the
product, while the Hall--Drinfeld/Borcherds branch gives
$\kappa_{\mathrm{BKM}}(\Delta_5) = \left.c_N(0)/2\right|_{N=1} = 10/2 = 5$ from
the primitive Gritsenko form $\Delta_5$ (whose primitive Igusa square
is $\Phi_{10}^{\mathrm{un}}=\Delta_5^2$), by
Theorem~\ref{thm:borcherds-weight-kappa-BKM-universal}. The fibre invariant is the
Mukai-lattice rank $\kappa_{\mathrm{fiber}}=24$. The BKM-side object is
the Hall--Drinfeld double of $\mathrm{CoHA}(K3\times E)$, not the K3
Yangian; the K3 Yangian is the separate Mukai self-mirror branch. The
tempting additive decomposition
$\kappa_{\mathrm{BKM}} \stackrel{?}{=} \kappa_{\mathrm{ch}}^{\mathrm{Hodge}} + \chi(\cO_{\text{fiber}})
= 3 + 2 = 5$ is convention-mixing: the $3$ is the Heisenberg
specialisation, not the compact total-space supertrace
$\kappa_{\mathrm{ch}}^{\mathrm{Hodge}}(K3 \times E)=0$. Under the compact convention
the additive formula already fails at $N=1$, since
$5 \neq 0 + \chi(\cO_E) = 0$; across the CHL $\Z/N\Z$ orbifolds it
continues to fail because the Borcherds constant changes to
$c_N(0)/2$. The universal formula is
$\kappa_{\mathrm{BKM}}(\Phi_N)=c_N(0)/2$. The invariants are
controlled by four different functorialities: compact Hodge supertrace,
manifold Euler characteristic, BKM denominator, and fibre lattice. The arithmetic
identity $5 = 3 + 2$ at $N = 1$ is not a structural decomposition.
\end{remark}


\section{Chiral QG equivalence on the Koszul locus}
\label{sec:outlook-chiral-qg-koszul}

Which $\Eone$-chiral algebras admit a bar--cobar calculus fine
enough to compare the three quantum-group structures below? The
answer is a locus cut out by the existence of an ordered chiral
Koszul datum.

\begin{definition}[Ordered Koszul locus;
\ClaimStatusDefinitional]
\label{def:kosz-ord-vol3}
Let $X$ be a smooth algebraic curve. An \emph{ordered chiral
Koszul datum} for an augmented $\Eone$-chiral algebra $A$ on
$\Ran(X)$ (factorisation structure over the ordered configuration
spaces $\Conf^{\mathrm{ord}}_n(X)$) is a triple
$(C, \tau, F_\bullet)$:
\begin{enumerate}[label=(\roman*)]
 \item a complete conilpotent coalgebra $C$ in the ordered
 factorisation category on $\Ran(X)$;
 \item a twisting morphism $\tau \colon C \to A$ satisfying the
 Maurer--Cartan equation in the ordered convolution algebra;
 \item an exhaustive, separated filtration $F_\bullet$ on $C$
 compatible with $\tau$, each of whose conformal-weight/bar-length
 pieces is finite-dimensional over the ground field or locally
 free of finite rank over $\cO_X$, with transition maps
 satisfying the Mittag--Leffler condition,
\end{enumerate}
such that the coalgebra unit $C \to B^{\mathrm{ord}}_X(A)$ into
the ordered bar coalgebra
$B^{\mathrm{ord}}_X(A) = T^c(s^{-1}\bar{A})$ and the cobar counit
$\Omega^{\mathrm{ord}}_X(C) \to A$ are filtered
quasi-isomorphisms. The \emph{ordered Koszul locus}
\[
\mathrm{Kosz}^{\mathrm{ord}}(X)
\;\subset\;
\Eone\text{-}\mathrm{ChirAlg}(X)
\]
is the full subcategory of augmented $\Eone$-chiral algebras on
$\Ran(X)$ admitting an ordered chiral Koszul datum. Membership is
the existence of such a datum; the comparison maps in the
statements below are attached to a chosen datum. This is the
ordered ($\Eone$) counterpart of the symmetric chiral Koszul
pairs of Vol~I \cite{VolI}, restated here self-containedly.
\end{definition}

\begin{theorem}[Chiral QG determinations on coproduct-equipped
input; \ClaimStatusProvedElsewhere]
\label{thm:chiral-qg-equiv-datum-vol3}
Let $A \in \mathrm{Kosz}^{\mathrm{ord}}(X)$ underlie a chiral
Yangian datum $(A, S(z), \Delta^{\mathrm{ch}}, \{m_k\})$ of
Definition~\ref{def:chiral-yangian-datum-vol3}, so that the
coproduct is supplied as part of the datum. Then, on the
evaluation-generated core and for all simple types, the
$R$-matrix $r(z) \in \End(A \otimes A)((z))$ and the chain-level
$\Ainf$ operations $\{m_k\}_{k \geq 1}$ determine each other up
to chiral quasi-isomorphism, and each is compatible with the
supplied coproduct: this is the Vol~I content of MC3, proved via
the five-family mechanism \cite{VolI}. For Yangian input
$Y(\frakg)^{\mathrm{ch}}$, the compatibility of the pair
$(r(z), \{m_k\})$ with the supplied coproduct follows from
Drinfeld's second presentation together with Etingof--Kazhdan
flatness \cite{VolI}. The scope is exactly this: the coproduct is
input, never output.
\end{theorem}

\begin{conjecture}[Chiral QG equivalence on the ordered Koszul
locus; \ClaimStatusConjectured]
\label{thm:chiral-qg-equiv-vol3}
On the ordered Koszul locus $\mathrm{Kosz}^{\mathrm{ord}}(X)$ of
Definition~\ref{def:kosz-ord-vol3}, the three structures
\begin{enumerate}[label=(\arabic*)]
 \item $R$-matrix: $r(z) \in \End(A \otimes A)((z))$;
 \item $\Ainf$ structure: chain-level operations
 $\{m_k\}_{k \geq 1}$ on $A$;
 \item Ordered coproduct: $\Delta^{\mathrm{ord}} \colon A \to
 A \otimes A$ at depth $z = z_1 - z_2$
\end{enumerate}
mutually determine each other up to chiral quasi-isomorphism.
The determinations within the pair (1)--(2), and their
compatibility with a coproduct supplied as input, are
Theorem~\ref{thm:chiral-qg-equiv-datum-vol3}. The open content is
the reconstruction of (3) from (1)--(2): the ordered coproduct
probes the depth-two collision residue $r(z_1 - z_2)$ at
$z_1 \neq z_2$, while the $\Eone$-factorisation data underlying
(1)--(2) probe only the depth-one residue in the limit
$z_1 \to z_2$. Whether the depth-two datum is determined by the
depth-one datum is precisely the open forgetful-arrow problem of
Remark~\ref{rem:yangian-forgetful-arrow}; no case of this
reconstruction is proved.
\end{conjecture}

\begin{remark}[Concrete instances in Vol III]
\label{rem:vol3-concrete-qg-instances}
The chiral QG equivalence has three concrete Vol~III instances:
\begin{enumerate}[label=(\alph*)]
 \item \emph{$\fsl_2$ Yangian}: explicit $\alpha, \beta, \gamma$
 order-$\hbar$ computation including the triangle coherence
 $\alpha \circ \gamma \circ \beta = \mathrm{id}$ (Vol~I
 \cite{VolI});
 \item \emph{Affine $V_k(\fsl_2)$}: KL equivalence at
 $q = e^{\pi i / (k+2)}$ (Theorem~\ref{thm:kl-equivalence})
 gives the chiral QG datum at integer level $k$;
 \item \emph{Toroidal Yangian on $K3 \times E$}: the elliptic
 deformation, conditional on chain-level class-M topologization
 (Vol~I open frontier F3) and the K3 elliptic fibration
 structure of Chapter~\ref{ch:k3-quantum-toroidal}.
\end{enumerate}
\end{remark}

\begin{remark}[Vol I comparison theorems]
\label{rem:vol3-comparison-theorems}
The chapter cross-refers to:
\begin{enumerate}[label=(\alph*)]
 \item Conjecture~\ref{thm:chiral-qg-equiv-vol3} (chiral QG
 equivalence on the ordered Koszul locus,
 Definition~\ref{def:kosz-ord-vol3}), whose
 coproduct-equipped sub-case is
 Theorem~\ref{thm:chiral-qg-equiv-datum-vol3};
 \item Vol~II \cite{VolII} (formal deformation
 families identification for simple $\frakg$, and
 chain-level identification at fixed KZ associator for
 simple / reductive / solvable $\frakg$; covers $\fgl_N$,
 semisimple, nilpotent Heisenberg as subcases);
 \item Vol~I \cite{VolI}
 (gl$_N$ chiral QG via Feigin--Frenkel screening, all
 $N \geq 2$ unconditional);
 \item Vol~II \cite{VolII}
 (unified chiral quantum group $Q_g^{k,f,\mu}$ covering
 nine specialisation fibres);
\item the universal trace identity
$K(A) = -c_{\mathrm{ghost}}(\mathrm{BRST}(A))$ on the Vol~I
 side bridging to $\kappa_{\mathrm{BKM}}(\Phi_N) = c_N(0)/2$
 on the Vol~III side via $\Phi$ for K3-fibered Class A.
\end{enumerate}
The Vol~III incarnation of the chiral QG equivalence is the
chapter's own
Conjecture~\ref{conj:cy-c-constructed-loci} restricted to verified
cases (i)--(iv). Cases (v)--(vii) remain conditional on their comparison maps.
\end{remark}

\section{\texorpdfstring{$K3 \times E$}{K3 x E} Hall--Drinfeld double representation theory}
\label{sec:qgreps-supplement}

\begin{conjecture}[\texorpdfstring{$K3 \times E$}{K3 x E} Hall--Drinfeld double representation theory]
\label{rem:qgreps-hall-drinfeld-reps}
Assume the oriented critical CoHA positive half
$Y^+_{\mathrm{Hall}}(K3\times E)$ is constructed with its negative
half, Cartan part, nondegenerate Hall--Serre Hopf pairing, completion,
and compatibility with the Drinfeld-centre half-braiding of
$\mathbf{H}_{\Delta_5}$. Then the completed Hall--Drinfeld double
\[
 \mathcal{D}_\hbar\bigl(Y^+_{\mathrm{Hall}}(K3\times E)\bigr)
\]
admits evaluation modules $V_u$ on which the Siegel-elliptic dynamical
$R$-matrix $R_{\mathrm{Sieg,dyn}}(u-v;\tau,\rho,z)$ acts as the
half-braiding $\sigma_{V_u}(V_v)$ in
$\mathcal{Z}(\mathrm{Rep}^{\Eone}(\mathbf{H}_{\Delta_5}))$. At the
abelian Lie/Hopf level one sees the $24$ Mukai-lattice Heisenberg
directions. The non-abelian BKM structure belongs to the vertex closure
of the Borcherds denominator package, not to the Mukai self-mirror K3
Yangian branch.
\end{conjecture}


% ==========================================
% PART VI: Connections
% ==========================================

\part{Connections and Frontier}
\label{part:connections}

\chapter{The Bar-Cobar Bridge to Volume I}
\label{ch:bar-cobar-bridge}

\section{CY bar complex vs chiral bar complex}
\label{sec:cy-bar-vs-chiral-bar}

The cyclic bar complex $\mathrm{CC}_\bullet(\cC)$ of a CY category and the chiral bar complex $B(A)$ of the associated chiral algebra are related by a canonical quasi-isomorphism (Theorem CY-A(ii)).  This section spells out the dictionary.

\section{Shadow tower of a CY category}
\label{sec:cy-shadow-tower-bridge}

\section{Koszul duality: CY vs chiral}
\label{sec:koszul-cy-chiral}

\section{The five theorems in the CY setting}
\label{sec:five-theorems-cy}

\chapter{Modular Koszul Duality and CY Geometry}
\label{ch:modular-koszul-bridge}

The bar-cobar bridge of Chapter~\ref{ch:bar-cobar-bridge} translates between the CY cyclic bar complex $\mathrm{CC}_\bullet(\cC)$ and the chiral bar complex $B(A_\cC)$.  The modular Koszul duality of Volume~I operates on $B(A)$: the Verdier intertwining $D_{\Ran}(B(A)) \simeq B(A^!)$, the complementarity theorem (Theorem~C), and the modular characteristic $\kappa(A)$ controlling the genus expansion (Theorem~D).  This chapter constructs the CY-geometric incarnation of these structures: the modular convolution algebra for CY categories, the CY complementarity pairing, and the shadow CohFT derived from the CY trace.

\section{The modular convolution algebra for CY categories}
\label{sec:modular-conv-cy}

\section{CY complementarity}
\label{sec:cy-complementarity-bridge}

\section{The CY shadow CohFT}
\label{sec:cy-shadow-cohft}

\chapter{Geometric Langlands and CY Quantum Groups}
\label{ch:geometric-langlands}

\section{The geometric Langlands programme}
\label{sec:geom-langlands}

\section{CY categories in the Langlands correspondence}
\label{sec:cy-langlands}

\section{Quantum geometric Langlands and $E_2$-chiral algebras}
\label{sec:quantum-geom-langlands}

\section{The Feigin--Frenkel center and the Drinfeld center}
\label{sec:ff-drinfeld-centers}


% ==========================================
% Appendices
% ==========================================

\appendix

\chapter{Conventions}
\label{app:conventions}

\section{Grading conventions}
\label{sec:grading-conventions}

All grading is \textbf{cohomological}: differentials have degree $+1$.  The bar construction uses desuspension (matching Volumes~I--II).

\section{Sign conventions}
\label{sec:sign-conventions}

Signs follow the Koszul rule throughout.  For the curved $\Ainf$-structure: $\mu_1^2(a) = [\mu_0, a]$ (commutator, minus sign).

\section{Categorical conventions}
\label{sec:categorical-conventions}

\begin{itemize}
  \item $\DGCat$ denotes the $\infty$-category of small dg categories over $k$.
  \item $\CY_d\text{-}\Cat$ denotes the $\infty$-category of smooth $d$-dimensional CY categories.
  \item Tensor products of dg categories are derived unless stated otherwise.
  \item ``Equivalence'' means quasi-equivalence of dg categories (equivalence of associated $\infty$-categories).
\end{itemize}

\section{Operadic conventions}
\label{sec:operadic-conventions}

\begin{itemize}
  \item $E_n$ denotes the little $n$-disks operad (Boardman--Vogt).
  \item $\FM_k(\mathbb{C})$ denotes the Fulton--MacPherson compactification of $\mathrm{Conf}_k(\mathbb{C})$.
  \item $\SC^{\mathrm{ch,top}}$ denotes the Swiss-cheese operad from Volume~II.
\end{itemize}


% ------------------------------------------
% Back matter
% ------------------------------------------

\backmatter
\bibliographystyle{alpha}
% \bibliography{references}

\end{document}
